\documentclass[11pt,letterpaper,oneside]{report}

% Motor: lualatex (Unicode nativo, sin necesidad de inputenc/fontenc).
\usepackage{fontspec}
\usepackage[spanish,es-tabla]{babel}
\usepackage[margin=2.5cm, headheight=15pt]{geometry}
\usepackage{listings}
\usepackage{xcolor}
\usepackage{hyperref}
\usepackage{titlesec}
\usepackage{amsmath}
\usepackage{amssymb}
\usepackage{amsfonts}
% Corchetes dobles \llbracket ⟦u⟧ \rrbracket: notación del salto de
% desplazamientos en las specs de discontinuidad embebida y material cohesivo
% (ADR 0010). No están en amssymb; sin stmaryrd la compilación aborta con
% "Undefined control sequence".
\usepackage{stmaryrd}
\usepackage{booktabs}
\usepackage{tabularx}
\usepackage{array}
\usepackage{ragged2e}
\usepackage{fancyhdr}
\usepackage{microtype}

\newcolumntype{Y}{>{\RaggedRight\arraybackslash\footnotesize}X}

\definecolor{yamlkey}{RGB}{0, 102, 204}
\definecolor{yamlval}{RGB}{204, 0, 0}
\definecolor{yamlcomment}{RGB}{120, 120, 120}
\definecolor{bglight}{RGB}{245, 245, 245}
\definecolor{darkgray}{RGB}{50, 50, 50}

\hypersetup{
    colorlinks=true,
    linkcolor=yamlkey,
    urlcolor=yamlkey,
    pdftitle={Manual de Referencia - Solidum FEM},
    pdfauthor={Jaime Retama Velasco}
}

\titleformat{\chapter}[display]
  {\normalfont\huge\bfseries\color{yamlkey}}{\chaptertitlename\ \thechapter}{20pt}{\Huge}
\titleformat{\section}{\Large\bfseries\color{darkgray}}{\thesection}{1em}{}
\titleformat{\subsection}{\large\bfseries\color{darkgray}}{\thesubsection}{1em}{}
\titleformat{\subsubsection}{\normalsize\bfseries\color{darkgray}}{\thesubsubsection}{1em}{}

\pagestyle{fancy}
\fancyhf{}
\fancyhead[L]{\textbf{\color{darkgray}Solidum FEM --- Referencia}}
\fancyhead[R]{\color{darkgray}\leftmark}
\fancyfoot[C]{\thepage}
\renewcommand{\headrulewidth}{0.4pt}
\renewcommand{\footrulewidth}{0.4pt}

\lstdefinelanguage{yaml}{
  keywords={nodes, materials, elements, mesh, name, kind, status, interface, parameters, signature, conventions, validity, out_of_scope, acceptance, references, material_contract, dof_names, n_nodes, strain_dim, n_integration_points, type, required, desc, sign, voigt, node_orientation, configuration, expected_behaviour, strain_kind, primary_state_var, state_passthrough, numerical_caveats, setup, expect, tol_rel, tol_abs},
  keywordstyle=\color{yamlkey}\bfseries,
  ndkeywords={true, false, null},
  ndkeywordstyle=\color{yamlval}\bfseries,
  identifierstyle=\color{black},
  sensitive=false,
  comment=[l]{\#},
  commentstyle=\color{yamlcomment}\ttfamily,
  stringstyle=\color{yamlval}\ttfamily,
  morestring=[b]',
  morestring=[b]"
}

\lstset{
  basicstyle=\ttfamily\footnotesize,
  backgroundcolor=\color{bglight},
  frame=single,
  rulecolor=\color{lightgray},
  breaklines=true,
  showstringspaces=false,
  extendedchars=true,
  captionpos=b,
  literate={á}{{\'a}}1 {é}{{\'e}}1 {í}{{\'i}}1 {ó}{{\'o}}1 {ú}{{\'u}}1
           {Á}{{\'A}}1 {É}{{\'E}}1 {Í}{{\'I}}1 {Ó}{{\'O}}1 {Ú}{{\'U}}1
           {ñ}{{\~n}}1 {Ñ}{{\~N}}1 {ü}{{\"u}}1 {Ü}{{\"U}}1
           {§}{{\S{}}}1
           {ε}{{$\varepsilon$}}1 {σ}{{$\sigma$}}1 {α}{{$\alpha$}}1
           {β}{{$\beta$}}1 {γ}{{$\gamma$}}1 {δ}{{$\delta$}}1
           {θ}{{$\theta$}}1 {κ}{{$\kappa$}}1 {λ}{{$\lambda$}}1
           {μ}{{$\mu$}}1 {ν}{{$\nu$}}1 {π}{{$\pi$}}1
           {ρ}{{$\rho$}}1 {τ}{{$\tau$}}1 {φ}{{$\varphi$}}1
           {ω}{{$\omega$}}1 {Δ}{{$\Delta$}}1 {Ω}{{$\Omega$}}1
           {Φ}{{$\Phi$}}1
           {⇔}{{$\Leftrightarrow$}}1 {⇒}{{$\Rightarrow$}}1 {⇐}{{$\Leftarrow$}}1
           {→}{{$\to$}}1 {←}{{$\leftarrow$}}1 {↔}{{$\leftrightarrow$}}1
           {≤}{{$\le$}}1 {≥}{{$\ge$}}1 {≠}{{$\neq$}}1
           {≲}{{$\lesssim$}}1 {≳}{{$\gtrsim$}}1 {≈}{{$\approx$}}1
           {·}{{$\cdot$}}1 {×}{{$\times$}}1 {±}{{$\pm$}}1
           {°}{{$^{\circ}$}}1 {²}{{$^{2}$}}1 {³}{{$^{3}$}}1
           {₀}{{$_{0}$}}1 {₁}{{$_{1}$}}1 {₂}{{$_{2}$}}1
           {ᵀ}{{$^{\top}$}}1 {⁺}{{$^{+}$}}1 {⁻}{{$^{-}$}}1
           {⟨}{{$\langle$}}1 {⟩}{{$\rangle$}}1 {‖}{{$\|$}}1
           {—}{{---}}1 {–}{{--}}1 {−}{{-}}1
           {“}{{``}}1 {”}{{''}}1
           {ʹ}{{'}}1 {′}{{$'$}}1
           {…}{{\ldots{}}}1
           {ϕ}{{$\phi$}}1 {Σ}{{$\Sigma$}}1 {Λ}{{$\Lambda$}}1
           {η}{{$\eta$}}1 {ζ}{{$\zeta$}}1 {ξ}{{$\xi$}}1 {Ξ}{{$\Xi$}}1
           {ψ}{{$\psi$}}1 {Ψ}{{$\Psi$}}1 {Θ}{{$\Theta$}}1 {χ}{{$\chi$}}1
           {ℝ}{{$\mathbb{R}$}}1
           {∈}{{$\in$}}1 {∂}{{$\partial$}}1 {∇}{{$\nabla$}}1
           {∞}{{$\infty$}}1 {∫}{{$\int$}}1 {√}{{$\surd$}}1
           {₃}{{$_{3}$}}1 {ₙ}{{$_{n}$}}1 {ₜ}{{$_{t}$}}1
}

\begin{document}

\begin{titlepage}
    \centering
    \vspace*{2cm}
    {\Huge\bfseries\color{yamlkey} Solidum FEM \par}
    \vspace{1cm}
    {\LARGE Manual de Referencia \par}
    \vspace{0.5cm}
    \rule{\linewidth}{0.5mm} \par
    \vspace{2cm}
    {\large Especificaciones físicas, formulaciones numéricas y contratos\\
            de los componentes del programa\par}
    \vspace{2cm}
    {\large Generado automáticamente desde \texttt{docs/specs/}\par}
    \vspace{1cm}
    {\Large \textbf{Autor:} Jaime Retama Velasco \par}
    \vspace{0.5cm}
    {\large Facultad de Estudios Superiores Aragón \\ Universidad Nacional Autónoma de México \par}
    \vfill
    {\large \today \par}
\end{titlepage}

\pagenumbering{roman}
\setcounter{page}{1}
\tableofcontents
\newpage

\chapter*{Sobre este Manual}
\addcontentsline{toc}{chapter}{Sobre este Manual}
\noindent Este manual contiene la \textbf{referencia formal} de los componentes de Solidum FEM: cada elemento finito y cada modelo constitutivo se documenta con su especificación física (ecuaciones), su formulación numérica (matrices $\mathbf B$, rigidez tangente, integración) y su contrato YAML.

El contenido se genera automáticamente desde los archivos \texttt{docs/specs/*.md} del repositorio, que constituyen la \emph{fuente única de verdad} de cada componente. No editar este PDF manualmente; cualquier corrección debe hacerse sobre la spec correspondiente y regenerar el manual mediante:
\begin{center}\texttt{python manuals/build\_reference\_manual.py}\end{center}

Para una guía orientada al uso del programa (sintaxis YAML, ejemplos, post-procesamiento), consulte el \textbf{Manual de Usuario} en \texttt{manuals/User\_manual.pdf}.

\newpage
\pagenumbering{arabic}
\setcounter{page}{1}


\chapter{Elementos 1D — Armaduras}

\section{Truss2D}



\textit{Orden de trabajo. El usuario escribe la \textbf{especificación física}, la \textbf{formulación numérica} y el \textbf{contrato}. La IA rellena \textbf{implementación} y responde en \textbf{diálogo} durante el trabajo.}



\section{Especificación física}

\subsection{0. Descripción general}
Elemento sólido 1D de primer orden, inmerso en el plano. Dos nodos articulados; transmite únicamente fuerza axial. Aproximación $C^0$ del desplazamiento.

\subsection{1. Cinemática de desplazamientos}
Interpolación lineal del desplazamiento axial a lo largo del eje $s \in [0, L]$:
\[
u(s) = a_0 + a_1 s
\]

\subsection{2. Cinemática de deformaciones}
Única componente (axial, en el sistema local de la barra):
\[
\varepsilon_{ss}(s) = \frac{du}{ds}
\]

\subsection{3. Ecuación constitutiva}
\[
\sigma_{ss} = E\, \varepsilon_{ss}
\]

\subsection{4. Equilibrio --- forma fuerte}
\[
-\frac{d}{ds}\!\left(EA\,\frac{du}{ds}\right) = b(s), \quad s \in (0, L),
\]
con condiciones de contorno Dirichlet o Neumann ($\sigma A = \bar F$) en los extremos.

\subsection{5. Equilibrio --- forma débil}
\[
\int_0^L \frac{d\,\delta u}{ds}\, EA\, \frac{du}{ds}\, ds \;=\; \int_0^L \delta u\, b\, ds \;+\; \bigl[\delta u\, \bar F\bigr]_{\partial \Omega}, \quad \forall\, \delta u \in V_0.
\]


\section{Formulación numérica (FEM)}

\subsection{6. Discretización}
Dos nodos en $s = 0$ y $s = L$. Cosenos directores del eje en globales:
\[
c = \frac{x_2 - x_1}{L}, \qquad s_\theta = \frac{y_2 - y_1}{L}.
\]

\subsection{7. Funciones de forma}
\[
N_1(s) = 1 - \tfrac{s}{L}, \qquad N_2(s) = \tfrac{s}{L}.
\]

\subsection{8. Matriz deformación-desplazamiento}
Local (1 DOF axial por nodo): $\mathbf B_{\text{loc}} = \tfrac{1}{L}[-1,\; 1]$.
Global ($\mathbf u_e = [u_x^{(1)}, u_y^{(1)}, u_x^{(2)}, u_y^{(2)}]^\top$):
\[
\mathbf B = \tfrac{1}{L}\,[\,-c,\; -s_\theta,\; c,\; s_\theta\,], \qquad \varepsilon = \mathbf B\, \mathbf u_e.
\]

\subsection{9. Rigidez elemental}
Integración analítica exacta ($EA$ constante):
\[
\mathbf K_e = \int_0^L \mathbf B^\top EA\, \mathbf B\, ds = \frac{EA}{L}\, \mathbf d^\top \mathbf d, \qquad \mathbf d = [-c,\; -s_\theta,\; c,\; s_\theta].
\]

\subsection{10. Fuerzas internas}
\[
\mathbf F_{\text{int}} = \sigma A\, \mathbf d.
\]

\subsection{11. Cuadratura}
No aplica (forma cerrada).


\section{Contrato de implementación}

\begin{lstlisting}[language=yaml]
name: Truss2D
kind: element
status: validated        # draft → implemented → validated

interface:
  dof_names: [ux, uy]
  n_nodes: 2
  strain_dim: 1
  n_integration_points: 1

parameters:
  - { name: A, type: float, required: true, desc: Área de la sección transversal }

material_contract:
  signature: "compute_state(ε, state) -> (σ, E_tangent, state')"
  strain_kind: axial escalar

conventions:
  sign: "σ > 0 ⇔ ε > 0 (elongación)"
  voigt: "N/A (escalar)"
  node_orientation: "libre; K_e y F_int invariantes bajo permutación"

validity:
  - "|ε| ≲ 1e-2"
  - pequeños desplazamientos y rotaciones
  - cargas exclusivamente axiales

out_of_scope:
  - flexión, cortante, momentos en extremos
  - grandes desplazamientos
  - pandeo

acceptance:
  - name: barra axial bajo carga puntual
    setup: "empotrada en s=0, carga F axial en s=L"
    expect: "u(L) = F·L / (E·A)"
    tol_rel: 1.0e-10

references:
  - "Bathe K.-J., Finite Element Procedures, §4.2.1"

\end{lstlisting}


\section{Implementación}

\begin{itemize}
  \item \textbf{Archivo}: solidum/elements/truss.py
  \item \textbf{Clase}: \texttt{Truss2D} (registrada vía \texttt{@ElementRegistry.register})
  \item \textbf{Tests}:
  \item tests/test\_truss.py \ensuremath{\cdot} \texttt{TestTruss2D} --- verifica \texttt{L0}, cosenos directores, entradas de $\mathbf K_e$ (incluyendo simetría), y evaluación de $\mathbf F_{\text{int}}$ sobre desplazamientos conocidos. Los valores numéricos esperados coinciden con $\mathbf K_e = (EA/L)\,\mathbf d\mathbf d^\top$ y $\mathbf F_{\text{int}} = N\mathbf d$ para geometría de prueba $L=5$, $c=0.6$, $s=0.8$.
  \item tests/test\_integration.py \ensuremath{\cdot} \texttt{TestSolversIntegration} --- resuelve end-to-end el caso de aceptación (barra empotrada en un extremo, carga axial en el otro) con \texttt{NonlinearSolver} y \texttt{ArcLengthSolver}; en régimen elástico precedente a la fluencia reproduce $u(L)=FL/(EA)$.
  \item \textbf{Notas de traducción}:
  \item \texttt{L0}, \texttt{c}, \texttt{s} se calculan una vez en \texttt{\_\_init\_\_} sobre la configuración inicial y no se actualizan --- coherente con el régimen infinitesimal declarado.
  \item El contrato con el material es el estándar del proyecto: \texttt{material.compute\_state(\ensuremath{\varepsilon}, state) \ensuremath{\to} (\ensuremath{\sigma}, E\ensuremath{_{t}}, state')}.
  \item Para grandes desplazamientos/rotaciones usar \texttt{Truss2DCorot}, que hereda de esta clase.
\end{itemize}


\section{Diálogo}

\begin{itemize}
  \item \textbf{2026-04-21} \ensuremath{\cdot} Cierre de ciclo retroactivo. La clase \texttt{Truss2D} preexistía al flujo spec-first; la spec se creó documentando la formulación ya implementada. Los tests unitarios y de integración satisfacen el único criterio de \texttt{acceptance} declarado (analíticamente, vía los valores de $\mathbf K_e$ y end-to-end vía el solver). Promovido \texttt{status: validated}.
\end{itemize}


\newpage

\section{Truss2DCorot}



\textit{Orden de trabajo. Usuario escribe \textbf{especificación física}, \textbf{formulación numérica} y \textbf{contrato}; la IA rellena \textbf{implementación} y responde en \textbf{diálogo}.}



\section{Especificación física}

\subsection{0. Descripción general}
Barra articulada 2D equivalente a \texttt{Truss2D} pero en régimen de \textbf{grandes desplazamientos y rotaciones} con \textbf{pequeña deformación axial}. La cinemática se actualiza en configuración corriente (corotacional / Updated Lagrangian): el eje de la barra gira con el movimiento y la deformación se mide sobre la longitud actual. Rigidez tangente = rigidez material + rigidez geométrica debida a la fuerza axial.

\subsection{1. Cinemática}
Configuraciones:
\begin{itemize}
  \item Referencia: nodos $\mathbf X_1, \mathbf X_2$, longitud $L_0 = \|\mathbf X_2 - \mathbf X_1\|$.
  \item Corriente: $\mathbf x_i = \mathbf X_i + \mathbf u_i$, longitud $l = \|\mathbf x_2 - \mathbf x_1\|$.
\end{itemize}

Cosenos directores en configuración corriente:
\[
c_\theta = \frac{x_2 - x_1}{l}, \qquad s_\theta = \frac{y_2 - y_1}{l}.
\]

\subsection{2. Medida de deformación}
Ingenieril corotacional (válida para pequeña deformación aunque la rotación sea grande):
\[
\varepsilon = \frac{l - L_0}{L_0}.
\]

\subsection{3. Ecuación constitutiva}
Material 1D: $\sigma = E\,\varepsilon$ (elástico lineal); el contrato admite cualquier material con \texttt{STRAIN\_DIM = 1}.

\subsection{4. Equilibrio --- forma débil incremental}
Principio de trabajos virtuales completo:
\[
\underbrace{\int_V \delta\varepsilon\, \sigma\, dV}_{\delta W_{\text{int}}} \;=\; \underbrace{\int_V \delta\mathbf u^\top \mathbf b\, dV + \int_{\partial V_t} \delta\mathbf u^\top \bar{\mathbf t}\, dA}_{\delta W_{\text{ext}}}.
\]

El elemento calcula \textbf{solo el lado interno}:
\[
\delta W_{\text{int}} = \delta\mathbf u_e^\top \mathbf F_{\text{int}},
\]
con la rigidez tangente consistente separada en parte material y geométrica (ver §9). Las cargas externas (nodales y, si las hubiera, de cuerpo o tracción) las ensambla el solver a partir del input del usuario; este elemento \textbf{no} produce vector nodal equivalente de fuerzas de cuerpo distribuidas a lo largo del eje.


\section{Formulación numérica (FEM)}

\subsection{6. Discretización}
Dos nodos, 2 DOFs por nodo ($u_x, u_y$). Vector elemental $\mathbf u_e = [u_x^{(1)}, u_y^{(1)}, u_x^{(2)}, u_y^{(2)}]^\top$.

\subsection{7. Vector dirección actual}
\[
\mathbf d = [-c_\theta,\; -s_\theta,\; c_\theta,\; s_\theta]^\top, \qquad \mathbf n = [-s_\theta,\; c_\theta,\; s_\theta,\; -c_\theta]^\top,
\]
con $c_\theta, s_\theta$ evaluados en la configuración corriente.

\subsection{8. Matriz B corotacional}
\[
\mathbf B = \frac{1}{L_0}\,\mathbf d^\top, \qquad \varepsilon = \mathbf B\, \mathbf u_e.
\]

\subsection{9. Rigidez tangente}
\[
\mathbf K_T = \mathbf K_M + \mathbf K_G,
\]
\[
\mathbf K_M = \frac{E A}{L_0}\, \mathbf d\,\mathbf d^\top,
\]
\[
\mathbf K_G = \frac{N}{l}\, \mathbf n\,\mathbf n^\top, \qquad N = \sigma A.
\]

\subsection{10. Fuerzas internas}
\[
\mathbf F_{\text{int}} = N\, \mathbf d.
\]

\subsection{11. Cuadratura}
No aplica (forma cerrada; un único punto conceptual en el eje).


\section{Contrato de implementación}

\begin{lstlisting}[language=yaml]
name: Truss2DCorot
kind: element
status: validated

interface:
  dof_names: [ux, uy]
  n_nodes: 2
  strain_dim: 1
  n_integration_points: 1

parameters:
  - { name: A, type: float, required: true, desc: Área de la sección transversal }

material_contract:
  signature: "compute_state(ε, state) -> (σ, E_tangent, state')"
  strain_kind: axial escalar (deformación ingenieril corotacional)

conventions:
  sign: "σ > 0 ⇔ ε > 0 (elongación)"
  voigt: "N/A (escalar)"
  node_orientation: "libre; K_T y F_int invariantes bajo permutación"
  configuration: "Updated Lagrangian — c_θ, s_θ, l se recalculan en cada evaluación"

validity:
  - "|ε| ≲ 1e-2 (pequeña deformación axial)"
  - rotaciones y desplazamientos de cualquier magnitud
  - cargas exclusivamente axiales

out_of_scope:
  - flexión, cortante, momentos
  - pandeo por bifurcación (captura de snap-through requiere arc-length)
  - plasticidad con grandes deformaciones
  - "fuerzas de cuerpo distribuidas (peso propio sobre la barra): el elemento no calcula vector nodal equivalente; el usuario reparte masa a los nodos en el input"

acceptance:
  - name: equivalencia con Truss2D en régimen lineal
    setup: "barra axial empotrada, carga F pequeña (ε ~ 1e-6)"
    expect: "u(L) = F·L/(E·A) con tol_rel 1e-8"
    tol_rel: 1.0e-8

  - name: invariancia bajo rotación de cuerpo rígido
    setup: "aplicar rotación rígida de 45° a ambos nodos (sin deformar)"
    expect: "F_int = 0 y σ = 0"
    tol_abs: 1.0e-10

  - name: rigidez geométrica bajo tensión
    setup: "barra en tensión con N > 0; verificar autovalor transverso"
    expect: "autovalor de K_T en dirección n = N/l"
    tol_rel: 1.0e-10

references:
  - "Crisfield M.A., Non-linear Finite Element Analysis of Solids and Structures, vol.1, §3.3"
  - "Belytschko T., Nonlinear Finite Elements for Continua and Structures, §4.5"

\end{lstlisting}


\section{Implementación}

\begin{itemize}
  \item \textbf{Archivo}: solidum/elements/truss.py
  \item \textbf{Clase}: \texttt{Truss2DCorot} --- hereda de \texttt{Truss2D} (reutiliza \texttt{\_\_init\_\_}, \texttt{ClassVar}s y metadatos de registro) y sobrescribe \texttt{compute\_element\_state} y \texttt{compute\_internal\_forces} para evaluar cosenos directores y longitud en configuración corriente.
  \item \textbf{Tests}: tests/test\_truss.py \ensuremath{\cdot} \texttt{TestTruss2DCorot} --- tres tests, uno por cada \texttt{acceptance}:
  \item \texttt{test\_acceptance\_linear\_limit\_matches\_truss2d} (criterio 1).
  \item \texttt{test\_acceptance\_rigid\_body\_rotation} (criterio 2).
  \item \texttt{test\_acceptance\_geometric\_stiffness\_under\_traction} (criterio 3).
  \item \textbf{Notas de traducción}:
  \item La longitud corriente $l$ y los cosenos $c_\theta, s_\theta$ se recalculan en cada evaluación (no se cachean entre llamadas del solver, como exige Updated Lagrangian).
  \item El vector $\mathbf d$ de §7 se construye con signos $[-c_\theta, -s_\theta, c_\theta, s_\theta]$; el transverso $\mathbf n$ como $[-s_\theta, c_\theta, s_\theta, -c_\theta]$. Ambos con norma $\sqrt 2$ --- el factor se absorbe en los coeficientes $EA/L_0$ y $N/l$.
  \item La rigidez tangente se devuelve ya sumada $\mathbf K_M + \mathbf K_G$; el solver no distingue las dos contribuciones.
  \item El test 3 confronta $\mathbf K_T$ con la descomposición esperada y además verifica $\mathbf K_G \mathbf n = (2N/l)\mathbf n$ (autovalor sobre el vector sin normalizar, que es $N/l$ sobre el versor --- así queda la ambigüedad resuelta explícitamente).
\end{itemize}


\section{Diálogo}

\begin{itemize}
  \item \textbf{2026-04-21} \ensuremath{\cdot} Implementación inicial tras validar el diseño spec-first. Se optó por herencia \texttt{Truss2DCorot(Truss2D)} para reutilizar el constructor (L0, cosenos iniciales) y los \texttt{ClassVar}s del registro; los métodos que dependen de configuración corriente se sobrescriben. La bandera \texttt{large\_strains} que tenía \texttt{Truss2D} previamente se elimina en el mismo commit --- el régimen corotacional vive ahora como elemento independiente.
  \item \textbf{2026-04-21} \ensuremath{\cdot} Normalización del autovalor transverso (criterio 3): la fórmula del libro da $N/l$ sobre el \textbf{versor} $\hat{\mathbf n} = \mathbf n/\sqrt 2$; sobre el vector $\mathbf n$ directo (norma $\sqrt 2$) el autovalor es $2N/l$. El test verifica la segunda forma para evitar raíces cuadradas y la nota de traducción documenta la equivalencia.
\end{itemize}


\newpage

\section{Truss3D}



\textit{Orden de trabajo. El usuario escribe la \textbf{especificación física}, la \textbf{formulación numérica} y el \textbf{contrato}. La IA rellena \textbf{implementación} y responde en \textbf{diálogo} durante el trabajo.}



\section{Especificación física}

\subsection{0. Descripción general}
Elemento sólido 1D de primer orden, inmerso en el espacio tridimensional. Dos nodos articulados; transmite únicamente fuerza axial. Aproximación $C^0$ del desplazamiento. Régimen estrictamente de \textbf{linealidad geométrica}: pequeños desplazamientos, pequeñas rotaciones, pequeñas deformaciones.

\subsection{1. Cinemática de desplazamientos}
Interpolación lineal del desplazamiento axial a lo largo del eje $s \in [0, L]$:
\[
u(s) = a_0 + a_1 s.
\]

\subsection{2. Cinemática de deformaciones}
Única componente (axial, en el sistema local de la barra):
\[
\varepsilon_{ss}(s) = \frac{du}{ds}.
\]

\subsection{3. Ecuación constitutiva}
\[
\sigma_{ss} = E\,\varepsilon_{ss}.
\]
El elemento admite cualquier material 1D con \texttt{STRAIN\_DIM = 1} (elástico lineal, plástico 1D, daño 1D).

\subsection{4. Equilibrio --- forma fuerte}
\[
-\frac{d}{ds}\!\left(EA\,\frac{du}{ds}\right) = b(s), \quad s \in (0, L),
\]
con condiciones de contorno Dirichlet o Neumann ($\sigma A = \bar F$) en los extremos.

\subsection{5. Equilibrio --- forma débil}
\[
\int_0^L \frac{d\,\delta u}{ds}\,EA\,\frac{du}{ds}\,ds \;=\; \int_0^L \delta u\,b\,ds \;+\; \bigl[\delta u\,\bar F\bigr]_{\partial\Omega}, \quad \forall\,\delta u \in V_0.
\]

El elemento calcula \textbf{solo el lado interno}; las cargas externas las ensambla el solver a partir del input del usuario.


\section{Formulación numérica (FEM)}

\subsection{6. Discretización}
Dos nodos en $s = 0$ y $s = L$. Cosenos directores del eje en coordenadas globales:
\[
c_x = \frac{x_2 - x_1}{L}, \quad c_y = \frac{y_2 - y_1}{L}, \quad c_z = \frac{z_2 - z_1}{L},
\]
con $L = \|\mathbf X_2 - \mathbf X_1\|$ (configuración inicial, fija).

\subsection{7. Funciones de forma}
\[
N_1(s) = 1 - \tfrac{s}{L}, \qquad N_2(s) = \tfrac{s}{L}.
\]

\subsection{8. Matriz deformación-desplazamiento}
Local (1 DOF axial por nodo): $\mathbf B_{\text{loc}} = \tfrac{1}{L}[-1,\; 1]$.
Global con $\mathbf u_e = [u_x^{(1)}, u_y^{(1)}, u_z^{(1)}, u_x^{(2)}, u_y^{(2)}, u_z^{(2)}]^\top$:
\[
\mathbf B = \tfrac{1}{L}\,[-c_x,\; -c_y,\; -c_z,\; c_x,\; c_y,\; c_z], \qquad \varepsilon = \mathbf B\,\mathbf u_e.
\]

\subsection{9. Rigidez elemental}
Integración analítica exacta ($EA$ constante, $\mathbf B$ constante):
\[
\mathbf K_e = \int_0^L \mathbf B^\top EA\,\mathbf B\,ds = \frac{EA}{L}\,\mathbf d\,\mathbf d^\top, \qquad \mathbf d = [-c_x,\; -c_y,\; -c_z,\; c_x,\; c_y,\; c_z]^\top.
\]

$\mathbf K_e$ es una matriz $6\times 6$ de rango 1.

\subsection{10. Fuerzas internas}
\[
\mathbf F_{\text{int}} = \sigma A\,\mathbf d = N\,\mathbf d.
\]

\subsection{11. Cuadratura}
No aplica (forma cerrada, $\mathbf B$ y $\sigma$ uniformes por elemento). El campo \texttt{n\_integration\_points: 1} se refiere al número de estados materiales conceptuales (uno por elemento), no a puntos de Gauss.


\section{Contrato de implementación}

\begin{lstlisting}[language=yaml]
name: Truss3D
kind: element
status: validated

interface:
  dof_names: [ux, uy, uz]
  n_nodes: 2
  strain_dim: 1
  n_integration_points: 1

parameters:
  - { name: A, type: float, required: true, desc: Área de la sección transversal }

material_contract:
  signature: "compute_state(ε, state) -> (σ, E_tangent, state')"
  strain_kind: axial escalar

conventions:
  sign: "σ > 0 ⇔ ε > 0 (elongación)"
  voigt: "N/A (escalar)"
  node_orientation: "libre; K_e y F_int invariantes bajo permutación"
  configuration: "inicial fija — L, c_x, c_y, c_z no se actualizan con los desplazamientos"

validity:
  - "|ε| ≲ 1e-2"
  - pequeños desplazamientos y rotaciones
  - cargas exclusivamente axiales

out_of_scope:
  - flexión, cortante, momentos
  - grandes desplazamientos o rotaciones (para ese régimen usar `Truss3DCorot`)
  - pandeo
  - "fuerzas de cuerpo distribuidas sobre el eje: el usuario reparte masa a los nodos"

acceptance:
  - name: barra axial bajo carga puntual (3D)
    setup: "empotrada en s=0, carga F axial en s=L, orientación arbitraria en el espacio"
    expect: "u(L) = F·L / (E·A) proyectado sobre el eje"
    tol_rel: 1.0e-10

references:
  - "Bathe K.-J., Finite Element Procedures, §4.2.1"
  - "Cook R.D., Concepts and Applications of FEA, §2.3"

\end{lstlisting}


\section{Implementación}

\begin{itemize}
  \item \textbf{Archivo}: solidum/elements/truss.py
  \item \textbf{Clase}: \texttt{Truss3D} --- hereda directamente de \texttt{Element} (sin relación de herencia con \texttt{Truss2D} ni con \texttt{Truss2DCorot}).
  \item \textbf{Tests}: tests/test\_truss.py \ensuremath{\cdot} \texttt{TestTruss3D} --- verifica \texttt{L0}, cosenos directores $(c_x, c_y, c_z)$, entradas de $\mathbf K_e$ (incluyendo simetría, dimensiones $6\times 6$) y evaluación de $\mathbf F_{\text{int}}$ sobre desplazamientos conocidos con geometría de prueba $L=13$, $(c_x,c_y,c_z)=(3/13, 4/13, 12/13)$.
  \item \textbf{Notas de traducción}:
  \item \texttt{L0}, \texttt{cx}, \texttt{cy}, \texttt{cz} se calculan una vez en \texttt{\_\_init\_\_} sobre la configuración inicial y no se actualizan --- coherente con el régimen geométricamente lineal declarado.
  \item El constructor tolera nodos con 2 ó 3 coordenadas (completa con $z=0$ si falta), para facilitar transición de problemas planos embebidos.
  \item El contrato con el material es el estándar del proyecto: \texttt{material.compute\_state(\ensuremath{\varepsilon}, state) \ensuremath{\to} (\ensuremath{\sigma}, E\ensuremath{_{t}}, state')}.
  \item Para grandes desplazamientos/rotaciones en 3D usar \texttt{Truss3DCorot}, que hereda de esta clase.
\end{itemize}


\section{Diálogo}

\begin{itemize}
  \item \textbf{2026-04-21} \ensuremath{\cdot} Apertura de spec retroactiva. La clase \texttt{Truss3D} preexistía al flujo spec-first; la spec se crea documentando la formulación ya implementada y estableciendo explícitamente su alcance: régimen de linealidad geométrica, sin bandera ni variante corotacional. Los tests unitarios existentes cubren el criterio de \texttt{acceptance} (vía los valores de $\mathbf K_e = (EA/L)\mathbf d\mathbf d^\top$ que implican algebraicamente $u(L)=FL/(EA)$). Promovido \texttt{status: validated}.
\end{itemize}


\newpage

\section{Truss3DCorot}



\textit{Orden de trabajo. Usuario escribe \textbf{especificación física}, \textbf{formulación numérica} y \textbf{contrato}; la IA rellena \textbf{implementación} y responde en \textbf{diálogo}.}



\section{Especificación física}

\subsection{0. Descripción general}
Barra articulada 3D equivalente a \texttt{Truss3D} pero en régimen de \textbf{grandes desplazamientos y rotaciones} con \textbf{pequeña deformación axial}. Updated Lagrangian / corotacional: el eje de la barra se redefine en cada iteración sobre la configuración corriente y la deformación se mide sobre la longitud actual. La rigidez tangente suma material + geométrica.

\subsection{1. Cinemática}
Configuraciones:
\begin{itemize}
  \item Referencia: $\mathbf X_1, \mathbf X_2 \in \mathbb{R}^3$, longitud $L_0 = \|\mathbf X_2 - \mathbf X_1\|$.
  \item Corriente: $\mathbf x_i = \mathbf X_i + \mathbf u_i$, longitud $l = \|\mathbf x_2 - \mathbf x_1\|$.
\end{itemize}

Vector unitario del eje corriente:
\[
\hat{\mathbf e} = \frac{\mathbf x_2 - \mathbf x_1}{l} = (c_x, c_y, c_z), \qquad c_x^2 + c_y^2 + c_z^2 = 1.
\]

\subsection{2. Medida de deformación}
Ingenieril corotacional:
\[
\varepsilon = \frac{l - L_0}{L_0}.
\]

\subsection{3. Ecuación constitutiva}
Material 1D con \texttt{STRAIN\_DIM = 1}: $\sigma = \sigma(\varepsilon)$ (el elemento no hardcodea la ley; la provee el material).

\subsection{4. Equilibrio --- forma débil incremental}
PTV estándar. El elemento calcula solo el lado interno:
\[
\delta W_{\text{int}} = \delta\mathbf u_e^\top\,\mathbf F_{\text{int}}, \qquad \mathbf K_T = \frac{\partial\mathbf F_{\text{int}}}{\partial\mathbf u_e} = \mathbf K_M + \mathbf K_G.
\]


\section{Formulación numérica (FEM)}

\subsection{6. Discretización}
Dos nodos, 3 DOFs por nodo ($u_x, u_y, u_z$). Vector elemental de 6 componentes:
\[
\mathbf u_e = [u_x^{(1)}, u_y^{(1)}, u_z^{(1)}, u_x^{(2)}, u_y^{(2)}, u_z^{(2)}]^\top.
\]

\subsection{7. Vector dirección actual}
\[
\mathbf d = [-c_x,\; -c_y,\; -c_z,\; c_x,\; c_y,\; c_z]^\top,
\]
con los cosenos evaluados en configuración \textbf{corriente}. $\|\mathbf d\|^2 = 2$.

\subsection{8. Matriz B corotacional}
\[
\mathbf B = \frac{1}{L_0}\,\mathbf d^\top, \qquad \varepsilon = \mathbf B\,\mathbf u_e.
\]

\subsection{9. Rigidez tangente}
\textbf{Material}:
\[
\mathbf K_M = \frac{E_t\,A}{L_0}\,\mathbf d\,\mathbf d^\top \quad (6\times 6,\ \text{rango 1}).
\]

\textbf{Geométrica}: en 3D la dirección perpendicular al eje es un \textbf{plano} (no un vector único). Se usa el proyector perpendicular $\mathbf P_3 = \mathbf I_3 - \hat{\mathbf e}\hat{\mathbf e}^\top$ (matriz $3\times 3$, rango 2). La rigidez geométrica en el elemento de 6 DOFs:
\[
\mathbf K_G = \frac{N}{l}\,\begin{bmatrix}\mathbf P_3 & -\mathbf P_3 \\ -\mathbf P_3 & \mathbf P_3\end{bmatrix} \quad (6\times 6,\ \text{rango 2}), \qquad N = \sigma A.
\]

\subsection{10. Fuerzas internas}
\[
\mathbf F_{\text{int}} = N\,\mathbf d.
\]

\subsection{11. Cuadratura}
No aplica (forma cerrada).


\section{Contrato de implementación}

\begin{lstlisting}[language=yaml]
name: Truss3DCorot
kind: element
status: validated

interface:
  dof_names: [ux, uy, uz]
  n_nodes: 2
  strain_dim: 1
  n_integration_points: 1

parameters:
  - { name: A, type: float, required: true, desc: Área de la sección transversal }

material_contract:
  signature: "compute_state(ε, state) -> (σ, E_tangent, state')"
  strain_kind: axial escalar (deformación ingenieril corotacional)

conventions:
  sign: "σ > 0 ⇔ ε > 0 (elongación)"
  voigt: "N/A (escalar)"
  node_orientation: "libre"
  configuration: "Updated Lagrangian — l, c_x, c_y, c_z se recalculan en cada evaluación"

validity:
  - "|ε| ≲ 1e-2"
  - rotaciones y desplazamientos de cualquier magnitud en el espacio
  - cargas exclusivamente axiales

out_of_scope:
  - flexión, cortante, momentos
  - pandeo por bifurcación
  - plasticidad con grandes deformaciones
  - "fuerzas de cuerpo distribuidas: el usuario reparte masa a los nodos"

acceptance:
  - name: equivalencia con Truss3D en régimen lineal
    setup: "barra 3D con orientación arbitraria, carga axial pequeña (ε ~ 1e-6)"
    expect: "K_T y F_int coinciden con Truss3D lineal a tolerancia rtol=1e-4"
    tol_rel: 1.0e-4

  - name: invariancia bajo rotación rígida 3D
    setup: "aplicar una rotación rígida arbitraria (p. ej. 30° sobre eje z, luego 45° sobre eje x)"
    expect: "F_int = 0 y σ = 0"
    tol_abs: 1.0e-10

  - name: rigidez geométrica transversa (plano perpendicular)
    setup: "barra en tensión con N > 0; aplicar K_G sobre dos vectores perpendiculares al eje linealmente independientes"
    expect: "K_G·v_perp = (N/l)·v_perp_struct, donde v_perp_struct es la extensión del vector transverso al espacio de 6 DOFs con signos opuestos en los dos nodos"
    tol_rel: 1.0e-10

references:
  - "Crisfield M.A., Non-linear Finite Element Analysis of Solids and Structures, vol.1, §3.3"
  - "Belytschko T., Nonlinear Finite Elements for Continua and Structures, §4.5"

\end{lstlisting}


\section{Implementación}

\begin{itemize}
  \item \textbf{Archivo}: solidum/elements/truss.py
  \item \textbf{Clase}: \texttt{Truss3DCorot} --- hereda de \texttt{Truss3D} (reutiliza \texttt{\_\_init\_\_}, tolerancia a nodos 2D/3D, metadatos de registro) y sobrescribe \texttt{compute\_element\_state} y \texttt{compute\_internal\_forces} para evaluar longitud y cosenos directores en configuración corriente.
  \item \textbf{Tests}: tests/test\_truss.py \ensuremath{\cdot} \texttt{TestTruss3DCorot} --- tres tests, uno por cada \texttt{acceptance}:
  \item \texttt{test\_acceptance\_linear\_limit\_matches\_truss3d} (criterio 1).
  \item \texttt{test\_acceptance\_rigid\_body\_rotation} (criterio 2 --- dos rotaciones rígidas distintas, verifica \ensuremath{\sigma}=0 y F\_int=0 en ambas).
  \item \texttt{test\_acceptance\_geometric\_stiffness\_transverse\_plane} (criterio 3 --- aplica $\mathbf K_G$ a \textbf{dos} direcciones linealmente independientes del plano perpendicular al eje).
  \item \textbf{Notas de traducción}:
  \item La clave respecto a 2D es el proyector $\mathbf P_3 = \mathbf I - \hat{\mathbf e}\hat{\mathbf e}^\top$ (3\ensuremath{\times}3, rango 2): el plano perpendicular al eje es bidimensional. $\mathbf K_G$ resulta entonces de rango 2, no rango 1.
  \item El proyector se construye con \texttt{solidum.math.geometry.perpendicular\_projector(ê)}, helper compartido con \texttt{Cable3DCorot} (operación geométrica pura, no específica de armaduras ni de cables).
  \item La longitud corriente $l$ y los cosenos $c_x, c_y, c_z$ se recalculan en cada evaluación (Updated Lagrangian); no se cachean entre llamadas del solver.
  \item Autovalores de $\mathbf K_G$ sobre modos transversos (rotación de la barra alrededor de su centro): $2N/l$ sobre el vector sin normalizar --- igual que en 2D, el factor 2 viene de que el vector de 6 componentes tiene norma $\sqrt 2$. Sobre el versor el autovalor es $N/l$.
\end{itemize}


\section{Diálogo}

\begin{itemize}
  \item \textbf{2026-04-21} \ensuremath{\cdot} Implementación tras validar el diseño spec-first. Se optó por herencia \texttt{Truss3DCorot(Truss3D)} por paralelismo con \texttt{Truss2DCorot(Truss2D)} y reutilización del constructor. La diferencia estructural con el caso 2D es el proyector $\mathbf P$: en 2D tenía rango 1 y generaba un $\mathbf K_G$ de rango 1 (vector $\mathbf n$ único); en 3D tiene rango 2 y $\mathbf K_G$ rango 2 (plano perpendicular al eje). Los tests del criterio 3 verifican explícitamente \textbf{dos} direcciones transversas linealmente independientes, capturando este cambio dimensional.
  \item \textbf{2026-04-21} \ensuremath{\cdot} Test de invariancia bajo rotación rígida 3D: se verifican dos rotaciones distintas (plano x-y y plano x-z) para descartar coincidencias accidentales en ejes canónicos.
\end{itemize}


\newpage

\chapter{Elementos 1D — Cables}

\section{Cable2DCorot}



\textit{Orden de trabajo. El usuario escribe \textbf{especificación física}, \textbf{formulación numérica} y \textbf{contrato}; la IA rellena \textbf{implementación} y responde en \textbf{diálogo}.}



\section{Especificación física}

\subsection{0. Descripción general}
Elemento 1D inmerso en el plano 2D que modela un \textbf{cable}: dos nodos articulados que solo transmiten \textbf{fuerza axial de tensión}. En compresión o estado sin tensar, el elemento no transmite nada --- se destensa. Cinemática \textbf{corotacional} (Updated Lagrangian): el eje rota con el movimiento, la longitud se mide sobre la configuración corriente.

La propiedad de unilateralidad vive en el \textbf{material}; el elemento es la maquinaria cinemática que aplica al material la deformación correcta. Para obtener el comportamiento propio de cable se empareja este elemento con un material unilateral como \texttt{CableMaterial1D}. Con un material elástico clásico, el elemento se comporta como una barra articulada corotacional (útil para pruebas, no para modelado realista de cables).

\subsection{1. Cinemática}
Configuraciones:
\begin{itemize}
  \item Referencia: $\mathbf X_1, \mathbf X_2 \in \mathbb{R}^2$, longitud $L_0 = \|\mathbf X_2 - \mathbf X_1\|$.
  \item Corriente: $\mathbf x_i = \mathbf X_i + \mathbf u_i$, longitud $l = \|\mathbf x_2 - \mathbf x_1\|$.
\end{itemize}

Cosenos directores del eje \textbf{corriente}:
\[
c_\theta = \frac{x_2 - x_1}{l}, \qquad s_\theta = \frac{y_2 - y_1}{l}.
\]

\subsection{2. Medida de deformación}
Ingenieril corotacional:
\[
\varepsilon = \frac{l - L_0}{L_0}.
\]

\subsection{3. Ecuación constitutiva}
Delegada al material con \texttt{STRAIN\_DIM = 1}. El elemento invoca \texttt{material.compute\_state(\ensuremath{\varepsilon}, state)} y recibe $(\sigma, E_t, \text{nuevo estado})$ sin inspeccionar su naturaleza. La unilateralidad (si la hay) es responsabilidad del material --- el elemento nunca filtra $\sigma$ ni $\varepsilon$ por signo.

\subsection{4. Equilibrio --- forma débil incremental}
Principio de trabajos virtuales, lado interno:
\[
\delta W_{\text{int}} = \int_V \delta\varepsilon\,\sigma\,dV = \delta\mathbf u_e^\top\,\mathbf F_{\text{int}}, \qquad \mathbf K_T = \frac{\partial\mathbf F_{\text{int}}}{\partial\mathbf u_e} = \mathbf K_M + \mathbf K_G.
\]

Las cargas externas las ensambla el solver. El elemento no calcula vector nodal equivalente de fuerzas de cuerpo.


\section{Formulación numérica (FEM)}

\subsection{6. Discretización}
Dos nodos, 2 DOFs por nodo ($u_x, u_y$). Vector elemental:
\[
\mathbf u_e = [u_x^{(1)},\; u_y^{(1)},\; u_x^{(2)},\; u_y^{(2)}]^\top.
\]

\subsection{7. Vectores de dirección (configuración corriente)}
\[
\mathbf d = [-c_\theta,\; -s_\theta,\; c_\theta,\; s_\theta]^\top, \qquad \mathbf n = [-s_\theta,\; c_\theta,\; s_\theta,\; -c_\theta]^\top.
\]

$\mathbf d$ apunta a lo largo del eje corriente (norma $\sqrt 2$); $\mathbf n$ al perpendicular (norma $\sqrt 2$). Son ortogonales: $\mathbf d^\top\mathbf n = 0$.

\subsection{8. Matriz B corotacional}
\[
\mathbf B = \frac{1}{L_0}\,\mathbf d^\top, \qquad \varepsilon = \mathbf B\,\mathbf u_e.
\]

\subsection{9. Rigidez tangente}
\[
\mathbf K_T = \mathbf K_M + \mathbf K_G,
\]
\[
\mathbf K_M = \frac{E_t\,A}{L_0}\,\mathbf d\,\mathbf d^\top, \qquad \mathbf K_G = \frac{N}{l}\,\mathbf n\,\mathbf n^\top, \qquad N = \sigma A.
\]

\textbf{Caso destensado} (material unilateral con $\varepsilon \leq 0$): el material devuelve $\sigma = 0$, $E_t = 0$, luego $N = 0$ y \textbf{ambas contribuciones se anulan} --- $\mathbf K_T = \mathbf 0$. El elemento aporta cero al sistema global, como físicamente corresponde a un cable flojo.

\subsection{10. Fuerzas internas}
\[
\mathbf F_{\text{int}} = N\,\mathbf d.
\]

Cuando el cable está destensado ($N = 0$), $\mathbf F_{\text{int}} = \mathbf 0$.

\subsection{11. Cuadratura}
No aplica (forma cerrada).


\section{Contrato de implementación}

\begin{lstlisting}[language=yaml]
name: Cable2DCorot
kind: element
status: validated

interface:
  dof_names: [ux, uy]
  n_nodes: 2
  strain_dim: 1
  n_integration_points: 1

parameters:
  - { name: A, type: float, required: true, desc: "Área de la sección transversal del cable" }

material_contract:
  signature: "compute_state(ε, state) -> (σ, E_tangent, state')"
  strain_kind: axial escalar (deformación ingenieril corotacional)
  expected_behaviour: "material unilateral (p. ej. CableMaterial1D) para obtener respuesta física de cable; cualquier material con STRAIN_DIM=1 es técnicamente aceptado pero producirá comportamiento de barra"

conventions:
  sign: "σ > 0 ⇔ ε > 0 (elongación); la responsabilidad de anular σ en compresión es del material"
  voigt: "N/A (escalar)"
  node_orientation: "libre; K_T y F_int invariantes bajo permutación"
  configuration: "Updated Lagrangian — l, c_θ, s_θ se recalculan en cada evaluación"

validity:
  - "|ε| ≲ 1e-2 en régimen tensado"
  - rotaciones y desplazamientos de cualquier magnitud
  - cargas axiales (las transversales se transmiten por el destensado-retensado del cable)

out_of_scope:
  - flexión, cortante, momentos
  - dinámica con inercia distribuida (sin matriz de masa)
  - pandeo (irrelevante: el cable no transmite compresión)
  - "fuerzas de cuerpo distribuidas: el usuario reparte masa/peso a los nodos"
  - pretensión mediante parámetro del elemento (modelar con longitud inicial ajustada)

numerical_caveats:
  - "Un cable completamente destensado aporta K_T = 0 al sistema global. Si otros elementos no garantizan rigidez suficiente en los DOFs compartidos, la matriz tangente global puede ser singular o mal condicionada."
  - "En transiciones tensado ↔ destensado (ε cruza 0), Newton-Raphson puede oscilar por el salto discontinuo en E_t. Usar arc-length o pasos de carga finos si el problema atraviesa destensiones."

acceptance:
  - name: cable tensado coincide con barra corotacional
    setup: "cable horizontal con material lineal (Elastic1D, E=1000), u_x nodo 2 = 1e-3 ⇒ ε = 1e-3 > 0"
    expect: "σ = E·ε, E_t = E, K_T = K_M + K_G con valores de barra corotacional"
    tol_rel: 1.0e-10

  - name: cable destensado produce aportación nula
    setup: "cable horizontal con CableMaterial1D, u_x nodo 2 = -1e-3 ⇒ ε < 0"
    expect: "σ = 0, F_int = 0, K_T = matriz cero"
    tol_abs: 1.0e-12

  - name: invariancia bajo rotación rígida
    setup: "cable inicialmente horizontal, rotación rígida de 45° de ambos nodos"
    expect: "σ = 0 y F_int = 0 (independiente del material)"
    tol_abs: 1.0e-10

  - name: cruce por cero
    setup: "u_e tal que ε = 0 exactamente, material unilateral"
    expect: "σ = 0, E_t = 0, F_int = 0, K_T = 0"
    tol_abs: 1.0e-12

references:
  - "Crisfield M.A., Non-linear Finite Element Analysis of Solids and Structures, vol.1, §3.3"
  - "Irvine H.M., Cable Structures, MIT Press, §1.2"

\end{lstlisting}


\section{Implementación}

\begin{itemize}
  \item \textbf{Archivo}: solidum/elements/cable.py --- archivo propio, no referencia a los módulos de armadura ni a ningún otro elemento.
  \item \textbf{Clase}: \texttt{Cable2DCorot} --- hereda directamente de \texttt{Element} (no de \texttt{Truss2DCorot} ni de ningún elemento de armadura). La maquinaria cinemática corotacional se reimplementa íntegra dentro de la clase.
  \item \textbf{Tests}: tests/test\_cable\_elements.py \ensuremath{\cdot} \texttt{TestCable2DCorot} --- los cuatro criterios de \texttt{acceptance} más un test de registro:
  \item \texttt{test\_acceptance\_cable\_tensado\_como\_barra\_corot} (criterio 1).
  \item \texttt{test\_acceptance\_cable\_destensado\_aportacion\_nula} (criterio 2).
  \item \texttt{test\_acceptance\_rotacion\_rigida} (criterio 3).
  \item \texttt{test\_acceptance\_cruce\_por\_cero} (criterio 4).
  \item \texttt{test\_registro\_en\_registry} --- autodiscover.
  \item \textbf{Notas de traducción}:
  \item La clase duplica deliberadamente la lógica de \texttt{Truss2DCorot}. No hay herencia. Es el precio de la independencia: si mañana se toca la armadura corotacional, este cable no se mueve.
  \item \texttt{\_current\_geometry(u\_e)} devuelve \texttt{(l, c\_\ensuremath{\theta}, s\_\ensuremath{\theta})} --- método privado para aislar la única operación que depende de configuración corriente.
  \item En estado destensado, $\mathbf K_M = \mathbf 0$ (porque $E_t = 0$ viene del material) y $\mathbf K_G = \mathbf 0$ (porque $N = 0$); el elemento aporta literalmente ceros al ensamblaje, como debe ser físicamente un cable flojo. Los tests lo verifican con tolerancia absoluta.
  \item El elemento no valida que el material sea unilateral: es responsabilidad del usuario emparejar con \texttt{CableMaterial1D}. Emparejar con \texttt{Elastic1D} produce respuesta de barra corotacional y se usa internamente en el test del criterio 1 para comparar contra valores analíticos.
\end{itemize}


\section{Diálogo}

\begin{itemize}
  \item \textbf{2026-04-21} \ensuremath{\cdot} Elemento creado como componente totalmente independiente de \texttt{Truss2DCorot} por decisión explícita del usuario: la maquinaria cinemática se duplica dentro de la clase de cable para desacoplar su evolución futura del elemento de armadura. El coste es \textasciitilde{}40 líneas duplicadas; la ganancia es que ninguna refactorización de las armaduras puede contaminar silenciosamente el comportamiento de los cables.
  \item \textbf{2026-04-21} \ensuremath{\cdot} Validación de material unilateral: se optó por \textbf{no} imponerla en \texttt{\_\_init\_\_}. Motivo: el contrato elemento\ensuremath{\leftrightarrow}material del proyecto es genérico sobre \texttt{STRAIN\_DIM=1}; introducir una verificación específica crearía un acoplamiento entre la clase del elemento y la identidad del material. Se documenta en \texttt{material\_contract.expected\_behaviour} del YAML para orientar al usuario.
\end{itemize}


\newpage

\section{Cable3DCorot}



\textit{Orden de trabajo. El usuario escribe \textbf{especificación física}, \textbf{formulación numérica} y \textbf{contrato}; la IA rellena \textbf{implementación} y responde en \textbf{diálogo}.}



\section{Especificación física}

\subsection{0. Descripción general}
Elemento 1D inmerso en el espacio 3D que modela un \textbf{cable}: dos nodos articulados que solo transmiten \textbf{fuerza axial de tensión}. En compresión o estado sin tensar, el elemento no transmite nada. Cinemática \textbf{corotacional} (Updated Lagrangian): el eje rota libremente en el espacio, la longitud se mide sobre la configuración corriente.

La unilateralidad vive en el \textbf{material}; el elemento es la maquinaria cinemática 3D. Para obtener respuesta física de cable, emparejar con \texttt{CableMaterial1D}. Con un material elástico clásico, el elemento se comporta como una barra corotacional 3D.

\subsection{1. Cinemática}
\begin{itemize}
  \item Referencia: $\mathbf X_1, \mathbf X_2 \in \mathbb{R}^3$, longitud $L_0 = \|\mathbf X_2 - \mathbf X_1\|$.
  \item Corriente: $\mathbf x_i = \mathbf X_i + \mathbf u_i$, longitud $l = \|\mathbf x_2 - \mathbf x_1\|$.
\end{itemize}

Vector unitario del eje corriente:
\[
\hat{\mathbf e} = \frac{\mathbf x_2 - \mathbf x_1}{l} = (c_x, c_y, c_z), \qquad c_x^2 + c_y^2 + c_z^2 = 1.
\]

\subsection{2. Medida de deformación}
Ingenieril corotacional:
\[
\varepsilon = \frac{l - L_0}{L_0}.
\]

\subsection{3. Ecuación constitutiva}
Delegada al material con \texttt{STRAIN\_DIM = 1}. El elemento llama \texttt{material.compute\_state(\ensuremath{\varepsilon}, state)} y recibe $(\sigma, E_t, \text{nuevo estado})$ sin inspeccionar su naturaleza. La unilateralidad (si la hay) es responsabilidad del material.

\subsection{4. Equilibrio --- forma débil incremental}
PTV estándar; el elemento calcula solo el lado interno:
\[
\delta W_{\text{int}} = \delta\mathbf u_e^\top\,\mathbf F_{\text{int}}, \qquad \mathbf K_T = \frac{\partial\mathbf F_{\text{int}}}{\partial\mathbf u_e} = \mathbf K_M + \mathbf K_G.
\]


\section{Formulación numérica (FEM)}

\subsection{6. Discretización}
Dos nodos, 3 DOFs por nodo ($u_x, u_y, u_z$). Vector elemental de 6 componentes:
\[
\mathbf u_e = [u_x^{(1)}, u_y^{(1)}, u_z^{(1)}, u_x^{(2)}, u_y^{(2)}, u_z^{(2)}]^\top.
\]

\subsection{7. Vector dirección actual}
\[
\mathbf d = [-c_x,\; -c_y,\; -c_z,\; c_x,\; c_y,\; c_z]^\top, \qquad \|\mathbf d\|^2 = 2,
\]
con los cosenos evaluados en configuración \textbf{corriente}.

\subsection{8. Matriz B corotacional}
\[
\mathbf B = \frac{1}{L_0}\,\mathbf d^\top, \qquad \varepsilon = \mathbf B\,\mathbf u_e.
\]

\subsection{9. Rigidez tangente}
\textbf{Material}:
\[
\mathbf K_M = \frac{E_t\,A}{L_0}\,\mathbf d\,\mathbf d^\top \quad (6\times 6,\ \text{rango 1}).
\]

\textbf{Geométrica}: en 3D la dirección perpendicular al eje es un plano bidimensional. Se usa el proyector perpendicular $\mathbf P_3 = \mathbf I_3 - \hat{\mathbf e}\hat{\mathbf e}^\top$ ($3\times 3$, rango 2):
\[
\mathbf K_G = \frac{N}{l}\,\begin{bmatrix}\mathbf P_3 & -\mathbf P_3 \\ -\mathbf P_3 & \mathbf P_3\end{bmatrix} \quad (6\times 6,\ \text{rango 2}), \qquad N = \sigma A.
\]

\textbf{Caso destensado} (material unilateral con $\varepsilon \leq 0$): el material devuelve $\sigma = 0$, $E_t = 0$, luego $N = 0$ y \textbf{ambas contribuciones se anulan} --- $\mathbf K_T = \mathbf 0$. El elemento aporta cero al sistema global.

\subsection{10. Fuerzas internas}
\[
\mathbf F_{\text{int}} = N\,\mathbf d.
\]

\subsection{11. Cuadratura}
No aplica (forma cerrada).


\section{Contrato de implementación}

\begin{lstlisting}[language=yaml]
name: Cable3DCorot
kind: element
status: validated

interface:
  dof_names: [ux, uy, uz]
  n_nodes: 2
  strain_dim: 1
  n_integration_points: 1

parameters:
  - { name: A, type: float, required: true, desc: "Área de la sección transversal del cable" }

material_contract:
  signature: "compute_state(ε, state) -> (σ, E_tangent, state')"
  strain_kind: axial escalar (deformación ingenieril corotacional)
  expected_behaviour: "material unilateral (p. ej. CableMaterial1D) para obtener respuesta física de cable; cualquier material con STRAIN_DIM=1 es técnicamente aceptado pero producirá comportamiento de barra"

conventions:
  sign: "σ > 0 ⇔ ε > 0 (elongación); la responsabilidad de anular σ en compresión es del material"
  voigt: "N/A (escalar)"
  node_orientation: "libre; K_T y F_int invariantes bajo permutación"
  configuration: "Updated Lagrangian — l, c_x, c_y, c_z se recalculan en cada evaluación"

validity:
  - "|ε| ≲ 1e-2 en régimen tensado"
  - rotaciones y desplazamientos de cualquier magnitud en el espacio
  - cargas axiales (las transversales se transmiten por destensado-retensado del cable)

out_of_scope:
  - flexión, cortante, momentos, torsión
  - dinámica con inercia distribuida (sin matriz de masa)
  - pandeo (irrelevante: el cable no transmite compresión)
  - "fuerzas de cuerpo distribuidas: el usuario reparte masa/peso a los nodos"
  - pretensión mediante parámetro del elemento

numerical_caveats:
  - "Cable destensado ⇒ K_T = 0 en los DOFs del elemento. Si otros elementos no garantizan rigidez en esos DOFs, la matriz tangente global puede ser singular."
  - "En transiciones tensado ↔ destensado (ε cruza 0), Newton-Raphson puede oscilar. Usar arc-length o pasos finos si el problema atraviesa destensiones."

acceptance:
  - name: cable tensado coincide con barra corotacional 3D
    setup: "cable con orientación arbitraria, material lineal (Elastic1D), desplazamiento axial pequeño (ε ~ 1e-6)"
    expect: "K_T y F_int coinciden con Truss3DCorot a tolerancia rtol = 1e-4"
    tol_rel: 1.0e-4

  - name: cable destensado produce aportación nula
    setup: "cable sobre eje x con CableMaterial1D, u_x nodo 2 = -1e-3 ⇒ ε < 0"
    expect: "σ = 0, F_int = 0, K_T = matriz cero"
    tol_abs: 1.0e-12

  - name: invariancia bajo rotación rígida 3D
    setup: "cable inicial sobre eje x, rotación rígida de 45° en plano x-z"
    expect: "σ = 0 y F_int = 0 (independiente del material)"
    tol_abs: 1.0e-10

  - name: cruce por cero
    setup: "u_e = 0 con material unilateral"
    expect: "σ = 0, E_t = 0, F_int = 0, K_T = 0"
    tol_abs: 1.0e-12

references:
  - "Crisfield M.A., Non-linear Finite Element Analysis of Solids and Structures, vol.1, §3.3"
  - "Belytschko T., Nonlinear Finite Elements for Continua and Structures, §4.5"
  - "Irvine H.M., Cable Structures, MIT Press, §1.2"

\end{lstlisting}


\section{Implementación}

\begin{itemize}
  \item \textbf{Archivo}: solidum/elements/cable.py --- comparte archivo con \texttt{Cable2DCorot} por convención temática del proyecto; no comparte código.
  \item \textbf{Clase}: \texttt{Cable3DCorot} --- hereda directamente de \texttt{Element} (no de \texttt{Cable2DCorot}, no de \texttt{Truss3DCorot}, no de ningún otro elemento). Maquinaria cinemática 3D autónoma.
  \item \textbf{Tests}: tests/test\_cable\_elements.py \ensuremath{\cdot} \texttt{TestCable3DCorot} --- los 4 criterios de \texttt{acceptance} más un test de registro:
  \item \texttt{test\_acceptance\_cable\_tensado\_como\_barra\_corot\_3d} (criterio 1).
  \item \texttt{test\_acceptance\_cable\_destensado\_aportacion\_nula} (criterio 2).
  \item \texttt{test\_acceptance\_rotacion\_rigida\_3d} (criterio 3).
  \item \texttt{test\_acceptance\_cruce\_por\_cero} (criterio 4).
  \item \texttt{test\_registro\_en\_registry} --- autodiscover.
  \item \textbf{Notas de traducción}:
  \item La clase duplica la maquinaria corotacional 3D de \texttt{Truss3DCorot}. Duplicación deliberada por la misma razón que en 2D: desacoplar la evolución futura de armaduras y cables.
  \item El proyector $\mathbf P_3 = \mathbf I_3 - \hat{\mathbf e}\hat{\mathbf e}^\top$ se construye con \texttt{solidum.math.geometry.perpendicular\_projector(ê)}, helper compartido con \texttt{Truss3DCorot} (operación geométrica pura, ortogonal al hecho de que el material sea unilateral o no).
  \item \texttt{\_current\_geometry} tolera nodos con 2 ó 3 coordenadas (completa con $z=0$ si falta), heredando la flexibilidad de \texttt{Truss3D} para problemas embebidos.
  \item En estado destensado, $\mathbf K_M = \mathbf 0$ y $\mathbf K_G = \mathbf 0$ simultáneamente; el elemento aporta literalmente ceros al ensamblaje.
  \item El test del criterio 1 compara directamente contra \texttt{Truss3DCorot} con material \texttt{Elastic1D}: si ambos producen $\mathbf K$ y $\mathbf F_{\text{int}}$ iguales a tolerancia, la cinemática del cable y de la barra son consistentes entre sí en el régimen tensado.
\end{itemize}


\section{Diálogo}

\begin{itemize}
  \item \textbf{2026-04-21} \ensuremath{\cdot} Elemento creado como pieza totalmente autónoma por decisión del usuario: no hereda de \texttt{Cable2DCorot} ni de \texttt{Truss3DCorot}. Compartir archivo con \texttt{Cable2DCorot} (\texttt{solidum/elements/cable.py}) es convención temática del proyecto (elementos del mismo dominio conviven en un archivo, como las 4 armaduras y los 2 frames 2D en \texttt{structural.py}), no implica dependencia de código.
  \item \textbf{2026-04-21} \ensuremath{\cdot} La única diferencia estructural con \texttt{Cable2DCorot} es el proyector 3D: en 2D la perpendicular es un vector único ($\mathbf n$); en 3D es un plano (dimensión 2). Esto se refleja en el rango de $\mathbf K_G$: 1 en 2D, 2 en 3D.
\end{itemize}


\newpage

\chapter{Elementos 1D — Marcos / Vigas}

\section{Frame2DEuler}



\textit{Orden de trabajo. El usuario escribe \textbf{especificación física}, \textbf{formulación numérica} y \textbf{contrato}; la IA rellena \textbf{implementación} y responde en \textbf{diálogo}.}



\section{Especificación física}

\subsection{0. Descripción general}
Elemento 1D inmerso en el plano 2D que modela una \textbf{viga esbelta} según la teoría de \textbf{Euler-Bernoulli}. Dos nodos rígidamente conectados (transmiten momento), 3 grados de libertad por nodo ($u_x, u_y, r_z$). Captura tres acciones internas: \textbf{fuerza axial}, \textbf{cortante transverso} y \textbf{momento flector}. Régimen de \textbf{linealidad geométrica} (pequeños desplazamientos y rotaciones).

\subsection{1. Hipótesis cinemática de Euler-Bernoulli}
\begin{itemize}
  \item Las secciones transversales se mantienen \textbf{planas} tras la deformación.
  \item Las secciones permanecen \textbf{perpendiculares al eje neutro deformado} (cizalladura transversal nula: $\gamma_{xy} = 0$).
\end{itemize}

Consecuencia: la rotación de la sección es igual a la pendiente de la elástica:
\[
\theta(s) = \frac{dv}{ds}.
\]

Válido solo para vigas \textbf{esbeltas} ($L/h \gtrsim 10$). En vigas peraltadas la cizalladura no es despreciable \ensuremath{\to} \texttt{Frame2DTimoshenko}.

\subsection{2. Campos de desplazamiento}
\begin{itemize}
  \item Axial: $u(s)$ a lo largo del eje local $s \in [0, L]$.
  \item Transverso: $v(s)$ perpendicular al eje local.
  \item Rotación de la sección: $\theta(s) = dv/ds$ (derivada de la flecha).
\end{itemize}

\subsection{3. Medidas de deformación}
\begin{itemize}
  \item Axial: $\varepsilon_{axial}(s) = du/ds$ (constante para interpolación lineal en $u$).
  \item Curvatura: $\kappa(s) = d^2v/ds^2$.
\end{itemize}

La deformación axial en un punto genérico de la sección es $\varepsilon(s, y) = \varepsilon_{axial}(s) - y\,\kappa(s)$; al integrar sobre la sección se obtienen axial ($N$) y momento flector ($M$) desacoplados.

\subsection{4. Ecuaciones constitutivas}
\begin{itemize}
  \item Axial: $N = EA\,\varepsilon_{axial}$ (o $N = A\,\sigma$ con $\sigma$ del material).
  \item Flexión: $M = EI\,\kappa$.
\end{itemize}

El material delega la respuesta axial; la rigidez a flexión escala con el \textbf{módulo tangente} $E_t$ devuelto por el material (aproximación: toda la sección entra en régimen no-elástico simultáneamente --- no hay plasticidad distribuida).

\subsection{5. Equilibrio --- forma débil}
Principio de trabajos virtuales. El elemento calcula solo el lado interno:
\[
\delta W_{\text{int}} = \int_0^L (N\,\delta\varepsilon_{axial} + M\,\delta\kappa)\,ds = \delta\mathbf u_e^\top\,\mathbf F_{\text{int}}.
\]


\section{Formulación numérica (FEM)}

\subsection{6. Discretización}
Dos nodos ($s = 0$ y $s = L$), 3 DOFs por nodo en coordenadas globales:
\[
\mathbf u_e = [u_x^{(1)}, u_y^{(1)}, r_z^{(1)}, u_x^{(2)}, u_y^{(2)}, r_z^{(2)}]^\top.
\]

\subsection{7. Sistema local y transformación}
Cosenos directores del eje (configuración inicial, fija): $c = (x_2 - x_1)/L$, $s_\theta = (y_2 - y_1)/L$. Matriz de transformación ortogonal $6\times 6$:
\[
\mathbf T = \begin{bmatrix}
 c & s_\theta & 0 & 0 & 0 & 0 \\
-s_\theta & c & 0 & 0 & 0 & 0 \\
 0 & 0 & 1 & 0 & 0 & 0 \\
 0 & 0 & 0 & c & s_\theta & 0 \\
 0 & 0 & 0 & -s_\theta & c & 0 \\
 0 & 0 & 0 & 0 & 0 & 1
\end{bmatrix}.
\]

Las rotaciones $r_z$ son invariantes ante rotación del sistema (escalares en el plano). $\mathbf u_{\text{local}} = \mathbf T\,\mathbf u_e$.

\subsection{8. Funciones de forma}
\begin{itemize}
  \item Axial: lineal, $N_1^u(s) = 1 - s/L$, $N_2^u(s) = s/L$.
  \item Transverso (Hermite cúbicos, cuatro funciones que interpolan $v$ y $\theta$):
\end{itemize}
\[
N_1^v = 1 - 3\xi^2 + 2\xi^3, \quad N_1^\theta = L(\xi - 2\xi^2 + \xi^3),
\]
\[
N_2^v = 3\xi^2 - 2\xi^3, \quad N_2^\theta = L(-\xi^2 + \xi^3),
\]
con $\xi = s/L$.

\subsection{9. Rigidez local}
En el sistema local, la matriz $6\times 6$ tiene forma cerrada por integración analítica:

\[
\mathbf K_{\text{local}} = \begin{bmatrix}
 \tfrac{EA}{L} & 0 & 0 & -\tfrac{EA}{L} & 0 & 0 \\
 0 & \tfrac{12EI}{L^3} & \tfrac{6EI}{L^2} & 0 & -\tfrac{12EI}{L^3} & \tfrac{6EI}{L^2} \\
 0 & \tfrac{6EI}{L^2} & \tfrac{4EI}{L} & 0 & -\tfrac{6EI}{L^2} & \tfrac{2EI}{L} \\
-\tfrac{EA}{L} & 0 & 0 & \tfrac{EA}{L} & 0 & 0 \\
 0 & -\tfrac{12EI}{L^3} & -\tfrac{6EI}{L^2} & 0 & \tfrac{12EI}{L^3} & -\tfrac{6EI}{L^2} \\
 0 & \tfrac{6EI}{L^2} & \tfrac{2EI}{L} & 0 & -\tfrac{6EI}{L^2} & \tfrac{4EI}{L}
\end{bmatrix}.
\]

En no-linealidad material, $E \to E_t(\varepsilon_{axial})$ devuelto por el material --- escala \textbf{todos} los términos de la matriz por igual. Esta es una aproximación: supone que la plasticidad/daño afecta uniformemente a la sección.

\subsection{10. Fuerzas internas}
En el sistema local:
\[
\mathbf F_{\text{int,local}} = \mathbf K_{\text{local}}\,\mathbf u_{\text{local}}.
\]

La componente axial se \textbf{corrige} con el esfuerzo axial verdadero del material: $F_1 = -\sigma A$, $F_4 = +\sigma A$. Esto permite que el material aporte $\sigma$ no-lineal sin que la formulación de flexión interfiera.

\subsection{11. Ensamblaje global}
\[
\mathbf K_{\text{global}} = \mathbf T^\top\,\mathbf K_{\text{local}}\,\mathbf T, \qquad \mathbf F_{\text{int}} = \mathbf T^\top\,\mathbf F_{\text{int,local}}.
\]

\subsection{12. Cuadratura}
No aplica (forma cerrada para la matriz; integración analítica de los polinomios de Hermite). El campo \texttt{n\_integration\_points: 1} se refiere al único estado material conceptual del elemento (la deformación axial media).


\section{Contrato de implementación}

\begin{lstlisting}[language=yaml]
name: Frame2DEuler
kind: element
status: validated

interface:
  dof_names: [ux, uy, rz]
  n_nodes: 2
  strain_dim: 1
  n_integration_points: 1

parameters:
  - { name: A, type: float, required: true, desc: "Área de la sección transversal" }
  - { name: I, type: float, required: true, desc: "Momento de inercia respecto al eje perpendicular al plano (z)" }

material_contract:
  signature: "compute_state(ε, state) -> (σ, E_tangent, state')"
  strain_kind: "axial escalar (deformación media axial ε = (u_2 - u_1)/L en sistema local)"
  nonlinearity_model: "E_tangent escala toda la matriz local — no hay distinción axial vs flexión en régimen no-elástico"

conventions:
  sign: "σ > 0 ⇔ ε > 0 (elongación) en el eje axial"
  rotation_sign: "r_z positivo antihorario (convención matemática estándar)"
  node_orientation: "el eje local va del nodo 1 al nodo 2; permutar nodos invierte el signo de las fuerzas internas"
  configuration: "inicial fija — L, c, s, T no se actualizan con los desplazamientos"

validity:
  - "vigas esbeltas: L/h ≳ 10 (h = canto de la sección)"
  - pequeños desplazamientos y pequeñas rotaciones
  - "|ε_axial| ≲ 1e-2"

out_of_scope:
  - vigas peraltadas con deformación por cortante significativa (usar Frame2DTimoshenko)
  - grandes rotaciones o desplazamientos (no hay variante corotacional en el catálogo actual)
  - plasticidad distribuida en la sección (fibras); el modelo escala rigidez global por E_t
  - pandeo (requiere rigidez geométrica, no implementada en esta viga lineal)
  - "fuerzas de cuerpo distribuidas: el usuario reparte carga equivalente a los nodos"

acceptance:
  - name: respuesta axial pura
    setup: "viga horizontal, empotrada en extremo izquierdo, carga axial F en extremo derecho"
    expect: "u_x (extremo libre) = F·L/(E·A); momentos y cortantes nulos"
    tol_rel: 1.0e-10

  - name: voladizo con carga transversal (flecha analítica)
    setup: "viga horizontal, empotrada en extremo izquierdo, carga vertical P en extremo derecho"
    expect: "v (extremo libre) = P·L³/(3·E·I); rotación θ = P·L²/(2·E·I)"
    tol_rel: 1.0e-10

  - name: simetría de K
    setup: "cualquier configuración"
    expect: "K_global = K_globalᵀ (elemento conservativo en régimen elástico)"
    tol_abs: 1.0e-12

references:
  - "Bathe K.-J., Finite Element Procedures, §4.2 (formulación isoparamétrica) y cap. 5 (vigas)"
  - "Cook R.D., Concepts and Applications of FEA, §2.7 (elementos de viga)"

\end{lstlisting}


\section{Implementación}

\begin{itemize}
  \item \textbf{Archivo}: solidum/elements/frame/euler.py --- submódulo del paquete \texttt{solidum/elements/frame/} que aloja también \texttt{Frame2DTimoshenko} y \texttt{Frame2DEulerCorot}.
  \item \textbf{Clase}: \texttt{Frame2DEuler} --- hereda directamente de \texttt{Element}. La construcción de la matriz de transformación $\mathbf T$ se delega a \texttt{build\_geometry\_2d} en solidum/elements/frame/\_shared.py, compartida con \texttt{Frame2DTimoshenko} (la versión corotacional reconstruye $\mathbf T$ desde \texttt{alpha0} por su lógica propia).
  \item \textbf{Tests}: tests/test\_frame.py \ensuremath{\cdot} \texttt{TestFrame2DEulerAcceptance} --- los tres criterios físicos y el registro:
  \item \texttt{test\_acceptance\_respuesta\_axial\_pura} (criterio 1) --- resuelve el voladizo con solver, verifica $u_x(L) = FL/(EA)$.
  \item \texttt{test\_acceptance\_voladizo\_carga\_transversal} (criterio 2) --- carga transversal $P$, verifica flecha $v = PL^3/(3EI)$ y rotación $\theta = PL^2/(2EI)$.
  \item \texttt{test\_acceptance\_simetria\_K} (criterio 3) --- viga oblicua (c=0.6, s=0.8) para evitar ejes canónicos, verifica $\mathbf K = \mathbf K^\top$.
  \item \texttt{test\_registro\_en\_registry} --- autodiscover.
  \item \textbf{Notas de traducción}:
  \item La matriz $\mathbf T$ se calcula una sola vez en \texttt{\_\_init\_\_} sobre la configuración inicial y se guarda como atributo --- coherente con el régimen geométricamente lineal.
  \item En \texttt{compute\_element\_state}, la componente axial de $\mathbf F_{\text{int,local}}$ se sobrescribe tras el producto $\mathbf K_{\text{local}}\mathbf u_{\text{local}}$ con el valor $\sigma A$ devuelto por el material. Esto permite usar materiales no-lineales en la rama axial sin que la parte lineal de flexión interfiera.
  \item La aproximación del modelo no-lineal es que $E_t$ escala \textbf{toda} la matriz local: no hay distinción entre rigidez axial y de flexión en régimen no-elástico. Documentado en \texttt{material\_contract.nonlinearity\_model}.
\end{itemize}


\section{Diálogo}

\begin{itemize}
  \item \textbf{2026-04-21} \ensuremath{\cdot} Elemento movido a archivo propio \texttt{solidum/elements/frame.py} y desacoplado del helper compartido \texttt{\_frame\_geometry}. El método estático \texttt{\_build\_geometry} reimplementa la construcción de $\mathbf T$ dentro de la clase. \texttt{Frame2DTimoshenko} sigue en \texttt{structural.py} usando su propio acceso al helper compartido; si mañana también se independiza, duplicará o internalizará su propia versión.
  \item \textbf{2026-04-21} \ensuremath{\cdot} Tests de aceptación cubren la flecha analítica del voladizo a 8 decimales ($v = PL^3/(3EI)$) usando el solver completo del proyecto --- validan no solo la cinemática del elemento sino también su integración con ensamblador y solver.
  \item \textbf{2026-05-13} \ensuremath{\cdot} \texttt{frame.py} se parte en paquete \texttt{solidum/elements/frame/\{euler,timoshenko,euler\_corot,\_shared\}.py}. La duplicación de \texttt{\_build\_geometry} entre Euler y Timoshenko se elimina: la construcción de $\mathbf T$ vuelve a vivir en un único lugar (\texttt{build\_geometry\_2d} en \texttt{\_shared.py}), pero esta vez sin función helper externa flotante --- pertenece al paquete \texttt{frame/}. Frame2DEulerCorot mantiene su propia reconstrucción de $\mathbf T$ desde \texttt{alpha0} porque su cinemática corotacional no se beneficia del helper común.
\end{itemize}


\newpage

\section{Frame2DEulerCorot}



\textit{Orden de trabajo. El usuario escribe \textbf{especificación física}, \textbf{formulación numérica} y \textbf{contrato}; la IA rellena \textbf{implementación} y responde en \textbf{diálogo}.}



\section{Especificación física}

\subsection{0. Descripción general}
Viga 2D esbelta según Euler-Bernoulli en formulación \textbf{corotacional} (Updated Lagrangian). Dos nodos rígidamente conectados, 3 DOFs por nodo ($u_x, u_y, r_z$). Captura grandes desplazamientos y grandes rotaciones del elemento como un todo con \textbf{pequeñas deformaciones axiales} ($|\varepsilon_{axial}| \lesssim 1\%$) y rotaciones nodales deformacionales moderadas.

Abre dominios que la viga lineal no captura: \textbf{pandeo por flexión}, \textbf{post-pandeo}, \textbf{snap-through} de arcos planos, brazos flexibles esbeltos.

\subsection{1. Cinemática corotacional}
Se separa el movimiento del elemento en:

\begin{enumerate}
  \item \textbf{Rotación rígida} $\alpha_e$: el ángulo que el eje de la barra ha girado desde la configuración inicial.
  \item \textbf{Rotación deformacional de cada sección} $\bar\theta_i$: la rotación nodal menos la rotación rígida (es lo que genera momento flector).
\end{enumerate}

Sean:
\[
\alpha_0 = \arctan2(Y_2 - Y_1,\; X_2 - X_1) \quad \text{(ángulo del eje en config. inicial, fijo)},
\]
\[
\alpha = \arctan2(y_2 - y_1,\; x_2 - x_1) \quad \text{(ángulo corriente)},
\]
\[
\alpha_e = \alpha - \alpha_0 \quad \text{(rotación rígida del elemento)}.
\]

Rotaciones deformacionales:
\[
\bar\theta_1 = \theta_1 - \alpha_e, \qquad \bar\theta_2 = \theta_2 - \alpha_e,
\]
donde $\theta_i = r_z^{(i)}$ son las rotaciones totales de los nodos.

\subsection{2. Medidas de deformación}
\begin{itemize}
  \item Axial: $\varepsilon_{axial} = (l - L_0)/L_0$ con $l$ longitud corriente.
  \item Curvatura efectiva: $\kappa = (\bar\theta_2 - \bar\theta_1)/L_0$ (constante si se usa interpolación cúbica).
\end{itemize}

\subsection{3. Ecuaciones constitutivas}
\begin{itemize}
  \item Axial: delegado al material ($\sigma, E_t$).
  \item Flexión: $M = E_t I\,\kappa$ (el mismo $E_t$ del material escala la rigidez de flexión, como en \texttt{Frame2DEuler}).
\end{itemize}

\subsection{4. Equilibrio --- forma débil}
PTV en configuración corriente. El elemento calcula solo el lado interno:
\[
\delta W_{\text{int}} = \delta\mathbf u_e^\top\,\mathbf F_{\text{int}}, \qquad \mathbf K_T = \mathbf K_{\text{material}} + \mathbf K_{\sigma}.
\]


\section{Formulación numérica (FEM)}

\subsection{5. DOFs locales deformacionales}
Se reduce la descripción del estado deformacional del elemento a \textbf{3 magnitudes escalares}:
\[
\mathbf q_l = [\,u_l,\; \bar\theta_1,\; \bar\theta_2\,]^\top, \qquad u_l = l - L_0.
\]

\subsection{6. Rigidez local}
En el sistema de coordenadas corrotado (eje de la barra corriente), la rigidez local $3\times 3$:
\[
\mathbf K_{ll} = \begin{bmatrix}
\tfrac{E_t A}{L_0} & 0 & 0 \\
0 & \tfrac{4 E_t I}{L_0} & \tfrac{2 E_t I}{L_0} \\
0 & \tfrac{2 E_t I}{L_0} & \tfrac{4 E_t I}{L_0}
\end{bmatrix}.
\]

Fuerzas locales $\mathbf F_l = \mathbf K_{ll}\,\mathbf q_l = [N,\; M_1,\; M_2]^\top$, con $N$ sustituido por $\sigma A$ directamente del material para permitir respuestas no-lineales.

\subsection{7. Matriz de transformación $\mathbf B$ (3\ensuremath{\times}6)}
Relación entre variaciones de $\mathbf q_l$ y de los DOFs globales $\delta\mathbf u_e$:
\[
\mathbf B = \begin{bmatrix}
-c & -s & 0 & c & s & 0 \\
-s/l & c/l & 1 & s/l & -c/l & 0 \\
-s/l & c/l & 0 & s/l & -c/l & 1
\end{bmatrix}, \qquad c = \cos\alpha,\; s = \sin\alpha.
\]

$c, s$ y $l$ se evalúan en \textbf{configuración corriente} --- esta es la esencia corotacional.

Fuerzas internas globales:
\[
\mathbf F_{\text{int}} = \mathbf B^\top\,\mathbf F_l.
\]

\subsection{8. Rigidez tangente}
Se descompone en contribución material más geométrica:

\[
\mathbf K_T = \mathbf K_{\text{material}} + \mathbf K_\sigma,
\]
\[
\mathbf K_{\text{material}} = \mathbf B^\top\,\mathbf K_{ll}\,\mathbf B,
\]
\[
\mathbf K_\sigma = \frac{N}{l}\,\mathbf z\,\mathbf z^\top + \frac{M_1 + M_2}{l^2}\,\bigl(\mathbf r\,\mathbf z^\top + \mathbf z\,\mathbf r^\top\bigr),
\]

con vectores geométricos de 6 componentes:
\[
\mathbf r = [-c,\;-s,\;0,\; c,\; s,\;0]^\top \quad \text{(dirección axial)},
\]
\[
\mathbf z = [-s,\; c,\;0,\; s,\;-c,\;0]^\top \quad \text{(perpendicular)}.
\]

La parte $(N/l)\mathbf z\mathbf z^\top$ es análoga al $\mathbf K_G$ de \texttt{Truss2DCorot}: captura la rigidización por tensión (y el ablandamiento precrítico por compresión, base del pandeo). La parte con momentos acopla rotaciones con traslaciones --- crítica para problemas de flexión con desplazamientos significativos.

\subsection{9. Cuadratura}
No aplica (forma cerrada con Hermite cúbicos).


\section{Contrato de implementación}

\begin{lstlisting}[language=yaml]
name: Frame2DEulerCorot
kind: element
status: validated

interface:
  dof_names: [ux, uy, rz]
  n_nodes: 2
  strain_dim: 1
  n_integration_points: 1

parameters:
  - { name: A, type: float, required: true, desc: "Área de la sección transversal" }
  - { name: I, type: float, required: true, desc: "Momento de inercia respecto al eje perpendicular al plano" }

material_contract:
  signature: "compute_state(ε, state) -> (σ, E_tangent, state')"
  strain_kind: "axial escalar corotacional (ε = (l − L₀)/L₀)"
  nonlinearity_model: "E_tangent escala K_material local; el K_σ geométrico depende solo de N, M₁, M₂ corrientes"

conventions:
  sign: "σ > 0 ⇔ ε > 0 (elongación); r_z positivo antihorario"
  node_orientation: "eje local del nodo 1 al nodo 2"
  configuration: "Updated Lagrangian — l, c, s, α se recalculan en cada evaluación"

validity:
  - "|ε_axial| ≲ 1e-2 (pequeña deformación axial)"
  - rotaciones rígidas del elemento de cualquier magnitud, pero acumuladas en |α_e| < π entre commits (rotaciones continuas > π requerirían tracking de vueltas — fuera de alcance)
  - rotaciones deformacionales de las secciones moderadas (|θ̄_i| ≲ 30°); para valores mayores la interpolación Hermite cúbica pierde precisión
  - vigas esbeltas (L/h ≳ 10)

out_of_scope:
  - rotaciones continuas del eje > π sin pasar por commit (pendulums; aliasing en atan2)
  - deformaciones axiales grandes (requiere medida logarítmica)
  - plasticidad distribuida en la sección (fibras)
  - cortante transverso significativo (usar versión Timoshenko, pendiente)
  - "fuerzas de cuerpo distribuidas: el usuario reparte carga a los nodos"

numerical_caveats:
  - "El ángulo α_e se unwrappea a (-π, π] por llamada. Si el problema tiene saltos de rotación > π entre iteraciones consecutivas del solver, la formulación produce resultados incorrectos. Usar pasos de carga finos."
  - "En problemas con pandeo, la rigidez tangente se hace singular al cruzar el punto crítico. Usar `ArcLengthSolver` para seguir la rama post-crítica."

acceptance:
  - name: límite lineal coincide con Frame2DEuler
    setup: "voladizo horizontal, cargas pequeñas (ε ~ 1e-8, v/L ~ 1e-6)"
    expect: "K_T y F_int coinciden con Frame2DEuler lineal a tolerancia rtol=1e-4"
    tol_rel: 1.0e-4

  - name: invariancia bajo rotación rígida
    setup: "viga inicialmente horizontal, rotada rígidamente 45° (ambos nodos)"
    expect: "F_int = 0, σ = 0 (ningún esfuerzo interno bajo rotación rígida pura)"
    tol_abs: 1.0e-10

  - name: coherencia tangente/fuerza interna (chequeo por diferencias finitas)
    setup: "configuración arbitraria no trivial; K_T analítica vs derivada numérica de F_int"
    expect: "‖K_T_analítica − K_T_numérica‖ / ‖K_T_analítica‖ ≲ 1e-5"
    tol_rel: 1.0e-5

  - name: simetría de K
    setup: "configuración arbitraria"
    expect: "K_T = K_Tᵀ"
    tol_abs: 1.0e-6

  - name: Bathe-Bolourchi — cuarto de círculo                 # añadido 2026-05-19
    setup: "voladizo recto, 10 elementos, momento puro M = (π/2)·EI/L en la punta, 8 pasos"
    expect: "tip se mueve a (R·sin(π/2)-L, R·(1-cos(π/2))) con R=EI/M; θ_tip = π/2"
    tol_abs: 0.005   # sobre L (cuerda ~0.2% en 10 elementos rectos)

  - name: Bathe-Bolourchi — medio círculo                     # añadido 2026-05-19
    setup: "voladizo recto, 20 elementos, momento puro M = π·EI/L en la punta, 20 pasos"
    expect: "tip a (-L, 2L/π) y θ_tip = π — la viga se enrosca en semicírculo cerrado"
    tol_abs: 0.01    # sobre L

references:
  - "Crisfield M.A., Non-linear Finite Element Analysis of Solids and Structures, vol.1, §7.3 (corotational beams)"
  - "Belytschko T., Nonlinear Finite Elements for Continua and Structures, §4.11"
  - "Wriggers P., Nonlinear Finite Element Methods, cap. 4"

\end{lstlisting}


\section{Implementación}

\begin{itemize}
  \item \textbf{Archivo}: solidum/elements/frame/euler\_corot.py --- convive con \texttt{Frame2DEuler} y \texttt{Frame2DTimoshenko} en el paquete \texttt{solidum/elements/frame/}, sin heredar de ninguno. Comparte con ellos los helpers libres de \_shared.py (carga de cuerpo, masa consistente, traducción a \texttt{ElementForces}), pero no \texttt{build\_geometry\_2d}: reconstruye \texttt{T} desde \texttt{alpha0} por su lógica corotacional propia.
  \item \textbf{Clase}: \texttt{Frame2DEulerCorot} --- hereda directamente de \texttt{Element}. Métodos internos \texttt{\_current\_geometry} (longitud, cosenos, ángulo corriente y rigid-body $\alpha_e$) y \texttt{\_unwrap} (reducción a $(-\pi, \pi]$).
  \item \textbf{Tests}: tests/test\_frame.py \ensuremath{\cdot} \texttt{TestFrame2DEulerCorotAcceptance} --- 4 criterios + registro:
  \item \texttt{test\_acceptance\_limite\_lineal\_coincide\_con\_frame2deuler}
  \item \texttt{test\_acceptance\_invariancia\_rotacion\_rigida}
  \item \texttt{test\_acceptance\_coherencia\_tangente\_fuerza\_interna} --- \textbf{finite-difference check} de $\mathbf K_T$ contra derivada numérica de $\mathbf F_{\text{int}}$, red de seguridad clave en esta formulación.
  \item \texttt{test\_acceptance\_simetria\_K}
  \item \texttt{test\_registro\_en\_registry}
  \item \textbf{Notas de traducción}:
  \item La formulación se basa en 3 DOFs deformacionales $[u, \bar\theta_1, \bar\theta_2]$ y una matriz $\mathbf B$ $3\times 6$ que los relaciona con los DOFs globales. Evaluada en configuración corriente.
  \item \textbf{Signo de la contribución de momentos en $\mathbf K_\sigma$}: tras derivar $\partial(\mathbf z/l)/\partial\mathbf u_e$ analíticamente se obtiene $-(1/l^2)(\mathbf r\mathbf z^\top + \mathbf z\mathbf r^\top)$. El signo es \textbf{negativo} --- un error común es ponerlo positivo. El test de diferencias finitas lo cazó en la primera implementación.
  \item \texttt{compute\_internal\_forces} proyecta $\mathbf F_{\text{int}}$ al sistema local corriente para reportar axial/cortante separados, utilizando la submatriz $2\times 2$ de rotación del eje corriente.
  \item \texttt{\_unwrap} lleva los ángulos al intervalo $(-\pi, \pi]$; la formulación es válida mientras la rotación acumulada por paso sea menor que $\pi$.
\end{itemize}


\section{Diálogo}

\begin{itemize}
  \item \textbf{2026-04-21} \ensuremath{\cdot} Implementación siguiendo Crisfield §7.3. Convive en \texttt{frame.py} con las vigas lineales por familia temática (todas son vigas 2D), sin herencia ni helpers compartidos.
  \item \textbf{2026-05-13} \ensuremath{\cdot} \texttt{frame.py} se parte en paquete \texttt{solidum/elements/frame/}. EulerCorot vive ahora en \texttt{euler\_corot.py}, sin herencia con las otras vigas 2D. Pasan a compartirse los helpers libres de carga de cuerpo, masa y traducción de fuerzas a través de \texttt{\_shared.py}; la lógica corotacional propia (reconstrucción de \texttt{T} desde \texttt{alpha0}, cinemática actual) se mantiene íntegra dentro de la clase.
  \item \textbf{2026-04-21} \ensuremath{\cdot} La primera versión tenía el signo equivocado en la parte de momentos de $\mathbf K_\sigma$ ($+$ en lugar de $-$). El test de coherencia tangente/fuerza interna por diferencias finitas lo detectó inmediatamente: $\|\mathbf K_{T,\text{analítica}} - \mathbf K_{T,\text{FD}}\|_{\text{rel}}$ del orden de 1e-1 en lugar de 1e-5. Tras derivar $\partial(\mathbf z/l)/\partial\mathbf u_e$ cuidadosamente (ver notas) el signo correcto es negativo. Este episodio justifica retrospectivamente incluir un test FD como criterio de aceptación en toda formulación no-lineal geométrica futura.
  \item \textbf{2026-04-21} \ensuremath{\cdot} El test del límite lineal usa $P=10$ N para obtener $v/L \sim 10^{-4}$: suficientemente pequeño para estar en régimen lineal pero suficientemente grande para que el residuo del solver converja por debajo de \texttt{tol=1e-8} (cargas menores producen desplazamientos del orden del ruido numérico y el Newton oscila en la tolerancia relativa).
\end{itemize}


\newpage

\section{Frame2DTimoshenko}



\textit{Orden de trabajo. El usuario escribe \textbf{especificación física}, \textbf{formulación numérica} y \textbf{contrato}; la IA rellena \textbf{implementación} y responde en \textbf{diálogo}.}



\section{Especificación física}

\subsection{0. Descripción general}
Elemento 1D inmerso en el plano 2D que modela una \textbf{viga} según la teoría de \textbf{Timoshenko}. Dos nodos rígidamente conectados, 3 DOFs por nodo ($u_x, u_y, r_z$). Transmite \textbf{fuerza axial}, \textbf{cortante transverso} y \textbf{momento flector}. Incluye explícitamente la \textbf{deformación por cortante transversal}. Régimen de \textbf{linealidad geométrica} (pequeños desplazamientos y rotaciones).

\subsection{1. Hipótesis cinemática de Timoshenko}
\begin{itemize}
  \item Las secciones transversales se mantienen \textbf{planas} tras la deformación (igual que Euler-Bernoulli).
  \item Las secciones \textbf{no} están obligadas a ser perpendiculares al eje neutro: se permite una \textbf{rotación adicional} respecto a la tangente.
\end{itemize}

Consecuencia: la rotación de la sección $\theta(s)$ y la pendiente de la elástica $dv/ds$ son \textbf{variables independientes}:
\[
\gamma(s) = \frac{dv}{ds} - \theta(s) \neq 0,
\]
donde $\gamma$ es la distorsión angular (deformación por cortante transverso). Euler-Bernoulli corresponde al caso $\gamma = 0$.

Apropiado para vigas \textbf{gruesas, cortas o peraltadas} ($L/h \lesssim 10$) donde la deformación por cortante no es despreciable.

\subsection{2. Campos de desplazamiento}
\begin{itemize}
  \item Axial: $u(s)$ a lo largo del eje local $s \in [0, L]$.
  \item Transverso: $v(s)$ perpendicular al eje local.
  \item Rotación de la sección: $\theta(s)$, \textbf{independiente} de $v$.
\end{itemize}

\subsection{3. Medidas de deformación}
\begin{itemize}
  \item Axial: $\varepsilon_{axial}(s) = du/ds$.
  \item Curvatura (flexión): $\kappa(s) = d\theta/ds$.
  \item Cortante: $\gamma(s) = dv/ds - \theta(s)$.
\end{itemize}

\subsection{4. Ecuaciones constitutivas}
\begin{itemize}
  \item Axial: $N = EA\,\varepsilon_{axial}$ (el material entrega $\sigma$ y $E_t$).
  \item Flexión: $M = EI\,\kappa$.
  \item Cortante: $V = G A_s\,\gamma$, con $G = E/[2(1+\nu)]$ (módulo de corte) y $A_s$ el \textbf{área efectiva de cortante} (menor que $A$ por la distribución no uniforme de la tensión tangencial en la sección).
\end{itemize}

\subsection{5. Factor de Phi (shear locking)}
La formulación estándar con interpolaciones lineales de $v$ y $\theta$ sufre \textbf{shear locking} para vigas esbeltas (la rigidez por cortante domina numéricamente aunque $\gamma$ físicamente sea pequeño). La matriz local incorpora un factor correctivo:
\[
\Phi = \frac{12\,E\,I}{G\,A_s\,L^2}.
\]
Para $\Phi \to 0$ (viga esbelta) la matriz se reduce a la de Euler-Bernoulli.

\subsection{6. Equilibrio --- forma débil}
Principio de trabajos virtuales (PTV). El elemento calcula solo el lado interno:
\[
\delta W_{\text{int}} = \int_0^L (N\,\delta\varepsilon_{axial} + M\,\delta\kappa + V\,\delta\gamma)\,ds = \delta\mathbf u_e^\top\,\mathbf F_{\text{int}}.
\]


\section{Formulación numérica (FEM)}

\subsection{7. Discretización}
Dos nodos ($s = 0$ y $s = L$), 3 DOFs por nodo en coordenadas globales:
\[
\mathbf u_e = [u_x^{(1)}, u_y^{(1)}, r_z^{(1)}, u_x^{(2)}, u_y^{(2)}, r_z^{(2)}]^\top.
\]

\subsection{8. Sistema local y transformación}
Cosenos directores del eje (configuración inicial, fija): $c = (x_2 - x_1)/L$, $s_\theta = (y_2 - y_1)/L$. Matriz de transformación ortogonal $6\times 6$:
\[
\mathbf T = \begin{bmatrix}
 c & s_\theta & 0 & 0 & 0 & 0 \\
-s_\theta & c & 0 & 0 & 0 & 0 \\
 0 & 0 & 1 & 0 & 0 & 0 \\
 0 & 0 & 0 & c & s_\theta & 0 \\
 0 & 0 & 0 & -s_\theta & c & 0 \\
 0 & 0 & 0 & 0 & 0 & 1
\end{bmatrix}, \qquad \mathbf u_{\text{local}} = \mathbf T\,\mathbf u_e.
\]

\subsection{9. Rigidez local con corrección por shear locking}
Coeficientes:
\[
a = \frac{12\,EI}{L^3(1+\Phi)}, \quad b = \frac{6\,EI}{L^2(1+\Phi)}, \quad c_m = \frac{(4+\Phi)\,EI}{L(1+\Phi)}, \quad d = \frac{(2-\Phi)\,EI}{L(1+\Phi)}.
\]

\[
\mathbf K_{\text{local}} = \begin{bmatrix}
 EA/L & 0 & 0 & -EA/L & 0 & 0 \\
 0 & a & b & 0 & -a & b \\
 0 & b & c_m & 0 & -b & d \\
-EA/L & 0 & 0 & EA/L & 0 & 0 \\
 0 & -a & -b & 0 & a & -b \\
 0 & b & d & 0 & -b & c_m
\end{bmatrix}.
\]

Para $\Phi \to 0$: $a \to 12EI/L^3$, $b \to 6EI/L^2$, $c_m \to 4EI/L$, $d \to 2EI/L$ --- se recupera Euler-Bernoulli.

\subsection{10. Fuerzas internas}
\[
\mathbf F_{\text{int,local}} = \mathbf K_{\text{local}}\,\mathbf u_{\text{local}}.
\]

La componente axial se corrige con el esfuerzo axial del material: $F_1 = -\sigma A$, $F_4 = +\sigma A$ (permite materiales no-lineales en la rama axial).

\subsection{11. Ensamblaje global}
\[
\mathbf K_{\text{global}} = \mathbf T^\top\,\mathbf K_{\text{local}}\,\mathbf T, \qquad \mathbf F_{\text{int}} = \mathbf T^\top\,\mathbf F_{\text{int,local}}.
\]

\subsection{12. Cuadratura}
No aplica (forma cerrada).


\section{Contrato de implementación}

\begin{lstlisting}[language=yaml]
name: Frame2DTimoshenko
kind: element
status: validated

interface:
  dof_names: [ux, uy, rz]
  n_nodes: 2
  strain_dim: 1
  n_integration_points: 1

parameters:
  - { name: A, type: float, required: true, desc: "Área de la sección transversal" }
  - { name: I, type: float, required: true, desc: "Momento de inercia respecto al eje perpendicular al plano" }
  - { name: As, type: float, required: true, desc: "Área efectiva de cortante" }
  - { name: nu, type: float, required: false, default: 0.3, desc: "Coeficiente de Poisson (si el material no lo expone)" }

material_contract:
  signature: "compute_state(ε, state) -> (σ, E_tangent, state')"
  strain_kind: "axial escalar (deformación media ε = (u_2 − u_1)/L en sistema local)"
  nonlinearity_model: "E_tangent escala toda la matriz local (axial, flexión y cortante por igual)"
  poisson_source: "si material.nu existe, se usa; en su defecto el parámetro `nu` del elemento con default 0.3 y advertencia por consola"

conventions:
  sign: "σ > 0 ⇔ ε > 0 (elongación) en el eje axial"
  rotation_sign: "r_z positivo antihorario"
  node_orientation: "eje local del nodo 1 al nodo 2"
  configuration: "inicial fija — L, c, s, T, Φ (en forma funcional) se calculan sin actualizar"

validity:
  - "vigas peraltadas o cortas (L/h ≲ 10) donde el cortante no es despreciable"
  - pequeños desplazamientos y pequeñas rotaciones
  - "|ε_axial| ≲ 1e-2"

out_of_scope:
  - grandes rotaciones o desplazamientos (no hay variante corotacional en el catálogo actual)
  - plasticidad distribuida en la sección
  - pandeo
  - "fuerzas de cuerpo distribuidas: el usuario reparte carga equivalente a los nodos"

acceptance:
  - name: convergencia a Euler-Bernoulli para viga esbelta
    setup: "voladizo con carga transversal P, viga esbelta (L/h ≳ 100) ⇒ Φ → 0"
    expect: "v(L) ≈ PL³/(3EI) (solución Euler-Bernoulli) con tolerancia relativa ≲ 1e-3"
    tol_rel: 1.0e-3

  - name: respuesta axial pura desacoplada
    setup: "voladizo con carga axial F"
    expect: "u_x(L) = F·L/(E·A); momentos y cortantes nulos"
    tol_rel: 1.0e-10

  - name: simetría de K
    setup: "viga oblicua"
    expect: "K_global = K_globalᵀ"
    tol_abs: 1.0e-12

references:
  - "Reddy J.N., An Introduction to the Finite Element Method, cap. 5 (vigas de Timoshenko)"
  - "Bathe K.-J., Finite Element Procedures, §5.4 (formulación con factor Φ)"

\end{lstlisting}


\section{Implementación}

\begin{itemize}
  \item \textbf{Archivo}: solidum/elements/frame/timoshenko.py --- submódulo del paquete \texttt{solidum/elements/frame/} que aloja también \texttt{Frame2DEuler} y \texttt{Frame2DEulerCorot}.
  \item \textbf{Clase}: \texttt{Frame2DTimoshenko} --- hereda directamente de \texttt{Element}, sin herencia con las otras vigas 2D. Comparte \texttt{build\_geometry\_2d} con \texttt{Frame2DEuler} en solidum/elements/frame/\_shared.py; además, los helpers libres de carga de cuerpo, masa consistente y traducción a \texttt{ElementForces} también viven en \texttt{\_shared.py}.
  \item \textbf{Tests}: tests/test\_frame.py \ensuremath{\cdot} \texttt{TestFrame2DTimoshenkoAcceptance}:
  \item \texttt{test\_acceptance\_convergencia\_euler\_en\_viga\_esbelta} (criterio 1) --- viga con $L/h$ muy grande, verifica que la flecha tiende a $PL^3/(3EI)$ con tolerancia relativa $10^{-3}$.
  \item \texttt{test\_acceptance\_respuesta\_axial\_pura} (criterio 2) --- carga axial pura, verifica $u_x = FL/(EA)$.
  \item \texttt{test\_acceptance\_simetria\_K} (criterio 3) --- viga oblicua, verifica $\mathbf K = \mathbf K^\top$.
  \item \texttt{test\_registro\_en\_registry} --- autodiscover.
  \item \textbf{Notas de traducción}:
  \item Para no colisionar con la variable cinemática $c$ (coseno director) dentro de \texttt{compute\_element\_state}, el coeficiente $c_m$ de la matriz local se llama \texttt{c\_coef} en el código y \texttt{d\_coef} su análogo (antes eran \texttt{c} y \texttt{d} inline, ambiguos).
  \item El Poisson $\nu$ se toma de \texttt{material.nu} si el material lo expone; en su defecto, del parámetro \texttt{nu} del elemento con default 0.3 y una advertencia por consola la primera vez que se construye con ese default --- atajo a errores silenciosos de usuarios que olvidan especificarlo.
  \item La matriz $\mathbf T$ se calcula una sola vez en \texttt{\_\_init\_\_} y se guarda como atributo (régimen geométricamente lineal).
  \item El factor $\Phi$ se recalcula en cada llamada porque depende de $E_t$ devuelto por el material (en régimen no-lineal $\Phi$ cambia con el estado).
\end{itemize}


\section{Diálogo}

\begin{itemize}
  \item \textbf{2026-04-21} \ensuremath{\cdot} Elemento movido a archivo propio \texttt{solidum/elements/frame.py} y desacoplado del helper \texttt{\_frame\_geometry}. Con este movimiento, el helper compartido se elimina completamente del repositorio: \texttt{Frame2DEuler} y \texttt{Frame2DTimoshenko} replican la construcción de $\mathbf T$ como método estático \texttt{\_build\_geometry}, cada una en su clase. La duplicación (\textasciitilde{}15 líneas) es el precio aceptado de la independencia mutua.
  \item \textbf{2026-04-21} \ensuremath{\cdot} Test del criterio 1: se eligió $L/h$ grande en lugar de reproducir un caso específico de libro porque la convergencia Timoshenko \ensuremath{\to} Euler es el comportamiento \textbf{esperado por construcción} (factor $\Phi \to 0$). Verificarlo directamente con el solver completo valida simultáneamente la cinemática de Timoshenko, el tratamiento del shear locking y la integración con ensamblador + solver.
  \item \textbf{2026-05-13} \ensuremath{\cdot} \texttt{frame.py} se parte en paquete \texttt{solidum/elements/frame/}. La duplicación de \texttt{\_build\_geometry} ya no se justifica: ambas vigas comparten \texttt{build\_geometry\_2d} en \texttt{\_shared.py}, helper interno al paquete (no flotante). Las clases siguen sin herencia entre sí.
\end{itemize}


\newpage

\section{Frame3D}



\textit{Orden de trabajo. El usuario escribe \textbf{especificación física}, \textbf{formulación numérica} y \textbf{contrato}; la IA rellena \textbf{implementación} y responde en \textbf{diálogo}.}



\section{Especificación física}

\subsection{0. Descripción general}
Elemento 1D inmerso en el espacio tridimensional que modela una \textbf{viga esbelta 3D}. Dos nodos rígidamente conectados, 6 DOFs por nodo ($u_x, u_y, u_z, r_x, r_y, r_z$). Captura seis acciones internas:

\begin{itemize}
  \item \textbf{Axial} (fuerza normal $N$ a lo largo del eje local $x$).
  \item \textbf{Cortante en dos planos} ($V_y$ en el plano $xy$, $V_z$ en el plano $xz$).
  \item \textbf{Flexión en dos planos} ($M_z$ alrededor del eje local $z$, $M_y$ alrededor del eje local $y$).
  \item \textbf{Torsión} ($M_x$ alrededor del eje longitudinal).
\end{itemize}

Régimen de \textbf{linealidad geométrica} (pequeños desplazamientos y rotaciones). Hipótesis de Euler-Bernoulli: secciones planas perpendiculares al eje deformado (sin cizalladura transversal).

\subsection{1. Hipótesis}
\begin{itemize}
  \item Secciones planas y perpendiculares al eje neutro deformado (igual que Frame2DEuler).
  \item Flexiones en los dos planos \textbf{desacopladas} entre sí (ejes principales de inercia).
  \item Torsión de Saint-Venant pura (sin alabeo). Válida para secciones macizas o cerradas; abiertas de pared delgada requerirían formulación con alabeo (fuera de alcance).
  \item Material isótropo 1D; $G = E/[2(1+\nu)]$.
\end{itemize}

\subsection{2. Sistema de ejes locales}
Tres ejes ortonormales en cada elemento:
\begin{itemize}
  \item \textbf{$\hat{\mathbf x}$ local}: dirección del eje de la barra ($\hat{\mathbf x} = (\mathbf x_2 - \mathbf x_1)/L$).
  \item \textbf{$\hat{\mathbf y}$ local} y \textbf{$\hat{\mathbf z}$ local}: ejes principales de inercia de la sección, ortogonales al eje.
\end{itemize}

La orientación de $\hat{\mathbf y}, \hat{\mathbf z}$ no está determinada solo por la dirección del eje --- hace falta un \textbf{vector de referencia} proporcionado por el usuario (o un default). Es la diferencia estructural respecto al caso 2D, donde el plano era único.

\subsection{3. Medidas de deformación / acciones internas}
\begin{itemize}
  \item Axial: $\varepsilon_{axial} = du/ds$, tensión $N = EA\,\varepsilon_{axial}$ (o $N = A\,\sigma$ con $\sigma$ del material).
  \item Flexión en plano $xy$: curvatura $\kappa_z = d^2v/ds^2$, momento $M_z = EI_z\,\kappa_z$.
  \item Flexión en plano $xz$: curvatura $\kappa_y = d^2w/ds^2$, momento $M_y = EI_y\,\kappa_y$.
  \item Torsión: $d\theta_x/ds$, momento torsor $M_x = GJ\,d\theta_x/ds$.
\end{itemize}

\subsection{4. Equilibrio --- forma débil}
PTV. El elemento calcula solo el lado interno:
\[
\delta W_{\text{int}} = \int_0^L \!\left(N\,\delta\varepsilon + M_z\,\delta\kappa_z + M_y\,\delta\kappa_y + M_x\,\delta\theta_x'\right)\,ds.
\]


\section{Formulación numérica (FEM)}

\subsection{5. Discretización}
Dos nodos, 6 DOFs por nodo:
\[
\mathbf u_e = [u_x^{(1)}, u_y^{(1)}, u_z^{(1)}, r_x^{(1)}, r_y^{(1)}, r_z^{(1)},\; u_x^{(2)}, u_y^{(2)}, u_z^{(2)}, r_x^{(2)}, r_y^{(2)}, r_z^{(2)}]^\top.
\]

Vector de 12 componentes \ensuremath{\to} matrices $12\times 12$.

\subsection{6. Construcción de ejes locales}
Dados $\mathbf X_1, \mathbf X_2$ y un vector de referencia $\mathbf v_{\text{ref}}$ proporcionado:

\begin{enumerate}
  \item $\hat{\mathbf x} = (\mathbf X_2 - \mathbf X_1)/L$.
  \item Proyectar $\mathbf v_{\text{ref}}$ al plano perpendicular a $\hat{\mathbf x}$:
\end{enumerate}
   \[
\tilde{\mathbf y} = \mathbf v_{\text{ref}} - (\mathbf v_{\text{ref}} \cdot \hat{\mathbf x})\,\hat{\mathbf x}.
\]
\begin{enumerate}
  \item $\hat{\mathbf y} = \tilde{\mathbf y}/\|\tilde{\mathbf y}\|$ (falla si $\|\tilde{\mathbf y}\| \approx 0$: $\mathbf v_{\text{ref}}$ paralelo al eje).
  \item $\hat{\mathbf z} = \hat{\mathbf x} \times \hat{\mathbf y}$.
\end{enumerate}

\textbf{Default de $\mathbf v_{\text{ref}}$} (cuando el usuario no lo proporciona): $[0, 0, 1]$ (eje $z$ global como "vertical"). Si el eje del elemento es cercanamente paralelo a $z$ global (viga vertical), se usa $[1, 0, 0]$ como fallback y se emite una advertencia.

\subsection{7. Matriz de transformación global \ensuremath{\to} local}
\[
\boldsymbol{\lambda} = \begin{bmatrix} \hat{\mathbf x}^\top \\ \hat{\mathbf y}^\top \\ \hat{\mathbf z}^\top \end{bmatrix} \quad (3\times 3,\ \text{ortogonal}).
\]

\[
\mathbf T = \text{bloques diag}(\boldsymbol{\lambda}, \boldsymbol{\lambda}, \boldsymbol{\lambda}, \boldsymbol{\lambda}) \quad (12\times 12).
\]

Aplicada a los 4 bloques de 3 DOFs: desplazamientos nodo 1, rotaciones nodo 1, desplazamientos nodo 2, rotaciones nodo 2. Las rotaciones infinitesimales en 3D transforman como vectores.

\[
\mathbf u_{\text{local}} = \mathbf T\,\mathbf u_e.
\]

\subsection{8. Matriz de rigidez local}
Desacoplada en cuatro contribuciones: axial, torsión, flexión $xy$ (plano con $u_y$ y $r_z$, inercia $I_z$), flexión $xz$ (plano con $u_z$ y $r_y$, inercia $I_y$).

Definiciones:
\[
a_z = \tfrac{12EI_z}{L^3},\ b_z = \tfrac{6EI_z}{L^2},\ c_z = \tfrac{4EI_z}{L},\ d_z = \tfrac{2EI_z}{L},
\]
\[
a_y = \tfrac{12EI_y}{L^3},\ b_y = \tfrac{6EI_y}{L^2},\ c_y = \tfrac{4EI_y}{L},\ d_y = \tfrac{2EI_y}{L}.
\]

La rigidez $\mathbf K_{\text{local}}$ es simétrica, con los bloques desacoplados (0s fuera de su plano correspondiente). Los signos del bloque de flexión $xz$ se invierten respecto al $xy$ porque $u_z$ y $r_y$ forman una pareja dextrógira opuesta (consecuencia del producto vectorial $\hat{\mathbf x}\times\hat{\mathbf y}=\hat{\mathbf z}$).

Forma concreta en términos de los 12 DOFs locales (la matriz completa vive en la implementación).

\subsection{9. Fuerzas internas}
\[
\mathbf F_{\text{int,local}} = \mathbf K_{\text{local}}\,\mathbf u_{\text{local}},
\]
con la componente axial corregida por el $\sigma$ que devuelve el material: $F_1 = -\sigma A$, $F_7 = +\sigma A$. El resto de componentes viene del producto lineal (permite materiales no-lineales en la rama axial).

\subsection{10. Ensamblaje global}
\[
\mathbf K_{\text{global}} = \mathbf T^\top\,\mathbf K_{\text{local}}\,\mathbf T, \qquad \mathbf F_{\text{int}} = \mathbf T^\top\,\mathbf F_{\text{int,local}}.
\]

\subsection{11. Cuadratura}
No aplica (forma cerrada; integración analítica de los polinomios de Hermite).


\section{Contrato de implementación}

\begin{lstlisting}[language=yaml]
name: Frame3D
kind: element
status: validated

interface:
  dof_names: [ux, uy, uz, rx, ry, rz]
  n_nodes: 2
  strain_dim: 1
  n_integration_points: 1

parameters:
  - { name: A,  type: float, required: true, desc: "Área de la sección transversal" }
  - { name: Iy, type: float, required: true, desc: "Momento de inercia alrededor del eje local y" }
  - { name: Iz, type: float, required: true, desc: "Momento de inercia alrededor del eje local z" }
  - { name: J,  type: float, required: true, desc: "Constante torsional de Saint-Venant" }
  - { name: nu, type: float, required: false, default: 0.3, desc: "Coeficiente de Poisson (si el material no lo expone)" }
  - { name: ref_vector, type: list, required: false, default: "[0, 0, 1]",
      desc: "Vector 3D de referencia que fija la orientación de los ejes locales y, z de la sección. Default: eje z global." }

material_contract:
  signature: "compute_state(ε, state) -> (σ, E_tangent, state')"
  strain_kind: "axial escalar (ε = (u_2 − u_1)/L en sistema local)"
  nonlinearity_model: "E_tangent escala toda la matriz local; G se recalcula como E_t/[2(1+ν)]"
  poisson_source: "material.nu si existe; en su defecto parámetro `nu` del elemento"

conventions:
  sign: "σ > 0 ⇔ ε > 0 (elongación); momentos siguen regla de la mano derecha respecto a los ejes locales"
  local_axes: "x_local = eje de la barra; y_local, z_local = ejes principales de inercia definidos por ref_vector"
  node_orientation: "eje local del nodo 1 al nodo 2"
  configuration: "inicial fija — L, T, ejes locales no se actualizan"

validity:
  - "vigas esbeltas (L/h ≳ 10 en ambos planos de flexión)"
  - "pequeños desplazamientos y pequeñas rotaciones en el espacio"
  - "|ε_axial| ≲ 1e-2"
  - "sección con ejes principales de inercia bien definidos por ref_vector"

out_of_scope:
  - grandes rotaciones espaciales (requeriría Frame3DCorot; pendiente)
  - torsión con alabeo restringido (secciones abiertas de pared delgada)
  - plasticidad distribuida en la sección (fibras)
  - pandeo por bifurcación
  - deformación por cortante transversal significativa (Timoshenko 3D, fuera de alcance)
  - "fuerzas de cuerpo distribuidas: el usuario reparte carga a los nodos"

numerical_caveats:
  - "Si ref_vector es paralelo al eje de la barra, la proyección al plano perpendicular se anula y el constructor aborta con mensaje claro — el usuario debe elegir otro vector."
  - "Para barras verticales (eje casi paralelo a z global) con ref_vector default, el constructor usa [1, 0, 0] como fallback y emite advertencia por consola."

acceptance:
  - name: respuesta axial pura
    setup: "voladizo alineado con eje x, carga F axial en extremo libre"
    expect: "u_x = F·L/(E·A); demás componentes nulas"
    tol_rel: 1.0e-10

  - name: flexión en plano xy (coherencia con Frame2D)
    setup: "voladizo alineado con eje x, carga P en dirección y, ref_vector = [0,0,1]"
    expect: "v_y = P·L³/(3·E·Iz), θ_z = P·L²/(2·E·Iz)"
    tol_rel: 1.0e-10

  - name: flexión en plano xz
    setup: "voladizo alineado con eje x, carga P en dirección z, ref_vector = [0,1,0]"
    expect: "v_z = P·L³/(3·E·Iy), θ_y = -P·L²/(2·E·Iy) (signo por mano derecha)"
    tol_rel: 1.0e-10

  - name: torsión pura
    setup: "voladizo con momento T aplicado alrededor del eje de la barra en el extremo libre"
    expect: "θ_x = T·L/(G·J); demás componentes nulas"
    tol_rel: 1.0e-10

  - name: simetría de K
    setup: "viga en orientación arbitraria con ref_vector genérico"
    expect: "K_global = K_globalᵀ"
    tol_abs: 1.0e-8

references:
  - "Przemieniecki J.S., Theory of Matrix Structural Analysis, Tabla 11.3"
  - "Cook R.D., Concepts and Applications of FEA, §2.8"
  - "Bathe K.-J., Finite Element Procedures, cap. 5"

\end{lstlisting}


\section{Implementación}

\begin{itemize}
  \item \textbf{Archivo}: solidum/elements/frame3d.py --- archivo propio, sin dependencia de ningún otro elemento.
  \item \textbf{Clase}: \texttt{Frame3D} --- hereda directamente de \texttt{Element}. Métodos privados \texttt{\_build\_local\_frame} (longitud, matriz $\boldsymbol{\lambda}$ 3\ensuremath{\times}3 y $\mathbf T$ 12\ensuremath{\times}12) y \texttt{\_build\_local\_stiffness} (matriz $\mathbf K_{\text{local}}$ 12\ensuremath{\times}12 con bloques axial, torsión y flexión en ambos planos).
  \item \textbf{Tests}: tests/test\_frame3d.py \ensuremath{\cdot} \texttt{TestFrame3DAcceptance} --- cubre los 5 criterios + validación de \texttt{ref\_vector} paralelo al eje + registro:
  \item \texttt{test\_acceptance\_respuesta\_axial\_pura}
  \item \texttt{test\_acceptance\_flexion\_xy\_coincide\_con\_frame2d}
  \item \texttt{test\_acceptance\_flexion\_xz}
  \item \texttt{test\_acceptance\_torsion\_pura}
  \item \texttt{test\_acceptance\_simetria\_K}
  \item \texttt{test\_rechazo\_ref\_vector\_paralelo}
  \item \texttt{test\_registro\_en\_registry}
  \item \textbf{Notas de traducción}:
  \item La matriz $\boldsymbol{\lambda}$ 3\ensuremath{\times}3 se construye como $[\hat{\mathbf x}^\top; \hat{\mathbf y}^\top; \hat{\mathbf z}^\top]$ --- las filas son los ejes locales expresados en coordenadas globales. $\mathbf T$ 12\ensuremath{\times}12 es un \texttt{blkdiag} de cuatro copias de $\boldsymbol{\lambda}$ (traslaciones y rotaciones transforman como vectores en 3D infinitesimal).
  \item El default de \texttt{ref\_vector} es \texttt{[0, 0, 1]}; el fallback \texttt{[1, 0, 0]} se activa cuando $|\hat{\mathbf x} \cdot \hat{\mathbf z}_{\text{global}}| > 0.99$ (barra cercanamente vertical). Ambas rutas emiten advertencia por consola la primera vez que se usan con default.
  \item La matriz $\mathbf K_{\text{local}}$ se construye rellenando entradas no nulas en sus 4 bloques desacoplados (axial, torsión, flexión $xy$, flexión $xz$). Los signos de la flexión $xz$ están invertidos respecto a la flexión $xy$ porque la pareja $(u_z, r_y)$ tiene orientación dextrógira opuesta a $(u_y, r_z)$: este punto se verifica explícitamente en \texttt{test\_acceptance\_flexion\_xz}.
  \item \texttt{compute\_internal\_forces} devuelve cortantes, torsión y momentos en ambos extremos separados --- útil para post-proceso estructural (diagramas de fuerzas internas).
\end{itemize}


\section{Diálogo}

\begin{itemize}
  \item \textbf{2026-04-21} \ensuremath{\cdot} Elemento creado como componente totalmente autónomo: no hereda de \texttt{Frame2DEuler} ni de ninguna otra viga. La maquinaria cinemática 3D se implementa íntegra dentro de \texttt{frame3d.py}, coherente con la filosofía aplicada a cables.
  \item \textbf{2026-04-21} \ensuremath{\cdot} Decisión sobre \texttt{ref\_vector} vs "tercer nodo de orientación". Se optó por \texttt{ref\_vector} (vector 3D proyectado al plano perpendicular al eje) por ser más compacto en YAML y no requerir nodos fantasma en la malla. El default \texttt{[0, 0, 1]} cubre la mayoría de casos de ingeniería estructural (estructuras terrestres con gravedad en $-z$); el fallback para barras verticales usa \texttt{[1, 0, 0]} con advertencia explícita.
  \item \textbf{2026-04-21} \ensuremath{\cdot} Verificación cruzada con \texttt{Frame2DEuler}: el criterio 2 aplica una viga alineada con el eje $x$ global bajo carga en $y$; los desplazamientos resultantes deben coincidir con la solución 2D clásica $v = PL^3/(3EI_z)$, $\theta = PL^2/(2EI_z)$. El test lo verifica a 8 decimales. Esto blinda la coherencia dimensional: el 3D se reduce al 2D sin error acumulado cuando el problema es efectivamente plano.
\end{itemize}


\newpage

\chapter{Elementos 2D — Sólidos}

\section{Tri3}



\textit{Documentación retroactiva del elemento ya implementado.}



\section{Especificación física}

\subsection{0. Descripción general}
Elemento sólido bidimensional plano de primer orden, triangular de tres nodos. Aproximación lineal $C^0$ del campo de desplazamientos en coordenadas naturales $(\xi, \eta)$. Conocido en la literatura como elemento \textit{constant strain triangle} (CST): la matriz $\mathbf{B}$ es constante sobre el elemento, por lo que la deformación $\boldsymbol{\varepsilon}$ y la tensión $\boldsymbol{\sigma}$ son uniformes dentro de cada Tri3. Régimen de pequeños desplazamientos y deformaciones.

\subsection{1. Cinemática de desplazamientos}
\[
u_x(\xi, \eta) = \sum_{i=1}^{3} N_i(\xi, \eta)\, u_x^{(i)}, \qquad u_y(\xi, \eta) = \sum_{i=1}^{3} N_i(\xi, \eta)\, u_y^{(i)}.
\]

\subsection{2. Cinemática de deformaciones}
Notación Voigt 2D $[\varepsilon_{xx}, \varepsilon_{yy}, \gamma_{xy}]$, idéntica a la de \texttt{Quad4}. Por la linealidad de $N_i$, las derivadas $\partial N_i / \partial x$ son constantes y $\boldsymbol{\varepsilon}$ es uniforme en el elemento.

\subsection{3. Ecuación constitutiva}
Delegada al material 2D (\texttt{STRAIN\_DIM = 3}). Soporta \texttt{Elastic2D}, \texttt{IsotropicDamage2D}, \texttt{VonMises2D} (este último solo \texttt{plane\_strain}).

\subsection{4. Equilibrio --- forma fuerte}
\[
-\nabla \cdot \boldsymbol{\sigma} = \mathbf{b} \quad \text{en } \Omega, \quad \boldsymbol{\sigma} \cdot \mathbf{n} = \bar{\mathbf{t}} \text{ en } \partial\Omega_N, \quad \mathbf{u} = \bar{\mathbf{u}} \text{ en } \partial\Omega_D.
\]

\subsection{5. Equilibrio --- forma débil}
Idéntica a la del \texttt{Quad4}; el principio variacional no depende de la forma del elemento.


\section{Formulación numérica (FEM)}

\subsection{6. Discretización}
Tres nodos en orden antihorario sobre el plano $(x, y)$. Mapeo isoparamétrico desde el triángulo de referencia con vértices en $(0, 0)$, $(1, 0)$, $(0, 1)$. Jacobiano $\mathbf{J} = \partial(x, y) / \partial(\xi, \eta)$ es constante sobre el elemento; el código aborta con error si $\det \mathbf{J} \le \text{tol}$ (nodos colineales o invertidos).

\subsection{7. Funciones de forma}
Coordenadas baricéntricas:
\[
N_1 = 1 - \xi - \eta, \qquad N_2 = \xi, \qquad N_3 = \eta.
\]
Sus derivadas en coordenadas naturales son constantes:
\[
\partial N_1 / \partial \xi = -1, \quad \partial N_2 / \partial \xi = 1, \quad \partial N_3 / \partial \xi = 0;
\]
\[
\partial N_1 / \partial \eta = -1, \quad \partial N_2 / \partial \eta = 0, \quad \partial N_3 / \partial \eta = 1.
\]

\subsection{8. Matriz deformación-desplazamiento}
$\mathbf{B}$ tiene tamaño $3 \times 6$ y, una vez transformadas las derivadas a globales vía $\mathbf{J}^{-1}$, sus entradas son \textbf{constantes} sobre el elemento:
\[
\mathbf{B}_i = \begin{bmatrix} \partial N_i / \partial x & 0 \\ 0 & \partial N_i / \partial y \\ \partial N_i / \partial y & \partial N_i / \partial x \end{bmatrix}.
\]

\subsection{9. Rigidez elemental}
\[
\mathbf{K}_e = \mathbf{B}^\top\, \mathbf{D}\, \mathbf{B}\, A_e\, t,
\]
con $A_e = \tfrac{1}{2}\, \det \mathbf{J}$ el área del triángulo y $t$ el espesor. La integración es exacta con un único punto central porque $\mathbf{B}$ es constante.

\subsection{10. Fuerzas internas}
\[
\mathbf{F}_{\text{int}} = \mathbf{B}^\top\, \boldsymbol{\sigma}\, A_e\, t.
\]

\subsection{11. Cuadratura}
Un punto central (en el baricentro) con peso $w = 1/2$ sobre el triángulo de referencia. Es exacto para formas constantes de $\mathbf{B}$ y $\boldsymbol{\sigma}$.

\subsection{12. Cargas distribuidas consistentes}

\textbf{Fuerza de cuerpo $\mathbf b$}. Con $N_i$ lineales y $\mathbf b$ uniforme la integral es exacta:
\[
\mathbf f_e^{b} = \int_{\Omega_e} \mathbf N^\top\, \mathbf b\, t\, dA \;\Longrightarrow\; \text{cada nodo recibe } \tfrac{1}{3}\, \mathbf b\, A_e\, t.
\]

\textbf{Tracción de superficie $\bar{\mathbf t}$} sobre un borde recto. Con $\bar{\mathbf t}$ constante el reparto es $\tfrac{L}{2}\, \bar{\mathbf t}\, t$ en cada uno de los dos nodos del borde. Bordes numerados según la conectividad:

\begin{center}
\begin{tabularx}{\textwidth}{YY}
\hline
\textbf{Edge} & \textbf{Nodos} \\
\hline
0 & (0, 1) \\
1 & (1, 2) \\
2 & (2, 0) \\
\hline
\end{tabularx}
\end{center}
$\bar{\mathbf t}$ se especifica en coordenadas globales $(t_x, t_y)$. Tracción variable o presión normal no soportadas en este paso.

\subsection{13. Salida por punto de Gauss}

\texttt{compute\_gauss\_state(U)} devuelve la misma estructura que en \texttt{Quad4} pero con un único punto (centroide). Para Tri3 \ensuremath{\varepsilon} y \ensuremath{\sigma} son constantes sobre el elemento, así que el dato del punto coincide con el promedio.


\section{Limitación crítica --- \textit{shear locking}}

El Tri3 es notoriamente rígido en flexión: tres DOFs de desplazamiento no bastan para representar el modo de flexión puro de una viga discretizada en triángulos, por lo que el elemento responde con cortante espurio (\textit{shear locking}) y subestima sistemáticamente la deflexión. La severidad crece con la relación $L/h$ del problema. Recomendación operativa: usar \texttt{Quad4} por defecto; reservar \texttt{Tri3} para zonas de transición de malla en regiones donde el campo de tensiones varía suavemente.


\section{Contrato de implementación}

\begin{lstlisting}[language=yaml]
name: Tri3
kind: element
status: validated

interface:
  dof_names: [ux, uy]
  n_nodes: 3
  strain_dim: 3                 # Voigt 2D [ε_xx, ε_yy, γ_xy]
  n_integration_points: 1

parameters:
  - { name: thickness, type: float, required: false, default: 1.0, desc: Espesor para estado plano }

material_contract:
  signature: "compute_state(ε, state) -> (σ, C_tangent, state')"
  strain_kind: "Voigt 2D [ε_xx, ε_yy, γ_xy]"

conventions:
  sign: "tensión positiva = tracción"
  voigt: "[xx, yy, xy] con γ_xy = 2·ε_xy"
  node_orientation: "antihorario (asegura det J > 0)"

validity:
  - "pequeños desplazamientos y deformaciones"
  - "campos de tensión que varían suavemente sobre el elemento"

out_of_scope:
  - "flexión dominante (shear locking severo)"
  - "elementos de alto orden (Tri6)"
  - "grandes deformaciones"

acceptance:
  verification:
    - name: patch_test_macneal_harder
      setup: "4 Tri3 distorsionados alrededor de un nodo interior descentrado
             con campo lineal u = 1e-3·(x + y/2), v = 1e-3·(y + x/2)
             impuesto en los 4 nodos del contorno del cuadrado unitario"
      expect: "nodo interior adopta el campo lineal y ε = (1e-3, 1e-3,
              1e-3) en cada elemento; σ uniforme"
      tol_rel: 1.0e-10
  specific:
    - name: parche de deformación constante sobre malla triangular regular
      setup: "malla triangular con campo de desplazamiento lineal y material Elastic2D"
      expect: "ε constante (CST reproduce campos lineales por construcción)"
      tol_rel: 1.0e-12
    - name: jacobiano degenerado abortado
      setup: "tres nodos colineales"
      expect: "ValueError en _compute_kinematics_tri3"
    - name: cargas consistentes — body load uniforme
      setup: "Tri3 con b uniforme"
      expect: "cada nodo recibe (1/3)·b·A_e·t; Σ = b·A_e·t"
      tol_rel: 1.0e-12
    - name: cargas consistentes — tracción uniforme en un borde
      setup: "Tri3 con tracción constante en un borde"
      expect: "Σf = t̄·L·t y reparto L/2 a cada nodo del borde, 0 en el opuesto"
      tol_rel: 1.0e-12

references:
  - "Cook, Malkus, Plesha, Witt, Concepts and Applications of FEA, §3.6 (CST)"
  - "Zienkiewicz O.C., The Finite Element Method, vol. 1, §5 (limitations of CST)"

\end{lstlisting}


\section{Implementación}

\begin{itemize}
  \item \textbf{Archivo}: solidum/elements/solid\_2d/tri3.py
  \item \textbf{Clase}: \texttt{Tri3} (registrada vía \texttt{@ElementRegistry.register})
  \item \textbf{Función núcleo}: \texttt{\_compute\_kinematics\_tri3(coords)}, decorada con \texttt{@njit}, en \_shared.py. Reutiliza \texttt{\_compute\_integrands} (también en \texttt{\_shared}, compartido con Quad4) con peso $w = 0{.}5$.
  \item \textbf{Tests}:
  \item tests/test\_solid\_2d.py --- patch test sobre malla triangular regular y modos de cuerpo rígido.
  \item tests/test\_patch\_solid\_2d.py --- patch test de MacNeal-Harder con nodo interior libre (NAFEMS, V\&V).
\end{itemize}


\section{Diálogo}

\begin{itemize}
  \item \textbf{2026-04-30} \ensuremath{\cdot} Spec retroactiva. \texttt{Tri3} precedía al protocolo spec-first; la spec se creó documentando la formulación ya implementada. Promovido \texttt{status: validated}.
  \item \textbf{2026-04-30} \ensuremath{\cdot} Añadido patch test de MacNeal-Harder (NAFEMS) en \texttt{acceptance.verification}. Cuatro triángulos alrededor de un nodo interior descentrado reproducen un campo lineal impuesto en el contorno con \ensuremath{\varepsilon} constante en cada elemento.
  \item \textbf{2026-05-04} \ensuremath{\cdot} Expuesta \texttt{compute\_gauss\_state(U)} (1 punto central). El \texttt{VtkExporter} consume \ensuremath{\sigma} por elemento y promedia a nodos (\texttt{Sigma\_*\_nodal}).
  \item \textbf{2026-05-04} \ensuremath{\cdot} Añadidas cargas distribuidas consistentes (\texttt{compute\_body\_load}, \texttt{compute\_edge\_traction}) con la misma semántica que \texttt{Quad4}. Para Tri3 ambos integrandos son exactos analíticamente: body load reparte 1/3 por nodo, tracción de borde reparte L/2 a cada uno de los dos nodos del borde.
\end{itemize}


\newpage

\section{Quad4}



\textit{Documentación retroactiva del elemento ya implementado. Cubre la formulación, el contrato de implementación y la trazabilidad con tests existentes.}



\section{Especificación física}

\subsection{0. Descripción general}
Elemento sólido bidimensional plano de primer orden, cuadrilátero de cuatro nodos. Aproximación bilineal $C^0$ del campo de desplazamientos en coordenadas naturales $(\xi, \eta) \in [-1, 1]^2$. Formulación isoparamétrica: la geometría y los desplazamientos comparten las mismas funciones de forma. Régimen de pequeños desplazamientos y deformaciones.

\subsection{1. Cinemática de desplazamientos}
\[
u_x(\xi, \eta) = \sum_{i=1}^{4} N_i(\xi, \eta)\, u_x^{(i)}, \qquad u_y(\xi, \eta) = \sum_{i=1}^{4} N_i(\xi, \eta)\, u_y^{(i)}.
\]

\subsection{2. Cinemática de deformaciones}
Tensor infinitesimal en notación Voigt 2D $[\varepsilon_{xx}, \varepsilon_{yy}, \gamma_{xy}]$:
\[
\varepsilon_{xx} = \frac{\partial u_x}{\partial x}, \quad \varepsilon_{yy} = \frac{\partial u_y}{\partial y}, \quad \gamma_{xy} = \frac{\partial u_x}{\partial y} + \frac{\partial u_y}{\partial x}.
\]

\subsection{3. Ecuación constitutiva}
Delegada al material 2D (\texttt{STRAIN\_DIM = 3}). El elemento llama \texttt{material.compute\_state(\ensuremath{\varepsilon}, state) \ensuremath{\to} (\ensuremath{\sigma}, C\_tangent, state')} por punto de Gauss y soporta cualquier ley constitutiva 2D (\texttt{Elastic2D}, \texttt{IsotropicDamage2D}, \texttt{VonMises2D}).

\subsection{4. Equilibrio --- forma fuerte}
\[
-\nabla \cdot \boldsymbol{\sigma} = \mathbf{b} \quad \text{en } \Omega, \qquad \boldsymbol{\sigma} \cdot \mathbf{n} = \bar{\mathbf{t}} \quad \text{en } \partial\Omega_N, \qquad \mathbf{u} = \bar{\mathbf{u}} \quad \text{en } \partial\Omega_D.
\]

\subsection{5. Equilibrio --- forma débil}
\[
\int_\Omega \delta \boldsymbol{\varepsilon}^\top \boldsymbol{\sigma}\, dV = \int_\Omega \delta \mathbf{u}^\top \mathbf{b}\, dV + \int_{\partial\Omega_N} \delta \mathbf{u}^\top \bar{\mathbf{t}}\, dS, \qquad \forall\, \delta \mathbf{u} \in V_0.
\]


\section{Formulación numérica (FEM)}

\subsection{6. Discretización}
Cuatro nodos en orden antihorario sobre el plano $(x, y)$. Mapeo isoparamétrico desde el cuadrado de referencia $[-1, 1]^2$:
\[
x(\xi, \eta) = \sum_{i=1}^{4} N_i(\xi, \eta)\, x^{(i)}, \qquad y(\xi, \eta) = \sum_{i=1}^{4} N_i(\xi, \eta)\, y^{(i)}.
\]
Jacobiano $\mathbf{J} = \partial(x, y) / \partial(\xi, \eta)$ y su determinante $\det \mathbf{J}$ entran en el cambio de variable de la integración. El código aborta con error si $\det \mathbf{J} \le \text{tol}$ (elemento degenerado o nodos en orden invertido).

\subsection{7. Funciones de forma}
Producto tensorial bilineal:
\[
N_i(\xi, \eta) = \tfrac{1}{4}(1 + \xi_i\,\xi)(1 + \eta_i\,\eta), \qquad (\xi_i, \eta_i) = (\pm 1, \pm 1).
\]
Convención de numeración (antihorario): $(\xi_1, \eta_1) = (-1, -1)$, $(\xi_2, \eta_2) = (+1, -1)$, $(\xi_3, \eta_3) = (+1, +1)$, $(\xi_4, \eta_4) = (-1, +1)$.

\subsection{8. Matriz deformación-desplazamiento}
Las derivadas en globales se obtienen vía $\partial N_i / \partial x = \mathbf{J}^{-1}\, \partial N_i / \partial \xi$. La matriz $\mathbf{B}$ de tamaño $3 \times 8$ se ensambla por nodo:
\[
\mathbf{B}_i = \begin{bmatrix} \partial N_i / \partial x & 0 \\ 0 & \partial N_i / \partial y \\ \partial N_i / \partial y & \partial N_i / \partial x \end{bmatrix}, \qquad \boldsymbol{\varepsilon} = \mathbf{B}\, \mathbf{u}_e.
\]

\subsection{9. Rigidez elemental}
\[
\mathbf{K}_e = \int_\Omega \mathbf{B}^\top\, \mathbf{D}\, \mathbf{B}\, t\, dA = \sum_{p=1}^{n_g} \mathbf{B}(\xi_p, \eta_p)^\top\, \mathbf{D}_p\, \mathbf{B}(\xi_p, \eta_p)\, \det \mathbf{J}_p\, w_p\, t,
\]
donde $t$ es el espesor (\texttt{thickness}) y $\mathbf{D}_p$ es el tensor constitutivo tangente devuelto por el material en cada punto.

\subsection{10. Fuerzas internas}
\[
\mathbf{F}_{\text{int}} = \sum_{p=1}^{n_g} \mathbf{B}(\xi_p, \eta_p)^\top\, \boldsymbol{\sigma}_p\, \det \mathbf{J}_p\, w_p\, t.
\]

\subsection{11. Cuadratura}
Por defecto Gauss--Legendre $2 \times 2$ (cuatro puntos): orden de exactitud 3 --- integra exactamente $\mathbf{K}_e$ para Jacobiano constante (geometría de paralelogramo) y captura el comportamiento dominante en geometrías irregulares moderadas. Soporta esquemas alternativos vía el parámetro \texttt{quadrature} (lista nombrada en \texttt{QuadratureRegistry}).

\textbf{Aviso sobre integración reducida}: el esquema $1 \times 1$ está disponible pero el código emite advertencia explícita; un único punto central no detecta los modos de hourglass (energía nula con desplazamientos de signo alternado), que pueden contaminar la respuesta sin estabilización adicional. No se implementa estabilización tipo Flanagan-Belytschko.

\subsection{12. Cargas distribuidas consistentes}

\textbf{Fuerza de cuerpo $\mathbf b$} (e.g. peso propio $\mathbf b = (0, -\rho g)$, fuerza centrífuga). Vector nodal equivalente:
\[
\mathbf f_e^{b} = \int_{\Omega_e} \mathbf N^\top\, \mathbf b\, t\, dA = \sum_{p=1}^{n_g} \mathbf N(\xi_p, \eta_p)^\top\, \mathbf b\, \det\mathbf J_p\, w_p\, t,
\]
donde $\mathbf N$ es la matriz $2 \times 8$ de funciones de forma. Se reutiliza la cuadratura del elemento (Gauss $2 \times 2$ por defecto) --- exacta para $\mathbf b$ uniforme y geometrías de paralelogramo, y adecuada para geometrías irregulares moderadas.

\textbf{Tracción de superficie $\bar{\mathbf t}$} sobre un borde $\Gamma_e \subset \partial\Omega_N$:
\[
\mathbf f_e^{t} = \int_{\Gamma_e} \mathbf N^\top\, \bar{\mathbf t}\, t\, dS.
\]
Cada borde es una recta entre dos nodos; sobre él las funciones de forma son lineales y, para $\bar{\mathbf t}$ constante, la integral se reduce a un reparto exacto de $\tfrac{L}{2}\, \bar{\mathbf t}\, t$ a cada uno de los dos nodos del borde, donde $L$ es la longitud del segmento. Bordes numerados según la conectividad nodal:

\begin{center}
\begin{tabularx}{\textwidth}{YY}
\hline
\textbf{Edge} & \textbf{Nodos} \\
\hline
0 & (0, 1) \\
1 & (1, 2) \\
2 & (2, 3) \\
3 & (3, 0) \\
\hline
\end{tabularx}
\end{center}
$\bar{\mathbf t}$ se especifica en \textbf{coordenadas globales} $(t_x, t_y)$; presiones normales se obtienen multiplicando previamente por la normal exterior del borde. Tracción variable a lo largo del borde no está soportada en este paso.

\subsection{13. Salida por punto de Gauss}

\texttt{compute\_gauss\_state(U)} devuelve un dict con $\boldsymbol\varepsilon$ y $\boldsymbol\sigma$ en cada uno de los $n_g$ puntos de Gauss del elemento, junto con sus coordenadas naturales y globales. Habilita post-proceso fino (mapas no promediados, suavizado nodal, extrapolación tipo Barlow). \texttt{compute\_internal\_forces} queda como atajo de promedio.


\section{Contrato de implementación}

\begin{lstlisting}[language=yaml]
name: Quad4
kind: element
status: validated

interface:
  dof_names: [ux, uy]
  n_nodes: 4
  strain_dim: 3                 # Voigt 2D [ε_xx, ε_yy, γ_xy]
  n_integration_points: 4       # default Gauss 2×2

parameters:
  - { name: thickness, type: float, required: false, default: 1.0, desc: Espesor para estado plano }
  - { name: quadrature, type: tuple, required: false, default: "2x2", desc: Regla de cuadratura desde QuadratureRegistry }

material_contract:
  signature: "compute_state(ε, state) -> (σ, C_tangent, state')"
  strain_kind: "Voigt 2D [ε_xx, ε_yy, γ_xy]"

conventions:
  sign: "tensión positiva = tracción"
  voigt: "[xx, yy, xy] con γ_xy = 2·ε_xy (engineering shear strain)"
  node_orientation: "antihorario (asegura det J > 0)"

validity:
  - "pequeños desplazamientos y deformaciones"
  - "geometría sin distorsión severa (det J > tol en todos los puntos de Gauss)"
  - "compatible con materiales Elastic2D, IsotropicDamage2D, VonMises2D (este último solo plane_strain)"

out_of_scope:
  - "grandes deformaciones, formulaciones corotacionales o totales lagrangianas"
  - "estabilización contra hourglass para integración reducida"
  - "elementos de alto orden (Q8, Q9)"

acceptance:
  verification:
    - name: patch_test_macneal_harder
      setup: "5 Quad4 distorsionados (4 nodos exteriores + 4 interiores)
             con campo lineal u = 1e-3·(x + y/2), v = 1e-3·(y + x/2)
             impuesto en los 4 nodos del contorno"
      expect: "nodos interiores adoptan el campo lineal y ε = (1e-3, 1e-3,
              1e-3) en todos los puntos de Gauss; σ uniforme"
      tol_rel: 1.0e-10
  specific:
    - name: parche de tracción uniaxial sobre malla regular
      setup: "malla regular con campo de desplazamiento lineal y material Elastic2D"
      expect: "ε_xx constante en todos los puntos de Gauss"
      tol_rel: 1.0e-12
    - name: jacobiano degenerado abortado
      setup: "elemento con nodos colineales o en orden inverso"
      expect: "ValueError en _compute_kinematics"
    - name: cargas consistentes — body load uniforme
      setup: "Quad4 regular y distorsionado con b uniforme; sumar componentes nodales"
      expect: "Σf_x = b_x·A·t y Σf_y = b_y·A·t (invariante de geometría)"
      tol_rel: 1.0e-10
    - name: cargas consistentes — tracción uniforme en un borde
      setup: "Quad4 con tracción constante en un borde"
      expect: "Σf = t̄·L·t y reparto L/2 a cada nodo del borde, 0 en los demás"
      tol_rel: 1.0e-12
    - name: patch físico de tracción uniaxial
      setup: "Quad4 cuadrado L×H, borde izquierdo con ux=0 (n0,n3) + uy=0 (n0); borde derecho con tracción uniforme (p,0)"
      expect: "u reproduce ux=p·x/E, uy=-ν·p·y/E (plane stress); σ_xx=p, σ_yy=σ_xy=0"
      tol_rel: 1.0e-10

references:
  - "Bathe K.-J., Finite Element Procedures, §5.3 (isoparametric elements)"
  - "Cook, Malkus, Plesha, Witt, Concepts and Applications of FEA, §6 (plane stress/strain)"

\end{lstlisting}


\section{Implementación}

\begin{itemize}
  \item \textbf{Archivo}: solidum/elements/solid\_2d/quad4.py
  \item \textbf{Clase}: \texttt{Quad4} (registrada vía \texttt{@ElementRegistry.register})
  \item \textbf{Funciones núcleo}: \texttt{\_compute\_kinematics(xi, eta, coords)} y \texttt{\_compute\_integrands(B, C, \ensuremath{\sigma}, detJ, w, t)}, ambas decoradas con \texttt{@njit} (Numba) para evaluar $\mathbf{B}$, $\det \mathbf{J}$ y los integrandos a velocidad cercana a Fortran. Viven en solidum/elements/solid\_2d/\_shared.py y se comparten con Tri3.
  \item \textbf{Tests}:
  \item tests/test\_solid\_2d.py --- ensamblaje de $\mathbf{B}$, simetría de $\mathbf{K}_e$, modos de cuerpo rígido y tracción uniaxial sobre malla regular.
  \item tests/test\_patch\_solid\_2d.py --- patch test de MacNeal-Harder sobre 5 Quad4 distorsionados (NAFEMS, V\&V).
\end{itemize}


\section{Diálogo}

\begin{itemize}
  \item \textbf{2026-04-30} \ensuremath{\cdot} Spec retroactiva. \texttt{Quad4} precedía al protocolo spec-first; la spec se creó documentando la formulación ya implementada y verificada. Promovido \texttt{status: validated}.
  \item \textbf{2026-04-30} \ensuremath{\cdot} Añadido patch test de MacNeal-Harder (NAFEMS) en \texttt{acceptance.verification}. Cinco Quad4 distorsionados con cuatro nodos interiores libres reproducen exactamente un campo lineal impuesto en el contorno, con \ensuremath{\varepsilon} constante e igual al gradiente analítico en todos los puntos de Gauss.
  \item \textbf{2026-05-04} \ensuremath{\cdot} Expuesta \texttt{compute\_gauss\_state(U)} con \ensuremath{\sigma} y \ensuremath{\varepsilon} por punto de Gauss y coordenadas naturales/globales. \texttt{compute\_internal\_forces} queda como promedio derivado. El \texttt{VtkExporter} añade campos nodales \texttt{Sigma\_XX\_nodal}, \texttt{Sigma\_YY\_nodal}, \texttt{Tau\_XY\_nodal} y \texttt{Von\_Mises\_nodal} por promedio simple sobre los elementos contiguos a cada nodo (suavizado de orden 0).
  \item \textbf{2026-05-04} \ensuremath{\cdot} Añadidas cargas distribuidas consistentes (\texttt{compute\_body\_load}, \texttt{compute\_edge\_traction}). Para tracciones uniformes en bordes rectos el reparto es exacto (L/2 a cada nodo del borde) y se evita la cuadratura 1D; las fuerzas de cuerpo se integran con la cuadratura del elemento. Solo tracciones \textbf{constantes por borde} y en \textbf{coordenadas globales} en este paso; tracción variable y presión normal quedan para una iteración posterior.
\end{itemize}


\newpage

\section{Tri6}



\textit{Triángulo isoparamétrico de 6 nodos. Cura el shear locking severo del \texttt{Tri3} y reproduce campos cuadráticos exactamente. Útil cuando la malla es triangular por geometría.}



\section{Especificación física}

Idéntica al \texttt{Tri3} salvo por la cinemática y la cuadratura.

\subsection{7. Funciones de forma}
Coordenadas baricéntricas $L_1 = 1 - \xi - \eta$, $L_2 = \xi$, $L_3 = \eta$:
\[
N_i^{vert} = L_i(2 L_i - 1), \quad i = 1, 2, 3
\]
\[
N_4 = 4 L_1 L_2, \quad N_5 = 4 L_2 L_3, \quad N_6 = 4 L_3 L_1.
\]

Numeración: 0,1,2 vértices en (0,0), (1,0), (0,1); 3 medio del borde 0-1, 4 medio del 1-2, 5 medio del 2-0.

\subsection{11. Cuadratura}
Cuadratura triangular de 3 puntos (puntos medios de los lados; peso 1/6 cada uno) --- exacta para polinomios hasta grado 2 sobre el triángulo de referencia.

\subsection{12. Cargas distribuidas}
\begin{itemize}
  \item Body load con la cuadratura del elemento.
  \item Edge traction uniforme: reparto 1/6, 4/6, 1/6 a (vértice, medio, vértice).
\end{itemize}


\section{Contrato de implementación}

\begin{lstlisting}[language=yaml]
name: Tri6
kind: element
status: validated

interface:
  dof_names: [ux, uy]
  n_nodes: 6
  strain_dim: 3
  n_integration_points: 3

acceptance:
  verification:
    - name: patch_cuadratico
      setup: "u = (x², 0) impuesto en los 6 nodos"
      expect: "ε_xx = 2x exacto en cada Gauss"
      tol_rel: 1.0e-10
    - name: cooks_membrane_4x4
      setup: "Trapezoide de Cook triangulado 4×4 (32 Tri6, diagonal c1→c3 con center node compartido)"
      expect: "u_y(48, 52) ≈ 23.91 (Bathe; Hughes)"
      tol_abs: 1.5
    - name: cooks_membrane_refinamiento
      setup: "Mallas 2×2 → 4×4 → 6×6 con la misma triangulación"
      expect: "Error |u_y − 23.91| decreciente monótonamente"
      tol_rel: 0.0

\end{lstlisting}


\section{Implementación}

\begin{itemize}
  \item \textbf{Archivo}: solidum/elements/solid\_2d/tri6.py \ensuremath{\cdot} clase \texttt{Tri6} (subclase de la base interna \texttt{\_HigherOrderSolid2D} en \_shared.py, compartida con Quad8 y Quad9). Declara \texttt{\_MASS\_QUADRATURE = "tri\_6"} (Dunavant 6 puntos, orden 4) porque la cuadratura del elemento (\texttt{tri\_3}, orden 2) subintegra el producto cuadrático\ensuremath{\times}cuadrático de la masa consistente.
  \item \textbf{Tests}: tests/test\_higher\_order\_solid\_2d.py (patch cuadrático), tests/test\_cooks\_membrane.py (benchmark Bathe/Hughes con malla triangulada 4\ensuremath{\times}4 + refinamiento monótono, añadido Fase B 2026-05-19).
\end{itemize}


\newpage

\section{Quad8}



\textit{Elemento sólido 2D de orden cuadrático con interpolación serendípita.}

\textit{Reproduce exactamente campos hasta $Q_2^{seren}$ (todos los polinomios de grado total \ensuremath{\le} 2 más algunos de orden 3 sin el término $\xi^2\eta^2$). Mucho más preciso que \texttt{Quad4} en flexión y geometrías curvas con el mismo grado de malla.}



\section{Especificación física}

\subsection{0. Descripción general}
Cuadrilátero plano de 8 nodos: 4 vértices + 4 nodos en los puntos medios de los bordes. Pequeñas deformaciones, isoparamétrico.

\subsection{1. Cinemática de desplazamientos}
\[
u_x(\xi, \eta) = \sum_{i=1}^{8} N_i(\xi, \eta)\, u_x^{(i)}, \qquad u_y(\xi, \eta) = \sum_{i=1}^{8} N_i(\xi, \eta)\, u_y^{(i)}.
\]

\subsection{2. Cinemática de deformaciones}
Voigt 2D $[\varepsilon_{xx}, \varepsilon_{yy}, \gamma_{xy}]$, idéntica al \texttt{Quad4}.

\subsection{3. Ecuación constitutiva}
Material 2D (\texttt{STRAIN\_DIM = 3}).

\subsection{4-5. Equilibrio}
Idénticas a \texttt{Quad4}.


\section{Formulación numérica}

\subsection{6. Discretización}
Nodos en orden estándar: 0-3 vértices antihorarios desde $(-1, -1)$; 4 medio del borde 0-1, 5 medio del 1-2, 6 medio del 2-3, 7 medio del 3-0. Mapeo isoparamétrico.

\subsection{7. Funciones de forma serendípitas}
Vértices ($i = 0, 1, 2, 3$ con $(\xi_i, \eta_i) \in \{\pm 1\}^2$):
\[
N_i = \tfrac{1}{4}(1 + \xi_i\xi)(1 + \eta_i\eta)(\xi_i\xi + \eta_i\eta - 1).
\]

Medios de borde:
\begin{itemize}
  \item Bordes con $\xi = 0$ (nodos 4, 6): $N_i = \tfrac{1}{2}(1 - \xi^2)(1 + \eta_i\eta)$ con $\eta_i = \mp 1$.
  \item Bordes con $\eta = 0$ (nodos 5, 7): $N_i = \tfrac{1}{2}(1 + \xi_i\xi)(1 - \eta^2)$ con $\xi_i = \pm 1$.
\end{itemize}

\subsection{8. Matriz $\mathbf B$}
Ensamblaje estándar a partir de $\partial N_i / \partial x = \mathbf J^{-1} \partial N_i / \partial \xi$. Tamaño $3 \times 16$.

\subsection{9-10. Rigidez y fuerzas internas}
\[
\mathbf K_e = \int \mathbf B^\top \mathbf D\, \mathbf B\, t\, dA, \qquad \mathbf F_{\text{int}}^e = \int \mathbf B^\top \boldsymbol\sigma\, t\, dA.
\]
Cuadratura Gauss $3 \times 3$ por defecto.

\subsection{11. Cuadratura}
Gauss--Legendre $3 \times 3$ (9 puntos) --- exacta para polinomios hasta grado 5, suficiente para integrar $\mathbf B^\top \mathbf D \mathbf B$ del Quad8 sobre Jacobiano constante. La opción \texttt{2x2} queda disponible (reducida) pero introduce modos espurios; emite advertencia.

\subsection{12. Cargas distribuidas consistentes}
\begin{itemize}
  \item \textbf{Body load}: $\int N^\top \mathbf b\, t\, dA$ con la cuadratura del elemento (3\ensuremath{\times}3).
  \item \textbf{Edge traction} uniforme: en un borde recto con N cuadráticas la integral analítica reparte $L/6$ a cada nodo extremo y $4L/6$ al nodo medio (regla 1-4-1, Simpson). Bordes:
\end{itemize}

\begin{center}
\begin{tabularx}{\textwidth}{YYY}
\hline
\textbf{Edge} & \textbf{Nodos extremos} & \textbf{Nodo medio} \\
\hline
0 & (0, 1) & 4 \\
1 & (1, 2) & 5 \\
2 & (2, 3) & 6 \\
3 & (3, 0) & 7 \\
\hline
\end{tabularx}
\end{center}
\subsection{13. Salida por punto de Gauss}
\texttt{compute\_gauss\_state(U)} con la misma estructura que \texttt{Quad4}: 9 puntos por defecto.


\section{Contrato de implementación}

\begin{lstlisting}[language=yaml]
name: Quad8
kind: element
status: validated

interface:
  dof_names: [ux, uy]
  n_nodes: 8
  strain_dim: 3
  n_integration_points: 9       # default Gauss 3×3

parameters:
  - { name: thickness, type: float, required: false, default: 1.0 }
  - { name: quadrature, type: tuple, required: false, default: "3x3" }

material_contract:
  signature: "compute_state(ε, state) -> (σ, C_tangent, state')"
  strain_kind: "Voigt 2D [ε_xx, ε_yy, γ_xy]"

conventions:
  sign: "tensión positiva = tracción"
  voigt: "[xx, yy, xy] con γ_xy = 2·ε_xy"
  node_orientation: "vértices antihorarios; medios en el orden de los bordes"

validity:
  - "pequeños desplazamientos y deformaciones"

out_of_scope:
  - "grandes deformaciones"
  - "estabilización contra modos espurios bajo integración reducida"

acceptance:
  verification:
    - name: patch_lineal
      setup: "campo u = a₀ + a₁·x + a₂·y impuesto en el contorno"
      expect: "ε constante en todos los Gauss e igual al gradiente analítico"
      tol_rel: 1.0e-10
    - name: patch_cuadratico
      setup: "campo u_x = x²·1e-3 impuesto en los 8 nodos"
      expect: "ε_xx(x, y) = 2e-3·x exacto en todos los puntos de Gauss"
      tol_rel: 1.0e-10
  specific:
    - name: cargas consistentes — body load uniforme
      setup: "Σf = b·A·t (regular y distorsionado)"
      tol_rel: 1.0e-10
    - name: cargas consistentes — tracción uniforme en un borde
      setup: "Σf = t̄·L·t y reparto 1/6, 4/6, 1/6 en (vértice, medio, vértice)"
      tol_rel: 1.0e-12
    - name: Cook's membrane (Cook 1974)                              # 2026-05-19
      setup: "trapezoid (0,0)-(48,44)-(48,60)-(0,44), E=1, ν=1/3, plane stress, cortante total F=1 distribuido en borde derecho; malla 4×4 Q8"
      expect: "u_y en (48,52) ≈ 23.91 ± 0.5; refinamiento 2×2→4×4→6×6 reduce error monotónicamente"
      tol_abs: 0.5

references:
  - "Bathe K.-J., Finite Element Procedures, §5.3"
  - "Cook, Malkus, Plesha, Witt, Concepts and Applications of FEA, §6"
  - "Zienkiewicz O.C., The Finite Element Method, vol. 1, §8"

\end{lstlisting}


\section{Implementación}

\begin{itemize}
  \item \textbf{Archivo}: solidum/elements/solid\_2d/quad8.py \ensuremath{\cdot} clase \texttt{Quad8} (subclase de la base interna \texttt{\_HigherOrderSolid2D} en \_shared.py, compartida con Quad9 y Tri6).
  \item \textbf{Tests}: tests/test\_higher\_order\_solid\_2d.py.
\end{itemize}


\newpage

\section{Quad9}



\textit{Variante Lagrangiana del Quad8: añade un noveno nodo central interior.}

\textit{Espacio polinómico completo $Q_2$ (incluye el término $\xi^2\eta^2$ que el serendípito omite).}



\section{Especificación física}

Idéntica a \texttt{Quad8} salvo por las funciones de forma.

\subsection{7. Funciones de forma}
Producto tensorial de los polinomios de Lagrange 1D de orden 2:
\[
L_0(\xi) = \tfrac{1}{2}\xi(\xi - 1), \quad L_1(\xi) = 1 - \xi^2, \quad L_2(\xi) = \tfrac{1}{2}\xi(\xi + 1).
\]
$N_i(\xi, \eta) = L_a(\xi)\,L_b(\eta)$ con $(a, b)$ asociado al nodo $i$ según la numeración:

\begin{itemize}
  \item 0..3: vértices (-1,-1), (+1,-1), (+1,+1), (-1,+1).
  \item 4: medio del borde 0-1 (\ensuremath{\eta} = -1).
  \item 5: medio del borde 1-2 (\ensuremath{\xi} = +1).
  \item 6: medio del borde 2-3 (\ensuremath{\eta} = +1).
  \item 7: medio del borde 3-0 (\ensuremath{\xi} = -1).
  \item 8: nodo central interior (\ensuremath{\xi}, \ensuremath{\eta}) = (0, 0) --- \textit{bubble}.
\end{itemize}

\subsection{11. Cuadratura}
Gauss $3 \times 3$ por defecto.

\subsection{12. Cargas distribuidas}
Body load con la cuadratura del elemento. Tracción en bordes idéntica a \texttt{Quad8} (1/6, 4/6, 1/6 --- el nodo 8 no participa, no toca bordes).


\section{Contrato de implementación}

\begin{lstlisting}[language=yaml]
name: Quad9
kind: element
status: validated

interface:
  dof_names: [ux, uy]
  n_nodes: 9
  strain_dim: 3
  n_integration_points: 9

acceptance:
  verification:
    - name: patch_cuadratico
      setup: "u = (x², 0) impuesto en los 9 nodos"
      expect: "ε_xx = 2x exacto en todos los Gauss"
      tol_rel: 1.0e-10
    - name: Cook's membrane (Cook 1974)                              # 2026-05-19
      setup: "trapezoid (0,0)-(48,44)-(48,60)-(0,44), E=1, ν=1/3, plane stress, cortante total F=1 distribuido en borde derecho; malla 4×4 Q9"
      expect: "u_y en (48,52) ≈ 23.91 ± 0.5"
      tol_abs: 0.5

\end{lstlisting}


\section{Implementación}

\begin{itemize}
  \item \textbf{Archivo}: solidum/elements/solid\_2d/quad9.py \ensuremath{\cdot} clase \texttt{Quad9} (subclase de la base interna \texttt{\_HigherOrderSolid2D} en \_shared.py, compartida con Quad8 y Tri6).
  \item \textbf{Tests}: tests/test\_higher\_order\_solid\_2d.py.
\end{itemize}


\newpage

\chapter{Elementos 2D — Discontinuidades embebidas}

\section{CST\_Embedded2D}



\textit{Orden de trabajo. Pre-redactada por la IA como borrador para validación del usuario (autor de la formulación, Retama 2010). La cinemática KOS y la fórmula \texttt{l\_d = (A/h)\ensuremath{\cdot}cos(\ensuremath{\theta}−\ensuremath{\alpha})} son traducción directa de los Caps. 2, 5, 6 y 7 de la tesis; la mecánica de implementación (Voigt, registry, signatura, contrato \texttt{Element}, condensación) sigue las convenciones del proyecto. El usuario revisa, corrige donde proceda, y aprueba antes de que la IA toque código.}



\section{Especificación física}

\subsection{0. Descripción general}

Elemento finito sólido 2D triangular de tres nodos (CST padre) \textbf{enriquecido} con una \textbf{discontinuidad interior embebida} \texttt{Γ\_d} que cruza el dominio del elemento cuando se cumple un criterio de activación. Es la primera implementación de la familia de elementos con DOFs enriquecidos elementales introducida por \textbf{ADR 0010}, y la \textbf{fase 2} de la hoja de ruta de discontinuidades embebidas.

Estructura del elemento:

\begin{itemize}
  \item \textbf{Estado intacto} (\texttt{discontinuity\_state is None}): se comporta exactamente como un \texttt{Tri3} estándar --- CST con un único punto de Gauss, B constante, integración exacta del bulk con peso \texttt{A\_e\ensuremath{\cdot}t}.
  \item \textbf{Estado agrietado} (\texttt{discontinuity\_state is not None}): el campo de desplazamientos del bulk se enriquece con un salto \texttt{[[u]] \ensuremath{\in} \ensuremath{\mathbb{R}}\ensuremath{^{2}}} localizado en \texttt{Γ\_d}. El bulk vecino descarga elásticamente; la disipación se concentra en \texttt{Γ\_d} y la gobierna un material cohesivo (\texttt{CohesiveMaterial}, ADR 0010 §1).
\end{itemize}

Aproximación adoptada: \textbf{discrete approach} (Retama 2010), \textbf{KOS} (Kinematically Optimal Symmetric) según la clasificación de Jirásek (2000). Tracking trivial: la normal \texttt{n} y la posición de \texttt{Γ\_d} se fijan al activar el elemento (perpendicular a \texttt{\ensuremath{\sigma}\_I} por criterio Rankine; \texttt{Γ\_d} pasa por el centroide) y no se reorientan después. Modo-I puro (consistente con \texttt{CohesiveDamageIsotropic} fase 1); modo mixto y KSON quedan diferidos (fases G y H del ADR 0010).

Régimen: pequeños desplazamientos y deformaciones, cuasi-estático. Pensado para hormigón, roca, cerámicos cuasi-frágiles. Bulk \textbf{elástico} por construcción de la formulación (la disipación está localizada en \texttt{Γ\_d}); materiales no lineales del bulk quedan \texttt{out\_of\_scope} en fase 1.

\subsection{1. Cinemática de desplazamientos (formulación KOS)}

Campo total de desplazamientos en el elemento:

\[
\mathbf u(\mathbf x) = \underbrace{\sum_{i=1}^{3} N_i(\mathbf x)\,\mathbf d_i}_{\mathbf u^{\text{std}}(\mathbf x)} \;+\; \underbrace{\bigl(H_{\Gamma_d}(\mathbf x) - \varphi(\mathbf x)\bigr)\,\mathbf G \,\llbracket\mathbf u\rrbracket}_{\mathbf u^{\text{enr}}(\mathbf x)}
\]

donde:

\begin{itemize}
  \item \texttt{N\_i(x)} son las funciones de forma estándar del CST (coordenadas baricéntricas; ver \texttt{Tri3} §7).
  \item \texttt{d\_i \ensuremath{\in} \ensuremath{\mathbb{R}}\ensuremath{^{2}}} son los desplazamientos nodales estándar (DOFs globales del modelo, 2 por nodo).
  \item \texttt{[[u]] \ensuremath{\in} \ensuremath{\mathbb{R}}\ensuremath{^{2}}} es el salto de desplazamientos a través de \texttt{Γ\_d}, expresado en el \textbf{frame local} \texttt{(n, s)} de la discontinuidad. Es un \textbf{DOF enriquecido elemental} --- vive en el elemento, no se ensambla al sistema global (ADR 0010 §2).
  \item \texttt{H\_\{Γ\_d\}(x)} es la función escalón de Heaviside relativa a \texttt{Γ\_d} (0 en \texttt{\ensuremath{\Omega}\ensuremath{^{-}}}, 1 en \texttt{\ensuremath{\Omega}\ensuremath{^{+}}}). \texttt{\ensuremath{\Omega}\ensuremath{^{+}}} es el subdominio del lado al que apunta \texttt{n}.
  \item \texttt{\ensuremath{\varphi}(x)} es una función de soporte continua que vale \texttt{1} en los nodos de \texttt{\ensuremath{\Omega}\ensuremath{^{+}}} y \texttt{0} en los de \texttt{\ensuremath{\Omega}\ensuremath{^{-}}}. Para CST con \texttt{Γ\_d} cortando dos lados, \texttt{\ensuremath{\varphi}(x)} es la combinación lineal \texttt{\ensuremath{\Sigma} N\_i\textasciicircum{}+(x)} sobre los nodos del lado positivo (típicamente un único nodo solitario, ver §6).
  \item \texttt{G} es la matriz que transforma el salto del frame local al sistema cartesiano global:
\end{itemize}

\[
\mathbf G = \begin{bmatrix} n_x & s_x \\ n_y & s_y \end{bmatrix}
\]

La construcción \texttt{(H\_\{Γ\_d\} - \ensuremath{\varphi})} garantiza que \texttt{u\textasciicircum{}enr} sea \textbf{continuo en el contorno del elemento} (porque \texttt{H\_\{Γ\_d\} = \ensuremath{\varphi}} en los nodos) y que el salto correcto \texttt{[[u]] = G\textasciicircum{}T\ensuremath{\cdot}[[u]]\_global} ocurra únicamente sobre \texttt{Γ\_d} interior. Es la elección \textbf{kinematically optimal}: el campo discontinuo está bien localizado y el salto no se transmite a elementos vecinos (a diferencia de XFEM, donde el salto sí cruza fronteras inter-elemento).

\subsection{2. Cinemática de deformaciones}

Deformación en el bulk (parte regular del campo):

\[
\boldsymbol\varepsilon^{\text{bulk}}(\mathbf x) = \mathbf B_{\text{std}}\,\mathbf d \;-\; \mathbf B^\varphi(\mathbf x)\,\mathbf G\,\llbracket\mathbf u\rrbracket
\]

donde \texttt{B\_std} es la matriz \texttt{B} del CST estándar (Voigt 2D, 3\ensuremath{\times}6, constante, ver \texttt{Tri3} §8) y \texttt{B\textasciicircum{}\ensuremath{\varphi}(x)} es la matriz Voigt construida con \texttt{\ensuremath{\nabla}\ensuremath{\varphi}}:

\[
\mathbf B^\varphi = \begin{bmatrix} \partial\varphi/\partial x & 0 \\ 0 & \partial\varphi/\partial y \\ \partial\varphi/\partial y & \partial\varphi/\partial x \end{bmatrix}
\]

Como \texttt{\ensuremath{\varphi} = \ensuremath{\Sigma} N\_i\textasciicircum{}+}, las derivadas \texttt{\ensuremath{\partial}\ensuremath{\varphi}/\ensuremath{\partial}x}, \texttt{\ensuremath{\partial}\ensuremath{\varphi}/\ensuremath{\partial}y} son \textbf{constantes sobre el elemento} (CST lineal). Por tanto \texttt{B\textasciicircum{}\ensuremath{\varphi}} es constante.

Salto cinemático en \texttt{Γ\_d} (en frame local):

\[
\llbracket\mathbf u(\mathbf x_{\Gamma_d})\rrbracket = \mathbf G^\top\,\bigl[\,\mathbf u(\mathbf x^+) - \mathbf u(\mathbf x^-)\,\bigr] \;=\; \llbracket\mathbf u\rrbracket \qquad (\text{i.e. los propios DOFs enriquecidos})
\]

por construcción de la cinemática KOS --- el salto en \texttt{Γ\_d} coincide exactamente con la variable enriquecida, sin contribución de \texttt{d} (continuidad del campo estándar atravesando \texttt{Γ\_d}).

\subsection{3. Ecuación constitutiva}

\textbf{Bulk} (en todo el dominio del elemento \texttt{\ensuremath{\Omega}\_e}):
\[
\boldsymbol\sigma(\mathbf x) = \mathbf C_e \,\boldsymbol\varepsilon^{\text{bulk}}(\mathbf x)
\]
con \texttt{C\_e} la matriz constitutiva elástica isótropa (parámetros \texttt{E}, \texttt{\ensuremath{\nu}}, hipótesis plane stress / plane strain del material \texttt{Elastic2D}).

\textbf{Discontinuidad} (sobre \texttt{Γ\_d}):
\[
\mathbf t(\mathbf x_{\Gamma_d}) = \mathbf T_{\text{coh}}\bigl(\llbracket\mathbf u\rrbracket, \text{estado cohesivo}\bigr)
\]
con \texttt{T\_coh} el operador del material cohesivo (en frame local; ver spec \texttt{CohesiveDamageIsotropic}).

\subsection{4. Equilibrio --- forma fuerte}

Balance estático en cada subdominio (\texttt{\ensuremath{\Omega}\ensuremath{^{-}}}, \texttt{\ensuremath{\Omega}\ensuremath{^{+}}}) más continuidad de tracciones a través de \texttt{Γ\_d}:

\[
\nabla\cdot\boldsymbol\sigma + \mathbf b = \mathbf 0 \quad \text{en } \Omega \setminus \Gamma_d, \qquad \boldsymbol\sigma\cdot\mathbf n = \mathbf t \quad \text{en } \Gamma_d.
\]

\subsection{5. Equilibrio --- forma débil (residuales acoplados)}

Aplicando el principio de trabajos virtuales con variaciones independientes \texttt{\ensuremath{\delta}d} y \texttt{\ensuremath{\delta}[[u]]}, se obtienen \textbf{dos residuales acoplados}:

\[
\mathbf R^d = \int_{\Omega_e} \mathbf B_{\text{std}}^\top\,\boldsymbol\sigma\,d\Omega \;-\; \mathbf f^{\text{ext}}_d = \mathbf 0
\]

\[
\mathbf R^{\llbracket u\rrbracket} = -\int_{\Omega_e} \mathbf G^\top\,(\mathbf B^\varphi)^\top\,\boldsymbol\sigma\,d\Omega \;+\; l_d\,\mathbf t = \mathbf 0
\]

con \texttt{l\_d} la longitud efectiva de \texttt{Γ\_d} dentro del elemento (ver §11). El primer residual se ensambla al sistema global; el segundo se resuelve \textbf{localmente} dentro del elemento (ver §10, condensación estática).


\section{Formulación numérica (FEM)}

\subsection{6. Discretización del CST padre}

Idéntica a \texttt{Tri3}: tres nodos en orden antihorario, mapeo isoparamétrico desde el triángulo de referencia con vértices \texttt{(0,0)}, \texttt{(1,0)}, \texttt{(0,1)}. Jacobiano constante; aborta con error si \texttt{det J \ensuremath{\le} tol}.

\textbf{Identificación del nodo solitario.} Cuando \texttt{Γ\_d} corta dos lados del triángulo, exactamente un nodo queda aislado en \texttt{\ensuremath{\Omega}\ensuremath{^{+}}} (lado al que apunta \texttt{n}); los otros dos quedan en \texttt{\ensuremath{\Omega}\ensuremath{^{-}}}. El nodo solitario es el \textbf{único} que tiene \texttt{\ensuremath{\varphi} = 1}; los otros dos tienen \texttt{\ensuremath{\varphi} = 0}. La identificación se hace al activar el elemento, comparando el producto \texttt{(x\_i − x\_c)\ensuremath{\cdot}n} para cada nodo (\texttt{x\_c} = centroide), donde el signo del producto determina el lado.

\subsection{7. Funciones de forma + función \texttt{\ensuremath{\varphi}}}

\texttt{N\_i(x)} idénticas a \texttt{Tri3} §7. Para \texttt{\ensuremath{\varphi}}:

\[
\varphi(\mathbf x) = N_{i^*}(\mathbf x)
\]

donde \texttt{i*} es el índice del nodo solitario. Es lineal y \texttt{C\textasciicircum{}0} continua, vale 1 en \texttt{i*} y 0 en los otros dos. Sus derivadas son \texttt{(\ensuremath{\partial}N\_\{i*\}/\ensuremath{\partial}x, \ensuremath{\partial}N\_\{i*\}/\ensuremath{\partial}y)}, \textbf{constantes} sobre el elemento.

\subsection{8. Matrices B estándar + B\textasciicircum{}\ensuremath{\varphi}}

\texttt{B\_std} es la matriz B clásica del CST (constante, 3\ensuremath{\times}6); ver \texttt{Tri3} §8.

\texttt{B\textasciicircum{}\ensuremath{\varphi}} es la matriz Voigt 3\ensuremath{\times}2 construida con \texttt{(\ensuremath{\partial}\ensuremath{\varphi}/\ensuremath{\partial}x, \ensuremath{\partial}\ensuremath{\varphi}/\ensuremath{\partial}y)}:

\[
\mathbf B^\varphi = \begin{bmatrix} \partial\varphi/\partial x & 0 \\ 0 & \partial\varphi/\partial y \\ \partial\varphi/\partial y & \partial\varphi/\partial x \end{bmatrix}
\]

Por construcción, \texttt{B\textasciicircum{}\ensuremath{\varphi}} coincide con la \textbf{submatriz B\_std correspondiente al nodo solitario} (las columnas \texttt{2\ensuremath{\cdot}i*} y \texttt{2\ensuremath{\cdot}i*+1} de \texttt{B\_std}). Eso permite implementación trivial sin recomputar derivadas.

\subsection{9. Sistema acoplado --- matrices elementales}

Linealización del residual acoplado conduce al sistema 8\ensuremath{\times}8 local (6 DOFs estándar + 2 enriquecidos):

\[
\begin{bmatrix} \mathbf K_{dd} & \mathbf K_{d\llbracket u\rrbracket} \\ \mathbf K_{\llbracket u\rrbracket d} & \mathbf K_{\llbracket u\rrbracket\llbracket u\rrbracket} \end{bmatrix} \begin{bmatrix} \Delta\mathbf d \\ \Delta\llbracket\mathbf u\rrbracket \end{bmatrix} = \begin{bmatrix} -\mathbf R^d \\ -\mathbf R^{\llbracket u\rrbracket} \end{bmatrix}
\]

con (todas las cantidades evaluadas en el único punto de Gauss central, peso \texttt{A\_e\ensuremath{\cdot}t}):

\begin{itemize}
  \item \texttt{K\_dd = B\_std\textasciicircum{}T \ensuremath{\cdot} C\_e \ensuremath{\cdot} B\_std \ensuremath{\cdot} A\_e \ensuremath{\cdot} t} (6\ensuremath{\times}6) --- rigidez estándar del CST.
  \item \texttt{K\_\{d[[u]]\} = -B\_std\textasciicircum{}T \ensuremath{\cdot} C\_e \ensuremath{\cdot} B\textasciicircum{}\ensuremath{\varphi} \ensuremath{\cdot} G \ensuremath{\cdot} A\_e \ensuremath{\cdot} t} (6\ensuremath{\times}2) --- acoplamiento.
  \item \texttt{K\_\{[[u]] d\} = -G\textasciicircum{}T \ensuremath{\cdot} (B\textasciicircum{}\ensuremath{\varphi})\textasciicircum{}T \ensuremath{\cdot} C\_e \ensuremath{\cdot} B\_std \ensuremath{\cdot} A\_e \ensuremath{\cdot} t} (2\ensuremath{\times}6) --- traspuesta de \texttt{K\_\{d[[u]]\}} en KOS (simétrica por construcción).
  \item \texttt{K\_\{[[u]][[u]]\} = G\textasciicircum{}T \ensuremath{\cdot} (B\textasciicircum{}\ensuremath{\varphi})\textasciicircum{}T \ensuremath{\cdot} C\_e \ensuremath{\cdot} B\textasciicircum{}\ensuremath{\varphi} \ensuremath{\cdot} G \ensuremath{\cdot} A\_e \ensuremath{\cdot} t + l\_d \ensuremath{\cdot} T\_coh} (2\ensuremath{\times}2) --- rigidez del salto, suma de contribución del bulk (descarga elástica del modo enriquecido) y de la rigidez tangente cohesiva \texttt{T\_coh ≡ \ensuremath{\partial}t/\ensuremath{\partial}[[u]]} proyectada sobre \texttt{Γ\_d}.
\end{itemize}

Convención de signos: el \texttt{-} en \texttt{K\_\{d[[u]]\}} y \texttt{K\_\{[[u]]d\}} viene del signo del término enriquecido en \texttt{B\textasciicircum{}enr = -B\textasciicircum{}\ensuremath{\varphi}\ensuremath{\cdot}G} (recordar \texttt{(H - \ensuremath{\varphi})} con \texttt{\ensuremath{\nabla}H = \ensuremath{\delta}\_\{Γ\_d\}\ensuremath{\cdot}n}).

\subsection{10. Condensación estática local}

El sistema 8\ensuremath{\times}8 se condensa eliminando los DOFs enriquecidos \texttt{\ensuremath{\Delta}[[u]]} antes de devolver \texttt{K\_e} al ensamblador (ADR 0010 §3):

\[
\Delta\llbracket\mathbf u\rrbracket = -\mathbf K_{\llbracket u\rrbracket\llbracket u\rrbracket}^{-1}\,\bigl(\mathbf R^{\llbracket u\rrbracket} + \mathbf K_{\llbracket u\rrbracket d}\,\Delta\mathbf d\bigr)
\]

Sustituyendo en la primera ecuación:

\[
\boxed{\;\mathbf K_{\text{cond}} = \mathbf K_{dd} - \mathbf K_{d\llbracket u\rrbracket}\,\mathbf K_{\llbracket u\rrbracket\llbracket u\rrbracket}^{-1}\,\mathbf K_{\llbracket u\rrbracket d}\;}
\]

\[
\boxed{\;\mathbf F_{\text{int}}^{\text{cond}} = \mathbf F_{\text{int}}^d - \mathbf K_{d\llbracket u\rrbracket}\,\mathbf K_{\llbracket u\rrbracket\llbracket u\rrbracket}^{-1}\,\mathbf R^{\llbracket u\rrbracket}\;}
\]

donde \texttt{F\_int\textasciicircum{}d = B\_std\textasciicircum{}T \ensuremath{\cdot} \ensuremath{\sigma} \ensuremath{\cdot} A\_e \ensuremath{\cdot} t} es la fuerza interna estándar antes de condensar.

\textbf{Inversión de \texttt{K\_\{[[u]][[u]]\}} (2\ensuremath{\times}2)}: cerrada analíticamente. El despachador algebraico (ADR 0003) no participa al ser una matriz elemental densa de 2\ensuremath{\times}2.

\textbf{Recuperación local} tras converger el Newton global: \texttt{[[u]]\_committed += \ensuremath{\Delta}[[u]]} con la fórmula de arriba aplicada a \texttt{\ensuremath{\Delta}d = d\_n+1 − d\_n}. Esta recuperación se hace en \texttt{commit\_state()}, no en \texttt{compute\_element\_state()}.

\subsection{11. Activación: criterio Rankine + dirección \texttt{n} + \texttt{l\_d} del Cap. 6}

\textbf{Criterio de activación}. Al inicio de cada paso de carga (no dentro del Newton --- ver §14), el elemento intacto evalúa su tensión principal mayor \texttt{\ensuremath{\sigma}\_I} en el centroide:

\[
\sigma_I = \frac{\sigma_{xx} + \sigma_{yy}}{2} + \sqrt{\Bigl(\frac{\sigma_{xx} - \sigma_{yy}}{2}\Bigr)^2 + \sigma_{xy}^2}
\]

Si \texttt{\ensuremath{\sigma}\_I > \ensuremath{\sigma}\_t0} (resistencia a tracción del material cohesivo), se activa:

\begin{itemize}
  \item \texttt{n} = autovector unitario de \texttt{\ensuremath{\sigma}\_I} (perpendicular a la dirección de tracción máxima, criterio Rankine).
  \item \texttt{s} = ortogonal a \texttt{n} orientada por regla de la mano derecha con \texttt{+z} saliendo del plano (consistente con Reglas.md §5).
  \item \texttt{Γ\_d} pasa por el centroide del CST (decisión cerrada: la posición no es DOF; se fija en el activación, igual que \texttt{n}).
  \item Identificación del nodo solitario \texttt{i*} (§6).
\end{itemize}

\textbf{Longitud efectiva \texttt{l\_d} --- Cap. 6 de Retama (2010)}. Fórmula cerrada:

\[
\boxed{\;l_d = \frac{A_e}{h}\,\cos(\theta - \alpha)\;}
\]

donde:
\begin{itemize}
  \item \texttt{A\_e} es el área del triángulo.
  \item \texttt{h} es la altura desde el nodo solitario \texttt{i*} al lado opuesto.
  \item \texttt{\ensuremath{\theta}} es el ángulo de \texttt{Γ\_d} respecto a \texttt{x} global (i.e. \texttt{\ensuremath{\theta} = atan2(s\_y, s\_x)} con \texttt{s = (s\_x, s\_y)} la tangente a \texttt{Γ\_d}).
  \item \texttt{\ensuremath{\alpha}} es el ángulo del lado opuesto al nodo solitario respecto a \texttt{x} global.
\end{itemize}

Esta fórmula es \textbf{decisión cerrada en el ADR 0010} (§6); no se relitiga aquí. La alternativa ingenua \texttt{l\_d = A\_e/h} (independiente de \texttt{\ensuremath{\theta}}) produce stress locking cuando la grieta no es paralela al lado opuesto; ver fase 3 del ADR 0010 para test numérico.

\subsection{12. Estado de la discontinuidad --- \texttt{DiscontinuityState}}

Dataclass paralela a \texttt{ElementState} (ADR 0010 §4), específica de la discontinuidad. Vive en \texttt{solidum/core/discontinuity\_state.py}:

\begin{lstlisting}[language=Python]
@dataclass
class DiscontinuityState:
    normal: np.ndarray        # n ∈ ℝ², unitario, fijado al activar
    tangent: np.ndarray       # s ∈ ℝ², s = (−n_y, n_x)
    centroid: np.ndarray      # punto por el que pasa Γ_d
    solitary_node: int        # índice (0, 1 o 2) del nodo en Ω⁺
    l_d: float                # (A_e/h)·cos(θ − α)
    jump_committed: np.ndarray   # [[u]] al cierre del último paso convergido
    jump_trial: np.ndarray       # [[u]] dentro de la iteración Newton
    cohesive_state_committed: dict  # estado interno del CohesiveMaterial (committed)
    cohesive_state_trial: dict      # estado interno del CohesiveMaterial (trial)

\end{lstlisting}

Semántica \textbf{trial / commit} consistente con \texttt{ElementState}: \texttt{jump\_trial} y \texttt{cohesive\_state\_trial} evolucionan en cada iteración del Newton global; al converger el paso, \texttt{commit\_state()} copia \texttt{trial \ensuremath{\to} committed} y aplica la recuperación local de \texttt{\ensuremath{\Delta}[[u]]} (§10).

\subsection{13. Cuadratura}

Un único punto de Gauss central, peso \texttt{A\_e\ensuremath{\cdot}t} --- idéntico a \texttt{Tri3}. Las matrices \texttt{B\_std}, \texttt{B\textasciicircum{}\ensuremath{\varphi}} y la deformación del bulk son constantes; la integración es exacta. El término cohesivo \texttt{l\_d\ensuremath{\cdot}t} se evalúa directamente sobre \texttt{Γ\_d} (no requiere cuadratura porque \texttt{t} es constante a lo largo de \texttt{Γ\_d} por la cinemática KOS --- el salto es constante).

\subsection{14. Activación al principio del paso --- interacción con el solver}

ADR 0010 §5 fija que la activación se evalúa \textbf{al inicio de cada paso de carga, no dentro del Newton}. Razón: dentro del Newton, \texttt{\ensuremath{\sigma}\_I} oscila por el predictor lineal y puede cruzar/descruzar el umbral varias veces antes de converger --- \textit{chattering}.

\textbf{Implementación propuesta} (a validar --- punto abierto en §Diálogo): añadir al contrato de \texttt{Element} un método opcional \texttt{prepare\_step(U\_committed)} que el \texttt{NonlinearSolver}/\texttt{ArcLengthSolver} invoca \textbf{una vez por paso} antes del primer Newton. La base lo hace no-op (la mayoría de elementos no necesitan preparación); \texttt{CST\_Embedded2D} lo sobreescribe para evaluar la activación con el estado convergido del paso anterior. Una vez activado, el elemento queda agrietado \textbf{irreversiblemente} (no revierte si el siguiente paso lo descarga).


\section{Caveats numéricos}

\begin{itemize}
  \item \textbf{Singularidad de \texttt{K\_\{[[u]][[u]]\}} en Modo-I puro saturado}. Cuando \texttt{\ensuremath{\omega} \ensuremath{\to} 1} en el cohesivo, la rigidez tangente cohesiva \texttt{T\_coh} se aproxima a la diagonal \texttt{[(1−DAMAGE\_MAX)\ensuremath{\cdot}K\_e, 0]} (la componente tangencial es cero por construcción del Modo-I). La suma con la contribución del bulk \texttt{G\textasciicircum{}T\ensuremath{\cdot}(B\textasciicircum{}\ensuremath{\varphi})\textasciicircum{}T\ensuremath{\cdot}C\_e\ensuremath{\cdot}B\textasciicircum{}\ensuremath{\varphi}\ensuremath{\cdot}G} introduce rigidez en la dirección tangencial vía el bulk; la matriz 2\ensuremath{\times}2 sigue siendo invertible. Verificable analíticamente.
  \item \textbf{Activación al principio del paso}. Si el incremento es muy grande, \texttt{\ensuremath{\sigma}\_I} puede saltar muy por encima de \texttt{\ensuremath{\sigma}\_t0} antes de la activación, generando un \ensuremath{\Delta}[[u]] grande en el primer Newton. Mitigación: pasos suficientemente pequeños o \texttt{ArcLengthSolver} cerca del pico.
  \item \textbf{Bulk descarga elástica}. Por la cinemática KOS, el bulk evalúa \texttt{\ensuremath{\varepsilon}\textasciicircum{}bulk = B\_std\ensuremath{\cdot}d − B\textasciicircum{}\ensuremath{\varphi}\ensuremath{\cdot}G\ensuremath{\cdot}[[u]]} --- descarga conforme \texttt{[[u]]} crece. Esto es físicamente correcto (la disipación va a \texttt{Γ\_d}) y consistente con la discrete approach de Retama 2010. El bulk \textbf{no acumula daño} propio (out\_of\_scope fase 1).
  \item \textbf{\texttt{Γ\_d} paralelo a un lado}. Caso degenerado: \texttt{cos(\ensuremath{\theta}−\ensuremath{\alpha}) = 1}, \texttt{l\_d = A\_e/h}. La fórmula del Cap. 6 sigue siendo correcta, pero numéricamente el nodo solitario debe identificarse con tolerancia para evitar oscilación entre dos nodos casi en el mismo plano.
  \item \textbf{\texttt{Γ\_d} casi vertical en CST con jacobiano malacondicionado}. Sin patología nueva: hereda los caveats del CST padre. El test \texttt{det J > tol} de \texttt{Tri3} cubre el caso.
  \item \textbf{Stress locking}. La fase 3 del ADR 0010 valida numéricamente que \texttt{l\_d = (A\_e/h)\ensuremath{\cdot}cos(\ensuremath{\theta}−\ensuremath{\alpha})} evita el locking que aparece con \texttt{l\_d = A\_e/h} ingenuo. Test cubierto en \texttt{tests/test\_ld\_chapter\_6\_validation.py}: con \texttt{Γ\_d} oblicua respecto al lado opuesto al nodo solitario (\texttt{cos(\ensuremath{\theta}−\ensuremath{\alpha}) = \ensuremath{\surd}3/2} en el setup), la versión correcta produce una apertura \texttt{[[u\_n]]} \ensuremath{\approx} 5.45 \% mayor que la versión ingenua para una misma carga, en rama de softening exponencial activo. La versión ingenua "infla" la rigidez cohesiva (\texttt{l\_d \ensuremath{\cdot} T\_coh}) por un factor \texttt{1/cos(\ensuremath{\theta}−\ensuremath{\alpha}) > 1}, restringe la apertura y el elemento se comporta como si fuera más rígido --- eso es el \textit{stress locking}. Cuando \texttt{Γ\_d} es paralela al lado opuesto (\texttt{cos(\ensuremath{\theta}−\ensuremath{\alpha}) = \ensuremath{\pm}1}), las dos fórmulas coinciden exactamente y no hay locking.
\end{itemize}


\section{Contrato de implementación}

\begin{lstlisting}[language=yaml]
name: CST_Embedded2D
kind: element
status: validated

interface:
  dof_names: [ux, uy]           # 2 DOFs globales por nodo (los enriquecidos NO son globales)
  n_nodes: 3
  strain_dim: 3                 # Voigt 2D [ε_xx, ε_yy, γ_xy]
  n_integration_points: 1       # CST con cuadratura central

parameters:
  - { name: thickness,         type: float, required: false, default: 1.0,
      desc: "Espesor para estado plano (idéntico a Tri3)." }
  - { name: cohesive_material, type: str,   required: true,
      desc: "Nombre (en CohesiveMaterialRegistry) del material cohesivo que gobierna Γ_d cuando se activa." }
  - { name: activation_criterion, type: str, required: false, default: rankine,
      desc: "Criterio de activación: 'rankine' (única opción fase 1)." }

material_contract:
  bulk:
    signature: "compute_state(ε, state) -> (σ, C_tangent, state')"
    strain_kind: "Voigt 2D [ε_xx, ε_yy, γ_xy]"
    accepted: ["Elastic2D"]      # fase 1: bulk elástico; otros → out_of_scope
  cohesive:
    signature: "compute_traction([[u]], state) -> (t, T_tan, state')"
    jump_kind: "frame local — [[u]] = ([[u_n]], [[u_s]]) en ejes (n, s) de Γ_d"
    accepted: ["CohesiveDamageIsotropic"]  # cualquier CohesiveMaterial con JUMP_DIM=2

conventions:
  sign: "tracción interna positiva; convención stress-resultant del proyecto (Reglas.md §5)."
  voigt: "[xx, yy, xy] con γ_xy = 2·ε_xy."
  node_orientation: "antihorario (det J > 0)."
  frame_local_discontinuity: "n perpendicular a σ_I al activarse; s = (−n_y, n_x) (RHR con +z saliendo del plano)."
  activation: "irreversible; evaluada al principio de cada paso con el estado convergido previo."

validity:
  - "Pequeños desplazamientos y deformaciones (no large strain)."
  - "Bulk elástico (Elastic2D) y cohesivo Modo-I (JUMP_DIM = 2)."
  - "Γ_d cortando exactamente dos lados del triángulo (configuración geométrica estándar)."

out_of_scope:
  - "Bulk no lineal (J2, daño continuo, Drucker-Prager): la disipación debe localizarse en Γ_d, no co-existir con disipación de bulk."
  - "Modo mixto I-II en el cohesivo (fase G del ADR 0010)."
  - "Reorientación de n tras la activación (tracking no trivial, fase F)."
  - "Múltiples discontinuidades por elemento."
  - "Contacto unilateral en compresión (cierre rígido)."
  - "Elementos de orden superior con discontinuidad embebida (Tri6_Embedded, Quad8_Embedded; fase J)."
  - "3D (fase I)."

acceptance:
  verification:
    - name: estado_intacto_reproduce_Tri3
      setup: "elemento sin activar (discontinuity_state=None), bulk Elastic2D; aplicar campo lineal de desplazamientos."
      expect: "K_e y F_int_e idénticos (bit-exact en doble) a los de un Tri3 con el mismo bulk."
      tol_rel: 1.0e-14

    - name: patch_test_intacto
      setup: "patch MacNeal-Harder con 4 CST_Embedded2D distorsionados, ninguno activado, σ < σ_t0 en todo el dominio."
      expect: "campo lineal exacto en nodos interiores; ε constante por elemento."
      tol_rel: 1.0e-10

    - name: activacion_rankine_correcta
      setup: "tracción uniaxial pura en x sobre un único CST con σ_t0 conocido; incrementar carga hasta cruzar σ_t0."
      expect: "activación en el paso donde σ_xx > σ_t0; n = (1, 0) y s = (0, 1) con error angular < 1e-10."
      tol_rel: 1.0e-10

    - name: bulk_descarga_elasticamente_post_activacion
      setup: "activar elemento, después fijar [[u]] e imponer Δd; verificar que la deformación del bulk responde como elástica con C_e."
      expect: "σ_bulk = C_e · (B_std·d − B^φ·G·[[u]]); coincide con C_e aplicado a la ε neta."
      tol_rel: 1.0e-12

    - name: condensacion_simetrica_KOS
      setup: "elemento activado con cohesivo en carga activa simétrica (Modo-I puro)."
      expect: "K_cond simétrica: ‖K_cond − K_cond^T‖_F / ‖K_cond‖_F < 1e-12."
      tol_rel: 1.0e-12

    - name: ld_formula_cerrada
      setup: "varias orientaciones de Γ_d (θ ∈ {0, π/6, π/4, π/3, π/2}) sobre un triángulo conocido."
      expect: "l_d coincide con (A_e/h)·cos(θ−α) calculado analíticamente."
      tol_rel: 1.0e-14

    - name: recuperacion_local_del_salto
      setup: "carga monótona hasta varios pasos post-activación; tras cada commit, [[u]]_committed debe satisfacer R^{[[u]]} = 0."
      expect: "‖R^{[[u]]}‖ < 1e-8 tras commit_state."
      tol_rel: 1.0e-8

  specific:
    - name: bulk_no_lineal_rechazado
      setup: "instanciar CST_Embedded2D con bulk_material=IsotropicDamage2D."
      expect: "ValueError clara en construcción, mensaje listando los bulks aceptados (fase 1: Elastic2D)."

    - name: jump_dim_validado
      setup: "instanciar con cohesive_material cuyo JUMP_DIM=3 (futuro 3D)."
      expect: "ValueError con mensaje sobre incompatibilidad dimensional (elemento 2D requiere JUMP_DIM=2)."

    - name: activacion_no_chattering_dentro_de_newton
      setup: "elemento cerca del umbral; correr iteraciones Newton sin invocar prepare_step."
      expect: "estado discontinuity_state no cambia durante las iteraciones; sólo prepare_step modifica activación."

    - name: identificacion_nodo_solitario
      setup: "construir CST con coordenadas tales que solo un nodo cae en Ω⁺ para una n dada."
      expect: "DiscontinuityState.solitary_node identifica correctamente el índice; los otros dos tienen φ=0."

    - name: ld_invariante_bajo_reflejo_de_n
      setup: "activar con n y comparar con activar con −n (mismo problema físico)."
      expect: "l_d, K_cond, F_int^cond idénticos (la convención de n positivo está bien definida — apunta hacia el nodo solitario)."
      tol_rel: 1.0e-14

  arch:
    - name: registro_en_ElementRegistry
      expect: "CST_Embedded2D registrado con kind='element'."

    - name: STRAIN_DIM_y_DOF_NAMES_correctos
      expect: "STRAIN_DIM=3, DOF_NAMES=['ux','uy'], N_NODES=3, N_INTEGRATION_POINTS=1."

    - name: prepare_step_es_no_op_en_Element_base
      expect: "Element.prepare_step existe y no hace nada por defecto; CST_Embedded2D lo sobreescribe."

references:
  - "Retama Velasco, J. (2010). Formulation and Approximation to Problems in Solids by Embedded Discontinuity Models. Tesis Doctoral, UNAM. **Caps. 2 (cinemática KOS), 5 (formulación variacional discreta), 6 (l_d correcto), 7 (condensación estática)**."
  - "Jirásek, M. (2000). Comparative study on finite elements with embedded discontinuities. CMAME 188, 307-330. — Clasificación SOS/KOS/SKON."
  - "Oliver, J., Huespe, A.E., Sánchez, P.J. (2006). A comparative study on finite elements for capturing strong discontinuities. CMAME 195, 4732-4752."
  - "Linder, C., Armero, F. (2007). Finite elements with embedded strong discontinuities. IJNME 72, 1391-1433."
  - "ADR 0010 — Discontinuidades interiores embebidas (hoja de ruta; fase 2 = este elemento)."
  - "Spec [CohesiveDamageIsotropic](CohesiveDamageIsotropic.md) — material cohesivo consumido en Γ_d."
  - "Spec [Tri3](Tri3.md) — CST padre del elemento; mismo bulk en estado intacto."

\end{lstlisting}


\section{Implementación}

\begin{itemize}
  \item Archivo: \texttt{solidum/elements/solid\_2d/embedded\_cst.py}
  \item Clase: \texttt{CST\_Embedded2D} (registrada en \texttt{ElementRegistry})
  \item Estado de la discontinuidad: \texttt{solidum/core/discontinuity\_state.py} (\texttt{DiscontinuityState})
  \item Hook de paso en clase base: \texttt{Element.prepare\_step(U\_committed)} (no-op por defecto)
  \item Integración en solvers:
  \item \texttt{NonlinearSolver}: invoca \texttt{self.assembler.prepare\_all\_steps(U\_current)} al inicio de cada paso, antes del Newton.
  \item \texttt{ArcLengthSolver}: idem, antes del predictor.
  \item Parser YAML: nueva sección \texttt{cohesive\_materials:} + campo \texttt{cohesive\_material:} por elemento.
  \item Tests: \texttt{tests/test\_cst\_embedded.py} --- 24 tests cubriendo \texttt{acceptance} completo (verification, specific, arch) + integración end-to-end del YAML.
  \item Notas de traducción:
  \item Estado intacto (\texttt{discontinuity\_state is None}) \ensuremath{\to} bit-exact con \texttt{Tri3} (mismo \texttt{B}, mismo \texttt{\_compute\_integrands}).
  \item Estado agrietado: Newton local sobre \texttt{[[u]]} hasta \texttt{R\textasciicircum{}\{[[u]]\} = 0}, después condensación estática local cerrada (\texttt{K\_jj} es 2\ensuremath{\times}2). Tolerancia interna: \texttt{\_LOCAL\_JUMP\_RTOL = 1e-10}, \texttt{\_LOCAL\_JUMP\_MAX\_ITER = 30}.
  \item \texttt{B\textasciicircum{}\ensuremath{\varphi}} se obtiene como la submatriz \texttt{B\_std[:, 2\ensuremath{\cdot}i*:2\ensuremath{\cdot}i*+2]} correspondiente al nodo solitario (consistente con \texttt{\ensuremath{\varphi} = N\_\{i*\}}).
  \item Activación: \texttt{n} apunta hacia el nodo solitario (convención fija); si el autovector de \texttt{\ensuremath{\sigma}\_I} deja 2 nodos en lado positivo, se invierte el signo de \texttt{n}.
  \item \texttt{\_compute\_ld} usa \texttt{|cos(\ensuremath{\theta} − \ensuremath{\alpha})|} (valor absoluto): la longitud es positiva y simétrica bajo \texttt{n \ensuremath{\to} −n} (test \texttt{ld\_invariante\_bajo\_reflejo\_de\_n}).
  \item \texttt{compute\_gauss\_state} extendido: en estado agrietado añade clave \texttt{'discontinuity': \{normal, tangent, centroid, solitary\_node, l\_d, jump, traction, damage\}} para post-proceso VTK.
  \item Bulks aceptados en fase 1: \texttt{\_ACCEPTED\_BULKS = ('Elastic2D',)}. Restricción declarativa, ampliación futura es aditiva.
\end{itemize}


\section{Diálogo}

\begin{itemize}
  \item \textbf{2026-05-18} \ensuremath{\cdot} Spec pre-redactada por la IA como \textbf{orden de trabajo} sobre los Caps. 2, 5, 6 y 7 de Retama 2010. La cinemática KOS, la condensación estática y el \texttt{l\_d = (A/h)\ensuremath{\cdot}cos(\ensuremath{\theta}−\ensuremath{\alpha})} son traducción directa de la tesis; las decisiones arquitecturales del ADR 0010 (familia paralela, condensación local, irreversibilidad, activación al principio del paso) están fijadas. Pendiente de validación del usuario.
\end{itemize}

\begin{itemize}
  \item \textbf{2026-05-18} \ensuremath{\cdot} Usuario delega las decisiones a la IA. Cierre con justificación arquitectural:
  \item \textbf{A --- Bulk restringido a \texttt{Elastic2D} en fase 1.} Adoptado. Razón: la \textit{discrete approach} del Cap. 3-7 de la tesis presupone bulk elástico (la disipación se localiza en \texttt{Γ\_d}); mezclar plasticidad continua del bulk con cohesivo de \texttt{Γ\_d} es una formulación distinta no cubierta por el ADR 0010. Validación cruzada en \texttt{\_\_init\_\_} con mensaje listando bulks aceptados. Apertura futura a otros materiales es aditiva.
  \item \textbf{B --- Hook \texttt{prepare\_step(U\_committed)} en \texttt{Element} base, no-op por defecto.} Adoptado. Lo prescribe el ADR 0010 §5 explícitamente; aditivo (no rompe llamadores) y testeable directamente. Alternativa de meter activación en \texttt{compute\_element\_state} con guarda \texttt{\_step\_id} descartada por ser frágil y mezclar conceptos (preparación vs cálculo).
  \item \textbf{C --- \texttt{INTERNAL\_DOF\_NAMES} no se añade al contrato declarativo todavía.} Resistir abstracción especulativa (Reglas.md §1: justificar con dos casos reales). Cuando entre el segundo elemento enriquecido (Quad4\_Embedded o Tet4\_Embedded), se centraliza la convención.
  \item \textbf{D --- Post-proceso vía \texttt{compute\_gauss\_state(U)} extendido.} Adoptado. En estado intacto se comporta como \texttt{Tri3}; en estado agrietado añade la clave \texttt{'discontinuity': \{'normal', 'jump', 'traction', 'damage', 'l\_d'\}}. Consistente con el patrón actual del \texttt{Tri3} y aditivo para \texttt{VtkExporter}. Un método separado \texttt{crack\_state()} duplicaría el patrón sin beneficio.
  \item \textbf{E --- Constructor \texttt{\_\_init\_\_(element\_id, nodes, material, cohesive\_material, thickness=1.0)}.} Mantengo el nombre \texttt{material} para el bulk (consistente con \texttt{Element.\_\_init\_\_}, que ya lo recibe así); \texttt{cohesive\_material} es el parámetro nuevo. El YAML mapea \texttt{bulk\_material: <name>} \ensuremath{\to} \texttt{material=...} y \texttt{cohesive\_material: <name>} \ensuremath{\to} \texttt{cohesive\_material=...}.
  \item \textbf{F --- Tolerancia \texttt{1e-14} (bit-exact)} en \texttt{estado\_intacto\_reproduce\_Tri3}. Si en la implementación emerge un desliz numérico documentado (por canceladas algebraicas tipo \texttt{G\ensuremath{\cdot}G\textasciicircum{}T = I}), se relaja con razón explícita en el Diálogo. Empezar estricta es mejor que pre-laxar.
  \item \textbf{2026-05-18} \ensuremath{\cdot} Status \ensuremath{\to} en draft hasta que el código + tests cierren; se promoverá a \texttt{validated} cuando los \texttt{acceptance} pasen. La IA arranca implementación.
\end{itemize}


\newpage

\chapter{Elementos 3D — Sólidos}

\section{Tet4}



\textit{Elemento sólido 3D de primer orden, espejo natural de \texttt{Tri3}. Pieza de la Etapa 7. Convención Voigt 6D y numeración de caras por ADR 0012.}



\section{Especificación física}

\subsection{0. Descripción general}

Elemento sólido tridimensional de primer orden, tetraedro de cuatro nodos. Aproximación \textbf{lineal} $C^0$ del campo de desplazamientos en coordenadas baricéntricas. Análogo 3D del \texttt{Tri3}: la matriz $\mathbf B$ es \textbf{constante sobre el elemento}, por lo que $\boldsymbol\varepsilon$ y $\boldsymbol\sigma$ son uniformes dentro de cada \texttt{Tet4} (\textit{constant strain tetrahedron} en la literatura). Régimen de pequeños desplazamientos y deformaciones.

\subsection{1. Cinemática de desplazamientos}

\[
u_x = \sum_{i=1}^{4} N_i\, u_x^{(i)}, \qquad u_y = \sum N_i\, u_y^{(i)}, \qquad u_z = \sum N_i\, u_z^{(i)}.
\]

\subsection{2. Cinemática de deformaciones}

Notación Voigt 6D del proyecto (ADR 0012):

\[
\boldsymbol\varepsilon = [\varepsilon_{xx},\ \varepsilon_{yy},\ \varepsilon_{zz},\ \gamma_{xy},\ \gamma_{yz},\ \gamma_{xz}]^\top
\]

Por la linealidad de $N_i$, las derivadas $\partial N_i / \partial \mathbf x$ son constantes y $\boldsymbol\varepsilon$ es uniforme en el elemento. Idéntico patrón conceptual al \texttt{Tri3} en 2D.

\subsection{3. Ecuación constitutiva}

Delegada al material 3D (\texttt{STRAIN\_DIM = 6}). En esta etapa solo \texttt{Elastic3D} está disponible.

\subsection{4. Equilibrio --- forma fuerte}

Idéntica a la del \texttt{Hex8}. El principio variacional no depende de la forma del elemento.

\subsection{5. Equilibrio --- forma débil}

Idéntica a la del \texttt{Hex8}.


\section{Formulación numérica (FEM)}

\subsection{6. Discretización}

Cuatro nodos en orden estándar de tetraedro VTK. En el tetraedro de referencia con vértices:

\begin{center}
\begin{tabularx}{\textwidth}{YY}
\hline
\textbf{Nodo local} & \textbf{$(\xi_i, \eta_i, \zeta_i)$} \\
\hline
0 & $(0, 0, 0)$ \\
1 & $(1, 0, 0)$ \\
2 & $(0, 1, 0)$ \\
3 & $(0, 0, 1)$ \\
\hline
\end{tabularx}
\end{center}
El orden de nodos asegura que el vector $(\mathbf x_1 - \mathbf x_0) \times (\mathbf x_2 - \mathbf x_0)$ apunte hacia $\mathbf x_3$ (volumen orientado positivo). Mapeo isoparamétrico desde el tetraedro de referencia; jacobiano $\mathbf J$ constante en todo el elemento; $\det\mathbf J = 6\,V_e$ con $V_e$ el volumen del tetraedro. El código aborta con \texttt{ValueError} si $\det\mathbf J \le \text{tol}$ (nodos coplanares o en orden invertido).

\subsection{7. Funciones de forma}

Coordenadas baricéntricas / tetraédricas:

\[
N_1 = 1 - \xi - \eta - \zeta, \qquad N_2 = \xi, \qquad N_3 = \eta, \qquad N_4 = \zeta.
\]

Sus derivadas en coordenadas naturales son constantes:

\[
\partial N_i / \partial \xi = (-1, 1, 0, 0),\quad \partial N_i / \partial \eta = (-1, 0, 1, 0),\quad \partial N_i / \partial \zeta = (-1, 0, 0, 1).
\]

\subsection{8. Matriz deformación-desplazamiento}

$\mathbf B$ tiene tamaño $6 \times 12$ y, una vez transformadas las derivadas a globales vía $\mathbf J^{-1}$, sus entradas son \textbf{constantes} sobre el elemento:

\[
\mathbf B_i = \begin{bmatrix}
\partial N_i / \partial x & 0 & 0 \\
0 & \partial N_i / \partial y & 0 \\
0 & 0 & \partial N_i / \partial z \\
\partial N_i / \partial y & \partial N_i / \partial x & 0 \\
0 & \partial N_i / \partial z & \partial N_i / \partial y \\
\partial N_i / \partial z & 0 & \partial N_i / \partial x
\end{bmatrix}, \qquad \boldsymbol\varepsilon = \mathbf B\,\mathbf u_e.
\]

Mismo patrón filas/columnas que \texttt{Hex8} (ver §8 de su spec).

\subsection{9. Rigidez elemental}

\[
\mathbf K_e = \mathbf B^\top\,\mathbf D\,\mathbf B\, V_e,
\]

con $V_e = \tfrac{1}{6}\,\det\mathbf J$ el volumen del tetraedro. La integración es \textbf{exacta con un único punto central} porque $\mathbf B$ y $\boldsymbol\sigma$ son constantes.

\subsection{10. Fuerzas internas}

\[
\mathbf F_{\text{int}} = \mathbf B^\top\,\boldsymbol\sigma\, V_e.
\]

\subsection{11. Cuadratura}

Un punto central (baricentro del tetraedro) con peso $w = 1/6$ sobre el tetraedro de referencia. Es exacto para formas constantes de $\mathbf B$ y $\boldsymbol\sigma$.

\subsection{12. Cargas distribuidas consistentes}

\textbf{Fuerza de cuerpo $\mathbf b$}. Con $N_i$ lineales y $\mathbf b$ uniforme la integral es exacta:

\[
\mathbf f_e^{b} = \int_{\Omega_e} \mathbf N^\top\,\mathbf b\, dV \;\Longrightarrow\; \text{cada nodo recibe } \tfrac{1}{4}\,\mathbf b\,V_e.
\]

\textbf{Tracción de superficie $\bar{\mathbf t}$} sobre una cara triangular. Caras numeradas con normal saliente (ADR 0012):

\begin{center}
\begin{tabularx}{\textwidth}{YYY}
\hline
\textbf{Cara} & \textbf{Nodos} & \textbf{Nodo opuesto} \\
\hline
0 & (1, 2, 3) & 0 \\
1 & (0, 3, 2) & 1 \\
2 & (0, 1, 3) & 2 \\
3 & (0, 2, 1) & 3 \\
\hline
\end{tabularx}
\end{center}
Sobre una cara triangular con tracción constante el reparto es exacto:

\[
\mathbf f_e^{t} = \tfrac{1}{3}\,\bar{\mathbf t}\,A_\text{cara} \text{ por nodo del triángulo}, \qquad 0 \text{ al nodo opuesto}.
\]

Cuadratura 1 punto centroide de la cara (suficiente por linealidad de las funciones de forma sobre la cara). $\bar{\mathbf t}$ se especifica en coordenadas globales $(t_x, t_y, t_z)$. Tracción variable o presión normal no soportadas en este paso.

\subsection{13. Salida por punto de Gauss}

\texttt{compute\_gauss\_state(U)} devuelve la misma estructura que \texttt{Hex8} pero con un único punto (centroide). Para \texttt{Tet4} $\boldsymbol\varepsilon$ y $\boldsymbol\sigma$ son constantes sobre el elemento, así que el dato del punto coincide con el promedio del elemento. Por ADR 0012, \texttt{internal\_forces} devuelve \texttt{None}.


\section{Limitación crítica --- \textit{shear locking}}

El \texttt{Tet4} es notoriamente rígido en flexión y dominado por cortante, exactamente como su análogo \texttt{Tri3} en 2D: cuatro DOFs de desplazamiento por dirección no bastan para representar el modo de flexión puro de una viga discretizada en tetraedros, por lo que el elemento responde con cortante espurio y subestima sistemáticamente la deflexión. La severidad crece con la relación $L/h$ del problema.

\textbf{Recomendación operativa}: usar \texttt{Hex8} por defecto en mallas hexaédricas o mixtas; reservar \texttt{Tet4} para zonas de transición de malla donde la generación automática de hexaedros falla y se requiere malla tetraédrica de baja regularidad. Para problemas dominados por flexión se necesitan tetraedros cuadráticos \texttt{Tet10} (sub-etapa posterior). Locking volumétrico con $\nu \to 0.5$ es \textbf{aún peor} que en \texttt{Hex8} por la misma razón: el espacio lineal es muy pobre.


\section{Contrato de implementación}

\begin{lstlisting}[language=yaml]
name: Tet4
kind: element
status: validated

interface:
  dof_names: [ux, uy, uz]
  n_nodes: 4
  strain_dim: 6                 # Voigt 3D [ε_xx, ε_yy, ε_zz, γ_xy, γ_yz, γ_xz]
  n_integration_points: 1

parameters: []

material_contract:
  signature: "compute_state(ε, state) -> (σ, C_tangent, state')"
  strain_kind: "Voigt 3D [ε_xx, ε_yy, ε_zz, γ_xy, γ_yz, γ_xz] con γ_ij = 2·ε_ij"

conventions:
  sign: "σ > 0 ⇔ tracción (Reglas §5)"
  voigt: "[xx, yy, zz, xy, yz, xz] con γ_ij = 2·ε_ij (ADR 0012)"
  node_orientation: "convención VTK_TETRA; nodos 0-2 base con normal hacia 3 (volumen orientado positivo); asegura det J > 0"
  face_numbering: "ADR 0012 — 4 caras con normal saliente; cara i opuesta al nodo i"

validity:
  - "pequeños desplazamientos y deformaciones"
  - "campos de tensión que varían suavemente sobre el elemento"
  - "compatible con materiales con STRAIN_DIM = 6 (en esta etapa: Elastic3D)"

out_of_scope:
  - "flexión dominante (shear locking severo)"
  - "casi-incompresibilidad (ν → 0.5 — locking volumétrico aún peor que Hex8)"
  - "elementos de alto orden (Tet10 — sub-etapa posterior)"
  - "grandes deformaciones"
  - "internal_forces (ADR 0012: sólidos exponen compute_gauss_state)"

acceptance:
  verification:
    - name: patch_test_3d
      setup: "malla tetraédrica con al menos un nodo interior (e.g.
             división de un cubo en 5 tetraedros con un nodo central
             libre); campo lineal u_i = a_ij·x_j impuesto en el contorno"
      expect: "nodo interior adopta el campo lineal exacto y ε constante
              e igual a 0.5·(a_ij + a_ji) en cada elemento; σ uniforme"
      tol_rel: 1.0e-10
    - name: modos_de_cuerpo_rigido
      setup: "evaluar K_e sobre un Tet4 cualquiera y proyectar sobre los 6
             modos de cuerpo rígido"
      expect: "K_e · u_rigid ≈ 0 en cada uno de los 6 modos"
      tol_abs: 1.0e-9
    - name: simetria_K_e
      setup: "evaluar K_e con Elastic3D"
      expect: "K_e == K_e.T exacto"
      tol_abs: 1.0e-12
  specific:
    - name: parche_deformacion_constante
      setup: "malla tetraédrica regular con campo lineal y Elastic3D"
      expect: "ε constante (CST 3D reproduce campos lineales por construcción)"
      tol_rel: 1.0e-12
    - name: jacobiano_degenerado_abortado
      setup: "cuatro nodos coplanares"
      expect: "ValueError en _compute_kinematics_tet4"
    - name: body_load_uniforme
      setup: "Tet4 con b uniforme"
      expect: "cada nodo recibe (1/4)·b·V_e; Σ = b·V_e"
      tol_rel: 1.0e-12
    - name: traccion_uniforme_cara
      setup: "Tet4 con tracción constante en una cara"
      expect: "Σf = t̄·A_cara; reparto A_cara/3·t̄ a cada uno de los 3
              nodos de la cara, 0 al nodo opuesto"
      tol_rel: 1.0e-12
    - name: cubo_lame_3d_malla_tetraedrica
      setup: "cubo unitario discretizado con tetraedros (e.g. 5 o 6 Tet4
             por cubo); tracción uniaxial σ_xx = p en cara +ξ"
      expect: "u_x(L) ≈ p·L/E con error decreciente al refinar; orden
              de convergencia O(h) en norma energética"
      tol_rel: "documentar convergencia, sin tolerancia absoluta"

references:
  - "Cook R.D., Malkus D.S., Plesha M.E., Witt R.J. (2002). Concepts and Applications of FEA. §6.8 (CST 3D)."
  - "Zienkiewicz O.C., The Finite Element Method, vol. 1, §6.7 (limitations of linear tets)."
  - "Bathe K.J. (2014). Finite Element Procedures. §6.6.2."
  - "ADR 0012 — Sólidos 3D: convención Voigt 6D y caras."

\end{lstlisting}


\section{Implementación}

\begin{itemize}
  \item \textbf{Archivo}: solidum/elements/solid\_3d/tet4.py.
  \item \textbf{Clase}: \texttt{Tet4} (registrada vía \texttt{@ElementRegistry.register}).
  \item \textbf{Función núcleo}: \texttt{\_compute\_kinematics\_tet4(coords)} con \texttt{@njit} en solidum/elements/solid\_3d/\_shared.py. Reutiliza \texttt{\_compute\_integrands\_3d} (también en \texttt{\_shared}, compartido con Hex8) con peso \texttt{1/6} (volumen del tet de referencia).
  \item \textbf{Masa consistente analítica}: `\texttt{M\_ij = \ensuremath{\rho}\ensuremath{\cdot}V\_e\ensuremath{\cdot}(1 + \ensuremath{\delta}\_ij)/20}\texttt{ (matriz escalar 4\ensuremath{\times}4, expandida 3D vía }\_expand\_scalar\_mass\_3d`). Evita cuadratura insuficiente en el centroide.
  \item \textbf{Tests}:
  \item tests/test\_solid\_3d.py --- clase \texttt{TestTet4Element} (10 tests: dimensiones/DOFs, simetría exacta de K (no numérica, por fórmula cerrada), volumen, patch tracción uniaxial, jacobiano degenerado abortado, body load, face traction (cara opuesta a nodo recibe cero, reparto A/3 a los 3 nodos restantes), masa consistente y lumped HRZ (reparto equitativo \ensuremath{\rho}V/4 por DOF)).
  \item tests/test\_rigid\_body\_modes.py --- \texttt{TestRBMTet4} (5 tests: traslación, 3 rotaciones, rank-deficiency = 6).
  \item tests/validation/test\_cube\_lame\_3d.py --- \texttt{test\_uniaxial\_traction\_tet4\_mesh\_exact} (cubo dividido en 5 tetraedros con tracción uniaxial; solución exacta nodal).
\end{itemize}


\section{Diálogo}

\begin{itemize}
  \item \textbf{2026-05-19} \ensuremath{\cdot} Spec inicial. Elemento espejo natural de \texttt{Tri3}; convención Voigt 6D y numeración de caras fijadas por ADR 0012. Sin elementos de alto orden en esta etapa. Shear locking y locking volumétrico declarados como limitaciones; \texttt{Hex8} es la opción preferente cuando la malla lo permite.
\end{itemize}


\newpage

\section{Hex8}



\textit{Elemento sólido 3D de primer orden, espejo natural de \texttt{Quad4}. Pieza central de la Etapa 7. Convención Voigt 6D y caras fijadas por ADR 0012.}



\section{Especificación física}

\subsection{0. Descripción general}

Elemento sólido tridimensional de primer orden, hexaedro de ocho nodos. Aproximación \textbf{trilineal} $C^0$ del campo de desplazamientos en coordenadas naturales $(\xi, \eta, \zeta) \in [-1, 1]^3$. Formulación isoparamétrica: la geometría y los desplazamientos comparten las mismas funciones de forma. Régimen de pequeños desplazamientos y deformaciones.

\subsection{1. Cinemática de desplazamientos}

\[
u_x(\xi, \eta, \zeta) = \sum_{i=1}^{8} N_i(\xi, \eta, \zeta)\, u_x^{(i)}, \qquad u_y = \sum N_i\, u_y^{(i)}, \qquad u_z = \sum N_i\, u_z^{(i)}.
\]

\subsection{2. Cinemática de deformaciones}

Tensor infinitesimal en notación Voigt 6D del proyecto (ADR 0012):

\[
\boldsymbol\varepsilon = [\varepsilon_{xx},\ \varepsilon_{yy},\ \varepsilon_{zz},\ \gamma_{xy},\ \gamma_{yz},\ \gamma_{xz}]^\top
\]

con

\[
\varepsilon_{xx} = \tfrac{\partial u_x}{\partial x},\quad \varepsilon_{yy} = \tfrac{\partial u_y}{\partial y},\quad \varepsilon_{zz} = \tfrac{\partial u_z}{\partial z},
\]

\[
\gamma_{xy} = \tfrac{\partial u_x}{\partial y} + \tfrac{\partial u_y}{\partial x},\quad \gamma_{yz} = \tfrac{\partial u_y}{\partial z} + \tfrac{\partial u_z}{\partial y},\quad \gamma_{xz} = \tfrac{\partial u_x}{\partial z} + \tfrac{\partial u_z}{\partial x}.
\]

\subsection{3. Ecuación constitutiva}

Delegada al material 3D (\texttt{STRAIN\_DIM = 6}). El elemento llama \texttt{material.compute\_state(\ensuremath{\varepsilon}, state) \ensuremath{\to} (\ensuremath{\sigma}, C\_tangent, state')} por punto de Gauss. En esta etapa solo \texttt{Elastic3D} está disponible; el contrato deja la puerta abierta a futuros \texttt{VonMises3D}, \texttt{DruckerPrager3D}, \texttt{IsotropicDamage3D}.

\subsection{4. Equilibrio --- forma fuerte}

\[
-\nabla\cdot\boldsymbol\sigma = \mathbf b \quad \text{en } \Omega, \qquad \boldsymbol\sigma\cdot\mathbf n = \bar{\mathbf t} \quad \text{en } \partial\Omega_N, \qquad \mathbf u = \bar{\mathbf u} \quad \text{en } \partial\Omega_D.
\]

\subsection{5. Equilibrio --- forma débil}

\[
\int_\Omega \delta\boldsymbol\varepsilon^\top \boldsymbol\sigma\, dV = \int_\Omega \delta\mathbf u^\top \mathbf b\, dV + \int_{\partial\Omega_N} \delta\mathbf u^\top \bar{\mathbf t}\, dS, \qquad \forall\, \delta\mathbf u \in V_0.
\]


\section{Formulación numérica (FEM)}

\subsection{6. Discretización}

Ocho nodos en orden estándar de hexaedro VTK (compatible con la mayoría de pre-procesadores). En el cubo de referencia $[-1, 1]^3$:

\begin{center}
\begin{tabularx}{\textwidth}{YY}
\hline
\textbf{Nodo local} & \textbf{$(\xi_i, \eta_i, \zeta_i)$} \\
\hline
0 & $(-1, -1, -1)$ \\
1 & $(+1, -1, -1)$ \\
2 & $(+1, +1, -1)$ \\
3 & $(-1, +1, -1)$ \\
4 & $(-1, -1, +1)$ \\
5 & $(+1, -1, +1)$ \\
6 & $(+1, +1, +1)$ \\
7 & $(-1, +1, +1)$ \\
\hline
\end{tabularx}
\end{center}
Cara inferior $\zeta = -1$ recorrida \texttt{(0, 1, 2, 3)} antihorario visto desde dentro; cara superior $\zeta = +1$ recorrida \texttt{(4, 5, 6, 7)} antihorario visto desde fuera. Mapeo isoparamétrico desde el cubo de referencia:

\[
x(\xi, \eta, \zeta) = \sum_{i=0}^{7} N_i\, x^{(i)}, \quad y = \sum N_i\, y^{(i)}, \quad z = \sum N_i\, z^{(i)}.
\]

Jacobiano $\mathbf J = \partial(x, y, z) / \partial(\xi, \eta, \zeta)$ es una matriz 3\ensuremath{\times}3 con determinante $\det\mathbf J$. El código aborta con \texttt{ValueError} si $\det\mathbf J \le \text{tol}$ (elemento degenerado o nodos en orden invertido).

\subsection{7. Funciones de forma}

Producto tensorial trilineal:

\[
N_i(\xi, \eta, \zeta) = \tfrac{1}{8}(1 + \xi_i\,\xi)(1 + \eta_i\,\eta)(1 + \zeta_i\,\zeta), \qquad (\xi_i, \eta_i, \zeta_i) \in \{-1, +1\}^3.
\]

\subsection{8. Matriz deformación-desplazamiento}

Las derivadas en globales se obtienen vía $\partial N_i / \partial \mathbf x = \mathbf J^{-1}\,\partial N_i / \partial \boldsymbol\xi$. La matriz $\mathbf B$ tiene tamaño \textbf{6\ensuremath{\times}24} y se ensambla por nodo $i \in \{0,\dots,7\}$:

\[
\mathbf B_i = \begin{bmatrix}
\partial N_i / \partial x & 0 & 0 \\
0 & \partial N_i / \partial y & 0 \\
0 & 0 & \partial N_i / \partial z \\
\partial N_i / \partial y & \partial N_i / \partial x & 0 \\
0 & \partial N_i / \partial z & \partial N_i / \partial y \\
\partial N_i / \partial z & 0 & \partial N_i / \partial x
\end{bmatrix}, \qquad \boldsymbol\varepsilon = \mathbf B\,\mathbf u_e.
\]

El orden de las filas reproduce exactamente el orden Voigt del proyecto (ADR 0012). El orden de las columnas por nodo es \texttt{[ux, uy, uz]}, agrupado por nodo $i$ para \texttt{i \ensuremath{\in} \{0..7\}}.

\subsection{9. Rigidez elemental}

\[
\mathbf K_e = \int_\Omega \mathbf B^\top\,\mathbf D\,\mathbf B\, dV = \sum_{p=1}^{n_g} \mathbf B(\boldsymbol\xi_p)^\top\,\mathbf D_p\,\mathbf B(\boldsymbol\xi_p)\,\det\mathbf J_p\, w_p.
\]

A diferencia del 2D, no hay parámetro de espesor --- el "espesor" está incluido en el volumen integrado.

\subsection{10. Fuerzas internas}

\[
\mathbf F_{\text{int}} = \sum_{p=1}^{n_g} \mathbf B(\boldsymbol\xi_p)^\top\,\boldsymbol\sigma_p\,\det\mathbf J_p\, w_p.
\]

\subsection{11. Cuadratura}

Por defecto \textbf{Gauss--Legendre 2\ensuremath{\times}2\ensuremath{\times}2} (ocho puntos): orden de exactitud 3 en cada dirección --- integra exactamente $\mathbf K_e$ para Jacobiano constante (geometría de paralelepípedo) y captura el comportamiento dominante en geometrías irregulares moderadas. Soporta esquemas alternativos vía el parámetro \texttt{quadrature} desde \texttt{QuadratureRegistry} (e.g. \texttt{hex\_3x3x3} para problemas con materiales no lineales severos).

\textbf{Aviso sobre integración reducida}: el esquema $1 \times 1 \times 1$ (un solo punto central) está disponible pero el código emite advertencia explícita; un único punto no detecta los modos de hourglass (energía nula con desplazamientos alternados) que en 3D son \textbf{doce} modos espurios (no cuatro como en 2D). No se implementa estabilización tipo Flanagan-Belytschko. Recomendación operativa: usar el default $2 \times 2 \times 2$ salvo necesidad explícita.

\subsection{12. Cargas distribuidas consistentes}

\textbf{Fuerza de cuerpo $\mathbf b$} (e.g. peso propio $\mathbf b = (0, 0, -\rho g)$). Vector nodal equivalente:

\[
\mathbf f_e^{b} = \int_{\Omega_e} \mathbf N^\top\,\mathbf b\, dV = \sum_{p=1}^{n_g} \mathbf N(\boldsymbol\xi_p)^\top\,\mathbf b\,\det\mathbf J_p\, w_p,
\]

donde $\mathbf N$ es la matriz $3 \times 24$ de funciones de forma. Se reutiliza la cuadratura del elemento (Gauss $2 \times 2 \times 2$ por defecto). Exacto para $\mathbf b$ uniforme y geometría de paralelepípedo.

\textbf{Tracción de superficie $\bar{\mathbf t}$} sobre una cara $\Gamma_e \subset \partial\Omega_N$. Caras numeradas con normal saliente (ADR 0012):

\begin{center}
\begin{tabularx}{\textwidth}{YYY}
\hline
\textbf{Cara} & \textbf{Nodos} & \textbf{Normal saliente} \\
\hline
0 & (0, 3, 2, 1) & −\ensuremath{\zeta} (inferior) \\
1 & (4, 5, 6, 7) & +\ensuremath{\zeta} (superior) \\
2 & (0, 1, 5, 4) & −\ensuremath{\eta} (frontal) \\
3 & (1, 2, 6, 5) & +\ensuremath{\xi} (derecha) \\
4 & (2, 3, 7, 6) & +\ensuremath{\eta} (trasera) \\
5 & (3, 0, 4, 7) & −\ensuremath{\xi} (izquierda) \\
\hline
\end{tabularx}
\end{center}
\[
\mathbf f_e^{t} = \int_{\Gamma_e} \mathbf N^\top\,\bar{\mathbf t}\, dS.
\]

Cuadratura 2D sobre la cara: Gauss $2 \times 2$ (cuatro puntos por cara). Exacto para tracción constante sobre cara plana. $\bar{\mathbf t}$ se especifica en \textbf{coordenadas globales} $(t_x, t_y, t_z)$; presiones normales se obtienen multiplicando previamente por la normal exterior de la cara. Tracción variable a lo largo de la cara no soportada en este paso.

\subsection{13. Salida por punto de Gauss}

\texttt{compute\_gauss\_state(U)} devuelve un dict con $\boldsymbol\varepsilon$ y $\boldsymbol\sigma$ en cada uno de los $n_g$ puntos de Gauss del elemento (8 por defecto), junto con coordenadas naturales y globales. Habilita post-proceso fino (mapas no promediados, suavizado nodal, recovery superconvergente futuro). Por ADR 0012, \texttt{internal\_forces} devuelve \texttt{None} (sólidos no tienen fuerzas internas seccionales discretas --- son su campo $\boldsymbol\sigma(x)$ continuo).


\section{Limitaciones declaradas}

\subsection{Locking volumétrico con $\nu \to 0.5$}

\texttt{Hex8} con integración completa $2\times 2\times 2$ y material casi-incompresible (\texttt{Elastic3D} con $\nu \to 0.5$, o un futuro \texttt{VonMises3D} en régimen plástico desarrollado) sufre \textbf{locking volumétrico}: la rigidez crece artificialmente porque el espacio de funciones trilineal no contiene desplazamientos puramente isocóricos suficientes. El elemento se queda atrapado en respuestas erróneamente rígidas.

Tratamiento en este proyecto: declarar limitación, \textbf{blindar con test} (\texttt{test\_volumetric\_locking\_3d.py}, espejo del análogo 2D) que documente cuantitativamente el colapso, \textbf{no implementar mitigación} (B-bar, F-bar, mixed u-p). Política idéntica al \texttt{Quad4} (ver \texttt{STATUS.md} "Limitaciones declaradas").

\subsection{Hourglass con integración reducida}

\texttt{Hex8} integrado con un solo punto central tiene \textbf{12 modos de hourglass} (energía nula con desplazamientos alternados). Reducción disponible vía \texttt{quadrature} con warning; estabilización Flanagan-Belytschko no implementada.

\subsection{Mass lumping en orientación arbitraria}

\texttt{compute\_mass\_matrix(lumping="lumped")} usa HRZ canónico (\texttt{solidum/math/mass\_lumping.py::lump\_hrz}). Diagonal en ejes locales del elemento. En sólidos 3D la matriz lumped traslacional es estrictamente diagonal porque \texttt{m\_t\ensuremath{\cdot}I\ensuremath{_{3}}} es invariante bajo SO(3) --- sin el caveat de \texttt{Frame3D} que sí aparece en su bloque rotacional (no hay rotaciones nodales en sólidos).


\section{Contrato de implementación}

\begin{lstlisting}[language=yaml]
name: Hex8
kind: element
status: validated

interface:
  dof_names: [ux, uy, uz]
  n_nodes: 8
  strain_dim: 6                 # Voigt 3D [ε_xx, ε_yy, ε_zz, γ_xy, γ_yz, γ_xz]
  n_integration_points: 8       # default Gauss 2×2×2

parameters:
  - { name: quadrature, type: str, required: false, default: "hex_2x2x2",
      desc: "Regla de cuadratura desde QuadratureRegistry (hex_2x2x2 default, hex_3x3x3 alternativa)" }

material_contract:
  signature: "compute_state(ε, state) -> (σ, C_tangent, state')"
  strain_kind: "Voigt 3D [ε_xx, ε_yy, ε_zz, γ_xy, γ_yz, γ_xz] con γ_ij = 2·ε_ij"

conventions:
  sign: "σ > 0 ⇔ tracción (Reglas §5)"
  voigt: "[xx, yy, zz, xy, yz, xz] con γ_ij = 2·ε_ij (ADR 0012)"
  node_orientation: "convención VTK_HEXAHEDRON; nodos 0-3 cara inferior antihorario visto desde dentro, 4-7 cara superior antihorario visto desde fuera; asegura det J > 0"
  face_numbering: "ADR 0012 — 6 caras con normal saliente: 0 (−ζ), 1 (+ζ), 2 (−η), 3 (+ξ), 4 (+η), 5 (−ξ)"

validity:
  - "pequeños desplazamientos y deformaciones"
  - "geometría sin distorsión severa (det J > tol en todos los puntos de Gauss)"
  - "compatible con materiales con STRAIN_DIM = 6 (en esta etapa: Elastic3D)"

out_of_scope:
  - "grandes deformaciones, formulaciones corotacionales o lagrangianas totales"
  - "estabilización contra hourglass para integración reducida"
  - "elementos de alto orden (Hex20, Hex27 — sub-etapa posterior)"
  - "tracción variable sobre cara o presión normal"
  - "internal_forces (ADR 0012: sólidos exponen compute_gauss_state)"

acceptance:
  verification:
    - name: patch_test_3d
      setup: "8 Hex8 alrededor de un nodo interior (cubo 2×2×2 partido en
             ocho octantes) con campo lineal u_i = a_ij·x_j impuesto en
             los 26 nodos del contorno; nodo interior libre"
      expect: "nodo interior adopta el campo lineal exacto y ε constante e
              igual a 0.5·(a_ij + a_ji) en todos los puntos de Gauss; σ
              uniforme"
      tol_rel: 1.0e-10
    - name: modos_de_cuerpo_rigido
      setup: "evaluar K_e sobre un Hex8 cualquiera y proyectar sobre los 6
             modos de cuerpo rígido (3 traslaciones + 3 rotaciones)"
      expect: "K_e · u_rigid ≈ 0 en cada uno de los 6 modos"
      tol_abs: 1.0e-9
    - name: simetria_K_e
      setup: "evaluar K_e con Elastic3D"
      expect: "K_e == K_e.T exacto"
      tol_abs: 1.0e-12
  specific:
    - name: cubo_lame_3d
      setup: "cubo unitario con Elastic3D; tracción uniaxial σ_xx = p en
             cara +ξ; cara opuesta restringida en ux; restricciones
             mínimas para evitar movimiento rígido"
      expect: "u_x(L, y, z) = p·L/E; u_y = -ν·p·y/E; u_z = -ν·p·z/E
              exactos en todos los nodos"
      tol_rel: 1.0e-10
    - name: traccion_uniaxial_sobre_cara
      setup: "Hex8 con tracción constante (p, 0, 0) sobre cara 3 (+ξ)"
      expect: "Σf_x = p·A_cara, Σf_y = Σf_z = 0; reparto p·A_cara/4 a cada
              nodo de la cara para Hex8 regular"
      tol_rel: 1.0e-12
    - name: body_load_uniforme
      setup: "Hex8 regular y distorsionado con b uniforme"
      expect: "Σf = b·V_e (invariante de geometría); por simetría en Hex8
              regular cada nodo recibe b·V_e / 8"
      tol_rel: 1.0e-10
    - name: jacobiano_degenerado_abortado
      setup: "Hex8 con cuatro nodos coplanares de la cara superior sobre
             la inferior (det J ≈ 0)"
      expect: "ValueError en _compute_kinematics"
    - name: locking_volumetrico_documentado
      setup: "viga esbelta empotrada-libre con Hex8 plane strain forzado
             (cara z restringida); ν ∈ {0.3, 0.49, 0.4999}"
      expect: "deflexión vs Frame3D decrece monotonamente con ν → 0.5
              (locking documentado, no mitigado)"
      tol_rel: "documentar valores, sin tolerancia de aceptación"

references:
  - "Bathe K.J. (2014). Finite Element Procedures. §6.6 (sólidos 3D isoparamétricos)."
  - "Cook R.D., Malkus D.S., Plesha M.E., Witt R.J. (2002). Concepts and Applications of FEA. §6.8."
  - "Zienkiewicz O.C., Taylor R.L. (2005). The Finite Element Method, vol. 1, §6.7."
  - "ADR 0012 — Sólidos 3D: convención Voigt 6D y caras."

\end{lstlisting}


\section{Implementación}

\begin{itemize}
  \item \textbf{Archivo}: solidum/elements/solid\_3d/hex8.py.
  \item \textbf{Clase}: \texttt{Hex8} (registrada vía \texttt{@ElementRegistry.register}).
  \item \textbf{Funciones núcleo}: \texttt{\_compute\_kinematics\_hex8(xi, eta, zeta, coords)}, \texttt{\_det\_jacobian\_hex8(...)} y \texttt{\_shape\_functions\_hex8(...)} con \texttt{@njit} (Numba) en solidum/elements/solid\_3d/\_shared.py. Integrandos comunes 3D (\texttt{\_compute\_integrands\_3d}) y expansor masa traslacional (\texttt{\_expand\_scalar\_mass\_3d}) también en \texttt{\_shared}.
  \item \textbf{Tests}:
  \item tests/test\_solid\_3d.py --- clase \texttt{TestHex8Element} (10 tests: dimensiones/DOFs, simetría K, patch tracción uniaxial, jacobiano negativo abortado, body load, face traction (balance + nodos activos + índice fuera de rango), masa consistente y lumped HRZ, gauss state).
  \item tests/test\_rigid\_body\_modes.py --- \texttt{TestRBMHex8} (5 tests: traslación, 3 rotaciones independientes en (x, y, z), rank-deficiency = 6 modos rígidos).
  \item tests/validation/test\_cube\_lame\_3d.py --- \texttt{test\_uniaxial\_traction\_hex8\_exact}, \texttt{test\_hydrostatic\_compression\_hex8} (solución exacta a precisión máquina).
  \item tests/validation/test\_macneal\_beam\_3d.py --- shear locking documentado en malla coarse + convergencia h monótona.
  \item tests/test\_volumetric\_locking\_3d.py --- locking volumétrico declarado y blindado por test.
\end{itemize}


\section{Diálogo}

\begin{itemize}
  \item \textbf{2026-05-19} \ensuremath{\cdot} Spec inicial. Elemento espejo natural de \texttt{Quad4}; convención Voigt 6D y numeración de caras fijadas por ADR 0012. Sin variantes corotacionales ni de alto orden en esta etapa. Locking volumétrico declarado como limitación con test que lo blinda; sin mitigación implementada (política idéntica al 2D).
\end{itemize}


\newpage

\section{Tet10}



\textit{Elemento sólido 3D de orden cuadrático sobre el simplex tetraédrico. Análogo 3D del `\texttt{Tri6}\texttt{ 2D. Reproduce campos cuadráticos completos en (\ensuremath{\xi}, \ensuremath{\eta}, \ensuremath{\zeta}). Mitiga drásticamente el shear locking y el locking volumétrico del }\texttt{Tet4}\texttt{ (CST 3D). Sub-fase 3 de A.ter; subclase de }\texttt{\_HigherOrderSolid3D}`.}



\section{Especificación física}

\subsection{0. Descripción general}

Tetraedro tridimensional de segundo orden con \textbf{10 nodos}: 4 vértices + 6 medios de arista. Aproximación cuadrática $C^0$ en coordenadas baricéntricas $L_1, L_2, L_3, L_4$ con $\sum_i L_i = 1$. Formulación isoparamétrica. Régimen de pequeños desplazamientos y deformaciones.

\subsection{1-5. Cinemática, ecuación constitutiva, equilibrio}

Idénticas al `\texttt{Hex20}\texttt{/}\texttt{Hex27}` --- la diferencia con los hexaedros está exclusivamente en la geometría de referencia (simplex tetraédrico vs cubo).


\section{Formulación numérica (FEM)}

\subsection{6. Discretización --- convención VTK\_QUADRATIC\_TETRA}

Los 10 nodos sobre el tetraedro de referencia con vértices $(0,0,0)$, $(1,0,0)$, $(0,1,0)$, $(0,0,1)$:

\begin{center}
\begin{tabularx}{\textwidth}{YYY}
\hline
\textbf{Nodo} & \textbf{Tipo} & \textbf{Coords naturales $(\xi, \eta, \zeta)$} \\
\hline
0 & vértice (idéntico Tet4) & $(0, 0, 0)$ \\
1 & vértice (idéntico Tet4) & $(1, 0, 0)$ \\
2 & vértice (idéntico Tet4) & $(0, 1, 0)$ \\
3 & vértice (idéntico Tet4) & $(0, 0, 1)$ \\
4 & medio arista 0-1 & $(0.5, 0, 0)$ \\
5 & medio arista 1-2 & $(0.5, 0.5, 0)$ \\
6 & medio arista 2-0 & $(0, 0.5, 0)$ \\
7 & medio arista 0-3 & $(0, 0, 0.5)$ \\
8 & medio arista 1-3 & $(0.5, 0, 0.5)$ \\
9 & medio arista 2-3 & $(0, 0.5, 0.5)$ \\
\hline
\end{tabularx}
\end{center}
Mapeo isoparamétrico desde el tetraedro de referencia con $L_1 = 1 - \xi - \eta - \zeta$, $L_2 = \xi$, $L_3 = \eta$, $L_4 = \zeta$.

\subsection{7. Funciones de forma}

\textbf{Vértices} ($i \in \{1, 2, 3, 4\}$ --- corresponde a los nodos locales $\{0, 1, 2, 3\}$):

\[
N_i = L_i\,(2\,L_i - 1).
\]

\textbf{Medios de arista} entre vértices $i, j$:

\[
N_{ij} = 4\,L_i\,L_j.
\]

Las seis aristas se etiquetan en el orden VTK\_QUADRATIC\_TETRA: $(1,2), (2,3), (3,1), (1,4), (2,4), (3,4)$ --- nodos locales $(0,1), (1,2), (2,0), (0,3), (1,3), (2,3)$. Garantizan partición de la unidad $\sum N_i = 1$ y propiedad delta de Kronecker.

\textbf{Espacio reproducido}: el `\texttt{Tet10}` reproduce \textbf{completamente} el espacio cuadrático en barycentric coords --- todos los polinomios de grado total $\le 2$ en $(\xi, \eta, \zeta)$. 10 monomios, idéntico al número de nodos.

\subsection{8. Matriz $\mathbf B$}

Tamaño \textbf{6\ensuremath{\times}30}, ensamblada con la misma plantilla 6\ensuremath{\times}3 del `\texttt{Hex8}\texttt{/}\texttt{Hex20}\texttt{/}\texttt{Hex27}` por nodo. Derivadas en coords naturales obtenidas con regla del producto: para vértices $\partial N_i / \partial x_k = (4 L_i - 1)\,\partial L_i / \partial x_k$; para medios $\partial N_{ij}/\partial x_k = 4(L_j\,\partial L_i/\partial x_k + L_i\,\partial L_j/\partial x_k)$.

\subsection{9-10. Rigidez y fuerzas internas}

Implementadas en la base `\texttt{\_HigherOrderSolid3D}`. Tamaño 30\ensuremath{\times}30 para $\mathbf K_e$.

\subsection{11. Cuadratura}

Por defecto \textbf{Stroud 4 puntos} (`\texttt{tet\_4}`): orden de exactitud 2. Para Tet10 con Jacobiano constante el integrando de $\mathbf B^\top \mathbf D\, \mathbf B$ es a lo sumo de grado 2 (B es lineal en barycentric), así que la regla es \textbf{exacta} para K.

Para problemas con Jacobiano variable (geometría distorsionada, mid-edges fuera del centro real) el integrando puede subir de grado; en esos casos el usuario puede seleccionar \texttt{tet\_15} (Keast 15 puntos, orden 5) vía el parámetro \texttt{quadrature}.

\textbf{Cuadratura de masa}: el `\texttt{Tet10}\texttt{ declara }\texttt{\_MASS\_QUADRATURE = "tet\_15"}\texttt{ (Keast 15 puntos, orden 5) **independientemente** de la cuadratura de K. El integrando @@MINL38@@ es de grado 4 (cuadrático \ensuremath{\times} cuadrático) y la regla de 4 puntos lo subintegra; }\texttt{tet\_15}\texttt{ lo integra exactamente. Esto garantiza que la masa total sea correcta y que análisis modal/transitorio con }Tet10` no sufran errores de cuadratura.

\subsection{12. Cargas distribuidas consistentes}

\textbf{Body load}: `\texttt{\ensuremath{\int} N\textasciicircum{}T b dV}\texttt{ con la cuadratura del elemento. Reparto no uniforme entre vértices y medios; suma global }\texttt{b\ensuremath{\cdot}V\_e}\texttt{ exacta para }\texttt{b}` uniforme y geometría sin distorsión severa.

\textbf{Face traction}: cada cara del `\texttt{Tet10}\texttt{ tiene **6 nodos** (3 vértices + 3 medios de arista) y se aproxima por funciones cuadráticas Tri6 sobre el triángulo de referencia @@MINL39@@-@@MINL40@@-@@MINL41@@. Cuadratura }\texttt{tri\_3}` (3 puntos, orden 2) --- exacta para tracción uniforme sobre cara plana.

\begin{center}
\begin{tabularx}{\textwidth}{YYYY}
\hline
\textbf{Cara} & \textbf{Vértices Tet10} & \textbf{Medios Tet10 (en orden Tri6)} & \textbf{Normal saliente} \\
\hline
0 & (1, 2, 3) & (5, 9, 8) & opuesta al nodo 0 \\
1 & (0, 3, 2) & (7, 9, 6) & opuesta al nodo 1 \\
2 & (0, 1, 3) & (4, 8, 7) & opuesta al nodo 2 \\
3 & (0, 2, 1) & (6, 5, 4) & opuesta al nodo 3 \\
\hline
\end{tabularx}
\end{center}
El orden dentro de cada cara reproduce la convención de `\texttt{\_N\_tri6}\texttt{ del paquete 2D: vértices }(0, 1, 2)\texttt{ antihorarios, luego medios }(0-1, 1-2, 2-0)`.

\subsection{13. Salida por punto de Gauss}

\texttt{compute\_gauss\_state(U)} con 4 puntos por defecto. Sin \texttt{internal\_forces} (ADR 0012).


\section{Limitaciones declaradas}

\subsection{Locking volumétrico atenuado pero presente}

\texttt{Tet10} con cuadratura `\texttt{tet\_4}\texttt{ y material casi-incompresible (}Elastic3D\texttt{ con @@MINL42@@, o }VonMises3D\texttt{ plástico desarrollado) sufre **locking volumétrico** mucho menor que el }Tet4\texttt{ (cuyo shear locking severo es la motivación principal del }Tet10\texttt{), pero todavía presente. Sin mitigación implementada. Blindado por test (mismo patrón que }Hex20\texttt{/}Hex27`).

\subsection{Distorsión severa}

Para Tet10 con mid-edges desplazados fuera del centro real de la arista (mid-edges curvos para representar geometría curva), el Jacobiano deja de ser constante y la cuadratura `\texttt{tet\_4}\texttt{ puede subintegrar K. Recomendación: usar }\texttt{tet\_15}` en mallas con curvatura severa.


\section{Contrato de implementación}

\begin{lstlisting}[language=yaml]
name: Tet10
kind: element
status: validated

interface:
  dof_names: [ux, uy, uz]
  n_nodes: 10
  strain_dim: 6
  n_integration_points: 4       # default Stroud tet_4

parameters:
  - { name: quadrature, type: str, required: false, default: "tet_4",
      desc: "Cuadratura tetraédrica desde QuadratureRegistry. Default
             tet_4 (Stroud 4 puntos, orden 2, suficiente para K con J
             constante). Alternativa tet_15 (Keast 15 puntos, orden 5)
             para geometrías distorsionadas." }

material_contract:
  signature: "compute_state(ε, state) -> (σ, C_tangent, state')"
  strain_kind: "Voigt 3D [ε_xx, ε_yy, ε_zz, γ_xy, γ_yz, γ_xz] con γ_ij = 2·ε_ij"

conventions:
  sign: "σ > 0 ⇔ tracción (Reglas §5)"
  voigt: "[xx, yy, zz, xy, yz, xz] con γ_ij = 2·ε_ij (ADR 0012)"
  node_orientation: "VTK_QUADRATIC_TETRA; 0-3 vértices (Tet4 con volumen
                     positivo), 4-9 medios de arista en orden (0-1, 1-2,
                     2-0, 0-3, 1-3, 2-3)"
  face_numbering: "ADR 0012 — 4 caras triangulares con normal saliente,
                   paritarias con Tet4 en los vértices; cada cara añade
                   3 medios de arista en orden Tri6"

validity:
  - "pequeños desplazamientos y deformaciones"
  - "geometría sin distorsión severa (det J > tol en todos los Gauss)"
  - "compatible con materiales con STRAIN_DIM = 6"

out_of_scope:
  - "grandes deformaciones, formulaciones corotacionales"
  - "internal_forces (ADR 0012)"
  - "mid-edges severamente curvos (usar tet_15 si es necesario)"

acceptance:
  verification:
    - name: patch_test_lineal
      setup: "campo lineal u_i = a_ij·x_j en los 10 nodos del contorno"
      expect: "ε constante = sym(a) en todos los Gauss"
      tol_rel: 1.0e-12
    - name: patch_test_cuadratico
      setup: "campo u_x = c·x² en los 10 nodos"
      expect: "ε_xx(x) = 2c·x exacto en todos los Gauss"
      tol_rel: 1.0e-12
    - name: modos_de_cuerpo_rigido
      setup: "K_e proyectado sobre 6 modos rígidos 3D"
      expect: "exactamente 6 autovalores nulos"
      tol_abs: 1.0e-9
    - name: simetria_K_e
      setup: "K_e con Elastic3D"
      expect: "K_e == K_e.T exacto"
      tol_abs: 1.0e-12
  specific:
    - name: cubo_lame_3d_uniaxial
      setup: "malla 5 tetraedros del cubo unitario, refinados a Tet10,
              con tracción uniaxial; el campo analítico es lineal y
              Tet10 lo reproduce exacto"
      expect: "u_x = p·x/E etc. exacto a precisión máquina"
      tol_rel: 1.0e-10
    - name: macneal_3d_tet10_dominates_tet4
      setup: "cantilever de Tet10 con la misma malla que Tet4 sufre
              dramáticamente menos shear locking"
      expect: "Tet10 alcanza > 80% u_EB con ~12 tetraedros (Tet4 << 20%)"
    - name: body_load_uniforme
      setup: "Tet10 ref con b = (0, 0, -g)"
      expect: "Σf_z = -g·V_e = -g/6"
      tol_rel: 1.0e-12
    - name: face_traction_balance
      setup: "tracción uniforme sobre una cara del Tet10"
      expect: "Σf = A_face · t̄ exacto; nodos fuera de la cara reciben 0"
      tol_rel: 1.0e-12
    - name: mass_consistent_total_with_tet15
      setup: "Tet10 cubo ref con ρ=7850, masa consistente con tet_15"
      expect: "Σ M = 3·ρ·V_e = 3·ρ/6"
      tol_rel: 1.0e-9
    - name: mass_lumped_HRZ_positivo
      setup: "masa lumped HRZ del Tet10"
      expect: "diagonal estricta, todas las entradas positivas,
              conservación de masa total"
      tol_abs: 1.0e-12

references:
  - "Bathe K.J. (2014). Finite Element Procedures. §5.3.3 (tetraedro 3D)."
  - "Cook R.D., Malkus D.S., Plesha M.E., Witt R.J. (2002). Concepts and
    Applications of FEA. §6.5."
  - "Keast P. (1986). Moderate-degree tetrahedral quadrature formulas.
    CMAME 55(3), 339-348 (regla tet_15 de masa)."
  - "Stroud A.H. (1971). Approximate calculation of multiple integrals
    (regla tet_4 de K)."
  - "ADR 0012 — Sólidos 3D: convención Voigt 6D y caras."

\end{lstlisting}


\section{Implementación}

\begin{itemize}
  \item \textbf{Archivo}: solidum/elements/solid\_3d/tet10.py.
  \item \textbf{Clase}: \texttt{Tet10} (registrada vía \texttt{@ElementRegistry.register}); subclase de \texttt{\_HigherOrderSolid3D}.
  \item \textbf{Funciones núcleo}: \texttt{\_N\_tet10(xi, eta, zeta)}, \texttt{\_dN\_tet10(xi, eta, zeta)} en solidum/elements/solid\_3d/\_shared.py. Implementadas en numpy puro vía coordenadas baricéntricas $L_i$.
  \item \textbf{Cuadraturas registradas}: \texttt{tet\_4} (Stroud, orden 2) y \texttt{tet\_15} (Keast, orden 5) en solidum/math/integration.py. El \texttt{tet\_4} cubre la integración de K; el \texttt{tet\_15} se usa exclusivamente para la masa consistente (declarada en \texttt{Tet10.\_MASS\_QUADRATURE}).
  \item \textbf{Tests}:
  \item \texttt{tests/test\_solid\_3d\_higher\_order.py} --- \texttt{TestTet10Element} (patch lineal + cuadrático, simetría K, body load, face traction, masa consistente con tet\_15 y lumped HRZ, gauss state).
  \item \texttt{tests/test\_rigid\_body\_modes.py} --- \texttt{TestRBMTet10} (3 trans + 3 rot + rank-deficiency = 6).
  \item \texttt{tests/validation/test\_cube\_lame\_3d.py} --- extensión con malla Tet10 del cubo unitario; uniaxial e hidrostática exactas a precisión máquina.
\end{itemize}


\section{Diálogo}

\begin{itemize}
  \item \textbf{2026-05-27} \ensuremath{\cdot} Spec inicial. Sub-fase 3 de A.ter --- cierra el catálogo cuadrático 3D con la rama tetraédrica. Subclase de \texttt{\_HigherOrderSolid3D} igual que \texttt{Hex20}/\texttt{Hex27}, pero con cara triangular (\texttt{Tri6} shape functions importadas del paquete 2D, cuadratura \texttt{tri\_3}) y cuadratura volumétrica tetraédrica nueva (\texttt{tet\_4} + \texttt{tet\_15}).
  \item \textbf{2026-05-27} \ensuremath{\cdot} Implementado y validado. Patch test lineal y cuadrático exactos a precisión máquina, 6 modos rígidos, body load suma exacta \texttt{b\ensuremath{\cdot}V\_e}, face traction balance exacto sobre la cara triangular, masa consistente con \texttt{tet\_15} exacta a precisión máquina (conservación \texttt{3\ensuremath{\cdot}\ensuremath{\rho}\ensuremath{\cdot}V}), lumped HRZ con todas las entradas positivas. Spec a \texttt{status: validated}.
\end{itemize}


\newpage

\section{Hex20}



\textit{Elemento sólido 3D de orden cuadrático con interpolación serendípita.}

\textit{Reproduce campos cuadráticos completos en 3D (todos los polinomios de grado total \ensuremath{\le} 2 más algunos de orden 3 sin los términos triquadráticos). Es al \texttt{Hex8} lo que el \texttt{Quad8} al \texttt{Quad4}: mucho más preciso en flexión y geometrías curvas con el mismo número de elementos. Primer elemento de la sub-etapa A.ter.}



\section{Especificación física}

\subsection{0. Descripción general}

Hexaedro tridimensional de segundo orden con \textbf{20 nodos}: 8 vértices + 12 medios de arista, sin nodos de cara ni centroide. Aproximación serendípita $C^0$ en coordenadas naturales $(\xi, \eta, \zeta) \in [-1, 1]^3$. Formulación isoparamétrica: la geometría y los desplazamientos comparten las mismas funciones de forma. Régimen de pequeños desplazamientos y deformaciones.

\subsection{1. Cinemática de desplazamientos}

\[
u_x(\xi, \eta, \zeta) = \sum_{i=0}^{19} N_i(\xi, \eta, \zeta)\, u_x^{(i)}, \qquad u_y = \sum N_i\, u_y^{(i)}, \qquad u_z = \sum N_i\, u_z^{(i)}.
\]

\subsection{2. Cinemática de deformaciones}

Tensor infinitesimal en notación Voigt 6D del proyecto (ADR 0012):

\[
\boldsymbol\varepsilon = [\varepsilon_{xx},\ \varepsilon_{yy},\ \varepsilon_{zz},\ \gamma_{xy},\ \gamma_{yz},\ \gamma_{xz}]^\top
\]

con $\gamma_{ij} = 2\,\varepsilon_{ij}$ \textit{engineering}. Idéntica al \texttt{Hex8}.

\subsection{3. Ecuación constitutiva}

Delegada al material 3D (\texttt{STRAIN\_DIM = 6}). El elemento llama \texttt{material.compute\_state(\ensuremath{\varepsilon}, state) \ensuremath{\to} (\ensuremath{\sigma}, C\_tangent, state')} por punto de Gauss. Compatible con todo el catálogo 3D actual: \texttt{Elastic3D}, \texttt{VonMises3D}, \texttt{DruckerPrager3D}, \texttt{IsotropicDamage3D}.

\subsection{4-5. Equilibrio}

Forma fuerte y débil idénticas al \texttt{Hex8} --- la diferencia entre \texttt{Hex8} y \texttt{Hex20} está exclusivamente en el espacio de aproximación discreto.


\section{Formulación numérica (FEM)}

\subsection{6. Discretización}

Veinte nodos en orden VTK\_QUADRATIC\_HEXAHEDRON (estándar gmsh/VTK/Abaqus). En el cubo de referencia $[-1, 1]^3$:

\begin{center}
\begin{tabularx}{\textwidth}{YYY}
\hline
\textbf{Nodo local} & \textbf{$(\xi_i, \eta_i, \zeta_i)$} & \textbf{Tipo} \\
\hline
0 & $(-1, -1, -1)$ & vértice \\
1 & $(+1, -1, -1)$ & vértice \\
2 & $(+1, +1, -1)$ & vértice \\
3 & $(-1, +1, -1)$ & vértice \\
4 & $(-1, -1, +1)$ & vértice \\
5 & $(+1, -1, +1)$ & vértice \\
6 & $(+1, +1, +1)$ & vértice \\
7 & $(-1, +1, +1)$ & vértice \\
8 & $(0, -1, -1)$ & medio arista 0-1 (eje \ensuremath{\xi}, cara inferior) \\
9 & $(+1, 0, -1)$ & medio arista 1-2 (eje \ensuremath{\eta}, cara inferior) \\
10 & $(0, +1, -1)$ & medio arista 2-3 (eje \ensuremath{\xi}, cara inferior) \\
11 & $(-1, 0, -1)$ & medio arista 3-0 (eje \ensuremath{\eta}, cara inferior) \\
12 & $(0, -1, +1)$ & medio arista 4-5 (eje \ensuremath{\xi}, cara superior) \\
13 & $(+1, 0, +1)$ & medio arista 5-6 (eje \ensuremath{\eta}, cara superior) \\
14 & $(0, +1, +1)$ & medio arista 6-7 (eje \ensuremath{\xi}, cara superior) \\
15 & $(-1, 0, +1)$ & medio arista 7-4 (eje \ensuremath{\eta}, cara superior) \\
16 & $(-1, -1, 0)$ & medio arista 0-4 (eje \ensuremath{\zeta}, vertical) \\
17 & $(+1, -1, 0)$ & medio arista 1-5 (eje \ensuremath{\zeta}, vertical) \\
18 & $(+1, +1, 0)$ & medio arista 2-6 (eje \ensuremath{\zeta}, vertical) \\
19 & $(-1, +1, 0)$ & medio arista 3-7 (eje \ensuremath{\zeta}, vertical) \\
\hline
\end{tabularx}
\end{center}
Mapeo isoparamétrico desde el cubo de referencia:

\[
x(\xi, \eta, \zeta) = \sum_{i=0}^{19} N_i\, x^{(i)},\quad y = \sum N_i\, y^{(i)},\quad z = \sum N_i\, z^{(i)}.
\]

Jacobiano $\mathbf J = \partial(x, y, z) / \partial(\xi, \eta, \zeta)$ es una matriz 3\ensuremath{\times}3 con determinante $\det\mathbf J$. El código aborta con \texttt{ValueError} si $\det\mathbf J \le \text{tol}$ en cualquier punto de Gauss (elemento severamente distorsionado o nodos en orden invertido).

\subsection{7. Funciones de forma serendípitas}

\textbf{Vértices} ($i \in \{0,\dots,7\}$, $(\xi_i, \eta_i, \zeta_i) \in \{\pm 1\}^3$):

\[
N_i = \frac{1}{8}(1 + \xi_i\xi)(1 + \eta_i\eta)(1 + \zeta_i\zeta)(\xi_i\xi + \eta_i\eta + \zeta_i\zeta - 2).
\]

El multiplicador $(\xi_i\xi + \eta_i\eta + \zeta_i\zeta - 2)$ vale $1$ en el propio vértice y $0$ en los demás vértices y en todos los medios de arista --- garantía $\delta_{ij}$ del esquema.

\textbf{Medios de arista paralelos a $\xi$} (nodos 8, 10, 12, 14; $\xi_i = 0$, $\eta_i, \zeta_i \in \{\pm 1\}$):

\[
N_i = \frac{1}{4}(1 - \xi^2)(1 + \eta_i\eta)(1 + \zeta_i\zeta).
\]

\textbf{Medios de arista paralelos a $\eta$} (nodos 9, 11, 13, 15; $\eta_i = 0$):

\[
N_i = \frac{1}{4}(1 + \xi_i\xi)(1 - \eta^2)(1 + \zeta_i\zeta).
\]

\textbf{Medios de arista paralelos a $\zeta$} (nodos 16, 17, 18, 19; $\zeta_i = 0$):

\[
N_i = \frac{1}{4}(1 + \xi_i\xi)(1 + \eta_i\eta)(1 - \zeta^2).
\]

Espacio reproducido: $\{1, \xi, \eta, \zeta, \xi^2, \eta^2, \zeta^2, \xi\eta, \eta\zeta, \xi\zeta, \xi^2\eta, \xi\eta^2, \eta^2\zeta, \eta\zeta^2, \xi^2\zeta, \xi\zeta^2, \xi\eta\zeta, \xi^2\eta\zeta, \xi\eta^2\zeta, \xi\eta\zeta^2\}$ --- 20 monomios, omitiendo los puramente cúbicos ($\xi^3$, etc.) y el término $\xi^2\eta^2\zeta^2$.

\subsection{8. Matriz deformación-desplazamiento}

Las derivadas en globales se obtienen vía $\partial N_i / \partial \mathbf x = \mathbf J^{-1}\,\partial N_i / \partial \boldsymbol\xi$. La matriz $\mathbf B$ tiene tamaño \textbf{6\ensuremath{\times}60} y se ensambla por nodo $i \in \{0,\dots,19\}$ con la misma plantilla 6\ensuremath{\times}3 del \texttt{Hex8}:

\[
\mathbf B_i = \begin{bmatrix}
\partial N_i / \partial x & 0 & 0 \\
0 & \partial N_i / \partial y & 0 \\
0 & 0 & \partial N_i / \partial z \\
\partial N_i / \partial y & \partial N_i / \partial x & 0 \\
0 & \partial N_i / \partial z & \partial N_i / \partial y \\
\partial N_i / \partial z & 0 & \partial N_i / \partial x
\end{bmatrix}, \qquad \boldsymbol\varepsilon = \mathbf B\,\mathbf u_e.
\]

El orden de las filas reproduce exactamente el orden Voigt del proyecto (ADR 0012). El orden de las columnas por nodo es \texttt{[ux, uy, uz]}, agrupado por nodo $i$ para \texttt{i \ensuremath{\in} \{0..19\}}.

\textbf{Derivadas explícitas} (cerradas vía regla del producto):

Vértices ($i \in \{0,\dots,7\}$, con $u = \xi_i\xi$, $v = \eta_i\eta$, $w = \zeta_i\zeta$):

\[
\frac{\partial N_i}{\partial \xi} = \frac{\xi_i}{8}(1 + v)(1 + w)(2u + v + w - 1),
\]

\[
\frac{\partial N_i}{\partial \eta} = \frac{\eta_i}{8}(1 + u)(1 + w)(u + 2v + w - 1),
\]

\[
\frac{\partial N_i}{\partial \zeta} = \frac{\zeta_i}{8}(1 + u)(1 + v)(u + v + 2w - 1).
\]

Medio paralelo a $\xi$ (nodos 8, 10, 12, 14):

\[
\frac{\partial N_i}{\partial \xi} = -\frac{\xi}{2}(1 + \eta_i\eta)(1 + \zeta_i\zeta),\quad \frac{\partial N_i}{\partial \eta} = \frac{\eta_i}{4}(1 - \xi^2)(1 + \zeta_i\zeta),\quad \frac{\partial N_i}{\partial \zeta} = \frac{\zeta_i}{4}(1 - \xi^2)(1 + \eta_i\eta).
\]

Medio paralelo a $\eta$ (nodos 9, 11, 13, 15):

\[
\frac{\partial N_i}{\partial \xi} = \frac{\xi_i}{4}(1 - \eta^2)(1 + \zeta_i\zeta),\quad \frac{\partial N_i}{\partial \eta} = -\frac{\eta}{2}(1 + \xi_i\xi)(1 + \zeta_i\zeta),\quad \frac{\partial N_i}{\partial \zeta} = \frac{\zeta_i}{4}(1 + \xi_i\xi)(1 - \eta^2).
\]

Medio paralelo a $\zeta$ (nodos 16, 17, 18, 19):

\[
\frac{\partial N_i}{\partial \xi} = \frac{\xi_i}{4}(1 + \eta_i\eta)(1 - \zeta^2),\quad \frac{\partial N_i}{\partial \eta} = \frac{\eta_i}{4}(1 + \xi_i\xi)(1 - \zeta^2),\quad \frac{\partial N_i}{\partial \zeta} = -\frac{\zeta}{2}(1 + \xi_i\xi)(1 + \eta_i\eta).
\]

\subsection{9. Rigidez elemental}

\[
\mathbf K_e = \int_\Omega \mathbf B^\top\,\mathbf D\,\mathbf B\, dV = \sum_{p=1}^{n_g} \mathbf B(\boldsymbol\xi_p)^\top\,\mathbf D_p\,\mathbf B(\boldsymbol\xi_p)\,\det\mathbf J_p\, w_p.
\]

Sin parámetro de espesor --- el volumen ya está incluido en $\det\mathbf J_p\cdot w_p$.

\subsection{10. Fuerzas internas}

\[
\mathbf F_{\text{int}} = \sum_{p=1}^{n_g} \mathbf B(\boldsymbol\xi_p)^\top\,\boldsymbol\sigma_p\,\det\mathbf J_p\, w_p.
\]

\subsection{11. Cuadratura}

Por defecto \textbf{Gauss--Legendre 3\ensuremath{\times}3\ensuremath{\times}3} (27 puntos): orden de exactitud 5 en cada dirección. Suficiente para integrar exactamente $\mathbf B^\top \mathbf D\, \mathbf B$ del \texttt{Hex20} sobre Jacobiano constante (el integrando es a lo sumo de grado 3 en cada dirección) y para integrar exactamente la masa consistente $\rho \mathbf N^\top \mathbf N$ (grado 4 en cada dirección).

\textbf{Aviso sobre integración reducida}: el esquema $2 \times 2 \times 2$ (\texttt{hex\_2x2x2}, 8 puntos) está disponible y reduce el coste computacional, pero introduce \textbf{6 modos de hourglass espurios} sobre un elemento aislado (verificación numérica directa del proyecto; Belytschko-Liu-Moran §8.4.4). Estos modos \textbf{no propagan} a través de mallas ensambladas con condiciones de Dirichlet razonables --- la rigidez global ensamblada queda rank-sufficient en la mayoría de casos prácticos, pero un elemento aislado o una capa con BC pobres puede mostrar hourglass. El código permite seleccionar la cuadratura reducida vía el parámetro \texttt{quadrature} pero \textbf{no estabiliza} los modos espurios. Recomendación operativa: usar el default $3 \times 3 \times 3$ salvo necesidad experimental. La matriz de masa se integra con $3 \times 3 \times 3$ independientemente de la cuadratura elegida para $\mathbf K$, para evitar subintegrarla.

\subsection{12. Cargas distribuidas consistentes}

\textbf{Fuerza de cuerpo $\mathbf b$}. Vector nodal equivalente:

\[
\mathbf f_e^{b} = \int_{\Omega_e} \mathbf N^\top\,\mathbf b\, dV = \sum_{p=1}^{n_g} \mathbf N(\boldsymbol\xi_p)^\top\,\mathbf b\,\det\mathbf J_p\, w_p,
\]

donde $\mathbf N$ es la matriz $3 \times 60$ de funciones de forma. Se reutiliza la cuadratura del elemento. Para $\mathbf b$ uniforme y geometría regular el resultado es exacto. \textbf{Para \texttt{Hex20} el reparto no es uniforme entre nodos}: la fórmula serendípita asigna pesos distintos a vértices y medios (los medios reciben más carga total que los vértices, con suma global $\rho\cdot V_e$ correcta).

\textbf{Tracción de superficie $\bar{\mathbf t}$} sobre una cara. Cada cara del \texttt{Hex20} tiene \textbf{8 nodos} (4 vértices + 4 medios de arista) y se aproxima por funciones serendípitas Quad8 sobre el parámetro 2D $(s, t) \in [-1, 1]^2$. La numeración de caras conserva la convención del \texttt{Hex8} (ADR 0012) extendida con los medios:

\begin{center}
\begin{tabularx}{\textwidth}{YYYY}
\hline
\textbf{Cara} & \textbf{Vértices} & \textbf{Medios de arista} & \textbf{Normal saliente} \\
\hline
0 & (0,3,2,1) & (11, 10, 9, 8) & −\ensuremath{\zeta} (inferior) \\
1 & (4,5,6,7) & (12, 13, 14, 15) & +\ensuremath{\zeta} (superior) \\
2 & (0,1,5,4) & (8, 17, 12, 16) & −\ensuremath{\eta} (frontal) \\
3 & (1,2,6,5) & (9, 18, 13, 17) & +\ensuremath{\xi} (derecha) \\
4 & (2,3,7,6) & (10, 19, 14, 18) & +\ensuremath{\eta} (trasera) \\
5 & (3,0,4,7) & (11, 16, 15, 19) & −\ensuremath{\xi} (izquierda) \\
\hline
\end{tabularx}
\end{center}
El orden dentro de cada cara reproduce la convención VTK\_QUADRATIC\_HEXAHEDRON: primero los 4 vértices, luego los 4 medios en el mismo orden cíclico (medio entre vértice $k$ y vértice $k+1$).

\[
\mathbf f_e^{t} = \int_{\Gamma_e} \mathbf N^\top\,\bar{\mathbf t}\, dS.
\]

Cuadratura 2D sobre la cara: Gauss $3 \times 3$ (nueve puntos) --- exacta para tracción constante sobre cara plana (integrando degree 2 en $(s,t)$). $\bar{\mathbf t}$ se especifica en \textbf{coordenadas globales} $(t_x, t_y, t_z)$; presiones normales se obtienen multiplicando previamente por la normal exterior de la cara. Tracción variable sobre la cara no soportada en este paso. \textbf{Para \texttt{Hex20} el reparto entre vértices y medios no es trivial} (la integral cuadrática de Quad8 sobre cara plana asigna fracciones específicas: los vértices reciben $-1/12$ de la fuerza total cada uno y los medios $4/12$ cada uno, con suma global exacta --- ver test de balance).

\subsection{13. Salida por punto de Gauss}

\texttt{compute\_gauss\_state(U)} devuelve un dict con $\boldsymbol\varepsilon$ y $\boldsymbol\sigma$ en cada uno de los $n_g$ puntos de Gauss del elemento (27 por defecto), junto con coordenadas naturales y globales. Habilita post-proceso fino (mapas no promediados, suavizado nodal, recovery superconvergente futuro). Por ADR 0012, \texttt{internal\_forces} devuelve \texttt{None} (sólidos no tienen fuerzas internas seccionales discretas).


\section{Limitaciones declaradas}

\subsection{Locking volumétrico con $\nu \to 0.5$}

\texttt{Hex20} con integración completa $3\times 3\times 3$ y material casi-incompresible (\texttt{Elastic3D} con $\nu \to 0.5$, o \texttt{VonMises3D} en régimen plástico desarrollado) sufre \textbf{locking volumétrico} menos severo que el \texttt{Hex8} pero todavía presente: el espacio serendípito triquadrático tampoco contiene desplazamientos puramente isocóricos suficientes en mallas distorsionadas (Cook-Malkus-Plesha §13.6). 

Tratamiento en este proyecto: declarar limitación, \textbf{blindar con test cuantitativo} (\texttt{test\_volumetric\_locking\_hex20.py}, espejo del análogo \texttt{Hex8}) que documente la mitigación parcial respecto al lineal, \textbf{no implementar B-bar/F-bar/mixed u-p}. Política idéntica a \texttt{Quad4}/\texttt{Quad8} y \texttt{Hex8} (ver \texttt{STATUS.md} "Limitaciones declaradas").

\subsection{Hourglass con integración reducida}

\texttt{Hex20} integrado con \texttt{hex\_2x2x2} tiene \textbf{6 modos de hourglass espurios} por elemento aislado (vs 12 del \texttt{Hex8} reducido). En mallas ensambladas con Dirichlet razonable los modos suelen no propagar --- la rigidez global queda rank-sufficient --- pero un elemento aislado o una capa con BC pobres lo mostraría. Reducción disponible vía \texttt{quadrature}; estabilización no implementada. Modo recomendado: full integration $3\times 3\times 3$.

\subsection{Mass lumping}

\texttt{compute\_mass\_matrix(lumping="lumped")} usa HRZ canónico (\texttt{solidum/math/mass\_lumping.py::lump\_hrz}). Para \texttt{Hex20} (con nodos de medio de arista) HRZ es la única opción razonable: row-sum produciría masas \textbf{negativas} en los vértices (los pesos de Gauss son negativos para los vértices en la cuadratura serendípita Lobatto equivalente), inadmisible para análisis dinámico (Bathe FEP §9.2.4). HRZ preserva masa total escalando la diagonal consistente por un factor común.


\section{Contrato de implementación}

\begin{lstlisting}[language=yaml]
name: Hex20
kind: element
status: validated

interface:
  dof_names: [ux, uy, uz]
  n_nodes: 20
  strain_dim: 6                 # Voigt 3D [ε_xx, ε_yy, ε_zz, γ_xy, γ_yz, γ_xz]
  n_integration_points: 27      # default Gauss 3×3×3

parameters:
  - { name: quadrature, type: str, required: false, default: "hex_3x3x3",
      desc: "Regla de cuadratura desde QuadratureRegistry (hex_3x3x3 default,
             hex_2x2x2 reducida con 4 modos de hourglass espurios)" }

material_contract:
  signature: "compute_state(ε, state) -> (σ, C_tangent, state')"
  strain_kind: "Voigt 3D [ε_xx, ε_yy, ε_zz, γ_xy, γ_yz, γ_xz] con γ_ij = 2·ε_ij"

conventions:
  sign: "σ > 0 ⇔ tracción (Reglas §5)"
  voigt: "[xx, yy, zz, xy, yz, xz] con γ_ij = 2·ε_ij (ADR 0012)"
  node_orientation: "convención VTK_QUADRATIC_HEXAHEDRON; vértices 0-7 idénticos
                     al Hex8; medios 8-11 bottom edges, 12-15 top edges, 16-19
                     vertical edges; asegura det J > 0 en todos los Gauss"
  face_numbering: "ADR 0012 — 6 caras con normal saliente, idéntica a Hex8 en
                   los vértices; cada cara añade 4 nodos medios en el mismo
                   orden cíclico (Quad8 serendipity)"

validity:
  - "pequeños desplazamientos y deformaciones"
  - "geometría sin distorsión severa (det J > tol en todos los puntos de Gauss)"
  - "compatible con materiales con STRAIN_DIM = 6 (Elastic3D, VonMises3D,
     DruckerPrager3D, IsotropicDamage3D)"

out_of_scope:
  - "grandes deformaciones, formulaciones corotacionales o lagrangianas totales"
  - "estabilización contra hourglass para integración reducida hex_2x2x2"
  - "tracción variable sobre cara o presión normal con gradiente"
  - "internal_forces (ADR 0012: sólidos exponen compute_gauss_state)"
  - "Hex27 lagrangiano (sub-fase 2 de A.ter)"

acceptance:
  verification:
    - name: patch_test_lineal
      setup: "Hex20 con campo lineal u_i = a_ij·x_j impuesto en los 20 nodos
             del contorno; sin nodos interiores libres (todos están en la
             frontera de un elemento aislado, pero al menos los medios no
             están fijados en su posición referencial)"
      expect: "ε constante e igual a sym(a) en todos los puntos de Gauss; σ
              uniforme"
      tol_rel: 1.0e-12
    - name: patch_test_cuadratico
      setup: "Hex20 con campo u_x = c·x², u_y = u_z = 0 impuesto en los 20 nodos
             del contorno (todos)"
      expect: "ε_xx(x) = 2c·x exacto en todos los puntos de Gauss"
      tol_rel: 1.0e-12
    - name: modos_de_cuerpo_rigido
      setup: "K_e del Hex20 proyectado sobre los 6 modos de cuerpo rígido"
      expect: "K_e · u_rigid ≈ 0 (autovalores ≈ 0: exactamente 6)"
      tol_abs: 1.0e-9
    - name: simetria_K_e
      setup: "evaluar K_e con Elastic3D"
      expect: "K_e == K_e.T exacto"
      tol_abs: 1.0e-12
  specific:
    - name: cubo_lame_3d
      setup: "cubo unitario con un Hex20; tracción uniaxial σ_xx = p en cara +ξ;
             cara opuesta restringida en ux; restricciones mínimas para evitar
             movimiento rígido"
      expect: "u_x(L, y, z) = p·L/E; u_y = −ν·p·y/E; u_z = −ν·p·z/E exactos en
              todos los nodos"
      tol_rel: 1.0e-10
    - name: body_load_uniforme
      setup: "Hex20 regular cubo unitario con b = (0, 0, −g)"
      expect: "Σf_z = −g·V_e = −g; reparto no uniforme entre vértices y medios"
      tol_rel: 1.0e-12
    - name: face_traction_balance
      setup: "Hex20 cubo unitario con tracción (p, 0, 0) sobre cara +ξ (face 3)"
      expect: "Σf_x = p·A_cara = p·1; vértices reciben −p/12 cada uno y medios
              4p/12 cada uno (suma = 4·(−1/12) + 4·(4/12) = 12/12 = 1, OK)"
      tol_rel: 1.0e-12
    - name: jacobiano_degenerado_abortado
      setup: "Hex20 con un nodo medio desplazado fuera del rango cuasi-recto"
      expect: "ValueError en _compute_kinematics_hex20 cuando det J ≤ tol"
    - name: convergencia_cantilever_macneal
      setup: "cantilever esbelto L/h = 30 con 1×1 sección + N elementos a lo
             largo; flecha vs Euler-Bernoulli"
      expect: "Hex20 con 1×1×6 alcanza >0.9 · u_EB (Hex8 con 1×1×6 solo ≈ 0.5);
              convergencia h monotónica al refinar"
      tol_rel: "documentar — sin tolerancia estricta"
    - name: locking_volumetrico_atenuado
      setup: "cantilever esbelto con Hex20; ν ∈ {0.3, 0.4999}"
      expect: "ratio u_tip(ν=0.4999) / u_tip(ν=0.3) significativamente mayor
              que el del Hex8 (Hex20 < 0.6 → Hex20 > 0.8 esperado al menos)"
      tol_rel: "documentar valores; cota relativa Hex20 vs Hex8"
    - name: mass_consistent_total
      setup: "Hex20 cubo unitario con ρ = 7850, masa consistente"
      expect: "Σ M (todas las entradas) = 3·ρ·V = 3·7850·1 (conservación 3D)"
      tol_rel: 1.0e-9
    - name: mass_lumped_total_HRZ
      setup: "Hex20 cubo unitario con ρ = 7850, masa lumped HRZ"
      expect: "diag(M) suma = 3·ρ·V; matriz estrictamente diagonal; todas las
              entradas estrictamente positivas (HRZ canónico, no row-sum)"
      tol_abs: 1.0e-12

references:
  - "Bathe K.J. (2014). Finite Element Procedures. §5.3.2 (serendipity 3D)."
  - "Cook R.D., Malkus D.S., Plesha M.E., Witt R.J. (2002). Concepts and
    Applications of FEA. §6.5, §13.6 (locking volumétrico Hex20)."
  - "Zienkiewicz O.C., Taylor R.L. (2005). The Finite Element Method, vol. 1,
    §8.7 (familia serendípita 3D)."
  - "ADR 0012 — Sólidos 3D: convención Voigt 6D y caras."

\end{lstlisting}


\section{Implementación}

\begin{itemize}
  \item \textbf{Archivo}: solidum/elements/solid\_3d/hex20.py.
  \item \textbf{Clase}: \texttt{Hex20} (registrada vía \texttt{@ElementRegistry.register}); subclase de \texttt{\_HigherOrderSolid3D} desde \textbf{sub-fase 2} de A.ter --- la entrada del \texttt{Hex27} disparó la centralización (regla de los dos casos reales). El cuerpo de \texttt{hex20.py} quedó como pura declaración de atributos: funciones de forma, cuadraturas, FACE\_NODES y funciones de cara Quad8 (\texttt{\_N\_quad8}/\texttt{\_dN\_quad8} del paquete 2D).
  \item \textbf{Funciones núcleo}: \texttt{\_N\_hex20(xi, eta, zeta)}, \texttt{\_dN\_hex20(xi, eta, zeta)} en solidum/elements/solid\_3d/\_shared.py. Numpy puro (no \texttt{@njit}) por la complejidad de la rama por nodo serendipity, paritario con \texttt{\_N\_quad8}/\texttt{\_dN\_quad8} del 2D. El kinematics genérico `\texttt{\_kinematics\_higher\_order\_3d}\texttt{ (también en }\_shared.py\texttt{) es reutilizado por la base con cualquier par }\texttt{(shape\_fn, grad\_fn)}` 3D.
  \item \textbf{Tests}:
  \item \texttt{tests/test\_solid\_3d\_higher\_order.py} --- clase \texttt{TestHex20Element} (dimensiones/DOFs, patch lineal, patch cuadrático, simetría K, jacobiano degenerado abortado, body load suma, face traction suma + reparto vértice/medio, masa consistente y lumped HRZ, gauss state).
  \item \texttt{tests/test\_rigid\_body\_modes.py} --- \texttt{TestRBMHex20} (3 trans + 3 rot + rank-deficiency = 6).
  \item \texttt{tests/validation/test\_cube\_lame\_3d.py} (extender) --- \texttt{test\_uniaxial\_traction\_hex20\_exact}.
  \item \texttt{tests/validation/test\_macneal\_beam\_3d.py} (extender) --- convergencia Hex20 1\ensuremath{\times}1 sección vs Hex8 1\ensuremath{\times}1.
  \item \texttt{tests/test\_volumetric\_locking\_3d.py} (extender) --- \texttt{TestHex20VolumetricLocking} documentando mitigación parcial.
\end{itemize}


\section{Diálogo}

\begin{itemize}
  \item \textbf{2026-05-27} \ensuremath{\cdot} Spec inicial. Sub-fase 1 de la sub-etapa A.ter (sólidos 3D cuadráticos). Hex20 standalone --- sin base interna compartida en esta sub-fase; el segundo caso (Hex27) entra en sub-fase 2 y dispara la centralización en \texttt{\_HigherOrderSolid3D} (regla de los dos casos reales aplicada).
  \item \textbf{2026-05-27} \ensuremath{\cdot} Implementado y validado. Patch test lineal y cuadrático exactos a precisión máquina, 6 modos rígidos exactos, cubo Lamé 3D uniaxial e hidrostático exactos. Demostración cuantitativa de la mitigación del shear locking sobre MacNeal 3D: Hex20 6\ensuremath{\times}1\ensuremath{\times}1 alcanza 97\% de la flecha analítica EB, frente a < 55\% de Hex8 12\ensuremath{\times}1\ensuremath{\times}1. Mitigación parcial del locking volumétrico documentada: ratio u(\ensuremath{\nu}=0.4999)/u(\ensuremath{\nu}=0.3) \ensuremath{\approx} 0.80 (vs < 0.6 del Hex8). Conteo de modos espurios con cuadratura reducida \texttt{hex\_2x2x2}: 12 = 6 rígidos + 6 hourglass por elemento aislado (corregido desde el "4" inicial tras verificación numérica directa; coincide con Belytschko-Liu-Moran §8.4.4). Spec actualizada a \texttt{status: validated}.
  \item \textbf{2026-05-27} \ensuremath{\cdot} Refactor a la base \texttt{\_HigherOrderSolid3D} con la entrada del \texttt{Hex27} (sub-fase 2). El cuerpo del archivo \texttt{hex20.py} queda como declaración de atributos (\texttt{\_SHAPE\_FN}, \texttt{\_GRAD\_FN}, \texttt{\_DEFAULT\_QUADRATURE}, \texttt{\_MASS\_QUADRATURE}, \texttt{\_FACE\_N\_FN}, \texttt{\_FACE\_DN\_FN}, \texttt{\_FACE\_QUADRATURE}, \texttt{FACE\_NODES}). La maquinaria de cálculo (bucle de Gauss para K, gauss state, body load, face traction, masa consistente, HRZ) ahora vive en la base --- paritaria con la centralización 2D \texttt{\_HigherOrderSolid2D} (Quad8/Quad9/Tri6). Los 26 tests previos pasan sin modificaciones.
\end{itemize}


\newpage

\section{Hex27}



\textit{Elemento sólido 3D de orden cuadrático con interpolación \textbf{Lagrangiana completa} --- producto tensorial de Lagrange cuadrático en las tres direcciones. Reproduce \textbf{todos} los polinomios triquadráticos `\texttt{\ensuremath{\xi}\textasciicircum{}a \ensuremath{\eta}\textasciicircum{}b \ensuremath{\zeta}\textasciicircum{}c}\texttt{ con }\texttt{a, b, c \ensuremath{\in} \{0, 1, 2\}}\texttt{, incluyendo los puramente triquadráticos (}\texttt{\ensuremath{\xi}\ensuremath{^{2}}\ensuremath{\eta}\ensuremath{^{2}}\ensuremath{\zeta}\ensuremath{^{2}}}\texttt{ y combinaciones) que faltan en el serendípito }\texttt{Hex20}\texttt{. Sub-fase 2 de A.ter; dispara la centralización en }\texttt{\_HigherOrderSolid3D}` (regla de los dos casos reales).}



\section{Especificación física}

\subsection{0. Descripción general}

Hexaedro tridimensional de segundo orden con \textbf{27 nodos}: 8 vértices + 12 medios de arista + 6 centros de cara + 1 centro del cuerpo. Aproximación Lagrangiana $C^0$ completa en coordenadas naturales $(\xi, \eta, \zeta) \in [-1, 1]^3$. Formulación isoparamétrica. Régimen de pequeños desplazamientos y deformaciones.

\subsection{1-5. Cinemática, ecuación constitutiva, equilibrio}

Idénticas al `\texttt{Hex20}` --- la diferencia entre serendipity y Lagrange está exclusivamente en el espacio de aproximación discreto, no en la formulación variacional.


\section{Formulación numérica (FEM)}

\subsection{6. Discretización --- convención VTK\_TRIQUADRATIC\_HEXAHEDRON}

Los 27 nodos en el cubo de referencia $[-1, 1]^3$:

\begin{center}
\begin{tabularx}{\textwidth}{YYY}
\hline
\textbf{Rango} & \textbf{Tipo} & \textbf{Posiciones $(\xi_i, \eta_i, \zeta_i)$} \\
\hline
0-7 & vértices & idénticos al \texttt{Hex8}/\texttt{Hex20} \\
8-19 & medios de arista & idénticos al \texttt{Hex20} \\
20 & centro de cara & $(-1, 0, 0)$ --- cara 5 (−\ensuremath{\xi}) \\
21 & centro de cara & $(+1, 0, 0)$ --- cara 3 (+\ensuremath{\xi}) \\
22 & centro de cara & $(0, -1, 0)$ --- cara 2 (−\ensuremath{\eta}) \\
23 & centro de cara & $(0, +1, 0)$ --- cara 4 (+\ensuremath{\eta}) \\
24 & centro de cara & $(0, 0, -1)$ --- cara 0 (−\ensuremath{\zeta}) \\
25 & centro de cara & $(0, 0, +1)$ --- cara 1 (+\ensuremath{\zeta}) \\
26 & centro del cuerpo & $(0, 0, 0)$ \\
\hline
\end{tabularx}
\end{center}
\subsection{7. Funciones de forma --- producto tensorial de Lagrange cuadrático}

Para cada nodo se asigna una terna de índices $(i, j, k) \in \{-1, 0, +1\}^3$ correspondiente a su posición natural. Las funciones de forma son productos tensoriales:

\[
N_n(\xi, \eta, \zeta) = L_{i_n}(\xi)\,L_{j_n}(\eta)\,L_{k_n}(\zeta),
\]

con los tres polinomios de Lagrange cuadrático 1D:

\[
L_{-1}(x) = \tfrac{1}{2}\,x(x - 1),\qquad L_0(x) = 1 - x^2,\qquad L_{+1}(x) = \tfrac{1}{2}\,x(x + 1).
\]

Cumplen $L_{-1}(-1) = L_0(0) = L_{+1}(+1) = 1$ y cero en los demás. Garantiza propiedad delta de Kronecker $N_m(\xi_n, \eta_n, \zeta_n) = \delta_{mn}$ y partición de la unidad $\sum_n N_n = 1$.

\textbf{Espacio reproducido}: el `\texttt{Hex27}\texttt{ reproduce **completamente** el espacio triquadrático @@MINL20@@ --- 27 monomios, idéntico al número de nodos. La diferencia con el }\texttt{Hex20}\texttt{ serendípito está en los términos puramente triquadráticos }\texttt{\ensuremath{\xi}\ensuremath{^{2}}\ensuremath{\eta}\ensuremath{^{2}}}\texttt{, }\texttt{\ensuremath{\eta}\ensuremath{^{2}}\ensuremath{\zeta}\ensuremath{^{2}}}\texttt{, }\texttt{\ensuremath{\xi}\ensuremath{^{2}}\ensuremath{\zeta}\ensuremath{^{2}}}\texttt{, }\texttt{\ensuremath{\xi}\ensuremath{^{2}}\ensuremath{\eta}\ensuremath{^{2}}\ensuremath{\zeta}}\texttt{, }\texttt{\ensuremath{\xi}\ensuremath{^{2}}\ensuremath{\eta}\ensuremath{\zeta}\ensuremath{^{2}}}\texttt{, }\texttt{\ensuremath{\xi}\ensuremath{\eta}\ensuremath{^{2}}\ensuremath{\zeta}\ensuremath{^{2}}}\texttt{, }\texttt{\ensuremath{\xi}\ensuremath{^{2}}\ensuremath{\eta}\ensuremath{^{2}}\ensuremath{\zeta}\ensuremath{^{2}}}\texttt{ que }\texttt{Hex20}` no captura.

\subsection{8. Matriz $\mathbf B$}

Tamaño \textbf{6\ensuremath{\times}81}, ensamblada con la misma plantilla 6\ensuremath{\times}3 del `\texttt{Hex8}\texttt{/}\texttt{Hex20}` por nodo. Derivadas obtenidas vía regla del producto sobre los polinomios de Lagrange:

\[
\frac{\partial N_n}{\partial \xi} = L'_{i_n}(\xi)\,L_{j_n}(\eta)\,L_{k_n}(\zeta),
\]

con

\[
L'_{-1}(x) = x - \tfrac{1}{2},\qquad L'_0(x) = -2x,\qquad L'_{+1}(x) = x + \tfrac{1}{2}.
\]

\subsection{9-10. Rigidez y fuerzas internas}

Idénticas a `\texttt{Hex8}\texttt{/}\texttt{Hex20}` en estructura. El tamaño es 81\ensuremath{\times}81.

\subsection{11. Cuadratura}

Por defecto \textbf{Gauss--Legendre 3\ensuremath{\times}3\ensuremath{\times}3} (27 puntos): orden de exactitud 5 en cada dirección. Suficiente para integrar exactamente $\mathbf B^\top \mathbf D\, \mathbf B$ del `\texttt{Hex27}` sobre Jacobiano constante (el integrando es a lo sumo de grado 3 en cada dirección --- derivada cuadrático\ensuremath{\times}cuadrático) y la masa consistente $\rho \mathbf N^\top \mathbf N$ (grado 4 en cada dirección).

\textbf{Aviso sobre integración reducida}: el esquema $2 \times 2 \times 2$ (\texttt{hex\_2x2x2}, 8 puntos) está disponible pero \textbf{introduce muchos modos espurios} sobre un elemento aislado: verificación numérica directa del proyecto cuenta \textbf{27 modos de hourglass} (= 81 DOFs − 8 Gauss \ensuremath{\times} 6 strain − 6 rigid; jump espectral limpio entre eigs[32] \ensuremath{\approx} 1e-10 y eigs[33] \ensuremath{\approx} 7e3). Es muchísimo peor que en el \texttt{Hex20} reducido (6 hourglass) --- el espacio extra del Lagrangiano completo es subintegrado severamente por 2\ensuremath{\times}2\ensuremath{\times}2. La cita clásica de "1 modo bubble" (Cook-Malkus-Plesha-Witt §11.6) corresponde al \textbf{único modo que persiste tras ensamblar} en mallas estructuradas; en un elemento aislado el conteo real es el de arriba. Sin estabilización implementada. Recomendación operativa estricta: usar el default $3 \times 3 \times 3$ con \texttt{Hex27}.

\subsection{12. Cargas distribuidas consistentes}

\textbf{Fuerza de cuerpo}: integración con cuadratura del elemento, reparto no uniforme entre vértices, medios, centros de cara y centro del cuerpo. Suma global $\mathbf b\cdot V_e$ exacta para $\mathbf b$ uniforme.

\textbf{Tracción de superficie}: cada cara del `\texttt{Hex27}` tiene \textbf{9 nodos} (4 vértices + 4 medios de arista + 1 centro de cara) que se aproximan por funciones Lagrangianas Quad9 sobre el parámetro 2D $(s, t) \in [-1, 1]^2$.

\begin{center}
\begin{tabularx}{\textwidth}{YYYYY}
\hline
\textbf{Cara} & \textbf{Vértices} & \textbf{Medios de arista} & \textbf{Centro de cara} & \textbf{Normal saliente} \\
\hline
0 & (0,3,2,1) & (11, 10, 9, 8) & 24 & −\ensuremath{\zeta} \\
1 & (4,5,6,7) & (12, 13, 14, 15) & 25 & +\ensuremath{\zeta} \\
2 & (0,1,5,4) & (8, 17, 12, 16) & 22 & −\ensuremath{\eta} \\
3 & (1,2,6,5) & (9, 18, 13, 17) & 21 & +\ensuremath{\xi} \\
4 & (2,3,7,6) & (10, 19, 14, 18) & 23 & +\ensuremath{\eta} \\
5 & (3,0,4,7) & (11, 16, 15, 19) & 20 & −\ensuremath{\xi} \\
\hline
\end{tabularx}
\end{center}
El orden dentro de cada cara reproduce la convención de `\texttt{\_N\_quad9}` del paquete 2D: primero los 4 vértices, después los 4 medios en orden cíclico, y finalmente el nodo central. Cuadratura 2D sobre la cara: Gauss $3 \times 3$ (exacta para tracción constante sobre cara plana).

\subsection{13. Salida por punto de Gauss}

\texttt{compute\_gauss\_state(U)} devuelve dict con $\boldsymbol\varepsilon$ y $\boldsymbol\sigma$ en cada uno de los 27 puntos de Gauss del elemento, junto con coordenadas naturales y globales. Sin \texttt{internal\_forces} (ADR 0012).


\section{Limitaciones declaradas}

\subsection{Coste computacional vs `\texttt{Hex20}`}

\texttt{Hex27} tiene \textbf{81 DOFs por elemento} vs 60 del \texttt{Hex20} (35\% más). El espacio extra que captura (términos puramente triquadráticos) \textbf{rara vez gobierna la solución} en problemas físicos estándar --- el \texttt{Hex20} serendípito es la elección por defecto en la mayoría de aplicaciones (Bathe FEP §5.3.2). El \texttt{Hex27} se prefiere cuando:

\begin{itemize}
  \item La geometría tiene curvatura severa y se mapea con elementos isoparamétricos cuadráticos: el nodo central de cuerpo ayuda a representar la curvatura interior.
  \item La sensibilidad al locking es crítica y se prefiere intentar el Q\_2 completo antes que recurrir a B-bar/F-bar.
  \item Se busca convergencia óptima $O(h^{p+1})$ con $p = 2$ en problemas donde el contenido espectral cubre el espacio completo.
\end{itemize}

Para problemas con flexión simple y geometría rectilínea, \texttt{Hex20} y \texttt{Hex27} dan resultados prácticamente idénticos a un coste menor para el primero.

\subsection{Locking volumétrico con $\nu \to 0.5$}

Comparable al `\texttt{Hex20}`: el espacio extra triquadrático añade flexibilidad sin eliminar la incompresibilidad --- formulación mixta u-p o B-bar/F-bar son necesarias para eliminarlo. Declarado como limitación, blindado por test cuantitativo.

\subsection{Hourglass con integración reducida}

\texttt{Hex27} integrado con \texttt{hex\_2x2x2} tiene \textbf{27 modos de hourglass espurios por elemento aislado} (vs 6 del \texttt{Hex20} reducido y 12 del \texttt{Hex8} reducido). El espacio Lagrangiano completo es subintegrado severamente por 8 puntos. En mallas estructuradas la mayoría no propagan tras ensamblar, pero el conteo a nivel de elemento es alto. Sin estabilización.


\section{Contrato de implementación}

\begin{lstlisting}[language=yaml]
name: Hex27
kind: element
status: validated

interface:
  dof_names: [ux, uy, uz]
  n_nodes: 27
  strain_dim: 6
  n_integration_points: 27      # default Gauss 3×3×3

parameters:
  - { name: quadrature, type: str, required: false, default: "hex_3x3x3",
      desc: "Regla de cuadratura desde QuadratureRegistry; hex_2x2x2
             reducida con 1 modo de hourglass espurio (bubble) por
             elemento aislado, sin estabilización" }

material_contract:
  signature: "compute_state(ε, state) -> (σ, C_tangent, state')"
  strain_kind: "Voigt 3D [ε_xx, ε_yy, ε_zz, γ_xy, γ_yz, γ_xz] con γ_ij = 2·ε_ij"

conventions:
  sign: "σ > 0 ⇔ tracción (Reglas §5)"
  voigt: "[xx, yy, zz, xy, yz, xz] con γ_ij = 2·ε_ij (ADR 0012)"
  node_orientation: "VTK_TRIQUADRATIC_HEXAHEDRON; 0-7 vértices (Hex8), 8-19
                     medios de arista (Hex20), 20-25 centros de cara en orden
                     (-x, +x, -y, +y, -z, +z), 26 centro del cuerpo"
  face_numbering: "ADR 0012 — 6 caras con normal saliente, paritarias con
                   Hex8/Hex20 en los vértices y medios; cada cara añade el
                   nodo de centro correspondiente (índice 20-25)"

validity:
  - "pequeños desplazamientos y deformaciones"
  - "geometría sin distorsión severa (det J > tol en todos los puntos de Gauss)"
  - "compatible con materiales con STRAIN_DIM = 6 (Elastic3D, VonMises3D,
     DruckerPrager3D, IsotropicDamage3D)"

out_of_scope:
  - "grandes deformaciones, formulaciones corotacionales"
  - "estabilización contra hourglass para integración reducida"
  - "tracción variable sobre cara"
  - "internal_forces (ADR 0012)"
  - "Tet10 cuadrático (sub-fase 3)"

acceptance:
  verification:
    - name: patch_test_lineal
      setup: "campo lineal u_i = a_ij·x_j en los 27 nodos del contorno"
      expect: "ε constante = sym(a) en todos los Gauss; σ uniforme"
      tol_rel: 1.0e-12
    - name: patch_test_triquadratico
      setup: "campo u_x = c·x²·y²·z² impuesto en los 27 nodos"
      expect: "ε_xx(x,y,z) = 2c·x·y²·z² exacto en todos los Gauss (el espacio
              Q_2 lagrangiano completo lo contiene; Hex20 serendípito NO)"
      tol_rel: 1.0e-12
    - name: modos_de_cuerpo_rigido
      setup: "evaluar K_e y proyectar sobre 6 modos rígidos 3D"
      expect: "K_e · u_rigid ≈ 0; exactamente 6 autovalores nulos"
      tol_abs: 1.0e-9
    - name: simetria_K_e
      setup: "K_e con Elastic3D"
      expect: "K_e == K_e.T exacto"
      tol_abs: 1.0e-12
  specific:
    - name: cubo_lame_3d
      setup: "1 Hex27 en cubo unitario con tracción uniaxial y BCs mínimas"
      expect: "u_x = p·x/E etc. exacto a precisión máquina (campo lineal)"
      tol_rel: 1.0e-10
    - name: body_load_uniforme
      setup: "Hex27 cubo unitario con b = (0, 0, -g)"
      expect: "Σf_z = -g·V_e = -g"
      tol_rel: 1.0e-12
    - name: face_traction_balance
      setup: "tracción (p, 0, 0) sobre cara +ξ del cubo unitario"
      expect: "Σf_x = p·A_cara = p; reparto consistente Lagrange Q_2 (vértices,
              medios y centro de cara con sus pesos respectivos)"
      tol_rel: 1.0e-12
    - name: reduced_integration_introduces_hourglass
      setup: "Hex27 con cuadratura reducida hex_2x2x2"
      expect: "81 - 8·6 = 33 modos cero totales (6 rígidos + 27 hourglass).
              Verificación con jump espectral limpio entre eigs[32] ≈ 1e-10
              y eigs[33] ≈ O(10^3) (escala del módulo de cortante × volumen)"
    - name: mass_consistent_total
      setup: "Hex27 cubo unitario con ρ=7850, masa consistente"
      expect: "Σ M = 3·ρ·V = 23550"
      tol_rel: 1.0e-9
    - name: mass_lumped_total_HRZ
      setup: "Hex27 cubo unitario con ρ=7850, masa lumped HRZ"
      expect: "diag(M) suma = 3·ρ·V; matriz diagonal; todas las entradas
              positivas (HRZ canónico)"
      tol_abs: 1.0e-12

references:
  - "Bathe K.J. (2014). Finite Element Procedures. §5.3.2 (Lagrangiano 3D)."
  - "Cook R.D., Malkus D.S., Plesha M.E., Witt R.J. (2002). Concepts and
    Applications of FEA. §6.6, §11.6 (Q2/E27 con 2×2×2 reducida)."
  - "Zienkiewicz O.C., Taylor R.L. (2005). The Finite Element Method, vol. 1,
    §8.6 (familia Lagrangiana 3D)."
  - "ADR 0012 — Sólidos 3D: convención Voigt 6D y caras."

\end{lstlisting}


\section{Implementación}

\begin{itemize}
  \item \textbf{Archivo}: solidum/elements/solid\_3d/hex27.py.
  \item \textbf{Clase}: \texttt{Hex27} (registrada vía \texttt{@ElementRegistry.register}); subclase de \texttt{\_HigherOrderSolid3D}.
  \item \textbf{Funciones núcleo}: \texttt{\_N\_hex27(xi, eta, zeta)}, \texttt{\_dN\_hex27(xi, eta, zeta)} en solidum/elements/solid\_3d/\_shared.py. Implementadas en numpy puro vía productos tensoriales de los polinomios Lagrange 1D \texttt{\_L\_lagrange\_quad} y \texttt{\_dL\_lagrange\_quad}.
  \item \textbf{Centralización (sub-fase 2)}: \texttt{\_HigherOrderSolid3D} abstrae el bucle de Gauss para \texttt{K}, \texttt{F\_int}, \texttt{gauss\_state}, \texttt{body\_load}, \texttt{face\_traction} y masa consistente. Comparte con \texttt{Hex20} la misma kinematics y la misma maquinaria; sólo cambian las funciones de forma 3D, la cuadratura, las funciones de forma 2D de cara y los \texttt{FACE\_NODES}.
  \item \textbf{Tests}:
  \item \texttt{tests/test\_solid\_3d\_higher\_order.py} --- \texttt{TestHex27Element} (patch lineal, patch \textbf{triquadrático}, simetría K, jacobiano degenerado, body load, face traction balance, masa consistente y lumped HRZ, reduced integration 1 hourglass).
  \item \texttt{tests/test\_rigid\_body\_modes.py} --- \texttt{TestRBMHex27} (3 trans + 3 rot + rank-deficiency = 6).
  \item \texttt{tests/validation/test\_cube\_lame\_3d.py} --- \texttt{test\_uniaxial\_traction\_hex27\_exact}, \texttt{test\_hydrostatic\_compression\_hex27}.
  \item \texttt{tests/validation/test\_macneal\_beam\_3d.py} --- \texttt{test\_macneal\_3d\_hex27\_*} (convergencia espejo del Hex20).
  \item \texttt{tests/test\_volumetric\_locking\_3d.py} --- \texttt{TestHex27VolumetricLocking} (mitigación comparable al Hex20).
\end{itemize}


\section{Diálogo}

\begin{itemize}
  \item \textbf{2026-05-27} \ensuremath{\cdot} Spec inicial. Sub-fase 2 de A.ter --- el segundo caso real (Hex27) dispara la centralización en \texttt{\_HigherOrderSolid3D} (regla de los dos casos reales aplicada en el momento canónico, igual que el \texttt{\_HigherOrderSolid2D} cuando entró el \texttt{Quad9} después del \texttt{Quad8}).
  \item \textbf{2026-05-27} \ensuremath{\cdot} Implementado y validado. Patch test lineal, bilineal mixto (\texttt{u\_x = c\ensuremath{\cdot}xy}) y triquadrático completo (\texttt{u\_x = c\ensuremath{\cdot}x\ensuremath{^{2}}y\ensuremath{^{2}}z\ensuremath{^{2}}}) exactos a precisión máquina --- esta última es la capacidad distintiva del Hex27 respecto al Hex20 serendípito. K simétrica, 6 modos rígidos, masas consistente y lumped HRZ con todas las entradas positivas, balance de cargas distribuidas correcto. Spec a \texttt{status: validated}.
\end{itemize}


\newpage

\chapter{Elementos Térmicos — 2D}

\section{Quad4Thermal}



\textit{Spec pre-redactada por la IA (2026-08-25). \textbf{El usuario revisa la física; el plumbing está resuelto.} Los puntos abiertos de §Diálogo son físicos o de alcance, nunca arquitecturales.}

\textit{> Segundo componente de la Etapa 8 --- C1 núcleo mínimo. Primer elemento de la familia térmica. Consume \texttt{ThermalConduction}.}



\section{Especificación física}

\subsection{0. Descripción general}

Cuadrilátero isoparamétrico de 4 nodos que resuelve \textbf{conducción de calor en 2D}. Misma geometría, mismas funciones de forma bilineales y misma cuadratura que el \texttt{Quad4} mecánico; lo que cambia es el campo que interpola y la ecuación que discretiza.

\textbf{Un grado de libertad por nodo}: la temperatura $T$. \texttt{DOF\_NAMES = ['T']}.

Es una \textbf{clase distinta} del \texttt{Quad4} mecánico, no un modo de operación suyo. \texttt{DOF\_NAMES} es atributo de clase y el registro, la validación temprana y el cruce spec\ensuremath{\leftrightarrow}código dependen de poder leerlo sin construir el objeto; un elemento que decidiera sus DOFs según el material recibido rompería ese contrato declarativo. Es también la separación que adoptan los códigos de referencia (Abaqus \texttt{CPE4} vs \texttt{DC2D4}; ANSYS \texttt{PLANE182} vs \texttt{PLANE55}).

\subsection{1. Interpolación del campo}

\[
T(\xi,\eta) = \sum_{i=0}^{3} N_i(\xi,\eta)\,T_i, \qquad N_i = \tfrac14(1+\xi_i\xi)(1+\eta_i\eta)
\]

Las $N_i$ son \textbf{las mismas} del \texttt{Quad4} mecánico. Isoparamétrico: geometría y temperatura comparten interpolación.

\subsection{2. Gradiente de temperatura}

\[
\nabla T = \mathbf B\,\mathbf T_e, \qquad \mathbf B = \begin{bmatrix} \partial N_0/\partial x & \cdots & \partial N_3/\partial x \\[2pt] \partial N_0/\partial y & \cdots & \partial N_3/\partial y \end{bmatrix} \in \mathbb R^{2\times 4}
\]

\textbf{Diferencia esencial con el mecánico}: aquí $\mathbf B$ es literalmente el gradiente de las funciones de forma, sin reordenamiento. En el \texttt{Quad4} mecánico las mismas derivadas $\partial N_i/\partial x_j$ se reordenan en una matriz $3\times 8$ que produce la deformación en notación Voigt. La materia prima es idéntica ---el kernel \texttt{\_compute\_kinematics} ya la calcula---; el ensamblaje difiere.

Consecuencia: el gradiente $\nabla T$ es \textbf{bilineal por elemento} y su valor en cada punto de Gauss no es constante, a diferencia del CST.

\subsection{3. Ecuación discretizada --- forma débil}

Partiendo del balance $\rho c\,\dot T = \nabla\cdot(\mathbf k\nabla T) + Q$, multiplicando por una función de peso $w$ e integrando por partes:

\[
\int_\Omega \rho c\,w\,\dot T\,d\Omega + \int_\Omega (\nabla w)^\top \mathbf k \nabla T\,d\Omega = \int_\Omega w\,Q\,d\Omega - \int_{\partial\Omega} w\,\bar q\,d\Gamma
\]

que en forma matricial es el sistema semidiscreto

\[
\mathbf C\,\dot{\mathbf T} + \mathbf K\,\mathbf T = \mathbf F
\]

\textbf{de primer orden en el tiempo} --- a diferencia del mecánico $\mathbf M\ddot{\mathbf u} + \mathbf C\dot{\mathbf u} + \mathbf K\mathbf u = \mathbf F$, que es de segundo. Por eso Newmark no aplica y el transitorio requiere solver propio (spec separada).

En \textbf{régimen estacionario} el término $\mathbf C\dot{\mathbf T}$ se anula y queda $\mathbf K\mathbf T = \mathbf F$, que resuelve el \texttt{LinearSolver} existente sin modificación.

\subsection{4. Matriz de conductividad}

\[
\mathbf K_e = \int_\Omega \mathbf B^\top\,\mathbf k\,\mathbf B\;t\;d\Omega = \sum_{g} \mathbf B_g^\top\,\mathbf k\,\mathbf B_g\;t\;|\mathbf J_g|\,w_g
\]

$4\times4$, \textbf{simétrica} (por serlo $\mathbf k$) y \textbf{semidefinida positiva}: tiene un modo nulo, el campo de temperatura uniforme $\mathbf T = c\,[1,1,1,1]^\top$, que no produce flujo. Es el análogo térmico exacto de los modos de sólido rígido en mecánica, y como allí, el sistema global sólo se vuelve resoluble al imponer al menos una condición de Dirichlet.

El \textbf{espesor} $t$ interviene igual que en el mecánico: el problema plano representa una rebanada de profundidad $t$, y el flujo total escala con ella.

\subsection{5. Matriz de capacidad}

\[
\mathbf C_e = \int_\Omega \rho c\;\mathbf N^\top\mathbf N\;t\;d\Omega
\]

$4\times4$, simétrica y \textbf{definida positiva}. Estructuralmente idéntica a la matriz de masa consistente con $\rho c$ en lugar de $\rho$; de hecho $\mathbf C_e = \rho c\,\mathbf M_e^{escalar}/\rho$ sobre el mismo integrando.

\textbf{Por decisión de alcance, el default es \texttt{lumped}, al revés que en dinámica estructural.} La razón es física, no de conveniencia:

\begin{itemize}
  \item La capacidad \textbf{consistente} produce \textbf{oscilaciones espurias} en los primeros pasos cuando existe un frente térmico abrupto (un salto de temperatura impuesto en la frontera). El campo calculado oscila alrededor de la solución y puede dar temperaturas \textbf{fuera del rango de los datos} ---por debajo del mínimo o por encima del máximo impuestos---, lo que viola el \textbf{principio del máximo} de la ecuación de difusión: una propiedad matemática de la solución exacta, no un artefacto tolerable.
  \item La capacidad \textbf{lumped} amortigua esas oscilaciones y preserva la monotonía, a cambio de algo de precisión en régimen suave.
\end{itemize}

Ambas quedan disponibles con el mismo parámetro \texttt{lumping} del resto del catálogo. La coherencia arquitectural se preserva en el \textbf{mecanismo}; el default lo fija el criterio físico del dominio (Lewis-Nithiarasu-Seetharamu §6; Zienkiewicz-Taylor Vol.1 §7).

\subsection{6. Condiciones de frontera y cargas}

\textbf{Dirichlet --- temperatura impuesta.} $T = \bar T$ en $\Gamma_T$. Se impone por eliminación directa, con la maquinaria existente del ADR 0004, sin cambios: el mecanismo es agnóstico al significado del DOF.

\textbf{Neumann --- flujo impuesto.} $\bar q$ [W/m\ensuremath{^{2}}] normal a un borde. El vector consistente sobre el borde $\Gamma_q$:

\[
\mathbf f^{\bar q}_e = -\int_{\Gamma_q} \bar q\,\mathbf N^\top\,t\;d\Gamma
\]

Para $\bar q$ uniforme sobre un borde recto de longitud $L$, la integral da $\bar q\,L\,t/2$ a cada uno de los dos nodos del borde, y cero a los otros dos --- el mismo reparto mitad-mitad que \texttt{compute\_edge\_traction} produce en el mecánico, con la misma numeración de bordes \texttt{EDGE\_NODES = ((0,1),(1,2),(2,3),(3,0))}.

\textbf{Convención de signo}: $\bar q > 0$ significa flujo \textbf{saliente} del dominio (enfriamiento), coherente con que la normal $\mathbf n$ es saliente y con la deducción de la forma débil, donde el término de frontera aparece con signo negativo. Un flujo entrante se declara con $\bar q < 0$.

Un borde sin condición declarada queda con $\bar q = 0$ --- \textbf{frontera adiabática}, aislada. Es el default natural de Neumann homogéneo, exactamente como un borde mecánico sin tracción queda libre.

\textbf{Fuente volumétrica.} $Q$ [W/m\ensuremath{^{3}}] se aplica como carga del elemento (decisión de alcance; ver §Diálogo de \texttt{ThermalConduction}):

\[
\mathbf f^{Q}_e = \int_\Omega Q\;\mathbf N^\top\,t\;d\Omega
\]

Para $Q$ uniforme la suma de las componentes nodales es exactamente $Q\,A_e\,t$, invariante ante distorsión de la geometría --- el mismo criterio de verificación que se usa en \texttt{compute\_body\_load} del mecánico.

\subsection{7. Salida por punto de Gauss}

\texttt{compute\_gauss\_state(T\_global)} devuelve, en paralelo al contrato de los sólidos (ADR 0012):

\begin{itemize}
  \item \texttt{points\_natural}, \texttt{points\_global} --- posición de los puntos de Gauss.
  \item \texttt{grad\_T} --- gradiente $(n_g, 2)$.
  \item \texttt{flux} --- vector de flujo $\mathbf q = -\mathbf k\nabla T$, $(n_g, 2)$.
\end{itemize}

No expone \texttt{internal\_forces}: por el cierre por dominio del ADR 0012, ese método corresponde a elementos estructurales 1D. La familia térmica sigue el mismo criterio que los sólidos.


\section{Formulación numérica}

\subsection{8. Cuadratura}

Gauss $2\times2$ (4 puntos) por defecto, la misma del \texttt{Quad4} mecánico, mediante \texttt{QuadratureRegistry}. Es \textbf{exacta} para $\mathbf K_e$ y $\mathbf C_e$ con jacobiano constante (elemento paralelogramo).

La reducción $1\times1$ queda disponible por el mismo parámetro, pero \textbf{no se recomienda}: al igual que en el mecánico introduce modos espurios. En el caso térmico el rango de $\mathbf K_e$ con un punto es 2, frente a los 3 necesarios (4 DOFs menos el modo de temperatura uniforme), de modo que aparece \textbf{un modo hourglass} de temperatura no detectado por la matriz. Se documenta y se blinda con test de conteo de modos nulos.

\subsection{9. Reutilización de kernels}

\texttt{\_shape\_functions\_quad4}, \texttt{\_det\_jacobian\_quad4} y las derivadas de \texttt{\_compute\_kinematics} ya existen compiladas con Numba en \texttt{solidum/elements/solid\_2d/\_shared.py} y se reutilizan sin duplicar. Lo específico térmico es el ensamblaje $\mathbf B^\top\mathbf k\,\mathbf B$, que es más simple que el mecánico.

\subsection{10. Caveats numéricos}

\begin{itemize}
  \item \textbf{Modo nulo de temperatura uniforme}: $\mathbf K$ global es singular sin Dirichlet. El sistema debe tener al menos un nodo con temperatura impuesta, o el solver falla con matriz singular. Es el análogo exacto de un modelo mecánico sin apoyos, y el diagnóstico debe nombrarlo como tal.
  \item \textbf{Sin locking}: el problema escalar no tiene el análogo de la incompresibilidad ni del bloqueo por cortante. No aplica ninguna de las limitaciones que arrastra el \texttt{Quad4} mecánico.
  \item \textbf{Número de Péclet}: sin término convectivo en la ecuación (fuera de alcance C1), no hay riesgo de oscilaciones por advección dominante. Se anota porque será la primera consideración a revisar si en el futuro entra convección forzada del medio.
  \item \textbf{Distorsión de la malla}: como en el mecánico, un jacobiano casi degenerado degrada la solución; se valida $|\mathbf J| > 0$ en todos los puntos de Gauss.
\end{itemize}


\section{Contrato de implementación}

\begin{lstlisting}[language=yaml]
name: Quad4Thermal
kind: element
status: draft            # draft → implemented → validated

interface:
  field: temperature            # campo primario (Etapa 8) — no displacement
  dof_names: [T]
  n_nodes: 4
  flux_dim: 2                   # dimensión de ∇T y de q en 2D (no hay strain_dim)
  n_integration_points: 4       # default Gauss 2×2

parameters:
  - { name: thickness,  type: float, required: false, default: 1.0,
      desc: "Espesor de la rebanada plana [m]; el flujo total escala con él" }
  - { name: quadrature, type: str,   required: false, default: "2x2",
      desc: "Regla desde QuadratureRegistry. 1x1 disponible pero introduce
             un modo hourglass de temperatura — sin estabilización" }

material_contract:
  signature: "compute_flux(∇T) -> (q, k)"
  gradient_kind: "gradiente físico [∂T/∂x, ∂T/∂y] — sin notación Voigt"
  accepted: [ThermalConduction]

conventions:
  sign_flux: "q̄ > 0 ⇔ flujo SALIENTE del dominio (enfriamiento); normal exterior"
  adiabatic_default: "borde sin condición declarada ⇒ q̄ = 0 (aislado)"
  node_orientation: "antihorario (asegura det J > 0) — idéntico al Quad4 mecánico"
  edge_numbering: "EDGE_NODES = ((0,1),(1,2),(2,3),(3,0)) — paritario con Quad4"
  capacity_default: "lumping='lumped' por default, al contrario que en dinámica
                     estructural; razón física en §5"

validity:
  - "conducción lineal; k independiente de T"
  - "material con field='temperature' (familia ThermalMaterial)"
  - "det J > 0 en todos los puntos de Gauss"
  - "el modelo global debe imponer al menos un Dirichlet de temperatura"

out_of_scope:
  - "convección (Robin) y radiación en la frontera — Etapa 8 posterior"
  - "acoplamiento termomecánico y deformación térmica"
  - "conducción con advección (término convectivo del medio)"
  - "cambio de fase / calor latente"
  - "estabilización del modo hourglass con cuadratura reducida"

acceptance:
  verification:
    - name: patch_test_termico_lineal
      setup: "5 Quad4Thermal distorsionados (4 exteriores + 1 interior);
             imponer T = a + b·x + c·y en los nodos del contorno"
      expect: "los nodos interiores adoptan el campo lineal exacto y ∇T =
              (b, c) constante en TODOS los puntos de Gauss; q = -k·(b,c)"
      tol_rel: 1.0e-10
    - name: modo_nulo_temperatura_uniforme
      setup: "K_e de un elemento aislado; campo T = [1,1,1,1]"
      expect: "K_e·T = 0 exacto; rango(K_e) = 3; exactamente un autovalor nulo"
      tol_abs: 1.0e-12
    - name: simetria_y_semidefinicion
      setup: "K_e y C_e con k isótropo y ortótropo, elemento regular y distorsionado"
      expect: "K_e == K_e.T y C_e == C_e.T exactas; autovalores de K_e ≥ 0 con un
               único cero; autovalores de C_e > 0 estrictos"
      tol_abs: 1.0e-12
  specific:
    - name: conduccion_1d_pared_plana
      setup: "franja de Quad4Thermal entre x=0 y x=L; T(0)=T_1, T(L)=T_2;
             bordes superior e inferior adiabáticos; estacionario"
      expect: "perfil lineal exacto T(x) = T_1 + (T_2-T_1)·x/L en todos los nodos;
              flujo uniforme q_x = k·(T_1-T_2)/L en todos los Gauss"
      tol_rel: 1.0e-12
    - name: cilindro_hueco_logaritmico
      setup: "corona circular r ∈ [a, b] mallada con Quad4Thermal;
             T(a)=T_a, T(b)=T_b; estacionario"
      expect: "perfil logarítmico analítico T(r) = T_a + (T_b-T_a)·ln(r/a)/ln(b/a);
              convergencia con refinamiento h. Es el análogo térmico exacto del
              benchmark de Lamé que el proyecto ya usa en mecánica"
      tol_rel: 5.0e-3
    - name: fuente_volumetrica_uniforme
      setup: "Quad4Thermal regular y distorsionado con Q uniforme; sumar
             componentes nodales de f^Q"
      expect: "Σf = Q·A_e·t exacto, invariante ante distorsión de la geometría"
      tol_rel: 1.0e-12
    - name: flujo_neumann_en_un_borde
      setup: "Quad4Thermal con q̄ constante sobre un borde"
      expect: "Σf = -q̄·L·t con reparto L/2 a cada nodo del borde y 0 en los demás"
      tol_rel: 1.0e-12
    - name: balance_energetico_estacionario
      setup: "dominio con fuente Q y Dirichlet en parte de la frontera; estacionario"
      expect: "el calor generado iguala al evacuado por las reacciones nodales:
              Σ Q·A·t + Σ reacciones = 0"
      tol_rel: 1.0e-10
    - name: capacidad_lumped_vs_consistente
      setup: "C_e en ambos modos para elemento regular"
      expect: "ambas suman ρ·c·A_e·t (capacidad total conservada); la lumped es
              diagonal y estrictamente positiva"
      tol_rel: 1.0e-12
    - name: anisotropia_desvia_el_flujo
      setup: "k = diag(k1, k2) con k1 ≠ k2; campo lineal a 45° de los ejes"
      expect: "q no es antiparalelo a ∇T; coincide con -k·∇T calculado a mano"
      tol_rel: 1.0e-12
    - name: jacobiano_degenerado_abortado
      setup: "elemento con nodos colineales o en orden inverso"
      expect: "ValueError con mensaje claro"
    - name: rechazo_material_mecanico
      setup: "construir Quad4Thermal con un material de la familia mecánica (Elastic2D)"
      expect: "TypeError/ValueError identificando la incompatibilidad de familia;
               nunca aceptación silenciosa"

references:
  - "Lewis R.W., Nithiarasu P., Seetharamu K.N. (2004). Fundamentals of the FEM for Heat and Fluid Flow. Wiley. §3 y §6 (capacidad lumped vs consistente, oscilaciones)."
  - "Zienkiewicz O.C., Taylor R.L. (2000). The Finite Element Method, Vol. 1. §7 (problemas de campo escalar)."
  - "Reddy J.N., Gartling D.K. (2010). The Finite Element Method in Heat Transfer and Fluid Dynamics. CRC Press. §4."
  - "Carslaw H.S., Jaeger J.C. (1959). Conduction of Heat in Solids. Oxford. §7.2 (cilindro hueco, solución logarítmica)."
  - "Bathe K.-J. (2014). Finite Element Procedures. §7 (formulación de campo escalar; analogía con elasticidad)."

\end{lstlisting}


\section{Implementación}

\textit{Rellena la IA tras programar.}

\begin{itemize}
  \item Archivo: ---
  \item Clase: ---
  \item Tests:
  \item ---
  \item Notas de traducción: ---
\end{itemize}


\section{Diálogo}

\textit{Puntos abiertos para el usuario. Sólo física y alcance.}

\textbf{2026-08-25 --- IA \ensuremath{\to} usuario. Dos puntos:}

\begin{enumerate}
  \item \textasciitilde{}\textasciitilde{}\textbf{Convención de signo del flujo de frontera.}\textasciitilde{}\textasciitilde{} ✅ \textbf{Resuelto 2026-08-25 --- $\bar q > 0$ es flujo SALIENTE del dominio} (decisión del usuario).
\end{enumerate}

   Es la que resulta de la deducción de la forma débil con normal exterior $\mathbf n$, donde el término de frontera aparece con signo negativo, y coincide con Abaqus y con la mayoría de la literatura de MEF. Un flujo entrante (calentamiento) se declara con $\bar q < 0$.

   \textbf{Consecuencia}: es convención del proyecto y debe registrarse en \texttt{Reglas.md §5} junto a las de signos mecánicas, al cerrar la etapa. Coherente además con \texttt{q = -k\ensuremath{\cdot}\ensuremath{\nabla}T} del material: ambos signos salen de la misma deducción, no son elecciones independientes.

\begin{enumerate}
  \item \textasciitilde{}\textasciitilde{}\textbf{Espesor \texttt{thickness}.}\textasciitilde{}\textasciitilde{} ✅ \textbf{Resuelto 2026-08-25 --- se mantiene, con \texttt{default = 1.0}} (decisión del usuario).
\end{enumerate}

   Misma semántica que en el \texttt{Quad4} mecánico: el problema plano representa una rebanada de profundidad $t$ y el flujo total escala con ella. Tres razones:

\begin{itemize}
  \item \textbf{Consistencia} con los cinco sólidos planos del catálogo, que ya lo declaran.
  \item \textbf{Balances energéticos en unidades absolutas}: el test \texttt{balance\_energetico\_estacionario} da watts reales, no "watts por metro de profundidad".
  \item \textbf{Extensión futura}: si entra el acoplamiento termomecánico, el elemento térmico y el mecánico deben compartir profundidad; tenerlo desde el principio evita una asimetría posterior.
\end{itemize}

   El coste es nulo: quien no lo declare obtiene exactamente el comportamiento "por unidad de profundidad" que asume la literatura térmica. Omitirlo habría sido particularizar al caso de hoy.


\newpage

\chapter{Elementos Térmicos — 3D}

\section{Hex8Thermal}



\textit{Spec pre-redactada por la IA (2026-08-25). \textbf{Spec de extensión} (\texttt{Reglas.md §4}): documenta \textbf{sólo lo que cambia} respecto a \texttt{Quad4Thermal}, que es el hermano 2D y donde vive la formulación completa de la familia térmica. Sin ADR propio: reusa las decisiones de alcance de la Etapa 8.}

\textit{> Tercer componente de la Etapa 8 --- C1 núcleo mínimo. Consume \texttt{ThermalConduction}.}



\section{Qué comparte con \texttt{Quad4Thermal}}

Todo lo esencial de la formulación, que \textbf{no se repite aquí}:

\begin{itemize}
  \item Un DOF por nodo, la temperatura $T$. \texttt{DOF\_NAMES = ['T']}.
  \item Sistema semidiscreto $\mathbf C\dot{\mathbf T} + \mathbf K\mathbf T = \mathbf F$, de primer orden en el tiempo. Estacionario $\mathbf K\mathbf T = \mathbf F$ con el \texttt{LinearSolver} existente.
  \item $\mathbf K_e = \int \mathbf B^\top\mathbf k\,\mathbf B\,d\Omega$ simétrica y \textbf{semidefinida positiva}, con el modo nulo de temperatura uniforme.
  \item $\mathbf C_e = \int \rho c\,\mathbf N^\top\mathbf N\,d\Omega$, con \textbf{\texttt{lumped} por defecto} y la justificación física del principio del máximo.
  \item Convención de signo: $\bar q > 0$ es flujo \textbf{saliente} del dominio; frontera sin declarar es \textbf{adiabática}.
  \item Fuente $Q$ como carga del elemento.
  \item Sin locking de ningún tipo; sin \texttt{internal\_forces} (cierre por dominio del ADR 0012).
  \item Clase separada del \texttt{Hex8} mecánico, por el mismo argumento del \texttt{ClassVar DOF\_NAMES}.
\end{itemize}


\section{Qué cambia respecto a \texttt{Quad4Thermal}}

\subsection{1. Dimensión del gradiente y de la matriz B}

\[
\nabla T = \mathbf B\,\mathbf T_e, \qquad \mathbf B = \begin{bmatrix} \partial N_0/\partial x & \cdots & \partial N_7/\partial x \\[2pt] \partial N_0/\partial y & \cdots & \partial N_7/\partial y \\[2pt] \partial N_0/\partial z & \cdots & \partial N_7/\partial z \end{bmatrix} \in \mathbb R^{3\times 8}
\]

\texttt{flux\_dim = 3}. La conductividad $\mathbf k$ es $3\times3$; el material \texttt{ThermalConduction} la expande a $k\mathbf I_3$ cuando recibe un escalar.

$\mathbf K_e$ y $\mathbf C_e$ son $8\times8$. El modo nulo sigue siendo único (temperatura uniforme), luego $\text{rango}(\mathbf K_e) = 7$.

\subsection{2. Sin espesor}

\textbf{\texttt{Hex8Thermal} no lleva \texttt{thickness}.} El volumen sale de la geometría, igual que en los sólidos 3D mecánicos. Es la diferencia visible más inmediata para el usuario del YAML respecto al hermano 2D.

\subsection{3. Interpolación trilineal}

\[
N_i(\xi,\eta,\zeta) = \tfrac18(1+\xi_i\xi)(1+\eta_i\eta)(1+\zeta_i\zeta)
\]

Las mismas del \texttt{Hex8} mecánico, orden de nodos VTK\_HEXAHEDRON. El gradiente resulta \textbf{trilineal} por elemento.

\subsection{4. Cuadratura}

Gauss $2\times2\times2$ (8 puntos) por defecto --- \texttt{hex\_2x2x2} del \texttt{QuadratureRegistry}, idéntica al \texttt{Hex8} mecánico. Exacta para $\mathbf K_e$ y $\mathbf C_e$ con jacobiano constante.

Reducida \texttt{hex\_1x1x1} disponible por el mismo parámetro pero \textbf{no recomendada}: con un solo punto el rango de $\mathbf K_e$ es 3 frente a los 7 necesarios, dejando \textbf{4 modos hourglass} de temperatura sin detectar (frente al único del \texttt{Quad4Thermal} reducido). Se blinda con test de conteo.

\subsection{5. Flujo de frontera por cara, no por borde}

El flujo Neumann se aplica sobre \textbf{caras} con normal saliente, no sobre bordes:

\[
\mathbf f^{\bar q}_e = -\int_{\Gamma_q} \bar q\,\mathbf N^\top\,d\Gamma
\]

Numeración \textbf{paritaria con el \texttt{Hex8} mecánico} --- misma \texttt{FACE\_NODES}, mismas 6 caras: 0 ($-\zeta$), 1 ($+\zeta$), 2 ($-\eta$), 3 ($+\xi$), 4 ($+\eta$), 5 ($-\xi$). Reutiliza la maquinaria de integración sobre caras que ya existe.

Para $\bar q$ uniforme sobre una cara plana de área $A$, el reparto es $\bar q A/4$ a cada uno de los 4 nodos de la cara y cero a los otros cuatro.

\subsection{6. Validación específica 3D}

Los benchmarks del hermano 2D se extienden, más uno propio de la dimensión:

\begin{itemize}
  \item \textbf{Pared plana 3D}: perfil lineal exacto; verifica que la tercera dimensión no introduce error.
  \item \textbf{Cilindro hueco 3D}: perfil logarítmico, extensión del 2D --- análogo térmico del Lamé 3D que el proyecto ya valida en mecánica.
  \item \textbf{Cross-check contra \texttt{Quad4Thermal}}: la misma pared plana resuelta en 2D y con una capa de \texttt{Hex8Thermal} con las caras $z$ adiabáticas debe dar campos idénticos nodo a nodo. Es el análogo del cross-check 3D\ensuremath{\leftrightarrow}2D \texttt{plane\_strain} que en A.bis blindó los materiales a 10-14 decimales, y \textbf{cubre los errores en la tercera dimensión sin necesidad de mallador}.
\end{itemize}

La \textbf{esfera hueca} ---único caso con flujo genuinamente tridimensional--- queda \textbf{diferida} de esta etapa (decisión del usuario, 2026-08-25): exige mallar un octante de corona esférica, y la cobertura que aportaría ya la da el cross-check anterior. Ver §Diálogo.


\section{Contrato de implementación}

\begin{lstlisting}[language=yaml]
name: Hex8Thermal
kind: element
status: draft            # draft → implemented → validated

interface:
  field: temperature
  dof_names: [T]
  n_nodes: 8
  flux_dim: 3                   # ∇T y q en 3D
  n_integration_points: 8       # default Gauss 2×2×2

parameters:
  - { name: quadrature, type: str, required: false, default: "hex_2x2x2",
      desc: "Regla desde QuadratureRegistry. hex_1x1x1 disponible pero deja 4
             modos hourglass de temperatura — sin estabilización" }
  # Sin `thickness`: el volumen sale de la geometría (igual que los sólidos 3D).

material_contract:
  signature: "compute_flux(∇T) -> (q, k)"
  gradient_kind: "gradiente físico [∂T/∂x, ∂T/∂y, ∂T/∂z] — sin notación Voigt"
  accepted: [ThermalConduction]

conventions:
  sign_flux: "q̄ > 0 ⇔ flujo SALIENTE del dominio; normal exterior"
  adiabatic_default: "cara sin condición declarada ⇒ q̄ = 0 (aislada)"
  node_orientation: "orden VTK_HEXAHEDRON, volumen positivo — idéntico al Hex8"
  face_numbering: "FACE_NODES paritaria con Hex8: 0(−ζ) 1(+ζ) 2(−η) 3(+ξ) 4(+η) 5(−ξ)"
  capacity_default: "lumping='lumped' por default; razón física en Quad4Thermal §5"

validity:
  - "conducción lineal; k independiente de T"
  - "material con field='temperature' (familia ThermalMaterial)"
  - "det J > 0 en todos los puntos de Gauss"
  - "el modelo global debe imponer al menos un Dirichlet de temperatura"

out_of_scope:
  - "convección (Robin) y radiación en la frontera — Etapa 8 posterior"
  - "acoplamiento termomecánico"
  - "conducción con advección; cambio de fase"
  - "estabilización del hourglass con cuadratura reducida"

acceptance:
  verification:
    - name: patch_test_termico_lineal_3d
      setup: "malla de Hex8Thermal distorsionados con nodo(s) interior(es);
             imponer T = a + b·x + c·y + d·z en los nodos del contorno"
      expect: "los nodos interiores adoptan el campo lineal exacto y ∇T =
              (b, c, d) constante en TODOS los puntos de Gauss"
      tol_rel: 1.0e-10
    - name: modo_nulo_temperatura_uniforme
      setup: "K_e de un elemento aislado; campo T = ones(8)"
      expect: "K_e·T = 0 exacto; rango(K_e) = 7; exactamente un autovalor nulo"
      tol_abs: 1.0e-12
    - name: hourglass_con_integracion_reducida
      setup: "K_e con quadrature='hex_1x1x1' sobre elemento aislado"
      expect: "rango(K_e) = 3; 5 autovalores nulos (1 físico + 4 hourglass).
               Documenta la limitación, no la corrige"
      tol_abs: 1.0e-12
    - name: simetria_y_semidefinicion
      setup: "K_e y C_e con k isótropo y ortótropo, elemento regular y distorsionado"
      expect: "K_e == K_e.T y C_e == C_e.T exactas; autovalores de K_e ≥ 0 con un
               único cero; autovalores de C_e > 0 estrictos"
      tol_abs: 1.0e-12
  specific:
    - name: conduccion_1d_pared_plana_3d
      setup: "bloque de Hex8Thermal entre x=0 y x=L; T(0)=T_1, T(L)=T_2;
             las otras cuatro caras adiabáticas; estacionario"
      expect: "perfil lineal exacto en todos los nodos; q uniforme y sin
              componentes transversales (q_y = q_z = 0)"
      tol_rel: 1.0e-12
    - name: cilindro_hueco_logaritmico_3d
      setup: "sector de corona circular extruido; T(a)=T_a, T(b)=T_b;
             caras planas adiabáticas; estacionario"
      expect: "perfil logarítmico T(r) = T_a + (T_b−T_a)·ln(r/a)/ln(b/a);
              convergencia con refinamiento h"
      tol_rel: 5.0e-3
    # Esfera hueca (perfil T(r) = T_a + (T_b−T_a)·(1/a − 1/r)/(1/a − 1/b)):
    # DIFERIDA de esta etapa por decisión del usuario (2026-08-25). Es el único
    # caso con flujo genuinamente tridimensional, pero exige mallar un octante
    # de corona esférica —trabajo geométrico que ya difirió la campaña de
    # validación 3D de A.bis por el mismo motivo— y la cobertura que aportaría
    # sobre la tercera dimensión ya la da `consistencia_con_quad4thermal`, que
    # no necesita mallador. Retomar si entra un mallador 3D nativo (gmsh API)
    # o si aparece un caso de uso con flujo radial esférico.
    - name: fuente_volumetrica_uniforme
      setup: "Hex8Thermal regular y distorsionado con Q uniforme"
      expect: "Σf = Q·V_e exacto, invariante ante distorsión"
      tol_rel: 1.0e-12
    - name: flujo_neumann_en_una_cara
      setup: "Hex8Thermal con q̄ constante sobre cada una de las 6 caras"
      expect: "Σf = -q̄·A_cara con reparto A/4 a cada nodo de la cara y 0 en el
              resto; verificado en las 6 caras"
      tol_rel: 1.0e-12
    - name: balance_energetico_estacionario
      setup: "dominio con fuente Q y Dirichlet en parte de la frontera"
      expect: "Σ Q·V + Σ reacciones = 0"
      tol_rel: 1.0e-10
    - name: capacidad_lumped_vs_consistente
      setup: "C_e en ambos modos para elemento regular"
      expect: "ambas suman ρ·c·V_e; la lumped es diagonal y estrictamente positiva"
      tol_rel: 1.0e-12
    - name: anisotropia_desvia_el_flujo_3d
      setup: "k = diag(k1, k2, k3) distintos; campo lineal oblicuo a los ejes"
      expect: "q coincide con -k·∇T calculado a mano; no antiparalelo a ∇T"
      tol_rel: 1.0e-12
    - name: consistencia_con_quad4thermal
      setup: "problema de pared plana resuelto con Quad4Thermal (2D) y con una
             capa de Hex8Thermal con las caras z adiabáticas"
      expect: "campos de temperatura idénticos nodo a nodo; es el cross-check
              dimensional análogo al que el proyecto usa entre materiales 3D y
              2D plane_strain"
      tol_rel: 1.0e-12
    - name: jacobiano_degenerado_abortado
      setup: "elemento con nodos coplanares o numeración invertida"
      expect: "ValueError con mensaje claro"
    - name: rechazo_material_mecanico
      setup: "construir Hex8Thermal con un material de la familia mecánica"
      expect: "TypeError/ValueError identificando la incompatibilidad de familia"

references:
  - "Lewis R.W., Nithiarasu P., Seetharamu K.N. (2004). Fundamentals of the FEM for Heat and Fluid Flow. Wiley. §3, §6."
  - "Zienkiewicz O.C., Taylor R.L. (2000). The Finite Element Method, Vol. 1. §7."
  - "Carslaw H.S., Jaeger J.C. (1959). Conduction of Heat in Solids. Oxford. §7.2 (cilindro hueco). §9.3 cubre la esfera hueca, diferida de esta etapa."
  - "Reddy J.N., Gartling D.K. (2010). The Finite Element Method in Heat Transfer and Fluid Dynamics. CRC Press. §4."

\end{lstlisting}


\section{Implementación}

\textit{Rellena la IA tras programar.}

\begin{itemize}
  \item Archivo: ---
  \item Clase: ---
  \item Tests:
  \item ---
  \item Notas de traducción: ---
\end{itemize}


\section{Diálogo}

\textit{Puntos abiertos para el usuario. Sólo física y alcance.}

\textbf{2026-08-25 --- IA \ensuremath{\to} usuario. Un punto:}

\begin{enumerate}
  \item \textasciitilde{}\textasciitilde{}\textbf{Benchmark de la esfera hueca.}\textasciitilde{}\textasciitilde{} ✅ \textbf{Resuelto 2026-08-25 --- DIFERIDA de esta etapa} (decisión del usuario).
\end{enumerate}

   La esfera hueca es el único caso propuesto con flujo genuinamente tridimensional: su perfil $\propto 1/r$ no lo producen ni la pared plana ni el cilindro, donde la temperatura varía en una sola dirección. Pero exige \textbf{mallar un octante de corona esférica}, trabajo geométrico que ya difirió la campaña de validación 3D de A.bis por el mismo motivo.

   Se difiere porque \textbf{la cobertura que aportaría está cubierta por otra vía}: \texttt{consistencia\_con\_quad4thermal} detecta errores en la tercera dimensión comparando el mismo problema resuelto en 2D y en 3D, y no necesita mallador. La esfera daría confirmación adicional, no una comprobación ausente.

   Nota técnica: cuando se retome, el \texttt{Hex8Thermal} \textbf{no} sufrirá el problema que bloquea al \texttt{Tet10} sobre superficie curva (deuda técnica \#8) --- al ser trilineal no tiene nodos intermedios que puedan quedar sobre aristas rectas. El bloqueante es sólo de mallado.

   \textbf{Criterio de retoma}: si entra un mallador 3D nativo (gmsh API) o aparece un caso de uso con flujo radial esférico.


\newpage

\chapter{Modelos Constitutivos — 1D}

\section{CableMaterial1D}



\textit{Orden de trabajo. El usuario escribe \textbf{especificación física}, \textbf{formulación} y \textbf{contrato}; la IA rellena \textbf{implementación} y responde en \textbf{diálogo}.}



\section{Especificación física}

\subsection{0. Descripción general}
Material 1D memoryless que modela la respuesta axial de un cable: \textbf{elástico lineal en tensión y nulo en compresión}. No tiene variables internas ni historia --- la respuesta depende exclusivamente de la deformación instantánea. Captura la propiedad esencial de un cable: que solo transmite esfuerzo cuando está tensado.

\subsection{1. Ecuación constitutiva}
\[
\sigma(\varepsilon) \;=\; E\,\langle\varepsilon\rangle^+ \;=\; \begin{cases} E\,\varepsilon & \varepsilon > 0 \\ 0 & \varepsilon \leq 0 \end{cases}
\]

Rampa lineal de pendiente $E$ en tensión ($\varepsilon > 0$), cero en compresión o estado sin deformación ($\varepsilon \leq 0$). $\langle\cdot\rangle^+$ denota la parte positiva (operador de MacAuley).

\subsection{2. Módulo tangente}
\[
E_t(\varepsilon) \;=\; \frac{d\sigma}{d\varepsilon} \;=\; \begin{cases} E & \varepsilon > 0 \\ 0 & \varepsilon \leq 0 \end{cases}
\]

\textbf{Discontinuo en $\varepsilon = 0$}. Por convención el valor en el punto de corte se asigna al régimen compresivo ($E_t = 0$), lo que evita generar rigidez espuria cuando el cable está exactamente sin tensión.

\subsection{3. Variables internas}
No hay. La respuesta es puramente función de $\varepsilon$ actual; el estado del material es el vacío.

\subsection{4. Interpretación física}
\begin{itemize}
  \item $\varepsilon > 0$: cable tensado --- se comporta como barra elástica.
  \item $\varepsilon < 0$: cable destensado --- no transmite nada, la carga pasa por otros caminos del modelo.
  \item $\varepsilon = 0$: frontera neutra --- no hay tensión, no hay rigidez.
\end{itemize}

La no linealidad es \textbf{constitutiva} (en la respuesta del material), no geométrica. Un elemento con este material puede tener cualquier cinemática (lineal o corotacional); la unilateralidad solo vive aquí.


\section{Contrato de implementación}

\begin{lstlisting}[language=yaml]
name: CableMaterial1D
kind: material
status: validated

interface:
  strain_dim: 1
  primary_state_var: null      # memoryless, no exporta variable principal

parameters:
  - { name: E, type: float, required: true, desc: "Módulo de Young (en tensión)" }

signature:
  compute_state: "(ε: float, state_vars=None) -> (σ: float, E_tangent: float, state_vars)"
  strain_kind: axial escalar

conventions:
  sign: "σ > 0 ⇔ ε > 0 (elongación); ε ≤ 0 ⇒ σ = 0"
  state_passthrough: "state_vars se devuelve intacto (el material no almacena historia)"

validity:
  - "E > 0"
  - "|ε| ≲ 1e-2 en el régimen tensado (elasticidad lineal)"
  - "IS_UNILATERAL = True ⇒ solo admitido en elementos con ACCEPTS_UNILATERAL = True (Cable2DCorot, Cable3DCorot). La base Element rechaza la combinación con Truss* en construcción."

out_of_scope:
  - plasticidad, fluencia, fatiga
  - viscoelasticidad / dependencia temporal
  - pretensión inicial (el usuario la modela con una longitud de referencia ajustada)
  - regularización numérica en compresión (no hay rigidez residual en ε ≤ 0)
  - anisotropía direccional; la respuesta es puramente axial escalar

numerical_caveats:
  - "En transiciones tensado ↔ destensado (ε cruza 0), E_tangent salta de E a 0 y Newton-Raphson puede oscilar. Mitigación: usar arc-length o pasos de carga finos si el problema tiene destensiones reales."
  - "Un cable completamente destensado tiene K_M = 0 y K_G = 0 ⇒ matriz tangente singular en los DOFs del elemento. El usuario debe garantizar que otros elementos del modelo proporcionen rigidez suficiente, o pretensar el cable antes de aplicar cargas que puedan destensarlo."

acceptance:
  - name: respuesta en tensión pura
    setup: "ε = 1e-4"
    expect: "σ = E · 1e-4, E_tangent = E"
    tol_rel: 1.0e-12

  - name: respuesta en compresión pura
    setup: "ε = -1e-4"
    expect: "σ = 0, E_tangent = 0"
    tol_abs: 1.0e-15

  - name: punto de corte
    setup: "ε = 0"
    expect: "σ = 0, E_tangent = 0 (régimen compresivo por convención)"
    tol_abs: 1.0e-15

  - name: state_vars se devuelve intacto
    setup: "state_vars = {'dummy': 42}, ε arbitrario"
    expect: "el tercer retorno es idéntico al state_vars de entrada"

references:
  - "Irvine H.M., Cable Structures, MIT Press, §1.2 (modelo elástico del cable)"
  - "MacAuley M.A., on the deflection of beams, 1919 (operador parte positiva)"

\end{lstlisting}


\section{Implementación}

\begin{itemize}
  \item \textbf{Archivo}: solidum/materials/cable\_1d.py --- archivo propio, sin dependencia de otros módulos de material salvo la clase base \texttt{Material}.
  \item \textbf{Clase}: \texttt{CableMaterial1D} --- hereda directamente de \texttt{Material} (no de \texttt{Elastic1D} ni de ningún otro material), registrada vía \texttt{@MaterialRegistry.register}.
  \item \textbf{Tests}: tests/test\_cable\_material.py \ensuremath{\cdot} \texttt{TestCableMaterial1D} --- seis tests:
  \item \texttt{test\_acceptance\_tension\_pura} (criterio 1).
  \item \texttt{test\_acceptance\_compresion\_pura} (criterio 2).
  \item \texttt{test\_acceptance\_punto\_de\_corte} (criterio 3).
  \item \texttt{test\_acceptance\_state\_vars\_passthrough} (criterio 4).
  \item \texttt{test\_registro\_en\_registry} --- verifica autodiscover.
  \item \texttt{test\_validacion\_E\_positivo} --- rechaza $E \leq 0$.
  \item \textbf{Notas de traducción}:
  \item La ramificación en \texttt{compute\_state} es un único \texttt{if strain > 0.0}. El \texttt{>} estricto implementa la convención del punto de corte (el cero se asigna al régimen compresivo). El \texttt{else} cubre tanto $\varepsilon < 0$ como $\varepsilon = 0$ con las mismas salidas.
  \item \texttt{state\_vars} se devuelve como llega (identidad, no copia). El test \texttt{test\_acceptance\_state\_vars\_passthrough} usa \texttt{assertIs} para confirmar identidad.
  \item Validación de $E > 0$ al construir, con mensaje claro. Atrapa typos del input YAML temprano.
  \item El material acepta \texttt{**kwargs} en \texttt{compute\_state} por compatibilidad con elementos que pudieran pasar argumentos adicionales (convención del proyecto).
\end{itemize}


\section{Diálogo}

\begin{itemize}
  \item \textbf{2026-04-21} \ensuremath{\cdot} Material creado como pieza independiente para la futura implementación de elementos de cable (\texttt{Cable2DCorot}, \texttt{Cable3DCorot}). Deliberadamente \textbf{no hereda} de \texttt{Elastic1D} aunque su rama en tensión sea idéntica: mantener el material autocontenido simplifica su lectura y evita acoplamientos futuros si la rama elástica cambia en otro archivo.
  \item \textbf{2026-04-21} \ensuremath{\cdot} Discontinuidad en $\varepsilon = 0$: documentada explícitamente en \texttt{out\_of\_scope: regularización numérica}. Si aparece un caso de uso con destensiones frecuentes donde Newton-Raphson oscile, se añadirá una variante \texttt{CableMaterial1DRegularized} con $E_c \ll E$ en compresión --- \textbf{no} se modificará este material (contrato limpio).
  \item \textbf{2026-05-12} \ensuremath{\cdot} Restricción de emparejamiento elevada a contrato: \texttt{IS\_UNILATERAL = True} en el material; \texttt{ACCEPTS\_UNILATERAL = True} en \texttt{Cable2DCorot}/\texttt{Cable3DCorot}. La base \texttt{Element.\_validate\_material\_compatibility} rechaza la combinación con \texttt{Truss*} (cuyo tangente colapsa la rigidez global sin K\_G que lo compense). Antes el contrato lo permitía y solo se desaconsejaba en docstring; ahora falla limpio al construir.
\end{itemize}


\newpage

\section{Elastic1D}



\textit{Spec \textbf{retroactiva}: el material existe desde el origen del proyecto y está validado por tests preexistentes. Esta spec lo documenta sin cambiar comportamiento --- H-5.3 de la auditoría 2026-05-18.}



\section{Especificación física}

\subsection{0. Descripción general}

Modelo elástico lineal isótropo unidimensional. Captura la relación $\sigma$--$\varepsilon$ axial de barras, cables traccionados (no unilaterales --- para unilateralidad ver \texttt{CableMaterial1D}), y la respuesta elástica de cualquier elemento 1D del catálogo (truss, frame Euler en su componente axial). Independiente de la velocidad, sin variables históricas.

Es el modelo más sencillo del catálogo: \textbf{una sola constante elástica} (módulo de Young $E$) y una sola ecuación constitutiva.

\subsection{1. Descomposición de la deformación}

No aplica --- no hay parte plástica. Toda la deformación es elástica:

\[
\varepsilon = \varepsilon^e
\]

\subsection{2. Ley elástica (Hooke 1D)}

\[
\sigma = E\,\varepsilon
\]

con $E > 0$ módulo de Young. Tangente constante $E_t = E$ en todo régimen.

\subsection{3. Superficie de fluencia / criterio de daño}

No aplica --- el material es indefinidamente elástico.

\subsection{4. Regla de flujo}

No aplica.

\subsection{5. Endurecimiento}

No aplica.

\subsection{6. Condiciones de Kuhn-Tucker}

No aplica.

\subsection{7. Variables internas}

Ninguna --- \texttt{state\_vars = None} se acepta y se devuelve intacto. Sin \texttt{PRIMARY\_STATE\_VAR}.

\subsection{8. Notación}

Esfuerzo y deformación escalares en convención axial: $\sigma > 0 \Leftrightarrow$ tracción, $\varepsilon > 0 \Leftrightarrow$ alargamiento (Reglas §5).


\section{Formulación numérica}

\subsection{9. Esquema temporal}

No aplica --- el material no tiene historia. Cada evaluación de \texttt{compute\_state} es independiente.

\subsection{10. Return mapping / actualización}

No aplica. La evaluación se reduce a una multiplicación escalar:

\[
\sigma_{n+1} = E\,\varepsilon_{n+1}, \qquad E_t = E
\]

\subsection{11. Tangente algorítmica consistente}

\[
C^{\text{alg}} = E
\]

Coincide con la tangente continua y es simétrica trivialmente. \texttt{IS\_SYMMETRIC = True} (por default del contrato base).

\subsection{12. Caveats numéricos}

Ninguno --- el modelo es perfectamente bien condicionado para $E > 0$.


\section{Contrato de implementación}

\begin{lstlisting}[language=yaml]
name: Elastic1D
kind: material
status: validated

interface:
  strain_dim: 1
  primary_state_var: null    # sin variables internas
  is_symmetric: true

parameters:
  - { name: E,       type: float, required: true,  desc: "Módulo de Young (>0)" }
  - { name: density, type: float, required: false, default: null,
      desc: "Densidad (kg/m³). Opcional al construir; obligatoria solo si se ensambla peso propio o masa (ADR 0008)" }

signature:
  compute_state: "(ε: float, state_vars=None, **kwargs) -> (σ: float, E_t: float, state_vars)"
  strain_kind: "axial escalar"

state_variables: []        # ninguna

conventions:
  sign: "σ > 0 ⇔ tracción, ε > 0 ⇔ alargamiento (Reglas §5)"
  units: "[E] = [σ] / [ε] = Pa (o consistente con el sistema)"

validity:
  - "E > 0"
  - "régimen elástico indefinido — sin umbral superior"
  - "pequeñas deformaciones (|ε| ≲ 1e-2 por hipótesis cinemática global)"

out_of_scope:
  - "endurecimiento / fluencia ⇒ usar Elastoplastic1D"
  - "softening por daño ⇒ usar IsotropicDamage1D"
  - "unilateralidad (sin compresión) ⇒ usar CableMaterial1D"
  - "anisotropía o dependencia de temperatura"
  - "viscoelasticidad"

acceptance:
  # Cubierto por tests existentes de elementos 1D (Truss, Frame, Cable, etc.)
  # que usan Elastic1D como material por default. La verificación se hace
  # implícita en cada test de elemento.
  verification:
    - name: ley_de_hooke_exacta
      setup: "evaluar compute_state con ε arbitrario; comparar σ = E·ε"
      expect: "σ = E·ε exacto; tangente = E"
      tol_rel: 1.0e-14
    - name: tangente_constante
      setup: "evaluar dos veces con strain distintos"
      expect: "E_t devuelto idéntico en ambas evaluaciones"
      tol_rel: 1.0e-14
    - name: ciclo_carga_descarga_sin_historia
      setup: "ε creciente luego decreciente; verificar reversibilidad σ↔ε"
      expect: "σ(ε)=E·ε para todo ε (sin histéresis, sin state_vars retornados modificados)"
      tol_rel: 1.0e-14
    - name: rechazo_E_no_positivo
      setup: "construir Elastic1D(E=0) y Elastic1D(E=-1)"
      expect: "ValueError con mensaje claro"
    - name: rechazo_density_negativa
      setup: "construir Elastic1D(E=210e9, density=-1)"
      expect: "ValueError con mensaje claro"

references:
  - "Cualquier texto elemental de mecánica de materiales (Beer, Hibbeler, Timoshenko §1)."

\end{lstlisting}


\section{Implementación}

\begin{itemize}
  \item \textbf{Archivo}: solidum/materials/elastic.py.
  \item \textbf{Clase}: \texttt{Elastic1D}, registrada vía \texttt{@MaterialRegistry.register}.
  \item \textbf{Estado}: \texttt{state\_vars=None} permitido; se acepta y se devuelve sin tocar (compatibilidad con elementos no lineales que pasan diccionarios históricos).
  \item \textbf{\texttt{admissibility\_scale}}: hereda el comportamiento por defecto del contrato \texttt{Material} (no usado --- sin criterio de admisibilidad).
  \item \textbf{Tests}: ningún test unitario propio de \texttt{Elastic1D} (es trivial); su correctitud se ejercita implícitamente en todos los tests de elementos 1D que lo usan como material (decenas de tests en \texttt{tests/test\_truss*.py}, \texttt{tests/test\_frame*.py}, \texttt{tests/test\_cable*.py}, \texttt{tests/test\_dynamic*.py}).
\end{itemize}


\section{Diálogo}

\begin{itemize}
  \item \textbf{2026-05-19} \ensuremath{\cdot} Spec creada retroactivamente para cerrar el hueco H-5.3 de la auditoría 2026-05-18. El material es anterior a la convención de specs validadas; no se modifica código. La aceptación se considera cumplida por cobertura indirecta en tests de elementos.
\end{itemize}


\newpage

\section{Elastoplastic1D}



\textit{Spec \textbf{retroactiva}: el material existe desde la fase inicial del proyecto y está validado por tests unitarios y de integración. Esta spec lo documenta sin cambiar comportamiento --- H-5.3 de la auditoría 2026-05-18.}



\section{Especificación física}

\subsection{0. Descripción general}

Modelo elastoplástico unidimensional independiente de la velocidad ("rate-independent") con criterio de fluencia $|\sigma| \leq \sigma_y + H\,\alpha$ y \textbf{endurecimiento isótropo lineal}. Es el análogo 1D de \texttt{VonMises2D} --- la versión axial sobre la que se valida y benchmarkean los pipelines plásticos. Se usa en barras (\texttt{Truss2D/3D}) y como componente axial conceptual de marcos plásticos (pendiente \texttt{FiberSection}).

\subsection{1. Descomposición aditiva}

\[
\varepsilon = \varepsilon^e + \varepsilon^p
\]

con $\varepsilon^p$ histórica (memoria del material).

\subsection{2. Ley elástica}

\[
\sigma = E\,\varepsilon^e = E\,(\varepsilon - \varepsilon^p)
\]

\subsection{3. Criterio de fluencia}

\[
f(\sigma, \alpha) = |\sigma| - (\sigma_y + H\,\alpha) \;\leq\; 0
\]

donde $\sigma_y > 0$ es la fluencia inicial y $H \geq 0$ el módulo de endurecimiento isótropo lineal.

\subsection{4. Regla de flujo (asociada)}

\[
\dot\varepsilon^p = \dot\gamma\,\operatorname{sign}(\sigma)
\]

\subsection{5. Endurecimiento isótropo lineal}

\[
\dot\alpha = \dot\gamma
\]

$\alpha$ es la deformación plástica acumulada equivalente, escalar monótona no decreciente.

\subsection{6. Condiciones de Kuhn-Tucker}

\[
\dot\gamma \geq 0, \qquad f \leq 0, \qquad \dot\gamma\,f = 0
\]

\subsection{7. Variables internas}

\begin{center}
\begin{tabularx}{\textwidth}{YYY}
\hline
\textbf{Símbolo} & \textbf{Tipo} & \textbf{Significado} \\
\hline
$\varepsilon^p$ & escalar & deformación plástica acumulada con signo \\
$\alpha$ & escalar & deformación plástica acumulada equivalente (\ensuremath{\ge}0) \\
\hline
\end{tabularx}
\end{center}
\texttt{alpha} es la variable principal exportable (\texttt{PRIMARY\_STATE\_VAR = 'alpha'}).

\subsection{8. Notación}

Escalares en convención axial: $\sigma > 0 \Leftrightarrow$ tracción, $\varepsilon > 0 \Leftrightarrow$ alargamiento (Reglas §5). $\varepsilon^p$ con signo (puede ser negativa tras carga compresiva).


\section{Formulación numérica}

\subsection{9. Esquema temporal --- Backward Euler implícito}

Dado el estado convergido $\{\varepsilon^p_n, \alpha_n\}$ al paso $n$ y la deformación total $\varepsilon_{n+1}$, resolver para $\{\sigma_{n+1}, \varepsilon^p_{n+1}, \alpha_{n+1}\}$.

\subsection{10. Return mapping --- forma cerrada}

\textbf{Predictor elástico}:

\[
\sigma^{\text{trial}} = E\,(\varepsilon_{n+1} - \varepsilon^p_n)
\]

\textbf{Función de fluencia trial}:

\[
f^{\text{trial}} = |\sigma^{\text{trial}}| - (\sigma_y + H\,\alpha_n)
\]

\begin{itemize}
  \item Si $f^{\text{trial}} \leq \text{tol}$ \ensuremath{\Rightarrow} \textbf{paso elástico}: $\sigma_{n+1} = \sigma^{\text{trial}}$, $E_t = E$, estado intacto.
  \item Si $f^{\text{trial}} > \text{tol}$ \ensuremath{\Rightarrow} \textbf{paso plástico}, retorno radial cerrado:
\end{itemize}

\[
\Delta\gamma = \frac{f^{\text{trial}}}{E + H}, \qquad \sigma_{n+1} = \sigma^{\text{trial}} - \Delta\gamma\,E\,\operatorname{sign}(\sigma^{\text{trial}})
\]

\[
\varepsilon^p_{n+1} = \varepsilon^p_n + \Delta\gamma\,\operatorname{sign}(\sigma^{\text{trial}}), \qquad \alpha_{n+1} = \alpha_n + \Delta\gamma
\]

Sin Newton local --- la forma cerrada es exacta en 1D.

\subsection{11. Tangente algorítmica consistente}

En régimen elástico: $E_t = E$. En régimen plástico:

\[
E_t = \frac{E\,H}{E + H}
\]

con $H = 0$ recuperando $E_t = 0$ (plasticidad perfecta, descenso de carga indefinido bajo control de desplazamiento).

\texttt{IS\_SYMMETRIC = True} (trivial en 1D).

\subsection{12. Caveats numéricos}

\begin{itemize}
  \item \textbf{Plasticidad perfecta} ($H = 0$): tangente $E_t = 0$ tras la fluencia. Newton-Raphson global puede oscilar o requerir line search en problemas multibarra (modos de mecanismo). Mitigación: usar incremento pequeño cerca del umbral o activar \texttt{line\_search=True} en el solver.
  \item \textbf{Switching exacto trial elástico \ensuremath{\leftrightarrow} plástico}: la tolerancia \texttt{is\_admissible} (ADR 0006) usa \texttt{admissibility\_scale = \ensuremath{\sigma}\_y + H\ensuremath{\cdot}\ensuremath{\alpha}} y patrón atol + rtol\ensuremath{\cdot}escala --- invariancia bajo cambio de unidades.
\end{itemize}


\section{Contrato de implementación}

\begin{lstlisting}[language=yaml]
name: Elastoplastic1D
kind: material
status: validated

interface:
  strain_dim: 1
  primary_state_var: alpha   # deformación plástica acumulada equivalente
  is_symmetric: true

parameters:
  - { name: E,       type: float, required: true,  desc: "Módulo de Young (>0)" }
  - { name: sigma_y, type: float, required: true,  desc: "Esfuerzo de fluencia inicial (>0)" }
  - { name: H,       type: float, required: false, default: 0.0,
      desc: "Módulo de endurecimiento isótropo lineal (≥0). H=0 ⇒ plasticidad perfecta" }
  - { name: density, type: float, required: false, default: null,
      desc: "Densidad (kg/m³). Opcional al construir; obligatoria solo si se ensambla peso propio o masa (ADR 0008)" }

signature:
  compute_state: "(ε: float, state_vars=None) -> (σ: float, E_t: float, state_vars')"
  strain_kind: "axial escalar"

state_schema:
  eps_p: "float — deformación plástica acumulada con signo"
  alpha: "float ≥ 0 — deformación plástica acumulada equivalente"

conventions:
  sign: "σ > 0 ⇔ tracción, ε^p con signo (puede ser negativa)"
  units: "[E]=[σ_y]=[H]=Pa; α adimensional"

validity:
  - "E > 0, σ_y > 0, H ≥ 0"
  - "carga proporcional o no proporcional (sin restricción de trayectoria)"
  - "pequeñas deformaciones (|ε| ≲ 1e-2)"

out_of_scope:
  - "endurecimiento cinemático (back-stress) ⇒ requiere variante 1D dedicada"
  - "endurecimiento isótropo no lineal (Voce, exponencial)"
  - "ablandamiento (H<0) ⇒ requiere regularización"
  - "viscoplasticidad / dependencia de la velocidad"
  - "grandes deformaciones (descomposición multiplicativa F = F^e · F^p)"

acceptance:
  # Cubierto por tests/test_materials_unit.py y tests/test_truss_plastic.py
  unit:
    - name: paso_elastico
      setup: "ε = 0.5·σ_y/E (mitad del umbral); state_vars=None"
      expect: "σ = E·ε exacto, α=0, ε^p=0, tangente=E"
      tol_rel: 1.0e-12
    - name: fluencia_uniaxial_traccion
      setup: "ε = 1.5·σ_y/E (1.5× el umbral elástico)"
      expect: "σ = σ_y (con H=0); α=0.5·σ_y/E; ε^p > 0; tangente=0"
      tol_rel: 1.0e-10
    - name: endurecimiento_lineal
      setup: "carga monotónica hasta ε = 2·σ_y/E con H = E/10"
      expect: "σ_final = σ_y + H·α; tangente = E·H/(E+H) en régimen plástico"
      tol_rel: 1.0e-10
    - name: ciclo_carga_descarga
      setup: "ε creciente hasta plástico, luego decreciente"
      expect: "descarga elástica con tangente E; α permanece constante durante descarga; histéresis correcta"
      tol_rel: 1.0e-10
    - name: cambio_de_signo
      setup: "tracción plástica luego compresión hasta fluencia inversa"
      expect: "fluencia inversa en σ = -(σ_y + H·α) (endurecimiento isótropo es simétrico)"
      tol_rel: 1.0e-10
    - name: monotonicidad_alpha
      setup: "trayectoria arbitraria"
      expect: "α no decrece nunca"

  # Cubierto por tests de integración con barras plásticas
  integration:
    - name: truss_uniaxial_plastico
      setup: "Truss2D con Elastoplastic1D, carga axial creciente"
      expect: "respuesta fuerza-desplazamiento bilineal exacta (E hasta fluencia, E_t después)"
      tol_rel: 1.0e-6

references:
  - "Simó J.C., Hughes T.J.R. (1998). Computational Inelasticity. Springer. §1 (J2 1D return mapping)."
  - "de Souza Neto E.A., Perić D., Owen D.R.J. (2008). Computational Methods for Plasticity. Wiley. §6.2 (uniaxial)."
  - "Lubliner J. (1990). Plasticity Theory. Macmillan. §3.2."

\end{lstlisting}


\section{Implementación}

\begin{itemize}
  \item \textbf{Archivo}: solidum/materials/plastic\_1d.py.
  \item \textbf{Clase}: \texttt{Elastoplastic1D}, registrada vía \texttt{@MaterialRegistry.register}.
  \item \textbf{\texttt{admissibility\_scale}}: $\sigma_y + H\,\alpha$ (ADR 0006). El switching elástico/plástico usa el patrón atol + rtol\ensuremath{\cdot}escala.
  \item \textbf{State schema}: \texttt{\{'eps\_p': float, 'alpha': float\}}. \texttt{state\_vars=None} se interpreta como estado virgen ($\varepsilon^p = 0$, $\alpha = 0$).
  \item \textbf{Tests}:
  \item tests/test\_materials\_unit.py \ensuremath{\cdot} clase \texttt{TestElastoplastic1D} (carga, descarga, endurecimiento, plasticidad perfecta).
  \item Integración: cualquier test de truss con material plástico.
\end{itemize}


\section{Diálogo}

\begin{itemize}
  \item \textbf{2026-05-19} \ensuremath{\cdot} Spec creada retroactivamente para cerrar el hueco H-5.3 de la auditoría 2026-05-18. Material anterior a la convención de specs; no se modifica código. La aceptación se considera cumplida por los tests unitarios existentes.
\end{itemize}


\newpage

\section{IsotropicDamage1D}



\textit{Orden de trabajo. El usuario escribe \textbf{especificación física}, \textbf{formulación} y \textbf{contrato}; la IA rellena \textbf{implementación} y responde en \textbf{diálogo}.}

\textit{> Spec \textbf{retroactiva con ampliación}: modelo implementado desde sesiones anteriores con tangente secante; esta spec lo documenta y añade la tangente algorítmica consistente, replicando la derivación de [[IsotropicDamage2D]] reducida a 1D escalar.}



\section{Especificación física}

Versión 1D de \texttt{IsotropicDamage2D}. Modelo de daño escalar isótropo en pequeñas deformaciones para barras, cables y elementos axiales. Captura degradación progresiva de la rigidez axial con ablandamiento exponencial.

\subsection{Ecuaciones}

\begin{itemize}
  \item \textbf{Constitutiva}: $\sigma = (1 - d)\,E\,\varepsilon$, con $\sigma^\text{eff} = E\varepsilon$.
  \item \textbf{Deformación equivalente}: $\varepsilon_\text{eq} = |\varepsilon|$.
  \item \textbf{Variable histórica}: $\kappa_{n+1} = \max(\kappa_n, |\varepsilon_{n+1}|)$, con $\kappa_0$ inicial.
  \item \textbf{Daño}: $d(\kappa) = \begin{cases}0 & \kappa \leq \kappa_0 \\ 1 - (\kappa_0/\kappa)\,e^{-\alpha(\kappa - \kappa_0)} & \kappa > \kappa_0\end{cases}$, saturado a \texttt{DAMAGE\_MAX}.
\end{itemize}

Todas las propiedades físicas (irreversibilidad, asintótica $d \to 1$, papel de $\alpha$) son las del modelo 2D restringidas a un escalar. No distingue tracción de compresión.


\section{Formulación numérica}

\subsection{Tangente algorítmica consistente}

\[
\frac{\partial \sigma}{\partial \varepsilon} = (1-d)\,E - E\varepsilon\,\frac{\partial d}{\partial \kappa}\,\frac{\partial \kappa}{\partial \varepsilon}
\]

con:

\begin{itemize}
  \item $\dfrac{\partial d}{\partial \kappa} = (1-d)\left(\dfrac{1}{\kappa} + \alpha\right)$ (idéntico a 2D).
  \item $\dfrac{\partial \kappa}{\partial \varepsilon} = \operatorname{sign}(\varepsilon)$ en \textbf{carga activa} ($|\varepsilon| > \kappa_n$); $0$ en descarga.
\end{itemize}

\textbf{Tangente final}:

\[
E^\text{alg} = \begin{cases}
E & \kappa \leq \kappa_0 \;(\text{sin daño}) \\
(1-d)\,E & \text{descarga} \\
(1-d)\,E\big[1 - (1/\kappa + \alpha)\,\kappa\big] = -(1-d)\,E\,\alpha\,\kappa & \text{carga activa, } d < d_\text{max} \\
(1-d_\text{max})\,E & d = d_\text{max} \;(\text{saturación})
\end{cases}
\]

donde se ha usado $(1/\kappa + \alpha)\kappa = 1 + \alpha\kappa$ y $E\varepsilon\,\operatorname{sign}(\varepsilon) = E|\varepsilon| = E\kappa$ en carga activa.

\subsection{Signo de la tangente consistente}

En carga activa \textbf{siempre es negativa}: $E^\text{alg} = -(1-d)\,E\,\alpha\,\kappa < 0$ porque $1-d > 0$, $E > 0$, $\alpha > 0$, $\kappa > 0$. Refleja físicamente el régimen de \textbf{softening} (ablandamiento): la curva $\sigma$--$\varepsilon$ desciende tras el umbral.

\textbf{Implicaciones numéricas}:

\begin{itemize}
  \item La rigidez tangente del elemento ($\mathbf K_e = \int B^\top E^\text{alg} B\, dV$) puede tener autovalores negativos. La matriz \textbf{sigue siendo simétrica} (escalar \ensuremath{\times} forma cuadrática) --- \texttt{IS\_SYMMETRIC = True} se mantiene.
  \item La matriz global puede dejar de ser \textbf{positiva definida}. El despachador algebraico (ADR 0003) detecta el fallo de Cholesky y degrada a LU automáticamente --- no requiere cambio declarativo.
  \item Control de carga puede no converger en problemas con suficiente daño para producir snap-back; en ese régimen se requiere control de longitud de arco (\texttt{ArcLengthSolver}).
\end{itemize}

\subsection{Por qué no tests de integración en esta spec}

A diferencia de [[IsotropicDamage2D]], no se añade benchmark de integración Truss2D + IsotropicDamage1D + NonlinearSolver: con un solo elemento axial el comportamiento global está dominado por la inestabilidad de softening ($E_{\text{alg}} < 0$) y la única vía limpia de seguir la respuesta post-pico es arc-length. La validación de la tangente consistente se cubre con tests unitarios (incluyendo diferencia finita centrada como verificación de la derivada cerrada). Si en el futuro aparece un caso real de problema con daño 1D ramificado, se añadirá un benchmark dedicado con \texttt{ArcLengthSolver}.


\section{Contrato de implementación}

\begin{lstlisting}[language=yaml]
name: IsotropicDamage1D
kind: material
status: validated

interface:
  strain_dim: 1
  primary_state_var: damage
  is_symmetric: true       # tangente escalar ⇒ contribución elemento simétrica (puede ser indefinida pero simétrica)

parameters:
  - { name: E,        type: float, required: true,  desc: "Módulo de Young intacto" }
  - { name: kappa_0,  type: float, required: true,  desc: "Umbral de deformación equivalente |ε|" }
  - { name: alpha,    type: float, required: true,  desc: "Velocidad de degradación exponencial (>0)" }
  - { name: density,  type: float, required: false, default: null, desc: "Densidad (ADR 0008)" }

signature:
  compute_state: "(ε: float, state_vars=None) -> (σ: float, E_tan: float, state_vars')"
  strain_kind: axial escalar

state_schema:
  kappa:  "float ≥ κ_0 — máxima |ε| histórica"
  damage: "float ∈ [0, DAMAGE_MAX]"

conventions:
  sign: "σ > 0 ⇔ tracción; el modelo no distingue tracción de compresión en la activación del daño"
  damage_irreversibility: "κ monótono no decreciente; d(κ) no decreciente"

validity:
  - "E > 0, κ_0 > 0, α > 0"
  - "|ε| ≲ 1e-2 (régimen de pequeñas deformaciones)"

out_of_scope:
  - "split tracción/compresión"
  - "fatiga / acumulación cíclica"
  - "regularización para mesh-dependency en problemas con varios elementos en serie"
  - "test de integración con NonlinearSolver — requiere arc-length para superar el snap-back, fuera de scope hasta caso real"

numerical_caveats:
  - "Tangente consistente NEGATIVA en carga activa (softening). La rigidez global puede no ser PD; el despachador algebraico (ADR 0003) degrada a LU automáticamente al detectar fallo de Cholesky."
  - "Control de carga con un elemento aislado tras alcanzar el pico se vuelve inestable. Usar arc-length para seguir la rama post-pico."
  - "Transición carga↔descarga discontinua en la tangente (rama loading vs unloading): aceptable para Newton convergente, problemático cerca del umbral en problemas patológicos."

acceptance:
  unit_existing_no_regression:
    - name: secant_response_below_threshold
      expect: "ε pequeño, κ ≤ κ_0 ⇒ d=0, σ=E·ε, tangente = E"
      tol_rel: 1.0e-12

    - name: damage_evolution_exponential
      expect: "ley exponencial reproducida con la fórmula"
      tol_rel: 1.0e-10

  unit_consistent_tangent:
    - name: tangent_equals_E_below_threshold
      setup: "ε < κ_0"
      expect: "E_tan = E (sin daño)"
      tol_rel: 1.0e-12

    - name: tangent_equals_secant_on_unloading
      setup: "cargar a κ > κ_0; descargar a |ε| < κ"
      expect: "E_tan = (1-d)·E exacto"
      tol_rel: 1.0e-12

    - name: tangent_negative_on_loading
      setup: "carga activa con κ > κ_0, d ∈ (0, DAMAGE_MAX)"
      expect: "E_tan = -(1-d)·E·α·κ < 0"
      tol_rel: 1.0e-12

    - name: tangent_equals_secant_at_saturation
      setup: "ε enorme tal que d = DAMAGE_MAX"
      expect: "E_tan = (1-DAMAGE_MAX)·E (sin término consistente al saturar)"
      tol_rel: 1.0e-12

    - name: tangent_matches_finite_difference
      setup: "carga activa; ∂σ/∂ε por FD centrada (h=1e-7) vs fórmula cerrada"
      expect: "error relativo < 1e-5"
      tol_rel: 1.0e-5

    - name: tangent_consistent_with_2D_when_uniaxial
      setup: "evaluar IsotropicDamage1D(E, κ_0, α) en ε=eps y comparar contra IsotropicDamage2D(E, ν=0, κ_0, α) plane stress en ε=[eps, 0, 0]; la primera componente de σ y la entrada [0,0] de C_alg deben coincidir"
      expect: "σ_1D = σ_xx_2D; E_tan_1D = C_alg_2D[0,0]"
      tol_rel: 1.0e-10

references:
  - "Lemaitre J., Chaboche J.-L. (1990). Mechanics of Solid Materials. Cambridge UP. Cap. 7."
  - "Simó J.C., Ju J.W. (1987). Strain- and stress-based continuum damage models — I. Int. J. Solids Struct. 23, 821-840."
  - "Spec hermana: [docs/specs/IsotropicDamage2D.md](IsotropicDamage2D.md) — modelo 2D, derivación detallada de la tangente consistente."

\end{lstlisting}


\section{Implementación}

\begin{itemize}
  \item \textbf{Archivo}: solidum/materials/damage\_1d.py.
  \item \textbf{Clase}: \texttt{IsotropicDamage1D}, registrada vía \texttt{@MaterialRegistry.register}. \texttt{IS\_SYMMETRIC} queda \texttt{True} (default heredado de \texttt{Material}) --- la contribución elemental es simétrica aunque pueda ser indefinida.
  \item \textbf{Algoritmo \texttt{compute\_state}}:
\end{itemize}
\begin{enumerate}
  \item Recuperar $\kappa_n$ y calcular $\varepsilon_\text{eq} = |\varepsilon|$.
  \item Si $\varepsilon_\text{eq} > \kappa_n$: carga activa, $\kappa_{n+1} = \varepsilon_\text{eq}$. Si no: descarga, $\kappa_{n+1} = \kappa_n$.
  \item Calcular $d$ por la ley exponencial, saturar a \texttt{DAMAGE\_MAX}.
  \item $\sigma = (1 - d)\,E\,\varepsilon$.
  \item Tangente: ramas (sin daño / descarga / saturación / carga activa) idénticas en estructura al modelo 2D, con la fórmula escalar resultante. En carga activa con daño efectivo no saturado, \texttt{E\_tan = -(1-d)\ensuremath{\cdot}E\ensuremath{\cdot}\ensuremath{\alpha}\ensuremath{\cdot}\ensuremath{\kappa}}.
\end{enumerate}
\begin{itemize}
  \item \textbf{Tests}:
  \item tests/test\_materials\_unit.py \ensuremath{\cdot} \texttt{TestIsotropicDamage1D} (existente, no regresión) + nueva clase \texttt{TestIsotropicDamage1DConsistentTangent}.
  \item Sin tests de integración (justificado arriba).
\end{itemize}


\section{Diálogo}

\begin{itemize}
  \item \textbf{2026-05-14} \ensuremath{\cdot} Spec creada retroactivamente sobre material existente (tangente secante) y extendida con tangente algorítmica consistente. Réplica directa de [[IsotropicDamage2D]] reducida a escalar. Sin ADR --- extensión local al material, no afecta a contratos.
  \item \textbf{2026-05-14} \ensuremath{\cdot} \texttt{IS\_SYMMETRIC} permanece \texttt{True}: la rigidez axial es escalar y la contribución elemental \texttt{B\textasciicircum{}T\ensuremath{\cdot}E\_tan\ensuremath{\cdot}B} es simétrica aunque pueda ser indefinida cuando \texttt{E\_tan < 0}. El sistema global puede dejar de ser PD; el fallback automático de Cholesky a LU del despachador algebraico (ADR 0003) lo gestiona sin cambios declarativos.
  \item \textbf{2026-05-14} \ensuremath{\cdot} No se añade test de integración. Justificación: con tangente consistente <0 en carga activa, un benchmark Truss2D bajo control de carga se vuelve inestable tras el pico; el seguimiento de la rama post-pico requiere arc-length que cae fuera del scope inmediato. Tests unitarios (incluyendo FD numérica y consistencia con la versión 2D reducida) cubren la corrección de la fórmula.
\end{itemize}


\newpage

\chapter{Modelos Constitutivos — 2D}

\section{DruckerPrager2D}



\textit{Orden de trabajo. El usuario escribe \textbf{especificación física}, \textbf{formulación numérica} y \textbf{contrato}; la IA rellena \textbf{implementación} y responde en \textbf{diálogo}.}



\section{Especificación física}

\subsection{0. Descripción general}

Modelo elastoplástico \textbf{friccional} independiente de la velocidad. Suelos, hormigón, granulares, materiales cuasi-frágiles cohesivo-friccionales. Suaviza la pirámide hexagonal de Mohr-Coulomb a un \textbf{cono circular} en el espacio de tensiones principales, eliminando las aristas que complican el return mapping.

Diferencia clave con J2:

\begin{itemize}
  \item \textbf{Dependencia de la presión}. La resistencia al corte aumenta con la presión de confinamiento (componente friccional).
  \item \textbf{Plasticidad no asociada}. La dilatancia volumétrica real del flujo plástico ($\psi$) es típicamente menor que la fricción del criterio ($\phi$). Forzar $\psi = \phi$ (asociada) sobreestima la dilatancia y el volumen post-fluencia.
  \item \textbf{Retorno al ápice}. El criterio define un cono con vértice; estados de prueba en zona puramente hidrostática extrema retornan al vértice (zona singular en la dirección de flujo).
\end{itemize}

Hipótesis cinemática soportada: \textbf{\texttt{plane\_strain}} únicamente en esta spec.

\subsection{1. Descomposición aditiva y elasticidad}

Pequeñas deformaciones: $\boldsymbol\varepsilon = \boldsymbol\varepsilon^e + \boldsymbol\varepsilon^p$. Ley elástica isótropa: $\boldsymbol\sigma = \mathbf C_e (\boldsymbol\varepsilon - \boldsymbol\varepsilon^p)$.

\subsection{2. Criterio de fluencia (Drucker-Prager)}

\[
f(\boldsymbol\sigma, \alpha) = \sqrt{J_2} + \eta_f\,I_1 - k(\alpha) \;\leq\; 0
\]

con:

\begin{itemize}
  \item $J_2 = \tfrac{1}{2}\,\mathbf s : \mathbf s$, segundo invariante del desviador.
  \item $I_1 = \operatorname{tr}(\boldsymbol\sigma) = \sigma_{xx} + \sigma_{yy} + \sigma_{zz}$, primer invariante.
  \item $\eta_f$, coeficiente friccional del cono (adimensional, $\geq 0$).
  \item $k(\alpha) = k_0 + H_k\,\alpha$, parámetro cohesivo con endurecimiento isótropo lineal.
\end{itemize}

La superficie es un cono con eje en la línea hidrostática ($\sigma_{xx} = \sigma_{yy} = \sigma_{zz}$), vértice en $I_1 = k/\eta_f$, abertura controlada por $\eta_f$.

\subsection{3. Calibración con Mohr-Coulomb}

Los parámetros $(\eta_f, k_0)$ se derivan de los parámetros físicos $(c_0, \phi)$ --- cohesión y ángulo de fricción interna --- según la variante del cono elegida:

\begin{center}
\begin{tabularx}{\textwidth}{YYY}
\hline
\textbf{variant} & \textbf{$\eta_f$} & \textbf{$k_0$} \\
\hline
\texttt{plane\_strain\_matched} (default) & $\tan(\phi)/\sqrt{9 + 12\tan^2(\phi)}$ & $3\,c_0/\sqrt{9 + 12\tan^2(\phi)}$ \\
\texttt{outer\_cone} & $2\sin(\phi)/[\sqrt 3\,(3 - \sin(\phi))]$ & $6\,c_0\cos(\phi)/[\sqrt 3\,(3 - \sin(\phi))]$ \\
\texttt{inner\_cone} & $2\sin(\phi)/[\sqrt 3\,(3 + \sin(\phi))]$ & $6\,c_0\cos(\phi)/[\sqrt 3\,(3 + \sin(\phi))]$ \\
\hline
\end{tabularx}
\end{center}
\texttt{plane\_strain\_matched} es la variante que coincide exactamente con Mohr-Coulomb en condiciones de plane strain (Drucker, 1953). Las otras dos circunscriben (outer) o inscriben (inner) el hexágono de MC en el plano \ensuremath{\pi}. Como esta spec solo soporta plane strain, \texttt{plane\_strain\_matched} es el default natural.

\subsection{4. Regla de flujo (no asociada por defecto)}

Potencial plástico $g$ con coeficiente de \textbf{dilatancia} $\eta_g$ (función del ángulo de dilatancia $\psi$, calibrado con la misma fórmula que $\eta_f$ pero usando $\psi$):

\[
g(\boldsymbol\sigma) = \sqrt{J_2} + \eta_g\,I_1
\]

\[
\dot{\boldsymbol\varepsilon}^p = \dot\gamma\,\frac{\partial g}{\partial \boldsymbol\sigma} = \dot\gamma\,\left(\frac{\mathbf s}{2\sqrt{J_2}} + \eta_g\,\mathbf I\right)
\]

Descomposición:
\begin{itemize}
  \item \textbf{Parte desviadora}: $\dot{\mathbf e}^p = \dot\gamma\,\mathbf s/(2\sqrt{J_2})$ --- análoga a J2.
  \item \textbf{Parte volumétrica}: $\operatorname{tr}(\dot{\boldsymbol\varepsilon}^p) = 3\dot\gamma\,\eta_g$ --- \textbf{dilatación plástica positiva} si $\eta_g > 0$.
\end{itemize}

Si $\psi = \phi$ (i.e. $\eta_g = \eta_f$) el modelo se vuelve \textbf{asociado}; tangente algorítmica simétrica. Para $\psi < \phi$ (típico en suelos: $\psi \approx \phi/2$ o incluso $\psi = 0$, "non-dilatant"), el modelo es \textbf{no asociado} y la tangente algorítmica es \textbf{asimétrica}.

\subsection{5. Endurecimiento isótropo lineal}

\[
\dot\alpha = \dot\gamma
\]

con $\alpha$ la \textbf{deformación plástica equivalente} (convención del proyecto para Drucker-Prager). La superficie crece manteniendo eje y abertura, solo $k$ aumenta linealmente.

Endurecimiento por fricción/dilatancia (cambio de $\eta_f$ o $\eta_g$ con $\alpha$) queda fuera de esta spec --- escenario que no necesitamos hoy y que añade no-linealidad fuerte al return mapping.

\subsection{6. Condiciones de Kuhn-Tucker}

\[
\dot\gamma \geq 0,\qquad f \leq 0,\qquad \dot\gamma\,f = 0
\]

\subsection{7. Variables internas}

\begin{center}
\begin{tabularx}{\textwidth}{YYY}
\hline
\textbf{Símbolo} & \textbf{Tipo} & \textbf{Significado} \\
\hline
$\boldsymbol\varepsilon^p$ & tensor 4 componentes \texttt{[xx, yy, zz, xy\_tensorial]} & deformación plástica \\
$\alpha$ & escalar $\geq 0$ & deformación plástica equivalente (= acumulado de $\Delta\gamma$) \\
\hline
\end{tabularx}
\end{center}
\texttt{alpha} es la variable principal exportable (\texttt{PRIMARY\_STATE\_VAR = 'alpha'}).


\section{Formulación numérica}

\subsection{8. Esquema implícito --- Backward Euler}

Dado $\{\boldsymbol\varepsilon^p_n, \alpha_n\}$ y $\boldsymbol\varepsilon_{n+1}$, predictor elástico:

\[
\boldsymbol\varepsilon^\text{trial} = \boldsymbol\varepsilon_{n+1} - \boldsymbol\varepsilon^p_n, \qquad \mathbf s^\text{trial} = 2G\,\mathbf e^\text{trial}_\text{dev},\qquad p^\text{trial} = K\,\operatorname{tr}(\boldsymbol\varepsilon^\text{trial})
\]

con $I_1^\text{trial} = 3\,p^\text{trial}$, $\sqrt{J_2^\text{trial}} = \|\mathbf s^\text{trial}\|/\sqrt{2}$.

Función de fluencia trial:

\[
f^\text{trial} = \sqrt{J_2^\text{trial}} + \eta_f\,I_1^\text{trial} - k(\alpha_n)
\]

Si $f^\text{trial} \leq \text{tol}$ \ensuremath{\to} paso elástico. Si no, intentar \textbf{return regular}.

\subsection{9. Return regular (cone surface)}

La dirección desviadora $\hat{\mathbf n} = \mathbf s^\text{trial}/\sqrt{2 J_2^\text{trial}} = \mathbf s/\sqrt{2 J_2}$ es invariante bajo carga radial. Las actualizaciones son lineales en $\Delta\gamma$:

\[
\sqrt{J_2^{n+1}} = \sqrt{J_2^\text{trial}} - G\,\Delta\gamma
\]
\[
p^{n+1} = p^\text{trial} - 3K\,\eta_g\,\Delta\gamma, \qquad I_1^{n+1} = 3\,p^{n+1}
\]

La consistencia $f(\sigma_{n+1}, \alpha_{n+1}) = 0$ con $\alpha_{n+1} = \alpha_n + \Delta\gamma$ y $k_{n+1} = k_0 + H_k(\alpha_n + \Delta\gamma)$ da:

\[
\boxed{\;\Delta\gamma = \frac{f^\text{trial}}{G + 9K\,\eta_f\,\eta_g + H_k}\;}
\]

\textbf{Forma cerrada} --- sin Newton local. Cada término del denominador:
\begin{itemize}
  \item $G$: cambio en $\sqrt{J_2}$ por la parte desviadora del flujo.
  \item $9K\,\eta_f\,\eta_g$: cambio en $\eta_f I_1$ por la parte volumétrica dilatante del flujo (multiplicado por $\eta_f$ del criterio y $\eta_g$ del potencial).
  \item $H_k$: endurecimiento.
\end{itemize}

\textbf{Actualizaciones}:

\[
\mathbf s^{n+1} = \frac{\sqrt{J_2^{n+1}}}{\sqrt{J_2^\text{trial}}}\,\mathbf s^\text{trial}
\]

\[
\boldsymbol\sigma^{n+1} = \mathbf s^{n+1} + p^{n+1}\,\mathbf I
\]

\[
\boldsymbol\varepsilon^p_{n+1} = \boldsymbol\varepsilon^p_n + \Delta\gamma\,\left(\frac{\mathbf s^{n+1}}{2\sqrt{J_2^{n+1}}} + \eta_g\,\mathbf I\right)
\]

\textbf{Condición de validez del return regular}: $\sqrt{J_2^{n+1}} \geq 0$. Si $\Delta\gamma_\text{reg} > \sqrt{J_2^\text{trial}}/G$, el desviador "se pasa de cero" y el algoritmo regular es inválido --- toca \textbf{return al ápice}.

\subsection{10. Return al ápice}

Cuando el estado de prueba cae en zona puramente hidrostática extrema (p\_trial muy traccionante respecto a la cohesión), la proyección al cono cae en su vértice, donde $\mathbf s = 0$ y la dirección desviadora $\hat{\mathbf n}$ no está definida.

En el ápice: $\boldsymbol\sigma^{n+1} = (k_{n+1}/\eta_f)\,\mathbf I$ (puramente hidrostático), $\sqrt{J_2^{n+1}} = 0$.

Por conservación volumétrica:

\[
p^{n+1} = p^\text{trial} - 3K\,\eta_g\,\Delta\gamma = \frac{k_{n+1}}{3\,\eta_f}
\]

con $k_{n+1} = k_0 + H_k(\alpha_n + \Delta\gamma)$. Despejando:

\[
\boxed{\;\Delta\gamma_\text{apex} = \frac{p^\text{trial}\,\eta_f - k(\alpha_n)/3 \cdot \eta_f \cdot 3}{3K\,\eta_f\,\eta_g + H_k} = \frac{3 p^\text{trial}\,\eta_f - k(\alpha_n)}{3K \cdot 3\eta_f\,\eta_g + H_k}\;}
\]

Reordeno: $3p^\text{trial}\,\eta_f - k(\alpha_n) = I_1^\text{trial}\,\eta_f - k(\alpha_n)$. Denominador: $9K\,\eta_f\,\eta_g + H_k$. (Equivalente a la fórmula regular sin el término $G$, porque la rama de apex es puramente volumétrica.)

\[
\Delta\gamma_\text{apex} = \frac{I_1^\text{trial}\,\eta_f - k(\alpha_n)}{9K\,\eta_f\,\eta_g + H_k}
\]

\textbf{Actualizaciones (apex)}:

\[
\mathbf s^{n+1} = 0,\qquad p^{n+1} = K(\operatorname{tr}(\boldsymbol\varepsilon_{n+1}) - 3\,\eta_g\,\Delta\gamma) \;\equiv\; \frac{k_{n+1}}{3\,\eta_f}
\]

\[
\boldsymbol\sigma^{n+1} = p^{n+1}\,\mathbf I
\]

\[
\boldsymbol\varepsilon^p_{n+1} = \boldsymbol\varepsilon^p_n + \Delta\gamma\,\eta_g\,\mathbf I + \boldsymbol\varepsilon^{p,\text{dev}}_\text{trial}
\]

(la parte desviadora plástica absorbe completamente el desviador de prueba, ya que $\mathbf s^{n+1} = 0$.)

\subsection{11. Detección del régimen de retorno}

Algoritmo de despacho:

\begin{enumerate}
  \item Calcular $\Delta\gamma_\text{reg}$ por la fórmula regular.
  \item Si $\sqrt{J_2^\text{trial}} - G\,\Delta\gamma_\text{reg} \geq 0$ \ensuremath{\to} return regular.
  \item Si no \ensuremath{\to} return al ápice con $\Delta\gamma_\text{apex}$.
\end{enumerate}

El criterio es geométricamente: el predictor cae más allá del ápice si la proyección desviadora exigiría reducir $\sqrt{J_2}$ por debajo de cero.

\subsection{12. Tangente algorítmica consistente}

\textbf{Return regular} ($\sqrt{J_2^{n+1}} > 0$):

Por linealización del algoritmo:

\[
\mathbf C^\text{alg} = K\,\mathbf 1\otimes\mathbf 1 + 2G\left(1 - \frac{G\,\Delta\gamma}{\sqrt{J_2^\text{trial}}}\right)\mathbf I_\text{dev} + 2G\,\Theta\,\hat{\mathbf n}\otimes\hat{\mathbf n} - 3K\,\eta_g\,\mathbf 1\otimes \mathbf a_f - 3K\,\eta_f\,\mathbf a_g\otimes\mathbf 1 + \text{corrección}
\]

donde:
\begin{itemize}
  \item $\hat{\mathbf n} = \mathbf s^{n+1}/(2\sqrt{J_2^{n+1}})$ (dirección de flujo desviador, normalizada por convención).
  \item $\mathbf 1 = (1,1,1,0)$ en notación 4 componentes (hidrostática).
  \item $\Theta$, $\mathbf a_f$, $\mathbf a_g$: términos cruzados que recogen el acoplamiento entre cambio de $\sqrt{J_2}$ y endurecimiento.
\end{itemize}

Forma compacta (de Souza Neto §8.3.4):

\[
\mathbf C^\text{alg} = K\,\mathbf 1\otimes\mathbf 1 + 2G\,\Big(1 - \tfrac{G\Delta\gamma}{\sqrt{J_2^\text{trial}}}\Big)\,\mathbf I_\text{dev} - \frac{1}{G + 9K\eta_f\eta_g + H_k}\,\mathbf b_g \otimes \mathbf b_f
\]

con $\mathbf b_g = G\hat{\mathbf n} + 3K\eta_g\,\mathbf 1$, $\mathbf b_f = G\hat{\mathbf n} + 3K\eta_f\,\mathbf 1$. \textbf{Asimétrica si $\eta_f \neq \eta_g$} (no asociada).

\textbf{Return al ápice} ($\sqrt{J_2^{n+1}} = 0$):

\[
\mathbf C^\text{alg} = K\,\Big(1 - \frac{3K\eta_f\eta_g}{3K\eta_f\eta_g + H_k/3}\Big)\,\mathbf 1\otimes\mathbf 1
\]

--- rigidez puramente volumétrica reducida (el material en el ápice se comporta como un fluido con compresibilidad finita; sin rigidez desviadora alguna).

Forma equivalente más limpia: $\mathbf C^\text{alg}_\text{apex} = \dfrac{K\,H_k}{9K\eta_f\eta_g + H_k}\,\mathbf 1\otimes\mathbf 1$ (con $H_k=0$, el ápice se vuelve completamente plástico volumétrico y la tangente colapsa; este caso queda manejado en el código devolviendo un escalar pequeño $>0$ para evitar singularidad).

\subsection{13. Tolerancia de fluencia}

ADR 0006: \texttt{admissibility\_scale = k(\ensuremath{\alpha}) = k\_0 + H\_k\ensuremath{\cdot}\ensuremath{\alpha}} (escala cohesiva corriente del cono). Tiene unidades de esfuerzo, igual que $f$.


\section{Contrato de implementación}

\begin{lstlisting}[language=yaml]
name: DruckerPrager2D
kind: material
status: validated

interface:
  strain_dim: 3
  primary_state_var: alpha
  is_symmetric: false      # tangente asimétrica si no asociada (ψ ≠ φ)

parameters:
  - { name: E,          type: float, required: true,  desc: "Módulo de Young" }
  - { name: nu,         type: float, required: true,  desc: "Coeficiente de Poisson" }
  - { name: cohesion,   type: float, required: true,  desc: "Cohesión c_0 inicial (esfuerzo)" }
  - { name: phi_deg,    type: float, required: true,  desc: "Ángulo de fricción interna en grados; típico suelos 20-40°, hormigón 30-37°" }
  - { name: psi_deg,    type: float, required: false, default: null, desc: "Ángulo de dilatancia en grados. Default = phi_deg (asociada). Típico no asociada: ψ < φ, e.g. 0 o φ/2." }
  - { name: H,          type: float, required: false, default: 0.0, desc: "Endurecimiento isótropo lineal en cohesión H_k (≥0). H=0 ⇒ perfectamente plástico." }
  - { name: hypothesis, type: str,   required: false, default: "plane_strain", desc: "plane_strain (único soportado)" }
  - { name: variant,    type: str,   required: false, default: "plane_strain_matched", desc: "plane_strain_matched | outer_cone | inner_cone — calibración con Mohr-Coulomb" }
  - { name: density,    type: float, required: false, default: null, desc: "Densidad (ADR 0008)" }

signature:
  compute_state: "(ε: ndarray(3,), state_vars=None) -> (σ: ndarray(3,), C_alg: ndarray(3,3), state_vars')"
  strain_kind: "plano Voigt — [ε_xx, ε_yy, γ_xy] con γ_xy = 2·ε_xy"

state_schema:
  eps_p:   "ndarray(4,) — [ε^p_xx, ε^p_yy, ε^p_zz, ε^p_xy_tensorial]"
  alpha:   "float ≥ 0 — multiplicador plástico acumulado"

conventions:
  sign: "σ > 0 ⇔ tracción (Reglas §5)"
  voigt: "γ_xy = 2·ε_xy en deformación; ε^p_xy en tensorial (sin factor 2)"
  associated_flow: "ψ = φ ⇒ asociada (tangente simétrica); ψ ≠ φ ⇒ no asociada"

validity:
  - "E > 0, c_0 > 0, 0 ≤ φ < 90°, 0 ≤ ψ ≤ φ"
  - "ν ∈ (-1, 0.5)"
  - "H ≥ 0"
  - "hypothesis == 'plane_strain' (único soportado en esta spec)"
  - "variant ∈ {'plane_strain_matched', 'outer_cone', 'inner_cone'}"

out_of_scope:
  - "plane_stress — la proyección de Drucker-Prager con σ_zz=0 acoplada a regla de flujo dilatante es notoriamente delicada; difiere hasta caso real"
  - "endurecimiento por cambio de φ y/o ψ con α — no lineal fuerte en el return mapping"
  - "endurecimiento cinemático (back-stress hidrostático y desviador)"
  - "cap (modelo cap-cone) para suelos compactables"
  - "ablandamiento (H<0) — requiere regularización"
  - "grandes deformaciones"

numerical_caveats:
  - "Tangente algorítmica asimétrica si ψ ≠ φ → IS_SYMMETRIC=False, el despachador algebraico usa LU (ADR 0003)."
  - "Modos de carga muy traccionantes pueden caer en la rama de ápice. La detección automática del régimen (regular vs apex) está en el algoritmo; si en el futuro aparece oscilación entre ramas en pasos cerca de la transición, considerar smoothing con vertex regularization."
  - "Con ψ = 0 (sin dilatancia) y H_k = 0, en la rama de ápice la tangente volumétrica se anula → singularidad. El código devuelve un valor de respaldo (rigidez volumétrica reducida no nula) para evitar matriz singular; típicamente sale del ápice en el siguiente paso."
  - "El return regular puede generar α creciente con incrementos pequeños cerca de la transición regular↔apex. Newton global converge igual; si en problema concreto se detecta oscilación, usar arc-length."

acceptance:
  unit:
    - name: parameter_validation
      expect: "constructor rechaza variant inválida, φ ≥ 90°, ψ > φ, hypothesis distinto de plane_strain"

    - name: calibration_matched_consistency
      expect: "para variant='plane_strain_matched' con ψ = φ, asociada → η_f = η_g; con ψ = 0, η_g = 0 (sin dilatación)"

    - name: paso_elastico_below_yield
      expect: "ε pequeño; σ = C_e·ε, eps_p y alpha intactos"
      tol_rel: 1.0e-10

    - name: regular_return_under_shear
      setup: "ε cortante puro suficiente para fluir (ε_xy = γ/2, sin parte hidrostática)"
      expect: "return regular activo, α > 0, √J₂ del estado final = √(2/3)·k - η_f·I₁ ≥ 0 (cumple criterio)"
      tol_rel: 1.0e-8

    - name: apex_return_under_pure_tension
      setup: "ε puramente hidrostática traccionante (ε_xx = ε_yy > 0, sin cortante)"
      expect: "return al ápice; σ_xx = σ_yy = k/(3η_f) tras converger; √J₂ = 0"
      tol_rel: 1.0e-8

    - name: associated_tangent_is_symmetric
      setup: "ψ = φ (asociada), estado plástico activo"
      expect: "C_alg simétrica (‖C_alg − C_alg^T‖ < tol)"
      tol_abs: 1.0e-8

    - name: nonassociated_tangent_is_asymmetric
      setup: "ψ = 0, φ = 30°, estado plástico activo no degenerado"
      expect: "C_alg − C_alg^T ≠ 0 con norma significativa"

    - name: tangent_matches_finite_difference
      setup: "estado plástico regular; comparar C_alg cerrada con diferencia finita centrada"
      expect: "error relativo < 1e-5"
      tol_rel: 1.0e-5

    - name: alpha_monotonic_under_increasing_load
      expect: "bajo carga creciente α no decrece"

    - name: unloading_recovers_C_e
      setup: "cargar a plástico, descargar elásticamente"
      expect: "tangente = C_e en el paso de descarga, α se conserva"
      tol_rel: 1.0e-10

    - name: unit_invariance_MPa_vs_Pa
      expect: "mismo path adimensional en (MPa,mm,N) y (Pa,m,N) ⇒ α idéntico"
      tol_rel: 1.0e-12

  integration:
    - name: quad4_pipeline_elastico
      setup: "Quad4 unitario plane strain, carga muy por debajo del yield"
      expect: "respuesta elástica pura, α=0"
      tol_rel: 1.0e-8

    - name: quad4_pipeline_plastico_converges
      setup: "Quad4 unitario plane strain, carga que produce plasticidad activa, 4 pasos, max_iter=8"
      expect: "Newton converge en cada paso, α > 0 al final"

    - name: quad4_yield_onset_confined_tension          # DP-1 (2026-05-19)
      setup: "Quad4 unitario plane strain con desplazamientos prescritos en los 4 nodos para ε_xx=k·ε_y, ε_yy=0, γ_xy=0"
      expect: "ε_xx=0.5·ε_y elástico (σ_xx coincide con (λ+2μ)·ε exacto, α=0); ε_xx=0.99·ε_y aún α=0; ε_xx=1.5·ε_y plástico con α>0 idéntico en los 4 Gauss points"
      tol_rel: 1.0e-8

    - name: quad4_flow_rule_invariant                   # DP-2 (2026-05-19)
      setup: "Quad4 unitario plane strain, carga monótona en rama regular (5 niveles 1.2·ε_y → 2.6·ε_y), ψ ≠ φ"
      expect: "invariante cinemática tr(ε_p) = 3·η_g·α en cada Gauss point y en cada nivel"
      tol_rel: 1.0e-8

    - name: quad4_apex_under_biaxial_tension            # DP-2bis (2026-05-19)
      setup: "Quad4 unitario plane strain, ε_xx = ε_yy ~ 1e-2 (predictor hidrostático extremo)"
      expect: "estado final cae en ápice: σ_xx ≈ σ_yy = k(α)/(3·η_f), σ_xy ≈ 0, en cada Gauss point"
      tol_rel: 1.0e-6

references:
  - "Drucker D.C., Prager W. (1952). Soil mechanics and plastic analysis or limit design. Quart. Appl. Math. 10, 157-165."
  - "Drucker D.C. (1953). Coulomb friction, plasticity, and limit loads. J. Appl. Mech. 21, 71-74. (calibración plane_strain_matched)"
  - "Simó J.C., Hughes T.J.R. (1998). Computational Inelasticity. Springer. §6 (cone plasticity, apex return)."
  - "de Souza Neto E.A., Perić D., Owen D.R.J. (2008). Computational Methods for Plasticity. Wiley. §8 (Drucker-Prager, tangentes algorítmicas detalladas)."
  - "Chen W.F., Han D.J. (1988). Plasticity for Structural Engineers. Springer. §7 (cohesivo-friccionales)."

\end{lstlisting}


\section{Implementación}

\begin{itemize}
  \item \textbf{Archivo}: solidum/materials/drucker\_prager\_2d.py.
  \item \textbf{Clase}: \texttt{DruckerPrager2D}, registrada vía \texttt{@MaterialRegistry.register}. \texttt{IS\_SYMMETRIC = False} declarativo (la tangente puede ser asimétrica; el despachador agrega el flag sobre todos los materiales del dominio).
  \item \textbf{Kernel Numba} \texttt{\_compute\_drucker\_prager\_plane\_strain}:
\end{itemize}
\begin{enumerate}
  \item Predictor elástico desviador-volumétrico (paralelo a J2 plane strain, $\varepsilon^p_{zz}$ libre).
  \item Check $f^\text{trial}$.
  \item Si elástico: devolver \ensuremath{\sigma}\_trial restaurando contribución hidrostática y volumétrica.
  \item Si no: intentar $\Delta\gamma_\text{reg}$ cerrado; comprobar $\sqrt{J_2^{n+1}} \geq 0$.
  \item Si regular válido: actualizar \ensuremath{\sigma}, $\boldsymbol\varepsilon^p$, $\alpha$, tangente regular (de Souza Neto §8.3.4).
  \item Si regular inválido: rama de ápice --- $\Delta\gamma_\text{apex}$ cerrado, \ensuremath{\sigma} puramente hidrostático, tangente volumétrica reducida.
\end{enumerate}
\begin{itemize}
  \item \textbf{Cálculo de $(\eta_f, k_0, \eta_g)$}: en \texttt{\_\_init\_\_} desde $(c_0, \phi, \psi, \text{variant})$ --- sin coste runtime. Se valida $\psi \leq \phi$ (físicamente sensato).
  \item \textbf{\texttt{admissibility\_scale}}: $k(\alpha) = k_0 + H \cdot \alpha$.
  \item \textbf{Tests}:
  \item tests/test\_materials\_unit.py \ensuremath{\cdot} \texttt{TestDruckerPrager2D}.
  \item tests/test\_solid\_2d\_drucker\_prager.py.
\end{itemize}


\section{Diálogo}

\begin{itemize}
  \item \textbf{2026-05-14} \ensuremath{\cdot} Spec creada. Decisiones clave del alcance MVP:
  \item Solo \texttt{plane\_strain} en esta entrega; \texttt{plane\_stress} con cono en \ensuremath{\sigma}\_zz=0 es notoriamente delicado y de uso geotécnico raro --- se difiere.
  \item Parámetros físicos \texttt{(c\_0, \ensuremath{\varphi}, \ensuremath{\psi})} en lugar de adimensionales \texttt{(k\_0, \ensuremath{\eta}\_f, \ensuremath{\eta}\_g)} --- más natural para el usuario de geotecnia. Calibración interna por \texttt{variant}.
  \item \textbf{No asociada por defecto} (\ensuremath{\psi} es parámetro independiente con default = \ensuremath{\varphi} que la fuerza a asociada solo si el usuario no especifica). Plasticidad asociada en suelos sobreestima dilatancia.
  \item Endurecimiento solo en cohesión (lineal); fricción/dilatancia constantes.
  \item \texttt{IS\_SYMMETRIC = False} declarativo a nivel material (independiente del caso particular de uso; conservador y consistente con el patrón del proyecto: el despachador agrega).
  \item \textbf{2026-05-14} \ensuremath{\cdot} Sin ADR. El modelo encaja en el patrón establecido por \texttt{VonMises2D} (kernel Numba, despacho hypothesis aunque solo soporte una, semántica trial/commit) y no rompe contratos transversales.
\end{itemize}


\newpage

\section{Elastic2D}



\textit{Spec \textbf{retroactiva}: el material existe desde el origen del subsistema 2D y está validado por tests preexistentes (todos los pipelines sólido 2D elástico). Esta spec lo documenta sin cambiar comportamiento --- H-5.3 de la auditoría 2026-05-18.}



\section{Especificación física}

\subsection{0. Descripción general}

Modelo elástico lineal isótropo bidimensional, parametrizado por $(E, \nu)$. Soporta dos hipótesis cinemáticas 2D mutuamente excluyentes seleccionadas en construcción:

\begin{itemize}
  \item \textbf{\texttt{plane\_stress}} (default) --- $\sigma_{zz} = \sigma_{xz} = \sigma_{yz} = 0$. Láminas o membranas delgadas en su plano (chapa traccionada, placa fina).
  \item \textbf{\texttt{plane\_strain}} --- $\varepsilon_{zz} = \varepsilon_{xz} = \varepsilon_{yz} = 0$. Cuerpos prismáticos largos cargados transversalmente (presa, túnel, cilindro largo).
\end{itemize}

Independiente de la velocidad, sin variables históricas. Es el material elástico canónico para todos los sólidos 2D del catálogo.

\subsection{1. Descomposición de la deformación}

No aplica --- toda la deformación es elástica:

\[
\boldsymbol\varepsilon = \boldsymbol\varepsilon^e
\]

\subsection{2. Ley elástica}

\[
\boldsymbol\sigma = \mathbf C_e\,\boldsymbol\varepsilon
\]

con $\boldsymbol\varepsilon = [\varepsilon_{xx},\,\varepsilon_{yy},\,\gamma_{xy}]^\top$ en notación Voigt del proyecto ($\gamma_{xy} = 2\varepsilon_{xy}$) y $\boldsymbol\sigma = [\sigma_{xx},\,\sigma_{yy},\,\sigma_{xy}]^\top$.

\textbf{Matriz constitutiva plane stress} (3\ensuremath{\times}3):

\[
\mathbf C^{ps}_e = \frac{E}{1-\nu^2}\begin{bmatrix} 1 & \nu & 0 \\ \nu & 1 & 0 \\ 0 & 0 & (1-\nu)/2 \end{bmatrix}
\]

\textbf{Matriz constitutiva plane strain} (3\ensuremath{\times}3):

\[
\mathbf C^{pe}_e = \frac{E}{(1+\nu)(1-2\nu)}\begin{bmatrix} 1-\nu & \nu & 0 \\ \nu & 1-\nu & 0 \\ 0 & 0 & (1-2\nu)/2 \end{bmatrix}
\]

Tangente constante $\mathbf C_e$ en todo régimen. Simétrica y positiva definida en el rango admisible de $\nu$.

\subsection{3. Superficie de fluencia / criterio de daño}

No aplica --- el material es indefinidamente elástico.

\subsection{4-6. Flujo, endurecimiento, Kuhn-Tucker}

No aplican.

\subsection{7. Variables internas}

Ninguna --- \texttt{state\_vars = None} se acepta y se devuelve intacto. Sin \texttt{PRIMARY\_STATE\_VAR}.

\subsection{8. Notación Voigt}

Convención del proyecto: $\boldsymbol\varepsilon = [\varepsilon_{xx},\,\varepsilon_{yy},\,\gamma_{xy}]$ con $\gamma_{xy} = 2\varepsilon_{xy}$ (deformación angular \textit{engineering}). Los esfuerzos son los componentes tensoriales reales $[\sigma_{xx},\,\sigma_{yy},\,\sigma_{xy}]$.


\section{Formulación numérica}

\subsection{9. Esquema temporal}

No aplica --- sin historia.

\subsection{10. Return mapping / actualización}

No aplica. Evaluación reducida a producto matriz-vector:

\[
\boldsymbol\sigma_{n+1} = \mathbf C_e\,\boldsymbol\varepsilon_{n+1}, \qquad \mathbf C^{\text{alg}} = \mathbf C_e
\]

La matriz $\mathbf C_e$ se precalcula y guarda en el constructor --- coste de \texttt{compute\_state} es un producto 3\ensuremath{\times}3 \ensuremath{\cdot} 3 sin alocación adicional.

\subsection{11. Tangente algorítmica consistente}

\[
\mathbf C^{\text{alg}} = \mathbf C_e
\]

Simétrica (\texttt{IS\_SYMMETRIC = True} por default). En plane stress la matriz es PD para $\nu \in (-1, 0.5)$; en plane strain también, con singularidad en $\nu \to 0.5$ (incompresible --- requiere formulación mixta, no soportada).

\subsection{12. Caveats numéricos}

\begin{itemize}
  \item \textbf{Plane strain $\nu \to 0.5$}: la matriz constitutiva diverge ($1-2\nu \to 0$). Recomendado $\nu \le 0.49$ para metales (típico $\nu=0.3$). El constructor rechaza $\nu \ge 0.5$ con \texttt{ValueError}.
  \item \textbf{Locking volumétrico} en plane strain con elementos de bajo orden (Quad4, Tri3) cuando $\nu \to 0.5$: declarado limitación arquitectural, sin mitigación implementada (B-bar, F-bar diferidas).
  \item Para $\nu < 0$ el material es admisible termodinámicamente (auxético) pero infrecuente en la práctica; el constructor lo permite hasta $\nu > -1$.
\end{itemize}


\section{Contrato de implementación}

\begin{lstlisting}[language=yaml]
name: Elastic2D
kind: material
status: validated

interface:
  strain_dim: 3
  primary_state_var: null
  is_symmetric: true

parameters:
  - { name: E,          type: float, required: true,  desc: "Módulo de Young (>0)" }
  - { name: nu,         type: float, required: true,  desc: "Coeficiente de Poisson ∈ (-1, 0.5)" }
  - { name: hypothesis, type: str,   required: false, default: "plane_stress",
      desc: "plane_stress | plane_strain" }
  - { name: density,    type: float, required: false, default: null,
      desc: "Densidad (kg/m³). Opcional al construir; obligatoria solo si se ensambla peso propio o masa (ADR 0008)" }

signature:
  compute_state: "(ε: ndarray(3,), state_vars=None) -> (σ: ndarray(3,), C_e: ndarray(3,3), state_vars)"
  strain_kind: "plano Voigt — [ε_xx, ε_yy, γ_xy] con γ_xy = 2·ε_xy"

state_variables: []

conventions:
  sign: "σ > 0 ⇔ tracción (Reglas §5)"
  voigt: "γ_xy = 2·ε_xy en deformación; σ_xy tensorial sin factor"
  hypothesis_dispatch: "selección por kwarg al construir; matriz C precalculada"

validity:
  - "E > 0"
  - "ν ∈ (-1, 0.5) estricto; ν = 0.5 rechazado (incompresible)"
  - "hypothesis ∈ {'plane_stress', 'plane_strain'}"
  - "pequeñas deformaciones (|ε| ≲ 1e-2)"

out_of_scope:
  - "anisotropía (ortótropo, monoclínico)"
  - "plasticidad ⇒ VonMises2D, DruckerPrager2D"
  - "daño ⇒ IsotropicDamage2D"
  - "incompresibilidad estricta (ν=0.5) ⇒ requiere formulación mixta u-p"
  - "grandes deformaciones"

acceptance:
  # Cubierto por tests existentes de elementos sólidos 2D que usan
  # Elastic2D como material por default.
  verification:
    - name: ley_de_hooke_3x3_exacta
      setup: "evaluar compute_state con ε arbitrario; comparar σ = C_e·ε en ambas hipótesis"
      expect: "σ = C_e·ε exacto componente a componente; tangente = C_e"
      tol_rel: 1.0e-14
    - name: simetria_C_e
      setup: "construir C_e en ambas hipótesis"
      expect: "C_e == C_e.T exacto"
      tol_abs: 1.0e-14
    - name: positividad_definida
      setup: "autovalores de C_e en ambas hipótesis con ν ∈ {0, 0.3, 0.49}"
      expect: "todos los autovalores estrictamente positivos"
      tol_abs: 0.0
    - name: plane_stress_traccion_uniaxial
      setup: "ε = (ε_0, -ν·ε_0, 0); plane_stress"
      expect: "σ_xx = E·ε_0, σ_yy = 0, σ_xy = 0 (uniaxial puro emergente)"
      tol_rel: 1.0e-12
    - name: plane_strain_confinamiento
      setup: "ε = (ε_0, 0, 0); plane_strain"
      expect: "σ_yy = ν/(1-ν) · σ_xx (reacción de confinamiento clásica)"
      tol_rel: 1.0e-12
    - name: rechazo_inputs_invalidos
      setup: "construir Elastic2D con E≤0, ν fuera de rango, hypothesis desconocida, density<0"
      expect: "ValueError con mensaje claro en cada caso"

references:
  - "Timoshenko S., Goodier J.N. (1970). Theory of Elasticity. McGraw-Hill. §3-4 (plane stress / plane strain)."
  - "Cook R.D., Malkus D.S., Plesha M.E., Witt R.J. (2002). Concepts and Applications of Finite Element Analysis. Wiley. §3.6."
  - "Bathe K.J. (2014). Finite Element Procedures. Prentice Hall. §4.2."

\end{lstlisting}


\section{Implementación}

\begin{itemize}
  \item \textbf{Archivo}: solidum/materials/elastic\_2d.py.
  \item \textbf{Clase}: \texttt{Elastic2D}, registrada vía \texttt{@MaterialRegistry.register}.
  \item \textbf{Despacho por hipótesis}: la matriz \texttt{C} se precalcula en el constructor según \texttt{hypothesis}. \texttt{compute\_state} solo hace \texttt{self.C @ strain} sin ramificaciones --- coste mínimo.
  \item \textbf{Validación temprana} en constructor: \texttt{E > 0}, \texttt{\ensuremath{\nu} \ensuremath{\in} (-1, 0.5)}, \texttt{hypothesis \ensuremath{\in} \{plane\_stress, plane\_strain\}}, \texttt{density \ensuremath{\ge} 0} si se pasa. Mensajes en español.
  \item \textbf{Tests}: cobertura implícita amplísima vía todos los pipelines de sólidos 2D (Quad4, Quad8, Quad9, Tri3, Tri6) que usan Elastic2D como material por default (\texttt{tests/test\_solid\_2d\_*.py}, decenas de tests).
\end{itemize}


\section{Diálogo}

\begin{itemize}
  \item \textbf{2026-05-19} \ensuremath{\cdot} Spec creada retroactivamente para cerrar el hueco H-5.3 de la auditoría 2026-05-18. Material anterior a la convención de specs validadas; no se modifica código. Aceptación cumplida por cobertura indirecta en pipelines existentes.
\end{itemize}


\newpage

\section{IsotropicDamage2D}



\textit{Orden de trabajo. El usuario escribe \textbf{especificación física}, \textbf{formulación} y \textbf{contrato}; la IA rellena \textbf{implementación} y responde en \textbf{diálogo}.}

\textit{> Spec \textbf{retroactiva con ampliación}: el modelo está implementado desde sesiones anteriores con \textbf{tangente secante}. Esta spec lo documenta y añade la \textbf{tangente algorítmica consistente} para recuperar convergencia cuadrática del Newton global.}



\section{Especificación física}

\subsection{0. Descripción general}

Modelo de \textbf{daño isótropo escalar} en pequeñas deformaciones para sólidos 2D. Captura degradación progresiva e isótropa de la rigidez en respuesta a la deformación máxima histórica, con ablandamiento exponencial. Pensado para hormigón en tracción uniforme, cerámicos, materiales cuasi-frágiles donde la microfisuración distribuida puede modelarse mediante un escalar de daño en lugar de microestructura explícita.

\subsection{1. Cinemática}

Pequeñas deformaciones, descomposición elástica pura: la deformación plástica no existe en este modelo (sin plasticidad). El daño se manifiesta como degradación de la rigidez efectiva.

\subsection{2. Esfuerzo nominal vs esfuerzo efectivo}

\[
\boldsymbol\sigma = (1 - d)\,\boldsymbol\sigma^\text{eff}, \qquad \boldsymbol\sigma^\text{eff} = \mathbf C_e\,\boldsymbol\varepsilon
\]

con $d \in [0, d_\text{max}]$ la variable de daño isótropa y $\mathbf C_e$ el tensor constitutivo elástico isótropo (2D, plane stress o plane strain según \texttt{hypothesis} heredada del \texttt{Elastic2D} interno).

El esfuerzo efectivo $\boldsymbol\sigma^\text{eff}$ es el que el material intacto soportaría con la deformación corriente. El esfuerzo nominal $\boldsymbol\sigma$ es lo que efectivamente transmite considerando la degradación.

\subsection{3. Deformación equivalente}

Cantidad escalar que cuantifica la "intensidad" de la deformación. Convención del proyecto, norma simétrica simple en notación Voigt $[\varepsilon_{xx}, \varepsilon_{yy}, \gamma_{xy}]$ con $\gamma_{xy} = 2\varepsilon_{xy}$:

\[
\varepsilon_\text{eq}(\boldsymbol\varepsilon) = \sqrt{\varepsilon_{xx}^2 + \varepsilon_{yy}^2 + \tfrac{1}{2}\gamma_{xy}^2} = \sqrt{\boldsymbol\varepsilon^\top \mathbf M\,\boldsymbol\varepsilon}
\]

con $\mathbf M = \operatorname{diag}(1, 1, 1/2)$. Equivalente a la norma de Frobenius del tensor deformación 2D plana ($\boldsymbol\varepsilon : \boldsymbol\varepsilon$ con $\varepsilon_{xy} = \gamma_{xy}/2$).

\textbf{Limitación}: no distingue tracción de compresión. Para hormigón a compresión se necesitaría un split tipo Mazars o de Vree; aquí se diferencia explícitamente como \textit{out-of-scope}.

\subsection{4. Variable histórica y condición de carga (Kuhn-Tucker)}

\[
\kappa_{n+1} = \max(\kappa_n, \varepsilon_\text{eq}(\boldsymbol\varepsilon_{n+1})), \qquad \kappa_0 \text{ inicial}
\]

Esto codifica:
\begin{itemize}
  \item \textbf{Carga activa} ($\dot\kappa > 0$): $\kappa$ crece con $\varepsilon_\text{eq}$; el daño puede aumentar.
  \item \textbf{Descarga / recarga} ($\dot\kappa = 0$): $\kappa$ queda congelado; el daño no cambia.
  \item \textbf{Irreversibilidad}: $\kappa$ es monótono no decreciente.
\end{itemize}

\subsection{5. Ley de evolución del daño}

\[
d(\kappa) = \begin{cases}
0 & \kappa \leq \kappa_0 \\
1 - \dfrac{\kappa_0}{\kappa}\,e^{-\alpha(\kappa - \kappa_0)} & \kappa > \kappa_0
\end{cases}
\]

con saturación numérica $d \leq d_\text{max} = $ \texttt{DAMAGE\_MAX} (\texttt{solidum.constants.DAMAGE\_MAX}, evita rigidez nula).

Propiedades:
\begin{itemize}
  \item $d(\kappa_0) = 0$ (continuidad).
  \item $d'(\kappa) > 0$ para $\kappa > \kappa_0$ (irreversibilidad).
  \item $\lim_{\kappa \to \infty} d = 1$ (asintótico).
  \item $\alpha$ controla la velocidad de degradación (ablandamiento más abrupto si $\alpha$ grande).
\end{itemize}


\section{Formulación numérica}

\subsection{6. Algoritmo en un paso}

Dado $\{\kappa_n\}$ y $\boldsymbol\varepsilon_{n+1}$:

\begin{enumerate}
  \item $\varepsilon_\text{eq} = \sqrt{\varepsilon_{xx}^2 + \varepsilon_{yy}^2 + \tfrac{1}{2}\gamma_{xy}^2}$.
  \item Si $\varepsilon_\text{eq} > \kappa_n$: \textbf{carga activa}, $\kappa_{n+1} = \varepsilon_\text{eq}$. Si no: \textbf{descarga}, $\kappa_{n+1} = \kappa_n$.
  \item Aplicar ley de daño con $\kappa_{n+1}$; saturar a $d_\text{max}$.
  \item $\boldsymbol\sigma_{n+1} = (1 - d)\,\mathbf C_e\,\boldsymbol\varepsilon_{n+1}$.
  \item Tangente: \textbf{secante} en descarga / sin daño / saturación; \textbf{algorítmica consistente} en carga activa con daño no saturado (ver §7).
\end{enumerate}

El algoritmo es \textbf{explícito}: no requiere Newton local. La actualización de $\kappa$ y $d$ es directa.

\subsection{7. Tangente algorítmica consistente (ampliación)}

\[
\mathbf C^\text{alg} = \frac{\partial \boldsymbol\sigma_{n+1}}{\partial \boldsymbol\varepsilon_{n+1}}
\]

Derivando $\boldsymbol\sigma = (1-d)\,\mathbf C_e\,\boldsymbol\varepsilon$:

\[
\mathbf C^\text{alg} = (1-d)\,\mathbf C_e - \boldsymbol\sigma^\text{eff} \otimes \frac{\partial d}{\partial \boldsymbol\varepsilon}
\]

con $\boldsymbol\sigma^\text{eff} = \mathbf C_e\,\boldsymbol\varepsilon$. La derivada del daño respecto a la deformación, por regla de la cadena:

\[
\frac{\partial d}{\partial \boldsymbol\varepsilon} = \frac{\partial d}{\partial \kappa}\,\frac{\partial \kappa}{\partial \boldsymbol\varepsilon}
\]

donde:

\begin{itemize}
  \item $\dfrac{\partial d}{\partial \kappa} = (1 - d)\left(\dfrac{1}{\kappa} + \alpha\right)$ para la ley exponencial. (Demostración: derivando $d = 1 - (\kappa_0/\kappa)e^{-\alpha(\kappa-\kappa_0)}$ y reagrupando vía $1 - d = (\kappa_0/\kappa)e^{-\alpha(\kappa-\kappa_0)}$.)
\end{itemize}

\begin{itemize}
  \item $\dfrac{\partial \kappa}{\partial \boldsymbol\varepsilon}$ depende del régimen:
  \item \textbf{Carga activa} ($\varepsilon_\text{eq} > \kappa_n$, $\kappa = \varepsilon_\text{eq}$):
\end{itemize}
    \[
\dfrac{\partial \kappa}{\partial \boldsymbol\varepsilon} = \dfrac{\partial \varepsilon_\text{eq}}{\partial \boldsymbol\varepsilon} = \dfrac{1}{\varepsilon_\text{eq}}\begin{bmatrix}\varepsilon_{xx} \\ \varepsilon_{yy} \\ \tfrac{1}{2}\gamma_{xy}\end{bmatrix}
\]
\begin{itemize}
  \item \textbf{Descarga} ($\varepsilon_\text{eq} \leq \kappa_n$): $\partial \kappa / \partial \boldsymbol\varepsilon = 0$.
\end{itemize}

\textbf{Tangente final}:

\[
\mathbf C^\text{alg} = \begin{cases}
\mathbf C_e & \kappa \leq \kappa_0 \;(\text{sin daño}) \\
(1-d)\,\mathbf C_e & \text{descarga} \\
(1-d)\,\mathbf C_e \;-\; \dfrac{(1-d)(1/\kappa + \alpha)}{\varepsilon_\text{eq}}\,(\mathbf C_e\boldsymbol\varepsilon)(\mathbf M\boldsymbol\varepsilon)^\top & \text{carga activa, } d < d_\text{max} \\
(1-d_\text{max})\,\mathbf C_e & d = d_\text{max} \;(\text{saturación})
\end{cases}
\]

con $\mathbf M = \operatorname{diag}(1, 1, 1/2)$.

\subsection{8. Pérdida de simetría}

El producto externo $(\mathbf C_e\boldsymbol\varepsilon)(\mathbf M\boldsymbol\varepsilon)^\top$ es una matriz $3\times 3$ que \textbf{no es simétrica en general}. $\mathbf C_e\boldsymbol\varepsilon$ y $\mathbf M\boldsymbol\varepsilon$ no son proporcionales salvo en estados de deformación muy particulares (e.g. uniaxial puro alineado con ejes).

Esto rompe la simetría heredada de la tangente secante $(1-d)\mathbf C_e$:

\begin{itemize}
  \item \textbf{Tangente secante}: simétrica. \texttt{IS\_SYMMETRIC = True} sería válido.
  \item \textbf{Tangente consistente}: no simétrica en general. \textbf{\texttt{IS\_SYMMETRIC = False} obligatorio.}
\end{itemize}

Consecuencia algebraica (ADR 0003): el despachador elige LU en lugar de Cholesky/LDL\ensuremath{^{\top}} para sistemas globales que incluyan este material. El coste por paso aumenta \textasciitilde{}30-50\% respecto a Cholesky, pero la convergencia cuadrática del Newton global con tangente consistente recupera ese coste con creces (típicamente 2-3 iteraciones en lugar de 5-10 con tangente secante).

\subsection{9. Decisión de régimen durante el Newton}

Durante una iteración del Newton global, el material recibe un $\boldsymbol\varepsilon$ trial y debe decidir carga/descarga comparando con $\kappa_n$ (el estado committed del paso anterior). La transición entre ramas (loading \ensuremath{\leftrightarrow} unloading) es \textbf{discontinua en la tangente}: al cruzar $\varepsilon_\text{eq} = \kappa_n$ el segundo término aparece/desaparece bruscamente. En la práctica esto no genera problemas porque el Newton converge antes de oscilar entre ramas; en problemas patológicos puede requerir continuation o arc-length.


\section{Contrato de implementación}

\begin{lstlisting}[language=yaml]
name: IsotropicDamage2D
kind: material
status: validated

interface:
  strain_dim: 3
  primary_state_var: damage
  is_symmetric: false      # tangente consistente no simétrica (ver §8)

parameters:
  - { name: E,          type: float, required: true,  desc: "Módulo de Young intacto" }
  - { name: nu,         type: float, required: true,  desc: "Coeficiente de Poisson" }
  - { name: kappa_0,    type: float, required: true,  desc: "Umbral de deformación equivalente; daño nulo mientras κ ≤ κ_0" }
  - { name: alpha,      type: float, required: true,  desc: "Velocidad de degradación exponencial (>0)" }
  - { name: hypothesis, type: str,   required: false, default: "plane_stress", desc: "plane_stress | plane_strain (delegado a Elastic2D interno)" }
  - { name: density,    type: float, required: false, default: null, desc: "Densidad (ADR 0008)" }

signature:
  compute_state: "(ε: ndarray(3,), state_vars=None) -> (σ: ndarray(3,), C_tan: ndarray(3,3), state_vars')"
  strain_kind: "plano Voigt — [ε_xx, ε_yy, γ_xy] con γ_xy = 2·ε_xy"

state_schema:
  kappa:  "float ≥ κ_0 — máxima deformación equivalente histórica"
  damage: "float ∈ [0, DAMAGE_MAX] — variable de daño escalar isótropa"

conventions:
  sign: "σ > 0 ⇔ tracción; el modelo no distingue tracción de compresión en la activación del daño"
  voigt: "γ_xy = 2·ε_xy en deformación; ε_eq usa norma M=diag(1,1,1/2)"
  damage_irreversibility: "κ monótono no decreciente vía max(κ_old, ε_eq); d(κ) creciente"

validity:
  - "E > 0, κ_0 > 0, α > 0"
  - "ν ∈ (-1, 0.5) (delegación a Elastic2D)"
  - "hypothesis ∈ {'plane_stress', 'plane_strain'}"

out_of_scope:
  - "split tracción/compresión (Mazars, de Vree, Mises modificado); el modelo daña por igual en cualquier régimen"
  - "anisotropía del daño (tensor de daño D vs escalar)"
  - "regularización para evitar mesh-dependency en regime de ablandamiento (longitud característica, no-local, gradient)"
  - "acoplamiento con plasticidad"
  - "fatiga / acumulación cíclica del daño"
  - "tangente consistente para IsotropicDamage1D — fuera del alcance de esta spec; misma deuda gemela"

numerical_caveats:
  - "Tangente consistente NO simétrica → IS_SYMMETRIC=False obliga al despachador algebraico a usar LU (ADR 0003). Costo por paso ~30-50% mayor que Cholesky pero convergencia cuadrática del Newton global lo compensa."
  - "Transición loading↔unloading es discontinua en la tangente. Newton converge antes de oscilar entre ramas en casos típicos. Para problemas patológicos cerca del umbral, considerar arc-length."
  - "Saturación d=DAMAGE_MAX corta la corrección consistente (∂d/∂κ deja de ser efectivamente (1-d)(1/κ+α)); el algoritmo devuelve tangente secante en ese caso para evitar términos espurios."
  - "Sin regularización, la solución es mesh-dependent en régimen de ablandamiento (localización en una banda de elementos). Para benchmark cuantitativo se requiere mismo refinamiento de malla; para comparaciones cualitativas (presencia de banda, dirección) basta con regularización de gradiente o longitud característica (fuera de scope)."

acceptance:
  # Cubierto por tests/test_materials_unit.py · TestIsotropicDamage2D (preexistentes)
  # y tests/test_damage2d_consistent_tangent.py (nuevos)
  unit_existing_no_regression:
    - name: secant_response_below_threshold
      expect: "κ ≤ κ_0 ⇒ d=0, σ = C_e·ε, tangente = C_e (sin cambio respecto a versión secante anterior)"
      tol_rel: 1.0e-12

    - name: damage_evolution_exponential
      expect: "para κ > κ_0, d cumple ley exponencial dentro de la fórmula"
      tol_rel: 1.0e-10

  unit_consistent_tangent:
    - name: tangent_equals_secant_on_unloading
      setup: "carga hasta κ>κ_0, descarga con ε de menor norma"
      expect: "C_tan = (1-d)·C_e exacto (segundo término se anula con ∂κ/∂ε = 0)"
      tol_rel: 1.0e-12

    - name: tangent_equals_secant_below_threshold
      setup: "ε pequeño, κ_old = κ_0"
      expect: "C_tan = C_e (sin daño activo, ni secante con d=0 ni corrección)"
      tol_rel: 1.0e-12

    - name: tangent_equals_secant_at_saturation
      setup: "carga hasta κ enorme tal que d alcanza DAMAGE_MAX"
      expect: "C_tan = (1-DAMAGE_MAX)·C_e (el término consistente se corta porque ∂d/∂κ efectivamente 0 al saturar)"
      tol_rel: 1.0e-12

    - name: tangent_consistent_term_on_loading
      setup: "carga activa con κ > κ_0 y d < DAMAGE_MAX"
      expect: "C_tan ≠ (1-d)·C_e — incluye el término rank-1 -(1-d)(1/κ+α)/ε_eq · (Ce·ε) ⊗ (M·ε)"
      verification: "compara contra fórmula explícita evaluada en el mismo estado"
      tol_rel: 1.0e-9

    - name: tangent_not_symmetric_in_loading
      setup: "carga activa con estado de deformación no degenerado (ε_xx ≠ ε_yy, γ_xy ≠ 0)"
      expect: "C_tan − C_tan^T ≠ 0 (al menos un elemento off-diagonal con diferencia significativa)"
      verification: "‖C_tan - C_tan^T‖_F / ‖C_tan‖_F ≥ 0.01 en el caso de prueba"

    - name: tangent_finite_difference_consistency
      setup: "carga activa; calcular tangente por diferencia finita centrada (perturbación 1e-7) y comparar con cerrada"
      expect: "‖C_tan_FD − C_tan_closed_form‖ / ‖C_tan‖ < 1e-5"
      tol_rel: 1.0e-5

  unit_arch:
    - name: is_symmetric_attribute_is_false
      expect: "IsotropicDamage2D.IS_SYMMETRIC == False (atributo declarado en clase)"

  integration:
    - name: quad4_damage_newton_converges_in_few_iter
      setup: "Quad4 unitario, IsotropicDamage2D plane stress, carga uniaxial hasta régimen plenamente dañado"
      expect: "Newton converge en ≤6 iteraciones por paso (evidencia de convergencia cuadrática con tangente consistente)"

references:
  - "Lemaitre J., Chaboche J.-L. (1990). Mechanics of Solid Materials. Cambridge UP. Cap. 7 (continuum damage mechanics)."
  - "Oliver J. (1989). A consistent characteristic length for smeared cracking models. IJNME 28, 461-474."
  - "Simó J.C., Ju J.W. (1987). Strain- and stress-based continuum damage models — I. Formulation. Int. J. Solids Struct. 23, 821-840 (tangente consistente para daño isótropo)."
  - "de Souza Neto E.A., Perić D., Owen D.R.J. (2008). Computational Methods for Plasticity. Wiley. §12 (damage models)."

\end{lstlisting}


\section{Implementación}

\begin{itemize}
  \item \textbf{Archivo}: solidum/materials/damage\_2d.py.
  \item \textbf{Clase}: \texttt{IsotropicDamage2D}, registrada vía \texttt{@MaterialRegistry.register}. Declara \texttt{IS\_SYMMETRIC = False} como ClassVar (sobreescribe el default \texttt{True} de \texttt{Material}).
  \item \textbf{Estado}: dict \texttt{\{'kappa': float, 'damage': float\}}. \texttt{PRIMARY\_STATE\_VAR = 'damage'}.
  \item \textbf{Algoritmo \texttt{compute\_state}}:
\end{itemize}
\begin{enumerate}
  \item Recuperar $\kappa_n$ y calcular $\varepsilon_\text{eq}(\boldsymbol\varepsilon)$.
  \item Si $\varepsilon_\text{eq} > \kappa_n$ y $\kappa_{n+1} > \kappa_0$: rama de carga activa.
  \item Calcular $d$ por la ley exponencial, saturar a \texttt{DAMAGE\_MAX}.
  \item Calcular $\boldsymbol\sigma = (1-d)\mathbf C_e\boldsymbol\varepsilon$.
  \item Tangente: rama según §7 (secante o consistente). En la rama consistente, calcular el término rank-1 una sola vez y restar a $(1-d)\mathbf C_e$.
\end{enumerate}
\begin{itemize}
  \item \textbf{\texttt{admissibility\_scale}}: sin cambios respecto a versión previa, \texttt{E\ensuremath{\cdot}\ensuremath{\kappa}\_0} (umbral inicial --- la escala no evoluciona con el estado porque el criterio se mide contra $\kappa_0$, no contra $\kappa$).
  \item \textbf{Compatibilidad}: el cambio rompe \texttt{IS\_SYMMETRIC} para este material. El despachador algebraico (ADR 0003) verifica \texttt{domain\_is\_symmetric} agregando \texttt{material.IS\_SYMMETRIC} sobre todos los materiales del dominio; basta un material con \texttt{IS\_SYMMETRIC=False} para que el solver elija LU.
  \item \textbf{Tests}:
  \item tests/test\_materials\_unit.py \ensuremath{\cdot} \texttt{TestIsotropicDamage2D} (existente, no regresión) + nueva clase \texttt{TestIsotropicDamage2DConsistentTangent}.
  \item tests/test\_solid\_2d\_damage.py nuevo: integración Quad4 + IsotropicDamage2D + NonlinearSolver con tangente consistente.
\end{itemize}


\section{Diálogo}

\begin{itemize}
  \item \textbf{2026-05-14} \ensuremath{\cdot} Spec creada retroactivamente sobre material existente (tangente secante) y extendida con tangente algorítmica consistente. Sin ADR: extiende un material existente añadiendo un término a la tangente; no rompe contratos arquitecturales transversales. Sí cambia \texttt{IS\_SYMMETRIC} para este material, lo cual el despachador algebraico ya soporta vía agregación sobre el dominio (no requiere refactor).
  \item \textbf{2026-05-14} \ensuremath{\cdot} Decisión: la transición loading\ensuremath{\leftrightarrow}unloading se decide con $\kappa_n$ committed (no con $\kappa$ trial dentro del Newton). Esto es lo estándar y evita oscilaciones del Newton entre ramas en pasos con ambigüedad. Si el paso converge, $\kappa$ se compromete y la decisión es consistente con la trayectoria.
  \item \textbf{2026-05-14} \ensuremath{\cdot} \texttt{IsotropicDamage1D} queda con tangente secante. Misma deuda gemela conceptual, pero fuera del alcance pactado (A2 era solo 2D). Si el usuario la prioriza luego, replicar es trivial: misma fórmula reducida a escalar.
\end{itemize}


\newpage

\section{VonMises2D}



\textit{Orden de trabajo. El usuario escribe \textbf{especificación física}, \textbf{formulación numérica} y \textbf{contrato}; la IA rellena \textbf{implementación} y responde en \textbf{diálogo}.}

\textit{> Spec \textbf{retroactiva con ampliación}: la hipótesis \texttt{plane\_strain} está implementada y testeada desde sesiones anteriores; esta spec la documenta y añade \texttt{plane\_stress} como hipótesis a implementar en esta sesión.}



\section{Especificación física}

\subsection{0. Descripción general}

Modelo elastoplástico independiente de la velocidad ("rate-independent") basado en el criterio \textbf{J2 / Von Mises} con regla de flujo \textbf{asociada} y \textbf{endurecimiento isótropo lineal}. Captura plasticidad de metales dúctiles en régimen de pequeñas deformaciones, sin efectos cinemáticos (Bauschinger), térmicos ni viscosos.

Soporta dos hipótesis cinemáticas 2D mutuamente excluyentes seleccionadas en construcción:

\begin{itemize}
  \item \textbf{\texttt{plane\_strain}} --- $\varepsilon_{zz} = \varepsilon_{xz} = \varepsilon_{yz} = 0$. Cuerpos prismáticos largos cargados transversalmente (presa, túnel, cilindro largo).
  \item \textbf{\texttt{plane\_stress}} --- $\sigma_{zz} = \sigma_{xz} = \sigma_{yz} = 0$. Láminas o membranas delgadas en su plano (chapa traccionada, placa fina).
\end{itemize}

\subsection{1. Descomposición aditiva}

En pequeñas deformaciones, el tensor de deformación se descompone aditivamente:

\[
\boldsymbol\varepsilon = \boldsymbol\varepsilon^e + \boldsymbol\varepsilon^p
\]

La parte elástica $\boldsymbol\varepsilon^e$ gobierna el esfuerzo vía la ley elástica isótropa; la parte plástica $\boldsymbol\varepsilon^p$ es \textbf{histórica} y evoluciona según la regla de flujo.

\subsection{2. Ley elástica}

\[
\boldsymbol\sigma = \mathbf C_e : \boldsymbol\varepsilon^e = \mathbf C_e : (\boldsymbol\varepsilon - \boldsymbol\varepsilon^p)
\]

con $\mathbf C_e$ tensor de rigidez elástica isótropa parametrizado por $(E, \nu)$ o equivalentemente $(K, G)$:

\[
K = \frac{E}{3(1-2\nu)}, \qquad G = \frac{E}{2(1+\nu)}
\]

\subsection{3. Superficie de fluencia (J2)}

\[
f(\boldsymbol\sigma, \alpha) = \|\mathbf s\| - \sqrt{\tfrac{2}{3}}\,\big(\sigma_y + H\,\alpha\big) \;\leq\; 0
\]

donde $\mathbf s = \boldsymbol\sigma - \tfrac{1}{3}\operatorname{tr}(\boldsymbol\sigma)\mathbf I$ es el tensor desviador, $\|\mathbf s\| = \sqrt{\mathbf s:\mathbf s}$, $\sigma_y > 0$ es la fluencia inicial y $H \geq 0$ el módulo de endurecimiento isótropo lineal. El factor $\sqrt{2/3}$ alinea $\sigma_y$ con la fluencia uniaxial estándar.

\subsection{4. Regla de flujo (asociada)}

\[
\dot{\boldsymbol\varepsilon}^p = \dot\gamma \,\frac{\partial f}{\partial \boldsymbol\sigma} = \dot\gamma \,\mathbf N, \qquad \mathbf N = \frac{\mathbf s}{\|\mathbf s\|}
\]

El flujo es puramente desviador ($\operatorname{tr}(\mathbf N) = 0$) --- \textbf{plasticidad incompresible}, propiedad central de J2.

\subsection{5. Endurecimiento isótropo lineal}

\[
\dot\alpha = \sqrt{\tfrac{2}{3}}\,\dot\gamma
\]

$\alpha$ es la \textbf{deformación plástica acumulada equivalente} (escalar, monótona no decreciente). La superficie de fluencia se expande uniformemente con $\alpha$ manteniendo su centro en el origen del espacio de tensiones desviadoras.

\subsection{6. Condiciones de Kuhn-Tucker}

\[
\dot\gamma \geq 0, \qquad f \leq 0, \qquad \dot\gamma\,f = 0
\]

Garantizan que $\dot\gamma > 0$ solo si el estado intenta salir de la superficie, y que el estado real siempre cumple $f \leq 0$.

\subsection{7. Variables internas}

\begin{center}
\begin{tabularx}{\textwidth}{YYY}
\hline
\textbf{Símbolo} & \textbf{Tipo} & \textbf{Significado} \\
\hline
$\boldsymbol\varepsilon^p$ & tensor (4 componentes Voigt extendido en plane strain, ver §8) & deformación plástica \\
$\alpha$ & escalar & deformación plástica acumulada equivalente \\
\hline
\end{tabularx}
\end{center}
\texttt{alpha} es la variable principal exportable (\texttt{PRIMARY\_STATE\_VAR = 'alpha'}); \texttt{eps\_p} es interna al algoritmo de return mapping.

\subsection{8. Notación Voigt y manejo de $\varepsilon^p_{zz}$}

Convención del proyecto para 2D: $\boldsymbol\varepsilon = [\varepsilon_{xx},\,\varepsilon_{yy},\,\gamma_{xy}]$ con $\gamma_{xy} = 2\varepsilon_{xy}$. Pero la \textbf{descomposición desviadora} requiere todas las componentes 3D, así que internamente el material maneja $\boldsymbol\varepsilon^p = [\varepsilon^p_{xx},\,\varepsilon^p_{yy},\,\varepsilon^p_{zz},\,\varepsilon^p_{xy}]$ (4 componentes, $\varepsilon^p_{xy}$ tensorial sin factor 2).

\begin{itemize}
  \item \textbf{Plane strain}: $\varepsilon_{zz} = 0$ se impone \textit{en deformación total}. La componente $\varepsilon^p_{zz}$ es no nula y evoluciona libremente (la incompresibilidad plástica acopla $\varepsilon^p_{zz} = -(\varepsilon^p_{xx} + \varepsilon^p_{yy})$). Por simetría, $\sigma_{zz}$ resulta no nulo pero no se expone en el API público 2D --- vive solo en el cómputo interno.
  \item \textbf{Plane stress}: $\sigma_{zz} = 0$ se impone \textit{en esfuerzo}. La componente $\varepsilon_{zz}$ es no nula (extensión transversal de Poisson + Poisson plástico), y aparece como variable adicional en el problema local de retorno (ver §10).
\end{itemize}


\section{Formulación numérica}

\subsection{9. Esquema temporal --- Backward Euler implícito}

Dado el estado convergido $\{\boldsymbol\varepsilon^p_n, \alpha_n\}$ al paso $n$ y la deformación total prescrita $\boldsymbol\varepsilon_{n+1}$, resolver para $\{\boldsymbol\sigma_{n+1}, \boldsymbol\varepsilon^p_{n+1}, \alpha_{n+1}\}$ con discretización implícita:

\[
\boldsymbol\varepsilon^p_{n+1} = \boldsymbol\varepsilon^p_n + \Delta\gamma\,\mathbf N_{n+1}, \qquad \alpha_{n+1} = \alpha_n + \sqrt{\tfrac{2}{3}}\,\Delta\gamma
\]

más $f(\boldsymbol\sigma_{n+1}, \alpha_{n+1}) \leq 0$ con la condición de complementariedad.

\subsection{10. Return mapping plane strain --- \textbf{algoritmo radial cerrado} (existente)}

Descomposición volumétrica-desviadora del estado de prueba:

\[
\boldsymbol\varepsilon^{\text{trial}}_{n+1} = \boldsymbol\varepsilon_{n+1} - \boldsymbol\varepsilon^p_n, \qquad \mathbf s^{\text{trial}} = 2G\,\mathbf e^{\text{trial}}, \qquad p^{\text{trial}} = K\,\operatorname{tr}(\boldsymbol\varepsilon_{n+1})
\]

donde $\mathbf e^{\text{trial}}$ es la parte desviadora de $\boldsymbol\varepsilon^{\text{trial}}_{n+1}$. Función de fluencia trial:

\[
f^{\text{trial}} = \|\mathbf s^{\text{trial}}\| - \sqrt{\tfrac{2}{3}}\,(\sigma_y + H\,\alpha_n)
\]

Si $f^{\text{trial}} \leq \text{tol}$ \ensuremath{\to} paso elástico, $\boldsymbol\sigma_{n+1} = \mathbf C_e:\boldsymbol\varepsilon_{n+1}$, estado intacto.

Si no, el flujo es radial (la dirección no cambia entre trial y final por isotropía + asociado en J2):

\[
\Delta\gamma = \frac{f^{\text{trial}}}{2G + \tfrac{2}{3}H} \qquad\text{(forma cerrada)}
\]

\[
\mathbf N = \frac{\mathbf s^{\text{trial}}}{\|\mathbf s^{\text{trial}}\|}, \quad \mathbf s_{n+1} = \mathbf s^{\text{trial}} - 2G\,\Delta\gamma\,\mathbf N, \quad \boldsymbol\sigma_{n+1} = \mathbf s_{n+1} + p^{\text{trial}}\,\mathbf I
\]

\textbf{Tangente consistente algorítmica} (derivada del algoritmo, no de la ley continua):

\[
\mathbf C^{\text{alg}} = K\,\mathbf 1\otimes\mathbf 1 + 2G(1-\beta)\,\mathbf I_{\text{dev}} - 2G\,\bar\gamma\,\mathbf N\otimes\mathbf N
\]

con $\beta = \dfrac{2G\,\Delta\gamma}{\|\mathbf s^{\text{trial}}\|}$, $\bar\gamma = \dfrac{1}{1 + H/(3G)} - \beta$.

\subsection{11. Return mapping plane stress --- \textbf{algoritmo proyectado} (a implementar)}

Referencia: Simó \& Hughes 1998, \textit{Computational Inelasticity}, §3.4.1 ("Plane stress projected algorithm").

La restricción $\sigma_{zz} = 0$ se impone \textbf{eliminando} $\varepsilon_{zz}$ del problema mediante un operador de proyección, no como ecuación extra. Resultado: el problema local sigue siendo escalar en $\Delta\gamma$, sin nested Newton.

\textbf{Matriz elástica plane stress} (3\ensuremath{\times}3, base Voigt $[\varepsilon_{xx},\varepsilon_{yy},\gamma_{xy}]$):

\[
\mathbf C^{ps}_e = \frac{E}{1-\nu^2}\begin{bmatrix} 1 & \nu & 0 \\ \nu & 1 & 0 \\ 0 & 0 & (1-\nu)/2 \end{bmatrix}
\]

\textbf{Predictor elástico} en el subespacio plane stress:

\[
\boldsymbol\sigma^{\text{trial}} = \mathbf C^{ps}_e\,(\boldsymbol\varepsilon_{n+1} - \mathbf P_\varepsilon\,\boldsymbol\varepsilon^p_n)
\]

con $\mathbf P_\varepsilon$ extrayendo las componentes plane $[\varepsilon^p_{xx},\varepsilon^p_{yy},2\varepsilon^p_{xy}]$ del tensor 4D.

\textbf{Operador $\mathbf P$ --- norma desviadora bajo $\sigma_{zz}=0$} (Simó-Hughes Box 3.1):

\[
\mathbf P = \frac{1}{3}\begin{bmatrix} 2 & -1 & 0 \\ -1 & 2 & 0 \\ 0 & 0 & 6 \end{bmatrix}
\]

Esta matriz codifica el producto escalar desviador en el subespacio plane stress; reemplaza $\|\mathbf s\|^2$ por $\boldsymbol\sigma^\top \mathbf P\,\boldsymbol\sigma$ con todas las componentes 3D ya eliminadas analíticamente.

\textbf{Función de fluencia plane stress proyectada}:

\[
\bar f(\boldsymbol\sigma, \alpha) = \tfrac{1}{2}\,\boldsymbol\sigma^\top \mathbf P\,\boldsymbol\sigma - \tfrac{1}{3}\big(\sigma_y + H\,\alpha\big)^2 \;\leq\; 0
\]

\textbf{Esquema de retorno}:

\[
\boldsymbol\sigma_{n+1} = \big[\mathbf C^{ps}_e\big]^{-1}\,\big(\boldsymbol\sigma^{\text{trial}}\big) - \Delta\gamma\,\mathbf P\,\boldsymbol\sigma_{n+1}
\]

reordenando:

\[
\boldsymbol\sigma_{n+1} = \mathbf A(\Delta\gamma)^{-1}\,\boldsymbol\sigma^{\text{trial}}, \qquad \mathbf A(\Delta\gamma) = \mathbf I + \Delta\gamma\,\mathbf C^{ps}_e\,\mathbf P
\]

La matriz $\mathbf A$ es 3\ensuremath{\times}3, diagonalizable en los autovectores ortonormales de $\mathbf C^{ps}_e\mathbf P$. Los autovalores son cerrados:

\[
\mu_1 = \frac{E}{3(1-\nu)}, \quad \mu_2 = 2G, \quad \mu_3 = 2G
\]

con autovectores $\mathbf v_1 = \tfrac{1}{\sqrt 2}[1,1,0]$ (hidrostático plano), $\mathbf v_2 = \tfrac{1}{\sqrt 2}[1,-1,0]$ (desviador plano), $\mathbf v_3 = [0,0,1]$ (cortante). Diagonalizan simultáneamente $\mathbf C^{ps}_e$ y $\mathbf P$, con $\mathbf v_i^\top\mathbf P\,\mathbf v_i \in \{1/3, 1, 2\}$.

\textbf{Actualización de la deformación plástica acumulada equivalente}. La regla $\dot\alpha = \sqrt{2/3}\,\dot\gamma$ usada en plane strain \textbf{no} es válida aquí: en plane stress projected $\dot\gamma$ tiene unidades $[1/\text{esfuerzo}]$ (porque $\dot{\boldsymbol\varepsilon}^p = \dot\gamma\,\mathbf P\boldsymbol\sigma$ con $\mathbf P\boldsymbol\sigma$ en unidades de esfuerzo). La fórmula físicamente correcta y dimensionalmente consistente, obtenida de $\dot\alpha = \sqrt{2/3}\,\|\dot{\boldsymbol\varepsilon}^p\|_\text{Frob}$:

\[
\alpha_{n+1} = \alpha_n + \Delta\gamma \,\sqrt{\tfrac{2}{3}\,\boldsymbol\sigma_{n+1}^\top\mathbf P\,\boldsymbol\sigma_{n+1}}
\]

\textbf{Newton local} sobre $\Delta\gamma$. En base autovectores $\sigma_i(\Delta\gamma) = a_i / (1+\Delta\gamma\,\mu_i)$, así $\boldsymbol\sigma^\top\mathbf P\,\boldsymbol\sigma$ se expresa cerrado:

\[
\boldsymbol\sigma^\top\mathbf P\,\boldsymbol\sigma(\Delta\gamma) = \frac{a_1^2}{3(1+\Delta\gamma\mu_1)^2} + \frac{a_2^2}{(1+\Delta\gamma\mu_2)^2} + \frac{2\,a_3^2}{(1+\Delta\gamma\mu_3)^2}
\]

donde $a_i$ son las proyecciones de $\boldsymbol\sigma^{\text{trial}}$ en los autovectores. La ecuación residual:

\[
\bar f(\Delta\gamma) = \tfrac{1}{2}\,\boldsymbol\sigma^\top\mathbf P\,\boldsymbol\sigma(\Delta\gamma) - \tfrac{1}{3}\big(\sigma_y + H\,\alpha(\Delta\gamma)\big)^2 = 0
\]

con $\alpha(\Delta\gamma) = \alpha_n + \Delta\gamma\sqrt{(2/3)\,\boldsymbol\sigma^\top\mathbf P\,\boldsymbol\sigma}$ es monótona decreciente; tangente cerrada con regla de la cadena. \textbf{Newton local converge} típicamente en 3-6 iteraciones desde $\Delta\gamma = 0$. Coste: O(1) por punto de Gauss, comparable a plane strain.

\textbf{Tangente consistente plane stress}. Sea $\mathbf D = \mathbf A^{-1}\mathbf C^{ps}_e$, $\mathbf q = \boldsymbol\sigma_{n+1}^\top\mathbf P\,\mathbf D\,\mathbf P\,\boldsymbol\sigma_{n+1}$, $w = \sqrt{(2/3)\boldsymbol\sigma_{n+1}^\top\mathbf P\,\boldsymbol\sigma_{n+1}}$, $m = (2/3)\,R\,H$, $n = 2\Delta\gamma/(3w)$. De la condición de consistencia ($\dot{\bar f}=0$) con la regla correcta de $\alpha$:

\[
\mathbf C^{\text{alg}}_{ps} = \mathbf D - \frac{(\mathbf D\,\mathbf P\boldsymbol\sigma_{n+1})\,(\boldsymbol\sigma_{n+1}^\top\mathbf P\,\mathbf D)}{q + m\,w / (1 - m\,n)}
\]

Si $H=0$ (plasticidad perfecta) el término $m$ se anula y el denominador se reduce a $q$. Forma cerrada --- sin diferenciación numérica.

\textbf{Predictor elástico con $\boldsymbol\varepsilon^p \neq 0$}. Tanto en plane strain como en plane stress, el predictor debe construirse como $\boldsymbol\sigma^{\text{trial}} = \mathbf C_e\,(\boldsymbol\varepsilon - \boldsymbol\varepsilon^p_n)$ (o equivalente desviador + presión), \textbf{no} como $\mathbf C_e\,\boldsymbol\varepsilon$, que ignoraría la deformación plástica acumulada. Importante en descargas elásticas desde estado plástico y al reevaluar el material con \texttt{compute\_gauss\_state(U\_final)} tras un análisis converged.

\subsection{12. Tolerancia de fluencia}

Vía ADR 0006 (\texttt{Material.admissibility\_tol}), patrón atol + rtol\ensuremath{\cdot}escala:

\begin{itemize}
  \item \textbf{Plane strain}: \texttt{admissibility\_scale = \ensuremath{\surd}(2/3)(\ensuremath{\sigma}\_y + H\ensuremath{\cdot}\ensuremath{\alpha})} (norma desviadora característica en la frontera).
  \item \textbf{Plane stress}: \texttt{admissibility\_scale = (\ensuremath{\sigma}\_y + H\ensuremath{\cdot}\ensuremath{\alpha})\ensuremath{^{2}}/3} (escala de $\bar f$ proyectada; los dos términos son del mismo orden).
\end{itemize}


\section{Contrato de implementación}

\begin{lstlisting}[language=yaml]
name: VonMises2D
kind: material
status: validated

interface:
  strain_dim: 3
  primary_state_var: alpha   # deformación plástica acumulada equivalente

parameters:
  - { name: E,          type: float, required: true,  desc: "Módulo de Young" }
  - { name: nu,         type: float, required: true,  desc: "Coeficiente de Poisson" }
  - { name: sigma_y,    type: float, required: true,  desc: "Tensión de fluencia inicial" }
  - { name: H,          type: float, required: false, default: 0.0, desc: "Módulo de endurecimiento isótropo lineal (≥0). H=0 ⇒ plasticidad perfecta" }
  - { name: hypothesis, type: str,   required: false, default: "plane_strain", desc: "plane_strain | plane_stress" }
  - { name: density,    type: float, required: false, default: null, desc: "Densidad (kg/m³). Opcional al construir; obligatoria solo si se ensambla peso propio o masa (ADR 0008)" }

signature:
  compute_state: "(ε: ndarray(3,), state_vars=None) -> (σ: ndarray(3,), C_alg: ndarray(3,3), state_vars')"
  strain_kind: "plano Voigt — [ε_xx, ε_yy, γ_xy] con γ_xy = 2·ε_xy"

state_schema:
  eps_p:   "ndarray(4,) — [ε^p_xx, ε^p_yy, ε^p_zz, ε^p_xy] (xy tensorial, sin factor 2)"
  alpha:   "float ≥ 0 — deformación plástica acumulada equivalente"

conventions:
  sign: "σ > 0 ⇔ tracción (Reglas §5)"
  voigt: "γ_xy = 2·ε_xy en deformación total; ε^p_xy almacenada en componente tensorial (sin factor 2)"
  hypothesis_dispatch: "selección por kwarg al construir; mutuamente excluyentes"

validity:
  - "E > 0, σ_y > 0, H ≥ 0"
  - "ν ∈ (-1, 0.5) en plane_strain; ν ∈ (-1, 0.5) en plane_stress (ν=0.5 no admitido por singularidad)"
  - "|ε| ≲ 1e-2 (régimen de pequeñas deformaciones)"
  - "hypothesis ∈ {'plane_strain', 'plane_stress'}"

out_of_scope:
  - "endurecimiento cinemático (back-stress); usar variante futura VonMises2DKinematic"
  - "endurecimiento isótropo no lineal (Voce, exponencial)"
  - "endurecimiento por ablandamiento (H<0); requiere regularización (longitud característica, no-local)"
  - "regla de flujo no asociada"
  - "acoplamiento con daño; usar modelo plástico-dañado dedicado"
  - "viscoplasticidad y dependencia de la velocidad de deformación"
  - "grandes deformaciones (descomposición multiplicativa F = F^e·F^p)"
  - "anisotropía (Hill, Barlat) — el modelo es estrictamente J2 isótropo"

numerical_caveats:
  - "Plane strain: incompresibilidad plástica genera locking volumétrico en elementos de bajo orden (Quad4, Tri3) cuando ν→0.5 o el campo plástico domina. Mitigaciones futuras: B-bar, F-bar, elementos mixtos u-p. No incluidas en esta spec."
  - "Plane stress: el Newton local sobre Δγ asume H ≥ 0; con H = 0 (perfectamente plástico) la ecuación residual sigue siendo monótona pero su pendiente al final del paso plástico tiende a 0 — convergencia más lenta. Mitigación: bisección de respaldo si Newton local diverge."
  - "Para ν cercano a 0.5 en plane strain, K → ∞ y el predictor elástico pierde precisión numérica. Recomendado ν ≤ 0.49 en metales (típico ν=0.3)."

acceptance:
  # Cubierto por tests/test_materials_unit.py · TestVonMises2D y TestVonMises2DPlaneStress
  unit_plane_strain:
    - name: paso_elastico_bajo_fluencia
      expect: "ε=[1e-5,0,0]: σ=C_e·ε exacto, alpha=0, tangente=C_e"
      tol_rel: 1.0e-10
    - name: fluencia_traccion_uniaxial_confinada
      expect: "fluencia activa con ε=[0.01,-0.01,0]: alpha>0 tras un solo paso"
    - name: monotonicidad_alpha
      expect: "α no decrece bajo carga creciente con eps_yy=-ε; α_final > 0"
    - name: no_further_yield_on_frontier
      expect: "repetir ε en frontera de fluencia no incrementa α (return mapping consistente)"
      tol_rel: 1.0e-10
    - name: hypothesis_validation
      expect: "constructor rechaza hypothesis distinto de 'plane_strain'/'plane_stress' con ValueError"

  unit_plane_stress:
    - name: paso_elastico_bajo_fluencia
      expect: "ε pequeño: σ=C_e^ps·ε, tangente=C_e^ps, alpha=0, eps_p=0"
      tol_rel: 1.0e-10
    - name: degeneracion_a_1d_traccion_pura_elastico
      expect: "ε=(eps,-ν·eps,0): σ_xx=E·eps, σ_yy≈0 (regimen elástico, no E/(1-ν²))"
      tol_rel: 1.0e-4
    - name: yield_uniaxial_at_sigma_y
      expect: "primer yield en σ_xx=σ_y bajo tracción uniaxial pura (no E/(1-ν²)·ε)"
      tol_rel: 1.0e-2
    - name: traccion_biaxial_isotropa
      expect: "fluencia biaxial isótropa exactamente en σ_xx=σ_yy=σ_y (||s||=σ·√(2/3) iguala √(2/3)σ_y); ε_yield=σ_y(1-ν)/E"
      tol_rel: 1.0e-2
    - name: incompresibilidad_plastica
      expect: "tr(ε^p) = ε^p_xx + ε^p_yy + ε^p_zz = 0 tras cualquier paso plástico"
      tol_abs: 1.0e-12
    - name: alpha_monotonic
      expect: "α no decrece bajo carga monotónica creciente"
    - name: descarga_recupera_C_e_ps
      expect: "tras estado plástico, paso siguiente con menor ε produce tangente=C_e^ps y α intacto"
      tol_rel: 1.0e-10
    - name: tangente_simetrica
      expect: "C_alg simétrica (J2 asociado ⇒ IS_SYMMETRIC=True)"
      tol_abs: 1.0e-8
    - name: unit_invariance_MPa_vs_Pa
      expect: "mismo path adimensional en (MPa,mm,N) y (Pa,m,N) ⇒ α y ε^p idénticos"
      tol_rel: 1.0e-12

  # Cubierto por tests/test_solid_2d_plasticity.py
  integration:
    - name: pipeline_plane_strain_elastico
      setup: "Quad4 unitario, confinamiento uy=0; tracción horizontal por debajo del yield"
      expect: "u_x = ε_xx_target exacto; α=0"
      tol_rel: 1.0e-8

    - name: pipeline_plane_strain_plastico_vs_material_directo
      setup: "mismo setup, carga 1.5×σ_xx_yield → régimen plástico; integrar 10 pasos"
      expect: "σ_xx FEM (gauss avg) coincide con material standalone en ε_xx_final bajo trayectoria monótona; α coincide"
      tol_rel: 1.0e-3

    - name: pipeline_plane_stress_elastico_uniaxial
      setup: "Quad4 unitario, lado izq apoyado en ux; nodo (0,0) y (1,0) en uy; nodo (1,1) libre en y. Tracción horizontal por debajo del yield"
      expect: "u_x = ε_xx_target; contracción transversal u_y(0,1) = -ν·ε_xx (Poisson); α=0"
      tol_rel: 1.0e-6

    - name: pipeline_plane_stress_plastico_vs_1d
      setup: "mismo setup, carga 1.2σ_y para entrar plástico; comparar contra Elastoplastic1D bajo mismo path ε_xx"
      expect: "σ_xx FEM = σ_1D, σ_yy FEM ≈ 0 (uniaxial puro emergente); α_FEM = α_1D"
      tol_rel: 1.0e-3

    - name: pipeline_converges_few_iter_plane_strain
      setup: "carga plástica significativa, 4 incrementos, max_iter=8"
      expect: "solver converge en todos los pasos sin agotar max_iter; alpha final > 0"

    - name: pipeline_converges_few_iter_plane_stress
      setup: "tracción uniaxial pura plane stress 1.5σ_y, 4 incrementos, max_iter=8"
      expect: "solver converge en todos los pasos; alpha final > 0"

references:
  - "Simó J.C., Hughes T.J.R. (1998). Computational Inelasticity. Springer. §3.3 (plane strain return mapping), §3.4 (plane stress projected algorithm), Box 3.1 (operador P)."
  - "Crisfield M.A. (1991). Non-linear Finite Element Analysis of Solids and Structures, Vol. 1. Wiley. Cap. 6 (J2 plasticity)."
  - "Chen W.F., Han D.J. (1988). Plasticity for Structural Engineers. Springer. Cap. 5 (cilindro a presión interna, sol. cerrada Hill)."
  - "de Souza Neto E.A., Perić D., Owen D.R.J. (2008). Computational Methods for Plasticity. Wiley. §9.4 (plane stress projection alternativa)."

\end{lstlisting}


\section{Implementación}

\begin{itemize}
  \item \textbf{Archivo}: solidum/materials/von\_mises\_2d.py.
  \item \textbf{Clase}: \texttt{VonMises2D}, registrada vía \texttt{@MaterialRegistry.register}. Despacho por \texttt{hypothesis} en construcción (sin coste runtime).
  \item \textbf{Kernels Numba}:
  \item \texttt{\_compute\_j2\_plane\_strain}: descomposición volumétrica-desviadora 3D ($\varepsilon^p_{zz}$ libre), return mapping radial cerrado, tangente algorítmica consistente cerrada. Predictor elástico construido como $\mathbf s_\text{trial} + p\mathbf I$ con $\mathbf s_\text{trial} = 2G(\mathbf e^\text{dev} - \boldsymbol\varepsilon^p_n)$ --- respeta $\boldsymbol\varepsilon^p_n$ en descargas/reevaluaciones.
  \item \texttt{\_compute\_j2\_plane\_stress}: predictor $\boldsymbol\sigma^\text{trial} = \mathbf C^{ps}_e(\boldsymbol\varepsilon - \boldsymbol\varepsilon^p_n)$, función de fluencia $\bar f$ proyectada, Newton local sobre $\Delta\gamma$ con tangente cerrada y máximo \texttt{\_PLANE\_STRESS\_MAX\_LOCAL\_ITER = 25} iteraciones. Construcción de $\mathbf A^{-1}$ y \ensuremath{\sigma}\_new vía base de autovectores conocida ($\mathbf v_1, \mathbf v_2, \mathbf v_3$); incompresibilidad plástica cierra $\varepsilon^p_{zz} = -(\varepsilon^p_{xx} + \varepsilon^p_{yy})$.
  \item \textbf{State schema}: \texttt{eps\_p} ndarray(4,) \texttt{[xx, yy, zz, xy\_tensorial]}, \texttt{alpha} escalar. Las cuatro componentes de \texttt{eps\_p} son obligatorias en ambas hipótesis para preservar invariantes (incompresibilidad, normas tensoriales).
  \item \textbf{\texttt{admissibility\_scale}}:
  \item plane\_strain: \texttt{\ensuremath{\surd}(2/3)\ensuremath{\cdot}(\ensuremath{\sigma}\_y + H\ensuremath{\cdot}\ensuremath{\alpha})}
  \item plane\_stress: \texttt{(\ensuremath{\sigma}\_y + H\ensuremath{\cdot}\ensuremath{\alpha})\ensuremath{^{2}}/3}
  \item \textbf{Tests}:
  \item tests/test\_materials\_unit.py \ensuremath{\cdot} \texttt{TestVonMises2D} (5 tests plane strain pre-existentes, no regresión) + \texttt{TestVonMises2DPlaneStress} (9 tests plane stress).
  \item tests/test\_solid\_2d\_plasticity.py \ensuremath{\cdot} 6 tests de integración del pipeline Quad4 + VonMises2D + NonlinearSolver en ambas hipótesis.
  \item tests/test\_density\_self\_weight.py \ensuremath{\cdot} \texttt{test\_vonmises2d\_density} --- densidad ADR 0008.
  \item \textbf{Notas de traducción}:
  \item El kernel plane stress trabaja con $\boldsymbol\sigma$ en Voigt mixto: las componentes $\boldsymbol\sigma = [\sigma_{xx},\sigma_{yy},\sigma_{xy}]$ son las "engineering" mientras que $\boldsymbol\varepsilon$ usa $\gamma_{xy} = 2\varepsilon_{xy}$. El operador $\mathbf P$ en esta convención tiene la forma con \texttt{6} en la entrada (3,3); $\mathbf P\boldsymbol\sigma$ devuelve la componente cortante como $2\sigma_{xy}$ (convención engineering), que se convierte a tensorial dividiendo por 2 antes de acumular en $\varepsilon^p_{xy}$.
  \item La actualización $\alpha_{n+1} = \alpha_n + \Delta\gamma\sqrt{(2/3)\boldsymbol\sigma^\top\mathbf P\boldsymbol\sigma}$ se evalúa \textbf{con $\boldsymbol\sigma$ final} del paso plástico (tras converger Newton local). Durante las iteraciones se usa la misma fórmula con el $\boldsymbol\sigma$ corriente.
  \item Validación de parámetros en constructor: \texttt{E>0}, \texttt{\ensuremath{\sigma}\_y>0}, \texttt{H\ensuremath{\ge}0}, \texttt{\ensuremath{\nu}\ensuremath{\in}(-1,0.5)}, \texttt{hypothesis \ensuremath{\in} \{plane\_strain, plane\_stress\}}. Mensajes en español describiendo la violación.
\end{itemize}

\subsection{Bug fixes detectados al implementar}

\begin{itemize}
  \item \textbf{Predictor elástico plane strain ignoraba $\boldsymbol\varepsilon^p$}. Versión previa devolvía $\boldsymbol\sigma = \mathbf C_e\boldsymbol\varepsilon$ cuando $f^\text{trial} \leq 0$ --- correcto solo si $\boldsymbol\varepsilon^p = 0$. En descargas elásticas desde estado plástico, o al reevaluar el material vía \texttt{compute\_gauss\_state(U\_final)} tras un análisis converged, devolvía \ensuremath{\sigma} inconsistente con el estado interno. Corregido a $\boldsymbol\sigma = \mathbf s_\text{trial} + p\mathbf I$. Detectado por el test de integración \texttt{test\_plastic\_regime\_matches\_standalone\_material}: \ensuremath{\sigma} FEM reportado coincidía con el valor de equilibrio (F/A) en Newton (donde sí entra a return mapping), pero \texttt{compute\_gauss\_state} post-converged daba el valor incorrecto del predictor elástico, divergiendo del material standalone.
\end{itemize}


\section{Diálogo}

\begin{itemize}
  \item \textbf{2026-05-14} \ensuremath{\cdot} Spec creada retroactivamente sobre material existente (J2 plane strain) y extendida con J2 plane stress. Decisión de algoritmo plane stress: \textbf{proyección de Simó-Hughes §3.4.1} en lugar de nested Newton sobre $\varepsilon_{zz}$. Justificación: forma cerrada del problema local (Newton local escalar sobre $\Delta\gamma$ con tangente cerrada, 3-6 iteraciones, O(1) por Gauss); Numba-compatible (kernel plano); tangente consistente cerrada vía operador $\mathbf P$. Nested Newton implicaría \textasciitilde{}5-10\ensuremath{\times} el coste por anidación de return mapping completo dentro de cada iteración externa sobre $\sigma_{zz}=0$, justificable solo con endurecimiento no lineal complejo (Voce, plasticidad mixta) que no es el caso. Registrado aquí, no en ADR --- afecta a un material, no a contratos arquitecturales transversales.
  \item \textbf{2026-05-14} \ensuremath{\cdot} Corrección durante la implementación: la regla heurística $\alpha_{n+1} = \alpha_n + \sqrt{2/3}\,\Delta\gamma$ heredada de plane strain rompe la invariancia bajo cambio de unidades en plane stress, porque allí $\Delta\gamma$ tiene unidades $[1/\text{esfuerzo}]$ (la regla de flujo $\dot{\boldsymbol\varepsilon}^p = \dot\gamma\,\mathbf P\boldsymbol\sigma$ tiene $\mathbf P\boldsymbol\sigma$ con dimensión de esfuerzo). La fórmula físicamente correcta es $\alpha_{n+1} = \alpha_n + \Delta\gamma\,\sqrt{(2/3)\boldsymbol\sigma^\top\mathbf P\boldsymbol\sigma}$ --- adimensional. Detectado por el test unitario \texttt{test\_unit\_invariance\_MPa\_vs\_Pa} que daba un factor 1e6 entre los dos sistemas. Implica recalcular la tangente consistente con $\partial\alpha/\partial\Delta\gamma \neq \sqrt{2/3}$ (ver §11).
  \item \textbf{2026-05-14} \ensuremath{\cdot} Corrección plane strain: el predictor elástico devolvía $\boldsymbol\sigma = \mathbf C_e\boldsymbol\varepsilon$, ignorando $\boldsymbol\varepsilon^p_n$. Bug latente preexistente sin tests que lo cubrieran (los unitarios solo probaban primer paso con $\boldsymbol\varepsilon^p = 0$). Detectado al validar el pipeline FEM: \texttt{compute\_gauss\_state(U\_final)} reportaba \ensuremath{\sigma} del predictor elástico (incorrecto) en lugar del \ensuremath{\sigma} post return-mapping. Corregido a $\boldsymbol\sigma = \mathbf s_\text{trial} + p\mathbf I$. La rama plane stress ya lo hacía bien desde el principio (predictor incluía la sustracción de $\boldsymbol\varepsilon^p$).
  \item \textbf{2026-05-14} \ensuremath{\cdot} Validación numérica de plane strain queda implícita en los tests unitarios existentes; los tests de integración \texttt{tests/test\_solid\_2d\_plasticity.py} abren por primera vez el pipeline completo Quad4 + VonMises2D + NonlinearSolver en ambas hipótesis --- hasta hoy no había cobertura de integración sólido 2D + material plástico.
  \item \textbf{2026-05-14} \ensuremath{\cdot} No se introduce ADR. La spec no rompe contratos: extiende un material existente con una rama nueva del despacho por \texttt{hypothesis}. Sí se actualiza \texttt{docs/catalogo\_materiales.md} al validar.
\end{itemize}


\newpage

\chapter{Modelos Constitutivos — 3D}

\section{DruckerPrager3D}



\textit{Orden de trabajo. El usuario revisa \textbf{especificación física}, \textbf{formulación numérica} y \textbf{contrato}; la IA propone, ejecuta y rellena \textbf{implementación} + \textbf{diálogo}.}

\textit{> \textbf{Pre-redactada por la IA} sobre la base de \texttt{DruckerPrager2D.md} (mismo modelo físico) y \texttt{VonMises3D.md} (mismo patrón Voigt 6D). Segunda entrega de la sub-etapa \textbf{A.bis} (materiales 3D no lineales).}



\section{Especificación física}

\subsection{0. Descripción general}

Modelo elastoplástico \textbf{friccional} independiente de la velocidad. Suelos, hormigón, granulares, materiales cuasi-frágiles cohesivo-friccionales en régimen 3D completo. Suaviza la pirámide hexagonal de Mohr-Coulomb a un \textbf{cono circular} en el espacio de tensiones principales, eliminando las aristas que complican el return mapping. Formulado íntegramente en Voigt 6D del proyecto (ADR 0012).

Diferencias clave con J2 3D (\texttt{VonMises3D}):

\begin{itemize}
  \item \textbf{Dependencia de la presión}. La resistencia al corte aumenta con la presión de confinamiento (componente friccional).
  \item \textbf{Plasticidad no asociada}. La dilatancia volumétrica real del flujo plástico ($\psi$) es típicamente menor que la fricción del criterio ($\phi$). Forzar $\psi = \phi$ (asociada) sobreestima la dilatancia.
  \item \textbf{Retorno al ápice}. El criterio define un cono con vértice; estados de prueba en zona puramente hidrostática extrema retornan al vértice (zona singular en la dirección de flujo).
\end{itemize}

Comparado con \texttt{DruckerPrager2D}:

\begin{itemize}
  \item \textbf{Sin variantes de hipótesis} (igual que VM3D): en 3D todas las componentes son activas.
  \item \textbf{El catálogo de calibraciones cambia}: la variante 2D \texttt{plane\_strain\_matched} no tiene análogo en 3D (Mohr-Coulomb en 3D es una pirámide hexagonal sin coincidencia circular global). Las calibraciones 3D disponibles son \texttt{outer\_cone} (circunscribe MC en el meridiano de compresión, \textbf{default}) e \texttt{inner\_cone} (inscribe MC en el meridiano de extensión).
\end{itemize}

\subsection{1. Descomposición aditiva y elasticidad}

Pequeñas deformaciones: $\boldsymbol\varepsilon = \boldsymbol\varepsilon^e + \boldsymbol\varepsilon^p$. Ley elástica isótropa con la matriz $\mathbf C_e$ 6\ensuremath{\times}6 de \texttt{Elastic3D}:

\[
\boldsymbol\sigma = \mathbf C_e\,(\boldsymbol\varepsilon - \boldsymbol\varepsilon^p)
\]

equivalente desviador + presión:

\[
\boldsymbol\sigma = \mathbf s + p\,\mathbf I, \qquad \mathbf s = 2G\,\mathbf e^e_\text{dev}, \qquad p = K\,\operatorname{tr}(\boldsymbol\varepsilon - \boldsymbol\varepsilon^p)
\]

A diferencia de J2, $\operatorname{tr}(\boldsymbol\varepsilon^p) \neq 0$ en general (dilatancia plástica si $\eta_g > 0$), por lo que $p$ depende de la deformación plástica acumulada vía $\operatorname{tr}(\boldsymbol\varepsilon^p)$.

\subsection{2. Criterio de fluencia (Drucker-Prager)}

\[
f(\boldsymbol\sigma, \alpha) = \sqrt{J_2} + \eta_f\,I_1 - k(\alpha) \;\leq\; 0
\]

con:

\begin{itemize}
  \item $J_2 = \tfrac{1}{2}\,\mathbf s : \mathbf s$, segundo invariante del desviador.
  \item $I_1 = \operatorname{tr}(\boldsymbol\sigma) = \sigma_{xx} + \sigma_{yy} + \sigma_{zz}$, primer invariante.
  \item $\eta_f$, coeficiente friccional del cono (adimensional, $\geq 0$).
  \item $k(\alpha) = k_0 + H_k\,\alpha$, parámetro cohesivo con endurecimiento isótropo lineal.
\end{itemize}

La superficie es un \textbf{cono circular} con eje en la línea hidrostática ($\sigma_{xx} = \sigma_{yy} = \sigma_{zz}$), vértice en $I_1 = k/\eta_f$, abertura controlada por $\eta_f$.

\textbf{Cómputo de $J_2$ en Voigt 6D del proyecto} (componentes cortantes tensoriales sin factor 2):

\[
J_2 = \tfrac{1}{2}\,\Big(s_{xx}^2 + s_{yy}^2 + s_{zz}^2 + 2\,(s_{xy}^2 + s_{yz}^2 + s_{xz}^2)\Big)
\]

El factor 2 contabiliza las dos posiciones simétricas del tensor.

\subsection{3. Calibración con Mohr-Coulomb (3D)}

Los parámetros $(\eta_f, k_0)$ se derivan de los parámetros físicos $(c_0, \phi)$ --- cohesión y ángulo de fricción interna --- según la variante del cono elegida. En 3D, \textbf{el cono no puede coincidir con la pirámide hexagonal de MC en todo el plano \ensuremath{\pi}} (las aristas de MC quedan fuera o dentro del círculo), por lo que las calibraciones eligen un meridiano de referencia:

\begin{center}
\begin{tabularx}{\textwidth}{YYYY}
\hline
\textbf{variant} & \textbf{$\eta_f$} & \textbf{$k_0$} & \textbf{Comportamiento} \\
\hline
\texttt{outer\_cone} (default) & $2\sin(\phi)/[\sqrt 3\,(3 - \sin(\phi))]$ & $6\,c_0\cos(\phi)/[\sqrt 3\,(3 - \sin(\phi))]$ & Circunscribe MC en el meridiano de compresión triaxial ($\theta_L = -30°$). Más conservador en compresión, menos conservador en extensión. \\
\texttt{inner\_cone} & $2\sin(\phi)/[\sqrt 3\,(3 + \sin(\phi))]$ & $6\,c_0\cos(\phi)/[\sqrt 3\,(3 + \sin(\phi))]$ & Inscribe MC en el meridiano de extensión triaxial ($\theta_L = +30°$). Cota inferior segura para diseño con MC. \\
\hline
\end{tabularx}
\end{center}
El default \texttt{outer\_cone} se elige porque la compresión triaxial es el modo dominante en problemas geotécnicos típicos (cimentaciones, muros de contención, taludes). Si la aplicación prioriza la extensión (e.g. excavaciones), \texttt{inner\_cone} es la elección segura.

\textbf{Caveat fundamental}: el cono no captura efectos de ángulo de Lode (la pirámide hexagonal de MC sí). Para problemas donde el modo de tensión rota fuera de los meridianos puros, considerar Mohr-Coulomb 3D directo (out-of-scope) o variantes con dependencia explícita del Lode angle (e.g. Matsuoka-Nakai, también out-of-scope).

\textbf{No se incluye \texttt{plane\_strain\_matched}} (variante 2D de DP2D): es un artefacto matemático que coincide con MC exactamente solo cuando $\varepsilon_{zz} = 0$ está impuesto cinemáticamente; no tiene análogo 3D porque MC en 3D depende del Lode angle activo.

\subsection{4. Regla de flujo (no asociada por defecto)}

Potencial plástico $g$ con coeficiente de \textbf{dilatancia} $\eta_g$ (función del ángulo de dilatancia $\psi$, calibrado con la misma fórmula que $\eta_f$ pero usando $\psi$):

\[
g(\boldsymbol\sigma) = \sqrt{J_2} + \eta_g\,I_1
\]

\[
\dot{\boldsymbol\varepsilon}^p = \dot\gamma\,\frac{\partial g}{\partial \boldsymbol\sigma} = \dot\gamma\,\left(\frac{\mathbf s}{2\sqrt{J_2}} + \eta_g\,\mathbf I\right)
\]

Descomposición:

\begin{itemize}
  \item \textbf{Parte desviadora}: $\dot{\mathbf e}^p = \dot\gamma\,\mathbf s/(2\sqrt{J_2})$ --- análoga a J2.
  \item \textbf{Parte volumétrica}: $\operatorname{tr}(\dot{\boldsymbol\varepsilon}^p) = 3\dot\gamma\,\eta_g$ --- \textbf{dilatación plástica positiva} si $\eta_g > 0$.
\end{itemize}

Si $\psi = \phi$ (i.e. $\eta_g = \eta_f$) el modelo se vuelve \textbf{asociado}; tangente algorítmica simétrica. Para $\psi < \phi$ (típico en suelos: $\psi \approx \phi/2$ o incluso $\psi = 0$, "non-dilatant"), el modelo es \textbf{no asociado} y la tangente algorítmica es \textbf{asimétrica}.

\subsection{5. Endurecimiento isótropo lineal}

\[
\dot\alpha = \dot\gamma
\]

con $\alpha$ la \textbf{deformación plástica equivalente} (convención del proyecto para Drucker-Prager). La superficie crece manteniendo eje y abertura, solo $k$ aumenta linealmente. Endurecimiento por fricción/dilatancia (cambio de $\eta_f$ o $\eta_g$ con $\alpha$) queda fuera de esta spec.

\subsection{6. Condiciones de Kuhn-Tucker}

\[
\dot\gamma \geq 0,\qquad f \leq 0,\qquad \dot\gamma\,f = 0
\]

\subsection{7. Variables internas}

\begin{center}
\begin{tabularx}{\textwidth}{YYY}
\hline
\textbf{Símbolo} & \textbf{Tipo} & \textbf{Significado} \\
\hline
$\boldsymbol\varepsilon^p$ & tensor 6 componentes \texttt{[xx, yy, zz, xy, yz, xz]} (cortantes tensoriales) & deformación plástica \\
$\alpha$ & escalar $\geq 0$ & deformación plástica equivalente (acumulado de $\Delta\gamma$) \\
\hline
\end{tabularx}
\end{center}
\texttt{alpha} es la variable principal exportable (\texttt{PRIMARY\_STATE\_VAR = 'alpha'}). $\boldsymbol\varepsilon^p$ \textbf{no es incompresible} en este modelo (a diferencia de J2): $\operatorname{tr}(\boldsymbol\varepsilon^p) = 3\eta_g\,\alpha$ en cualquier estado, lo que es un invariante cinemático verificable por test.

\subsection{8. Notación Voigt 6D}

Convención del proyecto (ADR 0012):

\[
\boldsymbol\varepsilon = [\varepsilon_{xx},\,\varepsilon_{yy},\,\varepsilon_{zz},\,\gamma_{xy},\,\gamma_{yz},\,\gamma_{xz}]^\top, \qquad \gamma_{ij} = 2\,\varepsilon_{ij}
\]

\textbf{Almacenamiento interno de $\boldsymbol\varepsilon^p$}: 6 componentes con cortantes \textbf{tensoriales sin factor 2}, idéntica convención a \texttt{VonMises3D} (consistencia entre materiales 3D no lineales del proyecto).

\subsection{9. Relación con \texttt{DruckerPrager2D} plane strain}

\texttt{DruckerPrager2D} plane strain con variante \texttt{outer\_cone} (o \texttt{inner\_cone}) es la restricción de este modelo bajo $\varepsilon_{zz} = \gamma_{yz} = \gamma_{xz} = 0$ con $\varepsilon^p_{zz} \neq 0$ libre (dilatancia volumétrica). Esta relación se verifica como test cruzado en §13.

\textbf{No se puede establecer la equivalencia con la variante \texttt{plane\_strain\_matched} de DP2D}: esa calibración es 2D-only por construcción. El test cruzado usa \texttt{outer\_cone} (o \texttt{inner\_cone}) que es la única variante compartida.


\section{Formulación numérica}

\subsection{10. Esquema implícito --- Backward Euler}

Dado $\{\boldsymbol\varepsilon^p_n, \alpha_n\}$ y $\boldsymbol\varepsilon_{n+1}$, predictor elástico. Conversión engineering\ensuremath{\to}tensorial dividiendo los cortantes por 2 (única vez al construir el desviador):

\[
\boldsymbol\varepsilon^\text{tens}_{n+1} = [\varepsilon_{xx},\,\varepsilon_{yy},\,\varepsilon_{zz},\,\tfrac{\gamma_{xy}}{2},\,\tfrac{\gamma_{yz}}{2},\,\tfrac{\gamma_{xz}}{2}]
\]

Deformación elástica de prueba (incluye toda la $\boldsymbol\varepsilon^p_n$, incluida su parte volumétrica por dilatancia):

\[
\boldsymbol\varepsilon^{e,\text{trial}}_{n+1} = \boldsymbol\varepsilon^\text{tens}_{n+1} - \boldsymbol\varepsilon^p_n
\]

Descomposición volumétrica-desviadora:

\[
\operatorname{tr}(\boldsymbol\varepsilon^{e,\text{trial}}_{n+1}) = (\varepsilon_{xx} + \varepsilon_{yy} + \varepsilon_{zz}) - (\varepsilon^p_{xx,n} + \varepsilon^p_{yy,n} + \varepsilon^p_{zz,n})
\]

\[
\mathbf e^\text{trial} = \operatorname{dev}(\boldsymbol\varepsilon^{e,\text{trial}}_{n+1}), \qquad \mathbf s^\text{trial} = 2G\,\mathbf e^\text{trial}, \qquad p^\text{trial} = K\,\operatorname{tr}(\boldsymbol\varepsilon^{e,\text{trial}}_{n+1})
\]

con $I_1^\text{trial} = 3\,p^\text{trial}$ y $\sqrt{J_2^\text{trial}} = \|\mathbf s^\text{trial}\|/\sqrt{2}$ (norma Frobenius en Voigt 6D tensorial).

Función de fluencia trial:

\[
f^\text{trial} = \sqrt{J_2^\text{trial}} + \eta_f\,I_1^\text{trial} - k(\alpha_n)
\]

Si $f^\text{trial} \leq \text{tol}$ \ensuremath{\to} paso elástico. Si no \ensuremath{\to} return mapping.

\subsection{11. Return regular (cone surface)}

La dirección desviadora $\hat{\mathbf n} = \mathbf s^\text{trial}/(2\sqrt{J_2^\text{trial}})$ es invariante bajo carga radial. Las actualizaciones son lineales en $\Delta\gamma$:

\[
\sqrt{J_2^{n+1}} = \sqrt{J_2^\text{trial}} - G\,\Delta\gamma
\]
\[
p^{n+1} = p^\text{trial} - 3K\,\eta_g\,\Delta\gamma, \qquad I_1^{n+1} = 3\,p^{n+1}
\]

La consistencia $f(\boldsymbol\sigma_{n+1}, \alpha_{n+1}) = 0$ con $\alpha_{n+1} = \alpha_n + \Delta\gamma$ y $k_{n+1} = k_0 + H_k(\alpha_n + \Delta\gamma)$ da:

\[
\boxed{\;\Delta\gamma = \frac{f^\text{trial}}{G + 9K\,\eta_f\,\eta_g + H_k}\;}
\]

\textbf{Forma cerrada} --- sin Newton local. Idéntica al 2D plane strain en estructura algebraica.

\textbf{Actualizaciones}:

\[
\mathbf s^{n+1} = \frac{\sqrt{J_2^{n+1}}}{\sqrt{J_2^\text{trial}}}\,\mathbf s^\text{trial}, \qquad \boldsymbol\sigma^{n+1} = \mathbf s^{n+1} + p^{n+1}\,\mathbf I
\]

\[
\boldsymbol\varepsilon^p_{n+1} = \boldsymbol\varepsilon^p_n + \Delta\gamma\,\left(\frac{\mathbf s^{n+1}}{2\sqrt{J_2^{n+1}}} + \eta_g\,\mathbf I\right)
\]

donde $\mathbf I = [1,1,1,0,0,0]$ en Voigt 6D (parte volumétrica) y la componente desviadora de $\Delta\boldsymbol\varepsilon^p$ se almacena en convención tensorial.

\textbf{Condición de validez}: $\sqrt{J_2^{n+1}} \geq 0$. Si $\Delta\gamma_\text{reg} > \sqrt{J_2^\text{trial}}/G$, el desviador "se pasa de cero" y el algoritmo regular es inválido \ensuremath{\to} toca \textbf{return al ápice}.

\subsection{12. Return al ápice}

Cuando el estado de prueba cae en zona puramente hidrostática extrema (predictor muy traccionante respecto a la cohesión), la proyección al cono cae en su vértice, donde $\mathbf s = 0$ y la dirección desviadora $\hat{\mathbf n}$ no está definida.

En el ápice: $\boldsymbol\sigma^{n+1} = (k_{n+1}/(3\eta_f))\,\mathbf I$ (puramente hidrostático), $\sqrt{J_2^{n+1}} = 0$.

\[
\boxed{\;\Delta\gamma_\text{apex} = \frac{I_1^\text{trial}\,\eta_f - k(\alpha_n)}{9K\,\eta_f\,\eta_g + H_k}\;}
\]

(Equivalente a la fórmula regular sin el término $G$, porque la rama de apex es puramente volumétrica.)

\textbf{Actualizaciones (apex)}:

\[
\mathbf s^{n+1} = 0,\qquad p^{n+1} = \frac{k_{n+1}}{3\,\eta_f}, \qquad \boldsymbol\sigma^{n+1} = p^{n+1}\,\mathbf I
\]

\[
\boldsymbol\varepsilon^p_{n+1} = \boldsymbol\varepsilon^p_n + \mathbf e^\text{trial} + \Delta\gamma_\text{apex}\,\eta_g\,\mathbf I
\]

(El desviador trial se absorbe completamente; la parte volumétrica del flujo añade $3\eta_g\Delta\gamma$ a la traza, distribuida $\eta_g\Delta\gamma$ por componente diagonal.)

\subsection{13. Detección del régimen de retorno}

Algoritmo de despacho (idéntico al 2D):

\begin{enumerate}
  \item Calcular $\Delta\gamma_\text{reg}$ por la fórmula regular.
  \item Si $\sqrt{J_2^\text{trial}} - G\,\Delta\gamma_\text{reg} \geq 0$ \ensuremath{\to} return regular.
  \item Si no \ensuremath{\to} return al ápice con $\Delta\gamma_\text{apex}$.
\end{enumerate}

\subsection{14. Tangente algorítmica consistente}

\textbf{Return regular} ($\sqrt{J_2^{n+1}} > 0$), forma compacta (de Souza Neto §8.3.4 extendida a 6D):

\[
\mathbf C^\text{alg} = K\,\mathbf v\otimes\mathbf v + 2G\,\Big(1 - \tfrac{G\Delta\gamma}{\sqrt{J_2^\text{trial}}}\Big)\,\mathbf I_\text{dev} + 4G\,\beta\,\hat{\mathbf n}\otimes\hat{\mathbf n} - \frac{1}{A}\,\mathbf b_g \otimes \mathbf b_f
\]

con:

\begin{itemize}
  \item $\mathbf v = [1, 1, 1, 0, 0, 0]$ (gradiente de $I_1$ en Voigt 6D engineering input).
  \item $\mathbf I_\text{dev}$: proyector desviador 6\ensuremath{\times}6 con $1/2$ en los cortantes (mapea engineering $\gamma_{ij}$ a tensorial $\varepsilon_{ij}$).
  \item $\hat{\mathbf n} = \mathbf s^\text{trial}/(2\sqrt{J_2^\text{trial}})$ en Voigt 6D tensorial.
  \item $\beta = G\,\Delta\gamma/\sqrt{J_2^\text{trial}}$.
  \item $\mathbf b_f = 2G\,\hat{\mathbf n} + 3K\,\eta_f\,\mathbf v$, $\mathbf b_g = 2G\,\hat{\mathbf n} + 3K\,\eta_g\,\mathbf v$.
  \item $A = G + 9K\,\eta_f\,\eta_g + H_k$.
\end{itemize}

\textbf{Asimétrica si $\eta_f \neq \eta_g$} (no asociada). El término $+4G\beta\,\hat{\mathbf n}\otimes\hat{\mathbf n}$ recoge el cambio de dirección de flujo con $\boldsymbol\varepsilon$ a través de $\mathbf s^\text{trial}$ (no solo el cambio de magnitud); sin él la tangente cierra mal en componentes de cortante alineados con el flujo (verificable por diferencia finita, idéntico al hallazgo del DP2D).

\textbf{Return al ápice} ($\sqrt{J_2^{n+1}} = 0$):

\[
\mathbf C^\text{alg}_\text{apex} = \frac{K\,H_k}{9K\eta_f\eta_g + H_k}\,\mathbf v\otimes\mathbf v
\]

Rigidez puramente volumétrica reducida; sin rigidez desviadora alguna. Con $H_k = 0$ y $\eta_g > 0$ la rigidez tangente colapsa por completo --- el código devuelve un escalado mínimo de $K$ (idéntico patrón al 2D) para evitar matriz globalmente singular hasta que el siguiente paso salga del ápice.

\subsection{15. Tolerancia de fluencia}

ADR 0006: \texttt{admissibility\_scale = k(\ensuremath{\alpha}) = k\_0 + H\_k\ensuremath{\cdot}\ensuremath{\alpha}} (escala cohesiva corriente del cono). Tiene unidades de esfuerzo, igual que $f$.


\section{Contrato de implementación}

\begin{lstlisting}[language=yaml]
name: DruckerPrager3D
kind: material
status: validated        # draft → implemented → validated

interface:
  strain_dim: 6
  primary_state_var: alpha
  is_symmetric: false      # tangente asimétrica si no asociada (ψ ≠ φ);
                           # override por instancia a True cuando ψ = φ (mismo
                           # patrón que DP2D — el despachador algebraico ADR
                           # 0003 elige Cholesky/LDLᵀ en problemas asociados).

parameters:
  - { name: E,          type: float, required: true,  desc: "Módulo de Young (>0)" }
  - { name: nu,         type: float, required: true,  desc: "Coeficiente de Poisson ∈ (-1, 0.5)" }
  - { name: cohesion,   type: float, required: true,  desc: "Cohesión c_0 inicial (esfuerzo, >0)" }
  - { name: phi_deg,    type: float, required: true,  desc: "Ángulo de fricción interna en grados; típico suelos 20-40°, hormigón 30-37°" }
  - { name: psi_deg,    type: float, required: false, default: null,
      desc: "Ángulo de dilatancia en grados. Default = phi_deg (asociada). Típico no asociada: ψ < φ, e.g. 0 o φ/2." }
  - { name: H,          type: float, required: false, default: 0.0,
      desc: "Endurecimiento isótropo lineal en cohesión H_k (≥0). H=0 ⇒ perfectamente plástico." }
  - { name: variant,    type: str,   required: false, default: "outer_cone",
      desc: "outer_cone (default, circunscribe MC en compresión) | inner_cone (inscribe MC en extensión)" }
  - { name: density,    type: float, required: false, default: null, desc: "Densidad (ADR 0008)" }

signature:
  compute_state: "(ε: ndarray(6,), state_vars=None) -> (σ: ndarray(6,), C_alg: ndarray(6,6), state_vars')"
  strain_kind: "3D Voigt — [ε_xx, ε_yy, ε_zz, γ_xy, γ_yz, γ_xz] con γ_ij = 2·ε_ij"

state_schema:
  eps_p: "ndarray(6,) — [ε^p_xx, ε^p_yy, ε^p_zz, ε^p_xy, ε^p_yz, ε^p_xz] (cortantes tensoriales, sin factor 2). tr(ε^p) = 3·η_g·α invariante cinemático."
  alpha: "float ≥ 0 — multiplicador plástico acumulado (= deformación plástica equivalente)"

conventions:
  sign:  "σ > 0 ⇔ tracción (Reglas §5)"
  voigt: "[xx, yy, zz, xy, yz, xz] (ADR 0012); γ_ij = 2·ε_ij engineering input;
          ε^p_ij almacenada tensorial; σ_ij tensorial sin factor"
  associated_flow: "ψ = φ ⇒ asociada (tangente simétrica); ψ ≠ φ ⇒ no asociada (tangente asimétrica)"

validity:
  - "E > 0, c_0 > 0, 0 ≤ φ < 90°, 0 ≤ ψ ≤ φ"
  - "ν ∈ (-1, 0.5) estricto"
  - "H ≥ 0"
  - "variant ∈ {'outer_cone', 'inner_cone'}"

out_of_scope:
  - "Mohr-Coulomb 3D con aristas (pirámide hexagonal): el cono circular DP no captura efectos de ángulo de Lode"
  - "Variantes con dependencia explícita del Lode angle (Matsuoka-Nakai, Lade-Duncan): formulaciones distintas"
  - "Calibración 'plane_strain_matched': sin análogo 3D (la coincidencia DP↔MC en plane strain es 2D-only)"
  - "Endurecimiento por cambio de φ y/o ψ con α (no lineal fuerte en return mapping)"
  - "Endurecimiento cinemático (back-stress hidrostático y desviador)"
  - "Cap (modelo cap-cone) para suelos compactables"
  - "Ablandamiento (H<0) — requiere regularización"
  - "Grandes deformaciones"

numerical_caveats:
  - "Tangente algorítmica asimétrica si ψ ≠ φ → IS_SYMMETRIC=False, el despachador algebraico usa LU (ADR 0003). Cuando ψ = φ se vuelve simétrica y el override por instancia activa Cholesky/LDLᵀ."
  - "Modos de carga muy traccionantes pueden caer en la rama de ápice. La detección automática del régimen (regular vs apex) está en el algoritmo."
  - "Con ψ = 0 (sin dilatancia) y H_k = 0, en la rama de ápice la tangente volumétrica se anula → singularidad. El código devuelve un valor de respaldo (rigidez volumétrica reducida no nula, K·1e-6) para evitar matriz singular; típicamente sale del ápice en el siguiente paso."
  - "Locking volumétrico en Hex8/Tet4 cuando ψ > 0 (flujo dilatante) combinado con ν moderado-alto. Mismo régimen documentado para Elastic3D/VonMises3D; mitigaciones B-bar/F-bar diferidas."
  - "El cono circular DP es independiente del ángulo de Lode. Para trayectorias de carga donde el modo rota fuera del meridiano de calibración (compresión triaxial → cortante → extensión triaxial), la diferencia DP vs MC puede ser hasta el 15-20% en la frontera. Es limitación inherente del modelo, no del algoritmo."

acceptance:
  unit:
    - name: parameter_validation
      setup: "constructor con variant inválida, φ ≥ 90°, ψ > φ, E≤0, c_0≤0, H<0, ν∉(-1,0.5)"
      expect: "ValueError con mensaje claro en cada caso"

    - name: calibration_outer_inner_consistency
      setup: "comparar η_f de outer_cone vs inner_cone para φ=30°"
      expect: "outer_cone η_f > inner_cone η_f (el outer es geométricamente mayor); ambos > 0; k_0 análogo"

    - name: paso_elastico_below_yield
      setup: "ε pequeño 6D, todas las componentes activas"
      expect: "σ = C_e·ε; eps_p y alpha intactos"
      tol_rel: 1.0e-10

    - name: regular_return_under_pure_shear
      setup: "ε cortante puro suficiente para fluir (γ_xy creciente, otras componentes 0)"
      expect: "return regular activo, α > 0; estado final cumple f ≤ tol"
      tol_rel: 1.0e-8

    - name: apex_return_under_hydrostatic_tension
      setup: "ε hidrostática traccionante (ε_xx = ε_yy = ε_zz > 0 suficientemente grande, sin cortante)"
      expect: "return al ápice; σ_xx = σ_yy = σ_zz = k(α)/(3η_f); √J₂ ≈ 0"
      tol_rel: 1.0e-8

    - name: tr_eps_p_invariant
      setup: "carga monotónica en rama regular (5 niveles), ψ ≠ φ (no asociada con η_g > 0)"
      expect: "tr(ε^p) = 3·η_g·α en cada paso plástico"
      tol_rel: 1.0e-10

    - name: associated_tangent_is_symmetric
      setup: "ψ = φ (asociada), estado plástico activo en rama regular"
      expect: "C_alg simétrica (‖C_alg − C_alg^T‖_∞ ≤ tol)"
      tol_abs: 1.0e-8

    - name: nonassociated_tangent_is_asymmetric
      setup: "ψ = 0, φ = 30°, estado plástico activo no degenerado"
      expect: "‖C_alg − C_alg^T‖_∞ con magnitud significativa relativa a ‖C_alg‖_∞"

    - name: tangent_matches_finite_difference
      setup: "estado plástico regular; comparar C_alg cerrada con diferencia finita centrada"
      expect: "error relativo < 1e-4 en todas las componentes"
      tol_rel: 1.0e-4

    - name: alpha_monotonic
      setup: "carga creciente con plasticidad sostenida"
      expect: "α no decrece nunca"

    - name: unloading_recovers_C_e
      setup: "cargar a plástico, descargar elásticamente"
      expect: "tangente = C_e en el paso de descarga; α se conserva"
      tol_rel: 1.0e-10

    - name: unit_invariance_MPa_vs_Pa
      setup: "mismo path adimensional con (MPa,mm,N) y (Pa,m,N)"
      expect: "α y eps_p (adimensionales) idénticos; σ/E idéntico"
      tol_rel: 1.0e-12

  cross_consistency:
    - name: equivalencia_plane_strain_outer_cone
      setup: "DP3D con variant='outer_cone' y ε = (ε_xx, ε_yy, 0, γ_xy, 0, 0);
             DP2D con variant='outer_cone' (plane_strain) y ε = (ε_xx, ε_yy, γ_xy);
             path multi-paso con plasticidad activa en rama regular"
      expect: "σ_xx, σ_yy, σ_xy y α coinciden entre 3D y 2D PS;
              componentes plásticas planas coinciden (xx, yy, zz, xy_tens);
              ε^p_yz = ε^p_xz = 0 en VM3D (no se activan)"
      tol_rel: 1.0e-10

  integration:
    - name: hex8_pipeline_elastico
      setup: "Hex8 unitario, carga muy por debajo del yield"
      expect: "respuesta elástica pura, α=0 en todos los Gauss points"
      tol_rel: 1.0e-8

    - name: hex8_pipeline_plastico_converges
      setup: "Hex8 unitario con confinamiento + cortante; carga plástica significativa; 4 pasos, max_iter=8"
      expect: "Newton converge en cada paso; α > 0 al final; tangente asimétrica activa LU"

    - name: hex8_apex_under_hydrostatic
      setup: "Hex8 unitario con ε hidrostática traccionante (todos los nodos del cubo desplazados radialmente desde el centro)"
      expect: "estado final cae en ápice en los 8 puntos de Gauss: σ_xx ≈ σ_yy ≈ σ_zz, √J₂ ≈ 0"
      tol_rel: 1.0e-6

    - name: tet4_basic_pipeline
      setup: "Tet4 con confinamiento + cortante, plasticidad activa"
      expect: "convergencia + α > 0 + tr(ε^p) = 3·η_g·α invariante en el único Gauss point"

references:
  - "Drucker D.C., Prager W. (1952). Soil mechanics and plastic analysis or limit design. Quart. Appl. Math. 10, 157-165."
  - "Simó J.C., Hughes T.J.R. (1998). Computational Inelasticity. Springer. §6 (cone plasticity, apex return)."
  - "de Souza Neto E.A., Perić D., Owen D.R.J. (2008). Computational Methods for Plasticity. Wiley. §8 (Drucker-Prager 3D, calibraciones outer/inner, tangentes algorítmicas detalladas)."
  - "Chen W.F., Han D.J. (1988). Plasticity for Structural Engineers. Springer. §7 (cohesivo-friccionales 3D)."
  - "ADR 0012 — Sólidos 3D: convención Voigt 6D del proyecto."
  - "ADR 0003 — Despachador algebraico: tangente asimétrica fuerza LU global."
  - "ADR 0006 — Tolerancias adimensionales: patrón atol + rtol·escala con escala física k(α)."

\end{lstlisting}


\section{Implementación}

\begin{itemize}
  \item \textbf{Archivo}: solidum/materials/drucker\_prager\_3d.py.
  \item \textbf{Clase}: \texttt{DruckerPrager3D}, registrada vía \texttt{@MaterialRegistry.register}. \texttt{IS\_SYMMETRIC = False} a nivel de clase; override por instancia a \texttt{True} cuando \texttt{\ensuremath{\psi} = \ensuremath{\varphi}} (asociada).
  \item \textbf{Kernel Numba} \texttt{\_compute\_drucker\_prager\_3d}: predictor elástico desviador-volumétrico en Voigt 6D tensorial (sustracción de $\boldsymbol\varepsilon^p_n$ antes del desviador --- \textit{necesario} en DP por la dilatancia plástica), detección automática regular/apex, tangente cerrada para cada rama.
  \item \textbf{Calibración}: reusa la función \texttt{\_calibrate\_drucker\_prager} de \texttt{drucker\_prager\_2d.py} por composición (mismas fórmulas para outer/inner; \texttt{plane\_strain\_matched} queda filtrada por el validador de variant antes de llegar a la función compartida).
  \item \textbf{State schema}: \texttt{eps\_p} ndarray(6,) \texttt{[xx, yy, zz, xy, yz, xz]} con cortantes tensoriales, \texttt{alpha} escalar. \textbf{No es incompresible}: \texttt{tr(eps\_p) = 3\ensuremath{\cdot}\ensuremath{\eta}\_g\ensuremath{\cdot}\ensuremath{\alpha}} por construcción del flujo.
  \item \textbf{\texttt{admissibility\_scale}}: \texttt{k(\ensuremath{\alpha}) = k\_0 + H\ensuremath{\cdot}\ensuremath{\alpha}} (cohesión efectiva corriente).
  \item \textbf{Tests}:
  \item tests/test\_materials\_unit.py \ensuremath{\cdot} \texttt{TestDruckerPrager3D} (17 tests unitarios) + \texttt{TestDruckerPrager3DvsPlaneStrain} (1 cruzado contra DP2D outer\_cone plane\_strain).
  \item tests/test\_solid\_3d\_plasticity.py \ensuremath{\cdot} 4 tests de integración: \texttt{TestHex8DruckerPrager3DElastic} (1 --- respuesta elástica), \texttt{TestHex8DruckerPrager3DRegularPlastic} (1 --- cortante puro + tangente asimétrica activa LU + invariante $\operatorname{tr}(\boldsymbol\varepsilon^p) = 3\eta_g\alpha$), \texttt{TestHex8DruckerPrager3DApex} (1 --- expansión hidrostática a 8 GP en rama ápice), \texttt{TestTet4DruckerPrager3DBasic} (1 --- smoke test Tet4).
\end{itemize}

\begin{itemize}
  \item \textbf{Notas de traducción}:
  \item Convención mixta engineering/tensorial idéntica a VM3D: entrada \texttt{strain} engineering (\texttt{\ensuremath{\gamma}\_ij = 2\ensuremath{\cdot}\ensuremath{\varepsilon}\_ij}), salida \texttt{sigma} con cortantes tensoriales, estado \texttt{eps\_p} con cortantes tensoriales. Conversión engineering\ensuremath{\to}tensorial una sola vez al construir el desviador (\texttt{strain[3..5]/2}).
  \item El predictor construye el desviador como $\mathbf s^\text{trial} = 2G\,\operatorname{dev}(\boldsymbol\varepsilon^\text{tens} - \boldsymbol\varepsilon^p_n)$, \textbf{no} como $2G\,(\operatorname{dev}(\boldsymbol\varepsilon^\text{tens}) - \boldsymbol\varepsilon^p_n)$. A diferencia de J2, $\boldsymbol\varepsilon^p$ en DP tiene parte volumétrica no nula (dilatancia plástica si $\eta_g > 0$), así que la sustracción debe hacerse \textit{antes} del operador desviador para que la parte volumétrica del flujo se contabilice correctamente.
  \item Tangente regular usa $\mathbf I_\text{dev}$ 6\ensuremath{\times}6 con $1/2$ en los cortantes (mapea engineering input a tensorial output), idéntica estructura al I\_dev de VM3D.
  \item Tangente apex: rigidez puramente volumétrica reducida con escalado mínimo $K \cdot 10^{-6}$ cuando $H_k = 0$ y $\eta_g$ pequeño (mismo patrón que DP2D para evitar matriz globalmente singular en el paso degenerado).
  \item Validación de parámetros en constructor: \texttt{E > 0}, \texttt{c\_0 > 0}, \texttt{0 \ensuremath{\le} \ensuremath{\varphi} < 90\ensuremath{^{\circ}}}, \texttt{0 \ensuremath{\le} \ensuremath{\psi} \ensuremath{\le} \ensuremath{\varphi}}, \texttt{H \ensuremath{\ge} 0}, \texttt{\ensuremath{\nu} \ensuremath{\in} (-1, 0.5)}, \texttt{variant \ensuremath{\in} \{outer\_cone, inner\_cone\}}, \texttt{density \ensuremath{\ge} 0}. Mensajes en español; rechazo explícito de \texttt{plane\_strain\_matched} con explicación de por qué es 2D-only.
\end{itemize}


\section{Diálogo}

\begin{itemize}
  \item \textbf{2026-05-21} \ensuremath{\cdot} Spec pre-redactada por la IA al avanzar la sub-etapa \textbf{A.bis} (materiales 3D no lineales) tras cerrar \texttt{VonMises3D} con 824 tests verdes. Base: \texttt{DruckerPrager2D.md} (modelo físico, calibraciones, two-branch return mapping) + \texttt{VonMises3D.md} (estructura Voigt 6D, patrón cross-consistency).
\end{itemize}

\begin{itemize}
  \item \textbf{2026-05-21} \ensuremath{\cdot} \textbf{Decisión de plumbing (decide la IA, Reglas §3)}: kernels Numba separados \texttt{\_compute\_drucker\_prager\_3d} (este material) vs \texttt{\_compute\_drucker\_prager\_plane\_strain} (existente en DP2D). \textbf{No} se extrae un \texttt{\_DPCore} compartido en esta entrega, por las mismas razones que motivaron mantener separados los kernels de VM2D plane strain y VM3D: Numba \textit{nopython} desfavorece arrays de tamaño variable, los kernels son cortos (\textasciitilde{}80 líneas cada uno), y el sub-patrón compartido (descomposición desviadora, branch detection, tangente cerrada) se documenta cruzadamente entre las specs.
\end{itemize}

\begin{itemize}
  \item \textbf{2026-05-21} \ensuremath{\cdot} \textbf{Decisión de alcance de variantes}: \texttt{DruckerPrager3D} ofrece \textbf{\texttt{outer\_cone}} (default) e \textbf{\texttt{inner\_cone}}, omitiendo \texttt{plane\_strain\_matched}. Razón: la variante 2D \texttt{plane\_strain\_matched} es un artefacto matemático que reproduce MC exactamente solo bajo $\varepsilon_{zz} = 0$ cinemático impuesto; no tiene análogo 3D porque MC en 3D depende del Lode angle activo. Default \texttt{outer\_cone} se elige por ser el modo dominante en problemas geotécnicos (cimentaciones, muros). Si la aplicación es extension-dominated (excavaciones), el usuario puede elegir \texttt{inner\_cone} explícitamente.
\end{itemize}

\begin{itemize}
  \item \textbf{2026-05-21} \ensuremath{\cdot} \textbf{Decisión de tipo de tests cruzados}: \texttt{cross\_consistency} compara DP3D (\texttt{variant='outer\_cone'}) contra DP2D (\texttt{variant='outer\_cone'}, plane\_strain), no contra \texttt{plane\_strain\_matched}. Razón: solo las dos calibraciones cono (outer/inner) existen en ambos modelos. La equivalencia matemática es exacta porque las fórmulas de calibración no dependen de la dimensión --- solo de $(c_0, \phi, \psi)$ --- y el cono 3D con $\varepsilon_{zz}=\gamma_{yz}=\gamma_{xz}=0$ impuesto se reduce a la formulación 2D plane strain en las componentes activas.
\end{itemize}

\begin{itemize}
  \item \textbf{2026-05-21} \ensuremath{\cdot} \textbf{Decisión de plumbing}: el predictor elástico construye el desviador como $\mathbf s^\text{trial} = 2G\,\operatorname{dev}(\boldsymbol\varepsilon^\text{tens} - \boldsymbol\varepsilon^p_n)$, \textbf{no} como $2G\,(\operatorname{dev}(\boldsymbol\varepsilon^\text{tens}) - \boldsymbol\varepsilon^p_n)$. Diferencia material: a diferencia de J2, $\boldsymbol\varepsilon^p$ en DP tiene parte volumétrica no nula ($\operatorname{tr}(\boldsymbol\varepsilon^p) = 3\eta_g\alpha$), así que la sustracción debe hacerse \textit{antes} del desviador para que la parte volumétrica del flujo plástico se contabilice correctamente. Esta es la lección documentada en el comentario del kernel de DP2D (línea 39--42) y se aplica desde el inicio en DP3D.
\end{itemize}

\begin{itemize}
  \item \textbf{2026-05-21} \ensuremath{\cdot} \textbf{Implementación completa, status \texttt{validated}}. Kernel + clase + 17 tests unitarios + 1 test cruzado + 4 tests de integración (Hex8 elástico, Hex8 rama regular con tangente asimétrica + invariante $\operatorname{tr}(\boldsymbol\varepsilon^p) = 3\eta_g\alpha$, Hex8 rama ápice bajo expansión hidrostática, Tet4 smoke test). \textbf{24 tests añadidos a la suite (824 \ensuremath{\to} 848 verdes; suite global sin regresiones).} Todos los tests pasaron a la primera, incluyendo el cruzado DP3D \ensuremath{\leftrightarrow} DP2D outer\_cone plane\_strain a 10 decimales --- señal fuerte de que la convención Voigt 6D, la calibración compartida y la sustracción de $\boldsymbol\varepsilon^p_n$ antes del desviador son correctas. La integración Hex8 con cortante puro y tangente asimétrica ($\psi \neq \phi$) confirma que el despachador algebraico (ADR 0003) selecciona LU correctamente vía el flag declarativo \texttt{IS\_SYMMETRIC = False}. Validación contra solución analítica cerrada (cilindro con falla DP, esfera hueca bajo presión interna, etc.) \textbf{diferida a la campaña 3D consolidada} que se hará al final de la sub-etapa A.bis (tras completar \texttt{IsotropicDamage3D}).
\end{itemize}


\newpage

\section{Elastic3D}



\textit{Material elástico isótropo en 3D, espejo natural de \texttt{Elastic2D} con la convención Voigt 6D fijada en ADR 0012.}



\section{Especificación física}

\subsection{0. Descripción general}

Modelo elástico lineal isótropo tridimensional, parametrizado por $(E, \nu)$. A diferencia de \texttt{Elastic2D}, \textbf{no tiene variantes de hipótesis cinemática}: en 3D no aplican \texttt{plane\_stress} ni \texttt{plane\_strain} --- toda la deformación es activa. Independiente de la velocidad, sin variables internas. Es el material elástico canónico para todos los sólidos 3D del catálogo.

\subsection{1. Descomposición de la deformación}

No aplica --- toda la deformación es elástica:

\[
\boldsymbol\varepsilon = \boldsymbol\varepsilon^e
\]

\subsection{2. Ley elástica}

\[
\boldsymbol\sigma = \mathbf C_e\,\boldsymbol\varepsilon
\]

con $\boldsymbol\varepsilon$ y $\boldsymbol\sigma$ en \textbf{notación Voigt 6D del proyecto} (ADR 0012, \texttt{Reglas.md §5}):

\[
\boldsymbol\varepsilon = [\varepsilon_{xx},\ \varepsilon_{yy},\ \varepsilon_{zz},\ \gamma_{xy},\ \gamma_{yz},\ \gamma_{xz}]^\top, \qquad \gamma_{ij} = 2\,\varepsilon_{ij}
\]

\[
\boldsymbol\sigma = [\sigma_{xx},\ \sigma_{yy},\ \sigma_{zz},\ \sigma_{xy},\ \sigma_{yz},\ \sigma_{xz}]^\top
\]

\textbf{Matriz constitutiva elástica 6\ensuremath{\times}6}:

\[
\mathbf C_e = \frac{E}{(1+\nu)(1-2\nu)}\begin{bmatrix}
1-\nu & \nu & \nu & 0 & 0 & 0 \\
\nu & 1-\nu & \nu & 0 & 0 & 0 \\
\nu & \nu & 1-\nu & 0 & 0 & 0 \\
0 & 0 & 0 & \tfrac{1-2\nu}{2} & 0 & 0 \\
0 & 0 & 0 & 0 & \tfrac{1-2\nu}{2} & 0 \\
0 & 0 & 0 & 0 & 0 & \tfrac{1-2\nu}{2}
\end{bmatrix}
\]

Tangente constante $\mathbf C_e$ en todo régimen. Simétrica y positiva definida en el rango admisible de $\nu$.

El factor $\tfrac{1-2\nu}{2}$ en los cortantes (no $1-2\nu$) es consecuencia del uso de $\gamma_{ij}$ \textit{engineering}: $\sigma_{ij} = G\,\gamma_{ij}$ con $G = \tfrac{E}{2(1+\nu)} = \tfrac{E}{(1+\nu)(1-2\nu)}\cdot\tfrac{1-2\nu}{2}$.

\subsection{3-6. Superficie de fluencia, flujo, endurecimiento, Kuhn-Tucker}

No aplican --- el material es indefinidamente elástico.

\subsection{7. Variables internas}

Ninguna --- \texttt{state\_vars = None} se acepta y se devuelve intacto. Sin \texttt{PRIMARY\_STATE\_VAR}.

\subsection{8. Notación Voigt}

Orden y factores fijados por \textbf{ADR 0012}:

\begin{itemize}
  \item Componentes diagonales primero: \texttt{[xx, yy, zz]}.
  \item Bloque cortante después: \texttt{[xy, yz, xz]} (extensión natural del orden 2D del proyecto \texttt{[xx, yy, xy]}).
  \item Deformaciones angulares \textit{engineering}: $\gamma_{ij} = 2\,\varepsilon_{ij}$.
  \item Esfuerzos tensoriales sin factor.
\end{itemize}

\textbf{Compatibilidad con códigos externos}: ABAQUS usa \texttt{[11, 22, 33, 12, 13, 23]}. La diferencia con la convención del proyecto es la permutación del bloque cortante (\texttt{yz \ensuremath{\leftrightarrow} xz}). Si se importan datos desde ABAQUS, aplicar permutación 5\ensuremath{\leftrightarrow}6 una vez en el preprocesador.


\section{Formulación numérica}

\subsection{9. Esquema temporal}

No aplica --- sin historia.

\subsection{10. Return mapping / actualización}

No aplica. Evaluación reducida a producto matriz-vector:

\[
\boldsymbol\sigma_{n+1} = \mathbf C_e\,\boldsymbol\varepsilon_{n+1}, \qquad \mathbf C^{\text{alg}} = \mathbf C_e
\]

La matriz $\mathbf C_e$ se precalcula y guarda en el constructor --- coste de \texttt{compute\_state} es un producto 6\ensuremath{\times}6 \ensuremath{\cdot} 6 sin alocación adicional.

\subsection{11. Tangente algorítmica consistente}

\[
\mathbf C^{\text{alg}} = \mathbf C_e
\]

Simétrica (\texttt{IS\_SYMMETRIC = True}). Positiva definida para $\nu \in (-1, 0.5)$ estricto; diverge en $\nu \to 0.5$ (incompresible --- requiere formulación mixta u-p, no soportada).

\subsection{12. Caveats numéricos}

\begin{itemize}
  \item \textbf{Incompresibilidad $\nu \to 0.5$}: la matriz constitutiva diverge ($1-2\nu \to 0$). El constructor rechaza $\nu \ge 0.5$ con \texttt{ValueError}. Recomendado $\nu \le 0.49$ para metales y polímeros casi-incompresibles.
  \item \textbf{Locking volumétrico} en \texttt{Hex8}/\texttt{Tet4} con $\nu \to 0.5$: declarado limitación arquitectural en la spec de cada elemento, sin mitigación implementada (B-bar/F-bar diferidas a caso de uso real, política idéntica al 2D).
  \item Para $\nu < 0$ el material es admisible termodinámicamente (auxético) pero infrecuente en la práctica; el constructor lo permite hasta $\nu > -1$.
\end{itemize}


\section{Contrato de implementación}

\begin{lstlisting}[language=yaml]
name: Elastic3D
kind: material
status: validated

interface:
  strain_dim: 6
  primary_state_var: null
  is_symmetric: true

parameters:
  - { name: E,       type: float, required: true,  desc: "Módulo de Young (>0)" }
  - { name: nu,      type: float, required: true,  desc: "Coeficiente de Poisson ∈ (-1, 0.5)" }
  - { name: density, type: float, required: false, default: null,
      desc: "Densidad (kg/m³). Opcional al construir; obligatoria solo si se ensambla peso propio o masa (ADR 0008)" }

signature:
  compute_state: "(ε: ndarray(6,), state_vars=None) -> (σ: ndarray(6,), C_e: ndarray(6,6), state_vars)"
  strain_kind: "3D Voigt — [ε_xx, ε_yy, ε_zz, γ_xy, γ_yz, γ_xz] con γ_ij = 2·ε_ij"

state_variables: []

conventions:
  sign: "σ > 0 ⇔ tracción (Reglas §5)"
  voigt: "[xx, yy, zz, xy, yz, xz] (ADR 0012); γ_ij = 2·ε_ij en deformación; σ_ij tensorial sin factor"
  hypothesis_dispatch: "no aplica en 3D"

validity:
  - "E > 0"
  - "ν ∈ (-1, 0.5) estricto; ν = 0.5 rechazado (incompresible)"
  - "pequeñas deformaciones (|ε| ≲ 1e-2)"

out_of_scope:
  - "anisotropía (ortótropo, monoclínico, totalmente anisótropo)"
  - "plasticidad ⇒ VonMises3D, DruckerPrager3D (sub-etapa posterior)"
  - "daño ⇒ IsotropicDamage3D (sub-etapa posterior)"
  - "incompresibilidad estricta (ν=0.5) ⇒ requiere formulación mixta u-p"
  - "grandes deformaciones"

acceptance:
  verification:
    - name: simetria_C_e
      setup: "construir Elastic3D y comparar C_e con su traspuesta"
      expect: "C_e == C_e.T exacto componente a componente"
      tol_abs: 1.0e-14
    - name: positividad_definida
      setup: "autovalores de C_e con ν ∈ {0, 0.3, 0.49}"
      expect: "todos los autovalores estrictamente positivos"
      tol_abs: 0.0
    - name: traccion_uniaxial
      setup: "ε = (ε_0, -ν·ε_0, -ν·ε_0, 0, 0, 0)"
      expect: "σ = (E·ε_0, 0, 0, 0, 0, 0) exacto"
      tol_rel: 1.0e-12
    - name: cortante_puro_xy
      setup: "ε = (0, 0, 0, γ_0, 0, 0) con γ_0 = 1e-3"
      expect: "σ = (0, 0, 0, G·γ_0, 0, 0) con G = E/[2(1+ν)]"
      tol_rel: 1.0e-12
    - name: cortante_puro_yz
      setup: "ε = (0, 0, 0, 0, γ_0, 0)"
      expect: "σ = (0, 0, 0, 0, G·γ_0, 0)"
      tol_rel: 1.0e-12
    - name: cortante_puro_xz
      setup: "ε = (0, 0, 0, 0, 0, γ_0)"
      expect: "σ = (0, 0, 0, 0, 0, G·γ_0)"
      tol_rel: 1.0e-12
    - name: compresion_hidrostatica
      setup: "ε = (ε_0, ε_0, ε_0, 0, 0, 0) con ε_0 = -1e-3 (compresión isótropa)"
      expect: "σ = (K·3ε_0, K·3ε_0, K·3ε_0, 0, 0, 0) con K = E/[3(1-2ν)] (módulo de compresibilidad)"
      tol_rel: 1.0e-12
    - name: rechazo_inputs_invalidos
      setup: "construir Elastic3D con E≤0, ν fuera de rango (-1, 0.5), density<0"
      expect: "ValueError con mensaje claro en cada caso"

references:
  - "Timoshenko S., Goodier J.N. (1970). Theory of Elasticity. McGraw-Hill. §6 (3D elasticity)."
  - "Bathe K.J. (2014). Finite Element Procedures. Prentice Hall. §6.6."
  - "Cook R.D., Malkus D.S., Plesha M.E., Witt R.J. (2002). Concepts and Applications of FEA. Wiley. §6.8."
  - "ADR 0012 — Sólidos 3D: convención Voigt 6D y cierre del contrato internal_forces."

\end{lstlisting}


\section{Implementación}

\begin{itemize}
  \item \textbf{Archivo}: solidum/materials/elastic\_3d.py.
  \item \textbf{Clase}: \texttt{Elastic3D}, registrada vía \texttt{@MaterialRegistry.register}.
  \item \textbf{Despacho}: matriz \texttt{C} precalculada en el constructor. \texttt{compute\_state} reduce a \texttt{self.C @ strain} sin ramificaciones --- coste mínimo.
  \item \textbf{Validación temprana} en constructor: \texttt{E > 0}, \texttt{\ensuremath{\nu} \ensuremath{\in} (-1, 0.5)}, \texttt{density \ensuremath{\ge} 0} si se pasa. Mensajes en español.
  \item \textbf{Tests}: tests/test\_solid\_3d.py --- clase \texttt{TestElastic3D} (11 tests: simetría C, positividad definida, tracción uniaxial, los tres cortantes puros, compresión hidrostática con K = E/[3(1-2\ensuremath{\nu})], rechazos de inputs inválidos, density opcional).
  \item \textbf{Validación de integración}: \texttt{tests/validation/test\_cube\_lame\_3d.py} consume Elastic3D end-to-end vía Hex8 y Tet4 (tracción uniaxial + compresión hidrostática con solución analítica exacta).
\end{itemize}


\section{Diálogo}

\begin{itemize}
  \item \textbf{2026-05-19} \ensuremath{\cdot} Spec creada como pieza inicial de la Etapa 7 (sólidos 3D, alcance acotado Hex8 + Tet4 + Elastic3D). Convención Voigt 6D fijada en ADR 0012. Sin variantes de hipótesis (no aplica plane\_stress/plane\_strain en 3D).
\end{itemize}


\newpage

\section{IsotropicDamage3D}



\textit{Orden de trabajo. El usuario revisa \textbf{especificación física}, \textbf{formulación numérica} y \textbf{contrato}; la IA propone, ejecuta y rellena \textbf{implementación} + \textbf{diálogo}.}

\textit{> \textbf{Pre-redactada por la IA} sobre la base de \texttt{IsotropicDamage2D.md} (mismo modelo físico) y \texttt{Elastic3D.md} / \texttt{VonMises3D.md} (patrón Voigt 6D). Tercera y última entrega de la sub-etapa \textbf{A.bis} (materiales 3D no lineales).}



\section{Especificación física}

\subsection{0. Descripción general}

Modelo de \textbf{daño isótropo escalar} en pequeñas deformaciones para sólidos 3D, formulado íntegramente en Voigt 6D del proyecto (ADR 0012). Captura degradación progresiva e isótropa de la rigidez en respuesta a la deformación máxima histórica, con ablandamiento exponencial. Pensado para hormigón en tracción uniforme, cerámicos, materiales cuasi-frágiles donde la microfisuración distribuida puede modelarse mediante un escalar.

Diferencias respecto a \texttt{IsotropicDamage2D}:

\begin{itemize}
  \item \textbf{Sin variantes de hipótesis} (igual que VM3D, DP3D, Elastic3D): en 3D todas las componentes son activas.
  \item \textbf{Algoritmo idéntico al 2D extendido a 6D}: misma ley de evolución exponencial, mismo régimen carga/descarga por $\kappa$, misma estructura de tangente algorítmica consistente. \textbf{Sin Newton local} (el algoritmo es explícito).
  \item \textbf{\texttt{IsotropicDamage2D} plane\_strain} es la restricción de este modelo bajo $\varepsilon_{zz} = \gamma_{yz} = \gamma_{xz} = 0$ con la misma fórmula de $\varepsilon_\text{eq}$ --- el test cruzado se hace contra esta variante (no contra plane\_stress, que es genuinamente distinta).
\end{itemize}

\subsection{1. Cinemática}

Pequeñas deformaciones, descomposición elástica pura: no hay deformación plástica (sin plasticidad). El daño se manifiesta como degradación de la rigidez efectiva. Sin variables internas tensoriales --- solo el historial escalar $\kappa$ y el daño $d$.

\subsection{2. Esfuerzo nominal vs esfuerzo efectivo}

\[
\boldsymbol\sigma = (1 - d)\,\boldsymbol\sigma^\text{eff}, \qquad \boldsymbol\sigma^\text{eff} = \mathbf C_e\,\boldsymbol\varepsilon
\]

con $d \in [0, d_\text{max}]$ la variable de daño isótropa y $\mathbf C_e$ el tensor constitutivo elástico isótropo 3D 6\ensuremath{\times}6 de \texttt{Elastic3D} (sin variantes de hipótesis en 3D).

El esfuerzo efectivo $\boldsymbol\sigma^\text{eff}$ es el que el material intacto soportaría con la deformación corriente. El esfuerzo nominal $\boldsymbol\sigma$ es lo que efectivamente transmite considerando la degradación.

\subsection{3. Deformación equivalente}

Cantidad escalar que cuantifica la "intensidad" de la deformación. Convención del proyecto, norma simétrica simple en notación Voigt 6D $[\varepsilon_{xx}, \varepsilon_{yy}, \varepsilon_{zz}, \gamma_{xy}, \gamma_{yz}, \gamma_{xz}]$ con $\gamma_{ij} = 2\varepsilon_{ij}$:

\[
\varepsilon_\text{eq}(\boldsymbol\varepsilon) = \sqrt{\varepsilon_{xx}^2 + \varepsilon_{yy}^2 + \varepsilon_{zz}^2 + \tfrac{1}{2}(\gamma_{xy}^2 + \gamma_{yz}^2 + \gamma_{xz}^2)} = \sqrt{\boldsymbol\varepsilon^\top \mathbf M\,\boldsymbol\varepsilon}
\]

con $\mathbf M = \operatorname{diag}(1, 1, 1, 1/2, 1/2, 1/2)$. Equivalente a la norma de Frobenius del tensor deformación 3D ($\boldsymbol\varepsilon : \boldsymbol\varepsilon$ con $\varepsilon_{ij} = \gamma_{ij}/2$):

\[
\boldsymbol\varepsilon:\boldsymbol\varepsilon = \sum_i \varepsilon_{ii}^2 + 2\sum_{i<j}\varepsilon_{ij}^2 = \varepsilon_{xx}^2 + \varepsilon_{yy}^2 + \varepsilon_{zz}^2 + \tfrac{1}{2}(\gamma_{xy}^2 + \gamma_{yz}^2 + \gamma_{xz}^2)
\]

Es \textbf{extensión natural} de la convención 2D del proyecto ($\mathbf M_\text{2D} = \operatorname{diag}(1, 1, 1/2)$). Bajo restricción plane strain ($\varepsilon_{zz} = \gamma_{yz} = \gamma_{xz} = 0$), $\varepsilon_\text{eq}^{3D}$ se reduce exactamente a $\varepsilon_\text{eq}^{2D}$ --- propiedad clave para el test cruzado con \texttt{IsotropicDamage2D} plane\_strain (§13).

\textbf{Limitación}: no distingue tracción de compresión. Para hormigón a compresión con resistencia distinta se necesitaría un split tipo Mazars o de Vree; explícitamente \textit{out-of-scope} (idéntica deuda gemela al 2D).

\subsection{4. Variable histórica y condición de carga (Kuhn-Tucker)}

\[
\kappa_{n+1} = \max(\kappa_n, \varepsilon_\text{eq}(\boldsymbol\varepsilon_{n+1})), \qquad \kappa_0 \text{ inicial}
\]

Esto codifica:

\begin{itemize}
  \item \textbf{Carga activa} ($\dot\kappa > 0$): $\kappa$ crece con $\varepsilon_\text{eq}$; el daño puede aumentar.
  \item \textbf{Descarga / recarga} ($\dot\kappa = 0$): $\kappa$ queda congelado; el daño no cambia.
  \item \textbf{Irreversibilidad}: $\kappa$ es monótono no decreciente.
\end{itemize}

\subsection{5. Ley de evolución del daño}

\[
d(\kappa) = \begin{cases}
0 & \kappa \leq \kappa_0 \\
1 - \dfrac{\kappa_0}{\kappa}\,e^{-\alpha(\kappa - \kappa_0)} & \kappa > \kappa_0
\end{cases}
\]

con saturación numérica $d \leq d_\text{max} = $ \texttt{DAMAGE\_MAX} (evita rigidez nula).

Propiedades idénticas al 2D: $d(\kappa_0) = 0$ (continuidad); $d'(\kappa) > 0$ para $\kappa > \kappa_0$ (irreversibilidad); $\lim_{\kappa \to \infty} d = 1$ (asintótico); $\alpha$ controla la velocidad de degradación. \textbf{Misma función de ablandamiento} que en 1D y 2D --- centralizada en \texttt{solidum.materials.\_softening.evaluate\_exponential\_damage}.


\section{Formulación numérica}

\subsection{6. Algoritmo en un paso}

Dado $\{\kappa_n\}$ y $\boldsymbol\varepsilon_{n+1}$:

\begin{enumerate}
  \item $\varepsilon_\text{eq} = \sqrt{\boldsymbol\varepsilon^\top \mathbf M\,\boldsymbol\varepsilon}$ (norma 6D).
  \item Si $\varepsilon_\text{eq} > \kappa_n$: \textbf{carga activa}, $\kappa_{n+1} = \varepsilon_\text{eq}$. Si no: \textbf{descarga}, $\kappa_{n+1} = \kappa_n$.
  \item Aplicar ley de daño con $\kappa_{n+1}$; saturar a $d_\text{max}$.
  \item $\boldsymbol\sigma_{n+1} = (1 - d)\,\mathbf C_e\,\boldsymbol\varepsilon_{n+1}$.
  \item Tangente: \textbf{secante} en descarga / sin daño / saturación; \textbf{algorítmica consistente} en carga activa con daño no saturado (ver §7).
\end{enumerate}

El algoritmo es \textbf{explícito}: no requiere Newton local. La actualización de $\kappa$ y $d$ es directa, idéntica al 2D.

\subsection{7. Tangente algorítmica consistente}

\[
\mathbf C^\text{alg} = \frac{\partial \boldsymbol\sigma_{n+1}}{\partial \boldsymbol\varepsilon_{n+1}}
\]

Derivando $\boldsymbol\sigma = (1-d)\,\mathbf C_e\,\boldsymbol\varepsilon$:

\[
\mathbf C^\text{alg} = (1-d)\,\mathbf C_e - \boldsymbol\sigma^\text{eff} \otimes \frac{\partial d}{\partial \boldsymbol\varepsilon}
\]

con $\boldsymbol\sigma^\text{eff} = \mathbf C_e\,\boldsymbol\varepsilon$ (vector 6\ensuremath{\times}1) y la derivada del daño:

\[
\frac{\partial d}{\partial \boldsymbol\varepsilon} = \frac{\partial d}{\partial \kappa}\,\frac{\partial \kappa}{\partial \boldsymbol\varepsilon}
\]

donde:

\begin{itemize}
  \item $\dfrac{\partial d}{\partial \kappa} = (1 - d)\left(\dfrac{1}{\kappa} + \alpha\right)$ para la ley exponencial (idéntica al 2D, derivación independiente de la dimensión).
\end{itemize}

\begin{itemize}
  \item $\dfrac{\partial \kappa}{\partial \boldsymbol\varepsilon}$ depende del régimen:
  \item \textbf{Carga activa} ($\varepsilon_\text{eq} > \kappa_n$, $\kappa = \varepsilon_\text{eq}$):
\end{itemize}
    \[
\dfrac{\partial \kappa}{\partial \boldsymbol\varepsilon} = \dfrac{\partial \varepsilon_\text{eq}}{\partial \boldsymbol\varepsilon} = \dfrac{1}{\varepsilon_\text{eq}}\,\mathbf M\,\boldsymbol\varepsilon = \dfrac{1}{\varepsilon_\text{eq}}\begin{bmatrix}\varepsilon_{xx} \\ \varepsilon_{yy} \\ \varepsilon_{zz} \\ \tfrac{1}{2}\gamma_{xy} \\ \tfrac{1}{2}\gamma_{yz} \\ \tfrac{1}{2}\gamma_{xz}\end{bmatrix}
\]
\begin{itemize}
  \item \textbf{Descarga} ($\varepsilon_\text{eq} \leq \kappa_n$): $\partial \kappa / \partial \boldsymbol\varepsilon = 0$.
\end{itemize}

\textbf{Tangente final}:

\[
\mathbf C^\text{alg} = \begin{cases}
\mathbf C_e & \kappa \leq \kappa_0 \;(\text{sin daño}) \\
(1-d)\,\mathbf C_e & \text{descarga} \\
(1-d)\,\mathbf C_e \;-\; \dfrac{(1-d)(1/\kappa + \alpha)}{\varepsilon_\text{eq}}\,(\mathbf C_e\boldsymbol\varepsilon)\otimes(\mathbf M\boldsymbol\varepsilon) & \text{carga activa, } d < d_\text{max} \\
(1-d_\text{max})\,\mathbf C_e & d = d_\text{max} \;(\text{saturación})
\end{cases}
\]

con $\mathbf M = \operatorname{diag}(1, 1, 1, 1/2, 1/2, 1/2)$ (6\ensuremath{\times}6).

\subsection{8. Pérdida de simetría}

El producto externo $(\mathbf C_e\boldsymbol\varepsilon)\otimes(\mathbf M\boldsymbol\varepsilon)$ es una matriz $6\times 6$ que \textbf{no es simétrica en general}. $\mathbf C_e\boldsymbol\varepsilon$ y $\mathbf M\boldsymbol\varepsilon$ no son proporcionales salvo en estados de deformación muy particulares (e.g. uniaxial puro alineado con ejes principales).

Consecuencia: \textbf{\texttt{IS\_SYMMETRIC = False} obligatorio}. El despachador algebraico (ADR 0003) elige LU en lugar de Cholesky/LDL\ensuremath{^{\top}} para sistemas globales que incluyan este material. Coste por paso \textasciitilde{}30-50\% mayor, compensado con creces por la convergencia cuadrática del Newton global (típicamente 2-3 iteraciones en lugar de 5-10 con tangente secante).

\subsection{9. Decisión de régimen durante el Newton}

Idéntica al 2D: la transición loading\ensuremath{\leftrightarrow}unloading se decide con $\kappa_n$ committed (no con $\kappa$ trial dentro del Newton). Estándar y evita oscilaciones del Newton entre ramas en pasos con ambigüedad. Si el paso converge, $\kappa$ se compromete y la decisión es consistente con la trayectoria.


\section{Contrato de implementación}

\begin{lstlisting}[language=yaml]
name: IsotropicDamage3D
kind: material
status: validated        # draft → implemented → validated

interface:
  strain_dim: 6
  primary_state_var: damage
  is_symmetric: false      # tangente consistente no simétrica (ver §8)

parameters:
  - { name: E,       type: float, required: true,  desc: "Módulo de Young intacto (>0)" }
  - { name: nu,      type: float, required: true,  desc: "Coeficiente de Poisson ∈ (-1, 0.5)" }
  - { name: kappa_0, type: float, required: true,  desc: "Umbral de deformación equivalente; daño nulo mientras κ ≤ κ_0 (>0)" }
  - { name: alpha,   type: float, required: true,  desc: "Velocidad de degradación exponencial (>0). α mayor → ablandamiento más abrupto" }
  - { name: density, type: float, required: false, default: null, desc: "Densidad (ADR 0008)" }

signature:
  compute_state: "(ε: ndarray(6,), state_vars=None) -> (σ: ndarray(6,), C_tan: ndarray(6,6), state_vars')"
  strain_kind: "3D Voigt — [ε_xx, ε_yy, ε_zz, γ_xy, γ_yz, γ_xz] con γ_ij = 2·ε_ij"

state_schema:
  kappa:  "float ≥ κ_0 — máxima deformación equivalente histórica"
  damage: "float ∈ [0, DAMAGE_MAX] — variable de daño escalar isótropa"

conventions:
  sign:  "σ > 0 ⇔ tracción; el modelo no distingue tracción de compresión en la activación del daño"
  voigt: "γ_ij = 2·ε_ij en deformación; ε_eq usa M = diag(1, 1, 1, 1/2, 1/2, 1/2)"
  damage_irreversibility: "κ monótono no decreciente vía max(κ_old, ε_eq); d(κ) creciente"

validity:
  - "E > 0, κ_0 > 0, α > 0"
  - "ν ∈ (-1, 0.5) estricto"

out_of_scope:
  - "split tracción/compresión (Mazars, de Vree, Mises modificado); el modelo daña por igual en cualquier régimen"
  - "anisotropía del daño (tensor de daño D vs escalar)"
  - "regularización para evitar mesh-dependency en regime de ablandamiento (longitud característica, no-local, gradient)"
  - "acoplamiento con plasticidad"
  - "fatiga / acumulación cíclica del daño"

numerical_caveats:
  - "Tangente consistente NO simétrica → IS_SYMMETRIC=False obliga al despachador algebraico a usar LU (ADR 0003). Costo por paso ~30-50% mayor que Cholesky pero convergencia cuadrática del Newton global lo compensa (Simó-Ju 1987)."
  - "Transición loading↔unloading discontinua en la tangente. Newton converge antes de oscilar entre ramas en casos típicos. Para problemas patológicos cerca del umbral, considerar arc-length."
  - "Saturación d=DAMAGE_MAX corta la corrección consistente; el algoritmo devuelve tangente secante en ese caso para evitar términos espurios."
  - "Sin regularización, la solución es mesh-dependent en régimen de ablandamiento (localización en una banda de elementos). Es limitación inherente del modelo continuo de daño isótropo escalar; out-of-scope esta entrega."

acceptance:
  unit:
    - name: paso_elastico_bajo_umbral
      setup: "ε pequeño 6D con todas las componentes activas, ε_eq < κ_0"
      expect: "σ = C_e·ε exacto; d=0; C_tan = C_e (sin corrección rank-1)"
      tol_rel: 1.0e-12

    - name: damage_evolution_exponential
      setup: "ε grande tal que ε_eq >> κ_0"
      expect: "d cumple la ley exponencial; estado coincide con la fórmula directa"
      tol_rel: 1.0e-10

    - name: rejects_loading_below_threshold
      setup: "ε_eq apenas por encima de κ_0 (banda elástica + bandera de carga activa)"
      expect: "d > 0 minúsculo; κ_new = ε_eq; tangente con corrección rank-1"

    - name: tangent_equals_secant_on_unloading
      setup: "carga hasta κ > κ_0; descarga con ε de menor norma"
      expect: "C_tan = (1-d)·C_e exacto (segundo término se anula con ∂κ/∂ε = 0)"
      tol_rel: 1.0e-12

    - name: tangent_equals_secant_below_threshold
      setup: "ε pequeño, κ_old = κ_0 (estado inicial intacto)"
      expect: "C_tan = C_e (sin daño activo)"
      tol_rel: 1.0e-12

    - name: tangent_equals_secant_at_saturation
      setup: "carga hasta κ enorme tal que d alcanza DAMAGE_MAX"
      expect: "C_tan = (1-DAMAGE_MAX)·C_e (el término consistente se corta)"
      tol_rel: 1.0e-12

    - name: tangent_not_symmetric_in_loading
      setup: "carga activa con estado de deformación no degenerado (ε_xx ≠ ε_yy ≠ ε_zz, cortantes activos)"
      expect: "‖C_tan - C_tan^T‖_F / ‖C_tan‖_F ≥ 0.01 (asimetría significativa)"

    - name: tangent_finite_difference_consistency
      setup: "carga activa; comparar C_alg cerrada con diferencia finita centrada"
      expect: "‖C_FD − C_alg‖ / ‖C_alg‖ < 1e-5"
      tol_rel: 1.0e-5

    - name: kappa_monotonic
      setup: "trayectoria multi-paso con carga, descarga, recarga"
      expect: "κ es monótono no decreciente en todos los pasos"

    - name: descarga_recarga_irreversibilidad
      setup: "cargar a κ_load, descargar elásticamente, recargar"
      expect: "d permanece constante durante toda la descarga; al recargar, d sigue creciendo desde el valor de descarga"

    - name: unit_invariance_MPa_vs_Pa
      setup: "mismo path adimensional con (MPa,mm,N) y (Pa,m,N)"
      expect: "d, κ y σ/E idénticos (κ_0 es una deformación, ε es adimensional → invariancia exacta)"
      tol_rel: 1.0e-12

    - name: rechazo_inputs_invalidos
      setup: "construir con E≤0, κ_0≤0, α≤0, ν∉(-1, 0.5), density<0"
      expect: "ValueError con mensaje claro en cada caso"

  cross_consistency:
    - name: equivalencia_plane_strain
      setup: "Damage3D con ε = (ε_xx, ε_yy, 0, γ_xy, 0, 0); Damage2D plane_strain con ε = (ε_xx, ε_yy, γ_xy);
             path multi-paso con carga, descarga, recarga"
      expect: "σ_xx, σ_yy, σ_xy, κ, d coinciden exactamente entre 3D y 2D PS;
              σ_zz en VM3D no nulo (consistente con C_e plane_strain interno);
              componentes 3D ausentes en 2D PS (σ_yz, σ_xz) permanecen nulas"
      tol_rel: 1.0e-12

  integration:
    - name: hex8_elastico
      setup: "Hex8 unitario con carga uniforme muy por debajo del umbral"
      expect: "d = 0 en todos los Gauss points; σ = C_e·ε exacto"
      tol_rel: 1.0e-8

    - name: hex8_damage_newton_converges_in_few_iter
      setup: "Hex8 unitario, carga uniaxial hasta régimen plenamente dañado (κ >> κ_0), 8 pasos"
      expect: "Newton converge en ≤6 iteraciones por paso (evidencia de convergencia cuadrática con tangente consistente)"

    - name: tet4_basic_pipeline
      setup: "Tet4 con confinamiento + tracción uniaxial, carga que activa daño"
      expect: "convergencia, d > 0 en el único Gauss point, κ ≥ κ_0"

references:
  - "Lemaitre J., Chaboche J.-L. (1990). Mechanics of Solid Materials. Cambridge UP. Cap. 7 (continuum damage mechanics)."
  - "Simó J.C., Ju J.W. (1987). Strain- and stress-based continuum damage models — I. Formulation. Int. J. Solids Struct. 23, 821-840 (tangente consistente para daño isótropo)."
  - "Oliver J. (1989). A consistent characteristic length for smeared cracking models. IJNME 28, 461-474."
  - "de Souza Neto E.A., Perić D., Owen D.R.J. (2008). Computational Methods for Plasticity. Wiley. §12 (damage models)."
  - "ADR 0012 — Sólidos 3D: convención Voigt 6D del proyecto."
  - "ADR 0003 — Despachador algebraico: tangente asimétrica fuerza LU global."
  - "ADR 0006 — Tolerancias adimensionales: patrón atol + rtol·escala con escala física E·κ_0."

\end{lstlisting}


\section{Implementación}

\begin{itemize}
  \item \textbf{Archivo}: solidum/materials/damage\_3d.py.
  \item \textbf{Clase}: \texttt{IsotropicDamage3D}, registrada vía \texttt{@MaterialRegistry.register}. Declara \texttt{IS\_SYMMETRIC = False} como ClassVar (tangente consistente asimétrica).
  \item \textbf{Composición}:
  \item \texttt{Elastic3D} como base elástica (acceso a $\mathbf C_e$ 6\ensuremath{\times}6 vía \texttt{self.elastic\_base.C}).
  \item \texttt{evaluate\_exponential\_damage} de \texttt{solidum.materials.\_softening} para la ley de softening (centralizada, compartida con \texttt{IsotropicDamage1D} e \texttt{IsotropicDamage2D}).
  \item \textbf{Estado}: dict \texttt{\{'kappa': float, 'damage': float\}}. \texttt{PRIMARY\_STATE\_VAR = 'damage'}.
  \item \textbf{Algoritmo \texttt{compute\_state}}: idéntica estructura a \texttt{IsotropicDamage2D} extendida a Voigt 6D. Sin Numba (operaciones numpy puras; el coste dominante es el producto matriz-vector $\mathbf C_e \cdot \boldsymbol\varepsilon$, no se beneficia significativamente de JIT).
  \item \textbf{\texttt{admissibility\_scale}}: \texttt{E\ensuremath{\cdot}\ensuremath{\kappa}\_0} constante respecto al estado (idéntico al patrón 1D/2D --- el criterio se mide contra \texttt{\ensuremath{\kappa}\_0}, no contra \texttt{\ensuremath{\kappa}} corriente).
  \item \textbf{Tests}:
  \item tests/test\_materials\_unit.py \ensuremath{\cdot} \texttt{TestIsotropicDamage3D} (12 tests unitarios) + \texttt{TestIsotropicDamage3DvsPlaneStrain} (1 cruzado contra Damage2D plane\_strain).
  \item tests/test\_solid\_3d\_plasticity.py \ensuremath{\cdot} 3 tests de integración: \texttt{TestHex8IsotropicDamage3DElastic} (1 --- sin daño), \texttt{TestHex8IsotropicDamage3DActiveLoad} (1 --- daño activo con tangente asimétrica), \texttt{TestTet4IsotropicDamage3DBasic} (1 --- smoke Tet4).
\end{itemize}

\begin{itemize}
  \item \textbf{Notas de traducción}:
  \item Matriz $\mathbf M = \operatorname{diag}(1, 1, 1, 1/2, 1/2, 1/2)$ aplicada implícitamente en código (factor $1/2$ sobre los slots cortantes de $\boldsymbol\varepsilon$) en lugar de allocar la 6\ensuremath{\times}6, idéntico patrón al 2D.
  \item Convención engineering input (\ensuremath{\gamma}\_ij): el factor $1/2$ en $\mathbf M$ y en $\partial\varepsilon_\text{eq}/\partial\boldsymbol\varepsilon$ corrige automáticamente para que la norma corresponda a la Frobenius del tensor (\ensuremath{\gamma}\_ij/2 = \ensuremath{\varepsilon}\_ij\_tens).
  \item Validación de parámetros en constructor: \texttt{E > 0}, \texttt{\ensuremath{\kappa}\_0 > 0}, \texttt{\ensuremath{\alpha} > 0}, \texttt{\ensuremath{\nu} \ensuremath{\in} (-1, 0.5)}, \texttt{density \ensuremath{\ge} 0} si se pasa. Mensajes en español.
\end{itemize}


\section{Diálogo}

\begin{itemize}
  \item \textbf{2026-05-21} \ensuremath{\cdot} Spec pre-redactada por la IA al cerrar el segundo material de A.bis (\texttt{DruckerPrager3D}). Base: \texttt{IsotropicDamage2D.md} (modelo físico, tangente algorítmica consistente, semánticas carga/descarga, cap DAMAGE\_MAX) + \texttt{VonMises3D.md}/\texttt{DruckerPrager3D.md} (estructura Voigt 6D, patrón cross-consistency).
\end{itemize}

\begin{itemize}
  \item \textbf{2026-05-21} \ensuremath{\cdot} \textbf{Decisión de plumbing (decide la IA, Reglas §3)}: reuso de la función centralizada \texttt{evaluate\_exponential\_damage} de \texttt{solidum.materials.\_softening}. Esta función ya está compartida entre \texttt{IsotropicDamage1D} e \texttt{IsotropicDamage2D} (auditoría H-3.4 cerrada el 2026-05-19); añadir \texttt{IsotropicDamage3D} como tercer consumidor satisface la regla de centralización del proyecto. Sin duplicación de la ley de softening.
\end{itemize}

\begin{itemize}
  \item \textbf{2026-05-21} \ensuremath{\cdot} \textbf{Decisión de plumbing}: el material usa \texttt{Elastic3D} internamente como base elástica (paralelo a cómo \texttt{IsotropicDamage2D} usa \texttt{Elastic2D}). Acceso a la matriz constitutiva 6\ensuremath{\times}6 vía \texttt{self.elastic\_base.C} --- idéntico patrón al 2D, sin duplicación de la construcción de $\mathbf C_e$.
\end{itemize}

\begin{itemize}
  \item \textbf{2026-05-21} \ensuremath{\cdot} \textbf{Decisión de plumbing}: la matriz $\mathbf M$ del cálculo de $\varepsilon_\text{eq}$ y de la derivada $\partial\kappa/\partial\boldsymbol\varepsilon$ se materializa implícitamente en código (multiplicaciones por $1/2$ sobre los slots cortantes) en lugar de allocar una 6\ensuremath{\times}6 --- ahorro mínimo de memoria pero código más legible y consistente con el patrón seguido en \texttt{damage\_2d.py}.
\end{itemize}

\begin{itemize}
  \item \textbf{2026-05-21} \ensuremath{\cdot} \textbf{Cross-consistency} contra \texttt{IsotropicDamage2D} plane\_strain, no plane\_stress. La equivalencia matemática es exacta bajo $\varepsilon_{zz} = \gamma_{yz} = \gamma_{xz} = 0$ impuesto (plane strain): la matriz $\mathbf C_e$ 3D restringida a \ensuremath{\varepsilon}\_zz=0 input se reduce a la matriz $\mathbf C_e^\text{plane strain}$ 2D, y $\varepsilon_\text{eq}^\text{3D}$ se reduce a $\varepsilon_\text{eq}^\text{2D}$ bajo la misma restricción. Plane stress 2D usa un $\mathbf C_e$ proyectado distinto y por tanto no es comparable directamente con 3D bajo restricción cinemática.
\end{itemize}

\begin{itemize}
  \item \textbf{2026-05-21} \ensuremath{\cdot} \textbf{Sin variantes de hipótesis} en 3D (igual que VM3D, DP3D, Elastic3D). Constructor más simple que el 2D.
\end{itemize}

\begin{itemize}
  \item \textbf{2026-05-21} \ensuremath{\cdot} \textbf{No se incluye consistent tangent para \texttt{IsotropicDamage1D}} (deuda gemela del 2D): fuera del alcance de esta entrega A.bis. La fórmula 1D es trivial reducción de la 3D si se prioriza después.
\end{itemize}

\begin{itemize}
  \item \textbf{2026-05-21} \ensuremath{\cdot} \textbf{Implementación completa, status \texttt{validated}}. Clase + 12 tests unitarios + 1 test cruzado + 3 tests de integración (Hex8 elástico, Hex8 con daño activo + tangente asimétrica, Tet4 smoke). \textbf{16 tests añadidos a la suite (848 \ensuremath{\to} 866 verdes; suite global sin regresiones).} Todos los tests pasaron a la primera, incluyendo:
  \item Cruzado Damage3D \ensuremath{\leftrightarrow} Damage2D plane\_strain a \textbf{14 decimales en \texttt{d} y \texttt{\ensuremath{\kappa}}} y 10 decimales en \ensuremath{\sigma} --- la equivalencia bajo $\varepsilon_{zz} = \gamma_{yz} = \gamma_{xz} = 0$ es exacta a precisión máquina (misma ley centralizada, mismo $\mathbf C_e$ plane\_strain bajo restricción cinemática).
  \item Diferencia finita de la tangente algorítmica vs cerrada con error rel. < 1e-5 en las 36 componentes 6\ensuremath{\times}6 en carga activa.
  \item Asimetría significativa de la tangente en estados no degenerados (norma rel. \ensuremath{\ge} 0.01).
  \item Invariancia bajo cambio de unidades (E en MPa vs Pa) --- \texttt{d}, \texttt{\ensuremath{\kappa}} idénticos a 12 decimales; \ensuremath{\sigma} escala linealmente con E.
\end{itemize}

  \textbf{Con esta entrega cierra la sub-etapa A.bis} (materiales 3D no lineales): VonMises3D + DruckerPrager3D + IsotropicDamage3D. Validación contra benchmark publicado (e.g. ensayo tensión-softening con valor cuantitativo de carga máxima en barra cargada axialmente con softening exponencial calibrado por G\_f efectivo) \textbf{diferida a la campaña 3D consolidada} que se hará al cierre de A.bis junto con Hill 3D para VM3D y esfera hueca para DP3D.


\newpage

\section{VonMises3D}



\textit{Orden de trabajo. El usuario revisa \textbf{especificación física}, \textbf{formulación numérica} y \textbf{contrato}; la IA propone, ejecuta y rellena \textbf{implementación} + \textbf{diálogo} durante el trabajo.}

\textit{> \textbf{Pre-redactada por la IA} sobre la base de \texttt{VonMises2D.md} y \texttt{Elastic3D.md}. Apertura de la sub-etapa \textbf{A.bis} (materiales 3D no lineales) tras cierre de la Etapa 7 (sólidos 3D acotados con ADR 0012).}



\section{Especificación física}

\subsection{0. Descripción general}

Modelo elastoplástico independiente de la velocidad ("rate-independent") basado en el criterio \textbf{J2 / Von Mises} con regla de flujo \textbf{asociada} y \textbf{endurecimiento isótropo lineal}, formulado \textbf{íntegramente en 3D} sobre la convención Voigt 6D del proyecto (ADR 0012). Captura plasticidad de metales dúctiles en régimen de pequeñas deformaciones, sin efectos cinemáticos (Bauschinger), térmicos ni viscosos.

A diferencia de \texttt{VonMises2D}, \textbf{no tiene variantes de hipótesis cinemática}: en 3D todas las componentes de $\boldsymbol\varepsilon$ y $\boldsymbol\sigma$ son activas; no hay proyección a un subespacio. La consecuencia es algorítmica: el return mapping es \textbf{único y radial cerrado}, sin Newton local proyectado.

\texttt{VonMises2D} en hipótesis \textbf{plane strain} es matemáticamente la restricción de este modelo bajo $\varepsilon_{zz} = \gamma_{yz} = \gamma_{xz} = 0$ (deformación total impuesta); las componentes plásticas $\varepsilon^p_{zz}, \varepsilon^p_{yz}, \varepsilon^p_{xz}$ son libres y evolucionan por incompresibilidad/isotropía. Esta relación se verifica como test cruzado en §13.

\subsection{1. Descomposición aditiva}

En pequeñas deformaciones, el tensor de deformación se descompone aditivamente:

\[
\boldsymbol\varepsilon = \boldsymbol\varepsilon^e + \boldsymbol\varepsilon^p
\]

La parte elástica $\boldsymbol\varepsilon^e$ gobierna el esfuerzo vía la ley elástica isótropa; la parte plástica $\boldsymbol\varepsilon^p$ es \textbf{histórica} y evoluciona según la regla de flujo. La parte plástica es \textbf{incompresible}: $\operatorname{tr}(\boldsymbol\varepsilon^p) = 0$.

\subsection{2. Ley elástica}

\[
\boldsymbol\sigma = \mathbf C_e : \boldsymbol\varepsilon^e = \mathbf C_e : (\boldsymbol\varepsilon - \boldsymbol\varepsilon^p)
\]

con $\mathbf C_e$ la matriz constitutiva isótropa 3D de \texttt{Elastic3D} (Voigt 6\ensuremath{\times}6), parametrizada por $(E, \nu)$ o equivalentemente $(K, G)$:

\[
K = \frac{E}{3(1-2\nu)}, \qquad G = \frac{E}{2(1+\nu)}
\]

Equivalente desviador + presión:

\[
\boldsymbol\sigma = \mathbf s + p\,\mathbf I, \qquad \mathbf s = 2G\,\mathbf e^e, \qquad p = K\,\operatorname{tr}(\boldsymbol\varepsilon)
\]

con $\mathbf e^e = \operatorname{dev}(\boldsymbol\varepsilon^e)$ y, por incompresibilidad plástica, $\operatorname{tr}(\boldsymbol\varepsilon^e) = \operatorname{tr}(\boldsymbol\varepsilon)$.

\subsection{3. Superficie de fluencia (J2)}

\[
f(\boldsymbol\sigma, \alpha) = \|\mathbf s\| - \sqrt{\tfrac{2}{3}}\,\big(\sigma_y + H\,\alpha\big) \;\leq\; 0
\]

donde $\mathbf s = \boldsymbol\sigma - \tfrac{1}{3}\operatorname{tr}(\boldsymbol\sigma)\mathbf I$ es el tensor desviador, $\|\mathbf s\| = \sqrt{\mathbf s : \mathbf s}$ la norma de Frobenius tensorial, $\sigma_y > 0$ la fluencia inicial y $H \geq 0$ el módulo de endurecimiento isótropo lineal. El factor $\sqrt{2/3}$ alinea $\sigma_y$ con la fluencia uniaxial estándar.

\textbf{Cómputo de $\|\mathbf s\|$ en Voigt 6D del proyecto} (componentes cortantes tensoriales sin factor 2):

\[
\|\mathbf s\|^2 = s_{xx}^2 + s_{yy}^2 + s_{zz}^2 + 2\,(s_{xy}^2 + s_{yz}^2 + s_{xz}^2)
\]

El factor 2 contabiliza las dos posiciones simétricas $s_{ij} = s_{ji}$ del tensor de 2º orden.

\subsection{4. Regla de flujo (asociada)}

\[
\dot{\boldsymbol\varepsilon}^p = \dot\gamma\,\frac{\partial f}{\partial\boldsymbol\sigma} = \dot\gamma\,\mathbf N, \qquad \mathbf N = \frac{\mathbf s}{\|\mathbf s\|}
\]

El flujo es puramente desviador ($\operatorname{tr}(\mathbf N) = 0$) --- \textbf{plasticidad incompresible}, propiedad central de J2. $\mathbf N$ es un tensor de 2º orden de norma unidad almacenado en Voigt 6D con cortantes tensoriales.

\subsection{5. Endurecimiento isótropo lineal}

\[
\dot\alpha = \sqrt{\tfrac{2}{3}}\,\dot\gamma
\]

$\alpha$ es la \textbf{deformación plástica acumulada equivalente} (escalar, monótona no decreciente). La superficie de fluencia se expande uniformemente con $\alpha$ manteniendo su centro en el origen del espacio de tensiones desviadoras. La fórmula es \textbf{idéntica a plane strain} porque la dimensión de $\mathbf N$ es la misma (norma unidad tensorial); contrasta con plane stress, donde $\Delta\gamma$ tiene unidades de $[1/\text{esfuerzo}]$ por la proyección.

\subsection{6. Condiciones de Kuhn-Tucker}

\[
\dot\gamma \geq 0, \qquad f \leq 0, \qquad \dot\gamma\,f = 0
\]

Garantizan que $\dot\gamma > 0$ solo si el estado intenta salir de la superficie, y que el estado real siempre cumple $f \leq 0$.

\subsection{7. Variables internas}

\begin{center}
\begin{tabularx}{\textwidth}{YYY}
\hline
\textbf{Símbolo} & \textbf{Tipo} & \textbf{Significado} \\
\hline
$\boldsymbol\varepsilon^p$ & tensor (6 componentes Voigt, ver §8) & deformación plástica \\
$\alpha$ & escalar & deformación plástica acumulada equivalente \\
\hline
\end{tabularx}
\end{center}
\texttt{alpha} es la variable principal exportable (\texttt{PRIMARY\_STATE\_VAR = 'alpha'}); \texttt{eps\_p} es interna al algoritmo de return mapping.

\subsection{8. Notación Voigt y manejo de $\boldsymbol\varepsilon^p$}

Convención del proyecto para 3D (ADR 0012):

\[
\boldsymbol\varepsilon = [\varepsilon_{xx},\,\varepsilon_{yy},\,\varepsilon_{zz},\,\gamma_{xy},\,\gamma_{yz},\,\gamma_{xz}]^\top, \qquad \gamma_{ij} = 2\,\varepsilon_{ij}
\]

\[
\boldsymbol\sigma = [\sigma_{xx},\,\sigma_{yy},\,\sigma_{zz},\,\sigma_{xy},\,\sigma_{yz},\,\sigma_{xz}]^\top
\]

\textbf{Almacenamiento interno de $\boldsymbol\varepsilon^p$}: 6 componentes con cortantes \textbf{tensoriales sin factor 2}:

\[
\boldsymbol\varepsilon^p = [\varepsilon^p_{xx},\,\varepsilon^p_{yy},\,\varepsilon^p_{zz},\,\varepsilon^p_{xy},\,\varepsilon^p_{yz},\,\varepsilon^p_{xz}]
\]

Decisión consistente con \texttt{VonMises2D} (que almacena $\varepsilon^p_{xy}$ tensorial en su slot 4). Razón: el flujo plástico $\mathbf N = \mathbf s/\|\mathbf s\|$ es un tensor de 2º orden con cortantes tensoriales; almacenar $\boldsymbol\varepsilon^p$ en la misma convención evita conversiones repetidas en el return mapping y mantiene la fórmula $\|\boldsymbol\varepsilon^p\|^2 = \sum_\text{diag} + 2\sum_\text{off}$ idéntica a la de $\mathbf s$.

\textbf{Traducción al entrar al kernel}: la deformación total llega con $\gamma$ \textit{engineering}; las componentes cortantes se dividen por 2 una sola vez al construir el tensor para deviator decomposition. Al salir, $\boldsymbol\sigma$ se devuelve en Voigt 6D \textit{engineering-compatible} (cortantes tensoriales sin factor --- esto es lo que el ensamblador espera).

\subsection{9. Relación con \texttt{VonMises2D} plane strain}

\texttt{VonMises2D} plane strain es la restricción de este modelo bajo:

\[
\varepsilon_{zz} = \gamma_{yz} = \gamma_{xz} = 0 \quad\text{(deformación total impuesta por la hipótesis 2D)}
\]

con $\varepsilon^p_{zz} \neq 0$ libre (incompresibilidad plástica). El kernel 2D-plane-strain mantiene 4 componentes de $\boldsymbol\varepsilon^p$ porque $\varepsilon^p_{yz} = \varepsilon^p_{xz} = 0$ por isotropía + ausencia de cortantes activos en esos planos.

\textbf{Implicación de testeo}: si se ejecuta \texttt{VonMises3D} con $\varepsilon = [\varepsilon_{xx}, \varepsilon_{yy}, 0, \gamma_{xy}, 0, 0]$ (path plane strain), debe reproducir $\sigma_{xx}, \sigma_{yy}, \sigma_{xy}, \alpha$ de \texttt{VonMises2D} plane strain exactamente, además de devolver el $\sigma_{zz}$ que \texttt{VonMises2D} calcula internamente pero no expone en el API público 2D. Ver §13.


\section{Formulación numérica}

\subsection{10. Esquema temporal --- Backward Euler implícito}

Dado el estado convergido $\{\boldsymbol\varepsilon^p_n, \alpha_n\}$ al paso $n$ y la deformación total prescrita $\boldsymbol\varepsilon_{n+1}$, resolver para $\{\boldsymbol\sigma_{n+1}, \boldsymbol\varepsilon^p_{n+1}, \alpha_{n+1}\}$ con discretización implícita:

\[
\boldsymbol\varepsilon^p_{n+1} = \boldsymbol\varepsilon^p_n + \Delta\gamma\,\mathbf N_{n+1}, \qquad \alpha_{n+1} = \alpha_n + \sqrt{\tfrac{2}{3}}\,\Delta\gamma
\]

más $f(\boldsymbol\sigma_{n+1}, \alpha_{n+1}) \leq 0$ con la condición de complementariedad.

\subsection{11. Return mapping 3D --- algoritmo radial cerrado}

Descomposición volumétrica-desviadora del estado de prueba. Conversión de $\boldsymbol\varepsilon$ a tensorial dividiendo por 2 los cortantes engineering:

\[
\boldsymbol\varepsilon^\text{tens} = [\varepsilon_{xx},\,\varepsilon_{yy},\,\varepsilon_{zz},\,\tfrac{\gamma_{xy}}{2},\,\tfrac{\gamma_{yz}}{2},\,\tfrac{\gamma_{xz}}{2}]
\]

(esta es la conversión que ya hace \texttt{Elastic3D} implícitamente vía la matriz $\mathbf C_e$ con sus factores $(1-2\nu)/2$ en los cortantes).

Deformación de prueba elástica y su descomposición:

\[
\boldsymbol\varepsilon^{e,\text{trial}}_{n+1} = \boldsymbol\varepsilon^\text{tens}_{n+1} - \boldsymbol\varepsilon^p_n
\]

\[
p^\text{trial} = K\,\operatorname{tr}(\boldsymbol\varepsilon_{n+1}) = K\,(\varepsilon_{xx} + \varepsilon_{yy} + \varepsilon_{zz})
\]

(la traza no depende de la sustracción de $\boldsymbol\varepsilon^p_n$ porque $\operatorname{tr}(\boldsymbol\varepsilon^p) = 0$.)

\[
\mathbf e^\text{trial} = \operatorname{dev}(\boldsymbol\varepsilon^{e,\text{trial}}_{n+1}), \qquad \mathbf s^\text{trial} = 2G\,\mathbf e^\text{trial}
\]

con $\mathbf e^\text{trial}$ y $\mathbf s^\text{trial}$ ambos en Voigt 6D tensorial.

Función de fluencia trial:

\[
f^\text{trial} = \|\mathbf s^\text{trial}\| - \sqrt{\tfrac{2}{3}}\,(\sigma_y + H\,\alpha_n)
\]

con $\|\mathbf s^\text{trial}\|^2 = \sum_\text{diag} s_i^2 + 2\sum_\text{off} s_{ij}^2$ (§3).

\textbf{Si $f^\text{trial} \leq \text{tol}$} \ensuremath{\to} paso elástico. $\boldsymbol\sigma_{n+1} = \mathbf C_e\,\boldsymbol\varepsilon^e_n + \mathbf C_e\,\Delta\boldsymbol\varepsilon$ equivalente, evaluado como $\mathbf s^\text{trial} + p^\text{trial}\,\mathbf I$. Estado interno intacto.

\textbf{Si $f^\text{trial} > 0$} \ensuremath{\to} corrector plástico radial. Por isotropía + asociado en J2, la dirección de flujo no cambia entre trial y final: $\mathbf N_{n+1} = \mathbf N^\text{trial} = \mathbf s^\text{trial}/\|\mathbf s^\text{trial}\|$. Forma cerrada:

\[
\boxed{\Delta\gamma = \frac{f^\text{trial}}{2G + \tfrac{2}{3}H}}
\]

Actualización (todo en Voigt 6D tensorial):

\[
\mathbf s_{n+1} = \mathbf s^\text{trial} - 2G\,\Delta\gamma\,\mathbf N, \qquad \boldsymbol\sigma_{n+1} = \mathbf s_{n+1} + p^\text{trial}\,\mathbf I
\]

\[
\boldsymbol\varepsilon^p_{n+1} = \boldsymbol\varepsilon^p_n + \Delta\gamma\,\mathbf N, \qquad \alpha_{n+1} = \alpha_n + \sqrt{\tfrac{2}{3}}\,\Delta\gamma
\]

Coste: O(1) por punto de Gauss. Sin Newton local. \textbf{Algoritmo idéntico al de plane strain extendido a 6D} --- solo cambian los índices de las componentes activas.

\subsection{12. Tangente algorítmica consistente}

Forma cerrada estándar (Simó-Hughes 1998, §3.3):

\[
\mathbf C^\text{alg} = K\,\mathbf 1\otimes\mathbf 1 + 2G(1-\beta)\,\mathbf I_\text{dev} - 2G\,\bar\gamma\,\mathbf N\otimes\mathbf N
\]

con

\[
\beta = \frac{2G\,\Delta\gamma}{\|\mathbf s^\text{trial}\|}, \qquad \bar\gamma = \frac{1}{1 + H/(3G)} - \beta
\]

En Voigt 6D del proyecto:

\begin{itemize}
  \item $\mathbf 1\otimes\mathbf 1$: matriz 6\ensuremath{\times}6 con \texttt{1} en las entradas $(i,j)$ con $i,j \in \{1,2,3\}$ y \texttt{0} resto.
  \item $\mathbf I_\text{dev}$: proyector desviador 6\ensuremath{\times}6. En la base tensorial, $\mathbf I_\text{dev} = \mathbf I - \tfrac{1}{3}\mathbf 1\otimes\mathbf 1$ en el bloque diagonal (3\ensuremath{\times}3) y un bloque identidad escalado en los cortantes con el factor adecuado para que $\mathbf I_\text{dev}\,\boldsymbol\varepsilon = \operatorname{dev}(\boldsymbol\varepsilon)$ cuando $\boldsymbol\varepsilon$ está en convención engineering. En la convención del proyecto (entrada engineering, salida engineering), el bloque cortante de $\mathbf I_\text{dev}$ es la identidad 3\ensuremath{\times}3 simple.
  \item $\mathbf N\otimes\mathbf N$: producto exterior 6\ensuremath{\times}6 de $\mathbf N$ (Voigt 6D tensorial) consigo mismo. \textbf{Cuidado con la convención}: el producto $\mathbf N\otimes\mathbf N$ del kernel debe ser consistente con $\mathbf C^\text{alg}\,\boldsymbol\varepsilon$ engineering. La forma simple: si $\mathbf N$ está almacenado tensorial (sin factor 2 en cortantes), entonces $\mathbf C^\text{alg}$ en convención engineering tiene $\mathbf N\otimes\mathbf N$ con un factor de 2 en las posiciones cortantes --- equivalente a multiplicar las componentes cortantes de $\mathbf N$ por $\sqrt 2$ antes del producto exterior (Mandel-like rescaling local al kernel, idéntico al patrón de VM2D plane strain).
\end{itemize}

\textbf{Simetría}: $\mathbf C^\text{alg}$ es simétrica (J2 asociado), \texttt{IS\_SYMMETRIC = True}.

\textbf{Coste}: precalcular los cuatro escalares $K, G, \beta, \bar\gamma$ y construir la 6\ensuremath{\times}6 sin loops anidados --- O(1) constante.

\subsection{13. Caveats numéricos}

\begin{itemize}
  \item \textbf{Locking volumétrico} en \texttt{Hex8}/\texttt{Tet4} cuando la plasticidad domina: la incompresibilidad plástica $\operatorname{tr}(\boldsymbol\varepsilon^p) = 0$ combinada con $\nu$ no muy lejano de 0.5 produce el mismo locking documentado para \texttt{Elastic3D} (limitación declarada, blindada por test en la Etapa 7). B-bar/F-bar diferidas a caso de uso real.
  \item \textbf{$\nu \to 0.5$}: $K \to \infty$ y el predictor elástico pierde precisión. Constructor rechaza $\nu \geq 0.5$. Recomendado $\nu \leq 0.49$ en metales.
  \item \textbf{$H = 0$ (plasticidad perfecta)}: $\Delta\gamma = f^\text{trial}/(2G)$, sin singularidades. Carga aplicada que supere la capacidad plástica produce divergencia del Newton global del solver (\texttt{LoadExceedsCapacityError}, ADR 0011); responsabilidad del solver, no del material.
  \item \textbf{Tangente algorítmica vs continua}: con $\Delta\gamma > 0$, $\mathbf C^\text{alg} \neq \mathbf C^\text{ep}_\text{continuo}$ por el término $\beta$. Usar la algorítmica garantiza convergencia cuadrática del Newton global.
  \item \textbf{Tolerancia de fluencia} vía ADR 0006: \texttt{admissibility\_scale = \ensuremath{\surd}(2/3)(\ensuremath{\sigma}\_y + H\ensuremath{\cdot}\ensuremath{\alpha})}. Idéntica a plane strain --- la escala física es la misma norma desviadora en la frontera de fluencia.
\end{itemize}


\section{Contrato de implementación}

\begin{lstlisting}[language=yaml]
name: VonMises3D
kind: material
status: validated        # draft → implemented → validated

interface:
  strain_dim: 6
  primary_state_var: alpha     # deformación plástica acumulada equivalente
  is_symmetric: true           # tangente algorítmica simétrica (J2 asociado)

parameters:
  - { name: E,       type: float, required: true,  desc: "Módulo de Young (>0)" }
  - { name: nu,      type: float, required: true,  desc: "Coeficiente de Poisson ∈ (-1, 0.5)" }
  - { name: sigma_y, type: float, required: true,  desc: "Esfuerzo de fluencia inicial (>0)" }
  - { name: H,       type: float, required: false, default: 0.0,
      desc: "Módulo de endurecimiento isótropo lineal (≥0). H=0 ⇒ plasticidad perfecta" }
  - { name: density, type: float, required: false, default: null,
      desc: "Densidad (kg/m³). Opcional al construir; obligatoria solo si se ensambla peso propio o masa (ADR 0008)" }

signature:
  compute_state: "(ε: ndarray(6,), state_vars=None) -> (σ: ndarray(6,), C_alg: ndarray(6,6), state_vars')"
  strain_kind: "3D Voigt — [ε_xx, ε_yy, ε_zz, γ_xy, γ_yz, γ_xz] con γ_ij = 2·ε_ij"

state_schema:
  eps_p: "ndarray(6,) — [ε^p_xx, ε^p_yy, ε^p_zz, ε^p_xy, ε^p_yz, ε^p_xz] (cortantes tensoriales, sin factor 2)"
  alpha: "float ≥ 0 — deformación plástica acumulada equivalente"

conventions:
  sign:  "σ > 0 ⇔ tracción (Reglas §5)"
  voigt: "[xx, yy, zz, xy, yz, xz] (ADR 0012); γ_ij = 2·ε_ij en deformación total engineering;
          ε^p_ij almacenada en componente tensorial (sin factor 2);
          σ_ij tensorial sin factor"

validity:
  - "E > 0, σ_y > 0, H ≥ 0"
  - "ν ∈ (-1, 0.5) estricto; ν = 0.5 rechazado (incompresible)"
  - "|ε| ≲ 1e-2 (régimen de pequeñas deformaciones)"

out_of_scope:
  - "endurecimiento cinemático (back-stress); usar variante futura VonMises3DKinematic"
  - "endurecimiento isótropo no lineal (Voce, exponencial, potencial)"
  - "ablandamiento (H<0); requiere regularización por longitud característica o no-local"
  - "regla de flujo no asociada"
  - "acoplamiento con daño; usar modelo plástico-dañado dedicado"
  - "viscoplasticidad y dependencia de la velocidad de deformación"
  - "grandes deformaciones (descomposición multiplicativa F = F^e·F^p)"
  - "anisotropía (Hill 1948, Barlat); el modelo es estrictamente J2 isótropo"
  - "incompresibilidad estricta (ν=0.5)"

numerical_caveats:
  - "Locking volumétrico en Hex8/Tet4 cuando la plasticidad domina (incompresibilidad plástica + ν moderado-alto). Mitigaciones B-bar/F-bar diferidas. Test específico documenta la limitación."
  - "ν cercano a 0.5: K diverge; recomendado ν ≤ 0.49 (típico ν=0.3 en metales)."
  - "H=0 (plasticidad perfecta): material correcto, divergencia bajo sobrecarga es responsabilidad del solver (ADR 0011)."

acceptance:
  # Cubierto por tests/test_materials_unit.py · TestVonMises3D (a crear).
  unit:
    - name: paso_elastico_bajo_fluencia
      setup: "ε = (1e-5, 0, 0, 0, 0, 0): tracción uniaxial por debajo del yield"
      expect: "σ = C_e·ε exacto; α = 0; eps_p = 0; tangente = C_e"
      tol_rel: 1.0e-12

    - name: yield_uniaxial_at_sigma_y
      setup: "tracción uniaxial monotónica creciente con ε_yy = ε_zz = -ν·ε_xx, ε_xy = ε_yz = ε_xz = 0"
      expect: "primer yield exactamente en σ_xx = σ_y; eps_yield = σ_y/E"
      tol_rel: 1.0e-10

    - name: traccion_biaxial_isotropa
      setup: "ε = (eps, eps, -2ν·eps/(1-ν), 0, 0, 0) — tracción biaxial isótropa en plano x-y con ε_zz libre por σ_zz=0"
      expect: "σ_xx = σ_yy convergen a yield uniforme; α evoluciona; σ_zz ≈ 0"
      tol_rel: 1.0e-4

    - name: cortante_puro_xy_yield
      setup: "γ_xy creciente con resto de componentes 0"
      expect: "primer yield en τ_xy = σ_y/√3 (criterio J2 en cortante puro)"
      tol_rel: 1.0e-10

    - name: incompresibilidad_plastica
      setup: "cualquier paso plástico monotónico"
      expect: "tr(eps_p) = ε^p_xx + ε^p_yy + ε^p_zz = 0 exacto"
      tol_abs: 1.0e-12

    - name: alpha_monotonic
      setup: "carga monotónica creciente con incrementos plásticos no triviales"
      expect: "α no decrece nunca; α_final > 0 si hubo plasticidad"

    - name: descarga_recupera_C_e
      setup: "tras estado plástico, paso con menor ε (descarga elástica)"
      expect: "σ devuelto coherente con eps_p_n acumulado; tangente devuelta = C_e; α intacto"
      tol_rel: 1.0e-10

    - name: tangente_simetrica
      setup: "evaluar C_alg en estado plástico arbitrario"
      expect: "máx |C_alg - C_alg.T| ≤ tol; IS_SYMMETRIC = True"
      tol_abs: 1.0e-8

    - name: unit_invariance_MPa_vs_Pa
      setup: "mismo path adimensional con (E=210, σ_y=0.25) y (E=210e9, σ_y=0.25e9)"
      expect: "α, eps_p (adimensionales) idénticos; σ/E idéntico"
      tol_rel: 1.0e-12

    - name: degeneracion_a_elasticidad_sin_plasticidad
      setup: "trayectoria que nunca alcanza la frontera de fluencia"
      expect: "todos los pasos devuelven C_alg = C_e; eps_p y α permanecen en cero"
      tol_rel: 1.0e-12

    - name: rechazo_inputs_invalidos
      setup: "construir VonMises3D con E≤0, ν fuera de rango, σ_y≤0, H<0"
      expect: "ValueError con mensaje claro en cada caso"

  # Test cruzado VM3D ↔ VM2D plane strain.
  cross_consistency:
    - name: equivalencia_plane_strain
      setup: "ejecutar VM3D con ε = (ε_xx, ε_yy, 0, γ_xy, 0, 0) y VM2D plane strain con ε = (ε_xx, ε_yy, γ_xy);
             mismo path multi-paso con plasticidad activa"
      expect: "σ_xx, σ_yy, σ_xy y α coinciden entre 3D y 2D PS;
              VM3D además devuelve σ_zz coherente con la incompresibilidad plástica"
      tol_rel: 1.0e-10

  # Cubierto por tests/test_solid_3d_plasticity.py (a crear, fuera de esta spec — pertenece a la integración Hex8/Tet4 + VM3D).
  integration:
    - name: pipeline_hex8_traccion_uniaxial_elastoplastica
      setup: "Hex8 unitario, condiciones de contorno equivalentes a tracción uniaxial pura
             (4 caras laterales restringidas en su dirección normal, una cara tirada en x).
             Carga monotónica que cruza el yield."
      expect: "σ_xx en el centro coincide con material standalone bajo el path equivalente;
              α coincide; convergencia Newton global en pocos pasos"
      tol_rel: 1.0e-6

    - name: pipeline_hill_cylinder_3d
      setup: "cilindro hueco bajo presión interna creciente discretizado con Hex8 (octante con simetrías);
             cargar hasta plasticidad parcial"
      expect: "frontera elasto-plástica radial coincide con solución analítica cerrada de Hill (1950)
              al avanzar la presión"
      tol_rel: 1.0e-2

references:
  - "Simó J.C., Hughes T.J.R. (1998). Computational Inelasticity. Springer.
     §3.3 (3D J2 return mapping radial cerrado, tangente algorítmica consistente)."
  - "Crisfield M.A. (1991). Non-linear Finite Element Analysis of Solids and Structures, Vol. 1. Wiley.
     Cap. 6 (J2 plasticity)."
  - "de Souza Neto E.A., Perić D., Owen D.R.J. (2008). Computational Methods for Plasticity. Wiley.
     §7 (3D J2 plasticity)."
  - "Hill R. (1950). The Mathematical Theory of Plasticity. Oxford University Press.
     Cap. V (cilindro hueco bajo presión interna, solución analítica)."
  - "Bathe K.J. (2014). Finite Element Procedures. Prentice Hall. §6.6, §6.8.4."
  - "ADR 0012 — Sólidos 3D: convención Voigt 6D del proyecto."
  - "ADR 0006 — Tolerancias adimensionales: patrón atol + rtol·escala con escala física."

\end{lstlisting}


\section{Implementación}

\begin{itemize}
  \item \textbf{Archivo}: solidum/materials/von\_mises\_3d.py.
  \item \textbf{Clase}: \texttt{VonMises3D}, registrada vía \texttt{@MaterialRegistry.register}. Sin despacho por hipótesis (no aplica en 3D).
  \item \textbf{Kernel Numba} \texttt{\_compute\_j2\_3d}: descomposición volumétrica-desviadora en Voigt 6D tensorial, return mapping radial cerrado, tangente algorítmica consistente cerrada. Predictor elástico construido como $\mathbf s_\text{trial} + p\,\mathbf I$ con $\mathbf s_\text{trial} = 2G(\mathbf e^\text{dev}_\text{trial} - \boldsymbol\varepsilon^p_n)$ --- respeta $\boldsymbol\varepsilon^p_n$ en descargas y reevaluaciones (aplicada desde el inicio la lección del bug detectado en VM2D plane strain el 2026-05-14).
  \item \textbf{State schema}: \texttt{eps\_p} ndarray(6,) \texttt{[xx, yy, zz, xy, yz, xz]} con cortantes tensoriales, \texttt{alpha} escalar.
  \item \textbf{\texttt{admissibility\_scale}}: \texttt{\ensuremath{\surd}(2/3)\ensuremath{\cdot}(\ensuremath{\sigma}\_y + H\ensuremath{\cdot}\ensuremath{\alpha})} --- idéntica a VM2D plane strain (misma escala física en la frontera J2).
  \item \textbf{Tests}:
  \item tests/test\_materials\_unit.py \ensuremath{\cdot} \texttt{TestVonMises3D} (12 tests unitarios) + \texttt{TestVonMises3DvsPlaneStrain} (1 test cruzado VM3D \ensuremath{\leftrightarrow} VM2D plane strain con path multi-paso).
  \item tests/test\_solid\_3d\_plasticity.py \ensuremath{\cdot} 5 tests de integración: \texttt{TestHex8VonMises3DUniaxialFree} (2 --- cubo libre con Poisson, yield uniaxial), \texttt{TestHex8VonMises3DConfinedVsStandalone} (1 --- cross-check contra material standalone), \texttt{TestHex8VonMises3DConvergence} (1 --- Newton converge en pocas iteraciones), \texttt{TestTet4VonMises3DBasic} (1 --- smoke test Tet4 + VM3D).
\end{itemize}

\begin{itemize}
  \item \textbf{Notas de traducción}:
  \item Convención mixta engineering/tensorial idéntica a VM2D: entrada \texttt{strain} engineering (\texttt{\ensuremath{\gamma}\_ij = 2\ensuremath{\cdot}\ensuremath{\varepsilon}\_ij}), salida \texttt{sigma} con cortantes tensoriales (lo que el ensamblador espera), estado \texttt{eps\_p} con cortantes tensoriales. La conversión engineering\ensuremath{\to}tensorial se hace una sola vez al construir el desviador (\texttt{strain[3..5]/2}).
  \item La tangente algorítmica $\mathbf C^\text{alg}$ usa el patrón $K\,\mathbf 1\otimes\mathbf 1 + 2G(1-\beta)\,\mathbf I_\text{dev} - 2G\,\bar\gamma\,\mathbf N\otimes\mathbf N$ con $\mathbf I_\text{dev}$ teniendo $1/2$ en los cortantes (mapea engineering input a tensorial output) y $\mathbf N\otimes\mathbf N$ con $\mathbf N$ tensorial. Esta combinación produce la convención correcta \texttt{engineering\ensuremath{\to}tensorial} en la salida sin factores adicionales --- verificado por el test cruzado VM3D \ensuremath{\leftrightarrow} VM2D plane strain (mismos $\sigma_{xx}, \sigma_{yy}, \sigma_{xy}, \alpha$ a 10 decimales) y por el test de integración Hex8 confinado vs material standalone.
  \item Validación de parámetros en constructor: \texttt{E > 0}, \texttt{\ensuremath{\sigma}\_y > 0}, \texttt{H \ensuremath{\ge} 0}, \texttt{\ensuremath{\nu} \ensuremath{\in} (-1, 0.5)}, \texttt{density \ensuremath{\ge} 0} si se pasa. Mensajes en español.
\end{itemize}


\section{Diálogo}

\begin{itemize}
  \item \textbf{2026-05-21} \ensuremath{\cdot} Spec pre-redactada por la IA al abrir la sub-etapa \textbf{A.bis} (materiales 3D no lineales) tras cierre de la Etapa 7 (ADR 0012). Base: \texttt{VonMises2D.md} (modelo físico) + \texttt{Elastic3D.md} (estructura Voigt 6D + sin variantes de hipótesis).
\end{itemize}

\begin{itemize}
  \item \textbf{2026-05-21} \ensuremath{\cdot} \textbf{Decisión de plumbing (decide la IA, Reglas §3)}: kernels Numba separados \texttt{\_compute\_j2\_3d} (este material) vs \texttt{\_compute\_j2\_plane\_strain} (existente en VonMises2D). \textbf{No} se extrae un \texttt{\_J2Core} compartido en esta entrega. Razones:
\end{itemize}
\begin{enumerate}
  \item La diferencia operativa entre los dos kernels es el número de componentes activas (4 vs 6), que en Numba \textit{nopython} requiere arrays de tamaño fijo --- un kernel "genérico" obligaría a zero-padding sistemático con coste y oscurecería la lectura.
  \item Cada kernel es corto (\textasciitilde{}40 líneas): la duplicación cabe en el mismo PR sin coste de mantenimiento desproporcionado.
  \item La regla de centralización post-dos-casos del proyecto se satisface ya con la documentación cruzada (esta spec + VonMises2D referencian el mismo algoritmo de Simó-Hughes §3.3). Si aparece un tercer caso (e.g. \texttt{VonMises3DCorot} o variante con anisotropía leve) será el momento natural de abrir \texttt{\_J2Core}.
  \item Plane stress 2D queda explícitamente fuera de un eventual \texttt{\_J2Core} por tener un algoritmo genuinamente distinto (Newton local proyectado).
\end{enumerate}

  Si al implementar la duplicación resulta más alta de lo previsto, se reconsidera y se eleva a zona gris.

\begin{itemize}
  \item \textbf{2026-05-21} \ensuremath{\cdot} \textbf{Decisión de plumbing}: el predictor elástico debe construirse como $\boldsymbol\sigma^\text{trial} = \mathbf C_e\,(\boldsymbol\varepsilon - \boldsymbol\varepsilon^p_n)$ (o equivalente desviador + presión), \textbf{no} como $\mathbf C_e\,\boldsymbol\varepsilon$. Lección de VM2D plane strain (bug detectado el 2026-05-14): en descargas elásticas desde estado plástico o reevaluaciones vía \texttt{compute\_gauss\_state(U\_final)}, ignorar $\boldsymbol\varepsilon^p_n$ devuelve \ensuremath{\sigma} inconsistente. Se aplica desde el inicio en VM3D y se cubre con \texttt{test\_descarga\_recupera\_C\_e} y con el test de integración \texttt{test\_pipeline\_hex8\_traccion\_uniaxial\_elastoplastica}.
\end{itemize}

\begin{itemize}
  \item \textbf{2026-05-21} \ensuremath{\cdot} \textbf{Tests cruzados VM3D \ensuremath{\leftrightarrow} VM2D plane strain} (\texttt{cross\_consistency}): garantía adicional de que VM3D no introduce regresión silenciosa sobre el comportamiento ya validado de VM2D. Es la red de seguridad principal para confirmar que la formulación 6D no contiene errores de signo/factor en la convención Voigt --- los tests unitarios analíticos por sí solos no cubren toda la combinatoria de componentes activas.
\end{itemize}

\begin{itemize}
  \item \textbf{2026-05-21} \ensuremath{\cdot} \textbf{Validación contra Hill 3D}: el caso del cilindro hueco bajo presión interna con plasticidad perfecta tiene \textbf{solución analítica cerrada} (Hill 1950, §V) --- la presión a la que la frontera elasto-plástica radial avanza a un radio dado es expresable como función explícita de $\sigma_y, r_i, r_o$. Es el benchmark canónico para validar J2 3D y ya está disponible en 2D plane strain (\texttt{tests/validation/test\_hill\_cylinder\_j2.py}). El test 3D será su análogo discretizado por octante con simetrías. Se redactará y validará tras tener \texttt{Hex8 + VonMises3D} funcionando end-to-end.
\end{itemize}

\begin{itemize}
  \item \textbf{2026-05-21} \ensuremath{\cdot} \textbf{Implementación completa, status \texttt{validated}}. Kernel + clase + 12 tests unitarios + 1 test cruzado + 5 tests de integración (Hex8 uniaxial libre con Poisson, Hex8 confinado vs material standalone, Hex8 convergencia, Tet4 smoke test). \textbf{18 tests añadidos a la suite (804 \ensuremath{\to} 824 verdes; suite global sin regresiones).} Todos los tests pasaron a la primera, incluyendo el test cruzado VM3D \ensuremath{\leftrightarrow} VM2D plane strain a 10 decimales --- señal fuerte de que la convención Voigt 6D en el kernel y la tangente algorítmica son correctas. Validación Hill 3D (cilindro hueco bajo presión interna con plasticidad perfecta) \textbf{diferida a entrega separada} porque requiere setup de malla por octante con simetrías y carga incremental hasta plasticidad parcial; será el primer test de \texttt{tests/validation/} que ejercite VM3D contra solución analítica cerrada.
\end{itemize}


\newpage

\chapter{Modelos Constitutivos — Cohesivos}

\section{CohesiveDamageIsotropic}



\textit{Orden de trabajo. Pre-redactada por la IA como borrador para validación del usuario (autor de la formulación, Retama 2010). La física y la formulación numérica son traducción directa del Cap. 3 de la tesis; la mecánica de implementación (Voigt, registry, signatura, tests) sigue las convenciones del proyecto. El usuario revisa, corrige donde proceda, y aprueba antes de que la IA toque código.}



\section{Especificación física}

\subsection{0. Descripción general}

Material cohesivo \textbf{traction-jump} para la superficie de discontinuidad \texttt{Γ\_d} interior a un elemento finito sólido 2D. Modelo de \textbf{daño escalar isótropo} con activación tipo Rankine (Modo-I), ablandamiento gobernado por la energía de fractura \texttt{G\_F}, y bulk vecino que descarga elásticamente sin disipación adicional.

Distinto en dominio de los materiales continuos \texttt{Elastic2D}, \texttt{VonMises2D}, \texttt{IsotropicDamage2D}:

\begin{itemize}
  \item \textbf{No opera sobre puntos de Gauss del bulk.} Opera sobre el salto \texttt{[[u]] \ensuremath{\in} \ensuremath{\mathbb{R}}\ensuremath{^{2}}} definido sobre \texttt{Γ\_d} y devuelve tracciones \texttt{t \ensuremath{\in} \ensuremath{\mathbb{R}}\ensuremath{^{2}}} definidas sobre la misma superficie.
  \item \textbf{Unidades}: parámetros físicos son \texttt{\ensuremath{\sigma}\_t0} (Pa, resistencia a tracción) y \texttt{G\_F} (N/m, energía de fractura) --- no \texttt{E} y \texttt{\ensuremath{\nu}}.
  \item \textbf{Es la primera implementación} de la familia \texttt{CohesiveMaterial} introducida por ADR 0010.
\end{itemize}

Pensado para hormigón, roca, cerámicos en régimen cuasi-frágil bajo tracción dominante. Diseñado específicamente para alimentar el elemento \texttt{CST\_Embedded2D} (fase 2 del ADR 0010), pero la abstracción es reutilizable por futuros elementos de interfaz cohesiva clásica (CZM) si entran al proyecto.

\subsection{1. Descomposición del salto}

Sin descomposición elástica-plástica (modelo de daño puro, no plasticidad cohesiva). El salto total \texttt{[[u]]} se proyecta sobre el frame local \texttt{(n, s)} de \texttt{Γ\_d}:

\[
\llbracket\mathbf u\rrbracket = \llbracket u_n\rrbracket\,\mathbf n + \llbracket u_s\rrbracket\,\mathbf s
\]

donde \texttt{n} es la normal a \texttt{Γ\_d} (fijada al activarse el elemento, perpendicular a \texttt{\ensuremath{\sigma}\_I} por criterio Rankine) y \texttt{s} es la tangente.

\textbf{Modo-I (fase 1)}: sólo \texttt{\ensuremath{\langle}[[u\_n]]\ensuremath{\rangle}} (parte positiva del salto normal) gobierna la disipación. La componente tangencial \texttt{[[u\_s]]} desliza sin restricción y sin disipación. Modo mixto I-II queda diferido a la fase G del ADR 0010.

\subsection{2. Ley elástica del salto (sin daño)}

\[
\mathbf t = (1 - \omega)\,\mathbf T^{el}\,\llbracket\mathbf u\rrbracket
\]

con \texttt{T\textasciicircum{}el} la rigidez elástica del salto:

\[
\mathbf T^{el} = K_e\,(\mathbf n \otimes \mathbf n)
\]

es decir: rigidez \texttt{K\_e} en dirección normal, cero en dirección tangencial (consecuencia explícita de Modo-I).

\texttt{K\_e} es una rigidez de penalty --- debe escogerse suficientemente grande para que la fase elástica de la grieta sea geométricamente despreciable (\texttt{\ensuremath{\kappa}\_0 = \ensuremath{\sigma}\_t0/K\_e ≪} escala del problema), pero no tan grande que perjudique el condicionamiento del sistema. Guía típica: \texttt{K\_e \ensuremath{\approx} E\_bulk / ℓ\_c} con \texttt{ℓ\_c} una longitud característica del elemento. Ver §12.

\subsection{3. Variable equivalente y función de fluencia}

\[
\llbracket u\rrbracket_\text{eq} = \langle\llbracket u_n\rrbracket\rangle
\]

donde \texttt{\ensuremath{\langle}\ensuremath{\cdot}\ensuremath{\rangle}} son los brackets de McAuley (parte positiva). La función de fluencia (Eq. 3.13 de Retama 2010):

\[
f(\llbracket\mathbf u\rrbracket, \kappa) = \llbracket u\rrbracket_\text{eq} - \kappa \leq 0
\]

donde \texttt{\ensuremath{\kappa}} es el historial del salto equivalente máximo.

\textbf{Interpretación}: la grieta sólo "siente" la tracción (y por tanto sólo daña) cuando se abre (\texttt{[[u\_n]] > 0}). En compresión (\texttt{[[u\_n]] < 0}) o deslizamiento puro, \texttt{f = -\ensuremath{\kappa} < 0}, no hay evolución de daño.

\subsection{4. Variable histórica y evolución (Kuhn-Tucker)}

\[
\kappa_{n+1} = \max\bigl(\kappa_n,\ \langle\llbracket u_n\rrbracket_{n+1}\rangle\bigr), \qquad \kappa_0 = \frac{\sigma_{t0}}{K_e}
\]

Codifica:

\begin{itemize}
  \item \textbf{Carga activa} (\texttt{\ensuremath{\langle}[[u\_n]]\ensuremath{\rangle} > \ensuremath{\kappa}\_n}): \texttt{\ensuremath{\kappa}} crece; el daño puede aumentar.
  \item \textbf{Descarga / recarga elástica} (\texttt{\ensuremath{\langle}[[u\_n]]\ensuremath{\rangle} \ensuremath{\le} \ensuremath{\kappa}\_n}): \texttt{\ensuremath{\kappa}} queda congelado; daño no cambia (secante respecto al origen con rigidez reducida \texttt{(1−\ensuremath{\omega})\ensuremath{\cdot}K\_e}).
  \item \textbf{Compresión} (\texttt{[[u\_n]] < 0}): \texttt{\ensuremath{\langle}[[u\_n]]\ensuremath{\rangle} = 0 \ensuremath{\le} \ensuremath{\kappa}\_n}, \texttt{\ensuremath{\kappa}} congelado. El daño no crece. \textit{Caveat de cierre}: la rigidez efectiva en compresión sigue siendo \texttt{(1−\ensuremath{\omega})\ensuremath{\cdot}K\_e}, no se recupera el contacto rígido --- limitación documentada en §12.
  \item \textbf{Irreversibilidad}: \texttt{\ensuremath{\kappa}} monótono no decreciente; \texttt{\ensuremath{\omega}(\ensuremath{\kappa})} monótono no decreciente.
\end{itemize}

Condiciones de Kuhn-Tucker estándar (Eq. 3.15):

\[
\dot\omega \geq 0, \qquad f \leq 0, \qquad \dot\omega\,f = 0
\]

\subsection{5. Ley de evolución del daño --- softening lineal y exponencial}

El daño \texttt{\ensuremath{\omega}(\ensuremath{\kappa})} se define implícitamente para que la tracción en carga monótona siga una curva prescrita \texttt{T\_\textbackslash\{\}text\{soft\}(\ensuremath{\kappa})}:

\[
t_n(\kappa) = (1 - \omega(\kappa))\,K_e\,\kappa \;\stackrel{!}{=}\; T_\text{soft}(\kappa) \;\Rightarrow\; \omega(\kappa) = 1 - \frac{T_\text{soft}(\kappa)}{K_e\,\kappa}
\]

con \texttt{\ensuremath{\kappa} \ensuremath{\ge} \ensuremath{\kappa}\_0}. Para \texttt{\ensuremath{\kappa} \ensuremath{\le} \ensuremath{\kappa}\_0}, \texttt{\ensuremath{\omega} = 0} por definición.

\textbf{Softening lineal}:

\[
T_\text{soft}(\kappa) = \sigma_{t0}\left(1 - \frac{\kappa - \kappa_0}{w_c - \kappa_0}\right) = \frac{\sigma_{t0}\,(w_c - \kappa)}{w_c - \kappa_0}, \qquad w_c = \frac{2\,G_F}{\sigma_{t0}}
\]

válida en \texttt{\ensuremath{\kappa} \ensuremath{\in} [\ensuremath{\kappa}\_0, w\_c]}. Para \texttt{\ensuremath{\kappa} \ensuremath{\ge} w\_c}: \texttt{T\_\textbackslash\{\}text\{soft\} = 0}, \texttt{\ensuremath{\omega} = 1} (saturada a \texttt{DAMAGE\_MAX}). Forma cerrada de \texttt{\ensuremath{\omega}(\ensuremath{\kappa})}:

\[
\omega(\kappa) = 1 - \frac{\sigma_{t0}\,(w_c - \kappa)}{K_e\,\kappa\,(w_c - \kappa_0)} \qquad (\kappa \in [\kappa_0,\,w_c])
\]

Verificación en los extremos: \texttt{\ensuremath{\omega}(\ensuremath{\kappa}\_0) = 1 − \ensuremath{\sigma}\_\{t0\}/(K\_e\ensuremath{\cdot}\ensuremath{\kappa}\_0) = 0} (por \texttt{\ensuremath{\kappa}\_0 = \ensuremath{\sigma}\_\{t0\}/K\_e}), y \texttt{\ensuremath{\omega}(w\_c) = 1}.

\textbf{Softening exponencial}:

\[
T_\text{soft}(\kappa) = \sigma_{t0}\,\exp\left(-\frac{\sigma_{t0}\,(\kappa - \kappa_0)}{H}\right), \qquad H = G_F - \frac{\sigma_{t0}\,\kappa_0}{2}
\]

asintótica a cero; \texttt{\ensuremath{\omega}(\ensuremath{\kappa}) \ensuremath{\to} 1} cuando \texttt{\ensuremath{\kappa} \ensuremath{\to} \ensuremath{\infty}}, saturada a \texttt{DAMAGE\_MAX} en la práctica. \texttt{H} se deriva imponiendo que la integral total de tracción a lo largo de la curva sea \texttt{G\_F} (consistencia energética).

La validez de la rama exponencial exige \texttt{G\_F > \ensuremath{\sigma}\_\{t0\}\ensuremath{\cdot}\ensuremath{\kappa}\_0/2}, es decir \texttt{K\_e > \ensuremath{\sigma}\_\{t0\}\ensuremath{^{2}}/(2\ensuremath{\cdot}G\_F)} --- condición geométrica trivial cuando \texttt{K\_e} se elige como penalty.

\subsection{6. Variables internas}

\begin{itemize}
  \item \texttt{\ensuremath{\kappa} : float} --- historial del salto equivalente máximo. \textbf{PRIMARY\_STATE\_VAR alternativa} (más fundamental, controla todo el resto).
  \item \texttt{\ensuremath{\omega} : float \ensuremath{\in} [0, DAMAGE\_MAX]} --- daño escalar. \textbf{PRIMARY\_STATE\_VAR adoptada} (más interpretable físicamente; se exporta al post-proceso).
  \item \texttt{\ensuremath{\omega}} es función determinista de \texttt{\ensuremath{\kappa}} (ec. §5); se cachea en el estado para evitar recálculo y para diagnóstico.
\end{itemize}

\subsection{7. Frame local y notación}

Convención del proyecto:

\begin{itemize}
  \item \texttt{n}: normal a \texttt{Γ\_d}, unitaria. Se fija al activar el elemento (perpendicular a \texttt{\ensuremath{\sigma}\_I}).
  \item \texttt{s}: tangente, ortogonal a \texttt{n}, orientada por regla de la mano derecha con \texttt{z} saliendo del plano (consistente con Reglas.md §5).
  \item \texttt{[[u]] = [[u\_n]]\ensuremath{\cdot}n + [[u\_s]]\ensuremath{\cdot}s} con \texttt{[[u\_n]] = [[u]]\ensuremath{\cdot}n}, \texttt{[[u\_s]] = [[u]]\ensuremath{\cdot}s}.
  \item \texttt{t = t\_n\ensuremath{\cdot}n + t\_s\ensuremath{\cdot}s} con \texttt{t\_n = K\_e\ensuremath{\cdot}(1−\ensuremath{\omega})\ensuremath{\cdot}[[u\_n]]} en Modo-I, \texttt{t\_s = 0}.
\end{itemize}

El material recibe \texttt{[[u]]} y devuelve \texttt{t} en \textbf{ejes locales \texttt{(n, s)} de \texttt{Γ\_d}}. La transformación a ejes globales es responsabilidad del elemento que lo consume, \textbf{no} del material. Mismo patrón que en \texttt{VonMises2D} (recibe \texttt{\ensuremath{\varepsilon}} en Voigt, no se ocupa de transformaciones de elemento).


\section{Formulación numérica}

\subsection{8. Algoritmo en un paso (explícito, sin Newton local)}

Dado \texttt{\ensuremath{\kappa}\_n} y \texttt{[[u]]\_\{n+1\}} (en ejes locales):

\begin{enumerate}
  \item Calcular \texttt{[[u]]\_eq = \ensuremath{\langle}[[u\_n]]\ensuremath{\rangle} = max(0, [[u]]\ensuremath{\cdot}n\_local)}. \textit{(En el frame local, \texttt{n\_local = (1, 0)} y \texttt{[[u\_n]] = [[u]][0]}. Se mantiene la notación abstracta por claridad física.)}
  \item Si \texttt{[[u]]\_eq > \ensuremath{\kappa}\_n} y \texttt{[[u]]\_eq > \ensuremath{\kappa}\_0}: \textbf{carga activa}, \texttt{\ensuremath{\kappa}\_\{n+1\} = [[u]]\_eq}. Si no: \textbf{descarga / no daño}, \texttt{\ensuremath{\kappa}\_\{n+1\} = \ensuremath{\kappa}\_n}.
  \item Aplicar \texttt{\ensuremath{\omega}(\ensuremath{\kappa}\_\{n+1\})} por la fórmula §5; saturar a \texttt{DAMAGE\_MAX}.
  \item Tracción: \texttt{t\_n = (1 − \ensuremath{\omega})\ensuremath{\cdot}K\_e\ensuremath{\cdot}[[u\_n]]}, \texttt{t\_s = 0}.
  \item Tangente: ver §9.
\end{enumerate}

El algoritmo es \textbf{explícito}: no requiere Newton local. La actualización de \texttt{\ensuremath{\kappa}} y \texttt{\ensuremath{\omega}} es directa. Mismo patrón que \texttt{IsotropicDamage2D}.

\subsection{9. Tangente algorítmica consistente}

\[
\mathbf T^\text{alg} = \frac{\partial\,\mathbf t_{n+1}}{\partial\,\llbracket\mathbf u\rrbracket_{n+1}}
\]

Derivando \texttt{t = (1−\ensuremath{\omega})\ensuremath{\cdot}K\_e\ensuremath{\cdot}[[u\_n]]\ensuremath{\cdot}n} en el frame local con \texttt{n = (1,0)}:

\begin{itemize}
  \item \textbf{Sin daño activo} (\texttt{\ensuremath{\kappa} \ensuremath{\le} \ensuremath{\kappa}\_0}):
\end{itemize}

\[
\mathbf T^\text{alg} = K_e\,(\mathbf n \otimes \mathbf n) = K_e\begin{pmatrix}1 & 0 \\ 0 & 0\end{pmatrix}
\]

\begin{itemize}
  \item \textbf{Descarga} (\texttt{[[u]]\_eq \ensuremath{\le} \ensuremath{\kappa}\_n}):
\end{itemize}

\[
\mathbf T^\text{alg} = (1-\omega)\,K_e\,(\mathbf n \otimes \mathbf n)
\]

(secante reducida; \texttt{\ensuremath{\partial}\ensuremath{\kappa}/\ensuremath{\partial}[[u]] = 0} anula el segundo término).

\begin{itemize}
  \item \textbf{Carga activa} (\texttt{[[u]]\_eq > \ensuremath{\kappa}\_n}, \texttt{\ensuremath{\kappa} > \ensuremath{\kappa}\_0}, \texttt{\ensuremath{\omega} < DAMAGE\_MAX}):
\end{itemize}

\[
\mathbf T^\text{alg} = (1-\omega)\,K_e\,(\mathbf n \otimes \mathbf n) - \frac{d\omega}{d\kappa}\,K_e\,\llbracket u_n\rrbracket\,(\mathbf n \otimes \mathbf n)
\]

\[
= \left[(1-\omega) - \frac{d\omega}{d\kappa}\,\llbracket u_n\rrbracket\right]\,K_e\,(\mathbf n \otimes \mathbf n)
\]

\begin{itemize}
  \item \textbf{Cap numérico de la tangente} (\texttt{\ensuremath{\omega} \ensuremath{\ge} DAMAGE\_MAX}, o \texttt{\ensuremath{\kappa} \ensuremath{\ge} w\_c} en lineal):
\end{itemize}

\[
\mathbf T^\text{alg} = (1-\text{DAMAGE\_MAX})\,K_e\,(\mathbf n \otimes \mathbf n)
\]

El cap se aplica \textbf{sólo a la rigidez tangente} para evitar singularidad del Newton cuando \texttt{\ensuremath{\omega} \ensuremath{\to} 1}. La tracción se calcula con el \texttt{\ensuremath{\omega}} físico (que llega a \texttt{1.0} exacto en lineal con \texttt{\ensuremath{\kappa} \ensuremath{\ge} w\_c}, asintóticamente en exponencial) y por tanto \texttt{t\_n = 0} cuando la grieta está totalmente abierta, sin residuo numérico contaminando la energía disipada.

Las derivadas \texttt{d\ensuremath{\omega}/d\ensuremath{\kappa}} cerradas:

\begin{itemize}
  \item \textbf{Lineal}, en \texttt{\ensuremath{\kappa} \ensuremath{\in} [\ensuremath{\kappa}\_0, w\_c]}. Reescribiendo \texttt{\ensuremath{\omega}(\ensuremath{\kappa}) = 1 − C\ensuremath{\cdot}(w\_c − \ensuremath{\kappa})/\ensuremath{\kappa} = 1 − C\ensuremath{\cdot}w\_c/\ensuremath{\kappa} + C} con \texttt{C = \ensuremath{\sigma}\_\{t0\} / [K\_e\ensuremath{\cdot}(w\_c − \ensuremath{\kappa}\_0)]} constante:
\end{itemize}

  \[
\frac{d\omega}{d\kappa} = \frac{C\,w_c}{\kappa^2} = \frac{\sigma_{t0}\,w_c}{K_e\,\kappa^2\,(w_c - \kappa_0)}
\]

  Positiva y finita en todo el rango \texttt{[\ensuremath{\kappa}\_0, w\_c]}. Verificación en extremos: en \texttt{\ensuremath{\kappa} = \ensuremath{\kappa}\_0} (con \texttt{\ensuremath{\kappa}\_0 = \ensuremath{\sigma}\_\{t0\}/K\_e} \ensuremath{\Rightarrow} \texttt{K\_e\ensuremath{\cdot}\ensuremath{\kappa}\_0\ensuremath{^{2}} = \ensuremath{\sigma}\_\{t0\}\ensuremath{\cdot}\ensuremath{\kappa}\_0}) toma el valor \texttt{w\_c / [\ensuremath{\kappa}\_0\ensuremath{\cdot}(w\_c − \ensuremath{\kappa}\_0)]}; en \texttt{\ensuremath{\kappa} = w\_c} toma \texttt{\ensuremath{\sigma}\_\{t0\} / [K\_e\ensuremath{\cdot}w\_c\ensuremath{\cdot}(w\_c − \ensuremath{\kappa}\_0)]}. Dimensionalmente, \texttt{[d\ensuremath{\omega}/d\ensuremath{\kappa}] = m\ensuremath{^{-}}¹} (\ensuremath{\omega} adimensional sobre \ensuremath{\kappa} con unidad de longitud). \ensuremath{\checkmark}

\begin{itemize}
  \item \textbf{Exponencial}, en \texttt{\ensuremath{\kappa} > \ensuremath{\kappa}\_0}:
\end{itemize}
  \[
\frac{d\omega}{d\kappa} = \frac{T_\text{soft}(\kappa)}{K_e\,\kappa}\left(\frac{1}{\kappa} + \frac{\sigma_{t0}}{H}\right) = (1 - \omega(\kappa))\left(\frac{1}{\kappa} + \frac{\sigma_{t0}}{H}\right)
\]

  \textit{(Análoga a la forma de \texttt{IsotropicDamage2D} con \texttt{\ensuremath{\alpha} \ensuremath{\leftrightarrow} \ensuremath{\sigma}\_\{t0\}/H}.)}

\subsection{10. Simetría de la tangente}

A diferencia de \texttt{IsotropicDamage2D} (donde el producto externo \texttt{(C\_e\ensuremath{\cdot}\ensuremath{\varepsilon}) ⊗ (M\ensuremath{\cdot}\ensuremath{\varepsilon})} rompe la simetría), aquí ambos vectores que entran al rank-1 colapsan a \texttt{n}:

\[
\mathbf T^\text{alg} = \alpha(\kappa, \omega, \llbracket u_n\rrbracket)\,K_e\,(\mathbf n \otimes \mathbf n)
\]

con \texttt{\ensuremath{\alpha}} escalar. \texttt{n ⊗ n} es simétrica, por tanto \texttt{T\textasciicircum{}alg} es \textbf{simétrica} en Modo-I.

Consecuencia: \texttt{IS\_SYMMETRIC = True} para este material. El despachador algebraico (ADR 0003) puede usar Cholesky/LDL\ensuremath{^{\top}} si el resto del dominio también es simétrico. \textbf{Modo mixto (fase G) reintroducirá asimetría} --- se evaluará entonces.

\subsection{11. Decisión de régimen durante el Newton global}

Dentro de una iteración del Newton global, el material recibe \texttt{[[u]]\_trial} y decide carga/descarga comparando con \texttt{\ensuremath{\kappa}\_n} (estado committed del paso anterior). Mismo patrón que \texttt{IsotropicDamage2D} §9. La transición loading\ensuremath{\leftrightarrow}unloading es discontinua en la tangente; en problemas patológicos cerca del umbral, considerar arc-length (\texttt{ArcLengthSolver} ya disponible).

\subsection{12. Caveats numéricos}

\begin{itemize}
  \item \textbf{Elección de \texttt{K\_e} (penalty)}. Si \texttt{K\_e} muy pequeño: \texttt{\ensuremath{\kappa}\_0 = \ensuremath{\sigma}\_\{t0\}/K\_e} no es despreciable y la respuesta elástica deforma significativamente la curva (la grieta "se abre" antes del pico). Si \texttt{K\_e} muy grande: condicionamiento del sistema empeora (\texttt{K\_e ≫ E\_bulk\ensuremath{\cdot}L\_e} introduce diferencia de magnitudes en \texttt{K\_\{ũũ\}}). Regla práctica: \texttt{K\_e \ensuremath{\approx} 10\ensuremath{\cdot}E\_bulk/ℓ\_c} con \texttt{ℓ\_c} la longitud característica del elemento (su lado más corto). El usuario lo declara explícitamente en YAML; no hay default automático.
  \item \textbf{Cierre en compresión}. Cuando \texttt{[[u\_n]] < 0} el modelo devuelve \texttt{t\_n = (1−\ensuremath{\omega})\ensuremath{\cdot}K\_e\ensuremath{\cdot}[[u\_n]]} (compresión reducida). Físicamente, la grieta cerrada debería transmitir compresión a rigidez completa \texttt{K\_e} (contacto unilateral). En fase 1 esta limitación se acepta tal cual --- los benchmarks de la tesis (Van Vliet, viga SEN) son traccionados puros. Un modelo con corrección de cierre se contempla como \textbf{deuda explícita} para casos con cambio de signo en la grieta.
  \item \textbf{Cap numérico de la rigidez tangente}. \texttt{DAMAGE\_MAX < 1} se aplica \textbf{sólo} sobre \texttt{T\_tan} para evitar singularidad del Newton (\texttt{stiffness = max(1−\ensuremath{\omega}, 1−DAMAGE\_MAX)\ensuremath{\cdot}K\_e}). El \texttt{\ensuremath{\omega}} reportado en el estado y usado en \texttt{t\_n = (1−\ensuremath{\omega})\ensuremath{\cdot}K\_e\ensuremath{\cdot}[[u\_n]]} es físico (puede valer exactamente 1.0 en lineal con \texttt{\ensuremath{\kappa} \ensuremath{\ge} w\_c}). Diferencia respecto al patrón en \texttt{IsotropicDamage2D}: allí \texttt{K\_e ≡ E} es moderado y truncar \ensuremath{\omega} no introduce residuo significativo en \ensuremath{\sigma}; aquí \texttt{K\_e} es penalty (≫E) y \texttt{u\_n} recorre rangos macroscópicos, así que truncar \ensuremath{\omega} contaminaría la energía disipada en \textasciitilde{}10--30 \%.
  \item \textbf{Sensibilidad de malla} en régimen de softening. Sin regularización, el patrón de fallo localiza en función del tamaño y orientación de elementos. La integración con el elemento \texttt{CST\_Embedded2D} (fase 2 del ADR 0010) y con el \texttt{l\_d = (A/h)\ensuremath{\cdot}cos(\ensuremath{\theta}−\ensuremath{\alpha})} de Cap. 6 da objetividad parcial respecto a la longitud de la grieta --- pero la dirección de propagación sigue dependiendo de la dirección de tensiones principales en el elemento. Mitigación a nivel del \textit{elemento}, no del material.
  \item \textbf{Validación energética}. Por construcción \texttt{\ensuremath{\int}\_0\textasciicircum{}\{w\_c\} t(w)\ensuremath{\cdot}dw = G\_F}. Esto se verifica en los tests (§acceptance). Si la curva implementada no integra a \texttt{G\_F}, hay error en \texttt{w\_c} o en \texttt{H}.
\end{itemize}


\section{Contrato de implementación}

\begin{lstlisting}[language=yaml]
name: CohesiveDamageIsotropic
kind: cohesive_material         # nueva familia, ADR 0010
status: validated

interface:
  jump_dim: 2                   # [[u]] ∈ ℝ² en 2D (componentes n, s en frame local)
  primary_state_var: damage     # ω, escalar exportable al post-proceso
  is_symmetric: true            # tangente simétrica en Modo-I (ver §10)

parameters:
  - { name: sigma_t0,           type: float, required: true,  desc: "Resistencia a tracción [Pa]; iniciación del daño" }
  - { name: G_f,                type: float, required: true,  desc: "Energía de fractura [N/m]; área bajo la curva t–[[u_n]]" }
  - { name: K_e,                type: float, required: true,  desc: "Rigidez elástica del salto en dirección normal [Pa/m] (penalty); ver §12" }
  - { name: softening,          type: str,   required: true,  desc: "Forma de la curva de softening: 'linear' | 'exponential'" }

signature:
  compute_traction: "([[u]]: ndarray(2,), state_vars=None) -> (t: ndarray(2,), T_tan: ndarray(2,2), state_vars')"
  jump_kind: "frame local — [[u]] = ([[u_n]], [[u_s]]) en ejes (n, s) de Γ_d"

state_schema:
  kappa: "float ≥ κ_0 — historial del salto equivalente máximo"
  damage: "float ∈ [0, DAMAGE_MAX] — variable de daño escalar"

conventions:
  sign: "[[u_n]] > 0 ⇔ apertura (tracción); el modelo daña sólo en apertura (Modo-I)"
  frame: "ejes locales (n, s) de Γ_d; transformación a globales por el elemento consumidor"
  damage_irreversibility: "κ monótono no decreciente vía max(κ_old, ⟨[[u_n]]⟩); ω(κ) monótono creciente"
  units: "[σ_t0] = Pa, [G_F] = N/m, [K_e] = Pa/m (asumiendo SI; el sistema no convierte unidades)"

validity:
  - "σ_t0 > 0, G_F > 0, K_e > 0"
  - "K_e > σ_t0² / (2·G_F)  — consistencia energética para softening exponencial (§5)"
  - "softening ∈ {'linear', 'exponential'}"

out_of_scope:
  - "Modo mixto I-II (cizalla acoplada al daño); diferido a fase G del ADR 0010"
  - "Contacto / cierre unilateral en compresión; el modelo reduce rigidez al estar dañado, no recupera contacto rígido"
  - "Anisotropía del daño (tensor de daño D vs escalar ω)"
  - "Acoplamiento viscoso (rate-dependent); rate-independent puro en fase 1"
  - "Fatiga / acumulación cíclica del daño"
  - "Regularización para mesh-objectivity en propagación; se aborda a nivel de elemento (l_d correcto, Cap. 6) y no de material"

numerical_caveats:
  - "Penalty K_e: balance entre κ_0 despreciable (K_e grande) y condicionamiento (K_e moderado). Guía: K_e ≈ 10·E_bulk/ℓ_c. Declarado por el usuario en YAML, sin default automático."
  - "Cierre en compresión: rigidez reducida (1−ω)·K_e, no recuperación de contacto rígido. Aceptable para benchmarks puramente traccionados; deuda explícita para casos con cambio de signo."
  - "Cap por DAMAGE_MAX aplicado sólo a la rigidez tangente (stiffness = max(1−ω, 1−DAMAGE_MAX)·K_e) para evitar singularidad del Newton. La tracción se calcula con ω físico (1.0 exacto en lineal con κ≥w_c) y por tanto t_n = 0 cuando la grieta está totalmente abierta, sin sesgo en la energía disipada."
  - "Régimen elastic↔softening discontinuo en la tangente al cruzar κ_0; régimen loading↔unloading discontinuo al cruzar κ_n committed. Mismo patrón que IsotropicDamage2D."
  - "En la zona ω ≥ DAMAGE_MAX la tangente colapsa a la secante reducida (1−DAMAGE_MAX)·K_e, no a la consistente; el Newton degrada de convergencia cuadrática a lineal en esa rama (típicamente irrelevante porque la rama corresponde a un elemento prácticamente roto)."

acceptance:
  verification:
    - name: respuesta_elastica_bajo_umbral
      setup: "carga monótona en [[u_n]] ∈ [0, κ_0); [[u_s]] = 0"
      expect: "ω = 0; t_n = K_e·[[u_n]] exacto; t_s = 0; T_tan = K_e·(n⊗n)"
      tol_rel: 1.0e-12

    - name: inicio_dano_en_umbral
      setup: "[[u_n]] cruza κ_0 = σ_t0/K_e por encima por primera vez"
      expect: "ω salta de 0 a un valor positivo continuo en κ_0 (ω(κ_0) = 0+); t_n(κ_0) = σ_t0 exacto"
      tol_rel: 1.0e-10

    - name: energia_fractura_softening_lineal
      setup: "carga monótona [[u_n]] de 0 a w_c = 2·G_F/σ_t0 con softening='linear', K_e tal que κ_0 ≪ w_c"
      expect: "∫_0^{w_c} t_n([[u_n]])·d[[u_n]] = G_F  (cuadratura trapezoidal con ≥1000 puntos)"
      tol_rel: 1.0e-3
      note: "La tolerancia 1e-3 absorbe el error de cuadratura; la consistencia analítica es exacta por construcción de w_c."

    - name: energia_fractura_softening_exponencial
      setup: "carga monótona [[u_n]] de 0 a 20·G_F/σ_t0 con softening='exponential', K_e tal que κ_0 ≪ G_F/σ_t0"
      expect: "∫ t_n·d[[u_n]] ≈ G_F·(1 − ε_truncamiento), con ε_truncamiento < 1e-6 por la cola exponencial cortada"
      tol_rel: 1.0e-3

    - name: descarga_secante
      setup: "carga hasta κ ∈ (κ_0, w_c/2), después descarga [[u_n]] → 0"
      expect: "rama de descarga lineal con pendiente (1−ω(κ))·K_e; al volver a [[u_n]]=0, t_n = 0 y ω no ha cambiado"
      tol_rel: 1.0e-10

    - name: recarga_hasta_mismo_kappa_sin_dano_extra
      setup: "carga → descarga → recarga hasta [[u_n]] = κ (mismo κ histórico)"
      expect: "ω no cambia (κ_{n+1} = max(κ_n, [[u_n]]) = κ_n)"
      tol_rel: 1.0e-12

    - name: recarga_mas_alla_continua_danando
      setup: "carga → descarga → recarga hasta [[u_n]] > κ histórico"
      expect: "ω evoluciona desde el valor previo hacia ω([[u_n]]_nuevo) sin saltos"
      tol_rel: 1.0e-10

  specific:
    - name: tangente_simetrica_en_carga_y_descarga
      setup: "estados de carga activa, descarga, sin daño, saturado"
      expect: "T_tan = T_tan^T (simetría exacta en Modo-I); ‖T_tan − T_tan^T‖_F / ‖T_tan‖_F < 1e-14"
      tol_rel: 1.0e-14

    - name: tangente_diferencia_finita_consistente
      setup: "carga activa con [[u_n]] ∈ (κ_0, w_c/2); perturbar [[u]] con ε = 1e-7"
      expect: "‖T_tan_FD − T_tan_analitica‖ / ‖T_tan‖ < 1e-5"
      tol_rel: 1.0e-5

    - name: tangencial_no_dana_ni_transmite
      setup: "[[u_s]] arbitrario, [[u_n]] = 0"
      expect: "t_s = 0; t_n = 0; ω no cambia"
      tol_rel: 1.0e-14

    - name: compresion_no_dana
      setup: "[[u_n]] < 0 puro (penetración)"
      expect: "⟨[[u_n]]⟩ = 0 ⇒ f = −κ ≤ 0 ⇒ ω no cambia. t_n = (1−ω_previo)·K_e·[[u_n]] (caveat de cierre, §12)"
      tol_rel: 1.0e-12

    - name: saturacion_en_w_c_softening_lineal
      setup: "[[u_n]] = w_c·1.01 con softening='linear'"
      expect: "ω = 1.0 exacto, t_n = 0 exacto (grieta totalmente abierta). T_tan = (1−DAMAGE_MAX)·K_e (cap numérico sólo sobre la rigidez)."
      tol_rel: 1.0e-10

    - name: degeneracion_a_elasticidad_intacta
      setup: "σ_t0 → ∞ (10^20 Pa) o cualquier [[u_n]] tan pequeño que jamás active"
      expect: "ω = 0, t_n = K_e·[[u_n]], T_tan = K_e·(n⊗n) (material intacto puro)"
      tol_rel: 1.0e-12

  arch:
    - name: registry_kind_es_cohesive
      expect: "Registro vía CohesiveMaterialRegistry, no MaterialRegistry; kind='cohesive_material' en YAML"

    - name: is_symmetric_attribute
      expect: "CohesiveDamageIsotropic.IS_SYMMETRIC == True"

references:
  - "Retama Velasco, J. (2010). Formulation and Approximation to Problems in Solids by Embedded Discontinuity Models. Tesis Doctoral, UNAM. **Cap. 3 — Discrete damage models**. Fórmulas §3.1 (energía libre, ec. 3.1–3.4); §3.1.1 (tensor tangente, ec. 3.7–3.12); §3.1.2 (función de fluencia, ec. 3.13–3.14); §3.1.3 (Kuhn-Tucker, ec. 3.15–3.16)."
  - "Hillerborg, A., Modéer, M., Petersson, P.-E. (1976). Analysis of crack formation and crack growth in concrete by means of fracture mechanics and finite elements. *Cement and Concrete Research* 6, 773-782. **Origen del modelo cohesivo con energía de fractura G_F**."
  - "Barenblatt, G.I. (1962). The mathematical theory of equilibrium cracks in brittle fracture. *Advances in Applied Mechanics* 7, 55-129."
  - "Simó, J.C., Ju, J.W. (1987). Strain- and stress-based continuum damage models — I. Formulation. *Int. J. Solids Struct.* 23, 821-840."
  - "ADR 0010 — Discontinuidades interiores embebidas (familia CohesiveMaterial, hoja de ruta)."
  - "Spec [IsotropicDamage2D](IsotropicDamage2D.md) — patrón análogo de daño isótropo escalar para el caso continuo."

\end{lstlisting}


\section{Implementación}

\begin{itemize}
  \item Archivo: \texttt{solidum/cohesive\_materials/damage\_isotropic.py}
  \item Clase: \texttt{CohesiveDamageIsotropic} (registrada en \texttt{CohesiveMaterialRegistry})
  \item Base abstracta: \texttt{solidum/core/cohesive\_material.py} (\texttt{CohesiveMaterial})
  \item Registry: \texttt{solidum/registry.py} (\texttt{CohesiveMaterialRegistry}, paralelo a \texttt{MaterialRegistry})
  \item Autodiscover: \texttt{solidum.autodiscover.initialize()} descubre \texttt{solidum.cohesive\_materials} automáticamente.
  \item Tests:
  \item \texttt{tests/test\_cohesive\_damage\_isotropic.py} --- todos los casos de \texttt{acceptance} (verification + specific + arch).
  \item Notas de traducción:
  \item El algoritmo §8 es explícito (sin Newton local); no se usa \texttt{is\_admissible} en este material.
  \item \texttt{JUMP\_DIM = 2}, \texttt{PRIMARY\_STATE\_VAR = 'damage'}, \texttt{IS\_SYMMETRIC = True}.
  \item Frame local interno: \texttt{n = (1, 0)}, \texttt{s = (0, 1)}; el material no transforma a globales (responsabilidad del elemento consumidor).
  \item \texttt{t\_s} y \texttt{T\_tan[1,1]} son 0 por construcción (Modo-I). El elemento que consuma el material debe estar preparado para rigidez tangencial nula.
  \item \texttt{K\_e} obligatorio en el constructor (sin default). \texttt{kappa\_0}, \texttt{w\_c} (lineal) y \texttt{H} (exponencial) se derivan en \texttt{\_\_init\_\_}.
\end{itemize}


\section{Diálogo}

\begin{itemize}
  \item \textbf{2026-05-13} \ensuremath{\cdot} Spec pre-redactada por la IA como \textbf{orden de trabajo} sobre el Cap. 3 de Retama 2010 (tesis del usuario). La física es traducción directa de la tesis; la mecánica del proyecto (Voigt, registry, signatura, tests) sigue convenciones de \texttt{IsotropicDamage2D}. Pendiente de validación del usuario.
  \item \textbf{2026-05-18} \ensuremath{\cdot} Usuario delega las tres decisiones a la IA. Cierre con justificación arquitectural:
  \item \textbf{PRIMARY\_STATE\_VAR = \texttt{damage} (\ensuremath{\omega})} --- adoptado. Razones: (a) coherencia con \texttt{IsotropicDamage2D} (mismo criterio en toda la familia de daños), (b) \ensuremath{\omega} \ensuremath{\in} [0, DAMAGE\_MAX] es adimensional e interpretable físicamente y es lo que se colorea en el post-proceso VTK, (c) \ensuremath{\kappa} tiene unidades [m] y depende de \texttt{K\_e} (no es portable entre materiales). \ensuremath{\kappa} se almacena como variable interna secundaria, controla todo determinísticamente vía \texttt{\ensuremath{\omega}(\ensuremath{\kappa})} cerrada.
  \item \textbf{\texttt{d\ensuremath{\omega}/d\ensuremath{\kappa}} lineal cerrada en §9} --- hecho. Forma simplificada: \texttt{d\ensuremath{\omega}/d\ensuremath{\kappa} = \ensuremath{\sigma}\_t0\ensuremath{\cdot}w\_c / [K\_e\ensuremath{\cdot}\ensuremath{\kappa}\ensuremath{^{2}}\ensuremath{\cdot}(w\_c − \ensuremath{\kappa}\_0)]}. Documentada en la spec para que el docstring del código pueda referenciarla en vez de derivarla.
  \item \textbf{\texttt{K\_e} obligatorio en YAML, sin default automático} --- adoptado. ADR 0010 §1 separa explícitamente \texttt{bulk\_material} y \texttt{cohesive\_material}; derivar \texttt{K\_e} de \texttt{E\_bulk\ensuremath{\cdot}ℓ\_c} los acoplaría a nivel de material y ensuciaría el contrato (el material recibiría \texttt{E\_bulk} y \texttt{ℓ\_c} como inputs ocultos). Si en una iteración futura el YAML pide derivado automático, la fábrica vivirá en el elemento (\texttt{CST\_Embedded2D} construye el \texttt{CohesiveDamageIsotropic} con \texttt{K\_e} derivado si se omite), no en el material.
  \item \textbf{2026-05-18} \ensuremath{\cdot} Status \ensuremath{\to} \texttt{validated}. La IA arranca implementación.
  \item \textbf{2026-05-18} \ensuremath{\cdot} \textit{Nota tras implementación}: durante los tests el corte por \texttt{DAMAGE\_MAX} aplicado a \ensuremath{\omega} introducía \textasciitilde{}10 \% de exceso en \texttt{\ensuremath{\int}t\ensuremath{\cdot}d[[u\_n]]} respecto a \texttt{G\_F} (vs \textasciitilde{}0.001 esperado). Razón física: en cohesivos \texttt{K\_e} es penalty (≫ E\_bulk) y \texttt{u\_n} recorre rango macroscópico, así que \texttt{(1−DAMAGE\_MAX)\ensuremath{\cdot}K\_e\ensuremath{\cdot}u\_n} no es despreciable. Fix arquitectural: el cap por \texttt{DAMAGE\_MAX} se aplica \textbf{sólo a la rigidez tangente} (para mantener el Newton no singular), no al \texttt{\ensuremath{\omega}} reportado ni a \texttt{t\_n}. Actualizadas §9, §12 y \texttt{acceptance.saturacion\_en\_w\_c\_softening\_lineal} consecuentemente. Distinto del patrón en \texttt{IsotropicDamage2D} por la diferencia en escalas de \texttt{K\_e} y de la variable de entrada.
\end{itemize}


\newpage

\chapter{Modelos Constitutivos — Térmicos}

\section{ThermalConduction}



\textit{Spec pre-redactada por la IA (2026-08-25) sobre las seis decisiones de alcance cerradas con el usuario para la Etapa 8. \textbf{El usuario revisa la física; el plumbing está resuelto.} Los puntos abiertos de §Diálogo son físicos o de alcance, nunca arquitecturales.}

\textit{> Primer material de la familia térmica. Etapa 8 --- C1 núcleo mínimo: conducción pura, sin convección, radiación ni acoplamiento termomecánico.}



\section{Especificación física}

\subsection{0. Descripción general}

Modelo constitutivo de \textbf{conducción de calor} según la ley de Fourier. Relaciona el vector de flujo de calor $\mathbf q$ con el gradiente de temperatura $\nabla T$. Es el análogo térmico de \texttt{Elastic2D}/\texttt{Elastic3D}: lineal, sin variables internas, con "tangente" constante.

No pertenece a la familia \texttt{Material} mecánica. Los materiales mecánicos relacionan $\boldsymbol\sigma$ con $\boldsymbol\varepsilon$ en notación Voigt; éste relaciona dos vectores en el espacio físico, sin notación Voigt y sin tensor simétrico de por medio. Vive en una \textbf{familia paralela} \texttt{ThermalMaterial} con su propio registro, exactamente como \texttt{CohesiveMaterial} (ADR 0010) resolvió el mismo tipo de incompatibilidad semántica.

\subsection{1. Ley constitutiva --- Fourier}

\[
\mathbf q = -\,\mathbf k\, \nabla T
\]

con:

\begin{itemize}
  \item $\mathbf q$ --- vector de flujo de calor [W/m\ensuremath{^{2}}]. Apunta en la dirección del transporte de energía.
  \item $\nabla T$ --- gradiente de temperatura [K/m].
  \item $\mathbf k$ --- tensor de conductividad térmica [W/(m\ensuremath{\cdot}K)]. Matriz $2\times2$ en 2D, $3\times3$ en 3D.
\end{itemize}

\textbf{El signo negativo es la física, no una convención elegible}: el calor fluye de mayor a menor temperatura (segunda ley de la termodinámica). De ahí se sigue que $\mathbf k$ debe ser \textbf{definida positiva} --- si no lo fuera, existiría una dirección en la que el calor fluiría espontáneamente hacia la región más caliente.

\subsection{2. Conductividad tensorial e isótropa}

El contrato guarda siempre $\mathbf k$ como tensor (decisión 4 del alcance). El caso isótropo es el particular

\[
\mathbf k = k\,\mathbf I
\]

y el constructor acepta un escalar $k$, expandiéndolo internamente. Esto evita romper el contrato cuando entren materiales ortótropos (madera, composites, roca estratificada, hormigón fibroreforzado), habituales en el dominio del usuario.

\textbf{Requisitos sobre $\mathbf k$}:

\begin{itemize}
  \item \textbf{Simétrica}: $\mathbf k = \mathbf k^\top$. Lo exigen las relaciones recíprocas de Onsager para el transporte en medios sin campo magnético.
  \item \textbf{Definida positiva}: todos los autovalores $> 0$. Equivale a exigir $k > 0$ en el caso isótropo.
\end{itemize}

\subsection{3. Ecuación de balance que este material alimenta}

El material aporta $\mathbf k$ (y, para el transitorio, $\rho c$) a la ecuación de conducción:

\[
\rho c\,\frac{\partial T}{\partial t} = -\nabla\cdot\mathbf q + Q = \nabla\cdot(\mathbf k\nabla T) + Q
\]

En régimen \textbf{estacionario} el término temporal se anula:

\[
\nabla\cdot(\mathbf k\nabla T) + Q = 0
\]

con $Q$ la fuente volumétrica de calor [W/m\ensuremath{^{3}}].

\subsection{4. Capacidad calorífica volumétrica}

El producto $\rho c$ [J/(m\ensuremath{^{3}}\ensuremath{\cdot}K)] es lo que gobierna el término transitorio:

\begin{itemize}
  \item $\rho$ --- densidad [kg/m\ensuremath{^{3}}].
  \item $c$ --- calor específico [J/(kg\ensuremath{\cdot}K)].
\end{itemize}

Se declaran por separado ---no como producto--- para que $\rho$ tenga el mismo significado y las mismas unidades que en los materiales mecánicos, y para que el usuario pueda tomar $c$ directamente de tabla.

\textbf{Ambas son opcionales al construir}, siguiendo el precedente de \texttt{density} en el ADR 0008: un análisis \textbf{estacionario} no necesita ninguna de las dos. Si un solver transitorio encuentra \texttt{None}, falla con \texttt{ValueError} identificando el material --- nunca con capacidad cero silenciosa.

La \textbf{difusividad térmica} $\alpha = k/(\rho c)$ [m\ensuremath{^{2}}/s] es magnitud derivada, no parámetro: gobierna la velocidad de propagación del frente térmico y sirve para estimar el paso de tiempo característico. Se expone como propiedad calculada, no se almacena.

\subsection{5. Variables internas}

\textbf{Ninguna.} El modelo es lineal y sin memoria: $\mathbf q$ depende sólo del $\nabla T$ instantáneo. \texttt{PRIMARY\_STATE\_VAR = None}, igual que en \texttt{Elastic1D}/\texttt{Elastic2D}/\texttt{Elastic3D}.

La dependencia $\mathbf k(T)$ ---conductividad función de la temperatura, que sí introduciría no linealidad--- queda fuera de alcance (§out\_of\_scope).

\subsection{6. Unidades y consistencia}

El sistema \textbf{no convierte unidades} (\texttt{Reglas.md §5}); la consistencia es responsabilidad del usuario. En SI: $[k]$ = W/(m\ensuremath{\cdot}K), $[c]$ = J/(kg\ensuremath{\cdot}K), $[\rho]$ = kg/m\ensuremath{^{3}}, $[Q]$ = W/m\ensuremath{^{3}}, $[T]$ = K.

\textbf{Sobre la escala de temperatura}: la conducción pura sólo involucra \textbf{diferencias} de temperatura, así que Kelvin y Celsius dan resultados idénticos. La distinción sí importaría con radiación (donde $T^4$ exige escala absoluta), que está fuera de esta etapa. Se documenta para que la extensión futura no herede una ambigüedad.


\section{Formulación numérica}

\subsection{7. Contribución a la matriz de conductividad}

El elemento integra la forma débil de la ecuación de balance. El material aporta $\mathbf k$ al integrando:

\[
\mathbf K_e = \int_\Omega \mathbf B^\top\,\mathbf k\,\mathbf B\;d\Omega, \qquad \mathbf B = \nabla\mathbf N
\]

donde $\mathbf B$ es $2\times n_{nodos}$ en 2D y $3\times n_{nodos}$ en 3D --- \textbf{el gradiente de las funciones de forma, sin reordenamiento de Voigt}. Es el mismo $\partial N/\partial x$ que ya calculan los kernels mecánicos; cambia sólo cómo se ensambla.

$\mathbf K_e$ es \textbf{simétrica} por serlo $\mathbf k$, luego el despachador algebraico (ADR 0003) puede elegir Cholesky. \texttt{IS\_SYMMETRIC = True}.

\subsection{8. Contribución a la matriz de capacidad}

\[
\mathbf C_e = \int_\Omega \rho c\;\mathbf N^\top\mathbf N\;d\Omega
\]

Estructuralmente idéntica a la matriz de masa consistente, con $\rho c$ en lugar de $\rho$. Por la decisión 5 del alcance se ofrecen ambas formas, \textbf{con \texttt{lumped} por defecto} --- al revés que en dinámica estructural. El motivo es físico y se detalla en la spec del elemento: la capacidad consistente produce oscilaciones espurias ante un frente térmico abrupto y puede violar el principio del máximo de la ecuación de difusión.

\subsection{9. Interfaz de cómputo}

\[
\texttt{compute\_flux}(\nabla T) \;\longrightarrow\; (\mathbf q,\; \mathbf k)
\]

Devuelve el flujo y el tensor de conductividad ---análogo a \texttt{(\ensuremath{\sigma}, C\_tangent)} del contrato mecánico---. Sin \texttt{state\_vars}: el modelo no tiene memoria, y añadir el parámetro por simetría con \texttt{Material} sería ruido. Cuando entre un material térmico no lineal, la firma se ampliará entonces, con el caso real delante.

\subsection{10. Caveats numéricos}

\begin{itemize}
  \item \textbf{Ninguna singularidad}: el modelo es lineal y $\mathbf k$ definida positiva por validación al construir. No hay régimen degenerado análogo al ápice de Drucker-Prager o a $d \to 1$ en daño.
  \item \textbf{Condicionamiento con contraste alto de conductividad}: en un dominio multimaterial con varios órdenes de magnitud entre las $k$ (p. ej. metal contra aislante), $\mathbf K$ queda mal condicionada. Es propiedad del problema físico, no del modelo; se documenta como caveat operativo.
  \item \textbf{Anisotropía severa}: con $k_{\max}/k_{\min}$ muy grande el problema se vuelve fuertemente direccional y la malla debe alinearse razonablemente con los ejes de anisotropía, o la solución degrada. Igual que arriba: propiedad del problema.
\end{itemize}


\section{Contrato de implementación}

\begin{lstlisting}[language=yaml]
name: ThermalConduction
kind: thermal_material
status: validated        # draft → implemented → validated

interface:
  field: temperature     # campo primario que gobierna (1 DOF escalar por nodo)
  flux_dim: null         # 2 ó 3 — no es ClassVar: se fija por instancia al
                         # construir (escalar + dim, o tamaño del tensor k)
  primary_state_var: null
  is_symmetric: true

parameters:
  - { name: k,       type: float | array-like, required: true,
      desc: "Conductividad térmica [W/(m·K)]. Escalar ⇒ isótropo, se expande a k·I.
             Matriz 2×2 o 3×3 ⇒ ortótropo/anisótropo; debe ser simétrica y definida positiva" }
  - { name: c,       type: float, required: false, default: null,
      desc: "Calor específico [J/(kg·K)]. Opcional al construir; obligatorio solo
             en análisis transitorio (mismo criterio que density en ADR 0008)" }
  - { name: density, type: float, required: false, default: null,
      desc: "Densidad [kg/m³]. Opcional al construir; obligatoria solo en transitorio" }

signature:
  compute_flux: "(grad_T: ndarray(d,)) -> (q: ndarray(d,), k: ndarray(d,d))  con d ∈ {2, 3}"
  gradient_kind: "gradiente físico ∇T = [∂T/∂x, ∂T/∂y(, ∂T/∂z)] — sin notación Voigt"

state_variables: []

conventions:
  sign: "q = -k·∇T; el flujo va de mayor a menor temperatura (segunda ley)"
  units: "k [W/(m·K)] · c [J/(kg·K)] · density [kg/m³] · Q [W/m³] · T [K o °C]"
  temperature_scale: "sólo intervienen diferencias de T ⇒ K y °C equivalentes en
                      conducción pura. La distinción importará al entrar radiación (T⁴)"
  derived: "difusividad térmica α = k/(ρ·c) [m²/s] expuesta como propiedad calculada,
            no almacenada"

validity:
  - "k simétrica y definida positiva (autovalores > 0); escalar k > 0"
  - "c > 0 y density > 0 cuando se declaran"
  - "conductividad independiente de la temperatura (modelo lineal)"
  - "medio continuo homogéneo dentro de cada elemento"

out_of_scope:
  - "conductividad dependiente de la temperatura k(T) ⇒ introduce no linealidad; etapa futura"
  - "convección (Robin) y radiación ⇒ condiciones de frontera, no del material; fuera de C1"
  - "cambio de fase / calor latente (Stefan) ⇒ etapa futura"
  - "acoplamiento termomecánico y deformación térmica α·ΔT ⇒ decisión 2 del alcance"
  - "medios porosos, convección forzada interna, transporte de masa"

acceptance:
  verification:
    - name: ley_de_fourier_exacta
      setup: "compute_flux con ∇T arbitrario; k isótropo y k ortótropo"
      expect: "q = -k·∇T exacto componente a componente; segundo retorno = k"
      tol_rel: 1.0e-14
    - name: expansion_escalar_a_tensor
      setup: "construir con k escalar en 2D y en 3D"
      expect: "k interno == k·I (2×2 y 3×3 respectivamente)"
      tol_abs: 1.0e-14
    - name: simetria_y_positividad
      setup: "construir con k escalar y con k tensorial válido"
      expect: "k == k.T exacto; todos los autovalores estrictamente positivos"
      tol_abs: 1.0e-14
    - name: flujo_opuesto_al_gradiente
      setup: "k isótropo; ∇T unitario en varias direcciones"
      expect: "q antiparalelo a ∇T; q·∇T < 0 estricto (disipación termodinámica)"
      tol_rel: 1.0e-14
    - name: anisotropia_desvia_el_flujo
      setup: "k = diag(k1, k2) con k1 ≠ k2; ∇T a 45° de los ejes"
      expect: "q NO es antiparalelo a ∇T; se desvía hacia el eje de mayor conductividad.
               Verificación cuantitativa contra q = -k·∇T calculado a mano"
      tol_rel: 1.0e-12
    - name: difusividad_derivada
      setup: "k, c, density conocidos"
      expect: "alpha == k/(density·c) para el caso isótropo"
      tol_rel: 1.0e-14
    - name: rechazo_inputs_invalidos
      setup: "k ≤ 0; k tensorial no simétrico; k tensorial no definido positivo;
              c ≤ 0; density ≤ 0; k de dimensión inconsistente con el elemento"
      expect: "ValueError con mensaje claro identificando el parámetro en cada caso"
    - name: capacidad_ausente_falla_en_transitorio
      setup: "construir sin c ni density; solicitar matriz de capacidad"
      expect: "ValueError identificando el material y el parámetro faltante,
               nunca capacidad cero silenciosa (precedente ADR 0008)"
    - name: mensaje_de_error_accionable
      setup: "construir sin c y/o sin density; solicitar matriz de capacidad"
      expect: "el mensaje nombra el material, nombra el/los parámetro(s) que faltan
               e indica la salida ('declara ambas' | 'usa un solver estacionario').
               No basta con reportar el atributo ausente"
    - name: estacionario_no_exige_capacidad
      setup: "construir sólo con k; ensamblar la matriz de conductividad"
      expect: "K_e se ensambla sin error; c y density no se consultan en ningún punto
               del camino estacionario"
      tol_rel: 1.0e-14

references:
  - "Incropera F.P., DeWitt D.P. (2007). Fundamentals of Heat and Mass Transfer. Wiley. §2 (ley de Fourier, conductividad tensorial)."
  - "Carslaw H.S., Jaeger J.C. (1959). Conduction of Heat in Solids. Oxford. §1 (formulación general, requisitos sobre k)."
  - "Lewis R.W., Nithiarasu P., Seetharamu K.N. (2004). Fundamentals of the Finite Element Method for Heat and Fluid Flow. Wiley. §3."
  - "Zienkiewicz O.C., Taylor R.L. (2000). The Finite Element Method, Vol. 1. Butterworth-Heinemann. §7 (problemas de campo escalar)."

\end{lstlisting}


\section{Implementación}

\begin{itemize}
  \item Archivo: \texttt{solidum/materials/thermal\_conduction.py}
  \item Clase: \texttt{ThermalConduction}, registrada en \texttt{ThermalMaterialRegistry}
  \item Clase base: \texttt{solidum/core/thermal\_material.py} \ensuremath{\cdot} \texttt{ThermalMaterial}
  \item Tests: \texttt{tests/test\_thermal\_material.py} --- 31 verdes
\end{itemize}

\textbf{Notas de implementación:}

\begin{itemize}
  \item \textbf{\texttt{FLUX\_DIM} es propiedad de instancia, no \texttt{ClassVar}.} A diferencia de \texttt{STRAIN\_DIM} y \texttt{JUMP\_DIM}, la dimensión del flujo no la fija la clase sino el tensor con que se construye: el mismo \texttt{ThermalConduction} sirve en 2D y en 3D. Se deduce de \texttt{conductivity.shape[0]}, de modo que \texttt{k} matricial manda sobre el argumento \texttt{dim} (que sólo se usa para expandir un escalar). El contrato de la spec declara \texttt{flux\_dim: null} por esta razón.
\end{itemize}

\begin{itemize}
  \item \textbf{Validación de simetría con tolerancia escalada.} El criterio es \texttt{atol = 1e-12 \ensuremath{\cdot} max|k|}, adimensional respecto a las unidades del usuario (ADR 0006). Un \texttt{allclose} con tolerancia absoluta fija habría rechazado tensores legítimos expresados en W/(mm\ensuremath{\cdot}K) o aceptado asimetrías reales en W/(m\ensuremath{\cdot}K).
\end{itemize}

\begin{itemize}
  \item \textbf{Positividad vía \texttt{eigvalsh}, no vía el determinante.} El determinante positivo no implica definición positiva en dimensión \ensuremath{\ge} 2 (dos autovalores negativos lo dejan positivo). Se comprueban todos los autovalores, que además es barato para matrices 2\ensuremath{\times}2 y 3\ensuremath{\times}3.
\end{itemize}

\begin{itemize}
  \item \textbf{\texttt{volumetric\_capacity()} vive en la clase base}, no en la subclase: el criterio del ADR 0008 y el formato del mensaje accionable son de familia, no de este material concreto. Cualquier material térmico futuro lo hereda.
\end{itemize}

\begin{itemize}
  \item \textbf{Concordancia gramatical del mensaje de error.} El mensaje flexiona singular/plural según cuántos parámetros falten ("el parámetro 'density', que no se declaró\ldots{} Decláralo" vs "los parámetros 'density' y 'c', que no se declararon\ldots{} Decláralos"). Un mensaje mal concordado resta credibilidad justo cuando el usuario está diagnosticando.
\end{itemize}

\begin{itemize}
  \item \textbf{\texttt{thermal\_diffusivity} rechaza el caso anisótropo} en vez de devolver un valor plausible: con \texttt{k} tensorial la difusividad es un tensor, y un escalar sería una simplificación silenciosa. El mensaje remite a los autovalores de \texttt{conductivity} para estimar los extremos.
\end{itemize}


\section{Diálogo}

\textit{Puntos abiertos para el usuario. Sólo física y alcance.}

\textbf{2026-08-25 --- IA \ensuremath{\to} usuario. Tres puntos:}

\begin{enumerate}
  \item \textasciitilde{}\textasciitilde{}\textbf{Nombre del material.}\textasciitilde{}\textasciitilde{} ✅ \textbf{Resuelto 2026-08-25 --- \texttt{ThermalConduction}} (decisión del usuario). El nombre inicial \texttt{ThermalIsotropic} describía mal el contrato, que admite tensor completo; \texttt{ThermalConduction} nombra la física y no compromete la simetría del tensor. Un futuro material térmico no lineal ---conductividad $\mathbf k(T)$, cambio de fase--- entrará como nombre propio distinto, no como variante de éste.
\end{enumerate}

\begin{enumerate}
  \item \textasciitilde{}\textasciitilde{}\textbf{\texttt{c} y \texttt{density} opcionales.}\textasciitilde{}\textasciitilde{} ✅ \textbf{Resuelto 2026-08-25 --- se traslada el criterio del ADR 0008} (decisión del usuario). Ambas son opcionales al construir: un análisis \textbf{estacionario} sólo usa $\mathbf k$, y exigir propiedades que no entran en ninguna ecuación obligaría a inventar valores ficticios. El fallo silencioso ---el riesgo real de un parámetro opcional--- está cerrado por diseño: cuando un consumidor transitorio las encuentra \texttt{None}, lanza \texttt{ValueError} identificando el material.
\end{enumerate}

   \textbf{Requisito añadido en la decisión}: el mensaje debe ser \textbf{accionable}, no sólo descriptivo --- decir qué hacer, no sólo qué falta. Forma canónica:

   > \texttt{ThermalConduction (id=1): el análisis transitorio requiere 'c' y 'density'. Declara ambas en el material, o usa un solver estacionario (LinearSolver).}

   Es el espíritu de "validación temprana con mensajes claros" de \texttt{Reglas.md §1} aplicado al caso.

\begin{enumerate}
  \item \textasciitilde{}\textasciitilde{}\textbf{Fuente volumétrica $Q$.}\textasciitilde{}\textasciitilde{} ✅ \textbf{Resuelto 2026-08-25 --- $Q$ es carga aplicada en el elemento, no propiedad del material} (decisión del usuario).
\end{enumerate}

   El material define la \textbf{ley constitutiva}, no las acciones sobre el medio; es el mismo reparto que en la familia mecánica, donde \texttt{Elastic2D} no sabe nada del peso propio ni de las fuerzas de cuerpo.

   La razón que decide no es la simetría arquitectural sino la \textbf{variabilidad espacial y temporal}: la generación por hidratación del hormigón depende de la madurez, y el efecto Joule de la corriente. Ninguna es constante del medio. Modelada como carga, $Q$ puede variar por zona y por paso sin duplicar materiales; como propiedad del material, cada valor distinto exigiría un material propio.

   Consecuencia para la spec del elemento: \texttt{compute\_body\_source(Q)} integra $\int_\Omega Q\,\mathbf N^\top d\Omega$, paralelo a \texttt{compute\_body\_load(b)} del mecánico. Se especifica allí, no aquí.

   No cierra la puerta a un atajo futuro ---un material que declare generación propia--- si aparece el caso de uso que lo justifique.


\newpage

\chapter{Esquemas de Solución — Estáticos}

\section{ArcLengthSolver}



\textit{Spec \textbf{retroactiva}: el solver existe desde la fase de daño/softening y está validado por tests de snap-through, snap-back y daño post-pico. Esta spec lo documenta sin cambiar comportamiento --- H-5.3 de la auditoría 2026-05-18.}



\section{Especificación física}

\subsection{0. Descripción general}

Análisis estático \textbf{no lineal} que traza curvas de equilibrio $\mathbf U(\lambda)$ atravesando puntos límite (snap-through, snap-back) donde $d\lambda/dU$ cambia de signo. A diferencia de \texttt{NonlinearSolver} (control de carga, $\lambda$ creciente impuesto), el arc-length \textbf{trata $\lambda$ como incógnita adicional} y restringe el avance combinado en $(U, \lambda)$ a un arco fijo $dl$ en el espacio aumentado. Esto permite seguir la respuesta postcrítica de estructuras con softening, pandeo geométrico, y daño con rama descendente.

Algoritmo: \textbf{Crisfield cilíndrico} (1981) --- restricción cuadrática $\lVert\Delta\mathbf U\rVert^2 = dl^2$, $\lambda$ libre.

\subsection{1. Ecuación de equilibrio resuelta}

Forma fuerte:

\[
\mathbf F_{\text{int}}(\mathbf U) \;=\; \lambda\,\mathbf F_{\text{ext}}^{\text{ref}}
\]

con $\lambda$ incógnita, sujeta a la restricción cilíndrica de Crisfield en cada paso:

\[
\lVert\Delta\mathbf U^{\text{paso}}\rVert^2 \;=\; dl^2
\]

(versión \textbf{cilíndrica}: no incluye $\Delta\lambda^2$, equivalente a tomar $\psi = 0$ en el parámetro de escala carga-desplazamiento; la versión esférica $\lVert\Delta\mathbf U\rVert^2 + \psi^2\Delta\lambda^2\lVert\mathbf F_{\text{ext}}^{\text{ref}}\rVert^2 = dl^2$ no está implementada).

\subsection{2. Condiciones de contorno}

Idénticas a \texttt{NonlinearSolver}: Dirichlet (homogéneo / no homogéneo) y MPC vía ADR 0004. Las condiciones Dirichlet no homogéneas se aplican proporcionalmente a $\lambda$.

\subsection{3. Salidas físicas}

\begin{itemize}
  \item $\mathbf U$ final en $\lambda \to \lambda_{\max}$ (o cuando se alcance \texttt{max\_steps}).
  \item Historia disponible vía \texttt{step\_callback(step, U, lambda)}.
  \item Trazas $\lambda(U_i)$ para visualizar snap-through/snap-back.
\end{itemize}


\section{Formulación numérica}

\subsection{4. Esquema operativo}

\begin{lstlisting}
λ = 0
dl = initial_dl
mientras λ < max_lambda y step < max_steps:
    # Predictor tangente
    K_t, F_int = assemble_non_linear_system(U_iter)
    δu_t = K_t^{-1} · F_ext^ref          # desplazamiento tangente unitario
    sign = sign(δU_prev · δu_t)          # evitar revertir dirección
    Δλ_pred = sign · dl / ‖δu_t‖
    ΔU_pred = Δλ_pred · δu_t

    # Corrector iterativo
    para iter en range(max_iter):
        K_t, F_int = assemble_non_linear_system(U_iter)
        R = (λ + Δλ)·F_ext^ref - F_int
        δu_R = K_t^{-1} · R               # corrección por residuo
        δu_t = K_t^{-1} · F_ext^ref       # tangente (refresh)

        # Restricción cuadrática de Crisfield
        # ‖ΔU_new + ddλ·δu_t‖² = dl²
        # ⇒ resolver ecuación cuadrática en ddλ, elegir raíz por menor ángulo
        ddλ = raiz_de_menor_angulo(...)
        ΔU += δu_R + ddλ · δu_t
        Δλ += ddλ
        U_iter = U_current + ΔU
        si converge (ADR 0007):
            commit_all_states()
            λ += Δλ
            ajustar_dl_adaptativamente(iter)
            break
    si no converge:
        dl /= dl_shrink_factor

\end{lstlisting}

\subsection{5. Predictor / corrector}

\textbf{Predictor tangente}: $\delta\mathbf u_t = \mathbf K_t^{-1}\,\mathbf F_{\text{ext}}^{\text{ref}}$ (vector que mide la pendiente local del camino de equilibrio). $\Delta\lambda_{\text{pred}} = \pm dl / \lVert\delta\mathbf u_t\rVert$ con signo determinado por el producto escalar contra el incremento previo $\Delta\mathbf U^{\text{paso prev}}$ --- así no se retrocede en el camino tras pasar un punto límite.

\textbf{Corrector} (Newton modificado para arc-length, Crisfield): cada iteración resuelve dos sistemas (residuo y tangente), combina con $dd\lambda$ que satisface la restricción cuadrática:

\[
\lVert\Delta\mathbf U + dd\lambda\,\delta\mathbf u_t\rVert^2 = dl^2
\]

Ecuación cuadrática en $dd\lambda$ con dos raíces. \textbf{Selección de raíz}: la que produce \textbf{menor ángulo} con el incremento acumulado $\Delta\mathbf U$ (evita pivotear el sentido del arco).

\subsection{6. Criterio de convergencia (ADR 0007)}

Idéntico a \texttt{NonlinearSolver} --- patrón dual fuerza + desplazamiento con escala autoderivada. Compartido vía \texttt{ConvergenceCriterion}. La calibración del solver usa $\lVert\mathbf F_{\text{ext}}^{\text{ref}}\rVert$ como referencia de fuerza (no $\lambda\,\mathbf F_{\text{ext}}^{\text{ref}}$, porque $\lambda$ varía).

\subsection{7. Imposición de Dirichlet y MPC}

Idéntico al \texttt{NonlinearSolver} --- \texttt{Assembler.reduce(...)} con \texttt{U\_current} y \texttt{load\_factor=lambda\_iter}.

\subsection{8. Backend algebraico}

Régimen \textbf{postcrítico}: $\mathbf K_t$ puede ser indefinida (autovalores negativos tras bifurcación). El solver fuerza \texttt{is\_positive\_definite=False} al construir el \texttt{linalg}, así el despachador ADR 0003 elige LU automáticamente. Si el usuario fuerza \texttt{linear\_algebra="cholesky"} desde YAML y $\mathbf K_t$ se vuelve indefinida en algún paso, hay fallback automático a LU con warning.

\subsection{9. Adaptatividad --- ajuste de $dl$}

Política de auto-ajuste basada en iteraciones del paso:

\begin{itemize}
  \item Converge en $<$ \texttt{dl\_grow\_iter\_threshold} iteraciones \ensuremath{\Rightarrow} $dl \to dl \cdot $ \texttt{dl\_grow\_factor} (cota: \texttt{initial\_dl \ensuremath{\cdot} dl\_max\_factor}).
  \item Converge en $>$ \texttt{dl\_shrink\_iter\_threshold} iteraciones \ensuremath{\Rightarrow} $dl \to dl \cdot $ \texttt{dl\_shrink\_factor}.
  \item No converge en \texttt{max\_iter} \ensuremath{\Rightarrow} $dl /= 2$ y reintentar el paso.
  \item Cota inferior: $dl < $ \texttt{ARCLENGTH\_MIN\_DL\_FACTOR \ensuremath{\cdot} initial\_dl} \ensuremath{\Rightarrow} aborto.
\end{itemize}

\subsection{10. Caveats numéricos}

\begin{itemize}
  \item \textbf{Cancelación o bifurcación verdadera}: en bifurcaciones simétricas múltiples, la selección de raíz por menor ángulo puede saltar a un camino paralelo. Para análisis de bifurcación auténtica usar perturbación + post-procesado modal.
  \item \textbf{Raíces imaginarias}: si el discriminante de la cuadrática es negativo (paso demasiado grande, mal condicionamiento severo), el solver aborta el paso y bisecta $dl$.
  \item \textbf{Último paso}: cuando $\lambda + \Delta\lambda_{\text{pred}} \ge \lambda_{\max}$, el solver fija $\Delta\lambda = \lambda_{\max} - \lambda$ y corrige solo en desplazamientos (Newton puro), sin restricción cuadrática. Permite cerrar exactamente en el target.
  \item \textbf{Softening con penalty cohesivo stiff} (CST\_Embedded2D): no se atraviesa la transición elástico\ensuremath{\to}softening por la quasi-singularidad de $\mathbf K_t$ cerca de $\kappa_0$. Limitación específica al embedded.
  \item \textbf{Diagnóstico de bifurcación por inercia} (Sturm sequence): placeholder en \texttt{\_negative\_pivots()} --- retorna \texttt{None} hasta que se implemente un backend LDL\ensuremath{^{\top}} verdadero (Bunch-Kaufman) con conteo de pivots negativos.
\end{itemize}


\section{Contrato de implementación}

\begin{lstlisting}[language=yaml]
name: ArcLengthSolver
kind: solver
status: validated

interface:
  yaml_type: arclength
  output: SolveResult
  pipeline_kind: static

parameters:
  - { name: convergence,              type: ConvergenceCriterion, required: false, default: "ConvergenceCriterion()",
      desc: "Política dual fuerza+desplazamiento (ADR 0007), compartida con NonlinearSolver" }
  - { name: max_iter,                 type: int,   required: false, default: 20,
      desc: "Iteraciones Newton (corrector) máximas por paso" }
  - { name: max_lambda,               type: float, required: false, default: 1.0,
      desc: "Factor de carga objetivo" }
  - { name: initial_dl,               type: float, required: false, default: 0.1,
      desc: "Longitud de arco inicial" }
  - { name: max_steps,                type: int,   required: false, default: 100,
      desc: "Cota dura sobre el número de pasos" }
  - { name: dl_grow_factor,           type: float, required: false, default: 1.5 }
  - { name: dl_max_factor,            type: float, required: false, default: 5.0,
      desc: "dl máximo = initial_dl · dl_max_factor" }
  - { name: dl_shrink_factor,         type: float, required: false, default: 0.6 }
  - { name: dl_grow_iter_threshold,   type: int,   required: false, default: 4 }
  - { name: dl_shrink_iter_threshold, type: int,   required: false, default: 8 }
  - { name: linear_algebra,           type: str,   required: false, default: "auto",
      desc: "'auto' (LU para postcrítico), 'cholesky', 'lu'. En arc-length el default fuerza LU por indefinitud potencial" }

requirements:
  - "Modelo con potencial snap-through/snap-back (geometría inestable, softening del material)"
  - "Apoyos suficientes; sin modos rígidos"
  - "dl inicial razonable (orden de magnitud del desplazamiento característico del primer paso)"

conventions:
  units: "heredadas del modelo (ADR 0008)"
  stability: "incondicional respecto a puntos límite; bifurcaciones complejas pueden saltar de rama"

out_of_scope:
  - "Variante esférica del arc-length (ψ ≠ 0) — sólo cilíndrica implementada"
  - "Análisis dinámico ⇒ usar NewmarkSolver y derivados"
  - "Continuación de bifurcaciones múltiples — requiere perturbación explícita por el usuario"
  - "Buckling lineal (eigenproblema) ⇒ futuro BucklingSolver"
  - "Sign-of-pivot tracking (Sturm sequence) — placeholder hasta LDLᵀ verdadero"

acceptance:
  verification:
    - name: barra_traccion_elastica_reproduce_caso_lineal
      setup: "Truss2D elástica, carga axial; arc-length avanza hasta λ=1"
      expect: "U(λ=1) = U_lineal exacto; equilibrio en cada paso"
      tol_rel: 1.0e-10

    - name: snap_through_lee_frame
      setup: "frame Lee con dos barras corotacionales en V, carga vertical en el vértice"
      expect: "se atraviesa el punto límite; respuesta U-λ coincide con Crisfield 1991 Vol.1 Fig. 9.5"
      tol_rel: 1.0e-3

    - name: snap_back_truss_inestable
      setup: "armadura con elemento corotacional cuya rama descendente requiere d λ/dU < 0"
      expect: "se atraviesa el snap-back; signo de Δλ cambia coherentemente"
      tol_rel: 1.0e-3

    - name: damage_softening_post_pico
      setup: "Quad4 con IsotropicDamage2D, tracción monotónica que entra a rama descendente"
      expect: "rama post-pico trazada; convergencia paso a paso"
      tol_rel: 1.0e-3

    - name: ajuste_dl_adaptativo
      setup: "trayectoria con tramos elásticos rápidos y zona post-pico lenta"
      expect: "dl crece en tramos fáciles, se reduce cerca del pico; converge globalmente"

    - name: ultimo_paso_cierra_en_target
      setup: "max_lambda = 0.7 sobre estructura elástica"
      expect: "λ_final = 0.7 exacto (no por overshoot)"
      tol_abs: 1.0e-12

  specific:
    - name: seleccion_raiz_no_revierte
      setup: "paso tras un snap-back con dos raíces de signos opuestos"
      expect: "se elige la raíz que continúa el camino (menor ángulo con incremento previo)"

    - name: fallback_cholesky_lu_indefinitud
      setup: "linear_algebra='cholesky' forzado, modelo con punto límite"
      expect: "warning + LU al detectar no-PD; análisis continúa"

references:
  - "Crisfield M.A. (1981). A fast incremental/iterative solution procedure that handles snap-through. Computers and Structures 13, 55-62."
  - "Crisfield M.A. (1991). Non-linear Finite Element Analysis of Solids and Structures, Vol. 1. Wiley. §9 (arc-length cilíndrico)."
  - "Riks E. (1979). An incremental approach to the solution of snapping and buckling problems. Int. J. Solids and Structures 15, 529-551."
  - "de Borst R., Sluys L.J. (1999). Computational Methods in Non-linear Solid Mechanics. TU Delft. §3.5."
  - "ADR 0003, ADR 0007."

\end{lstlisting}


\section{Implementación}

\begin{itemize}
  \item \textbf{Archivo}: solidum/math/solvers/arclength.py.
  \item \textbf{Clase}: \texttt{ArcLengthSolver}, registrada vía \texttt{@SolverRegistry.register} con \texttt{PIPELINE\_KIND = "static"}.
  \item \textbf{Restricción cuadrática y selección de raíz}: implementada en \texttt{solve}. Comparación de ángulos \texttt{theta1 = \ensuremath{\Delta}U\_iter \ensuremath{\cdot} (\ensuremath{\Delta}U\_new + ddl1\ensuremath{\cdot}\ensuremath{\delta}u\_t)}, idem \texttt{theta2}, seleccionar el mayor \ensuremath{\to} menor ángulo con la dirección previa.
  \item \textbf{\texttt{\_make\_linalg}}: fuerza \texttt{is\_positive\_definite=False} independientemente de la simetría del dominio (régimen postcrítico).
  \item \textbf{Entrypoint público}: \texttt{solidum.run\_static(model, solver="arclength", ...)}.
  \item \textbf{Tests}:
  \item tests/test\_arclength.py \ensuremath{\cdot} snap-through Lee frame, snap-back truss, recuperación lineal.
  \item tests/test\_solver\_robustness.py \ensuremath{\cdot} \texttt{test\_arc\_length\_traverses\_damage\_softening}.
\end{itemize}


\section{Diálogo}

\begin{itemize}
  \item \textbf{2026-05-19} \ensuremath{\cdot} Spec creada retroactivamente para cerrar el hueco H-5.3. Solver anterior a la convención de specs. La spec recoge la implementación tal como está al cierre de la sesión de saneamiento post-auditoría.
\end{itemize}


\newpage

\section{DissipationArcLengthSolver}



\textit{\textbf{Spec corta tipo extensión} (Reglas.md §4). Variante de \texttt{ArcLengthSolver} que \textbf{solo cambia la restricción del paso}: pasa de cilíndrica ($\lVert\Delta\mathbf U\rVert^2 = dl^2$) a basada en \textbf{disipación de energía} ($g(\Delta\mathbf U, \Delta\lambda) = \tau$). Reusa todo el plumbing del padre (predictor tangente, corrector Newton, adaptatividad de paso, Dirichlet/MPC, backend algebraico). Sin ADR nuevo: reusa el contrato arquitectural ya aceptado del arc-length.}

\textit{> Motivación: el \texttt{ArcLengthSolver} cilíndrico no atraviesa la transición elástico\ensuremath{\to}softening con penalty stiff (\texttt{K\_e} del embedded discontinuity, ADR 0010 fase 4). La restricción cuadrática queda casi degenerada cerca del pico --- el discriminante de la cuadrática tiende a cero, el solver bisecta indefinidamente y aborta. La restricción de disipación es \textbf{lineal en $(\Delta\mathbf U, \Delta\lambda)$} y se activa \textbf{proporcionalmente a la energía consumida} por la grieta cohesiva, navegando la rama post-pico de forma estable. Esta variante desbloquea el benchmark de fractura faithful Van Vliet (\texttt{docs/specs/CST\_Embedded2D.md} §"Out of scope"), el balance energético end-to-end de $G_F$ (hueco abierto en [[project\_validacion\_fase\_b]]), y validación cuantitativa de cohesivos en pipeline real.}



\section{Especificación física}

\subsection{0. Descripción general}

Solver estático no lineal con \textbf{traza por disipación}: la magnitud que se mantiene constante de paso a paso ya no es la longitud euclídea $\lVert\Delta\mathbf U\rVert$ sino la \textbf{disipación incremental de energía} $\Delta E_d$ asociada al paso. Adecuado para:

\begin{itemize}
  \item Fractura cohesiva (Mode-I y, en el futuro, mixto) con penalty \texttt{K\_e} stiff.
  \item Daño con softening severo donde la rama post-pico tiene pendiente casi vertical en el plano $(U, \lambda)$.
  \item Localización de deformación plástica con bandas estrechas.
\end{itemize}

Régimen elástico (sin disipación): la restricción degenera a $g = 0$. El solver detecta este caso y \textbf{conmuta automáticamente} a arc-length cilíndrico durante el régimen elástico, retornando a control por disipación cuando se detecta inicio de disipación (\texttt{\ensuremath{\Delta}\ensuremath{\omega} > 0} en daño, $\Delta\alpha > 0$ en plasticidad, $\Delta\kappa > \kappa_0$ en cohesivo).

\subsection{1. Ecuación de equilibrio resuelta}

Idéntica al padre:

\[
\mathbf F_{\text{int}}(\mathbf U) = \lambda\,\mathbf F_{\text{ext}}^{\text{ref}}
\]

con $\lambda$ incógnita. La \textbf{restricción del paso} cambia respecto al padre.

\subsubsection{Restricción de disipación (Gutiérrez 2004, Verhoosel et al. 2009)}

La energía disipada por paso, en problemas conservativos en el componente elástico, es:

\[
\Delta E_d = \frac{1}{2}\left(\lambda_n\,\mathbf F_{\text{ext}}^{\text{ref}T}\,\Delta\mathbf U \;-\; \Delta\lambda\,\mathbf F_{\text{ext}}^{\text{ref}T}\,\mathbf U_n\right)
\]

donde $(\mathbf U_n, \lambda_n)$ son los valores committed al inicio del paso y $(\Delta\mathbf U, \Delta\lambda)$ los incrementos del paso. La restricción es:

\[
g(\Delta\mathbf U, \Delta\lambda) \;\equiv\; \frac{1}{2}\left(\lambda_n\,\mathbf F_{\text{ext}}^{\text{ref}T}\,\Delta\mathbf U \;-\; \mathbf U_n^T\,\mathbf F_{\text{ext}}^{\text{ref}}\,\Delta\lambda\right) \;=\; \tau
\]

con $\tau$ el \textbf{incremento de disipación objetivo} (parámetro del solver, papel análogo al \texttt{initial\_dl} del padre).

\textbf{Linealidad}: $g$ es lineal en $(\Delta\mathbf U, \Delta\lambda)$. Por eso el corrector tiene \textbf{una sola raíz}, no dos como en la cuadrática cilíndrica --- eliminando la selección por menor ángulo y la patología de raíces imaginarias en pasos grandes.

\subsection{2. Condiciones de contorno}

Idénticas al padre (\texttt{ArcLengthSolver} §2). Dirichlet homogéneo/no homogéneo y MPC vía ADR 0004.

\subsection{3. Salidas físicas}

Idénticas al padre --- \texttt{SolveResult} con $\mathbf U$ final y \texttt{step\_callback(step, U, \ensuremath{\lambda})} para historia. Adicionalmente, se expone vía el callback el valor de disipación acumulada $\sum_n \Delta E_d^{(n)}$ --- útil para validar \textbf{$G_F$ end-to-end} integrando $\int_{\Gamma_d} t\cdot d[\![u]\!] \approx \sum_n \Delta E_d^{(n)}$.


\section{Formulación numérica --- solo lo que cambia respecto al padre}

\subsection{4. Esquema operativo}

Estructura idéntica al padre con tres diferencias:

\begin{lstlisting}
λ = 0
τ = initial_tau          # análogo a initial_dl, pero en unidades de energía
modo = "cylindrical"      # arranque en cilíndrico mientras g ≈ 0
mientras λ < max_lambda y step < max_steps:
    # Predictor tangente (idéntico al padre)
    δu_t = K_t^{-1} · F_ext^ref
    sign = sign(δU_prev · δu_t)
    si modo == "cylindrical":
        Δλ_pred = sign · dl / ‖δu_t‖
    sino:                  # modo == "dissipation"
        # Predictor por disipación: Δλ_pred derivado de τ
        Δλ_pred = sign · τ / (½·|λ_n·(F_ext^ref·δu_t) − F_ext^ref·U_n|)
    ΔU_pred = Δλ_pred · δu_t

    # Corrector iterativo
    para iter en range(max_iter):
        K_t, F_int = assemble(...)
        R = (λ + Δλ)·F_ext^ref − F_int
        δu_R = K_t^{-1} · R
        δu_t = K_t^{-1} · F_ext^ref

        si modo == "cylindrical":
            # Restricción cuadrática, idéntica al padre (ArcLengthSolver §5)
            ddλ = raiz_de_menor_angulo(...)
        sino:               # modo == "dissipation"
            # Restricción lineal de Gutiérrez: g(ΔU_new, Δλ_new) = τ
            # ΔU_new = ΔU + δu_R + ddλ·δu_t
            # g(ΔU_new, Δλ + ddλ) = τ  ⇒  ddλ explícito:
            num = τ − ½·(λ_n·F_ext^ref·(ΔU+δu_R) − (Δλ)·F_ext^ref·U_n)
            den = ½·(λ_n·F_ext^ref·δu_t − F_ext^ref·U_n)
            ddλ = num / den
        ΔU += δu_R + ddλ·δu_t
        Δλ += ddλ
        U_iter = U_current + ΔU
        si converge: break

    si converge:
        commit_all_states()
        λ += Δλ
        actualizar_modo(...)            # ← switching cilíndrico ↔ disipación
        ajustar_tau_adaptativamente(iter)
    sino:
        τ /= shrink_factor

\end{lstlisting}

\subsection{5. Predictor / corrector --- punto clave del cambio}

\subsubsection{Switching automático cilíndrico \ensuremath{\leftrightarrow} disipación}

El solver mantiene un atributo \texttt{modo} con dos estados:

\begin{itemize}
  \item \texttt{"cylindrical"}: arc-length cilíndrico estándar; restricción $\lVert\Delta\mathbf U\rVert^2 = dl^2$. \textbf{Régimen elástico}, sin disipación activa.
  \item \texttt{"dissipation"}: arc-length por energía; restricción $g = \tau$. \textbf{Régimen disipativo}.
\end{itemize}

\textbf{Criterio de transición} (Verhoosel et al. 2009 §3.2):

\begin{itemize}
  \item Inicio de paso en \texttt{"cylindrical"}: tras commit, calcular $\Delta E_d$ del paso recién cerrado. Si $\Delta E_d > \varepsilon_{\text{diss}}$ (umbral pequeño, p. ej. $10^{-12}\cdot \lVert\mathbf F_{\text{ext}}^{\text{ref}}\rVert\cdot \lVert\mathbf U_n\rVert$) \ensuremath{\Rightarrow} siguiente paso entra en \texttt{"dissipation"}. Inicializa $\tau \leftarrow \Delta E_d$ del último paso para continuidad.
  \item Inicio de paso en \texttt{"dissipation"}: si el paso anterior tuvo iteraciones $\le$ \texttt{tau\_relax\_iter\_threshold} \textbf{y} el "predictor de disipación" arrojaría $\Delta\lambda$ desproporcionadamente grande (señal de retorno a elástico tras la rama post-pico), revertir a \texttt{"cylindrical"}. Patológico pero posible si hay descarga global tras grieta saturada.
\end{itemize}

\subsubsection{Sign-of-pivot tracking (mejora sobre el padre)}

\texttt{ArcLengthSolver.\_negative\_pivots()} es placeholder hoy. Esta variante \textbf{lo implementa} mediante factorización LDL\ensuremath{^{\top}} (Bunch-Kaufman) del backend algebraico, exponiendo el conteo de pivots negativos como diagnóstico de bifurcación. Uso interno: detectar paso del punto límite (cambio en el signo del primer pivot indefinido) y orientar la dirección del predictor con respecto al signo previo, no sólo al producto escalar con $\Delta\mathbf U_{\text{prev}}$.

Esto resuelve el caso patológico de bifurcación simétrica donde el "menor ángulo" elige una rama paralela en lugar de continuar el camino físico --- caveat documentado en \texttt{ArcLengthSolver} §10.

\subsection{6. Criterio de convergencia (ADR 0007)}

Idéntico al padre --- patrón dual fuerza + desplazamiento, escala autoderivada.

\subsection{7. Imposición de Dirichlet y MPC}

Idéntico al padre. La restricción de disipación opera sobre DOFs libres después de la reducción Dirichlet/MPC.

\subsection{8. Backend algebraico}

Idéntico al padre (\texttt{is\_positive\_definite=False}, despachador LU). \textbf{Adición}: para sign-of-pivot tracking auténtico se requiere LDL\ensuremath{^{\top}}; la implementación inicial puede usar la firma de la factorización LU (signo del determinante) como aproximación, con upgrade a LDL\ensuremath{^{\top}} Bunch-Kaufman cuando se priorice.

\subsection{9. Adaptatividad --- ajuste de $\tau$}

Análoga a la adaptatividad de $dl$ del padre, con $\tau$ en lugar de $dl$:

\begin{itemize}
  \item Converge en $<$ \texttt{tau\_grow\_iter\_threshold} iteraciones \ensuremath{\Rightarrow} $\tau \to \tau \cdot $ \texttt{tau\_grow\_factor} (cota: \texttt{initial\_tau \ensuremath{\cdot} tau\_max\_factor}).
  \item Converge en $>$ \texttt{tau\_shrink\_iter\_threshold} iteraciones \ensuremath{\Rightarrow} $\tau \to \tau \cdot $ \texttt{tau\_shrink\_factor}.
  \item No converge en \texttt{max\_iter} \ensuremath{\Rightarrow} $\tau /= 2$ y reintentar.
  \item Cota inferior: $\tau < $ \texttt{DISS\_MIN\_TAU\_FACTOR \ensuremath{\cdot} initial\_tau} \ensuremath{\Rightarrow} aborto.
\end{itemize}

En modo \texttt{"cylindrical"} durante régimen elástico inicial, la adaptatividad opera sobre $dl$ (idéntica al padre).

\subsection{10. Caveats numéricos}

Heredados del padre (\texttt{ArcLengthSolver} §10) más los específicos:

\begin{itemize}
  \item \textbf{Régimen disipativo no estrictamente monótono}: si el material vuelve transitoriamente a régimen elástico (descarga local), $\Delta E_d$ puede ser negativo. El solver tolera este caso pero conmuta a \texttt{"cylindrical"} si la disipación neta del paso cae por debajo del umbral.
  \item \textbf{Inicialización de $\tau$}: debe ser del orden de la disipación esperada por paso (típicamente fracción pequeña de $G_F \cdot l_d$ por elemento agrietado). Demasiado pequeño \ensuremath{\Rightarrow} pasos numerosos; demasiado grande \ensuremath{\Rightarrow} saltos sobre la rama post-pico.
  \item \textbf{Disipación de origen no físico}: si el modelo numérico introduce disipación espuria (algorithmic damping en transitorio, regularización no consistente), $\Delta E_d$ contiene esa componente. \textbf{Este solver es estático}: si el usuario invoca \texttt{DissipationArcLengthSolver} sobre un modelo con material no disipativo (elástico puro), el modo permanece en \texttt{"cylindrical"} indefinidamente --- comportamiento equivalente al padre. No es un caveat operacional, es la definición del switching.
  \item \textbf{MPC con offset $\mathbf g$ no homogéneo}: la restricción de disipación se evalúa sobre DOFs libres reducidos; si las MPC introducen offset, $g$ se evalúa sobre el sistema reducido y la disipación medida es consistente con la energía aportada por las MPC (tratamiento idéntico al residuo en el padre).
\end{itemize}


\section{Contrato de implementación}

\begin{lstlisting}[language=yaml]
name: DissipationArcLengthSolver
kind: solver
status: implemented      # draft → implemented → validated. Validación parcial — ver §"Implementación".

parent_spec: ArcLengthSolver

interface:
  yaml_type: dissipation_arclength
  output: SolveResult
  pipeline_kind: static

parameters:
  # Heredados del padre (todos con misma semántica):
  - { name: convergence,              type: ConvergenceCriterion, required: false, default: "ConvergenceCriterion()" }
  - { name: max_iter,                 type: int,   required: false, default: 20 }
  - { name: max_lambda,               type: float, required: false, default: 1.0 }
  - { name: max_steps,                type: int,   required: false, default: 100 }
  - { name: initial_dl,               type: float, required: false, default: 0.1,
      desc: "Longitud de arco inicial — usado en modo cilíndrico (régimen elástico)" }
  - { name: dl_grow_factor,           type: float, required: false, default: 1.5 }
  - { name: dl_max_factor,            type: float, required: false, default: 5.0 }
  - { name: dl_shrink_factor,         type: float, required: false, default: 0.6 }
  - { name: dl_grow_iter_threshold,   type: int,   required: false, default: 4 }
  - { name: dl_shrink_iter_threshold, type: int,   required: false, default: 8 }
  - { name: linear_algebra,           type: str,   required: false, default: "auto" }

  # Específicos de esta variante:
  - { name: initial_tau,              type: float, required: true,
      desc: "Incremento de disipación objetivo (unidades de energía: F·U). Calibrar al orden de G_F·l_d/n_pasos_post_pico" }
  - { name: tau_grow_factor,          type: float, required: false, default: 1.5 }
  - { name: tau_max_factor,           type: float, required: false, default: 5.0 }
  - { name: tau_shrink_factor,        type: float, required: false, default: 0.6 }
  - { name: tau_grow_iter_threshold,  type: int,   required: false, default: 4 }
  - { name: tau_shrink_iter_threshold,type: int,   required: false, default: 8 }
  - { name: dissipation_threshold,    type: float, required: false, default: 1.0e-12,
      desc: "Umbral relativo bajo el cual ΔE_d se considera 0 (mantiene modo cilíndrico)" }

requirements:
  - "Modelo con material disipativo (plasticidad, daño, cohesivo) — en modelo puramente elástico el solver opera equivalente a ArcLengthSolver cilíndrico"
  - "F_ext^ref no nulo (la restricción de disipación se evalúa sobre F_ext^ref·U)"
  - "Heredados del padre: apoyos suficientes; sin modos rígidos; dl_inicial razonable"

conventions:
  units: "heredadas del modelo (ADR 0008); τ en unidades de energía consistentes"
  stability: "incondicional respecto a puntos límite y rama post-pico de softening severo. La estabilidad ante bifurcaciones múltiples mejora vs el padre por el sign-of-pivot tracking"

out_of_scope:
  - "Bifurcaciones múltiples simétricas — la detección por inercia desambigua bifurcaciones simples; las múltiples siguen requiriendo perturbación explícita"
  - "Análisis dinámico — la formulación de disipación es válida solo en régimen cuasi-estático conservativo"
  - "Modo II y mixto I-II en cohesivo — heredado del scope ADR 0010 fases F+"
  - "3D — coherente con el scope global del proyecto en 2026-05"

acceptance:
  verification:
    - name: recuperacion_del_arclength_cilindrico_en_elastico
      setup: "Modelo elástico puro (sin disipación); ejecutar DissipationArcLengthSolver y ArcLengthSolver con los mismos parámetros y misma malla"
      expect: "El modo permanece 'cylindrical'; trayectoria U(λ) coincide a paridad de bits con el padre"
      tol_rel: 1.0e-12

    - name: snap_through_lee_frame
      setup: "Mismo benchmark que ArcLengthSolver §acceptance — frame de Lee, dos barras corotacionales"
      expect: "Atraviesa el punto límite; U-λ coincide con Crisfield 1991 Fig. 9.5"
      tol_rel: 1.0e-3

    - name: damage_softening_post_pico
      setup: "Quad4 con IsotropicDamage2D, tracción que entra en rama descendente"
      expect: "Modo conmuta a 'dissipation' al detectar Δω>0; rama post-pico trazada; coincide con ArcLengthSolver del padre dentro de tolerancia (ambos resuelven el mismo problema)"
      tol_rel: 1.0e-3

    - name: cohesive_embedded_softening_postpeak       # benchmark CRÍTICO — desbloquea fase 4 ADR 0010
      setup: "Probeta rectangular con CST_Embedded2D + CohesiveDamageIsotropic; tracción uniaxial monotónica que activa Rankine y entra en rama post-pico con penalty K_e stiff (1e13)"
      expect: "El solver atraviesa la transición elástico→softening (régimen donde el ArcLengthSolver cilíndrico aborta); cierra el ciclo de softening; daño cohesivo committed alcanza saturación"
      tol_rel: 1.0e-3

    - name: balance_energetico_GF
      setup: "Mismo modelo cohesivo+embedded del test anterior; integrar ΣΔE_d sobre toda la historia hasta saturación (ω→1) y comparar con G_F · longitud_grieta_abierta del material"
      expect: "Σ ΔE_d ≈ G_F · l_d a tolerancia de discretización temporal (1% típico)"
      tol_rel: 1.0e-2

    - name: convergencia_temporal_tau
      setup: "Reducir initial_tau en factor 2 sucesivamente sobre el benchmark cohesivo, observar U(λ_final)"
      expect: "Convergencia primer orden en τ (esquema lineal en restricción): error → 0 como O(τ)"
      tol_rel: 0.2

  specific:
    - name: switching_cilindrico_a_disipacion_en_pico
      setup: "Modelo con daño que cruza σ_t0 en el paso N"
      expect: "Modos: 'cylindrical' en pasos 1..N-1; 'dissipation' desde N+1 en adelante; sin oscilación de modo en pasos contiguos"

    - name: sign_of_pivot_tracking_bifurcacion_simple
      setup: "Modelo con un punto límite simple (snap-back); registrar el número de pivots negativos antes y después del paso del límite"
      expect: "Conteo de pivots negativos cambia por 1 al atravesar el límite; predictor mantiene la dirección física"

    - name: fallback_a_cilindrico_en_descarga_global
      setup: "Modelo con grieta saturada (ω≈1) y descarga: λ decrece, F_ext invertida"
      expect: "Modo conmuta de vuelta a 'cylindrical' al detectar ΔE_d < umbral"

references:
  - "Gutiérrez M.A. (2004). Energy release control for numerical simulations of failure in quasi-brittle solids. *Communications in Numerical Methods in Engineering* 20(1), 19-29. (Restricción lineal de disipación; fórmula central de esta spec.)"
  - "Verhoosel C.V., Remmers J.J.C., Gutiérrez M.A. (2009). A dissipation-based arc-length method for robust simulation of brittle and ductile failure. *International Journal for Numerical Methods in Engineering* 77(9), 1290-1321. (Versión robusta con switching y aplicación a cohesivos.)"
  - "Crisfield M.A. (1991). *Non-linear Finite Element Analysis of Solids and Structures*, Vol. 1, Wiley. §9.4 (alternative constraints to cylindrical/spherical)."
  - "de Borst R., Sluys L.J. (1999). *Computational Methods in Non-linear Solid Mechanics*, TU Delft. §3.5 (controlled descent)."
  - "ADR 0010 §Hoja de ruta (fase 4 diferida por solver)."
  - "Spec padre: ArcLengthSolver.md."

\end{lstlisting}


\section{Implementación}

\begin{itemize}
  \item \textbf{Archivo}: \texttt{solidum/math/solvers/dissipation\_arclength.py}.
  \item \textbf{Clase}: \texttt{DissipationArcLengthSolver(ArcLengthSolver)} --- subclase. Reusa predictor, corrector, ajuste de paso, manejo de Dirichlet/MPC y backend del padre. Override completo de \texttt{solve()} con bifurcación por \texttt{self.\_mode}.
  \item \textbf{Atributo \texttt{self.\_mode}}: \texttt{"cylindrical" | "dissipation"}. Transición tras commit según umbral `\texttt{dissipation\_threshold\ensuremath{\cdot}\ensuremath{\|}F\_ref\ensuremath{\|}\ensuremath{\cdot}\ensuremath{\|}U\ensuremath{\|}}`.
  \item \textbf{Sign-of-pivot tracking}: \texttt{\_negative\_pivots()} implementado vía signo del \texttt{slogdet} de \texttt{K\_t}. Aproximación válida para distinguir indefinitud simple (0 vs número impar de pivots negativos) --- base para detección de paso por punto límite simple. \textbf{LDL\ensuremath{^{\top}} Bunch-Kaufman queda como deuda técnica} para tracking exacto del conteo.
  \item **Salvaguarda contra `\texttt{final\_step}\texttt{ prematuro**: si }\texttt{|d\ensuremath{\lambda}\_pred|}\texttt{ excede }\texttt{3 \ensuremath{\times} (max\_lambda − lambda\_curr)}`, se bisecta dl o \ensuremath{\tau} antes de aceptar el paso. Esencial en pasos iniciales sin calibrar y tras switch cilíndrico\ensuremath{\to}disipación donde \ensuremath{\alpha} puede ser pequeño.
  \item \textbf{Detección de \ensuremath{\alpha}\ensuremath{\approx}0 vía threshold relativo}: `\texttt{|\ensuremath{\alpha}| < 1e-6\ensuremath{\cdot}½\ensuremath{\cdot}|\ensuremath{\lambda}\_n\ensuremath{\cdot}F\ensuremath{\cdot}du\_t|}` reverte a cilíndrico --- para problemas lineales monotónicos \ensuremath{\alpha}≡0 exactamente, no es un error sino la matemática del régimen sin disipación física.
  \item \textbf{Entrypoint público}: registrado vía \texttt{@SolverRegistry.register}, expuesto como \texttt{solidum.math.solvers.DissipationArcLengthSolver}.
\end{itemize}

\subsection{Tests (suite tests/test\_dissipation\_arclength.py)}

\textbf{7/7 verde}. Cubren:

\begin{itemize}
  \item ✅ Recuperación cilíndrica en régimen elástico puro (regresión a paridad de bits con \texttt{ArcLengthSolver}).
  \item ✅ Daño 1D con softening pronunciado (rama post-pico atravesada).
  \item ✅ Switching automático cilíndrico\ensuremath{\leftrightarrow}disipación observable.
  \item ✅ Daño 2D continuo (Quad4 + IsotropicDamage2D plane strain).
  \item ✅ Balance energético cualitativo (\ensuremath{\Sigma}\ensuremath{\Delta}E\_d \ensuremath{\ge} 0 con daño activo).
  \item ✅ Sign-of-pivot tracking via signo del determinante (0 para PD, 1 para indefinida).
\end{itemize}

\subsection{Validación parcial --- qué falta para status \texttt{validated}}

El acceptance `\texttt{cohesive\_embedded\_softening\_postpeak}` \textbf{no está cubierto} en esta primera implementación. Diagnóstico:

La activación discreta del cohesivo en \texttt{prepare\_step} (Rankine onset al cruzar \texttt{\ensuremath{\sigma}\_t0}) introduce un salto en \texttt{F\_int} y \texttt{K\_t} entre el paso pre-activación y el paso post-activación. El predictor por \texttt{sign = sign(dot(\ensuremath{\delta}U\_prev, du\_t))} puede heredar una orientación inconsistente cuando el Newton del paso de activación reorienta \texttt{dU\_iter} para acomodar el nuevo \texttt{K\_t} con el cohesivo activo. El paso siguiente entonces predice hacia atrás (descarga) en lugar de continuar la rama post-pico.

Resolverlo limpiamente requiere \textbf{una de las dos vías siguientes}, ambas más allá del scope incremental de esta spec:

\begin{enumerate}
  \item \textbf{Sign-of-pivot tracking real con LDL\ensuremath{^{\top}} Bunch-Kaufman}, no aproximación por signo del determinante. Permitiría detectar la inercia exacta de \texttt{K\_t} y orientar el predictor por cambio de signo de pivot en lugar del producto escalar \texttt{\ensuremath{\delta}U\ensuremath{\cdot}du\_t}. Esto desacopla la orientación del predictor de la cinemática perturbada por la activación. Requiere extender el backend algebraico (ADR 0003 fase 2) con un solver LDL\ensuremath{^{\top}} verdadero.
  \item \textbf{Control por CMOD/CTOD del salto} mediante MPC sobre el DOF interno del cohesivo, en lugar de la restricción de disipación global. Cambia la naturaleza del solver --- sería una spec separada (\texttt{CMODControlSolver}) en lugar de variante de arc-length.
\end{enumerate}

Recomendación: dejar esta validación como deuda técnica priorizada en STATUS.md §"Deuda técnica priorizada" hasta que el LDL\ensuremath{^{\top}} Bunch-Kaufman entre en otro contexto (lo amerita por más razones, no solo esta).


\section{Diálogo}

\begin{itemize}
  \item \textbf{2026-05-19 (tarde)} \ensuremath{\cdot} Spec promovida de \texttt{draft} a \texttt{implemented}. 7/7 tests verde sobre los casos de daño continuo (1D y 2D); el caso crítico cohesivo+embedded reveló una limitación más profunda del \texttt{sign} por producto escalar tras activación discreta del Rankine, que requiere LDL\ensuremath{^{\top}} Bunch-Kaufman real o control CMOD/CTOD --- ambos fuera del scope incremental de esta variante. Documentado en §"Implementación" como deuda técnica. La spec mantiene su valor: provee la arquitectura del switching cilíndrico\ensuremath{\leftrightarrow}disipación y la fórmula lineal de Gutiérrez funcional para todos los problemas con softening continuo (daño bulk 1D y 2D), que es donde más casos prácticos de Solidum se concentran.
\end{itemize}

\begin{itemize}
  \item \textbf{2026-05-19 (mañana)} \ensuremath{\cdot} Spec creada como variante de \texttt{ArcLengthSolver} para desbloquear ADR 0010 fase 4 (embedded discontinuity con softening severo) y validación cuantitativa de $G_F$ en pipeline real. Motivada por el cierre de la campaña de validación externa Tandas 12-14 ([[project\_validacion\_externa\_tandas\_12\_14]]) que identificó esta pieza como \textbf{el primer componente nuevo priorizado} tras la campaña. Sin ADR nuevo: reusa decisiones arquitecturales del padre (estructura del bucle predictor/corrector, manejo de Dirichlet/MPC, ADR 0003 algebraico, ADR 0007 convergencia, ADR 0008 unidades). El cambio se circunscribe a la \textbf{restricción del paso} (cuadrática \ensuremath{\to} lineal) más mejoras secundarias (switching automático cilíndrico\ensuremath{\leftrightarrow}disipación, sign-of-pivot tracking).
\end{itemize}

\begin{itemize}
  \item \textbf{Puntos abiertos para el usuario} (físicos/de alcance, no de plumbing --- ver \texttt{[[feedback\_filtrar\_puntos\_abiertos]]}):
\end{itemize}
\begin{enumerate}
  \item ¿La restricción de disipación de Gutiérrez 2004 es la fórmula que quieres, o prefieres una alternativa (p. ej. constraint sobre \texttt{K\_t}-norma de Crisfield 1991 §9.4)? La de Gutiérrez es la más extendida en literatura cohesiva.
  \item ¿El switching automático cilíndrico\ensuremath{\leftrightarrow}disipación es deseable, o prefieres dos solvers separados (\texttt{CylindricalArcLengthSolver} y \texttt{DissipationArcLengthSolver} puros) que el usuario elige según el caso? El switching es estándar en Verhoosel et al. 2009; dos solvers separados son más predecibles pero exigen que el usuario diagnostique cuándo cambiar.
  \item ¿Sign-of-pivot tracking se aborda en esta spec o se difiere a una spec independiente (\texttt{PivotTracker})? Está acoplado al backend algebraico (ADR 0003) --- si el LDL\ensuremath{^{\top}} Bunch-Kaufman se introduce como helper compartido, otros solvers se beneficiarían. Mi propuesta es \textbf{abordarlo aquí en versión "signo del determinante LU"} (aproximación válida para indefinitud simple) y dejar el LDL\ensuremath{^{\top}} verdadero como deuda técnica del proyecto.
\end{enumerate}


\newpage

\section{LinearSolver}



\textit{Spec \textbf{retroactiva}: el solver existe desde la primera fase del proyecto y está validado por la inmensa mayoría de tests del pipeline estático. Esta spec lo documenta sin cambiar comportamiento --- H-5.3 de la auditoría 2026-05-18.}



\section{Especificación física}

\subsection{0. Descripción general}

Análisis estático \textbf{lineal} de una estructura discretizada por FEM: resolución directa del sistema algebraico \texttt{K\ensuremath{\cdot}U = F}. Hipótesis: linealidad geométrica y material (tangente constante = rigidez secante), apoyos Dirichlet (homogéneos o no homogéneos), restricciones lineales multipunto opcionales (ADR 0004). Sin iteración. Sin historia de carga.

\subsection{1. Ecuación de equilibrio resuelta}

Forma fuerte semidiscreta:

\[
\mathbf K\,\mathbf U \;=\; \mathbf F_{\text{ext}}
\]

donde $\mathbf K$ es la matriz de rigidez global ensamblada, $\mathbf U$ los desplazamientos nodales (incluidos los Dirichlet) y $\mathbf F_{\text{ext}}$ las cargas externas (incluidas peso propio si está activado).

\subsection{2. Condiciones de contorno}

Apoyos Dirichlet (homogéneos $u_s = 0$ o no homogéneos $u_s = g$) eliminados directamente por ADR 0004 fase 1. Restricciones lineales multipunto admitidas vía transformación $\mathbf T$ + offset $\mathbf g$ (ADR 0004 fase 2). El sistema reducido se resuelve en DOFs libres:

\[
\mathbf K_{\text{red}} = \mathbf T^\top \mathbf K\,\mathbf T, \qquad \mathbf F_{\text{red}} = \mathbf T^\top\,(\mathbf F - \mathbf K\,\mathbf g)
\]

\subsection{3. Salidas físicas}

\begin{itemize}
  \item Vector de desplazamientos global $\mathbf U \in \mathbb R^{n_{\text{dof}}}$ (libres + prescritos consistentes).
  \item Reacciones se recuperan a posteriori vía \texttt{domain.reactions(U)}.
  \item Esfuerzos en puntos de Gauss vía \texttt{domain.recompute\_gauss\_state(U)}.
  \item El resultado se envuelve en \texttt{SolveResult} por la capa pública (\texttt{solidum.run\_static}).
\end{itemize}


\section{Formulación numérica}

\subsection{4. Esquema operativo}

Un solo paso:

\begin{enumerate}
  \item Ensamblar $\mathbf K$ (sparse COO \ensuremath{\to} CSR) con caché topológica del Assembler.
  \item Aplicar la reducción ADR 0004: $\mathbf K_{\text{red}}$, $\mathbf F_{\text{red}}$, operadores $\mathbf T$ y $\mathbf g$.
  \item Calcular $\mathbf F_{\text{dir}} = \mathbf T^\top\,(\mathbf K\,\mathbf g)$ (contribución de Dirichlet no homogéneo).
  \item Factorizar $\mathbf K_{\text{red}}$ con el backend algebraico seleccionado.
  \item Resolver $\mathbf U_{\text{red}} = \mathbf K_{\text{red}}^{-1}\,(\mathbf F_{\text{red}} - \mathbf F_{\text{dir}})$.
  \item Expandir $\mathbf U = \mathbf T\,\mathbf U_{\text{red}} + \mathbf g$.
\end{enumerate}

\subsection{5. Predictor / corrector}

No aplica --- solver no iterativo.

\subsection{6. Criterio de convergencia}

No aplica --- la solución directa es exacta (a precisión de máquina y condicionamiento de $\mathbf K$).

\subsection{7. Imposición de Dirichlet y MPC}

ADR 0004 fase 1 (eliminación directa de apoyos) y fase 2 (transformación $\mathbf T$ + offset $\mathbf g$ para MPC). Manejados uniformemente por \texttt{Assembler.reduce(...)}.

\subsection{8. Backend algebraico}

Despacho automático ADR 0003:

\begin{itemize}
  \item Si el dominio es \textbf{simétrico} (todos los materiales con tangente simétrica) \ensuremath{\Rightarrow} se asume PD y se factoriza con \textbf{Cholesky} (\texttt{scipy.sparse.linalg.splu} con \texttt{options=\{'SymmetricMode': True\}} o equivalente Cholesky sparse).
  \item Si Cholesky reporta no-positividad (puede ocurrir con casi-singular o mal condicionado) \ensuremath{\Rightarrow} \textbf{fallback automático} a LU.
  \item Si el dominio declara material \textbf{no simétrico} \ensuremath{\Rightarrow} LU directo.
\end{itemize}

Override manual: parámetro \texttt{linear\_algebra \ensuremath{\in} \{"auto", "cholesky", "lu"\}}.

\textbf{Cache de factorización (ADR 0003 fase 2)}: la primera llamada a \texttt{solve} ensambla $\mathbf K$, reduce y factoriza; la factorización se guarda en el solver. Llamadas posteriores con un $\mathbf F$ distinto reutilizan el factor sin reensamblaje --- el coste se reduce a una resolución triangular barata. Si el usuario modifica el modelo entre llamadas (nuevos elementos, BCs, materiales) debe invocar \texttt{invalidate\_cache()} explícitamente.

\subsection{9. Adaptatividad / control de paso}

No aplica.

\subsection{10. Caveats numéricos}

\begin{itemize}
  \item \textbf{Suposición de PD por simetría}: el despachador asume que un dominio con materiales simétricos tiene $\mathbf K$ positiva definida. Es cierto para elasticidad estándar con suficientes apoyos; falso para casos degenerados (apoyos insuficientes, modos rígidos no restringidos, casi-singularidad geométrica). El fallback Cholesky\ensuremath{\to}LU cubre estos casos.
  \item \textbf{Cache desactualizado}: si el modelo se modifica entre \texttt{solve} calls sin llamar a \texttt{invalidate\_cache}, los resultados serán inconsistentes. La política deliberada es \textbf{cache silenciosa} (no se detectan modificaciones del modelo automáticamente) por velocidad --- la responsabilidad de invalidar es del usuario.
  \item \textbf{No es no lineal}: si el problema es no lineal (material plástico, geometría corotacional con cargas significativas) usar \texttt{NonlinearSolver}. Este solver evaluará la tangente en estado virgen ($\mathbf U = 0$) y producirá una respuesta lineal incorrecta sin avisar.
\end{itemize}


\section{Contrato de implementación}

\begin{lstlisting}[language=yaml]
name: LinearSolver
kind: solver
status: validated

interface:
  yaml_type: linear
  output: SolveResult
  pipeline_kind: static

parameters:
  - { name: linear_algebra, type: str, required: false, default: "auto",
      desc: "Selección del backend algebraico: 'auto' (Cholesky→LU según simetría), 'cholesky', 'lu'" }

requirements:
  - "Modelo lineal: materiales con tangente constante, geometría no corotacional o cargas suficientemente pequeñas"
  - "Apoyos suficientes para descartar modos rígidos (estructura cinemáticamente estable)"
  - "Compatibilidad ensamblaje-elemento-material verificada por Domain"

conventions:
  units: "heredadas del modelo (ADR 0008)"
  stability: "trivial — un solo paso"

out_of_scope:
  - "Análisis no lineal ⇒ usar NonlinearSolver"
  - "Snap-through / snap-back ⇒ usar ArcLengthSolver"
  - "Análisis dinámico ⇒ usar NewmarkSolver, HHTSolver, etc."
  - "Detección automática de modos rígidos (apoyos insuficientes)"

acceptance:
  verification:
    - name: viga_en_voladizo_carga_puntual
      setup: "viga de Frame2DEuler en voladizo con carga puntual P en el extremo"
      expect: "δ = P·L³/(3·E·I) ± O(elementos)"
      tol_rel: 1.0e-8
    - name: barra_traccion_uniaxial
      setup: "Truss2D, longitud L, área A, carga axial F"
      expect: "δ = F·L/(E·A) exacto (sin discretización)"
      tol_rel: 1.0e-12
    - name: convergencia_h
      setup: "viga discretizada con N elementos, N ∈ {2, 4, 8, 16}"
      expect: "error decrece con orden teórico del elemento (4 para Hermite cúbicos)"
      tol_rel: 0.1
    - name: reactions_equilibrio
      setup: "estructura cargada en una dirección"
      expect: "Σ(reacciones) + Σ(cargas externas) ≈ 0 (Newton 3ª ley)"
      tol_abs: 1.0e-9
    - name: dirichlet_no_homogeneo
      setup: "barra con apoyo prescrito u_s = g ≠ 0 en un extremo"
      expect: "solución exacta esperada (translación rígida + alargamiento si carga)"
      tol_rel: 1.0e-12
    - name: cache_factorizacion
      setup: "dos llamadas consecutivas a solve con F distintos sin invalidar cache"
      expect: "segunda llamada reutiliza el factor (mensaje de log diferenciado); resultado igual al cómputo desde cero"
      tol_rel: 1.0e-12
    - name: fallback_cholesky_lu
      setup: "matriz simétrica pero casi-singular o no PD efectiva (mecanismo cinemático añadido y luego soltado)"
      expect: "Cholesky aborta, LU resuelve sin fallar"
    - name: mpc_simple
      setup: "Quad4 con MPC u_2 = u_1 (master/slave)"
      expect: "u_2 = u_1 en la solución, equilibrio respetado"
      tol_abs: 1.0e-12

references:
  - "Bathe K.J. (2014). Finite Element Procedures. Prentice Hall. §2."
  - "Cook R.D., Malkus D.S., Plesha M.E., Witt R.J. (2002). Concepts and Applications of FEA. Wiley. §2.4 (eliminación de DOFs Dirichlet)."
  - "ADR 0003 — selección automática de backend algebraico."
  - "ADR 0004 — manejo de Dirichlet y constraints lineales."

\end{lstlisting}


\section{Implementación}

\begin{itemize}
  \item \textbf{Archivo}: solidum/math/solvers/linear.py.
  \item \textbf{Clase}: \texttt{LinearSolver}, registrada vía \texttt{@SolverRegistry.register} con \texttt{PIPELINE\_KIND = "static"}.
  \item \textbf{Cache}: atributos \texttt{\_factor}, \texttt{\_T}, \texttt{\_g\_full}, \texttt{\_F\_dir}, \texttt{\_n\_free} se rellenan lazy en la primera llamada a \texttt{solve}. Método público \texttt{invalidate\_cache()} los pone a \texttt{None} para forzar reensamblaje.
  \item \textbf{Fallback Cholesky\ensuremath{\to}LU}: capturado en \texttt{\_build\_cache} vía \texttt{CholeskyNotPositiveDefiniteError} (ADR 0003 §5); registra un warning y reintenta con LU.
  \item \textbf{Entrypoint público}: \texttt{solidum.run\_static(model, solver="linear", ...)} (despacho declarativo por \texttt{PIPELINE\_KIND}, regla C de la auditoría).
  \item \textbf{Tests}: cobertura masiva implícita --- prácticamente todos los tests del pipeline estático lineal lo usan. Tests específicos del cache: \texttt{tests/test\_solver\_robustness.py::test\_linear\_solver\_cache\_reuse}.
\end{itemize}


\section{Diálogo}

\begin{itemize}
  \item \textbf{2026-05-19} \ensuremath{\cdot} Spec creada retroactivamente para cerrar el hueco H-5.3. Solver anterior a la convención de specs validadas. La cache de factorización se añadió en la sesión de saneamiento post-auditoría (commit \texttt{b517718}, H-4.2); esta spec recoge ese comportamiento como parte del contrato actual.
\end{itemize}


\newpage

\section{NonlinearSolver}



\textit{Spec \textbf{retroactiva}: el solver existe desde la primera fase del proyecto y está validado por tests de plasticidad, daño, corotacional, embedded. Esta spec lo documenta sin cambiar comportamiento --- H-5.3 de la auditoría 2026-05-18.}



\section{Especificación física}

\subsection{0. Descripción general}

Análisis estático \textbf{no lineal} con \textbf{control de carga} (load-controlled) en una estructura discretizada por FEM. La no-linealidad puede provenir de la geometría (formulación corotacional, grandes desplazamientos pequeñas deformaciones) o del material (plasticidad, daño, embedded discontinuity). Esquema \textbf{incremental-iterativo}: la carga objetivo se aplica por incrementos $\Delta\lambda$, y en cada incremento se resuelve el equilibrio por \textbf{Newton-Raphson} con tangente reensamblada y paso bisectado al fallar.

\subsection{1. Ecuación de equilibrio resuelta}

Forma fuerte semidiscreta, residuo:

\[
\mathbf R(\mathbf U, \lambda) \;=\; \lambda\,\mathbf F_{\text{ext}}^{\text{ref}} - \mathbf F_{\text{int}}(\mathbf U) \;=\; \mathbf 0
\]

con $\lambda \in [0, 1]$ factor de carga aplicado a la carga de referencia $\mathbf F_{\text{ext}}^{\text{ref}}$, y $\mathbf F_{\text{int}}(\mathbf U)$ fuerzas internas no lineales (dependientes del estado actual de los puntos de Gauss).

\subsection{2. Condiciones de contorno}

Idénticas a \texttt{LinearSolver}: Dirichlet (homogéneo / no homogéneo) y MPC vía ADR 0004. Las cargas Dirichlet no homogéneas se aplican proporcionalmente al factor de carga (incremento $\Delta\mathbf g = \mathbf g \cdot \Delta\lambda$).

\subsection{3. Salidas físicas}

\begin{itemize}
  \item $\mathbf U$ final converged en $\lambda = 1$.
  \item Historia disponible vía \texttt{step\_callback(step, U, load\_factor)} que el usuario puede pasar al solver.
  \item Diagnóstico de divergencia con excepciones tipadas (ADR 0011): \texttt{OscillatingNewtonError}, \texttt{SingularTangentError}, \texttt{LoadExceedsCapacityError}, \texttt{UnknownDivergenceError}, todas subclases de \texttt{SolverDivergedError}.
\end{itemize}


\section{Formulación numérica}

\subsection{4. Esquema operativo}

\begin{lstlisting}
λ = 0
Δλ = 1 / num_steps
mientras λ < 1:
    intentar avanzar λ → λ + Δλ:
        U_iter = U_current
        prepare_all_steps(U_current)              # hook ADR 0010 (embedded)
        para iter en range(max_iter):
            K_t, F_int = assemble_non_linear_system(U_iter)
            R = (λ+Δλ)·F_ext - F_int
            δU = K_t^{-1} · R                     # reducido ADR 0004
            U_iter, ‖R‖, F_int = line_search(...) # ADR 0011, opt-in
            si converge (ADR 0007 dual):
                commit_all_states()
                λ += Δλ
                si iter < threshold y Δλ < base:
                    Δλ ← min(Δλ · growth, base)   # adaptativo: acelerar
                break
        si no converge:
            Δλ /= 2                               # bisección
            si Δλ < min_delta_lambda:
                clasificar divergencia y raise    # ADR 0011

\end{lstlisting}

\subsection{5. Predictor / corrector}

\textbf{Predictor}: $\Delta\mathbf U^{(0)} = \mathbf 0$ --- el incremento empieza desde el estado convergido del paso anterior.

\textbf{Corrector} (Newton iter $k$): $\mathbf K_t^{(k)}\,\delta\mathbf U^{(k)} = \mathbf R^{(k)}$, $\mathbf U^{(k+1)} = \mathbf U^{(k)} + \alpha\,\delta\mathbf U^{(k)}$, con $\alpha$ del line search (default $\alpha = 1$).

\subsection{6. Criterio de convergencia (ADR 0007)}

Patrón dual fuerza + desplazamiento, con escalas autoderivadas:

\[
\lVert\mathbf R[\text{libres}]\rVert \le \text{atol}_F + \text{rtol}_F \cdot \max(\lVert\mathbf F_{\text{ext}}\rVert, \lVert\mathbf F_{\text{int}}\rVert)
\]

\[
\lVert\Delta\mathbf U\rVert \le \text{atol}_d + \text{rtol}_d \cdot \lVert\mathbf U\rVert
\]

Semántica \textbf{AND} (ambos simultáneamente). Calibración: en el primer ensamblaje del solve, \texttt{force\_scale = max(\ensuremath{\|}F\_ext\ensuremath{\|}, \ensuremath{\|}F\_int\ensuremath{\|}, 1)} y \texttt{disp\_scale = force\_scale / max(|diag(K)|)}.

\subsection{7. Imposición de Dirichlet y MPC}

Idéntica a \texttt{LinearSolver} --- ADR 0004. La carga Dirichlet no homogénea $\mathbf g$ se aplica proporcionalmente al factor de carga $\lambda$ a través del operador del Assembler (\texttt{U\_current=U\_iter, load\_factor=next\_load\_factor} en la llamada a \texttt{reduce}).

\subsection{8. Backend algebraico}

Despacho automático ADR 0003 idéntico a \texttt{LinearSolver}. Fallback Cholesky\ensuremath{\to}LU activado.

\textbf{Newton modificado (ADR 0003 fase 2)}: parámetro opcional \texttt{freeze\_tangent\_after\_iter}. Si es \texttt{int N}, el solver factoriza fresca durante las primeras $N$ iteraciones del paso y reutiliza la factorización en las siguientes iteraciones del mismo paso (no entre pasos --- la tangente cambia con $\mathbf U$). Default \texttt{None} = Newton estándar (factoriza cada iter).

\subsection{9. Adaptatividad / control de paso}

\textbf{Acelerado}: si un paso converge en $< $ \texttt{NEWTON\_ADAPTIVE\_GROWTH\_ITER\_THRESHOLD} iteraciones y $\Delta\lambda < 1/$ \texttt{num\_steps}, $\Delta\lambda$ se multiplica por \texttt{NEWTON\_ADAPTIVE\_GROWTH\_FACTOR} (cota: $1/$\texttt{num\_steps}).

\textbf{Bisección}: si un paso no converge en \texttt{max\_iter}, $\Delta\lambda \to \Delta\lambda / 2$ y se reintenta. Aborto si $\Delta\lambda < $ \texttt{min\_delta\_lambda}.

\textbf{Line search ADR 0011} (opt-in, default \texttt{False}): variante GLL (Grippo-Lampariello-Lucidi) --- backtracking simple por descenso no monótono. Justificación: line search Wolfe puro rechaza pasos Newton correctos en regímenes plásticos/dañados con residuo transitorio creciente. Documentado en ADR 0011.

\subsection{10. Caveats numéricos}

\begin{itemize}
  \item \textbf{Snap-through / snap-back}: NO soportado --- \texttt{NonlinearSolver} con control de carga falla en puntos límite con $dU/d\lambda \to \infty$ o $d\lambda/dU < 0$. Usar \texttt{ArcLengthSolver}. Excepción: el solver atraviesa puntos límite \textit{suaves} (sin cambio de signo en $d\lambda$) con bisección + cinemática corotacional, hallazgo empírico de la auditoría 2026-05-18.
  \item \textbf{Softening severo con penalty cohesivo stiff} (embedded): el solver no atraviesa la transición elástico\ensuremath{\to}softening con $K_t$ casi singular en $\kappa_0$. Limitación específica al embedded; el daño bulk continuo (1D/2D) sí se atraviesa.
  \item \textbf{Plasticidad perfecta con line search off}: puede oscilar (residuo periódico, no diverge ni converge). Mitigación: activar \texttt{line\_search=True}.
  \item \textbf{Tangente singular}: lanza \texttt{SingularTangentError} con el último estado.
  \item \textbf{Newton modificado}: con \texttt{freeze\_tangent\_after\_iter} no nulo, la convergencia cuadrática se degrada a lineal en las iteraciones congeladas. Útil cuando el coste de factorizar domina (\texttt{K\_t} grande, materiales con tangente cara de evaluar).
\end{itemize}


\section{Contrato de implementación}

\begin{lstlisting}[language=yaml]
name: NonlinearSolver
kind: solver
status: validated

interface:
  yaml_type: nonlinear
  output: SolveResult
  pipeline_kind: static

parameters:
  - { name: convergence,                 type: ConvergenceCriterion, required: false, default: "ConvergenceCriterion()",
      desc: "Política dual fuerza+desplazamiento (ADR 0007)" }
  - { name: max_iter,                    type: int,   required: false, default: 20,
      desc: "Iteraciones Newton máximas por paso antes de bisectar (keyword-only)" }
  - { name: num_steps,                   type: int,   required: false, default: 10,
      desc: "Número base de incrementos para alcanzar λ=1 (keyword-only)" }
  - { name: adaptive,                    type: bool,  required: false, default: true,
      desc: "Habilita acelerado + bisección de Δλ (keyword-only)" }
  - { name: min_delta_lambda,            type: float, required: false, default: NEWTON_DEFAULT_MIN_DELTA_LAMBDA,
      desc: "Cota inferior para Δλ; por debajo se aborta con SolverDivergedError (keyword-only)" }
  - { name: linear_algebra,              type: str,   required: false, default: "auto",
      desc: "'auto' (Cholesky→LU según simetría), 'cholesky', 'lu' (keyword-only)" }
  - { name: freeze_tangent_after_iter,   type: int|None, required: false, default: null,
      desc: "Newton modificado: factoriza fresca primeras N iter, congela después (keyword-only)" }
  - { name: line_search,                 type: bool,  required: false, default: false,
      desc: "Line search GLL ADR 0011 (keyword-only)" }

requirements:
  - "Modelo con no-linealidad bien definida (material con tangente válida en todo el rango de carga)"
  - "Cargas Dirichlet no homogéneas constantes en magnitud (se aplican proporcionalmente a λ)"
  - "Apoyos suficientes para descartar modos rígidos"

conventions:
  units: "heredadas del modelo (ADR 0008)"
  stability: "no garantizada — los puntos límite con d λ/dU<0 abortan vía bisección"

out_of_scope:
  - "Snap-through / snap-back con d λ/dU<0 ⇒ usar ArcLengthSolver"
  - "Análisis dinámico ⇒ usar NewtonNewmarkSolver, NewtonHHTSolver, etc."
  - "Buckling lineal (eigenproblema generalizado) ⇒ futuro BucklingSolver"

acceptance:
  verification:
    - name: recuperacion_caso_lineal
      setup: "modelo enteramente lineal (Elastic2D + Quad4 + sin corotacional)"
      expect: "resultado idéntico a LinearSolver; Newton converge en 1 iteración por paso"
      tol_rel: 1.0e-12

    - name: plasticidad_uniaxial_truss
      setup: "Truss2D con Elastoplastic1D, carga axial creciente hasta plástico"
      expect: "respuesta bilineal exacta (E hasta yield, E_t = E·H/(E+H) después)"
      tol_rel: 1.0e-6

    - name: corotacional_voladizo_grande
      setup: "viga Frame2DEulerCorot, carga lateral grande (rotación >30°)"
      expect: "δ obtenido coincide con benchmark Bathe-Bolourchi 1979"
      tol_rel: 1.0e-3

    - name: damage_bulk_softening
      setup: "Quad4 con IsotropicDamage2D, tracción monotónica hasta post-pico"
      expect: "rama softening trazada; convergencia paso a paso con bisección"
      tol_rel: 1.0e-3

    - name: divergencia_clasificada
      setup: "carga que excede la capacidad última (plasticidad perfecta con mecanismo)"
      expect: "raise OscillatingNewtonError o LoadExceedsCapacityError con métricas estructuradas (ADR 0011)"

    - name: paso_adaptativo_acelera_y_bisecta
      setup: "trayectoria con tramos elásticos rápidos y zona plástica lenta"
      expect: "Δλ crece en tramos fáciles y bisecta en zonas duras; alcanza λ=1"

  specific:
    - name: newton_modificado_factoriza_menos
      setup: "freeze_tangent_after_iter=2, 5 iter típicas por paso"
      expect: "factorizaciones por paso ≤ 3 en lugar de 5; resultado dentro de tolerancia"
      tol_rel: 1.0e-8

    - name: line_search_corrige_oscilacion_DP
      setup: "Drucker-Prager 2D perfectamente plástico con carga que provoca oscilación"
      expect: "line_search=True converge; line_search=False oscila o requiere bisección extra"

references:
  - "Crisfield M.A. (1991). Non-linear Finite Element Analysis of Solids and Structures, Vol. 1. Wiley. §1-§9."
  - "Bonet J., Wood R.D. (2008). Nonlinear Continuum Mechanics for Finite Element Analysis. Cambridge. §8-§9."
  - "de Borst R., Sluys L.J. (1999). Computational Methods in Non-linear Solid Mechanics. TU Delft. §3 (Newton-Raphson y variantes)."
  - "ADR 0003, ADR 0004, ADR 0006, ADR 0007, ADR 0011."

\end{lstlisting}


\section{Implementación}

\begin{itemize}
  \item \textbf{Archivo}: solidum/math/solvers/nonlinear.py.
  \item \textbf{Clase}: \texttt{NonlinearSolver}, registrada vía \texttt{@SolverRegistry.register} con \texttt{PIPELINE\_KIND = "static"}.
  \item \textbf{Argumentos keyword-only}: todos menos \texttt{assembler} y \texttt{convergence} son keyword-only (auditoría H-1.5, sesión 2026-05-19, commit \texttt{4e4ed54}).
  \item \textbf{Bisección y clasificación}: en \texttt{solve}, al fallar \texttt{\ensuremath{\Delta}\ensuremath{\lambda} < min\_delta\_lambda} se llama a \texttt{classify\_divergence(residual\_history, delta\_history, singular\_tangent\_detected)} que devuelve la subclase apropiada de \texttt{SolverDivergedError} con métricas tipadas.
  \item \textbf{Entrypoint público}: \texttt{solidum.run\_static(model, solver="nonlinear", ...)}.
  \item \textbf{Tests}: cobertura amplísima --- todos los tests de plasticidad/daño/corotacional/embedded del pipeline estático lo ejercitan. Tests específicos del comportamiento adaptativo y de divergencia tipada en \texttt{tests/test\_solver\_robustness.py} y \texttt{tests/test\_solver\_diagnostics.py}.
\end{itemize}


\section{Diálogo}

\begin{itemize}
  \item \textbf{2026-05-19} \ensuremath{\cdot} Spec creada retroactivamente para cerrar el hueco H-5.3. Solver anterior a la convención de specs. Comportamiento documenta el estado tras la sesión de saneamiento post-auditoría (commits \texttt{4e4ed54} para keyword-only, ADR 0011 ya integrado).
\end{itemize}


\newpage

\chapter{Esquemas de Solución — Modal y dinámicos}

\section{CentralDifferenceSolver}



\textit{Spec de \textbf{solver nuevo} (no variante). Cubre el integrador explícito de segundo orden para análisis dinámico transitorio lineal y no lineal. Aprovecha la masa lumped diagonal habilitada por la fase 2 del ADR 0009 para evitar resolver sistema lineal en cada paso.}



\section{Especificación física}

\subsection{0. Descripción general}

Análisis dinámico transitorio \textbf{explícito} de segundo orden. Resuelve \texttt{M\ensuremath{\cdot}ü + C\ensuremath{\cdot}u̇ + K\ensuremath{\cdot}u = F(t)} (lineal) o \texttt{M\ensuremath{\cdot}ü + C\ensuremath{\cdot}u̇ + F\_int(u) = F(t)} (no lineal) avanzando la aceleración en cada paso por inversión trivial de la masa diagonal lumped --- sin sistema lineal, sin Newton interno. Apropiado para wave propagation, impacto, dinámica de respuesta rápida.

Hipótesis cinemáticas globales: las mismas que admite el bloque de elementos del modelo (corotacional, pequeñas deformaciones, etc.). Material lineal o no lineal con historia --- en modo \texttt{nonlinear=True} el solver llama \texttt{Assembler.assemble\_non\_linear\_system(u\_n)} cada paso para reevaluar \texttt{F\_int}.

\subsection{1. Ecuación de equilibrio resuelta}

\[
\mathbf{M}\,\ddot{\mathbf{u}} + \mathbf{C}\,\dot{\mathbf{u}} + \mathbf{F}_\text{int}(\mathbf{u}) = \mathbf{F}_\text{ext}(t)
\]

con $\mathbf{F}_\text{int}(\mathbf{u}) = \mathbf{K}\,\mathbf{u}$ en el caso lineal.

\subsection{2. Condiciones iniciales / de contorno}

\begin{itemize}
  \item Apoyos Dirichlet constantes en el tiempo (eliminación directa, ADR 0004).
  \item $\mathbf{u}(0) = \mathbf{u}_0$, $\dot{\mathbf{u}}(0) = \dot{\mathbf{u}}_0$ (cero por defecto).
  \item Cargas externas $\mathbf{F}_\text{ext}(t)$ vía \texttt{F\_func(t)}.
  \item Amortiguamiento Rayleigh $\mathbf{C} = \alpha\,\mathbf{M} + \beta\,\mathbf{K}$ (mismo contrato que \texttt{NewmarkSolver}).
\end{itemize}

MPC no soportado en esta fase --- no aparece en uso real con dinámica explícita.

\subsection{3. Salidas físicas}

\texttt{TransientResult} con historiales:

\begin{itemize}
  \item \texttt{u\_history[ndof, n\_steps+1]} desplazamientos.
  \item \texttt{udot\_history[ndof, n\_steps+1]} velocidades.
  \item \texttt{uddot\_history[ndof, n\_steps+1]} aceleraciones.
  \item \texttt{t\_history[n\_steps+1]} instantes de tiempo.
\end{itemize}


\section{Formulación numérica}

\subsection{4. Esquema operativo --- leapfrog Belytschko-Liu-Moran}

Inicialización:

\[
\mathbf{a}_0 = \mathbf{M}^{-1}\big(\mathbf{F}(0) - \mathbf{F}_\text{int}(\mathbf{u}_0) - \mathbf{C}\,\mathbf{v}_0 - \mathbf{F}_\text{dir}\big)
\]

\[
\mathbf{v}_{1/2} = \mathbf{v}_0 + \tfrac{\Delta t}{2}\,\mathbf{a}_0
\]

Paso $n \to n+1$:

\[
\mathbf{u}_{n+1} = \mathbf{u}_n + \Delta t\,\mathbf{v}_{n+1/2}
\]

\[
\mathbf{F}_\text{int}^{n+1} = \mathbf{K}\,\mathbf{u}_{n+1} \quad\text{(lineal)} \quad\text{o}\quad \text{ensamble}(\mathbf{u}_{n+1}) \quad\text{(no lineal)}
\]

\[
\mathbf{a}_{n+1} = \mathbf{M}^{-1}\big(\mathbf{F}(t_{n+1}) - \mathbf{F}_\text{int}^{n+1} - \mathbf{C}\,\mathbf{v}_{n+1/2} - \mathbf{F}_\text{dir}\big)
\]

\[
\mathbf{v}_{n+1} = \mathbf{v}_{n+1/2} + \tfrac{\Delta t}{2}\,\mathbf{a}_{n+1} \quad\text{(salida de paso)}
\]

\[
\mathbf{v}_{n+3/2} = \mathbf{v}_{n+1} + \tfrac{\Delta t}{2}\,\mathbf{a}_{n+1} \quad\text{(para el paso siguiente)}
\]

Sin Newton interno --- la ecuación de aceleración es explícita en $\mathbf{F}_\text{int}^{n+1}$ que ya se evaluó en el estado actual.

\subsection{5. Predictor / corrector}

El esquema leapfrog \textbf{no tiene predictor/corrector} en el sentido implícito. La actualización $\mathbf{u}_{n+1}$ se calcula sólo con datos en $n$ y $n+1/2$; la velocidad $\mathbf{v}_{n+1}$ se obtiene del semipaso.

\subsection{6. Criterio de convergencia}

No aplica --- esquema no iterativo. La convergencia "del problema" se materializa en la condición CFL y en la detección de divergencia.

\subsection{7. Imposición de Dirichlet y MPC}

Eliminación directa idéntica a \texttt{NewmarkSolver} (operador $\mathbf{T}$ + offset $\mathbf{g}$). MPC fuera de alcance en esta fase.

\subsection{8. Backend algebraico}

\textbf{Ninguno} --- la única operación lineal es $\mathbf{M}^{-1} \cdot \mathbf{r}$ con $\mathbf{M}$ diagonal, implementada como $1/\text{diag}(\mathbf{M})$ y multiplicación elemento a elemento. El solver verifica al construir $\mathbf{M}_\text{red}$ que la diagonal es estrictamente positiva y que los off-diagonals son numéricamente nulos (tolerancia $10^{-12} \cdot \max|\text{diag}|$); si no, lanza \texttt{ValueError}.

\subsection{9. Adaptatividad / control de paso}

No adaptativo en esta fase --- paso temporal constante $\Delta t$. La adaptatividad explícita (e.g. estimación de $\omega_{\max}$ por power iteration con actualización por paso) queda diferida hasta que aparezca un caso de uso con $\omega_{\max}$ variable significativamente en el tiempo.

\subsection{10. Caveats numéricos}

\begin{itemize}
  \item \textbf{Estabilidad condicional}. La condición CFL exige $\Delta t < 2/\omega_{\max}$ donde $\omega_{\max}$ es la frecuencia natural máxima del sistema discreto. Para una barra axial elástica con elemento de longitud $L_e$ y velocidad de propagación $c_p = \sqrt{E/\rho}$: $\Delta t_\text{crit} \approx L_e/c_p$. Para sólidos 2D la fórmula análoga es $\Delta t_\text{crit} \approx h_e/c_d$ donde $c_d$ es la velocidad dilatacional. Si se excede, la solución diverge exponencialmente --- el solver lo detecta a posteriori (umbral configurable \texttt{divergence\_threshold}, default $10^8 \cdot |\mathbf{u}_0|$) y aborta con \texttt{RuntimeError} informativo.
\end{itemize}

\begin{itemize}
  \item \textbf{Frame3D oblicuo}. La fase 2 del ADR 0009 documenta que $\mathbf{M}_\text{lumped}$ es bloque-diagonal (no estrictamente diagonal) en Frame3D cuando el eje del elemento no coincide con un eje global. Para central differences esto significaría invertir bloques 6\ensuremath{\times}6 por nodo cada paso, \textbf{anulando} la ventaja del esquema explícito. El solver rechaza ese caso con \texttt{ValueError}.
\end{itemize}

\begin{itemize}
  \item \textbf{Lumping obligatorio}. \texttt{lumping="consistent"} se rechaza al construir el solver --- invertir una $\mathbf{M}$ consistente cada paso es caro y elimina la razón de ser del explícito. El usuario debe pasar \texttt{lumping="lumped"} (default).
\end{itemize}

\begin{itemize}
  \item \textbf{Modos espurios de altas frecuencias}. El esquema no introduce amortiguamiento numérico (a diferencia de HHT-\ensuremath{\alpha}); las altas frecuencias de la malla pueden contaminar la respuesta de bajas frecuencias en problemas de wave propagation. Mitigaciones: refinamiento de malla, lumping puro (mejor filtrado que consistent), o cambio a HHT-\ensuremath{\alpha} implícito.
\end{itemize}


\section{Contrato de implementación}

\begin{lstlisting}[language=yaml]
name: CentralDifferenceSolver
kind: solver
status: validated

interface:
  yaml_type: CentralDifferenceSolver
  output: TransientResult

parameters:
  - { name: t_end, type: float, required: true, desc: "tiempo final (s)" }
  - { name: dt, type: float, required: true, desc: "paso temporal constante (s); debe satisfacer Δt < Δt_crit" }
  - { name: nonlinear, type: bool, required: false, desc: "default False; True recalcula F_int cada paso vía ensamble" }
  - { name: rayleigh, type: dict | null, required: false, desc: "C = α·M + β·K; mismo contrato que NewmarkSolver" }
  - { name: u0, type: ndarray | null, required: false, desc: "desplazamiento inicial global" }
  - { name: u0_dot, type: ndarray | null, required: false, desc: "velocidad inicial global" }
  - { name: F_func, type: callable | null, required: false, desc: "F(t) → ndarray global" }
  - { name: lumping, type: str, required: false, desc: "default 'lumped'; 'consistent' rechazado" }
  - { name: divergence_threshold, type: float, required: false, desc: "default 1e8; multiplicador sobre escala inicial para abortar" }

requirements:
  - "Materiales con `density` declarada (ADR 0008)."
  - "Todos los elementos del modelo deben soportar `compute_mass_matrix(lumping='lumped')` (ADR 0009 fase 2)."
  - "Frame3D con eje oblicuo a globales NO soportado (M_red no es diagonal pura)."

conventions:
  units:    "heredadas del modelo (ADR 0008)."
  stability: "Condicional — Δt < 2/ω_max (CFL). Para barras: Δt < L_el/c_p. Sin power iteration en esta fase — responsabilidad del usuario calcular o estimar Δt_crit."

out_of_scope:
  - "Estimación automática de Δt_crit (power iteration sobre M⁻¹K) — diferida."
  - "Paso adaptativo (Δt variable) — diferido."
  - "MPC (restricciones lineales multipunto) — diferido."
  - "Frame3D con eje oblicuo (rechazado por arquitectura del lumping)."

acceptance:
  verification:
    - name: free_vibration_1dof_undamped
      setup: "Truss2D 1 elemento con E=25, ρ=2, A=1, L=1 (lumped). u₀=1, v₀=0. Δt = T/200."
      expect: "u(t) = cos(ωt) con ω=5 rad/s, error < 2%"
      tol_rel: 0.02
    - name: step_response_1dof
      setup: "Mismo modelo con F₀=50·H(t)"
      expect: "u(t) = (F₀/K)(1 − cos(ωt)), error < 2% del régimen permanente"
      tol_rel: 0.02
    - name: matches_newmark_at_small_dt
      setup: "Mismo oscilador, Newmark β=1/4 γ=1/2 con masa lumped vs CentralDiff con Δt = T/500"
      expect: "u_Newmark − u_CD máximo < 1%"
      tol_rel: 0.01
  specific:
    - name: cfl_violation_diverges
      setup: "Δt = 3·Δt_crit"
      expect: "RuntimeError con mensaje CFL informativo"
      tol_rel: 0
    - name: cfl_threshold_matches_2_over_omega_max_1dof
      setup: "Oscilador 1-GDL (ω=5, Δt_crit=2/ω=0.4): integrar con Δt = 0.95·Δt_crit vs 1.05·Δt_crit"
      expect: "0.95·Δt_crit produce trayectoria acotada (30 ciclos); 1.05·Δt_crit lanza RuntimeError CFL"
      tol_rel: 0
    - name: cfl_threshold_matches_2_over_omega_max_2dof
      setup: |
        Cadena 2-DOF (3 trusses E=A=L=1, ρ=1, lumped) con K=[[2,-1],[-1,2]], M=I →
        autovalores cerrados ω² ∈ {1, 3}, ω_max=√3, Δt_crit=2/√3. Estado inicial
        excita modo antisimétrico (1, -1) — puro ω_max.
      expect: |
        (1) Δt = 0.95·Δt_crit estable.
        (2) Δt = 1.05·Δt_crit lanza RuntimeError CFL.
        (3) Δt = Δt_crit/10 reproduce u(t) = (cos(√3·t), -cos(√3·t)) con error < 2%.
      tol_rel: 0.02
      ref: "tests/test_central_difference.py::TestCentralDifferenceCFLAnalytic"
    - name: linear_mode_matches_nonlinear_mode_for_linear_material
      setup: "Mismo problema con material lineal puro, solver con nonlinear=False vs nonlinear=True"
      expect: "Coincidencia exacta a precisión de máquina"
      tol_rel: 1.0e-9
    - name: consistent_lumping_rejected
      setup: "Construir solver con lumping='consistent'"
      expect: "ValueError"
      tol_rel: 0
    - name: frame3d_oblique_rejected
      setup: "Frame3D con eje en dirección (1,1,1)/√3 e Iy ≠ Iz"
      expect: "ValueError 'no es estrictamente diagonal'"
      tol_rel: 0

references:
  - "Belytschko, T., Liu, W. K., Moran, B. (2014). *Nonlinear Finite Elements for Continua and Structures*, §6.2."
  - "Hughes, T. J. R. (2000). *The Finite Element Method*, §9.1."
  - "Crisfield, M. A. (1997). *Non-Linear Finite Element Analysis of Solids and Structures*, vol. 2, §24.6."
  - "Bathe, K.-J. (2014). *Finite Element Procedures*, §9.2 (estabilidad y CFL)."

\end{lstlisting}


\section{Implementación}

\begin{itemize}
  \item Archivo: solidum/math/solvers/central\_difference.py
  \item Clase: \texttt{CentralDifferenceSolver} (registrada en \texttt{SolverRegistry})
  \item Atributo de despacho: \texttt{PIPELINE\_KIND = "transient"} --- \texttt{run\_yaml} despacha a \texttt{run\_transient} por el mecanismo declarativo de la regla C (aplicada 2026-05-18 al introducir este solver, primer caso fuera de la jerarquía de Newmark)
  \item Tests: tests/test\_central\_difference.py (10 tests verdes)
\end{itemize}

Notas de traducción: ninguna --- la convención de signos del proyecto (Reglas §5) ya es la stress-resultant en el interior, sin reinterpretación específica de central differences.


\section{Diálogo}

\begin{itemize}
  \item \textit{2026-05-18}: regla C aplicada al introducir este solver. Antes de él, \texttt{run\_yaml::isinstance(NewmarkSolver)} capturaba HHT-\ensuremath{\alpha} por subclasificación; CentralDifference no podía heredar de NewmarkSolver (flujo radicalmente distinto, sin sistema lineal en cada paso, estabilidad condicional). En lugar de añadir una tercera rama \texttt{isinstance}, se refactorizó el dispatch a \texttt{PIPELINE\_KIND} declarativo --- futuros solvers no clásicos (Harmonic, ResponseSpectrum, \ldots{}) no requieren tocar \texttt{entry.py}.
\end{itemize}


\newpage

\section{HHTSolver}



\textit{Variante de \texttt{NewmarkSolver} según la regla de variantes}

\textit{(Reglas.md §4). Documenta \textbf{solo lo que cambia} respecto al padre. Implementa}

\textit{el integrador HHT-\ensuremath{\alpha} (1977) para introducir **amortiguamiento numérico}

\textit{controlado** en altas frecuencias preservando segundo orden de precisión y}

\textit{estabilidad incondicional. Reusa todas las decisiones arquitecturales del}

\textit{ADR 0009 --- sin ADR nuevo.}



\section{Qué cambia respecto a \texttt{NewmarkSolver}}

\texttt{NewmarkSolver} integra \texttt{M\ensuremath{\cdot}ü + C\ensuremath{\cdot}u̇ + K\ensuremath{\cdot}u = F(t)} evaluando todas las fuerzas
en el instante final \texttt{t\_\{n+1\}} del paso. Con la elección \ensuremath{\beta}=1/4, \ensuremath{\gamma}=1/2 (regla
trapezoidal, default) el integrador es \textbf{no disipativo}: los modos de alta
frecuencia espurios introducidos por la discretización persisten
indefinidamente en la respuesta y pueden enmascarar la señal de interés.

HHT-\ensuremath{\alpha} modifica la ecuación de equilibrio temporal para evaluar las **fuerzas
viscosas, elásticas y externas** en un instante intermedio \texttt{t\_\{n+1-\ensuremath{\alpha}\}},
mientras la \textbf{fuerza inercial} se mantiene en \texttt{t\_\{n+1\}}. Esto introduce
disipación numérica controlada por un único parámetro escalar \texttt{\ensuremath{\alpha} \ensuremath{\in} [−1/3, 0]},
sin sacrificar la precisión de segundo orden global ni la estabilidad
incondicional. Con \ensuremath{\alpha} = 0 se recupera Newmark trapezoidal.

\subsection{Ecuación de movimiento HHT-\ensuremath{\alpha}}

\[
\mathbf M\,\ddot{\mathbf u}_{n+1}
  + (1+\alpha)\,\mathbf C\,\dot{\mathbf u}_{n+1} - \alpha\,\mathbf C\,\dot{\mathbf u}_n
  + (1+\alpha)\,\mathbf K\,\mathbf u_{n+1} - \alpha\,\mathbf K\,\mathbf u_n
  \;=\; (1+\alpha)\,\mathbf F(t_{n+1}) - \alpha\,\mathbf F(t_n)
\]

Para preservar el orden 2 y la estabilidad incondicional, los parámetros
Newmark se derivan automáticamente del \ensuremath{\alpha}:

\[
\beta = \frac{(1-\alpha)^2}{4}, \qquad \gamma = \frac{1-2\alpha}{2}
\]

El usuario pasa solo \texttt{alpha}; \ensuremath{\beta} y \ensuremath{\gamma} se autoderivan salvo override explícito.

\subsection{Sistema efectivo}

Sustituyendo los predictores y correctores Newmark estándar (idénticos a la
fase 3 con los \ensuremath{\beta}, \ensuremath{\gamma} autoderivados), el sistema efectivo a resolver en cada
paso es:

\[
\mathbf A_\text{eff} \,\ddot{\mathbf u}_{n+1} \;=\; \mathbf b
\]

con

\[
\mathbf A_\text{eff} \;=\; \mathbf M + (1+\alpha)\,\gamma\Delta t\,\mathbf C + (1+\alpha)\,\beta\Delta t^2\,\mathbf K
\]

\[
\mathbf b \;=\; (1+\alpha)\,\mathbf F_{n+1} - \alpha\,\mathbf F_n
              - \mathbf F_\text{dir}
              - (1+\alpha)\,\mathbf C\,\tilde{\dot{\mathbf u}} + \alpha\,\mathbf C\,\dot{\mathbf u}_n
              - (1+\alpha)\,\mathbf K\,\tilde{\mathbf u} + \alpha\,\mathbf K\,\mathbf u_n
\]

\texttt{A\_eff} es constante (depende solo de M, C, K y de los coeficientes), por lo
que se \textbf{factoriza una sola vez} al inicio del análisis (igual que Newmark).
Cada paso temporal es una resolución triangular barata.

\subsection{Disipación numérica controlada}

El radio espectral del operador de amplificación en alta frecuencia
(\texttt{\ensuremath{\Delta}t \ensuremath{\to} \ensuremath{\infty}}) es

\[
\rho_\infty = \frac{1+\alpha}{1-\alpha}
\]

Para \texttt{\ensuremath{\alpha} = 0}: \texttt{\ensuremath{\rho}\_\ensuremath{\infty} = 1} (sin disipación, persistencia indefinida del modo).
Para \texttt{\ensuremath{\alpha} = −0.05}: \texttt{\ensuremath{\rho}\_\ensuremath{\infty} \ensuremath{\approx} 0.905} (disipación ligera, modo decae 9.5\% por
periodo de alta frecuencia). Para \texttt{\ensuremath{\alpha} = −1/3} (límite teórico): \texttt{\ensuremath{\rho}\_\ensuremath{\infty} = 0.5}
(disipación máxima manteniendo orden 2).

La curva de disipación es \textbf{selectiva}: amortigua significativamente solo los
modos con \texttt{\ensuremath{\omega}\ensuremath{\Delta}t > \ensuremath{\pi}}, dejando casi intactos los modos de interés con
\texttt{\ensuremath{\omega}\ensuremath{\Delta}t « 1}. Es lo que distingue a HHT-\ensuremath{\alpha} de simplemente subir el \ensuremath{\gamma} Newmark
(que disipa en \textit{todas} las frecuencias por igual y baja el orden a 1).

\subsection{Recuperación del comportamiento Newmark}

Con \texttt{\ensuremath{\alpha} = 0}, los coeficientes auto-derivados quedan \texttt{\ensuremath{\beta} = 1/4, \ensuremath{\gamma} = 1/2} y la
ecuación de movimiento HHT colapsa exactamente a la de Newmark trapezoidal.
\texttt{HHTSolver(alpha=0.0)} reproduce \textbf{bit a bit} los resultados de
\texttt{NewmarkSolver} con \texttt{beta=0.25, gamma=0.5}. Esto se valida en los acceptance.

\subsection{Variante no lineal}

\texttt{NewtonHHTSolver} (subclase de \texttt{NewtonNewmarkSolver}) aplica el mismo cambio
de la ecuación de movimiento al residuo dinámico. El Newton interno y su
jacobiano se ven afectados análogamente:

\[
\mathbf R(\ddot{\mathbf u}_{n+1}) \;=\; (1+\alpha)\,\mathbf F_{n+1} - \alpha\,\mathbf F_n
  - [(1+\alpha)\,\mathbf F_\text{int}(\mathbf u_{n+1}) - \alpha\,\mathbf F_\text{int}(\mathbf u_n)]
  - [(1+\alpha)\,\mathbf C\,\dot{\mathbf u}_{n+1} - \alpha\,\mathbf C\,\dot{\mathbf u}_n]
  - \mathbf M\,\ddot{\mathbf u}_{n+1}
\]

\[
\mathbf J \;=\; \mathbf M + (1+\alpha)\,\gamma\Delta t\,\mathbf C + (1+\alpha)\,\beta\Delta t^2\,\mathbf K_\text{t}
\]

Todo lo demás (criterio de convergencia, line search opt-in del ADR 0011,
telemetría tipada de divergencia, commit/rollback de estado, amortiguamiento
Rayleigh constante en el tiempo) se hereda de \texttt{NewtonNewmarkSolver} sin
cambios.


\section{Contrato de implementación}

\begin{lstlisting}[language=yaml]
name: HHTSolver
kind: solver
status: validated
extends: NewmarkSolver       # variante según Reglas §4

interface:
  type_yaml: HHTSolver
  pipeline_kind: transient    # mismo que NewmarkSolver

parameters_extra_to_NewmarkSolver:
  - { name: alpha,  type: float, required: false, default: -0.05,
      desc: "Parámetro HHT en [-1/3, 0]. α=0 recupera Newmark trapezoidal.
             α=-0.05 (default) es disipación ligera estándar; α=-1/3 disipación máxima." }

parameters_overridden_with_autoderivation:
  - { name: beta,   type: float, required: false, default: "(1-alpha)^2/4",
      desc: "Auto-derivado del alpha si no se pasa; override explícito posible
             (no recomendado salvo experimentación)." }
  - { name: gamma,  type: float, required: false, default: "(1-2*alpha)/2",
      desc: "Auto-derivado del alpha si no se pasa." }

parameters_inherited:
  # Idénticos a NewmarkSolver: t_end, dt, rayleigh, u0, u0_dot, F_func,
  # linear_algebra, lumping. Ver docs/specs/NewmarkSolver.md.

acceptance:
  recovery:
    - name: alpha_zero_matches_NewmarkSolver
      setup: "Oscilador 1 GDL con HHTSolver(alpha=0.0) y NewmarkSolver(beta=0.25, gamma=0.5)"
      expect: "u_history coincide bit-a-bit (atol 1e-10)"

  damping:
    - name: high_frequency_dissipation
      setup: "Dos osciladores acoplados con ω₁=1 rad/s y ω₂=100 rad/s; condición inicial excita ambos modos"
      expect: "HHTSolver(alpha=-0.1) atenúa el modo alto en >50% tras 5 periodos del modo bajo,
               manteniendo amplitud del modo bajo dentro de 5% del valor inicial"

    - name: spectral_radius_at_high_frequency_limit
      setup: "Δt grande (>10 periodos del modo); HHTSolver con varios α"
      expect: "ρ_∞ observado = (1+α)/(1-α) dentro de tolerancia 5%"

    - name: amplification_matrix_spectral_radius_analytic
      setup: "Oscilador 1-GDL con HHTSolver(α∈{-0.05, -0.10, -1/3}); Δt tal que Ω=ω·Δt=2"
      expect: |
        Razón de contracción por paso medida en la simulación
        coincide con max|λ(A)|, donde A es la matriz de amplificación 3×3 HHT-α
        (Hughes 1987 §9.3, Bathe 2014 §9.5.4): D=1+β(1+α)Ω², A[0,0]=(1+αβΩ²)/D, …
        Tolerancia 1%.
      ref: "tests/test_hht.py::TestHHTNumericalDampingAnalytic"

  stability:
    - name: unconditional_stability_large_dt
      setup: "Sistema lineal con Δt = 100·T (T periodo natural), 10 pasos"
      expect: "u_history acotada (no diverge) para todo α ∈ [-1/3, 0]"

  nonlinear:
    - name: oscilador_plastico_con_disipacion_HHT
      setup: "Truss elastoplástico 1 GDL, vibración libre desde u_0 más allá del yield"
      expect: "NewtonHHTSolver(alpha=-0.1) reduce oscilación espuria post-yielding
               manteniendo decay físico de la disipación plástica"

references:
  - "Hilber, H. M., Hughes, T. J. R. & Taylor, R. L. (1977). Improved numerical
     dissipation for time integration algorithms in structural dynamics.
     Earthquake Eng. Struct. Dyn. 5(3), 283-292."
  - "Hughes, T.J.R. (2000). The Finite Element Method, §9.3 (HHT-α y generalized-α)."
  - "Chopra, A. K. (2017). Dynamics of Structures, §15.4 (integradores con disipación numérica)."
  - "ADR 0009 §Hoja de ruta — fila 3 (Newmark / HHT-α)."
  - "Spec padre: [docs/specs/NewmarkSolver.md](NewmarkSolver.md)."
  - "Spec hermana: [docs/specs/NewtonNewmarkSolver.md](NewtonNewmarkSolver.md) — padre del NewtonHHTSolver."

\end{lstlisting}


\section{Implementación}

\begin{itemize}
  \item \textbf{Archivo}: solidum/math/solvers/newmark.py --- junto a \texttt{NewmarkSolver} y \texttt{NewtonNewmarkSolver} (lógicamente cercanos; comparten predictores Newmark, reducción Dirichlet y amortiguamiento Rayleigh).
  \item \textbf{Clases}:
  \item \texttt{HHTSolver(NewmarkSolver)} --- variante lineal con disipación numérica.
  \item \texttt{NewtonHHTSolver(NewtonNewmarkSolver)} --- variante no lineal con Newton interno + disipación numérica.
\end{itemize}
  Ambas registradas con \texttt{@SolverRegistry.register}.
\begin{itemize}
  \item \textbf{Auto-derivación de \ensuremath{\beta}, \ensuremath{\gamma}}: en \texttt{\_\_init\_\_}, si el usuario no pasa \ensuremath{\beta} o \ensuremath{\gamma} explícitos, se computan desde \ensuremath{\alpha} antes de invocar \texttt{super().\_\_init\_\_}. Esto asegura que el padre vea la \ensuremath{\beta}, \ensuremath{\gamma} ya correctas y la mayor parte del flujo herede sin cambios.
  \item \textbf{Solve overrideado}: \texttt{solve()} se reescribe en ambas subclases para usar las ecuaciones HHT-\ensuremath{\alpha} (ver §"Sistema efectivo" y §"Variante no lineal" arriba). El estado anterior \texttt{F\_int\_n}, \texttt{F\_n}, \texttt{u̇\_n} necesario para el término \texttt{−\ensuremath{\alpha}\ensuremath{\cdot}X\_n} se mantiene en variables locales del bucle.
  \item \textbf{Despacho YAML}: \texttt{solver: type: HHTSolver} (o \texttt{NewtonHHTSolver}). Como \texttt{HHTSolver} hereda de \texttt{NewmarkSolver} y \texttt{NewtonHHTSolver} hereda de \texttt{NewtonNewmarkSolver}, el \texttt{isinstance} en \texttt{entry.run\_yaml} los detecta como \texttt{transient} automáticamente. \textbf{No dispara la regla C arquitectural} (que esperaría una rama no-clásica que NO sea subclase de las existentes).
  \item \textbf{Tests}: \texttt{tests/test\_hht.py} nuevo con los acceptance arriba.
\end{itemize}


\section{Diálogo}

\begin{itemize}
  \item \textbf{2026-05-18} \ensuremath{\cdot} Spec creada como \textbf{variante} de \texttt{NewmarkSolver} (y \texttt{NewtonHHTSolver} como variante de \texttt{NewtonNewmarkSolver}) aplicando la regla §4. Sin ADR nuevo: reusa la decisión del ADR 0009 fila 3 ("Transitorio lineal Newmark / HHT-\ensuremath{\alpha}"). Sin entrada propia en el manual estructural --- visible al usuario solo a través del \texttt{type} YAML y el parámetro \texttt{alpha}.
\end{itemize}

\begin{itemize}
  \item \textbf{2026-05-18} \ensuremath{\cdot} Decisión sobre el default de \texttt{alpha}: \textbf{−0.05} (disipación ligera estándar) en lugar de \texttt{0} (sin disipación). Razones físicas: (i) si el usuario pide \texttt{HHTSolver}, está pidiendo disipación --- default \texttt{\ensuremath{\alpha}=0} lo convierte en un alias trivial de \texttt{NewmarkSolver}, ruido en el catálogo; (ii) \texttt{−0.05} es el valor canónico de Hilber 1977 (artículo original) y de los códigos de referencia (Abaqus, OpenSees) por ser balance entre disipación útil (>9\% por ciclo en altas frecuencias) y precisión preservada (<0.5\% en modos de interés); (iii) no introduce regresiones --- \texttt{HHTSolver} es un solver nuevo, nadie lo usa hoy. Diferente del default \texttt{line\_search=False} del ADR 0011, que era \textit{no} poder romper código existente; aquí no hay código existente que romper.
\end{itemize}

\begin{itemize}
  \item \textbf{2026-05-18} \ensuremath{\cdot} Decisión sobre auto-derivación de \ensuremath{\beta}, \ensuremath{\gamma}: por defecto sí, desde \texttt{alpha}. Override explícito posible pero no recomendado (combinaciones \ensuremath{\alpha}, \ensuremath{\beta}, \ensuremath{\gamma} arbitrarias pueden perder estabilidad incondicional u orden 2). El parámetro de la API es \textbf{el} \ensuremath{\alpha}, no la terna; \ensuremath{\beta}, \ensuremath{\gamma} son consecuencia matemática de \ensuremath{\alpha} salvo experimentación expresa. Esto simplifica el uso típico y centraliza la responsabilidad: si el usuario solo conoce \ensuremath{\alpha}, el solver hace el resto.
\end{itemize}


\newpage

\section{HarmonicSolver}



\textit{Spec de \textbf{solver nuevo}. Resuelve \texttt{(−\ensuremath{\omega}\ensuremath{^{2}}M + i\ensuremath{\omega}C + K)\ensuremath{\cdot}û = F̂} para un barrido de frecuencias y devuelve la amplitud compleja \texttt{û(\ensuremath{\omega})}. Análisis lineal en frecuencia, sin transitorio.}



\section{Especificación física}

\subsection{0. Descripción general}

Análisis \textbf{lineal} de respuesta forzada armónica en estado estacionario. La excitación \texttt{F(t) = Re\{F̂\ensuremath{\cdot}e\textasciicircum{}\{i\ensuremath{\omega}t\}\}} produce una respuesta \texttt{u(t) = Re\{û(\ensuremath{\omega})\ensuremath{\cdot}e\textasciicircum{}\{i\ensuremath{\omega}t\}\}} también armónica con la misma frecuencia. El solver calcula \texttt{û(\ensuremath{\omega})} para un barrido de frecuencias. Útil para:

\begin{itemize}
  \item Funciones de transferencia (FRFs) entrada\ensuremath{\to}salida.
  \item Diagnóstico de resonancias previo a análisis transitorio caro.
  \item Respuesta forzada en estructuras con cargas periódicas (maquinaria rotante, sismo armónico, viento turbulento descompuesto en armónicos).
\end{itemize}

Hipótesis cinemáticas y material: lineales (\texttt{K} y \texttt{M} constantes). No se modelan grandes desplazamientos, plasticidad ni daño --- quien quiera no linealidad debe usar \texttt{NewtonNewmarkSolver} con excitación armónica.

\subsection{1. Ecuación de equilibrio resuelta}

\[
\big(-\omega^2\mathbf{M} + i\omega\mathbf{C} + \mathbf{K}\big)\,\hat{\mathbf{u}}(\omega) = \hat{\mathbf{F}}
\]

para cada \texttt{\ensuremath{\omega}} del barrido. \texttt{C = \ensuremath{\alpha}\ensuremath{\cdot}M + \ensuremath{\beta}\ensuremath{\cdot}K} Rayleigh.

\subsection{2. Condiciones iniciales / de contorno}

\begin{itemize}
  \item Apoyos Dirichlet constantes (eliminación directa, ADR 0004).
  \item \texttt{F̂} \ensuremath{\in} ℂ\textasciicircum{}ndof. Real puro si la carga lleva fase cero; complejo si lleva fase.
  \item \textbf{No} hay condiciones iniciales --- el estado estacionario es independiente de ellas (el transitorio se descarta por definición).
  \item MPC no soportado en esta fase.
\end{itemize}

\subsection{3. Salidas físicas}

\texttt{HarmonicResult} (inmutable):

\begin{itemize}
  \item \texttt{omega[n\_omega]} frecuencias evaluadas (rad/s).
  \item \texttt{frequencies\_hz} propiedad derivada.
  \item \texttt{u\_complex[ndof, n\_omega]} amplitud compleja del desplazamiento.
  \item \texttt{F\_complex[ndof]} amplitud compleja de la carga (informativa).
  \item Métodos auxiliares: \texttt{.amplitude()} \ensuremath{\to} \texttt{|û|}, \texttt{.phase()} \ensuremath{\to} \texttt{arg(û)}.
\end{itemize}


\section{Formulación numérica}

\subsection{4. Esquema operativo}

\begin{lstlisting}
Ensamblar K, M.
C = α·M + β·K (Rayleigh).
Reducir por Dirichlet: K_red, M_red, C_red, F_red.
Para cada ω en el barrido:
    Z(ω) = -ω²·M_red + iω·C_red + K_red   [matriz compleja]
    û_red(ω) = Z(ω)⁻¹ · F_red             [spsolve complejo]
    û(ω) = T · û_red(ω) [expand a globales]

\end{lstlisting}

\subsection{5. Predictor / corrector}

No aplica --- un solve directo por frecuencia.

\subsection{6. Criterio de convergencia}

No aplica --- esquema no iterativo. La precisión depende del backend
algebraico (\texttt{scipy.sparse.linalg.spsolve} LU complejo).

\subsection{7. Imposición de Dirichlet y MPC}

Eliminación directa idéntica a \texttt{NewmarkSolver}. En DOFs prescritos
\texttt{û = 0} (el apoyo no oscila). MPC fuera de alcance.

\subsection{8. Backend algebraico}

Factorización LU compleja por frecuencia vía
\texttt{scipy.sparse.linalg.spsolve(Z.tocsc(), F\_red)}. \textbf{No hay caché}: \texttt{Z}
cambia con \texttt{\ensuremath{\omega}}. Coste dominante del análisis es `O(n\_omega \ensuremath{\cdot}
coste\_factorizacion(Z))`.

Posible optimización futura (fuera de esta fase): superposición modal
--- ortogonalizar la respuesta sobre los primeros \texttt{n\_modes} modos del
problema generalizado, lo que reduce el barrido a \texttt{n\_omega} operaciones
escalares. Justifica el dispositivo cuando \texttt{n\_modes ≪ n\_dof} y el
operador es simétrico.

\subsection{9. Adaptatividad / control de paso}

Espaciado del barrido controlado por el usuario:

\begin{itemize}
  \item \texttt{scale="linear"}: \texttt{np.linspace(omega\_min, omega\_max, n\_omega)}.
  \item \texttt{scale="log"}: \texttt{np.geomspace(omega\_min, omega\_max, n\_omega)} ---
\end{itemize}
  apropiado para FRFs sobre rangos de varias décadas.
\begin{itemize}
  \item Vector explícito vía \texttt{omega=...} --- usuarios con criterios
\end{itemize}
  específicos (densificar alrededor de resonancias conocidas, etc.).

No hay refinamiento adaptativo de la malla de frecuencia.

\subsection{10. Caveats numéricos}

\begin{itemize}
  \item \textbf{Resonancia exacta sin amortiguamiento}: \texttt{Z(\ensuremath{\omega}\_n)} es singular ---
\end{itemize}
  \texttt{spsolve} puede fallar o devolver valores enormes. Mitigación:
  evitar \texttt{\ensuremath{\omega}} coincidente con un modo, o añadir amortiguamiento
  Rayleigh.
\begin{itemize}
  \item \textbf{Mal condicionamiento cerca de resonancia}: la factorización LU
\end{itemize}
  pierde precisión a medida que \texttt{Z} se aproxima a la singularidad. Es
  intrínseco al método directo; superposición modal es la alternativa
  estable.
\begin{itemize}
  \item \textbf{Cargas con fase}: \texttt{F̂} compleja requiere construcción desde código
\end{itemize}
  Python; YAML no soporta tipos complejos nativos. Si el usuario YAML
  necesita fase, debe construir el dominio programáticamente.
\begin{itemize}
  \item \textbf{No analiza estabilidad} del estado estacionario (no aplica al
\end{itemize}
  problema lineal --- el estado estacionario siempre existe y es único
  para \texttt{C} definida positiva y \texttt{K} regular).


\section{Contrato de implementación}

\begin{lstlisting}[language=yaml]
name: HarmonicSolver
kind: solver
status: validated

interface:
  yaml_type: HarmonicSolver
  output: HarmonicResult

parameters:
  - { name: omega, type: array-like | null, required: false, desc: "vector explícito (rad/s); excluye omega_min/omega_max/n_omega/scale" }
  - { name: omega_min, type: float, required: false, desc: "límite inferior del barrido (rad/s)" }
  - { name: omega_max, type: float, required: false, desc: "límite superior del barrido (rad/s)" }
  - { name: n_omega, type: int, required: false, desc: "número de puntos del barrido, default 100" }
  - { name: scale, type: str, required: false, desc: "'linear' (default) o 'log'" }
  - { name: F_amplitude, type: ndarray | null, required: false, desc: "amplitud compleja shape (ndof,); default ceros; YAML inyecta desde point_loads como real" }
  - { name: rayleigh, type: dict | null, required: false, desc: "C = α·M + β·K; sin amortiguamiento ⇒ pico singular en resonancia" }
  - { name: lumping, type: str, required: false, desc: "'consistent' (default) o 'lumped'" }

requirements:
  - "Materiales con `density` declarada (ADR 0008)."
  - "Problema lineal — `K` y `M` constantes en el barrido."

conventions:
  units:    "heredadas del modelo (ADR 0008). ω en rad/s — coherente con frecuencias de ModalSolver."
  stability: "Lineal en frecuencia — estable para C definida positiva. Mal condicionado en ω → ω_n sin amortiguamiento."

out_of_scope:
  - "Superposición modal — diferida; el barrido directo es más simple y suficiente para los tamaños actuales del catálogo."
  - "MPC (restricciones lineales multipunto) — diferido."
  - "No linealidad (geométrica o material) — usar NewtonNewmarkSolver con excitación armónica."
  - "Cargas con fase declaradas desde YAML — usar construcción desde código."

acceptance:
  verification:
    - name: transfer_function_no_damping
      setup: "1 GDL con K=25, M=1; barrido en ω lejos de resonancia"
      expect: "û_real(ω) = F̂/(K − ω²M); û_imag ≈ 0"
      tol_rel: 1.0e-10
    - name: transfer_function_rayleigh_damping
      setup: "1 GDL con α-mass Rayleigh"
      expect: "û = F̂/(K − ω²M + iωc) complejo"
      tol_rel: 1.0e-9
  specific:
    - name: resonance_peak_location
      setup: "1 GDL con ξ=2% Rayleigh; barrido fino alrededor de ω_n"
      expect: "pico en ω = ω_n·√(1 − 2ξ²)"
      tol_rel: 0.01
    - name: matches_newmark_steady_state
      setup: "Newmark con F(t)=F₀cos(ωt), suficiente t_end para descartar transitorio (8τ)"
      expect: "amplitud Newmark estacionaria = |û_harmonic|"
      tol_rel: 0.02
    - name: log_sweep_geometric_spacing
      setup: "scale='log', omega_min=1, omega_max=1000, n_omega=4"
      expect: "omega = geomspace(1, 1000, 4) = [1, 10, 100, 1000]"
      tol_rel: 1.0e-12

references:
  - "Bathe, K.-J. (2014). *Finite Element Procedures*, §9.5."
  - "Clough, R. W., Penzien, J. (1993). *Dynamics of Structures*, §4.6."
  - "Chopra, A. K. (2017). *Dynamics of Structures*, cap. 4."

\end{lstlisting}


\section{Implementación}

\begin{itemize}
  \item Archivo: solidum/math/solvers/harmonic.py
  \item Clase: \texttt{HarmonicSolver} (registrada en \texttt{SolverRegistry})
  \item Atributo de despacho: \texttt{PIPELINE\_KIND = "harmonic"} \ensuremath{\to} \texttt{solidum.entry.run\_harmonic} \ensuremath{\to} \texttt{HarmonicResult}.
  \item Tests: tests/test\_harmonic.py (11 tests verdes).
  \item Resultado: \texttt{HarmonicResult} --- dataclass inmutable con \texttt{omega}, \texttt{u\_complex}, \texttt{F\_complex}, propiedades \texttt{frequencies\_hz}, \texttt{n\_omega}, métodos \texttt{amplitude()} y \texttt{phase()}.
\end{itemize}


\section{Diálogo}

\begin{itemize}
  \item \textit{2026-05-18}: nuevo \texttt{PIPELINE\_KIND = "harmonic"} añadido al despacho declarativo de \texttt{run\_yaml} (regla C, ya aplicada en D2). Tercer valor del literal junto a \texttt{"modal"}, \texttt{"transient"}, \texttt{"static"}. El parser YAML inyecta automáticamente las \texttt{point\_loads} del bloque estándar como \texttt{F\_amplitude} real cuando el usuario no la especifica explícitamente --- convención coherente con el uso de \texttt{external\_forces} en análisis estáticos.
\end{itemize}


\newpage

\section{ModalSolver}



\textit{Orden de trabajo. La especificación física y la formulación numérica son patrimonio del usuario (revisa con detalle). La IA implementa y rellena las secciones marcadas.}



\section{Especificación física}

\subsection{0. Descripción general}
Análisis modal lineal de una estructura discretizada por FEM: cálculo de frecuencias naturales y modos de vibración no amortiguados. Hipótesis: linealidad geométrica y material, amortiguamiento nulo (\texttt{C = 0}), masa constante en el tiempo (lagrangeano total).

\subsection{1. Problema físico}
Vibración libre no amortiguada del sistema discreto:
\[
\mathbf M\,\ddot{\mathbf u}(t) + \mathbf K\,\mathbf u(t) = \mathbf 0.
\]
Solución armónica $\mathbf u(t) = \boldsymbol\phi\, e^{i\omega t}$. Sustituyendo:
\[
\bigl(\mathbf K - \omega^2 \mathbf M\bigr)\,\boldsymbol\phi = \mathbf 0.
\]

\subsection{2. Problema de autovalor generalizado}
Solución no trivial requiere $\det(\mathbf K - \omega^2 \mathbf M) = 0$. Equivalente:
\[
\mathbf K\,\boldsymbol\phi_n = \omega_n^2\, \mathbf M\,\boldsymbol\phi_n, \qquad n = 1, \ldots, N_{\text{dof}}.
\]
Autovalores $\lambda_n = \omega_n^2 \ge 0$; autovectores $\boldsymbol\phi_n$ son los modos de vibración.

\subsection{3. Propiedades del par $(\mathbf K, \mathbf M)$}
\begin{itemize}
  \item $\mathbf K$ simétrica, positiva (semi)definida tras Dirichlet.
  \item $\mathbf M$ simétrica, \textbf{positiva definida} estrictamente para masa consistente (toda fila/columna tiene aporte de algún elemento con $\rho > 0$).
  \item Espectro real no negativo: $\omega_n \in \mathbb R_{\ge 0}$.
\end{itemize}

\subsection{4. Salidas}
\begin{itemize}
  \item Frecuencias naturales $\omega_n$ (rad/s), $f_n = \omega_n/(2\pi)$ (Hz), períodos $T_n = 1/f_n$ (s).
  \item Modos $\boldsymbol\phi_n$ M-ortonormales: $\boldsymbol\phi_i^\top \mathbf M\, \boldsymbol\phi_j = \delta_{ij}$.
  \item Ordenación ascendente: $\omega_1 \le \omega_2 \le \ldots \le \omega_{n_\text{modes}}$.
\end{itemize}


\section{Formulación numérica}

\subsection{5. Reducción por Dirichlet}
Por eliminación directa (ADR 0004) con operador $\mathbf T \in \mathbb R^{N_\text{dof} \times N_\text{libre}}$ que selecciona DOFs libres:
\[
\mathbf K_\text{red} = \mathbf T^\top \mathbf K\, \mathbf T, \qquad \mathbf M_\text{red} = \mathbf T^\top \mathbf M\, \mathbf T.
\]
El problema reducido $\mathbf K_\text{red}\,\boldsymbol\phi_\text{red} = \omega^2\, \mathbf M_\text{red}\,\boldsymbol\phi_\text{red}$ se resuelve sobre los DOFs libres. El término $\mathbf g_\text{indep}$ de restricciones lineales no homogéneas no aplica: los modos son solución del problema homogéneo (los modos de la estructura con apoyos a desplazamiento prescrito no nulo coinciden con los de apoyos nulos; el campo prescrito solo entra en la respuesta estática). Tras resolver, expansión $\boldsymbol\phi = \mathbf T\,\boldsymbol\phi_\text{red}$; en DOFs prescritos $\phi_i = 0$.

\subsection{6. Algoritmo numérico: Lanczos con shift-invert}
ARPACK vía \texttt{scipy.sparse.linalg.eigsh}:
\[
\text{eigsh}(\mathbf K_\text{red},\; k = n_\text{modes},\; \mathbf M = \mathbf M_\text{red},\; \sigma,\; \text{which}\!=\!\text{`LM'},\; \text{tol}).
\]
\begin{itemize}
  \item \textbf{Shift-invert} transforma el problema en $\bigl(\mathbf K_\text{red} - \sigma\,\mathbf M_\text{red}\bigr)^{-1}\mathbf M_\text{red}\,\boldsymbol\phi = \mu\,\boldsymbol\phi$ con $\mu = 1/(\omega^2 - \sigma)$. Buscar \texttt{which="LM"} (mayor magnitud de $\mu$) equivale a buscar autovalores cercanos al shift $\sigma$.
  \item $\sigma = 0$ por defecto: las frecuencias más bajas, que son las relevantes en ingeniería estructural.
  \item Internamente ARPACK factoriza $(\mathbf K_\text{red} - \sigma\,\mathbf M_\text{red})$ una sola vez y la reusa en cada iteración Lanczos.
\end{itemize}

\subsection{7. Normalización}
\texttt{eigsh} con $\mathbf M$ explícita devuelve autovectores M-ortonormales sin reescalado adicional. Signo arbitrario (autovector definido módulo signo); convención de signo se documenta solo si emerge como preferencia (no parte del contrato).

\subsection{8. Post-procesado}
\begin{itemize}
  \item $\omega_n = \sqrt{\lambda_n}$ (clip a cero los autovalores ligeramente negativos por ruido numérico, $|\lambda| < 10^{-14}\,\lambda_\text{max}$).
  \item $f_n = \omega_n / (2\pi)$, $T_n = 1/f_n$ (con $T_n = \infty$ si $\omega_n = 0$, modo de cuerpo rígido).
  \item Reordenación ascendente por $\omega_n$.
\end{itemize}


\section{Contrato de implementación}

\begin{lstlisting}[language=yaml]
name: ModalSolver
kind: solver
status: validated       # draft → implemented → validated

interface:
  yaml_type: modal      # solver: { type: modal }
  output: ModalResult   # campos: frequencies_hz, frequencies_rad, periods, modes, n_modes

parameters:
  - { name: n_modes,   type: int,   required: true,  desc: "Número de modos a calcular" }
  - { name: sigma,     type: float, required: false, default: 0.0,     desc: "Shift (rad²/s²) para shift-invert; 0 → frecuencias más bajas" }
  - { name: which,     type: str,   required: false, default: "LM",    desc: "Estrategia ARPACK: LM | SM | LA | SA | BE" }
  - { name: tolerance, type: float, required: false, default: 1.0e-9,  desc: "Tolerancia ARPACK sobre el residuo del autovalor" }
  - { name: lumping,   type: str,   required: false, default: "consistent", desc: "Discretización de masa; solo 'consistent' en fase 1" }
  - { name: linear_algebra, type: str, required: false, default: "auto", desc: "Backend para la factorización de (K - σM) interna a shift-invert (ADR 0003)" }

requirements:
  - "Todos los materiales del modelo declaran `density > 0` (ADR 0008). Densidad cero genera M_red singular → eigsh con sigma=0 falla; el solver propaga el error con mensaje físico."
  - "Al menos un DOF libre tras Dirichlet."
  - "n_modes < N_libre (ARPACK exige k < n)."

conventions:
  units:        "heredadas del modelo (ADR 0008); ω en rad/s, f en Hz, T en s consistentes con kg/N/m o equivalentes."
  frequencies:  "ω y f expuestos simultáneamente; el usuario consulta cualquiera."
  modes:        "M-ortonormales (Φᵀ M Φ = I); signo arbitrario por modo."
  ordering:     "ascendente en ω: ω₁ ≤ ω₂ ≤ … ≤ ω_n_modes."
  rigid_body:   "modos de cuerpo rígido aparecen con ω≈0; T=∞ se reporta sin error."

out_of_scope:
  - "Modos complejos con amortiguamiento no proporcional."
  - "Análisis modal sobre estado deformado (modos de pequeña amplitud alrededor de configuración prestressed) — diferido."
  - "Masa lumped y factores de participación / masas efectivas — implementados (fase 2 y fase 7 del ADR 0009, cerradas 2026-05-18); accesibles vía `lumping='lumped'` y `ResponseSpectrumSolver` respectivamente."

acceptance:
  verification:
    - name: barra_axial_empotrada_libre
      setup: |
        Barra 1D elástica empotrada en x=0, libre en x=L. Discretización con
        N=20 elementos Truss2D alineados con el eje x. E=210e9 Pa, ρ=7850 kg/m³,
        A=1e-4 m², L=1.0 m.
      expect: |
        Frecuencias propias analíticas (ondas axiales en barra empotrada-libre):
          ω_n = ((2n-1)π / (2L)) · sqrt(E/ρ),    n = 1, 2, 3, …
        Los primeros 3 modos calculados coinciden con la solución analítica.
      tol_rel: 0.01
    - name: viga_bernoulli_simplemente_apoyada
      setup: |
        Viga Bernoulli-Euler simplemente apoyada en sus dos extremos.
        Discretización con N=20 elementos Frame2DEuler alineados con el eje x.
        E=210e9 Pa, ρ=7850 kg/m³, A=1e-4 m², I=8.33e-10 m⁴, L=1.0 m.
      expect: |
        Frecuencias propias analíticas (flexión transversal):
          ω_n = (nπ/L)² · sqrt(E·I / (ρ·A)),    n = 1, 2, 3, …
        Los primeros 3 modos coinciden con la solución analítica.
      tol_rel: 0.02
  specific:
    - name: ortonormalidad_de_masa
      setup: "Cualquier modelo bien formado tras resolver el problema modal."
      expect: "Φᵀ M Φ = I (n_modes × n_modes)."
      tol_abs: 1.0e-10
    - name: ortogonalidad_de_rigidez
      setup: "Mismo modelo."
      expect: "Φᵀ K Φ = diag(ω²₁, …, ω²_n)."
      tol_rel: 1.0e-8
    - name: dirichlet_anula_modos_en_prescritos
      setup: "Barra empotrada-libre; nodo x=0 con ux=0."
      expect: "Componente del modo en el DOF prescrito = 0 exacto."
      tol_abs: 1.0e-14
    - name: density_faltante_lanza_value_error
      setup: "Material sin density declarada en el modelo."
      expect: "ModalSolver.solve() lanza ValueError listando los materiales afectados (mismo patrón que assemble_self_weight, ADR 0008)."

references:
  - "Bathe K.-J., Finite Element Procedures (2014), cap. 9 (matriz de masa consistente) y cap. 11 (algoritmos de autovalor)."
  - "Cook, Malkus & Plesha, Concepts and Applications of Finite Element Analysis, §11.3–§11.5."
  - "Hughes T. J. R., The Finite Element Method, cap. 7."
  - "Wilson E. L., Three-Dimensional Static and Dynamic Analysis of Structures, cap. 10."
  - "Lehoucq, Sorensen, Yang — ARPACK Users' Guide (Lanczos con shift-invert)."
  - "ADR 0003 — Capa algebraica (factorización reutilizable de (K − σM))."
  - "ADR 0008 — Material.density."
  - "ADR 0009 — Análisis modal y dinámico (este solver es la fase 1: el análisis modal estricto, problema generalizado K·φ = ω²·M·φ; el análisis dinámico transitorio entra en fases 3-7)."

\end{lstlisting}


\section{Implementación}

\begin{itemize}
  \item \textbf{Archivo}: solidum/math/solvers/modal.py (clase \texttt{ModalSolver}, registrada vía \texttt{@SolverRegistry.register}).
  \item \textbf{Backend algebraico}: solidum/math/linalg/eigen.py (clase \texttt{EigenSolver} envolviendo \texttt{scipy.sparse.linalg.eigsh}).
  \item \textbf{Tipo de resultado}: \texttt{ModalResult} --- dataclass frozen con \texttt{frequencies\_rad}, \texttt{frequencies\_hz}, \texttt{periods}, \texttt{modes}, \texttt{n\_modes}, \texttt{converged}. Expone \texttt{free\_vibration(M, u0, u0\_dot, t)} para reconstruir analíticamente la respuesta temporal de vibración libre sin amortiguamiento por superposición modal --- solución cerrada de \texttt{M\ensuremath{\cdot}ü + K\ensuremath{\cdot}u = 0} proyectada sobre los modos calculados (modos rígidos tratados aparte con evolución lineal \texttt{q\_n(t) = a\_n + b\_n\ensuremath{\cdot}t}).
  \item \textbf{Entrypoint público}: \texttt{solidum.run\_modal} para uso programático; \texttt{solidum.run\_yaml} despacha automáticamente a \texttt{run\_modal} cuando el YAML declara \texttt{solver: type: ModalSolver}.
  \item \textbf{Pipeline}: \texttt{Assembler.assemble\_system()} \ensuremath{\to} \texttt{Assembler.assemble\_mass\_matrix()} \ensuremath{\to} \texttt{Assembler.reduce\_pair(K, M)} \ensuremath{\to} \texttt{EigenSolver.solve(K\_red, M\_red, n\_modes)} \ensuremath{\to} expansión \texttt{\ensuremath{\Phi} = T \ensuremath{\cdot} \ensuremath{\Phi}\_red} \ensuremath{\to} \texttt{ModalResult}.
  \item \textbf{Caché}: \texttt{Assembler} cachea \texttt{M\_global} con el lumping usado; reusos posteriores del mismo análisis no recomputan.
  \item \textbf{Validación de density}: pre-chequeo agregado en \texttt{Assembler.assemble\_mass\_matrix} con el mismo patrón de ADR 0008 --- \texttt{assemble\_self\_weight}.
  \item \textbf{Tests}:
  \item tests/test\_modal.py --- 13 tests cubriendo:
  \item \texttt{TestModalAxialBar}: barra empotrada-libre, 20 elementos Truss2D, frecuencias contra \texttt{\ensuremath{\omega}\_n = ((2n-1)\ensuremath{\pi}/(2L))\ensuremath{\cdot}\ensuremath{\surd}(E/\ensuremath{\rho})} (tol 1\%).
  \item \texttt{TestModalSimplySupportedBeam}: viga biapoyada, 20 elementos Frame2DEuler, frecuencias contra \texttt{\ensuremath{\omega}\_n = (n\ensuremath{\pi}/L)\ensuremath{^{2}}\ensuremath{\cdot}\ensuremath{\surd}(EI/(\ensuremath{\rho}A))} (tol 2\%).
  \item \texttt{TestModalAlgebraicProperties}: ortonormalidad \texttt{\ensuremath{\Phi}\ensuremath{^{\top}}M\ensuremath{\Phi} = I} (atol 1e-10) y ortogonalidad \texttt{\ensuremath{\Phi}\ensuremath{^{\top}}K\ensuremath{\Phi} = diag(\ensuremath{\omega}\ensuremath{^{2}})} (rtol 1e-8) en ambos modelos.
  \item \texttt{TestModalSolverContract}: errores agregados (\texttt{density} faltante, \texttt{run\_modal} sin solver ni n\_modes) y verificación de que \texttt{lumping="lumped"} corre tras ADR 0009 fase 2 (cerrada 2026-05-18).
  \item \texttt{TestModalYamlPipeline}: pipeline end-to-end desde YAML.
  \item \textbf{Notas de traducción}:
  \item El parámetro \texttt{linear\_algebra} está en el constructor por coherencia con los demás solvers, pero en fase 1 no se usa: ARPACK gestiona internamente la factorización LU de \texttt{(K − \ensuremath{\sigma}M)} para shift-invert. Cuando se enchufe a la capa algebraica (ADR 0003), se podrá inyectar Cholesky o LDL\ensuremath{^{\top}}.
  \item El bloque \texttt{convergence:} del YAML (ADR 0007) se omite silenciosamente para \texttt{ModalSolver} igual que para \texttt{LinearSolver}: no hay iteración Newton ni norma de residuo configurable; la tolerancia es la propia de ARPACK (\texttt{tolerance:}).
  \item Matriz de masa consistente "translacional invariante rotacional" para \texttt{Truss*} y \texttt{Cable*} (\texttt{M = kron([[2,1],[1,2]] \ensuremath{\cdot} \ensuremath{\rho}AL/6, I\_dim)}). Para frames se ensambla la masa local axial + Hermitiana cúbica y se rota con \texttt{T\textasciicircum{}T \ensuremath{\cdot} M\_local \ensuremath{\cdot} T}. Para Frame3D se incluye masa torsional con momento polar geométrico \texttt{Jp = Iy + Iz} (no la constante de Saint-Venant \texttt{J}, que rige la rigidez). Para sólidos 2D, integración por la cuadratura del propio elemento (Tri3 usa fórmula analítica exacta porque su cuadratura 1-punto sería insuficiente).
  \item Inercia rotacional propia de sección (\texttt{\ensuremath{\rho}I\ensuremath{\cdot}L}) \textbf{incluida} en los DOFs rotacionales de Frame2D (Euler, Timoshenko, EulerCorot) y Frame3D --- coherente con Timoshenko y con la rigidez \texttt{K\_local} que tiene términos rotacionales puros (ADR 0009 §1).
\end{itemize}


\section{Diálogo}

\begin{itemize}
  \item \textbf{2026-05-13} \ensuremath{\cdot} Validado. Los 13 tests de \texttt{tests/test\_modal.py} cubren los criterios de \texttt{acceptance} (verificación analítica + propiedades algebraicas + contrato de errores + pipeline YAML) y pasan con las tolerancias declaradas. Promovido \texttt{status: validated}.
\end{itemize}


\newpage

\section{NewmarkSolver}



\textit{Orden de trabajo. La especificación física y la formulación numérica son patrimonio del usuario (revisa con detalle). La IA implementa y rellena las secciones marcadas.}



\section{Especificación física}

\subsection{0. Descripción general}
Análisis dinámico transitorio \textbf{lineal} de una estructura discretizada por FEM: integración temporal directa de la ecuación semidiscreta de movimiento bajo cargas dependientes del tiempo. Hipótesis: linealidad geométrica y material (K constante), masa M constante (lagrangeano total), amortiguamiento Rayleigh proporcional \texttt{C = \ensuremath{\alpha}\ensuremath{\cdot}M + \ensuremath{\beta}\ensuremath{\cdot}K}. Apoyos Dirichlet constantes en el tiempo. La no-linealidad (Newton dentro de cada paso) se difiere a la fase 4 del ADR 0009.

\subsection{1. Ecuación de movimiento semidiscreta}
Tras la discretización espacial por FEM, el problema continuo
\[
\rho\,\ddot{\mathbf u} - \operatorname{div}\boldsymbol\sigma = \mathbf b(\mathbf x, t)
\]
se reduce al sistema acoplado de ODEs:
\[
\mathbf M\,\ddot{\mathbf u}(t) + \mathbf C\,\dot{\mathbf u}(t) + \mathbf K\,\mathbf u(t) = \mathbf F(t),
\]
con condiciones iniciales $\mathbf u(0) = \mathbf u_0$, $\dot{\mathbf u}(0) = \dot{\mathbf u}_0$ y condiciones de contorno Dirichlet $\mathbf u_s(t) = \mathbf g$ constantes.

\subsection{2. Amortiguamiento Rayleigh}
\[
\mathbf C = \alpha\,\mathbf M + \beta\,\mathbf K,
\]
con coeficientes calibrados a partir de dos pares modales $(\xi_1, \omega_1)$ y $(\xi_2, \omega_2)$:
\[
\alpha = \frac{2\,\omega_1\omega_2\,(\xi_1\omega_2 - \xi_2\omega_1)}{\omega_2^2 - \omega_1^2}, \qquad
  \beta = \frac{2\,(\xi_2\omega_2 - \xi_1\omega_1)}{\omega_2^2 - \omega_1^2}.
\]
La razón de amortiguamiento del modo $n$-ésimo resulta $\xi_n = \tfrac{1}{2}(\alpha/\omega_n + \beta\omega_n)$ --- exacta en $\omega_1, \omega_2$ y suavemente sobreamortiguada fuera del rango. El usuario también puede pasar \texttt{(\ensuremath{\alpha}, \ensuremath{\beta})} directos.

\subsection{3. Salidas}
\begin{itemize}
  \item $\mathbf u(t_k)$, $\dot{\mathbf u}(t_k)$, $\ddot{\mathbf u}(t_k)$ para cada paso temporal $k = 0, 1, \ldots, N_t$.
  \item Vector de instantes $t_k = k\,\Delta t$.
  \item Diagnósticos: $\alpha, \beta$ Rayleigh efectivos; número de pasos; tiempo de factorización vs tiempo de pasos.
\end{itemize}


\section{Formulación numérica}

\subsection{4. Método de Newmark-\ensuremath{\beta} (a-form)}
Familia paramétrica con $(\beta, \gamma)$. Predictores:
\[
\tilde{\mathbf u}_{n+1} = \mathbf u_n + \Delta t\,\dot{\mathbf u}_n + \tfrac{\Delta t^2}{2}(1 - 2\beta)\,\ddot{\mathbf u}_n,
\]
\[
\tilde{\dot{\mathbf u}}_{n+1} = \dot{\mathbf u}_n + \Delta t\,(1 - \gamma)\,\ddot{\mathbf u}_n.
\]

Sistema efectivo en aceleraciones:
\[
\bigl(\mathbf M + \gamma\Delta t\,\mathbf C + \beta\Delta t^2\,\mathbf K\bigr)\,\ddot{\mathbf u}_{n+1} \;=\; \mathbf F_{n+1} - \mathbf C\,\tilde{\dot{\mathbf u}}_{n+1} - \mathbf K\,\tilde{\mathbf u}_{n+1}.
\]

Correctores:
\[
\mathbf u_{n+1} = \tilde{\mathbf u}_{n+1} + \beta\Delta t^2\,\ddot{\mathbf u}_{n+1}, \qquad
  \dot{\mathbf u}_{n+1} = \tilde{\dot{\mathbf u}}_{n+1} + \gamma\Delta t\,\ddot{\mathbf u}_{n+1}.
\]

\subsection{5. Aceleración inicial}
$\ddot{\mathbf u}_0$ se calcula consistentemente con la ecuación de movimiento en $t = 0$:
\[
\mathbf M\,\ddot{\mathbf u}_0 = \mathbf F(0) - \mathbf C\,\dot{\mathbf u}_0 - \mathbf K\,\mathbf u_0.
\]

\subsection{6. Parámetros canónicos}
\begin{center}
\begin{tabularx}{\textwidth}{YYYY}
\hline
\textbf{Esquema} & \textbf{\ensuremath{\beta}} & \textbf{\ensuremath{\gamma}} & \textbf{Propiedades} \\
\hline
\textbf{Average acceleration (default)} & 1/4 & 1/2 & Incondicionalmente estable, sin amortiguamiento numérico, exactitud $O(\Delta t^2)$ \\
Linear acceleration & 1/6 & 1/2 & Condicionalmente estable ($\Delta t \le 2\sqrt{3}/\omega_{\max}$), $O(\Delta t^2)$ \\
Fox-Goodwin & 1/12 & 1/2 & Condicionalmente estable, $O(\Delta t^4)$ \\
Diferencias centradas & 0 & 1/2 & Explícito; condicionalmente estable ($\Delta t \le 2/\omega_{\max}$), $O(\Delta t^2)$ \\
\hline
\end{tabularx}
\end{center}
\subsection{7. Imposición de Dirichlet con apoyos constantes}
Mismo mecanismo que estática (ADR 0004 fase 1) con operador $\mathbf T$ que selecciona DOFs libres y $\mathbf g$ que recoge los valores prescritos: $\mathbf u = \mathbf T\,\mathbf u_\text{free} + \mathbf g$. Como $\mathbf g$ es constante, $\dot{\mathbf u} = \mathbf T\,\dot{\mathbf u}_\text{free}$ y $\ddot{\mathbf u} = \mathbf T\,\ddot{\mathbf u}_\text{free}$. El sistema efectivo se proyecta a DOFs libres:
\[
\mathbf A_\text{eff,red}\,\ddot{\mathbf u}_{\text{free},n+1} = \mathbf T^\top\bigl(\mathbf F_{n+1} - \mathbf K\,\mathbf g\bigr) - \mathbf C_\text{red}\,\tilde{\dot{\mathbf u}}_\text{free} - \mathbf K_\text{red}\,\tilde{\mathbf u}_\text{free},
\]
donde $\mathbf A_\text{eff,red} = \mathbf M_\text{red} + \gamma\Delta t\,\mathbf C_\text{red} + \beta\Delta t^2\,\mathbf K_\text{red}$. El término $\mathbf T^\top\,\mathbf K\,\mathbf g$ se precalcula una vez. La expansión final a globales es $\mathbf u_n = \mathbf T\,\mathbf u_{\text{free},n} + \mathbf g$.

\subsection{8. Factorización reutilizable}
Con $\Delta t$ constante y problema lineal, $\mathbf A_\text{eff,red}$ es constante. Se factoriza una vez al inicio (ADR 0003 --- \texttt{FactorizedSolver}) y cada paso es una resolución triangular barata. Para $N_t$ pasos sobre $N_\text{libre}$ DOFs: coste $\sim O(\text{factor})_\text{una vez} + N_t \cdot O(N_\text{libre})_\text{por paso}$.


\section{Contrato de implementación}

\begin{lstlisting}[language=yaml]
name: NewmarkSolver
kind: solver
status: validated       # draft → implemented → validated

interface:
  yaml_type: NewmarkSolver
  output: TransientResult   # campos: t_history, u_history, udot_history, uddot_history

parameters:
  - { name: t_end,     type: float, required: true,                desc: "Tiempo final del análisis (s)" }
  - { name: dt,        type: float, required: true,                desc: "Paso temporal constante (s)" }
  - { name: beta,      type: float, required: false, default: 0.25, desc: "Parámetro β de Newmark (default: average acceleration)" }
  - { name: gamma,     type: float, required: false, default: 0.5,  desc: "Parámetro γ de Newmark" }
  - { name: rayleigh,  type: dict,  required: false, default: null, desc: "Amortiguamiento Rayleigh: {alpha, beta} directos o {xi1, omega1, xi2, omega2} para calibración modal" }
  - { name: u0,        type: ndarray, required: false, default: null, desc: "Desplazamiento inicial (ndof,); cero por defecto" }
  - { name: u0_dot,    type: ndarray, required: false, default: null, desc: "Velocidad inicial (ndof,); cero por defecto" }
  - { name: F_func,    type: callable, required: false, default: null, desc: "Callback Python F_func(t) -> ndarray (ndof,); cero por defecto" }
  - { name: linear_algebra, type: str, required: false, default: "auto", desc: "Backend para factorizar A_eff (ADR 0003)" }
  - { name: lumping,   type: str, required: false, default: "consistent", desc: "Discretización de masa; solo 'consistent' en esta fase" }

requirements:
  - "Todos los materiales del modelo declaran `density > 0` (ADR 0008)."
  - "K lineal y constante (hipótesis de pequeñas deformaciones, material elástico)."
  - "Apoyos Dirichlet constantes en el tiempo."

conventions:
  units:   "heredadas del modelo (ADR 0008); t en s, ω en rad/s, F en N consistente con kg/N/m."
  damping: "Rayleigh proporcional: C = α·M + β·K. Si se pasan (ξ₁,ω₁),(ξ₂,ω₂), (α,β) se derivan automáticamente."
  stability: |
    Default (β=1/4, γ=1/2) es incondicionalmente estable. Para β=0 (diferencias
    centradas) o β=1/6 (lineal) el usuario es responsable de que Δt ≤ 2/ω_max;
    el solver no comprueba esta cota automáticamente.

out_of_scope:
  - "Apoyos dependientes del tiempo (multi-support excitation sísmica) — diferido."
  - "Δt adaptativo — diferido."
  - "Generalized-α, Bossak-α — diferidos (variantes adicionales de la familia con amortiguamiento numérico)."
  - "No-linealidad y disipación numérica controlada — implementados como variantes en clases hermanas (`NewtonNewmarkSolver`, `HHTSolver`, `NewtonHHTSolver`); fases 4 y 3-bis del ADR 0009, cerradas."

acceptance:
  verification:
    - name: oscilador_1gdl_no_amortiguado
      setup: |
        Sistema 1 GDL: M=1, K=ω²=25 (ω=5 rad/s), C=0. Condiciones iniciales
        u₀=1, u̇₀=0. F(t)=0. Δt=ω·T/100 (100 pasos por período). Integrar
        durante 4 períodos.
      expect: |
        u(t) = cos(ω·t) exacto. Error máximo en u durante el intervalo
        < 1% (β=1/4 tiene error de período O(Δt²) bajo control con 100
        pasos por período).
      tol_rel: 0.01
    - name: oscilador_1gdl_amortiguado_rayleigh
      setup: |
        Sistema 1 GDL: M=1, K=25, C=2·ξ·ω·M con ξ=0.05 (sub-amortiguado).
        u₀=1, u̇₀=0, F=0. Δt = T/100, t_end = 4·T. Construir C vía
        Rayleigh con un único par (ξ, ω).
      expect: |
        u(t) = e^(-ξωt)·(cos(ω_d·t) + (ξω/ω_d)·sin(ω_d·t)) con
        ω_d = ω·√(1-ξ²). Error máximo < 2 %.
      tol_rel: 0.02
    - name: viga_voladizo_carga_subita_modal_vs_newmark
      setup: |
        Voladizo Frame2DEuler (20 elementos) con carga puntual aplicada en
        el extremo en t=0 como step function. Δt = T₁/40 (T₁ período del
        primer modo), t_end = 2·T₁. Sin amortiguamiento.
      expect: |
        La reconstrucción modal de la respuesta (superposición de primeros
        5 modos forzados por F(t)=H(t)·F₀ con solución analítica conocida)
        coincide con Newmark a < 5 %.
      tol_rel: 0.05
    - name: orden_de_convergencia
      setup: |
        Oscilador 1 GDL no amortiguado integrado con β=1/4, γ=1/2 a
        Δt, Δt/2, Δt/4. Calcular el error máximo en u contra solución
        analítica.
      expect: |
        El cociente de errores e(Δt)/e(Δt/2) tiende a 4 (orden 2 exacto).
      tol_rel: 0.5  # tolerancia sobre el cociente: 4 ± 2
  specific:
    - name: rayleigh_calibration
      setup: "Dados (ξ₁=0.02, ω₁=10), (ξ₂=0.05, ω₂=100), calibrar (α, β)."
      expect: |
        α = 2·ω₁·ω₂·(ξ₁·ω₂ − ξ₂·ω₁)/(ω₂² − ω₁²) ≈ 0.303,
        β = 2·(ξ₂·ω₂ − ξ₁·ω₁)/(ω₂² − ω₁²) ≈ 9.6e-4.
        Y verificar ξ(ω_n) = (α/ω_n + β·ω_n)/2 reproduce ξ₁ en ω₁ y ξ₂ en ω₂.
      tol_rel: 1.0e-10

references:
  - "Newmark N. M. (1959), A method of computation for structural dynamics, J. Eng. Mech. Div. ASCE."
  - "Bathe K.-J., Finite Element Procedures (2014), cap. 9 (matriz de masa) y cap. 9.4 (Newmark)."
  - "Hughes T. J. R., The Finite Element Method, cap. 9 (algoritmos de integración temporal)."
  - "Chopra A. K., Dynamics of Structures, cap. 5 (1 GDL) y cap. 15 (MDOF Newmark)."
  - "ADR 0003 — Capa algebraica (factorización reutilizable de A_eff)."
  - "ADR 0007 — Convergencia y tolerancias."
  - "ADR 0009 — Análisis modal y dinámico (este solver es la fase 3)."

\end{lstlisting}


\section{Implementación}

\begin{itemize}
  \item \textbf{Archivo}: solidum/math/solvers/newmark.py (clase \texttt{NewmarkSolver}, registrada vía \texttt{@SolverRegistry.register}).
  \item \textbf{Amortiguamiento}: solidum/math/damping.py (\texttt{rayleigh\_from\_modes}, \texttt{rayleigh\_xi}).
  \item \textbf{Tipo de resultado}: \texttt{TransientResult} --- dataclass frozen con \texttt{t\_history}, \texttt{u\_history}, \texttt{udot\_history}, \texttt{uddot\_history}, \texttt{n\_steps}, \texttt{alpha\_rayleigh}, \texttt{beta\_rayleigh}, \texttt{converged}.
  \item \textbf{Entrypoint público}: \texttt{solidum.run\_transient} para uso programático; \texttt{solidum.run\_yaml} despacha automáticamente cuando el YAML declara \texttt{solver: type: NewmarkSolver}.
  \item \textbf{Pipeline}: \texttt{Assembler.assemble\_system()} \ensuremath{\to} \texttt{Assembler.assemble\_mass\_matrix()} \ensuremath{\to} \texttt{C = \ensuremath{\alpha}\ensuremath{\cdot}M + \ensuremath{\beta}\ensuremath{\cdot}K} \ensuremath{\to} reducción simultánea por Dirichlet (operador \texttt{T}) \ensuremath{\to} cálculo de \texttt{ü\ensuremath{_{0}}} \ensuremath{\to} factorización única de \texttt{A\_eff\_red = M\_red + \ensuremath{\gamma}\ensuremath{\Delta}t\ensuremath{\cdot}C\_red + \ensuremath{\beta}\ensuremath{\Delta}t\ensuremath{^{2}}\ensuremath{\cdot}K\_red} \ensuremath{\to} bucle de pasos con resolución triangular reutilizada.
  \item \textbf{Tests}:
  \item tests/test\_newmark.py --- 14 tests:
  \item \texttt{TestRayleighCalibration}: \texttt{(\ensuremath{\alpha}, \ensuremath{\beta})} reproduce \ensuremath{\xi} objetivo en \ensuremath{\omega}\ensuremath{_{1}}, \ensuremath{\omega}\ensuremath{_{2}}; errores agregados (\ensuremath{\omega} iguales, no positivos).
  \item \texttt{TestNewmark1DofUndamped}: vibración libre \texttt{u(t) = cos(\ensuremath{\omega}t)} a 1 \%; estado inicial preservado.
  \item \texttt{TestNewmark1DofDamped}: respuesta sub-amortiguada \texttt{e\textasciicircum{}(-\ensuremath{\xi}\ensuremath{\omega}t)\ensuremath{\cdot}(cos(\ensuremath{\omega}\_d\ensuremath{\cdot}t) + ...)} a 2 \%; coeficientes Rayleigh propagados al \texttt{TransientResult}.
  \item \texttt{TestNewmark1DofStepResponse}: \texttt{u(t) = (F\ensuremath{_{0}}/K)(1 − cos(\ensuremath{\omega}t))} a 0.1 \%.
  \item \texttt{TestNewmarkConvergenceOrder}: cociente de errores al refinar \ensuremath{\Delta}t \ensuremath{\in} {T/40, T/80, T/160} cae en \texttt{[3.5, 4.5]} \ensuremath{\to} orden 2 confirmado.
  \item \texttt{TestNewmarkContract}: validaciones de constructor y \texttt{run\_transient}.
  \item \texttt{TestNewmarkYamlPipeline}: end-to-end desde YAML.
  \item \textbf{Notas de traducción}:
  \item \textbf{Factorización única}: con \texttt{\ensuremath{\Delta}t} constante y problema lineal, \texttt{A\_eff\_red} es invariante; se factoriza al inicio (\texttt{A\_solver.factorize(A\_eff)}) y cada paso es una sustitución triangular barata (ADR 0003 --- \texttt{FactorizedSolver}).
  \item \textbf{\texttt{reduce\_pair} vs \texttt{reduce} triple}: en el problema modal \texttt{reduce\_pair(K, M)} basta porque \texttt{g} no entra en autovalores. En Newmark se manejan tres matrices (\texttt{K, M, C}) y un término constante \texttt{F\_dir = T\ensuremath{^{\top}}K\ensuremath{\cdot}g} para apoyos no nulos; se hace explícitamente en \texttt{NewmarkSolver.solve()} sin un helper unificado (no aporta valor centralizarlo con un solo consumidor).
  \item \textbf{Aceleración inicial consistente}: \texttt{M\ensuremath{\cdot}ü\ensuremath{_{0}} = F(0) − C\ensuremath{\cdot}u̇\ensuremath{_{0}} − K\ensuremath{\cdot}u\ensuremath{_{0}}} resuelve un sistema lineal aparte; se usa el dispatcher de la capa algebraica sobre \texttt{M\_red} (siempre SPD si la masa es consistente y la densidad estrictamente positiva).
  \item \textbf{\texttt{F\_func} opcional}: si es \texttt{None}, se trata como vector cero (vibración libre). El callback se invoca una vez por paso con \texttt{t\_\{n+1\}}; el usuario es responsable de la coherencia temporal (no se reusan valores entre pasos).
  \item \textbf{El bloque \texttt{convergence:} del YAML} (ADR 0007) se omite para \texttt{NewmarkSolver} igual que para \texttt{LinearSolver} y \texttt{ModalSolver}: la fase 3 es lineal y no hay iteración Newton interna; cuando entre la fase 4 (Newton dentro de cada paso), el bloque se reincorporará.
\end{itemize}


\section{Diálogo}

\begin{itemize}
  \item \textbf{2026-05-13} \ensuremath{\cdot} Validado. Los 14 tests de \texttt{tests/test\_newmark.py} cubren \texttt{acceptance.verification} (osciladores 1 GDL no amortiguado / amortiguado / step / orden de convergencia) y \texttt{acceptance.specific} (calibración Rayleigh). Promovido \texttt{status: validated}.
\end{itemize}


\newpage

\section{NewtonHHTSolver}



\textit{Variante de \texttt{NewtonNewmarkSolver} según la regla}

\textit{de variantes (Reglas.md §4). Aplica el esquema temporal HHT-\ensuremath{\alpha} (Hilber-}

\textit{Hughes-Taylor 1977) al residuo dinámico no lineal: amortiguamiento numérico}

\textit{controlado en altas frecuencias dentro del bucle Newton de cada paso.}

\textit{> Reusa todas las decisiones arquitecturales del ADR 0009 (esquema temporal,}

\textit{contrato \texttt{PIPELINE\_KIND}, telemetría tipada, line search opt-in) y de}

\textit{\texttt{HHTSolver} (formulación temporal, autoderivación de \ensuremath{\beta}/\ensuremath{\gamma}).}

\textit{\textbf{Sin ADR nuevo}. Documenta solo lo que cambia respecto a los dos padres.}



\section{Qué cambia respecto a \texttt{NewtonNewmarkSolver}}

\texttt{NewtonNewmarkSolver} integra \texttt{M\ensuremath{\cdot}ü + C\ensuremath{\cdot}u̇ + F\_int(u) = F(t)} con Newmark-\ensuremath{\beta}
clásico (\ensuremath{\beta}=1/4, \ensuremath{\gamma}=1/2 default, sin amortiguamiento numérico). Cada paso
abre un bucle Newton porque \texttt{F\_int(u)} es no lineal en \texttt{u} (plasticidad,
daño, corotacional, embedded).

\texttt{NewtonHHTSolver} reemplaza el residuo Newmark por el residuo HHT-\ensuremath{\alpha}: las
\textbf{fuerzas internas, viscosas y externas} se evalúan en un instante
intermedio \texttt{t\_\{n+1-\ensuremath{\alpha}\}} mientras la \textbf{fuerza inercial} se mantiene en
\texttt{t\_\{n+1\}}. El bucle Newton se ejecuta sobre este residuo modificado.

\subsection{Residuo dinámico HHT-\ensuremath{\alpha} no lineal}

\[
\mathbf R(\ddot{\mathbf u}_{n+1}) \;=\; (1+\alpha)\,\mathbf F_{n+1} - \alpha\,\mathbf F_n
   - \big[(1+\alpha)\,\mathbf F_{\text{int}}(\mathbf u_{n+1}) - \alpha\,\mathbf F_{\text{int}}(\mathbf u_n)\big]
   - \big[(1+\alpha)\,\mathbf C\,\dot{\mathbf u}_{n+1} - \alpha\,\mathbf C\,\dot{\mathbf u}_n\big]
   - \mathbf M\,\ddot{\mathbf u}_{n+1}
\]

El término no lineal $\mathbf F_{\text{int}}(\mathbf u_{n+1})$ depende del
estado actualizado por los correctores Newmark, así que el residuo es
implícito en $\ddot{\mathbf u}_{n+1}$ y requiere Newton.

\subsection{Jacobiano efectivo}

\[
\mathbf J \;=\; \mathbf M + (1+\alpha)\,\gamma\,\Delta t\,\mathbf C
                            + (1+\alpha)\,\beta\,\Delta t^2\,\mathbf K_t
\]

donde $\mathbf K_t = \partial\mathbf F_{\text{int}}/\partial\mathbf u$ es
la tangente material/geométrica reensamblada en cada iteración Newton.

\subsection{Parámetros heredados de \texttt{HHTSolver}}

\begin{itemize}
  \item Rango admisible $\alpha \in [-1/3, 0]$. Valores fuera lanzan \texttt{ValueError}.
  \item $\beta$ y $\gamma$ auto-derivados desde $\alpha$:
\end{itemize}

\[
\beta = \frac{(1-\alpha)^2}{4}, \qquad \gamma = \frac{1-2\alpha}{2}
\]

\begin{itemize}
  \item Override explícito de $\beta$/$\gamma$ permitido pero genera \texttt{RuntimeWarning}
\end{itemize}
  (combinaciones distintas pueden romper estabilidad incondicional u orden 2).
\begin{itemize}
  \item Radio espectral en alta frecuencia: $\rho_\infty = (1+\alpha)/(1-\alpha)$.
\end{itemize}
  Disipación selectiva en modos con $\omega\Delta t > \pi$.

\subsection{Comportamiento que NO cambia}

Todo lo heredado de \texttt{NewtonNewmarkSolver}:

\begin{itemize}
  \item Convergencia dual fuerza + desplazamiento (ADR 0007) con autocalibración.
  \item Line search GLL opt-in (ADR 0011, default \texttt{False}).
  \item Diagnóstico tipado de divergencia (\texttt{SolverDivergedError} y subclases).
  \item Rayleigh con $\mathbf K_0$ constante (matriz de rigidez inicial).
  \item Commit/rollback del estado material por paso.
  \item Detección de oscilación por ventana móvil.
  \item Despacho algebraico ADR 0003 con fallback Cholesky\ensuremath{\to}LU.
  \item Lumping de masa (consistent / lumped) ADR 0009 fase 2.
\end{itemize}

\subsection{Recuperación de \texttt{NewtonNewmark} con $\alpha = 0$}

Con $\alpha = 0$: $\beta = 1/4$, $\gamma = 1/2$, el residuo HHT colapsa al
de Newmark estándar. Verificado en tests:
\texttt{tests/test\_dynamic\_nonlinear\_hht.py::test\_alpha\_zero\_recovers\_newton\_newmark}.

\subsection{Recuperación de \texttt{HHTSolver} con materiales lineales}

Con materiales lineales ($\mathbf F_{\text{int}}(\mathbf u) = \mathbf K\,\mathbf u$),
el residuo de Newton se anula en una iteración y el solver reproduce
exactamente \texttt{HHTSolver}. Tangente constante = rigidez secante. Verificado:
\texttt{tests/test\_dynamic\_nonlinear\_hht.py::test\_linear\_material\_matches\_hht\_solver}.


\section{Contrato de implementación}

\begin{lstlisting}[language=yaml]
name: NewtonHHTSolver
kind: solver
status: validated

extends: NewtonNewmarkSolver       # variante por Reglas §4
also_references: HHTSolver         # formulación temporal heredada

interface:
  yaml_type: newton_hht
  output: TransientResult
  pipeline_kind: transient

parameters:
  # Solo se listan los parámetros propios o con default distinto al padre.
  # Resto heredados de NewtonNewmarkSolver.
  - { name: alpha, type: float, required: false, default: -0.05,
      desc: "Parámetro HHT-α ∈ [-1/3, 0]. α=0 recupera NewtonNewmark; α=-1/3 disipación máxima preservando orden 2" }
  - { name: beta,  type: float, required: false, default: null,
      desc: "Override del β autoderivado (no recomendado — emite RuntimeWarning)" }
  - { name: gamma, type: float, required: false, default: null,
      desc: "Override del γ autoderivado (no recomendado — emite RuntimeWarning)" }
  # heredados de NewtonNewmarkSolver:
  #   t_end, dt, convergence, max_iter, freeze_tangent_after_iter,
  #   line_search, rayleigh, u0, u0_dot, F_func, linear_algebra, lumping

requirements:
  - "Modelo dinámico con no-linealidad de material o geometría"
  - "M con density > 0 en todos los materiales activos (ADR 0008)"
  - "Δt razonable (incondicionalmente estable para sistemas lineales, no garantizado para fuertes no-linealidades)"

conventions:
  units: "heredadas del modelo (ADR 0008)"
  stability: "incondicionalmente estable en problemas lineales para α ∈ [-1/3, 0]; condicionalmente para no lineales severos"
  disipacion: "ρ_∞ = (1+α)/(1-α); selectiva en alta frecuencia"

out_of_scope:
  - "α fuera de [-1/3, 0] (pierde estabilidad u orden 2)"
  - "Generalized-α / Bossak (variantes con disipación distinta) — diferidos"
  - "Análisis estático ⇒ usar NonlinearSolver"

acceptance:
  verification:
    - name: alpha_cero_recupera_newton_newmark
      setup: "mismo problema con NewtonNewmarkSolver y NewtonHHTSolver(α=0)"
      expect: "respuesta idéntica paso a paso"
      tol_rel: 1.0e-12

    - name: material_lineal_recupera_hht
      setup: "mismo problema con HHTSolver y NewtonHHTSolver, ambos α=−0.1"
      expect: "respuesta idéntica; Newton converge en 1 iteración por paso"
      tol_rel: 1.0e-12

    - name: disipacion_modo_alta_frecuencia
      setup: "respuesta libre con modos altos espurios excitados; α=−0.1"
      expect: "amplitud de modos espurios decae ~ρ_∞ por periodo"
      tol_rel: 0.05

    - name: respuesta_plastica_dinamica
      setup: "barra Truss2D con Elastoplastic1D, carga dinámica que cruza el yield"
      expect: "estado plástico converged paso a paso; histéresis dinámica coherente"
      tol_rel: 1.0e-3

    - name: rechazo_alpha_fuera_de_rango
      setup: "NewtonHHTSolver(alpha=-0.5) o (alpha=0.1)"
      expect: "ValueError con mensaje claro"

    - name: warning_override_beta_gamma
      setup: "NewtonHHTSolver(alpha=-0.1, beta=0.3)"
      expect: "RuntimeWarning emitido al construir"

  specific:
    - name: tangente_singular_dispara_excepcion_tipada
      setup: "modelo con material que produce K_t singular en algún paso"
      expect: "raise SingularTangentError con métricas estructuradas"

references:
  - "Hilber H.M., Hughes T.J.R., Taylor R.L. (1977). Improved numerical dissipation for time integration algorithms in structural dynamics. Earthquake Eng. Struct. Dyn. 5, 283-292."
  - "Hughes T.J.R. (2000). The Finite Element Method. Dover. §9.3."
  - "Newmark N.M. (1959). A method of computation for structural dynamics. ASCE 85(EM3), 67-94."
  - "ADR 0009 (fase 4: HHT-α — formulación temporal en residuo no lineal)."
  - "ADR 0011 (excepciones tipadas, line search opt-in)."
  - "Spec padre: HHTSolver.md (formulación temporal lineal)."
  - "Spec padre: NewtonNewmarkSolver.md (Newton dentro de Newmark no lineal)."

\end{lstlisting}


\section{Implementación}

\begin{itemize}
  \item \textbf{Archivo}: solidum/math/solvers/newmark.py (clase \texttt{NewtonHHTSolver}).
  \item \textbf{Clase}: hereda de \texttt{NewtonNewmarkSolver}; añade el atributo \texttt{self.alpha} y reemplaza el cómputo del residuo y del jacobiano en \texttt{solve()}.
  \item \textbf{Auto-derivación} de \ensuremath{\beta} y \ensuremath{\gamma} vía \texttt{\_hht\_autoderive\_beta\_gamma(alpha)} (función helper compartida con \texttt{HHTSolver}).
  \item \textbf{Validación temprana} en \texttt{\_\_init\_\_}: \texttt{alpha \ensuremath{\in} [-1/3, 0]}, warning si override de \ensuremath{\beta}/\ensuremath{\gamma}.
  \item \textbf{Entrypoint público}: \texttt{solidum.run\_transient(model, solver="newton\_hht", ...)} (despacho por \texttt{PIPELINE\_KIND = "transient"}).
  \item \textbf{Tests}:
  \item tests/test\_dynamic\_nonlinear\_hht.py --- incluye recuperación de \texttt{NewtonNewmark} con \ensuremath{\alpha}=0, recuperación de \texttt{HHTSolver} con materiales lineales, respuesta plástica dinámica.
\end{itemize}


\section{Diálogo}

\begin{itemize}
  \item \textbf{2026-05-19} \ensuremath{\cdot} Spec corta tipo extensión creada para cerrar el hueco H-5.3 (variante de \texttt{HHTSolver}/\texttt{NewtonNewmarkSolver} por Reglas §4). Sin ADR nuevo: reúsa ADR 0009 (formulación temporal HHT) y ADR 0011 (Newton + diagnóstico). El comportamiento no se modifica.
\end{itemize}


\newpage

\section{NewtonNewmarkSolver}



\textit{Variante de \texttt{NewmarkSolver} según la regla de variantes}

\textit{(Reglas.md §4). Documenta \textbf{solo lo que cambia} respecto al padre. Implementa}

\textit{la \textbf{fase 4 del ADR 0009}; reusa todas las decisiones arquitecturales del}

\textit{ADR padre --- sin ADR nuevo.}



\section{Qué cambia respecto a \texttt{NewmarkSolver}}

\texttt{NewmarkSolver} integra \texttt{M\ensuremath{\cdot}ü + C\ensuremath{\cdot}u̇ + K\ensuremath{\cdot}u = F(t)} asumiendo \textbf{K constante}:
factoriza \texttt{A\_eff = M + \ensuremath{\gamma}\ensuremath{\Delta}t\ensuremath{\cdot}C + \ensuremath{\beta}\ensuremath{\Delta}t\ensuremath{^{2}}\ensuremath{\cdot}K} una vez y cada paso temporal es una
resolución triangular barata.

\texttt{NewtonNewmarkSolver} resuelve el mismo problema cuando **la rigidez no es
constante**: hay materiales con historia (plasticidad, daño) o no linealidad
geométrica. En cada paso temporal se introduce un bucle de \textbf{Newton-Raphson}
sobre el residuo dinámico, hasta cumplir el criterio de convergencia.

\subsection{Residuo dinámico en cada paso}

En el instante $t_{n+1}$ las incógnitas son $(\mathbf u_{n+1}, \dot{\mathbf u}_{n+1}, \ddot{\mathbf u}_{n+1})$
ligadas por los correctores Newmark:

\[
\mathbf u_{n+1} = \tilde{\mathbf u}_{n+1} + \beta\Delta t^2\,\ddot{\mathbf u}_{n+1}, \qquad
  \dot{\mathbf u}_{n+1} = \tilde{\dot{\mathbf u}}_{n+1} + \gamma\Delta t\,\ddot{\mathbf u}_{n+1}
\]

(con predictores idénticos a la fase 3). El residuo dinámico es:

\[
\mathbf R(\ddot{\mathbf u}_{n+1}) \;=\; \mathbf F_\text{ext}(t_{n+1}) \;-\; \mathbf F_\text{int}(\mathbf u_{n+1}) \;-\; \mathbf C\,\dot{\mathbf u}_{n+1} \;-\; \mathbf M\,\ddot{\mathbf u}_{n+1} \;-\; \mathbf F_\text{dir}
\]

donde $\mathbf F_\text{int}(\mathbf u)$ es el vector de fuerzas internas no lineales
(plástico, dañado, etc.) que el \texttt{Assembler.assemble\_non\_linear\_system} ya
calcula para los solvers estáticos (ADR 0003 + ADR 0006 + ADR 0008).

\subsection{Jacobiano dinámico tangente}

Derivando $\mathbf R$ respecto a $\ddot{\mathbf u}_{n+1}$ usando los correctores:

\[
\mathbf J \;=\; -\frac{\partial \mathbf R}{\partial \ddot{\mathbf u}_{n+1}} \;=\; \mathbf M \;+\; \gamma\Delta t\,\mathbf C \;+\; \beta\Delta t^2\,\mathbf K_\text{t}
\]

donde $\mathbf K_\text{t} = \partial \mathbf F_\text{int}/\partial \mathbf u$ es la
rigidez tangente del estado corriente. La iteración es:

\[
\mathbf J\,\delta \ddot{\mathbf u} = \mathbf R, \qquad \ddot{\mathbf u}_{n+1}^{(k+1)} = \ddot{\mathbf u}_{n+1}^{(k)} + \delta \ddot{\mathbf u}
\]

Los desplazamientos y velocidades se actualizan con los mismos correctores
Newmark tras cada update de $\ddot{\mathbf u}_{n+1}$.

\subsection{Amortiguamiento Rayleigh --- heredado, constante en el tiempo}

$\mathbf C = \alpha\,\mathbf M + \beta_R\,\mathbf K_0$ con $\mathbf K_0$ la
rigidez \textbf{elástica de referencia} (la primera ensamblada en $t = 0$, $\mathbf u = 0$).
No se actualiza con la plasticidad: el Rayleigh es modelo aproximado de
amortiguamiento estructural, no consecuencia de la cinemática elastoplástica.
Mantenerlo constante evita un acoplamiento ad-hoc que ningún código de
referencia (Abaqus, ANSYS, OpenSees) realiza por defecto.

\subsection{Reducción y commit de estado}

\begin{itemize}
  \item \textbf{Reducción por Dirichlet} heredada de la fase 3 (ADR 0004): mismo operador $\mathbf T$ y $\mathbf g$ constantes; el residuo y el jacobiano se proyectan a DOFs libres antes de cada solve.
  \item \textbf{Commit de variables internas}: al converger el paso, se llama \texttt{assembler.commit\_all\_states()} --- los estados $(\boldsymbol\varepsilon^p, \alpha, \kappa, d, \ldots)$ pasan de \textit{trial} a \textit{committed}. Mismo patrón que \texttt{NonlinearSolver} y \texttt{ArcLengthSolver}. Si el paso no converge, los estados trial se descartan (no se llama commit) y se reporta error.
\end{itemize}

\subsection{Convergencia (ADR 0007)}

Criterio dual fuerza + desplazamiento, calibrado al primer paso temporal con
la escala del ensamblaje inicial:

\begin{itemize}
  \item $\|\mathbf R\|/(\text{atol} + \text{rtol}\,\|\mathbf F\|)\,\leq 1$
  \item $\|\delta \mathbf u\|/(\text{atol} + \text{rtol}\,\|\mathbf u\|)\,\leq 1$
\end{itemize}

Idéntico al \texttt{NonlinearSolver} estático. La calibración se hace una vez al
inicio del análisis; las tolerancias son adimensionales (invariantes bajo
cambio de unidades) gracias a la calibración con \texttt{force\_scale}/\texttt{disp\_scale}.

\subsection{Factorización por iteración}

A diferencia de \texttt{NewmarkSolver} (que factoriza $\mathbf J$ una sola vez al
inicio), aquí $\mathbf K_\text{t}$ cambia con cada iteración, así que $\mathbf J$
debe re-factorizarse. Se mantiene el patrón de \textbf{Newton modificado opcional}
(ADR 0003 fase 2): si \texttt{freeze\_tangent\_after\_iter} es \texttt{int}, se factoriza fresco
las primeras N iteraciones del paso y se reusa la factorización en las
siguientes. \texttt{None} \ensuremath{\Rightarrow} Newton estándar (re-factoriza cada iteración).

\subsection{Recuperación del comportamiento lineal}

Si todos los materiales del modelo son lineales, $\mathbf K_\text{t} \equiv \mathbf K_0$
constante y el residuo se anula en una iteración de Newton (\ensuremath{\Delta}x exacto en el
primer solve). Por tanto **\texttt{NewtonNewmarkSolver} con materiales lineales
reproduce exactamente \texttt{NewmarkSolver}** (sin discretización extra ni error
adicional). Esto se valida en los acceptance.


\section{Contrato de implementación}

\begin{lstlisting}[language=yaml]
name: NewtonNewmarkSolver
kind: solver
status: validated
extends: NewmarkSolver       # variante según Reglas §4

interface:
  type_yaml: NewtonNewmarkSolver
  pipeline_kind: transient    # mismo que NewmarkSolver

parameters_extra_to_NewmarkSolver:
  - { name: convergence,                type: ConvergenceCriterion, required: false, desc: "Política ADR 0007. Default: ConvergenceCriterion() con rtols de solidum.constants." }
  - { name: max_iter,                   type: int,                 required: false, default: 20, desc: "Máximo iteraciones Newton por paso temporal." }
  - { name: freeze_tangent_after_iter,  type: int or None,          required: false, default: null, desc: "Si int, Newton modificado: factoriza fresco las primeras N iter y reusa después (ADR 0003 fase 2)." }

parameters_inherited:
  # Idénticos a NewmarkSolver: t_end, dt, beta, gamma, rayleigh, u0, u0_dot,
  # F_func, linear_algebra, lumping. Ver docs/specs/NewmarkSolver.md.

state_handling:
  commit: "assembler.commit_all_states() tras converger cada paso"
  rollback: "estado trial descartado si el paso no converge → RuntimeError"

acceptance:
  recovery_linear:
    - name: linear_material_matches_NewmarkSolver
      setup: "Mismo modelo y carga que un caso analítico de NewmarkSolver (oscilador 1 GDL, viga), pero con NewtonNewmarkSolver"
      expect: "u_history idéntico (rtol 1e-10) al de NewmarkSolver — convergencia en 1-2 iter por paso"

  nonlinear:
    - name: oscilador_1gdl_elastoplastico_libre
      setup: "1 nodo, Truss2D con Elastoplastic1D, masa concentrada, condición inicial u_0 más allá del yield, vibración libre (F=0)"
      expect: "u(t) refleja decaimiento por disipación plástica; α (deformación plástica acumulada) > 0 al final"

    - name: viga_voladizo_elastoplastica_pulso
      setup: "Frame2DEuler con varios elementos, masa distribuida, pulso de carga concentrada en el extremo libre"
      expect: "respuesta no lineal con redistribución plástica; converge en pocas iter por paso (max ≤6 típicamente)"

  convergence:
    - name: newton_converges_in_few_iter
      setup: "paso plástico activo en un benchmark"
      expect: "número de iteraciones Newton por paso ≤ max_iter; raramente excede 8 con tangente algorítmica consistente"

references:
  - "ADR 0009 §Fase 4 — Transitorio no lineal Newmark + Newton"
  - "Hughes T.J.R. (2000). The Finite Element Method, cap. 7-9 (Newmark, Newton dentro de paso temporal)"
  - "Crisfield M.A. (1991). Non-Linear Finite Element Analysis of Solids and Structures, vol. 2 cap. 24 (dinámica no lineal)"
  - "Spec padre: [docs/specs/NewmarkSolver.md](NewmarkSolver.md)"

\end{lstlisting}


\section{Implementación}

\begin{itemize}
  \item \textbf{Archivo}: solidum/math/solvers/newmark.py --- junto al padre.
  \item \textbf{Clase}: \texttt{NewtonNewmarkSolver(NewmarkSolver)}, registrada con \texttt{@SolverRegistry.register}. Hereda \texttt{\_\_init\_\_} con kwargs adicionales \texttt{convergence}, \texttt{max\_iter}, \texttt{freeze\_tangent\_after\_iter}. Sobrescribe \texttt{solve()}.
  \item \textbf{Flujo de \texttt{solve()}}:
\end{itemize}
\begin{enumerate}
  \item Ensamblar $\mathbf M$, ensamblar sistema no lineal inicial para obtener $\mathbf K_0$ y $\mathbf F_\text{int,0}$.
  \item Calibrar amortiguamiento Rayleigh con $\mathbf K_0$ (constante en el tiempo).
  \item Reducir por Dirichlet (mismo $\mathbf T$, $\mathbf g$, $\mathbf F_\text{dir}$ que la fase 3).
  \item Aceleración inicial consistente: $\mathbf M\,\ddot{\mathbf u}_0 = \mathbf F_\text{ext}(0) - \mathbf F_\text{int}(\mathbf u_0) - \mathbf C\,\dot{\mathbf u}_0 - \mathbf F_\text{dir}$.
  \item Bucle temporal con Newton interno: predictor \ensuremath{\to} iteración $(\mathbf R, \mathbf J = \mathbf M + \gamma\Delta t\,\mathbf C + \beta\Delta t^2\,\mathbf K_\text{t})$, factoriza/solve, actualiza correctores; converger; commit.
  \item Volcar historial; reportar éxito.
\end{enumerate}
\begin{itemize}
  \item \textbf{Refactor del padre}: la fase 3 (\texttt{NewmarkSolver}) sigue intacta. La subclase no extrae métodos del padre por ahora (el flujo no lineal es lo suficientemente distinto como para que la herencia sea estructural, no compartición de pasos internos). Cuando emergió la variante HHT-\ensuremath{\alpha} (2026-05-18, \texttt{NewtonHHTSolver} como subclase de \texttt{NewtonNewmarkSolver}) la herencia estructural funcionó tal cual --- los métodos \texttt{solve()} se sobrescriben íntegramente en cada subclase sin compartir pasos internos. Si en el futuro entran generalized-\ensuremath{\alpha} u otras variantes y se observa que comparten \textit{este} flujo no lineal específico, entonces se refactorizará al patrón con \texttt{\_solve\_step} virtual.
  \item \textbf{Despacho YAML}: \texttt{solver: type: NewtonNewmarkSolver} activado vía \texttt{SolverRegistry} y \texttt{run\_yaml}. Como en \texttt{NewmarkSolver}, \texttt{PIPELINE\_KIND = "transient"} declarativo: si futura regla disparo C reemplaza el \texttt{isinstance} en \texttt{entry.run\_yaml}, no se rompe.
  \item \textbf{Tests}:
  \item tests/test\_newmark\_nonlinear.py nuevo: acceptance.
  \item Cobertura de regresión de los tests de \texttt{NewmarkSolver} permanece intacta (la subclase no toca el padre).
\end{itemize}


\section{Diálogo}

\begin{itemize}
  \item \textbf{2026-05-14} \ensuremath{\cdot} Spec creada como \textbf{variante} de \texttt{NewmarkSolver} aplicando la nueva regla §4 (registrada en el mismo commit que esta spec). Sin ADR nuevo: reusa las decisiones de ADR 0009 fase 4 (residuo, jacobiano, conmutar Newton dentro del paso). Sin entrada propia en el manual estructural --- la subclase es invisible al usuario API (excepto por el \texttt{type} en YAML y los parámetros extra).
  \item \textbf{2026-05-14} \ensuremath{\cdot} Decisión: \textbf{subclase} sobre \texttt{NewmarkSolver} en lugar de bandera \texttt{nonlinear=True} en la misma clase. Razones: (i) más explícito en el código y en el YAML (\texttt{type} distinto para análisis distinto); (ii) evita condicionales que ensucian el flujo lineal de la fase 3; (iii) deja sitio para variantes adicionales (HHT-\ensuremath{\alpha} no lineal, generalized-\ensuremath{\alpha}) por composición. El ADR 0009 fila "Fase 4" decía "\texttt{NewmarkSolver} (mismo, con \texttt{nonlinear=True})" --- interpretación textualista vs. arquitectural: la regla §4 nueva legitima la elección de subclase, y el patrón conserva el espíritu (mismo análisis, mismas decisiones, mismo entrypoint) sin la implementación literal.
  \item \textbf{2026-05-14} \ensuremath{\cdot} Amortiguamiento Rayleigh queda fijado con $\mathbf K_0$ (rigidez elástica de referencia) y constante en el tiempo. Alternativa rechazada: Rayleigh con $\mathbf K_\text{t}$ corriente --- no es la convención estándar y crea coupling artificial entre la disipación viscosa y la plástica. Documentado en sección "Amortiguamiento Rayleigh --- heredado".
\end{itemize}


\newpage

\section{ResponseSpectrumSolver}



\textit{Spec de \textbf{solver nuevo}. Combina las respuestas máximas modales bajo un espectro de respuesta dado (SRSS o CQC). Cierra el ADR 0009 completo y aplica la regla D de refactor arquitectural (free\_vibration sale de \texttt{results.py} al nuevo \texttt{solidum/math/modal\_response.py}).}



\section{Especificación física}

\subsection{0. Descripción general}

Análisis sísmico clásico: el suelo se caracteriza por un \textbf{espectro de respuesta} (en aceleración \texttt{S\_a(T)} o desplazamiento \texttt{S\_d(T)} vs período), no por un acelerograma temporal. Para cada modo del sistema se calcula la respuesta máxima individual \texttt{u\_n\textasciicircum{}max = Γ\_n\ensuremath{\cdot}\ensuremath{\varphi}\_n\ensuremath{\cdot}S\_d(\ensuremath{\omega}\_n)} y se combinan las respuestas modales en una envolvente máxima. La fase temporal entre modos se pierde por construcción --- sólo se conservan amplitudes envolventes.

Es \textbf{lineal} y \textbf{estadístico}: aproxima los máximos absolutos sin reconstruir la historia temporal. Útil para diseño normativo donde el espectro reglamentario es la entrada y los máximos modales son la salida solicitada (e.g., ASCE 7, Eurocódigo 8, NCh 433, NTC-DS México).

Hipótesis: lineales (K, M constantes). Para no-linealidad usar transitorio (Newton-Newmark con acelerograma).

\subsection{1. Ecuación de equilibrio resuelta}

Para cada modo \texttt{n}:

\[
u_n^{\max} = \Gamma_n \cdot \phi_n \cdot S_d(\omega_n)
\]

con $\Gamma_n = \phi_n^T M\, r$ el factor de participación modal en la dirección de excitación \texttt{r}.

\textbf{Combinación SRSS} (Square Root of Sum of Squares):

\[
u^{\max}_k = \sqrt{\sum_n (u_{k,n}^{\max})^2}
\]

\textbf{Combinación CQC} (Complete Quadratic Combination, Der Kiureghian 1980):

\[
u^{\max}_k = \sqrt{\sum_i \sum_j \rho_{ij} \cdot u_{k,i}^{\max} \cdot u_{k,j}^{\max}}
\]

\[
\rho_{ij} = \frac{8\xi^2(1+r) r^{3/2}}{(1-r^2)^2 + 4\xi^2 r (1+r)^2}, \quad r = \omega_i/\omega_j
\]

Cuando $\xi \to 0$ o $r \to 0$: $\rho_{ij} \to \delta_{ij}$ y CQC degenera a SRSS.

\subsection{2. Condiciones iniciales / de contorno}

\begin{itemize}
  \item Apoyos Dirichlet constantes (eliminación directa, ADR 0004).
  \item \textbf{Dirección de excitación rígida} \texttt{r} (vector unitario o por DOF name): vector \texttt{r} con valor 1 en los DOFs traslacionales en la dirección sísmica, 0 en el resto.
  \item \textbf{No} hay condiciones iniciales --- la respuesta es envolvente máxima.
\end{itemize}

\subsection{3. Salidas físicas}

\texttt{ResponseSpectrumResult} (inmutable):

\begin{itemize}
  \item \texttt{u\_combined[ndof]} respuesta máxima envolvente combinada (positiva).
  \item \texttt{u\_per\_mode[ndof, n\_modes]} contribución modal individual con signo.
  \item \texttt{frequencies\_rad[n\_modes]} frecuencias naturales empleadas.
  \item \texttt{participation\_factors[n\_modes]} ($\Gamma_n$).
  \item \texttt{effective\_masses[n\_modes]} ($\Gamma_n^2$). Suma converge a la masa total proyectada en la dirección al incluir suficientes modos.
  \item \texttt{combination} (\texttt{"SRSS"} o \texttt{"CQC"}).
  \item \texttt{damping} (\ensuremath{\xi} usado para CQC; 0.0 si SRSS).
  \item \texttt{direction[ndof]} vector de excitación.
  \item Método \texttt{.cumulative\_effective\_mass\_ratio()} para verificar el porcentaje de masa modal capturado.
\end{itemize}


\section{Formulación numérica}

\subsection{4. Esquema operativo}

\begin{lstlisting}
1. Ensamble interno vía ModalSolver(n_modes, sigma, lumping):
   → modes, frequencies_rad
2. Resolver direction (np.ndarray o {"dof_name": "ux"}).
3. Resolver spectrum (callable o dict tabulado/constante).
4. Calcular Γ_n = φ_nᵀ M r (participation_factors).
5. Calcular S_d(ω_n) para cada modo.
6. u_per_mode[:, n] = Γ_n · φ_n · S_d(ω_n).
7. Combinar: SRSS (raíz cuadrada de suma de cuadrados elemento a elemento)
   o CQC (forma cuadrática con ρ_ij(ω, ξ)).
8. Devolver ResponseSpectrumResult.

\end{lstlisting}

\subsection{5. Predictor / corrector}

No aplica --- un solve por análisis.

\subsection{6. Criterio de convergencia}

No aplica como Newton --- pero la calidad del análisis depende del número de modos incluidos. Se verifica con
\texttt{.cumulative\_effective\_mass\_ratio() \ensuremath{\ge} 0.9} (objetivo normativo típico).

\subsection{7. Imposición de Dirichlet y MPC}

Eliminación directa heredada del \texttt{ModalSolver} interno. MPC fuera de alcance.

\subsection{8. Backend algebraico}

Delega en \texttt{ModalSolver} para el cálculo de modos (\texttt{scipy.sparse.linalg.eigsh} con shift-invert) y luego operaciones de álgebra densa para la combinación modal (los modos son densos por columna).

\subsection{9. Adaptatividad / control de paso}

No aplica. El usuario controla \texttt{n\_modes}. Recomendación: iterativamente subir \texttt{n\_modes} hasta que la masa efectiva acumulada supere el 90\% en la dirección de excitación.

\subsection{10. Caveats numéricos}

\begin{itemize}
  \item \textbf{Cuántos modos}: el método pierde precisión si los modos significativos en la dirección de excitación no están incluidos. Verifica con \texttt{cumulative\_effective\_mass\_ratio}.
  \item \textbf{CQC vs SRSS}: usar CQC cuando los modos están cercanos en frecuencia (cociente < 1.3 según práctica habitual). Para modos bien separados, SRSS es la opción estándar y CQC se degenera a SRSS automáticamente.
  \item \textbf{Espectro fuera de rango}: \texttt{spectrum\_tabulated} aplica clamp a los valores extremos (no extrapola). Si el barrido de frecuencias del modelo cae fuera del rango tabulado, el usuario debe extender la tabla.
  \item \textbf{Excitación multi-direccional}: este solver maneja una sola dirección por análisis. Para combinación direccional (CQC3, regla 100/30/30 de norma) se requieren dos o tres análisis independientes y una combinación adicional fuera del solver.
  \item \textbf{Cargas con fase}: no aplica al método espectral por construcción --- la fase se pierde en el cómputo de máximos modales.
\end{itemize}


\section{Contrato de implementación}

\begin{lstlisting}[language=yaml]
name: ResponseSpectrumSolver
kind: solver
status: validated

interface:
  yaml_type: ResponseSpectrumSolver
  output: ResponseSpectrumResult

parameters:
  - { name: n_modes, type: int, required: true, desc: "Número de modos a calcular y combinar" }
  - { name: direction, type: "ndarray | dict", required: true, desc: "Vector r o {'dof_name': 'ux'}" }
  - { name: spectrum, type: "callable | dict", required: true, desc: "S_d(ω) o {'type': 'tabulated'|'constant_acceleration', ...}" }
  - { name: combination, type: str, required: false, desc: "'SRSS' (default) o 'CQC'" }
  - { name: damping, type: float, required: false, desc: "ξ modal (default 0.05); sólo usado por CQC" }
  - { name: sigma, type: float, required: false, desc: "Shift para ARPACK (default 0)" }
  - { name: lumping, type: str, required: false, desc: "'consistent' (default) o 'lumped'" }

requirements:
  - "Materiales con `density` declarada (ADR 0008)."
  - "Problema lineal — K y M constantes."

conventions:
  units:    "heredadas del modelo (ADR 0008)."
  stability: "No iterativo; estabilidad numérica gobernada por ModalSolver interno."

out_of_scope:
  - "Excitación multi-direccional simultánea (CQC3) — usuario combina afuera."
  - "Espectros con dependencia explícita de ξ por modo (ξ_n distinto) — `damping` es global."
  - "MPC y no linealidad."

acceptance:
  verification:
    - name: single_mode_recovers_gamma_phi_Sd
      setup: "Modelo 2-DOF, n_modes=1, S_d constante=1"
      expect: "u_combined = |u_per_mode[:, 0]| (SRSS trivial sobre un modo)"
      tol_rel: 1.0e-12
    - name: laplaciano_1d_3dof_analitico
      setup: |
        Cadena 4-truss n1-n2-n3-n4-n5 con n1, n5 fijos. E=A=L=1, ρ=1, lumped →
        K Toeplitz tridiagonal [[2,-1,0],[-1,2,-1],[0,-1,2]], M=I. Autovalores
        cerrados del Laplaciano 1D discreto: ω_n² = 2 − 2·cos(nπ/4); modos
        φ_n[i] = √(2/4)·sin(n·i·π/4).
      expect: |
        (1) ω₁ = √(2−√2), ω₂ = √2 a precisión doble.
        (2) Con d = (1, 1/√2, 0): γ₁ = 1, γ₂ = 1/√2, γ₃ = 0 (φ₃ᵀd = 0 por construcción).
        (3) Para Sd ≡ 1 y SRSS: u_combined = (1/√2, 1/√2, 1/√2) exacto.
        (4) Con d = e_{ux₃} (centro): sólo modo 1 contribuye, u_combined = (c/(2√2), c/2, c/(2√2)).
      tol_rel: 1.0e-10
      ref: "tests/test_response_spectrum.py::TestResponseSpectrumAnalytic2DOF"
    - name: srss_manual_combination
      setup: "Modelo 3-DOF, 2 modos, comparar contra raíz cuadrada de suma de cuadrados componente a componente"
      expect: "Coincidencia exacta a precisión de máquina"
      tol_rel: 1.0e-12
    - name: well_separated_modes_cqc_equals_srss
      setup: "Modelo 3-DOF con cociente ω₂/ω₁ > 2"
      expect: "|u_CQC − u_SRSS| / |u_SRSS| < 1%"
      tol_rel: 0.01
  specific:
    - name: close_modes_cqc_differs_from_srss
      setup: "Sistema sintético con ω₂/ω₁ = 1.05 y modos no-ortogonales-en-componentes"
      expect: "Diferencia relativa CQC vs SRSS > 5%"
      tol_rel: 0
    - name: dof_name_builds_unit_vector
      setup: "direction={'dof_name': 'ux'}"
      expect: "direction[node.dofs['ux']] = 1 en todos los nodos con ux"
      tol_rel: 0
    - name: cumulative_effective_mass_monotonic
      setup: "n_modes=2 sobre 3-DOF chain"
      expect: "cumulative crece monotonamente y termina en 1.0"
      tol_rel: 1.0e-10
    - name: free_vibration_wrapper_unchanged
      setup: "Regla D aplicada — ModalResult.free_vibration sigue funcionando"
      expect: "Resultado idéntico a llamada directa a modal_response.free_vibration"
      tol_rel: 1.0e-12

references:
  - "Chopra, A. K. (2017). *Dynamics of Structures*, cap. 13."
  - "Wilson, E. L. (2002). *Three-Dimensional Static and Dynamic Analysis of Structures*, cap. 15."
  - "Der Kiureghian, A. (1980). 'Structural Response to Stationary Excitation'. *J. Eng. Mech. Div. ASCE*, 106(EM6), 1195-1213."

\end{lstlisting}


\section{Implementación}

\begin{itemize}
  \item Archivo solver: solidum/math/solvers/response\_spectrum.py
  \item Clase: \texttt{ResponseSpectrumSolver} (registrada en \texttt{SolverRegistry})
  \item Atributo de despacho: \texttt{PIPELINE\_KIND = "spectrum"} \ensuremath{\to} \texttt{solidum.entry.run\_response\_spectrum} \ensuremath{\to} \texttt{ResponseSpectrumResult}.
  \item Algoritmos centralizados: solidum/math/modal\_response.py (\texttt{response\_spectrum\_srss}, \texttt{response\_spectrum\_cqc}, \texttt{participation\_factors}, helpers \texttt{spectrum\_from\_sa}, \texttt{spectrum\_tabulated}).
  \item Resultado: \texttt{ResponseSpectrumResult} --- dataclass inmutable.
  \item Tests: tests/test\_response\_spectrum.py (19 tests verdes).
\end{itemize}


\section{Diálogo}

\begin{itemize}
  \item \textit{2026-05-18}: regla D aplicada al introducir este solver --- segundo método de cómputo sobre \texttt{ModalResult} (el primero fue \texttt{free\_vibration}). El algoritmo de free\_vibration se movió a \texttt{solidum/math/modal\_response.py}; \texttt{ModalResult.free\_vibration} queda como wrapper delgado de 3 líneas, preservando la API histórica. Los nuevos \texttt{response\_spectrum\_srss} y \texttt{response\_spectrum\_cqc} viven en el mismo módulo. \texttt{ResponseSpectrumSolver} orquesta \texttt{ModalSolver} interno + delegación a \texttt{modal\_response}.
  \item \textit{2026-05-18}: el cuarto valor de \texttt{PIPELINE\_KIND} (\texttt{"spectrum"}) extiende el despacho declarativo de \texttt{run\_yaml} introducido en D2. Cierra el ADR 0009 completo (fases 1-7 + variante HHT-\ensuremath{\alpha}).
\end{itemize}


\newpage

\section{ThetaMethodSolver}



\textit{Spec pre-redactada por la IA (2026-08-25). \textbf{El usuario revisa la física; el plumbing está resuelto.}}

\textit{> Cuarto y último componente de la Etapa 8 --- C1 núcleo mínimo. Resuelve el transitorio de conducción sobre \texttt{Quad4Thermal} y \texttt{Hex8Thermal}.}



\section{Especificación física}

\subsection{0. Descripción general y por qué es un solver nuevo}

Integrador temporal para sistemas semidiscretos de \textbf{primer orden}:

\[
\mathbf C\,\dot{\mathbf T} + \mathbf K\,\mathbf T = \mathbf F(t)
\]

\textbf{No es una variante de \texttt{NewmarkSolver}.} La familia Newmark integra la ecuación de \textbf{segundo orden} $\mathbf M\ddot{\mathbf u} + \mathbf C\dot{\mathbf u} + \mathbf K\mathbf u = \mathbf F$ mediante hipótesis sobre la aceleración dentro del paso. En conducción no existe segunda derivada temporal: no hay aceleración, no hay inercia, y las hipótesis de Newmark no tienen sobre qué aplicarse. Forzar Newmark con $\mathbf M = \mathbf 0$ degenera el esquema.

Es, por tanto, un solver de familia propia, no una subclase --- y no le corresponde spec corta de extensión sino spec completa (\texttt{Reglas.md §4}).

El \textbf{régimen estacionario} no lo usa: $\mathbf K\mathbf T = \mathbf F$ lo resuelve el \texttt{LinearSolver} existente sin modificación.

\subsection{1. Esquema \ensuremath{\theta}}

Se aproxima la ecuación en un punto intermedio del paso, ponderando el estado con el parámetro $\theta \in [0,1]$:

\[
\mathbf C\,\frac{\mathbf T_{n+1} - \mathbf T_n}{\Delta t} + \mathbf K\left[\theta\,\mathbf T_{n+1} + (1-\theta)\,\mathbf T_n\right] = \theta\,\mathbf F_{n+1} + (1-\theta)\,\mathbf F_n
\]

Reordenando, el sistema a resolver por paso:

\[
\underbrace{\left(\mathbf C + \theta\,\Delta t\,\mathbf K\right)}_{\mathbf A}\,\mathbf T_{n+1} = \left[\mathbf C - (1-\theta)\,\Delta t\,\mathbf K\right]\mathbf T_n + \Delta t\left[\theta\,\mathbf F_{n+1} + (1-\theta)\,\mathbf F_n\right]
\]

$\mathbf A$ es \textbf{simétrica y definida positiva} para $\theta > 0$ (suma de $\mathbf C$ definida positiva y $\mathbf K$ semidefinida, con Dirichlet impuesto), luego el despachador algebraico (ADR 0003) elige Cholesky. Con $\Delta t$ constante, $\mathbf A$ \textbf{se factoriza una sola vez} y se reutiliza en todos los pasos.

\subsection{2. Casos particulares de \ensuremath{\theta}}

\begin{center}
\begin{tabularx}{\textwidth}{YYYY}
\hline
\textbf{$\theta$} & \textbf{Nombre} & \textbf{Orden} & \textbf{Estabilidad} \\
\hline
$0$ & Euler explícito (forward) & 1 & \textbf{Condicional}: $\Delta t \le 2/\lambda_{\max}$ \\
$1/2$ & Crank-Nicolson & \textbf{2} & Incondicional (A-estable) \\
$2/3$ & Galerkin & 1 & Incondicional \\
$1$ & Euler implícito (backward) & 1 & Incondicional, \textbf{L-estable} \\
\hline
\end{tabularx}
\end{center}
Con $\theta \ge 1/2$ el esquema es \textbf{incondicionalmente estable}: ningún $\Delta t$ produce divergencia. Es la propiedad que distingue este solver del \texttt{CentralDifferenceSolver} mecánico, cuya estabilidad es siempre condicional (CFL).

\subsection{3. Estabilidad no es lo mismo que ausencia de oscilaciones}

Punto físico central de esta spec, y la razón de la elección de default.

Crank-Nicolson ($\theta = 1/2$) es de \textbf{segundo orden} y por tanto el más preciso de la tabla, pero \textbf{A-estable y no L-estable}: su factor de amplificación tiende a $-1$ para los modos de frecuencia alta, en lugar de a $0$. Consecuencia práctica: ante un \textbf{cambio brusco} ---un escalón de temperatura impuesto en la frontera, típico en el arranque de un análisis--- produce \textbf{oscilaciones amortiguadas lentamente} en los primeros pasos, con temperaturas que pueden salirse del rango de los datos.

Euler implícito ($\theta = 1$) es sólo de primer orden, pero \textbf{L-estable}: amortigua los modos altos por completo en un paso. Nunca oscila.

Esto es el mismo fenómeno físico que motivó el default \texttt{lumped} en la matriz de capacidad, atacado desde el otro lado: \textbf{el principio del máximo} de la ecuación de difusión dice que la solución exacta nunca excede los extremos de los datos iniciales y de frontera. Un esquema que lo viola no es simplemente impreciso --- produce un resultado cualitativamente imposible.

\textbf{Por eso el default propuesto es $\theta = 1$} (Euler implícito), robusto ante el arranque típico. Crank-Nicolson queda disponible para quien busque precisión de segundo orden y sepa que su problema es suave. Se documenta el criterio para que la elección sea informada y no una preferencia oculta.

\subsection{4. Condición inicial}

El campo $\mathbf T_0$ debe declararse. Dos formas admitidas:

\begin{itemize}
  \item \textbf{Uniforme}: un escalar, el caso más frecuente (cuerpo inicialmente a temperatura ambiente).
  \item \textbf{Por nodo}: vector completo, para arrancar de un estado no uniforme (p. ej. el resultado de un análisis estacionario previo).
\end{itemize}

A diferencia del mecánico, \textbf{no hay velocidad inicial}: la ecuación es de primer orden y $\mathbf T_0$ la determina por completo.

\textbf{Consistencia con Dirichlet}: si $\mathbf T_0$ contradice una temperatura impuesta en un nodo con Dirichlet, existe una discontinuidad en $t=0$. Es físicamente legítimo (un choque térmico es exactamente eso) pero es la situación que más excita las oscilaciones descritas en §3. El solver \textbf{avisa} cuando lo detecta, sin abortar: es modelización válida, no error.

\subsection{5. Selección del paso de tiempo}

El paso no se elige por estabilidad ---con $\theta \ge 1/2$ no hay límite--- sino por \textbf{precisión}. La escala física la da la difusividad $\alpha = k/(\rho c)$ [m\ensuremath{^{2}}/s] del material y el tamaño característico de elemento $h$:

\[
\Delta t_{\text{car}} \sim \frac{h^2}{\alpha}
\]

Es el tiempo que tarda el frente térmico en atravesar un elemento. Un $\Delta t \gg \Delta t_{\text{car}}$ es estable pero no resuelve el transitorio: salta por encima de la física que se quiere ver.

El solver \textbf{no calcula ni impone} este valor ---igual que el \texttt{CentralDifferenceSolver} no calcula el $\Delta t$ crítico---, pero sí lo \textbf{reporta como diagnóstico} al arrancar, para que el usuario contraste su elección. Es información, no restricción.

\subsection{6. Sin amortiguamiento de Rayleigh}

El amortiguamiento de Rayleigh $\mathbf C = \alpha\mathbf M + \beta\mathbf K$ del subsistema dinámico \textbf{no aplica}: es un modelo de disipación de la ecuación de segundo orden. En conducción la "disipación" es la propia conducción, ya contenida en $\mathbf K$. El parámetro no existe en este solver.


\section{Formulación numérica}

\subsection{7. Algoritmo}

\begin{lstlisting}
A = C + θ·Δt·K                     # una vez, si Δt constante
factor = factorize(A)              # Cholesky (A es SPD); reutilizada
B = C − (1−θ)·Δt·K                 # una vez

T = T_0
para cada paso n:
    rhs = B·T_n + Δt·[θ·F_{n+1} + (1−θ)·F_n]
    aplicar Dirichlet (eliminación, ADR 0004)
    T_{n+1} = factor.solve(rhs_reducido)
    almacenar según output_every

\end{lstlisting}

\textbf{Coste por paso}: una sustitución hacia adelante/atrás. La factorización es única mientras $\Delta t$ y las matrices no cambien --- el problema es lineal y sin historia, así que $\mathbf K$ y $\mathbf C$ se ensamblan una sola vez.

\subsection{8. Dirichlet dependiente del tiempo}

Una temperatura impuesta puede variar en el tiempo, $\bar T(t)$ --- es un caso de uso central (una superficie sometida a un ciclo térmico). La eliminación del ADR 0004 se aplica en cada paso con el valor correspondiente; el término de acoplamiento $\mathbf K_{fp}\bar{\mathbf T}_p$ se recalcula, la factorización de $\mathbf A_{ff}$ \textbf{no}.

\subsection{9. Resultado}

Devuelve un \texttt{ThermalTransientResult} con \texttt{t\_history}, \texttt{T\_history}, \texttt{n\_steps} y \texttt{converged}.

\textbf{Decisión de zona gris que se señala}: no se reutiliza \texttt{TransientResult}, cuyos campos \texttt{udot\_history}, \texttt{uddot\_history}, \texttt{alpha\_rayleigh} y \texttt{beta\_rayleigh} no tienen significado térmico. Rellenarlos con \texttt{None} o ceros produciría un objeto que miente sobre su contenido. Se sigue el criterio ya aplicado en el proyecto con \texttt{ModalResult}, \texttt{HarmonicResult} y \texttt{ResponseSpectrumResult}: un tipo de análisis con semántica propia tiene su dataclass propia.

\subsection{10. Caveats numéricos}

\begin{itemize}
  \item \textbf{Sin Dirichlet, $\mathbf A$ es singular} en el estacionario y mal condicionada en el transitorio con $\Delta t$ grande. El diagnóstico debe nombrar el modo de temperatura uniforme, no reportar un fallo algebraico genérico.
  \item \textbf{Oscilaciones con Crank-Nicolson ante escalón}: descrito en §3. No es un error del solver; se documenta y se cuantifica con test.
  \item \textbf{$\theta = 0$ (explícito)}: se admite por completitud pero requiere $\Delta t \le 2/\lambda_{\max}$, y con capacidad consistente $\lambda_{\max}$ es grande. Sin masa lumped es prácticamente inutilizable; se avisa.
  \item \textbf{Contraste alto de conductividad}: mal condicionamiento heredado del problema, no del esquema.
\end{itemize}


\section{Contrato de implementación}

\begin{lstlisting}[language=yaml]
name: ThetaMethodSolver
kind: solver
status: draft            # draft → implemented → validated

interface:
  pipeline_kind: thermal_transient
  field: temperature
  equation_order: 1      # primer orden en el tiempo (Newmark integra segundo)

parameters:
  - { name: theta, type: float, required: false, default: 1.0,
      desc: "Peso del esquema ∈ [0,1]. 1.0 Euler implícito (L-estable, default);
             0.5 Crank-Nicolson (2º orden, A-estable, puede oscilar ante escalón);
             2/3 Galerkin; 0.0 explícito (condicionalmente estable)" }
  - { name: dt, type: float, required: true,
      desc: "Paso de tiempo [s]. No limitado por estabilidad si θ ≥ 0.5, sí por
             precisión: contrastar con el Δt característico h²/α reportado" }
  - { name: n_steps, type: int, required: true, desc: "Número de pasos" }
  - { name: T_initial, type: float | array-like, required: true,
      desc: "Condición inicial. Escalar ⇒ campo uniforme; vector ⇒ valor por nodo" }
  - { name: lumping, type: str, required: false, default: "lumped",
      desc: "Forma de la matriz de capacidad. Default lumped por el mismo criterio
             físico que en los elementos térmicos (Quad4Thermal §5)" }
  - { name: output_every, type: int, required: false, default: 1,
      desc: "Almacenar el campo cada N pasos; limita la memoria en corridas largas" }
  - { name: linear_algebra, type: str, required: false, default: "auto",
      desc: "Backend algebraico (ADR 0003). A es SPD ⇒ Cholesky por defecto" }
  # Sin `rayleigh`: el amortiguamiento de Rayleigh es de la ecuación de 2º orden.

conventions:
  stability: "θ ≥ 0.5 incondicionalmente estable; θ = 1 además L-estable
              (amortigua modos altos por completo, nunca oscila)"
  dt_guidance: "Δt característico ~ h²/α con α = k/(ρc). Reportado como
                diagnóstico al arrancar; no impuesto"
  order_reporting: "el solver reporta el orden efectivo del esquema según θ
                    (2 si θ=0.5, 1 en el resto), para que el coste en precisión
                    del default robusto sea visible y no silencioso"
  factorization: "A = C + θ·Δt·K se factoriza UNA vez si Δt es constante"

validity:
  - "sistema lineal de primer orden: C·Ṫ + K·T = F"
  - "θ ∈ [0, 1]; dt > 0; n_steps ≥ 1"
  - "el modelo debe imponer al menos un Dirichlet de temperatura"
  - "material con c y density declaradas (ThermalConduction §4)"

out_of_scope:
  - "no linealidad: k(T), cambio de fase, radiación ⇒ requerirían Newton por paso"
  - "paso de tiempo adaptativo ⇒ diferido hasta caso de uso real"
  - "amortiguamiento de Rayleigh (no aplica a 1er orden)"
  - "acoplamiento termomecánico"

acceptance:
  verification:
    - name: decaimiento_exponencial_1dof
      setup: "sistema de 1 DOF c·Ṫ + k·T = 0 con T(0) = T_0; comparar contra
             la solución analítica T(t) = T_0·exp(−k·t/c)"
      expect: "error < tol para θ=1 y θ=0.5; el error de Crank-Nicolson decrece
              como O(Δt²) y el de Euler implícito como O(Δt) — verificar AMBAS
              tasas por refinamiento sucesivo de Δt"
      tol_rel: 1.0e-3
    - name: orden_de_convergencia_temporal
      setup: "mismo problema, Δt sucesivamente halvado"
      expect: "pendiente log-log ≈ 1 para θ=1 y ≈ 2 para θ=0.5.
               Es el test que distingue realmente los dos esquemas"
      tol_abs: 0.15
    - name: estabilidad_incondicional
      setup: "θ = 1 con Δt tres órdenes de magnitud mayor que h²/α"
      expect: "la solución NO diverge; converge monótonamente al estacionario.
               Con θ = 0 y el mismo Δt sí diverge (contraste explícito)"
    - name: convergencia_al_estacionario
      setup: "problema de pared plana con Dirichlet fijo en ambos extremos;
             integrar hasta t → ∞"
      expect: "el campo converge al perfil lineal que da el LinearSolver sobre
              el MISMO modelo; comparación nodo a nodo"
      tol_rel: 1.0e-8
  specific:
    - name: L_estabilidad_ante_escalon
      setup: "cuerpo a T_0 uniforme; escalón de Dirichlet en una cara;
             Δt grande; comparar θ=1 contra θ=0.5"
      expect: "θ=1 monótono, sin temperaturas fuera del rango [T_0, T_pared];
               θ=0.5 exhibe oscilación amortiguada en los primeros pasos.
               El test DOCUMENTA la diferencia, no la corrige"
    - name: principio_del_maximo
      setup: "θ=1, capacidad lumped, sin fuente; condiciones de frontera
             acotadas entre T_min y T_max"
      expect: "en TODO paso y TODO nodo, T ∈ [T_min, T_max]. Ningún valor
               fuera del rango de los datos"
      tol_abs: 1.0e-12
    - name: conduccion_transitoria_semi_infinita
      setup: "barra larga inicialmente a T_0, con la cara x=0 llevada
             súbitamente a T_s; comparar con la solución analítica de
             Carslaw-Jaeger T(x,t) = T_s + (T_0−T_s)·erf(x/(2√(αt)))"
      expect: "perfil coincidente en varios instantes mientras el frente no
              alcanza el extremo opuesto (hipótesis de medio semi-infinito)"
      tol_rel: 2.0e-2
    - name: balance_energetico_transitorio
      setup: "dominio aislado (todo Neumann homogéneo) con fuente Q constante"
      expect: "la energía acumulada ΣρcVΔT iguala a Q·V·t en cada paso —
              el esquema conserva la energía introducida"
      tol_rel: 1.0e-10
    - name: dirichlet_variable_en_el_tiempo
      setup: "temperatura impuesta con variación sinusoidal en el tiempo"
      expect: "el campo sigue la excitación; la factorización de A_ff se reutiliza
               (verificado por contador de llamadas) mientras el acoplamiento
               K_fp·T_p se recalcula por paso"
      tol_rel: 1.0e-10
    - name: condicion_inicial_escalar_y_vectorial
      setup: "arrancar con T_initial escalar y con vector equivalente"
      expect: "resultados idénticos; el escalar se expande correctamente"
      tol_abs: 1.0e-14
    - name: aviso_incompatibilidad_inicial_dirichlet
      setup: "T_initial contradice un Dirichlet impuesto en un nodo"
      expect: "el solver AVISA identificando los nodos, y continúa: un choque
               térmico es modelización legítima, no un error"
    - name: reporte_del_orden_efectivo
      setup: "construir el solver con θ = 1, 0.5 y 2/3"
      expect: "reporta orden 1, 2 y 1 respectivamente. Hace visible el coste en
               precisión del default robusto en vez de dejarlo implícito"
    - name: rechazo_inputs_invalidos
      setup: "θ fuera de [0,1]; dt ≤ 0; n_steps < 1; material sin c ni density"
      expect: "ValueError con mensaje accionable en cada caso"

references:
  - "Zienkiewicz O.C., Taylor R.L. (2000). The Finite Element Method, Vol. 1. Butterworth-Heinemann. §18 (esquemas de un paso para 1er orden; análisis de estabilidad)."
  - "Lewis R.W., Nithiarasu P., Seetharamu K.N. (2004). Fundamentals of the FEM for Heat and Fluid Flow. Wiley. §6 (θ-method, oscilaciones, criterios de Δt)."
  - "Hughes T.J.R. (2000). The Finite Element Method: Linear Static and Dynamic FEA. Dover. §8 (algoritmos para problemas parabólicos; A-estabilidad vs L-estabilidad)."
  - "Carslaw H.S., Jaeger J.C. (1959). Conduction of Heat in Solids. Oxford. §2.4 (sólido semi-infinito, solución con función error)."
  - "Wood W.L. (1990). Practical Time-Stepping Schemes. Oxford."

\end{lstlisting}


\section{Implementación}

\textit{Rellena la IA tras programar.}

\begin{itemize}
  \item Archivo: ---
  \item Clase: ---
  \item Tests:
  \item ---
  \item Notas de traducción: ---
\end{itemize}


\section{Diálogo}

\textit{Puntos abiertos para el usuario. Sólo física y alcance.}

\textbf{2026-08-25 --- IA \ensuremath{\to} usuario. Un punto:}

\begin{enumerate}
  \item \textasciitilde{}\textasciitilde{}\textbf{Valor por defecto de $\theta$.}\textasciitilde{}\textasciitilde{} ✅ \textbf{Resuelto 2026-08-25 --- $\theta = 1$ (Euler implícito)}. El usuario delegó la decisión en la IA por ser materia de criterio técnico; queda razonada aquí para que sea auditable.
\end{enumerate}

   \textbf{La disyuntiva.} Crank-Nicolson ($\theta = 1/2$) es de orden 2 y más preciso para un mismo $\Delta t$; Euler implícito ($\theta = 1$) es de orden 1 pero \textbf{L-estable}. Ante el arranque típico ---un escalón de temperatura en la frontera--- Crank-Nicolson produce oscilaciones que pueden dar \textbf{temperaturas fuera del rango de los datos}, violando el principio del máximo.

   \textbf{El criterio que decide} no es el orden de convergencia sino qué debe hacer el programa cuando el usuario no elige: un \textbf{default protege a quien no eligió}. Los dos modos de fallo no son simétricos --- un resultado impreciso se detecta refinando el paso y se corrige; un resultado \textit{imposible} desconcierta, y quien no conozca la teoría de A-estabilidad frente a L-estabilidad no tiene forma de diagnosticarlo. La precisión perdida es recuperable por el usuario; la confianza en un resultado cualitativamente absurdo, no.

   \textbf{Coherencia interna}: es el mismo criterio que fijó \texttt{lumped} como default de la matriz de capacidad, aplicado ahora al eje temporal en vez del espacial. Ambos atacan el mismo fenómeno ---oscilación espuria ante frente abrupto--- y que apunten en la misma dirección hace el sistema predecible.

   \textbf{El default no encierra a nadie}: Crank-Nicolson queda a un parámetro de distancia para quien sepa que su problema es suave y busque orden 2.

   \textbf{Requisito añadido con la decisión}: el solver \textbf{reporta el orden efectivo del esquema} según el $\theta$ en uso, para que el coste en precisión del default sea visible y no una penalización silenciosa. Es la contrapartida honesta de elegir robustez.


\newpage

\chapter{Elementos — catálogo}



\textit{Referencia rápida de los elementos implementados. Una entrada por elemento. Para detalles físicos/numéricos \ensuremath{\to} código fuente.}

\textit{> \textbf{Convenciones}: \texttt{STRAIN\_DIM} = dimensión Voigt esperada del material asociado (1 = axial escalar, 3 = 2D \texttt{[\ensuremath{\varepsilon}\_xx, \ensuremath{\varepsilon}\_yy, \ensuremath{\gamma}\_xy]}, 6 = 3D). DOFs por nodo = \texttt{DOF\_NAMES}.}



\section{Truss2D --- barra axial 2D}

Elemento sólido 1D de primer orden, inmerso en el plano. Dos nodos; transmite exclusivamente fuerza axial.

\subsection{Cinemática}
\begin{itemize}
  \item Interpolación lineal del desplazamiento entre los dos nodos.
  \item Deformación axial uniforme en el elemento.
\end{itemize}

\subsection{Formulación}
Cosenos directores del eje: \texttt{c = (x\ensuremath{_{2}}−x\ensuremath{_{1}})/L}, \texttt{s = (y\ensuremath{_{2}}−y\ensuremath{_{1}})/L}.

\begin{lstlisting}
ε     = B · u_e,          B = (1/L) · [−c, −s, c, s]
K_e   = (E·A / L) · dᵀd,  d = [−c, −s, c, s]
F_int = σ·A · d

\end{lstlisting}

\subsection{Convenciones}
\begin{itemize}
  \item Convención de signos: \texttt{\ensuremath{\sigma} > 0 \ensuremath{\Leftrightarrow} \ensuremath{\varepsilon} > 0 (elongación)}.
  \item \texttt{STRAIN\_DIM = 1}: la deformación que recibe el material es un escalar (no Voigt).
  \item La orientación de los nodos no afecta el resultado (\texttt{K\_e} y \texttt{F\_int} son invariantes bajo su permutación).
\end{itemize}

\subsection{Parámetros}
\begin{itemize}
  \item \texttt{A} --- área de la sección transversal.
\end{itemize}

\subsection{Cargas distribuidas}
\begin{itemize}
  \item \texttt{compute\_body\_load(b)} reparte \texttt{b \ensuremath{\cdot} A \ensuremath{\cdot} L\ensuremath{_{0}} / 2} por nodo en cada componente global (forma cerrada exacta para \texttt{b} uniforme). Útil para peso propio: \texttt{b = (0, -\ensuremath{\rho}\ensuremath{\cdot}g)}.
\end{itemize}

\subsection{Régimen de validez}
\begin{itemize}
  \item Pequeñas deformaciones (\texttt{|\ensuremath{\varepsilon}| \ensuremath{\lesssim} 10\ensuremath{^{-}}\ensuremath{^{2}}}) y pequeños desplazamientos.
  \item Cargas exclusivamente axiales. Si la carga transversal sobre la barra no es despreciable \ensuremath{\to} \texttt{Frame2DEuler} / \texttt{Frame2DTimoshenko}.
\end{itemize}

\subsection{Validación}
\begin{itemize}
  \item Tests: tests/test\_truss.py \ensuremath{\cdot} \texttt{TestTruss2D}.
  \item Archivo fuente: solidum/elements/truss.py \ensuremath{\cdot} clase \texttt{Truss2D}.
  \item Spec: docs/specs/Truss2D.md.
  \item Referencia: Bathe, \textit{Finite Element Procedures}, §4.2.1.
\end{itemize}


\section{Truss2DCorot --- armadura 2D corotacional}

Barra axial 2D en régimen de \textbf{grandes desplazamientos y rotaciones con pequeña deformación} (Updated Lagrangian). Hereda de \texttt{Truss2D}; comparte DOFs, parámetros y contrato con el material.

\subsection{Cinemática}
\begin{itemize}
  \item Longitud y cosenos directores se recalculan en configuración \textbf{corriente} en cada evaluación.
  \item Deformación ingenieril corotacional: \texttt{\ensuremath{\varepsilon} = (l − L\ensuremath{_{0}})/L\ensuremath{_{0}}}.
\end{itemize}

\subsection{Formulación}
\begin{lstlisting}
d = [−c_θ, −s_θ, c_θ, s_θ]      (dirección corriente del eje)
n = [−s_θ,  c_θ, s_θ, −c_θ]     (perpendicular al eje)
K_M   = (E·A / L₀) · d·dᵀ
K_G   = (N / l) · n·nᵀ           (rigidez geométrica)
K_T   = K_M + K_G
F_int = N · d

\end{lstlisting}

\subsection{Cargas distribuidas}
\begin{itemize}
  \item Heredado de \texttt{Truss2D}: \texttt{compute\_body\_load(b)} reparte \texttt{b \ensuremath{\cdot} A \ensuremath{\cdot} L\ensuremath{_{0}} / 2} por nodo, evaluado en geometría de referencia (aproximación estándar para cargas conservadoras en formulaciones corotacionales).
\end{itemize}

\subsection{Régimen de validez}
\begin{itemize}
  \item \texttt{|\ensuremath{\varepsilon}| \ensuremath{\lesssim} 10\ensuremath{^{-}}\ensuremath{^{2}}} (pequeña deformación axial).
  \item Desplazamientos y rotaciones de cualquier magnitud.
  \item Cargas exclusivamente axiales.
  \item \textbf{Fuera de alcance}: pandeo por bifurcación (se capta ablandamiento precrítico pero no el punto crítico; para snap-through usar arc-length).
\end{itemize}

\subsection{Validación}
\begin{itemize}
  \item Tests: tests/test\_truss.py \ensuremath{\cdot} \texttt{TestTruss2DCorot} (3 tests, uno por criterio de aceptación).
  \item Archivo fuente: solidum/elements/truss.py \ensuremath{\cdot} clase \texttt{Truss2DCorot}.
  \item Spec: docs/specs/Truss2DCorot.md.
  \item Referencias: Crisfield §3.3, Belytschko §4.5.
\end{itemize}


\section{Truss3D --- barra axial 3D}

Elemento sólido 1D de primer orden, inmerso en el espacio tridimensional. Dos nodos articulados; transmite exclusivamente fuerza axial. Régimen estrictamente de \textbf{linealidad geométrica}.

\subsection{Cinemática}
\begin{itemize}
  \item Interpolación lineal del desplazamiento entre los dos nodos.
  \item Deformación axial uniforme en el elemento.
  \item Configuración inicial fija: \texttt{L\ensuremath{_{0}}}, \texttt{cₓ}, \texttt{c\_y}, \texttt{c\_z} no se actualizan con los desplazamientos.
\end{itemize}

\subsection{Formulación}
Cosenos directores del eje: \texttt{cₓ = (x\ensuremath{_{2}}−x\ensuremath{_{1}})/L}, \texttt{c\_y = (y\ensuremath{_{2}}−y\ensuremath{_{1}})/L}, \texttt{c\_z = (z\ensuremath{_{2}}−z\ensuremath{_{1}})/L}.

\begin{lstlisting}
ε     = B · u_e,          B = (1/L) · [−cₓ, −c_y, −c_z, cₓ, c_y, c_z]
K_e   = (E·A / L) · d·dᵀ,  d = [−cₓ, −c_y, −c_z, cₓ, c_y, c_z]    (6×6, rango 1)
F_int = σ·A · d

\end{lstlisting}

\subsection{Convenciones}
\begin{itemize}
  \item Convención de signos: \texttt{\ensuremath{\sigma} > 0 \ensuremath{\Leftrightarrow} \ensuremath{\varepsilon} > 0 (elongación)}.
  \item \texttt{STRAIN\_DIM = 1}: la deformación que recibe el material es un escalar (no Voigt).
  \item La orientación de los nodos no afecta el resultado (\texttt{K\_e} y \texttt{F\_int} son invariantes bajo su permutación).
  \item Acepta nodos con 2 ó 3 coordenadas (completa con \texttt{z = 0} si la tercera falta).
\end{itemize}

\subsection{Parámetros}
\begin{itemize}
  \item \texttt{A} --- área de la sección transversal.
\end{itemize}

\subsection{Cargas distribuidas}
\begin{itemize}
  \item \texttt{compute\_body\_load(b)} reparte \texttt{b \ensuremath{\cdot} A \ensuremath{\cdot} L\ensuremath{_{0}} / 2} por nodo en cada componente global (forma cerrada exacta). Útil para peso propio: \texttt{b = (0, 0, -\ensuremath{\rho}\ensuremath{\cdot}g)}.
\end{itemize}

\subsection{Régimen de validez}
\begin{itemize}
  \item Pequeñas deformaciones (\texttt{|\ensuremath{\varepsilon}| \ensuremath{\lesssim} 10\ensuremath{^{-}}\ensuremath{^{2}}}), pequeños desplazamientos y pequeñas rotaciones.
  \item Cargas exclusivamente axiales.
  \item Para grandes rotaciones en 3D \ensuremath{\to} \texttt{Truss3DCorot}.
\end{itemize}

\subsection{Validación}
\begin{itemize}
  \item Tests: tests/test\_truss.py \ensuremath{\cdot} \texttt{TestTruss3D}.
  \item Archivo fuente: solidum/elements/truss.py \ensuremath{\cdot} clase \texttt{Truss3D}.
  \item Spec: docs/specs/Truss3D.md.
  \item Referencias: Bathe §4.2.1; Cook §2.3.
\end{itemize}


\section{Truss3DCorot --- barra axial 3D corotacional}

Barra axial 3D en régimen de \textbf{grandes desplazamientos y rotaciones con pequeña deformación} (Updated Lagrangian). Hereda de \texttt{Truss3D}; comparte DOFs, parámetros y contrato con el material.

\subsection{Cinemática}
\begin{itemize}
  \item Longitud y cosenos directores se recalculan en configuración \textbf{corriente} en cada evaluación.
  \item Deformación ingenieril corotacional: \texttt{\ensuremath{\varepsilon} = (l − L\ensuremath{_{0}})/L\ensuremath{_{0}}}.
  \item La dirección perpendicular al eje es un \textbf{plano} (dimensión 2), no un vector único.
\end{itemize}

\subsection{Formulación}
\begin{lstlisting}
d = [−cₓ, −c_y, −c_z, cₓ, c_y, c_z]       (dirección corriente del eje)
ê = (cₓ, c_y, c_z),  P = I₃ − ê·êᵀ        (proyector 3×3 al plano perpendicular)

K_M   = (E·A / L₀) · d·dᵀ                  (6×6, rango 1)
K_G   = (N / l) · [[P, −P], [−P, P]]       (6×6, rango 2)
K_T   = K_M + K_G
F_int = N · d

\end{lstlisting}

\subsection{Cargas distribuidas}
\begin{itemize}
  \item Heredado de \texttt{Truss3D}: \texttt{compute\_body\_load(b)} reparte \texttt{b \ensuremath{\cdot} A \ensuremath{\cdot} L\ensuremath{_{0}} / 2} por nodo, evaluado en geometría de referencia.
\end{itemize}

\subsection{Régimen de validez}
\begin{itemize}
  \item \texttt{|\ensuremath{\varepsilon}| \ensuremath{\lesssim} 10\ensuremath{^{-}}\ensuremath{^{2}}} (pequeña deformación axial).
  \item Desplazamientos y rotaciones de cualquier magnitud en el espacio.
  \item Cargas exclusivamente axiales.
\end{itemize}

\subsection{Validación}
\begin{itemize}
  \item Tests: tests/test\_truss.py \ensuremath{\cdot} \texttt{TestTruss3DCorot} (3 tests, uno por criterio de aceptación; el de rigidez geométrica verifica dos direcciones transversas independientes del plano perpendicular).
  \item Archivo fuente: solidum/elements/truss.py \ensuremath{\cdot} clase \texttt{Truss3DCorot}.
  \item Spec: docs/specs/Truss3DCorot.md.
  \item Referencias: Crisfield §3.3; Belytschko §4.5.
\end{itemize}


\section{Cable2DCorot --- cable 2D corotacional}

Elemento 1D inmerso en el plano que modela un cable: dos nodos, transmite solo tensión. Cinemática corotacional; la unilateralidad la aporta el material (típicamente \texttt{CableMaterial1D}).

\subsection{Cinemática}
\begin{itemize}
  \item Idéntica a una barra corotacional: longitud y cosenos directores se recalculan en configuración corriente en cada evaluación.
  \item Deformación ingenieril corotacional: \texttt{\ensuremath{\varepsilon} = (l − L\ensuremath{_{0}})/L\ensuremath{_{0}}}.
  \item La respuesta física de cable aparece \textbf{solo} si el material devuelve \texttt{\ensuremath{\sigma} = 0} en compresión.
\end{itemize}

\subsection{Formulación}
\begin{lstlisting}
d = [−c_θ, −s_θ, c_θ, s_θ]
n = [−s_θ,  c_θ, s_θ, −c_θ]

K_M   = (E_t · A / L₀) · d·dᵀ        (0 si material destensado)
K_G   = (N / l) · n·nᵀ                (0 si N = 0)
K_T   = K_M + K_G
F_int = N · d                         (0 si N = 0)

\end{lstlisting}

\subsection{Cargas distribuidas}
\begin{itemize}
  \item \texttt{compute\_body\_load(b)} reparte \texttt{b \ensuremath{\cdot} A \ensuremath{\cdot} L\ensuremath{_{0}} / 2} por nodo en cada componente global, evaluado en geometría de referencia. Útil para peso propio.
\end{itemize}

\subsection{Régimen de validez}
\begin{itemize}
  \item \texttt{|\ensuremath{\varepsilon}| \ensuremath{\lesssim} 10\ensuremath{^{-}}\ensuremath{^{2}}} en régimen tensado.
  \item Grandes desplazamientos y rotaciones en el plano.
  \item Cables completamente destensados aportan \texttt{K\_T = 0} al sistema global: otros elementos deben garantizar la estabilidad numérica.
\end{itemize}

\subsection{Independencia del diseño}
El elemento \textbf{no hereda} de \texttt{Truss2DCorot}. La cinemática corotacional se implementa dentro de la clase para que futuras modificaciones de las armaduras no afecten a los cables, y viceversa.

\subsection{Validación}
\begin{itemize}
  \item Tests: tests/test\_cable\_elements.py \ensuremath{\cdot} \texttt{TestCable2DCorot} (4 tests de aceptación: tensado, destensado, rotación rígida, cruce por cero).
  \item Archivo fuente: solidum/elements/cable.py \ensuremath{\cdot} clase \texttt{Cable2DCorot}.
  \item Spec: docs/specs/Cable2DCorot.md.
  \item Referencias: Crisfield §3.3; Irvine §1.2.
\end{itemize}


\section{Cable3DCorot --- cable 3D corotacional}

Elemento 1D inmerso en el espacio que modela un cable: dos nodos, transmite solo tensión, puede rotar libremente en el espacio. Cinemática corotacional 3D; unilateralidad aportada por el material (típicamente \texttt{CableMaterial1D}).

\subsection{Cinemática}
\begin{itemize}
  \item Idéntica en filosofía a una barra corotacional 3D: longitud y cosenos directores se recalculan en configuración corriente.
  \item Deformación ingenieril corotacional: \texttt{\ensuremath{\varepsilon} = (l − L\ensuremath{_{0}})/L\ensuremath{_{0}}}.
  \item La dirección perpendicular al eje es un \textbf{plano} (dimensión 2).
\end{itemize}

\subsection{Formulación}
\begin{lstlisting}
d = [−cₓ, −c_y, −c_z, cₓ, c_y, c_z]
ê = (cₓ, c_y, c_z),  P = I₃ − ê·êᵀ      (proyector 3×3)

K_M   = (E_t · A / L₀) · d·dᵀ            (0 si material destensado)
K_G   = (N / l) · [[P, −P], [−P, P]]     (0 si N = 0)
K_T   = K_M + K_G
F_int = N · d                            (0 si N = 0)

\end{lstlisting}

\subsection{Cargas distribuidas}
\begin{itemize}
  \item \texttt{compute\_body\_load(b)} reparte \texttt{b \ensuremath{\cdot} A \ensuremath{\cdot} L\ensuremath{_{0}} / 2} por nodo en cada componente global, evaluado en geometría de referencia.
\end{itemize}

\subsection{Régimen de validez}
\begin{itemize}
  \item \texttt{|\ensuremath{\varepsilon}| \ensuremath{\lesssim} 10\ensuremath{^{-}}\ensuremath{^{2}}} en régimen tensado.
  \item Grandes desplazamientos y rotaciones en el espacio.
  \item Cables completamente destensados aportan \texttt{K\_T = 0} al sistema global.
\end{itemize}

\subsection{Independencia del diseño}
No hereda de \texttt{Cable2DCorot} ni de \texttt{Truss3DCorot}. La maquinaria cinemática 3D se implementa íntegra dentro de la clase.

\subsection{Validación}
\begin{itemize}
  \item Tests: tests/test\_cable\_elements.py \ensuremath{\cdot} \texttt{TestCable3DCorot} (4 tests de aceptación).
  \item Archivo fuente: solidum/elements/cable.py \ensuremath{\cdot} clase \texttt{Cable3DCorot}.
  \item Spec: docs/specs/Cable3DCorot.md.
  \item Referencias: Crisfield §3.3; Belytschko §4.5; Irvine §1.2.
\end{itemize}


\section{Frame2DEuler --- marco/viga 2D Euler-Bernoulli}

Viga esbelta 2D con hipótesis de Euler-Bernoulli: secciones planas perpendiculares al eje neutro deformado. Dos nodos rígidamente conectados; transmite axial + cortante + momento flector.

\subsection{Cinemática}
\begin{itemize}
  \item Interpolación lineal del desplazamiento axial, Hermite cúbica del desplazamiento transverso.
  \item Pequeños desplazamientos y rotaciones (régimen geométricamente lineal).
  \item Configuración inicial fija: \texttt{L\ensuremath{_{0}}, c, s, T} se calculan una vez y no se actualizan.
\end{itemize}

\subsection{Formulación}
Matriz de transformación global\ensuremath{\to}local $\mathbf T$ ($6\times 6$) con cosenos directores del eje. Matriz local analítica:
\begin{lstlisting}
K_local = [[ EA/L,       0,       0, -EA/L,       0,       0 ],
           [    0, 12EI/L³,  6EI/L²,     0,-12EI/L³,  6EI/L²],
           [    0,  6EI/L²,   4EI/L,     0, -6EI/L²,   2EI/L],
           [-EA/L,       0,       0,  EA/L,       0,       0 ],
           [    0,-12EI/L³, -6EI/L²,     0, 12EI/L³, -6EI/L²],
           [    0,  6EI/L²,   2EI/L,     0, -6EI/L²,   4EI/L]]
K_global = Tᵀ · K_local · T
F_int    = Tᵀ · F_int_local    (componente axial corregida con σ·A del material)

\end{lstlisting}

\subsection{Cargas distribuidas}
\begin{itemize}
  \item \texttt{compute\_body\_load(b)} distribuye \texttt{q = A \ensuremath{\cdot} b} (carga por unidad de longitud) con la fórmula consistente Hermite cúbica: axial mitad-mitad por nodo, transversal \texttt{q\ensuremath{\cdot}L/2} en fuerza y \texttt{\ensuremath{\pm}q\ensuremath{\cdot}L\ensuremath{^{2}}/12} en momento (signo positivo en nodo i, negativo en j). Transformación local\ensuremath{\leftrightarrow}global vía la misma \texttt{T} del elemento. Útil para peso propio.
\end{itemize}

\subsection{Régimen de validez}
\begin{itemize}
  \item Vigas esbeltas (\texttt{L/h \ensuremath{\gtrsim} 10}); peraltadas \ensuremath{\to} \texttt{Frame2DTimoshenko}.
  \item Pequeños desplazamientos, pequeñas rotaciones.
  \item No captura \textbf{plasticidad por flexión} distribuida en la sección: \texttt{E\_tangent} escala toda la matriz de rigidez por igual. La fluencia que ocurre primero en la fibra inferior bajo flexión no se modela hasta que entre \texttt{FiberSection} (deuda técnica documentada).
\end{itemize}

\subsection{Independencia del diseño}
Vive en el paquete \texttt{solidum/elements/frame/} junto a \texttt{Frame2DTimoshenko} y \texttt{Frame2DEulerCorot}, sin herencia entre ellas. La construcción de longitud, cosenos directores y matriz de transformación 6\ensuremath{\times}6 se delega a \texttt{build\_geometry\_2d} en \texttt{solidum/elements/frame/\_shared.py}, compartida con \texttt{Frame2DTimoshenko} (la versión corotacional reconstruye T desde \texttt{alpha0} por su lógica propia).

\subsection{Validación}
\begin{itemize}
  \item Tests: tests/test\_frame.py \ensuremath{\cdot} \texttt{TestFrame2DEulerAcceptance} (3 tests de aceptación + registro, incluye flecha analítica del voladizo $PL^3/(3EI)$).
  \item Archivo fuente: solidum/elements/frame/euler.py \ensuremath{\cdot} clase \texttt{Frame2DEuler}.
  \item Spec: docs/specs/Frame2DEuler.md.
  \item Referencias: Bathe cap. 5; Cook §2.7.
\end{itemize}


\section{Frame2DTimoshenko --- marco/viga 2D Timoshenko}

Viga 2D con deformación por cortante transversal explícita. Dos nodos rígidamente conectados; transmite axial + cortante + momento flector. Apropiado para vigas peraltadas o cortas (\texttt{L/h \ensuremath{\lesssim} 10}).

\subsection{Cinemática}
\begin{itemize}
  \item Secciones planas, \textbf{no} perpendiculares al eje neutro deformado (rotación independiente de la pendiente de la elástica).
  \item Distorsión angular $\gamma = dv/ds - \theta$ tratada como campo independiente.
  \item Configuración inicial fija (régimen geométricamente lineal).
\end{itemize}

\subsection{Formulación}
Factor de corrección contra shear locking:
\begin{lstlisting}
Φ = 12 · E · I / (G · A_s · L²),      G = E / [2(1 + ν)]

\end{lstlisting}
Matriz local con coeficientes ajustados:
\begin{lstlisting}
a = 12·EI / [L³(1+Φ)],   b = 6·EI / [L²(1+Φ)]
c = (4+Φ)·EI / [L(1+Φ)], d = (2-Φ)·EI / [L(1+Φ)]

K_local = [[ EA/L,  0,  0, -EA/L,  0,  0],
           [    0,  a,  b,     0, -a,  b],
           [    0,  b,  c,     0, -b,  d],
           [-EA/L,  0,  0,  EA/L,  0,  0],
           [    0, -a, -b,     0,  a, -b],
           [    0,  b,  d,     0, -b,  c]]

\end{lstlisting}
Para \texttt{\ensuremath{\Phi} \ensuremath{\to} 0} (viga esbelta) se recupera la matriz Euler-Bernoulli.

\subsection{Parámetros}
\begin{itemize}
  \item \texttt{A}, \texttt{I}, \texttt{As} (área efectiva de cortante).
  \item \texttt{nu} (opcional): se toma de \texttt{material.nu} si está disponible; en su defecto, del parámetro del elemento con default 0.3 y aviso por consola.
\end{itemize}

\subsection{Cargas distribuidas}
\begin{itemize}
  \item \texttt{compute\_body\_load(b)} usa la misma fórmula consistente Hermite cúbica que \texttt{Frame2DEuler} (idéntica para carga uniforme --- la diferencia entre formulaciones aparece solo con cargas no uniformes o concentradas).
\end{itemize}

\subsection{Régimen de validez}
\begin{itemize}
  \item Vigas gruesas/peraltadas (\texttt{L/h \ensuremath{\lesssim} 10}); para esbeltas \ensuremath{\to} \texttt{Frame2DEuler} (más simple y sin shear locking).
  \item Pequeños desplazamientos, pequeñas rotaciones.
  \item No captura plasticidad distribuida: \texttt{E\_tangent} escala toda la matriz.
\end{itemize}

\subsection{Independencia del diseño}
Vive en el paquete \texttt{solidum/elements/frame/}, sin herencia con las otras vigas 2D. Comparte \texttt{build\_geometry\_2d} con \texttt{Frame2DEuler} en \texttt{solidum/elements/frame/\_shared.py} --- la construcción de la matriz de transformación 6\ensuremath{\times}6 ya no se duplica.

\subsection{Validación}
\begin{itemize}
  \item Tests: tests/test\_frame.py \ensuremath{\cdot} \texttt{TestFrame2DTimoshenkoAcceptance} (convergencia a Euler en viga esbelta, axial puro, simetría).
  \item Archivo fuente: solidum/elements/frame/timoshenko.py \ensuremath{\cdot} clase \texttt{Frame2DTimoshenko}.
  \item Spec: docs/specs/Frame2DTimoshenko.md.
  \item Referencias: Reddy cap. 5; Bathe §5.4.
\end{itemize}


\section{Frame2DEulerCorot --- viga 2D Euler-Bernoulli corotacional}

Viga 2D esbelta Euler-Bernoulli en formulación \textbf{corotacional} (Updated Lagrangian). Grandes desplazamientos y grandes rotaciones rígidas del elemento; deformaciones axiales pequeñas y rotaciones nodales deformacionales moderadas.

\subsection{Cinemática corotacional}
\begin{itemize}
  \item Rotación rígida $\alpha_e$ del eje separada de las rotaciones deformacionales de las secciones $\bar\theta_1, \bar\theta_2$.
  \item DOFs deformacionales locales: $[u_l = l - L_0, \bar\theta_1, \bar\theta_2]$.
  \item Todo en configuración corriente: $c, s, l, \alpha$ se recalculan por evaluación.
\end{itemize}

\subsection{Formulación}
\begin{lstlisting}
Rigidez local (3×3) en DOFs deformacionales:
  K_ll = diag-block([EA/L₀], [[4EI/L₀, 2EI/L₀], [2EI/L₀, 4EI/L₀]])

Matriz B (3×6) en configuración corriente acopla DOFs locales y globales.

Rigidez tangente:
  K_T = B ᵀ · K_ll · B + K_σ
  K_σ = (N/l)·z·zᵀ − ((M₁+M₂)/l²)·(r·zᵀ + z·rᵀ)

con r = [−c,−s,0,c,s,0], z = [−s,c,0,s,−c,0].

\end{lstlisting}

\subsection{Cargas distribuidas}
\begin{itemize}
  \item \texttt{compute\_body\_load(b)} evalúa la integral consistente en la configuración de referencia (misma fórmula Hermite cúbica que \texttt{Frame2DEuler}). Para grandes rotaciones con cargas conservadoras (gravedad) la diferencia frente a la integración sobre la geometría corriente es de segundo orden y se asume aceptable.
\end{itemize}

\subsection{Régimen de validez}
\begin{itemize}
  \item \texttt{|\ensuremath{\varepsilon}\_axial| \ensuremath{\lesssim} 10\ensuremath{^{-}}\ensuremath{^{2}}}.
  \item Rotaciones rígidas del elemento ilimitadas entre commits (hasta $|\alpha_e| < \pi$ por paso).
  \item Rotaciones deformacionales de las secciones moderadas (\texttt{|\ensuremath{\theta}̄| \ensuremath{\lesssim} 30\ensuremath{^{\circ}}}).
  \item Rotaciones continuas > \ensuremath{\pi} requerirían tracking de vueltas (fuera de alcance).
\end{itemize}

\subsection{Dominios que abre}
Pandeo por flexión de columnas esbeltas, post-pandeo, snap-through de arcos, brazos flexibles, cables con rigidez a flexión. Problemas que la viga lineal \texttt{Frame2DEuler} no puede atacar.

\subsection{Independencia del diseño}
Vive en el paquete \texttt{solidum/elements/frame/} junto a las otras vigas 2D, \textbf{sin herencia}: la cinemática corotacional se implementa íntegra dentro de la clase (reconstruye \texttt{T} desde \texttt{alpha0} en lugar de usar \texttt{build\_geometry\_2d} compartido). Sí comparte con \texttt{Frame2DEuler} y \texttt{Frame2DTimoshenko} los helpers libres de \texttt{\_shared.py} para carga de cuerpo (\texttt{\_frame2d\_consistent\_body\_load}), masa (\texttt{\_frame2d\_consistent\_mass\_local}) y traducción a \texttt{ElementForces} (\texttt{\_frame2d\_forces\_from\_local}).

\subsection{Validación}
\begin{itemize}
  \item Tests: tests/test\_frame.py \ensuremath{\cdot} \texttt{TestFrame2DEulerCorotAcceptance} (4 criterios físicos + \textbf{chequeo por diferencias finitas de $\mathbf K_T$ contra $\mathbf F_{\text{int}}$} + registro).
  \item Archivo fuente: solidum/elements/frame/euler\_corot.py \ensuremath{\cdot} clase \texttt{Frame2DEulerCorot}.
  \item Spec: docs/specs/Frame2DEulerCorot.md.
  \item Referencias: Crisfield §7.3; Belytschko §4.11; Wriggers cap. 4.
\end{itemize}


\section{Frame3D --- marco/viga 3D Euler-Bernoulli}

Viga 3D esbelta con 6 DOFs por nodo (\texttt{ux, uy, uz, rx, ry, rz}). Transmite axial + cortante en dos planos + flexión en dos planos + torsión. Régimen geométricamente lineal.

\subsection{Cinemática}
\begin{itemize}
  \item Secciones planas perpendiculares al eje neutro deformado (Euler-Bernoulli).
  \item Torsión de Saint-Venant pura (sin alabeo).
  \item Flexiones en ejes principales desacopladas.
  \item Configuración inicial fija (no hay variante corotacional 3D por ahora).
\end{itemize}

\subsection{Ejes locales y orientación de la sección}
El eje $\hat{\mathbf x}$ local va del nodo 1 al nodo 2. Los ejes $\hat{\mathbf y}, \hat{\mathbf z}$ de la sección se definen proyectando un \textbf{vector de referencia} (\texttt{ref\_vector}) sobre el plano perpendicular al eje. Default: \texttt{[0, 0, 1]} (eje z global); fallback a \texttt{[1, 0, 0]} con advertencia si la barra es casi vertical.

\subsection{Formulación}
Matriz de transformación $\mathbf T$ 12\ensuremath{\times}12 en bloques de $\boldsymbol{\lambda}$ (3\ensuremath{\times}3). Matriz local 12\ensuremath{\times}12 desacoplada:
\begin{lstlisting}
Axial      (2×2) : EA/L
Torsión    (2×2) : GJ/L       con G = E/[2(1+ν)]
Flexión xy (4×4) : 12EIz/L³, 6EIz/L², 4EIz/L, 2EIz/L
Flexión xz (4×4) : 12EIy/L³, 6EIy/L², 4EIy/L, 2EIy/L  (signos invertidos)

K_global = Tᵀ K_local T

\end{lstlisting}

\subsection{Parámetros}
\begin{itemize}
  \item \texttt{A}, \texttt{Iy}, \texttt{Iz} (área y momentos de inercia respecto a ejes locales y, z).
  \item \texttt{J} (constante torsional de Saint-Venant).
  \item \texttt{nu} (opcional, del material si lo expone, sino del elemento con default 0.3).
  \item \texttt{ref\_vector} (opcional, default \texttt{[0, 0, 1]}).
\end{itemize}

\subsection{Cargas distribuidas}
\begin{itemize}
  \item \texttt{compute\_body\_load(b)} proyecta \texttt{q = A \ensuremath{\cdot} b} sobre los ejes locales y reparte con la fórmula Hermite cúbica en ambos planos de flexión: \texttt{q \ensuremath{\cdot} L / 2} en fuerzas, \texttt{\ensuremath{\pm}q \ensuremath{\cdot} L\ensuremath{^{2}} / 12} en momentos (con signos \texttt{+} en nodo i y \texttt{−} en nodo j para \texttt{Mz}, y signos opuestos para \texttt{My} por la regla de la mano derecha). No genera torsión (carga uniforme sin excentricidad respecto al centro de cortadura). Útil para peso propio: \texttt{b = (0, 0, -\ensuremath{\rho}\ensuremath{\cdot}g)}.
\end{itemize}

\subsection{Régimen de validez}
\begin{itemize}
  \item Vigas esbeltas (\texttt{L/h \ensuremath{\gtrsim} 10}).
  \item Pequeños desplazamientos, pequeñas rotaciones.
  \item Sin alabeo (secciones macizas o cerradas).
  \item Para grandes rotaciones en 3D: pendiente \texttt{Frame3DCorot}.
\end{itemize}

\subsection{Masa lumped: limitación en orientación oblicua}
La masa lumped es estrictamente diagonal en ejes locales (\texttt{\ensuremath{\rho}AL/2} traslacional, \texttt{\ensuremath{\rho}I\ensuremath{\cdot}L/2} rotacional con \texttt{I} específica por DOF). Tras la rotación \texttt{T\textasciicircum{}T\ensuremath{\cdot}M\ensuremath{\cdot}T} a globales, el bloque traslacional 3\ensuremath{\times}3 por nodo permanece diagonal porque \texttt{m\_t\ensuremath{\cdot}I\ensuremath{_{3}}} es invariante bajo SO(3), pero el bloque rotacional 3\ensuremath{\times}3 queda \textbf{lleno} porque \texttt{m\_rx = \ensuremath{\rho}Jp\ensuremath{\cdot}L/2}, \texttt{m\_ry = \ensuremath{\rho}Iy\ensuremath{\cdot}L/2}, \texttt{m\_rz = \ensuremath{\rho}Iz\ensuremath{\cdot}L/2} son valores distintos para una sección real. \texttt{M\_global} resulta entonces \textbf{bloque-diagonal por nodo} (no estrictamente diagonal). Es la limitación estándar del lumping en frames 3D (Cook-Malkus-Plesha §11.4); aceptable para Newmark/HHT pero \texttt{CentralDifferenceSolver} rechaza con \texttt{ValueError} y orienta al usuario hacia Newmark o hacia un eje del elemento alineado con un eje global.

\subsection{Independencia del diseño}
Archivo propio. No hereda ni comparte helpers con \texttt{Frame2DEuler}, \texttt{Frame2DTimoshenko}, ni con los trusses 3D.

\subsection{Validación}
\begin{itemize}
  \item Tests: tests/test\_frame3d.py \ensuremath{\cdot} \texttt{TestFrame3DAcceptance} (5 criterios físicos + validación de \texttt{ref\_vector} + registro).
  \item Archivo fuente: solidum/elements/frame3d.py \ensuremath{\cdot} clase \texttt{Frame3D}.
  \item Spec: docs/specs/Frame3D.md.
  \item Referencias: Przemieniecki Tabla 11.3; Cook §2.8; Bathe cap. 5.
\end{itemize}


\section{Quad4 --- cuadrilátero bilineal 2D}

\begin{itemize}
  \item \textbf{Propósito}: continuo 2D isoparamétrico; problema mecánico plano.
  \item \textbf{DOFs por nodo}: \texttt{['ux', 'uy']} \ensuremath{\cdot} 4 nodos (antihorario) \ensuremath{\cdot} \texttt{STRAIN\_DIM = 3}.
  \item \textbf{Cinemática}: matriz \texttt{B(\ensuremath{\xi}, \ensuremath{\eta})} calculada por mapeo isoparamétrico estándar; deformación Voigt \texttt{[\ensuremath{\varepsilon}\_xx, \ensuremath{\varepsilon}\_yy, \ensuremath{\gamma}\_xy]}.
  \item \textbf{Integración}: por defecto Gauss 2\ensuremath{\times}2 (4 puntos). Configurable vía \texttt{quadrature} desde \texttt{QuadratureRegistry}. \textbf{Aviso}: integración reducida 1\ensuremath{\times}1 \ensuremath{\to} riesgo de hourglassing.
  \item \textbf{Parámetros}: \texttt{thickness} (espesor para estado plano), \texttt{quadrature} (opcional).
  \item \textbf{Cargas distribuidas}: \texttt{compute\_body\_load(b)} integra $\int \mathbf N^\top \mathbf b\, dV$ sobre el elemento; \texttt{compute\_edge\_traction(edge, t̄)} reparte una tracción uniforme sobre un borde (índices 0..3 según conectividad nodal). Tracción siempre en globales; sin soporte aún para tracción variable o presión normal.
  \item \textbf{Salida por Gauss}: \texttt{compute\_gauss\_state(U)} devuelve \ensuremath{\sigma}, \ensuremath{\varepsilon} y coordenadas (naturales y globales) en cada punto de Gauss. Habilita el suavizado nodal del exporter VTK (\texttt{Sigma\_*\_nodal}).
  \item \textbf{Implementación}: kernels \texttt{\_compute\_kinematics} y \texttt{\_compute\_integrands} con \texttt{@njit} (Numba) --- JIT en primera llamada. Viven en \texttt{solidum/elements/solid\_2d/\_shared.py}, compartidos con Tri3.
  \item \textbf{Limitaciones}: bloqueo volumétrico con materiales casi-incompresibles (\ensuremath{\nu} \ensuremath{\to} 0.5); en ese régimen usar formulación mixta (no implementada).
  \item \textbf{Archivo}: solidum/elements/solid\_2d/quad4.py
\end{itemize}


\section{Tri3 --- triángulo lineal 2D (CST)}

\begin{itemize}
  \item \textbf{Propósito}: triángulo de deformación constante; útil para transiciones de malla.
  \item \textbf{DOFs por nodo}: \texttt{['ux', 'uy']} \ensuremath{\cdot} 3 nodos \ensuremath{\cdot} \texttt{STRAIN\_DIM = 3}.
  \item \textbf{Cinemática}: matriz \texttt{B} constante (deformación uniforme dentro del elemento).
  \item \textbf{Integración}: 1 punto central, peso 0.5 (área triángulo en coordenadas naturales).
  \item \textbf{Parámetros}: \texttt{thickness}.
  \item \textbf{Cargas distribuidas}: \texttt{compute\_body\_load(b)} reparte 1/3 a cada nodo (exacto para b uniforme); \texttt{compute\_edge\_traction(edge, t̄)} reparte L/2 a cada uno de los 2 nodos del borde (índices 0..2). Tracción en globales.
  \item \textbf{Salida por Gauss}: \texttt{compute\_gauss\_state(U)} devuelve el único \ensuremath{\sigma}, \ensuremath{\varepsilon} y centroide del elemento (mismo formato que Quad4).
  \item \textbf{Limitaciones}: shear locking severo; convergencia lenta. \textbf{Preferir \texttt{Quad4}} salvo en transiciones donde Quad4 no encaja geométricamente.
  \item \textbf{Implementación}: kernel \texttt{\_compute\_kinematics\_tri3} con \texttt{@njit} en \texttt{solidum/elements/solid\_2d/\_shared.py}.
  \item \textbf{Archivo}: solidum/elements/solid\_2d/tri3.py
\end{itemize}


\section{Quad8 --- cuadrilátero serendípito 2D de orden 2}

\begin{itemize}
  \item \textbf{Propósito}: continuo 2D cuadrático sin shear locking severo en flexión; reproduce campos cuadráticos exactamente.
  \item \textbf{DOFs por nodo}: \texttt{['ux', 'uy']} \ensuremath{\cdot} 8 nodos (4 vértices antihorarios + 4 medios de borde) \ensuremath{\cdot} \texttt{STRAIN\_DIM = 3}.
  \item \textbf{Funciones de forma}: serendípitas (sin término \ensuremath{\xi}\ensuremath{^{2}}\ensuremath{\eta}\ensuremath{^{2}}).
  \item \textbf{Integración}: Gauss 3\ensuremath{\times}3 (9 puntos) por defecto. Reducida 2\ensuremath{\times}2 disponible vía \texttt{quadrature} (con riesgo de modos espurios).
  \item \textbf{Cargas}: body load por cuadratura; tracción de borde uniforme reparte 1/6, 4/6, 1/6 (vértice, medio, vértice).
  \item \textbf{Salida por Gauss}: \texttt{compute\_gauss\_state(U)} con 9 puntos por defecto.
  \item \textbf{Spec}: docs/specs/Quad8.md.
  \item \textbf{Archivo}: solidum/elements/solid\_2d/quad8.py (subclase de la base interna \texttt{\_HigherOrderSolid2D} definida en \texttt{\_shared.py}, compartida con Quad9 y Tri6).
\end{itemize}


\section{Quad9 --- cuadrilátero Lagrangiano 2D de orden 2}

\begin{itemize}
  \item \textbf{Propósito}: continuo 2D cuadrático con espacio polinómico completo $Q_2$; añade un noveno nodo central interior al Quad8.
  \item \textbf{DOFs por nodo}: \texttt{['ux', 'uy']} \ensuremath{\cdot} 9 nodos \ensuremath{\cdot} \texttt{STRAIN\_DIM = 3}.
  \item \textbf{Funciones de forma}: producto tensorial Lagrange 1D-1D.
  \item \textbf{Integración}: Gauss 3\ensuremath{\times}3.
  \item \textbf{Cargas y salida por Gauss}: idénticas a Quad8 (el nodo 8 es interior y no participa en bordes).
  \item \textbf{Spec}: docs/specs/Quad9.md.
  \item \textbf{Archivo}: solidum/elements/solid\_2d/quad9.py (subclase de la base interna \texttt{\_HigherOrderSolid2D}, compartida con Quad8 y Tri6).
\end{itemize}


\section{Tri6 --- triángulo 2D cuadrático completo P\ensuremath{_{2}}}

\begin{itemize}
  \item \textbf{Propósito}: triángulo isoparamétrico cuadrático que cura el shear locking del Tri3; reproduce campos cuadráticos exactamente.
  \item \textbf{DOFs por nodo}: \texttt{['ux', 'uy']} \ensuremath{\cdot} 6 nodos (3 vértices + 3 medios de borde) \ensuremath{\cdot} \texttt{STRAIN\_DIM = 3}.
  \item \textbf{Integración}: 3 puntos en los puntos medios (cuadratura \texttt{tri\_3}).
  \item \textbf{Cargas}: body load por cuadratura; tracción de borde reparte 1/6, 4/6, 1/6.
  \item \textbf{Spec}: docs/specs/Tri6.md.
  \item \textbf{Archivo}: solidum/elements/solid\_2d/tri6.py (subclase de la base interna \texttt{\_HigherOrderSolid2D} en \texttt{\_shared.py}, compartida con Quad8 y Quad9; declara \texttt{\_MASS\_QUADRATURE = "tri\_6"} para integrar exactamente la masa consistente, que la cuadratura del elemento subintegraría).
\end{itemize}


Sólidos tridimensionales con convención Voigt 6D del proyecto (\texttt{Reglas.md §5}):
`\texttt{[\ensuremath{\varepsilon}\_xx, \ensuremath{\varepsilon}\_yy, \ensuremath{\varepsilon}\_zz, \ensuremath{\gamma}\_xy, \ensuremath{\gamma}\_yz, \ensuremath{\gamma}\_xz]}\texttt{ con }\texttt{\ensuremath{\gamma}\_ij = 2\ensuremath{\cdot}\ensuremath{\varepsilon}\_ij}`. Esta familia
\textbf{no expone \texttt{internal\_forces}} --- la primitiva natural de salida en sólidos
es `\texttt{compute\_gauss\_state(U)}\texttt{, que devuelve }\texttt{\{stress, strain, points\_*\}}`
por punto de integración (ADR 0012: cierre del contrato \texttt{internal\_forces}
por dominio explícito).

\section{Hex8 --- hexaedro trilineal 3D}

\begin{itemize}
  \item \textbf{Propósito}: sólido 3D isoparamétrico de primer orden, espejo natural de Quad4.
  \item \textbf{DOFs por nodo}: \texttt{['ux', 'uy', 'uz']} \ensuremath{\cdot} 8 nodos (orden VTK\_HEXAHEDRON) \ensuremath{\cdot} \texttt{STRAIN\_DIM = 6} \ensuremath{\cdot} \texttt{N\_INTEGRATION\_POINTS = 8} (default Gauss 2\ensuremath{\times}2\ensuremath{\times}2).
  \item \textbf{Cinemática}: matriz \texttt{B(\ensuremath{\xi}, \ensuremath{\eta}, \ensuremath{\zeta})} 6\ensuremath{\times}24 por mapeo isoparamétrico estándar; deformación Voigt 3D.
  \item \textbf{Integración}: por defecto Gauss 2\ensuremath{\times}2\ensuremath{\times}2 (8 puntos). Configurable vía \texttt{quadrature} (\texttt{hex\_3x3x3} para no-lineales severos, \texttt{hex\_1x1x1} reducido con riesgo de hourglass).
  \item \textbf{Caras (ADR 0012)}: 6 caras numeradas con normal saliente --- 0 (−\ensuremath{\zeta}), 1 (+\ensuremath{\zeta}), 2 (−\ensuremath{\eta}), 3 (+\ensuremath{\xi}), 4 (+\ensuremath{\eta}), 5 (−\ensuremath{\xi}). Ver \texttt{Hex8.FACE\_NODES}.
  \item \textbf{Cargas distribuidas}: \texttt{compute\_body\_load(b)} integra \texttt{\ensuremath{\int} N\ensuremath{^{\top}}b dV}; \texttt{compute\_face\_traction(face, t̄)} reparte tracción uniforme sobre una cara con cuadratura 2D Gauss 2\ensuremath{\times}2 (4 puntos por cara). \texttt{t̄} en globales.
  \item \textbf{Salida por Gauss}: \texttt{compute\_gauss\_state(U)} devuelve \texttt{\{points\_natural, points\_global, strain (n\_g, 6), stress (n\_g, 6)\}}. Sin \texttt{internal\_forces} (ADR 0012).
  \item \textbf{Implementación}: kernels \texttt{\_compute\_kinematics\_hex8}, \texttt{\_det\_jacobian\_hex8}, \texttt{\_shape\_functions\_hex8} con \texttt{@njit} en \texttt{solidum/elements/solid\_3d/\_shared.py}. Mass lumping vía \texttt{lump\_hrz} con \texttt{n\_translational\_dirs=3} y \texttt{\_expand\_scalar\_mass\_3d}.
  \item \textbf{Limitaciones declaradas}:
  \item \textbf{Locking volumétrico} con \ensuremath{\nu} \ensuremath{\to} 0.5 --- sin mitigación (política idéntica al Quad4 2D). Blindado por \texttt{tests/test\_volumetric\_locking\_3d.py}.
  \item \textbf{Hourglass} con integración reducida \texttt{hex\_1x1x1}: 12 modos espurios; sin estabilización Flanagan-Belytschko.
  \item \textbf{Shear locking} con malla coarse 1 capa en sección: documentado en \texttt{tests/validation/test\_macneal\_beam\_3d.py}.
  \item \textbf{Spec}: docs/specs/Hex8.md.
  \item \textbf{Archivo}: solidum/elements/solid\_3d/hex8.py.
\end{itemize}


\section{Hex20 --- hexaedro serendípito 3D de orden 2 (20 nodos)}

\begin{itemize}
  \item \textbf{Propósito}: sólido 3D isoparamétrico de orden cuadrático con interpolación serendípita; análogo 3D natural del \texttt{Quad8}. Reproduce campos cuadráticos completos en 3D y mitiga drásticamente el shear locking del \texttt{Hex8} en problemas de flexión y geometrías curvas.
  \item \textbf{DOFs por nodo}: \texttt{['ux', 'uy', 'uz']} \ensuremath{\cdot} 20 nodos (orden VTK\_QUADRATIC\_HEXAHEDRON: 8 vértices + 12 medios de arista, sin nodos de cara ni centroide) \ensuremath{\cdot} \texttt{STRAIN\_DIM = 6} \ensuremath{\cdot} \texttt{N\_INTEGRATION\_POINTS = 27} (default Gauss 3\ensuremath{\times}3\ensuremath{\times}3).
  \item \textbf{Cinemática}: matriz \texttt{B(\ensuremath{\xi}, \ensuremath{\eta}, \ensuremath{\zeta})} 6\ensuremath{\times}60 sobre Voigt 6D del proyecto (ADR 0012). Funciones de forma serendípitas: vértices con multiplicador \texttt{(\ensuremath{\xi}\_i \ensuremath{\xi} + \ensuremath{\eta}\_i \ensuremath{\eta} + \ensuremath{\zeta}\_i \ensuremath{\zeta} - 2)} que anula la función en todos los demás nodos; medios \texttt{(1 - \ensuremath{\alpha}\ensuremath{^{2}})} con \texttt{\ensuremath{\alpha}} el eje del cual son punto medio.
  \item \textbf{Integración}: por defecto Gauss 3\ensuremath{\times}3\ensuremath{\times}3 (27 puntos, exacta para K con Jacobiano constante y para masa consistente cuadrática). Configurable vía \texttt{quadrature} (\texttt{hex\_2x2x2} reducida con 6 modos de hourglass espurios por elemento aislado --- sin estabilización implementada).
  \item \textbf{Caras (ADR 0012)}: 6 caras numeradas con normal saliente, \textbf{8 nodos por cara} (4 vértices + 4 medios). Numeración paritaria con \texttt{Hex8} en los vértices; cada cara añade los 4 medios en el mismo orden cíclico (medio entre vértice k y k+1). Ver \texttt{Hex20.FACE\_NODES}.
  \item \textbf{Cargas distribuidas}: \texttt{compute\_body\_load(b)} integra \texttt{\ensuremath{\int} N\ensuremath{^{\top}}b dV} con la cuadratura del elemento (reparto \textbf{no uniforme} entre vértices y medios serendípitos, suma global \texttt{b\ensuremath{\cdot}V\_e} exacta). \texttt{compute\_face\_traction(face, t̄)} reparte la tracción uniforme sobre la cara serendípita Quad8 con Gauss 3\ensuremath{\times}3 (exacto para tracción constante en cara plana): cada vértice recibe \texttt{-A\ensuremath{\cdot}t̄/12} y cada medio \texttt{4\ensuremath{\cdot}A\ensuremath{\cdot}t̄/12}; los signos negativos en vértices son el fenómeno serendípito conocido del Quad8 2D (Cook-Malkus-Plesha-Witt §6.5).
  \item \textbf{Salida por Gauss}: \texttt{compute\_gauss\_state(U)} devuelve \texttt{\{points\_natural, points\_global, strain (n\_g, 6), stress (n\_g, 6)\}}. Sin \texttt{internal\_forces} (ADR 0012).
  \item \textbf{Masa}: consistente con cuadratura \texttt{hex\_3x3x3} \textbf{fija} (independiente de la cuadratura elegida para K, para que la masa siempre quede exacta). Lumped por HRZ canónico --- para Hex20, row-sum produciría masas negativas en vértices serendípitos (Bathe FEP §9.2.4); HRZ es la única opción razonable.
  \item \textbf{Implementación}: funciones de forma y derivadas en \texttt{\_N\_hex20}/\texttt{\_dN\_hex20} (numpy puro, no \texttt{@njit} por la complejidad de la rama por tipo de nodo) en \texttt{solidum/elements/solid\_3d/\_shared.py}. Kinematics genérico de orden superior 3D en \texttt{\_kinematics\_higher\_order\_3d} (también \texttt{\_shared.py}), reutilizable por Hex27 y Tet10 en sub-fases siguientes. Cara serendípita integrada con \texttt{\_N\_quad8}/\texttt{\_dN\_quad8} importados del paquete 2D.
  \item \textbf{Limitaciones declaradas}:
  \item \textbf{Locking volumétrico} con \ensuremath{\nu} \ensuremath{\to} 0.5 \textbf{atenuado} respecto al \texttt{Hex8} pero aún presente --- sin mitigación implementada (política idéntica a \texttt{Quad8}/\texttt{Hex8}). Blindado por \texttt{tests/test\_volumetric\_locking\_3d.py::TestHex20VolumetricLocking} (ratio u(0.4999)/u(0.3) \ensuremath{\approx} 0.80 vs < 0.6 en \texttt{Hex8}).
  \item \textbf{Hourglass} con integración reducida \texttt{hex\_2x2x2}: 6 modos espurios por elemento aislado; en mallas ensambladas con BC razonable los modos no propagan, pero un elemento aislado o capa pobre los mostraría. Sin estabilización.
  \item \textbf{Shear locking} \textbf{drásticamente mitigado} respecto al \texttt{Hex8} en problemas de flexión: cuantificado en \texttt{tests/validation/test\_macneal\_beam\_3d.py} (Hex20 6\ensuremath{\times}1\ensuremath{\times}1 alcanza 97\% de u\_EB, vs < 55\% de Hex8 12\ensuremath{\times}1\ensuremath{\times}1).
  \item \textbf{Validación externa}: NAFEMS LE10 thick plate pressure (\texttt{tests/validation/test\_nafems\_le10.py}, 6\ensuremath{\times}6\ensuremath{\times}2 \ensuremath{\to} \ensuremath{\sigma}\_yy(D) \ensuremath{\approx} −5.38 MPa canónico dentro de banda 10\%). \textbf{Lamé thick cylinder 3D} (\texttt{tests/validation/test\_lame\_cylinder\_3d.py}, 4\ensuremath{\times}4\ensuremath{\times}1 \ensuremath{\to} error L\ensuremath{^{2}}-rel < 4\% en \ensuremath{\sigma} y < 2\% en u\_r contra solución cerrada de Timoshenko-Goodier §28; convergencia h monótona del error L\ensuremath{^{2}}(\ensuremath{\sigma}\_\ensuremath{\theta}\ensuremath{\theta}) verificada). Es la primera demostración cuantitativa del proyecto de \textbf{capacidad isoparamétrica curva} --- los mid-edges quedan automáticamente sobre la superficie cilíndrica por la no-linealidad cos/sin del mapeo angular.
  \item \textbf{Spec}: docs/specs/Hex20.md.
  \item \textbf{Archivo}: solidum/elements/solid\_3d/hex20.py.
\end{itemize}


\section{Hex27 --- hexaedro Lagrangiano 3D triquadrático (27 nodos)}

\begin{itemize}
  \item \textbf{Propósito}: sólido 3D isoparamétrico Lagrangiano completo de orden cuadrático; análogo 3D del \texttt{Quad9}. Reproduce \textbf{todos} los polinomios triquadráticos `\texttt{\ensuremath{\xi}\textasciicircum{}a \ensuremath{\eta}\textasciicircum{}b \ensuremath{\zeta}\textasciicircum{}c}\texttt{ con }\texttt{a, b, c \ensuremath{\in} \{0, 1, 2\}}\texttt{, incluyendo los términos puramente triquadráticos (}\texttt{\ensuremath{\xi}\ensuremath{^{2}}\ensuremath{\eta}\ensuremath{^{2}}\ensuremath{\zeta}\ensuremath{^{2}}}\texttt{ y combinaciones) que faltan en el }Hex20\texttt{ serendípito. Se prefiere sobre }Hex20\texttt{ cuando la geometría tiene curvatura severa y se busca convergencia óptima en problemas donde el contenido espectral cubre el espacio completo; en flexión simple con geometría rectilínea }Hex27\texttt{ y }Hex20` dan resultados prácticamente idénticos a un coste mayor para el primero (Bathe FEP §5.3.2).
  \item \textbf{DOFs por nodo}: \texttt{['ux', 'uy', 'uz']} \ensuremath{\cdot} 27 nodos (orden VTK\_TRIQUADRATIC\_HEXAHEDRON: 0-7 vértices, 8-19 medios de arista (paritarios con Hex20), 20-25 centros de cara en orden `\texttt{(-x, +x, -y, +y, -z, +z)}\texttt{, 26 centro de cuerpo) \ensuremath{\cdot} }STRAIN\_DIM = 6\texttt{ \ensuremath{\cdot} }N\_INTEGRATION\_POINTS = 27` (default Gauss 3\ensuremath{\times}3\ensuremath{\times}3).
  \item \textbf{Cinemática}: matriz \texttt{B(\ensuremath{\xi}, \ensuremath{\eta}, \ensuremath{\zeta})} 6\ensuremath{\times}81 sobre Voigt 6D del proyecto (ADR 0012). Funciones de forma como producto tensorial de Lagrange cuadrático 1D: `\texttt{N\_n = L\_\{i\_n\}(\ensuremath{\xi})\ensuremath{\cdot}L\_\{j\_n\}(\ensuremath{\eta})\ensuremath{\cdot}L\_\{k\_n\}(\ensuremath{\zeta})}\texttt{ con }\texttt{i\_n, j\_n, k\_n \ensuremath{\in} \{-1, 0, +1\}}`.
  \item \textbf{Integración}: por defecto Gauss 3\ensuremath{\times}3\ensuremath{\times}3 (27 puntos, exacta para K y masa consistente sobre Jacobiano constante). Reducida \texttt{hex\_2x2x2} disponible pero \textbf{introduce 27 modos de hourglass espurios por elemento aislado} --- la combinación más problemática del catálogo 3D. Recomendación operativa estricta: usar 3\ensuremath{\times}3\ensuremath{\times}3.
  \item \textbf{Caras (ADR 0012)}: 6 caras numeradas con normal saliente, \textbf{9 nodos por cara} (4 vértices + 4 medios + 1 centro de cara). Numeración paritaria con \texttt{Hex8}/\texttt{Hex20} en los vértices y medios; el centro de cara correspondiente se añade al final (índice 20-25 según la cara). Las funciones de forma 2D de cara son las del \texttt{Quad9} Lagrangiano (\texttt{\_N\_quad9}/\texttt{\_dN\_quad9} del paquete 2D).
  \item \textbf{Cargas distribuidas}: implementadas en la base \texttt{\_HigherOrderSolid3D} (paritaria con \texttt{\_HigherOrderSolid2D}). \texttt{compute\_body\_load(b)} integra \texttt{\ensuremath{\int} N\ensuremath{^{\top}}b dV}; \texttt{compute\_face\_traction(face, t̄)} integra sobre la cara Quad9 con cuadratura 3\ensuremath{\times}3. Suma global \texttt{b\ensuremath{\cdot}V\_e} (body) y \texttt{A\ensuremath{\cdot}t̄} (face) exactas en cara plana con campo uniforme.
  \item \textbf{Salida por Gauss}: \texttt{compute\_gauss\_state(U)} con 27 puntos. Sin \texttt{internal\_forces} (ADR 0012).
  \item \textbf{Masa}: consistente con cuadratura \texttt{hex\_3x3x3} fija (independiente de la cuadratura de K). Lumped HRZ canónico --- row-sum produciría masas negativas en nodos interiores (Bathe FEP §9.2.4).
  \item \textbf{Implementación}: funciones de forma y derivadas en \texttt{\_N\_hex27}/\texttt{\_dN\_hex27} (numpy puro, vía productos tensoriales de los polinomios Lagrange 1D \texttt{\_L\_lagrange\_quad}/\texttt{\_dL\_lagrange\_quad}). Subclase de la base \texttt{\_HigherOrderSolid3D} introducida con la entrada del \texttt{Hex27} (regla de los dos casos reales aplicada). El cuerpo de \texttt{hex27.py} queda como declaración de atributos.
  \item \textbf{Limitaciones declaradas}:
  \item \textbf{Coste computacional}: 81 DOFs/elemento vs 60 del \texttt{Hex20} (35\% más). Para flexión simple \texttt{Hex20} da resultados equivalentes a menor coste.
  \item \textbf{Locking volumétrico} con \ensuremath{\nu} \ensuremath{\to} 0.5 comparable al \texttt{Hex20} (ratio \ensuremath{\approx} 0.80 vs < 0.6 del \texttt{Hex8}). Sin mitigación implementada.
  \item \textbf{Hourglass} con integración reducida \texttt{hex\_2x2x2}: 27 modos espurios por elemento aislado (peor del catálogo 3D); en mallas estructuradas la mayoría no propaga tras ensamblar.
  \item \textbf{Validación externa}: NAFEMS LE10 thick plate pressure (\texttt{tests/validation/test\_nafems\_le10.py}, 5\ensuremath{\times}5\ensuremath{\times}2 \ensuremath{\to} \ensuremath{\sigma}\_yy(D) \ensuremath{\approx} −5.38 MPa canónico dentro de banda 10\%). \textbf{Lamé thick cylinder 3D} (\texttt{tests/validation/test\_lame\_cylinder\_3d.py}, 4\ensuremath{\times}4\ensuremath{\times}1 \ensuremath{\to} error L\ensuremath{^{2}}-rel < 4\% en \ensuremath{\sigma} y < 2\% en u\_r contra Timoshenko-Goodier §28; convergencia h monótona). Mid-edges automáticos sobre superficie cilíndrica por la no-linealidad del mapeo angular.
  \item \textbf{Spec}: docs/specs/Hex27.md.
  \item \textbf{Archivo}: solidum/elements/solid\_3d/hex27.py.
\end{itemize}


\section{Tet10 --- tetraedro cuadrático 3D (10 nodos)}

\begin{itemize}
  \item \textbf{Propósito}: sólido 3D isoparamétrico cuadrático sobre simplex tetraédrico. Análogo 3D del \texttt{Tri6} 2D. Mitiga drásticamente el shear locking y el locking volumétrico del \texttt{Tet4} (CST 3D). Adecuado para mallas no estructuradas y geometrías complejas donde el \texttt{Hex20}/\texttt{Hex27} requeriría mapeo estructurado.
  \item \textbf{DOFs por nodo}: \texttt{['ux', 'uy', 'uz']} \ensuremath{\cdot} 10 nodos (orden VTK\_QUADRATIC\_TETRA: 0-3 vértices idénticos al \texttt{Tet4}, 4-9 medios de arista en orden \texttt{(0-1, 1-2, 2-0, 0-3, 1-3, 2-3)}) \ensuremath{\cdot} \texttt{STRAIN\_DIM = 6} \ensuremath{\cdot} \texttt{N\_INTEGRATION\_POINTS = 4} (default Stroud \texttt{tet\_4}).
  \item \textbf{Cinemática}: matriz \texttt{B(\ensuremath{\xi}, \ensuremath{\eta}, \ensuremath{\zeta})} 6\ensuremath{\times}30 sobre Voigt 6D del proyecto (ADR 0012). Funciones de forma en coordenadas baricéntricas \texttt{L\_i}: vértices \texttt{N\_i = L\_i(2L\_i - 1)}; medios \texttt{N\_ij = 4\ensuremath{\cdot}L\_i\ensuremath{\cdot}L\_j}. Reproduce todos los polinomios cuadráticos completos en \texttt{(\ensuremath{\xi}, \ensuremath{\eta}, \ensuremath{\zeta})}.
  \item \textbf{Integración}: por defecto Stroud \texttt{tet\_4} (4 puntos, orden 2 --- exacta para K con Jacobiano constante, pues B es lineal en barycentric). Alternativa \texttt{tet\_15} (Keast 15 puntos, orden 5) para geometrías distorsionadas. \textbf{La masa consistente usa \texttt{tet\_15} fijo} (independiente de la cuadratura de K) para integrar exactamente el producto cuadrático\ensuremath{\times}cuadrático y garantizar correctitud en análisis modal/transitorio.
  \item \textbf{Caras (ADR 0012)}: 4 caras triangulares con normal saliente, paritarias con \texttt{Tet4} en los vértices. Cada cara añade los 3 medios de arista correspondientes en orden Tri6 (vértices antihorarios, medios entre 0-1, 1-2, 2-0). Cara i opuesta al vértice i. Las funciones de forma 2D de cara son las del \texttt{Tri6} (\texttt{\_N\_tri6}/\texttt{\_dN\_tri6} del paquete 2D).
  \item \textbf{Cargas distribuidas}: implementadas en la base \texttt{\_HigherOrderSolid3D}. \texttt{compute\_body\_load(b)} integra \texttt{\ensuremath{\int} N\ensuremath{^{\top}}b dV} con la cuadratura del elemento; reparto no uniforme entre vértices y medios, suma global \texttt{b\ensuremath{\cdot}V\_e} exacta. \texttt{compute\_face\_traction(face, t̄)} integra sobre la cara triangular Tri6 con cuadratura \texttt{tri\_3} (exacta para tracción uniforme sobre cara plana).
  \item \textbf{Salida por Gauss}: \texttt{compute\_gauss\_state(U)} con 4 puntos por defecto.
  \item \textbf{Implementación}: funciones de forma y derivadas en \texttt{\_N\_tet10}/\texttt{\_dN\_tet10} (numpy puro, vía coordenadas baricéntricas) en \texttt{solidum/elements/solid\_3d/\_shared.py}. Subclase de la base \texttt{\_HigherOrderSolid3D}. Las nuevas cuadraturas \texttt{tet\_4} y \texttt{tet\_15} viven en \texttt{solidum/math/integration.py}.
  \item \textbf{Limitaciones declaradas}:
  \item \textbf{Locking volumétrico} con \ensuremath{\nu} \ensuremath{\to} 0.5 atenuado respecto al \texttt{Tet4} pero presente; sin mitigación implementada.
  \item \textbf{Distorsión severa}: para Tet10 con mid-edges curvos, el Jacobiano varía espacialmente; el usuario debería seleccionar \texttt{tet\_15} vía el parámetro \texttt{quadrature}.
  \item \textbf{Superficie curva con descomposición ingenua hex\ensuremath{\to}5tets}: no converge. La cara del tet que toca una superficie curva tiene 3 mid-edges; con descomposición genérica de cada hex en 5 tets, sólo 1 de los 3 mid-edges está sobre una arista del grid paramétrico (curva), los otros 2 corresponden a aristas diagonales del hex que se vuelven aristas del tet con mid-edges rectos. La representación isoparamétrica se pierde y el error \ensuremath{\sigma}\_rr \textbf{crece con refinamiento h} (ver \texttt{tests/validation/test\_lame\_cylinder\_3d.py} --- tests Tet10 con \texttt{@pytest.mark.skip} y motivo). Para validar Tet10 sobre superficie curva se requiere mesher tetraédrico nativo (gmsh API) o descomposición específica del dominio en tets cuyas aristas estén todas sobre el grid paramétrico. Deuda técnica \#8 en STATUS. La capacidad del Tet10 sobre geometría plana está validada en \texttt{test\_cube\_lame\_3d.py} (cubo Lamé exacto a precisión máquina), así que NO es fallo del elemento.
  \item \textbf{Validación externa}: cubo Lamé 3D tracción uniaxial sobre malla 5-Tet (exacto a precisión máquina, \texttt{test\_cube\_lame\_3d.py}). Superficie curva diferida --- ver limitación arriba.
  \item \textbf{Spec}: docs/specs/Tet10.md.
  \item \textbf{Archivo}: solidum/elements/solid\_3d/tet10.py.
\end{itemize}


\section{Tet4 --- tetraedro lineal 3D (CST 3D)}

\begin{itemize}
  \item \textbf{Propósito}: sólido 3D isoparamétrico de primer orden, espejo natural de Tri3. Útil para mallas no estructuradas y transiciones.
  \item \textbf{DOFs por nodo}: \texttt{['ux', 'uy', 'uz']} \ensuremath{\cdot} 4 nodos (orden VTK\_TETRA, volumen positivo) \ensuremath{\cdot} \texttt{STRAIN\_DIM = 6} \ensuremath{\cdot} \texttt{N\_INTEGRATION\_POINTS = 1}.
  \item \textbf{Cinemática}: \texttt{B} constante 6\ensuremath{\times}12 (deformación uniforme dentro del elemento --- CST 3D). Reproduce campos lineales exactamente.
  \item \textbf{Integración}: 1 punto baricéntrico con peso 1/6 (exacto por linealidad de B).
  \item \textbf{Caras (ADR 0012)}: 4 caras triangulares numeradas con normal saliente; cara \texttt{i} opuesta al nodo \texttt{i}. Ver \texttt{Tet4.FACE\_NODES}.
  \item \textbf{Cargas distribuidas}: \texttt{compute\_body\_load(b)} reparte \texttt{b\ensuremath{\cdot}V\_e/4} a cada nodo (exacto); \texttt{compute\_face\_traction(face, t̄)} reparte \texttt{A\_cara\ensuremath{\cdot}t̄/3} a cada uno de los 3 nodos de la cara, cero al nodo opuesto.
  \item \textbf{Salida por Gauss}: \texttt{compute\_gauss\_state(U)} con 1 punto (centroide); \ensuremath{\sigma} y \ensuremath{\varepsilon} constantes sobre el elemento.
  \item \textbf{Implementación}: kernel \texttt{\_compute\_kinematics\_tet4} con \texttt{@njit} en \texttt{\_shared.py}. Masa consistente analítica \texttt{M\_ij = \ensuremath{\rho}\ensuremath{\cdot}V\_e\ensuremath{\cdot}(1+\ensuremath{\delta}\_ij)/20}; lumped HRZ canónico (= row-sum por simetría).
  \item \textbf{Limitaciones declaradas}:
  \item \textbf{Shear locking severo} en flexión --- peor que Hex8 por la pobreza del espacio lineal. Recomendación: usar Hex8 en mallas hexaédricas; Tet4 para transiciones de malla no estructurada o cuando geometría compleja impide hexaedros.
  \item \textbf{Locking volumétrico} con \ensuremath{\nu} \ensuremath{\to} 0.5 --- aún peor que en Hex8.
  \item \textbf{Spec}: docs/specs/Tet4.md.
  \item \textbf{Archivo}: solidum/elements/solid\_3d/tet4.py.
\end{itemize}


Subfamilia de elementos con \textbf{DOFs enriquecidos elementales} y \textbf{condensación estática local}: el ensamblador no ve los grados de libertad del salto. Cuando se cumple un criterio de activación (Rankine en fase 1), el elemento materializa una discontinuidad interna \texttt{Γ\_d} y enriquece su cinemática con un salto \texttt{[[u]]} gobernado por un material cohesivo (\texttt{CohesiveMaterial}, ver el catálogo de materiales §"Materiales cohesivos"). La activación se evalúa en el hook \texttt{Element.prepare\_step(U\_committed)} que los solvers no lineales invocan \textbf{una vez por paso}, antes del Newton (anti-chattering, ADR 0010 §5).

\section{CST\_Embedded2D --- triángulo CST con discontinuidad interior embebida (KOS)}

\begin{itemize}
  \item \textbf{Propósito}: introducir fractura computacional en aproximación discreta sobre el CST padre. Fiel a Retama (2010), Caps. 2, 5, 6 y 7: cinemática KOS, condensación estática local, longitud efectiva \texttt{l\_d = (A\_e/h)\ensuremath{\cdot}cos(\ensuremath{\theta}−\ensuremath{\alpha})}.
  \item \textbf{DOFs por nodo}: \texttt{['ux', 'uy']} \ensuremath{\cdot} 3 nodos \ensuremath{\cdot} \texttt{STRAIN\_DIM = 3} \ensuremath{\cdot} \texttt{N\_INTEGRATION\_POINTS = 1}.
  \item \textbf{DOFs enriquecidos (elementales, no globales)}: \texttt{[[u]] \ensuremath{\in} \ensuremath{\mathbb{R}}\ensuremath{^{2}}} en frame local \texttt{(n, s)} de \texttt{Γ\_d}. Se condensan dentro de \texttt{compute\_element\_state} y nunca llegan al ensamblador.
  \item \textbf{Estado intacto}: bit-exact con Tri3 (mismo \texttt{B}, mismo material, misma cuadratura).
  \item \textbf{Estado agrietado}: Newton local sobre \texttt{[[u]]} hasta \texttt{R\textasciicircum{}\{[[u]]\} = 0}; después la condensación \texttt{K\_cond = K\_dd − K\_du \ensuremath{\cdot} K\_\{[[u]][[u]]\}\ensuremath{^{-}}¹ \ensuremath{\cdot} K\_du\textasciicircum{}T} se devuelve al ensamblador. El bulk descarga elásticamente conforme \texttt{[[u]]} crece (discrete approach); la disipación va al cohesivo en \texttt{Γ\_d}.
  \item \textbf{Activación}: criterio Rankine (\texttt{\ensuremath{\sigma}\_I > \ensuremath{\sigma}\_t0} del cohesivo) en el centroide del CST, con el estado convergido del paso anterior. \textbf{Irreversible}: una vez activada, la discontinuidad persiste aunque el siguiente paso descargue.
  \item \textbf{Materiales aceptados}: bulk \textbf{solo \texttt{Elastic2D}} en fase 1 (la discrete approach presupone bulk elástico, ADR 0010); cohesivo cualquier \texttt{CohesiveMaterial} con \texttt{JUMP\_DIM = 2}.
  \item \textbf{Estado de la discontinuidad}: \texttt{DiscontinuityState} (en solidum/core/discontinuity\_state.py) --- paralela a \texttt{ElementState}, semántica trial/commit.
  \item \textbf{YAML}:
\end{itemize}
  \begin{lstlisting}[language=yaml]
  cohesive_materials:
    - {id: 1, type: CohesiveDamageIsotropic, sigma_t0: 2.5e6, G_f: 100.0,
       K_e: 1.0e13, softening: linear}
  elements:
    - {id: 1, type: CST_Embedded2D, nodes: [1, 2, 3],
       material: 1, cohesive_material: 1}
  
\end{lstlisting}
\begin{itemize}
  \item \textbf{Post-procesamiento}: \texttt{compute\_gauss\_state(U)} añade clave \texttt{'discontinuity'} con \texttt{\{normal, tangent, centroid, solitary\_node, l\_d, jump, traction, damage\}} cuando el elemento está agrietado.
  \item \textbf{Limitaciones declaradas} (\texttt{out\_of\_scope} en la spec): modo mixto I-II (fase G del ADR 0010), reorientación de \texttt{n} tras activación (tracking no trivial, fase F), múltiples discontinuidades por elemento, contacto unilateral en compresión, bulks no elásticos, orden superior (\texttt{Tri6\_Embedded}), 3D (\texttt{Tet4\_Embedded}).
  \item \textbf{Spec}: docs/specs/CST\_Embedded2D.md.
  \item \textbf{Archivo}: solidum/elements/solid\_2d/embedded\_cst.py.
\end{itemize}


\section{Cómo añadir un elemento nuevo}

\begin{enumerate}
  \item \textbf{Spec primero} --- el usuario crea \texttt{docs/specs/<Nombre>.md} a partir de \texttt{docs/specs/\_template\_element.md} (especificación física + formulación + contrato YAML). Sin spec, la IA no escribe código.
  \item \textbf{Scaffolding} --- \texttt{/solidum-new element <Nombre>} genera archivo en \texttt{solidum/elements/}, decorador \texttt{@ElementRegistry.register} y esqueleto de test.
  \item \textbf{Implementación + validación} --- la IA codifica contra la spec; los tests cubren los casos de \texttt{acceptance} declarados.
  \item \textbf{Catálogo} --- cuando la spec pasa a \texttt{status: validated}, se añade aquí una entrada breve siguiendo el formato de arriba (la spec sigue siendo la referencia detallada).
\end{enumerate}


\newpage

\chapter{Modelos constitutivos — catálogo}



\textit{Referencia rápida de los modelos constitutivos implementados. Una entrada por material. Para detalles del algoritmo \ensuremath{\to} código fuente.}

\textit{> \textbf{Convenciones}: \texttt{STRAIN\_DIM} = dimensión Voigt de la deformación de entrada (1 = escalar, 3 = 2D \texttt{[\ensuremath{\varepsilon}\_xx, \ensuremath{\varepsilon}\_yy, \ensuremath{\gamma}\_xy]}, 6 = 3D). \texttt{PRIMARY\_STATE\_VAR} = variable interna que el \texttt{VtkExporter} lleva al output.}

\textit{> \textbf{Densidad (ADR 0008)}: todos los materiales aceptan un parámetro opcional \texttt{density} (kg/m\ensuremath{^{3}} en las unidades del problema). No requerido al construir --- análisis estáticos sin peso propio funcionan sin declararla. Cuando un consumidor que requiere masa la pide (\texttt{Assembler.assemble\_self\_weight(g)} y futuras \texttt{assemble\_mass\_matrix}, etc.) y la encuentra como \texttt{None}, \textbf{falla con \texttt{ValueError}} listando los materiales sin densidad. Valor \texttt{0.0} declarado explícitamente cubre el caso legítimo de material sin masa física (penalty, restricción) y emite warning informativo.}



\section{Elastic1D --- elástico lineal 1D}

\begin{itemize}
  \item \textbf{Ley}: \texttt{\ensuremath{\sigma} = E \ensuremath{\cdot} \ensuremath{\varepsilon}} (Hooke 1D).
  \item \textbf{STRAIN\_DIM}: 1.
  \item \textbf{Parámetros}: \texttt{E} (módulo de Young).
  \item \textbf{Variables internas}: ninguna (sin historia).
  \item \textbf{Tangente}: constante = \texttt{E}.
  \item \textbf{Compatible con}: \texttt{Truss2D}, \texttt{Truss3D}, \texttt{Frame2DEuler}, \texttt{Frame2DTimoshenko}.
  \item \textbf{Spec}: docs/specs/Elastic1D.md
  \item \textbf{Archivo}: solidum/materials/elastic.py
\end{itemize}


\section{Elastic2D --- elástico lineal 2D}

\begin{itemize}
  \item \textbf{Ley}: \texttt{\ensuremath{\sigma} = C \ensuremath{\cdot} \ensuremath{\varepsilon}}, con \texttt{C} el tensor constitutivo isótropo en notación Voigt \texttt{[xx, yy, xy]}.
  \item \textbf{STRAIN\_DIM}: 3.
  \item \textbf{Parámetros}: \texttt{E}, \texttt{nu}, \texttt{hypothesis \ensuremath{\in} \{'plane\_stress', 'plane\_strain'\}}.
  \item \textbf{Variables internas}: ninguna.
  \item \textbf{Tangente}: constante = \texttt{C}.
  \item \textbf{Compatible con}: \texttt{Quad4}, \texttt{Tri3}, \texttt{Tri6}, \texttt{Quad8}, \texttt{Quad9} (todos los sólidos 2D con \texttt{STRAIN\_DIM = 3}).
  \item \textbf{Spec}: docs/specs/Elastic2D.md
  \item \textbf{Archivo}: solidum/materials/elastic\_2d.py
\end{itemize}


\section{Elastic3D --- elástico lineal isótropo 3D}

\begin{itemize}
  \item \textbf{Ley}: \texttt{\ensuremath{\sigma} = C \ensuremath{\cdot} \ensuremath{\varepsilon}}, con \texttt{C} el tensor constitutivo isótropo 6\ensuremath{\times}6 en notación Voigt 3D del proyecto \texttt{[xx, yy, zz, xy, yz, xz]} (ADR 0012, \texttt{Reglas.md §5}).
  \item \textbf{STRAIN\_DIM}: 6.
  \item \textbf{Parámetros}: \texttt{E} (>0), \texttt{nu \ensuremath{\in} (-1, 0.5)}, \texttt{density} (opcional, ADR 0008).
  \item \textbf{Variables internas}: ninguna.
  \item \textbf{Tangente}: constante = \texttt{C}.
  \item \textbf{Hipótesis}: no aplican \texttt{plane\_stress}/\texttt{plane\_strain} en 3D --- toda la deformación es activa.
  \item \textbf{Compatible con}: \texttt{Hex8}, \texttt{Tet4} (todos los sólidos 3D con \texttt{STRAIN\_DIM = 6}).
  \item \textbf{Caveat de compatibilidad}: ABAQUS usa orden Voigt \texttt{[11, 22, 33, 12, 13, 23]} (permutación \texttt{yz \ensuremath{\leftrightarrow} xz} respecto al proyecto). Importar datos de ABAQUS requiere intercambiar los componentes 5\ensuremath{\leftrightarrow}6 una vez en el preprocesador.
  \item \textbf{Spec}: docs/specs/Elastic3D.md
  \item \textbf{Archivo}: solidum/materials/elastic\_3d.py
\end{itemize}


\section{IsotropicDamage1D --- daño isótropo 1D con softening exponencial}

\begin{itemize}
  \item \textbf{Modelo}: daño escalar \texttt{d \ensuremath{\in} [0, DAMAGE\_MAX]}, esfuerzo \texttt{\ensuremath{\sigma} = (1 − d) \ensuremath{\cdot} E \ensuremath{\cdot} \ensuremath{\varepsilon}}. Versión 1D de \texttt{IsotropicDamage2D}.
  \item \textbf{Evolución del daño} (Kuhn-Tucker):
  \item Deformación equivalente: \texttt{\ensuremath{\varepsilon}\_eq = |\ensuremath{\varepsilon}|}.
  \item Variable histórica: \texttt{\ensuremath{\kappa} = max(\ensuremath{\kappa}\_old, \ensuremath{\varepsilon}\_eq)}.
  \item Si \texttt{\ensuremath{\kappa} \ensuremath{\le} \ensuremath{\kappa}\_0} \ensuremath{\to} \texttt{d = 0}. Si no: \texttt{d = 1 − (\ensuremath{\kappa}\_0/\ensuremath{\kappa}) \ensuremath{\cdot} exp(−\ensuremath{\alpha}\ensuremath{\cdot}(\ensuremath{\kappa} − \ensuremath{\kappa}\_0))}, saturado en \texttt{DAMAGE\_MAX}.
  \item \textbf{Tangente}: \textbf{algorítmica consistente} en carga activa con daño no saturado; \textbf{secante} \texttt{(1-d)\ensuremath{\cdot}E} en descarga, sin daño activo o al saturar.
  \item Fórmula consistente: \texttt{E\_tan = -(1-d)\ensuremath{\cdot}E\ensuremath{\cdot}\ensuremath{\alpha}\ensuremath{\cdot}\ensuremath{\kappa}} --- \textbf{negativa} durante toda la fase de softening (refleja la pendiente descendente \ensuremath{\sigma}-\ensuremath{\varepsilon} post-pico).
  \item Derivación: de \texttt{(1-d)\ensuremath{\cdot}E - E\ensuremath{\cdot}\ensuremath{\varepsilon}\ensuremath{\cdot}(\ensuremath{\partial}d/\ensuremath{\partial}\ensuremath{\kappa})\ensuremath{\cdot}sign(\ensuremath{\varepsilon})} con \texttt{\ensuremath{\partial}d/\ensuremath{\partial}\ensuremath{\kappa} = (1-d)(1/\ensuremath{\kappa} + \ensuremath{\alpha})} y \texttt{sign(\ensuremath{\varepsilon})\ensuremath{\cdot}\ensuremath{\varepsilon} = |\ensuremath{\varepsilon}| = \ensuremath{\kappa}} en carga.
  \item \textbf{STRAIN\_DIM}: 1 \ensuremath{\cdot} \textbf{PRIMARY\_STATE\_VAR}: \texttt{'damage'} \ensuremath{\cdot} \textbf{IS\_SYMMETRIC}: \texttt{True} (la tangente escalar produce contribución elemental simétrica aunque pueda ser indefinida; el despachador algebraico ADR 0003 detecta no-PD y degrada Cholesky\ensuremath{\to}LU automáticamente).
  \item \textbf{Parámetros}: \texttt{E}, \texttt{kappa\_0} (umbral elástico), \texttt{alpha} (velocidad de degradación), \texttt{density} (opcional, ADR 0008).
  \item \textbf{Variables internas}: \texttt{kappa}, \texttt{damage}.
  \item \textbf{Admisibilidad (ADR 0006)}: \texttt{admissibility\_scale = E \ensuremath{\cdot} \ensuremath{\kappa}\_0} --- esfuerzo umbral inicial. Garantiza invariancia bajo cambio de unidades del check \texttt{f \ensuremath{\le} atol + rtol \ensuremath{\cdot} escala}.
  \item \textbf{Limitaciones}: no distingue tracción/compresión en la activación del daño; control de carga global puede ser inestable tras el pico (requiere \texttt{ArcLengthSolver} para seguir la rama de softening).
  \item \textbf{Compatible con}: \texttt{Truss2D}, \texttt{Truss3D}.
  \item \textbf{Referencia}: ver \texttt{docs/specs/IsotropicDamage1D.md} y la hermana 2D para derivación detallada. Lemaitre \& Chaboche, \textit{Mechanics of Solid Materials}. Simó \& Ju (1987, IJSS 23) marco de tangente consistente.
  \item \textbf{Archivo}: solidum/materials/damage\_1d.py
\end{itemize}


\section{IsotropicDamage3D --- daño isótropo 3D con softening exponencial}

\begin{itemize}
  \item \textbf{Modelo}: extensión 3D de \texttt{IsotropicDamage2D}. Esfuerzo nominal \texttt{\ensuremath{\sigma} = (1 − d) \ensuremath{\cdot} C\_e \ensuremath{\cdot} \ensuremath{\varepsilon}} con \texttt{C\_e} 6\ensuremath{\times}6 isótropa 3D; daño escalar \texttt{d \ensuremath{\in} [0, DAMAGE\_MAX]}. \textbf{Sin variantes de hipótesis} --- en 3D todas las componentes son activas.
  \item \textbf{Deformación equivalente} (Voigt 6D del proyecto): \texttt{\ensuremath{\varepsilon}\_eq = \ensuremath{\surd}(\ensuremath{\varepsilon}\_xx\ensuremath{^{2}} + \ensuremath{\varepsilon}\_yy\ensuremath{^{2}} + \ensuremath{\varepsilon}\_zz\ensuremath{^{2}} + ½(\ensuremath{\gamma}\_xy\ensuremath{^{2}} + \ensuremath{\gamma}\_yz\ensuremath{^{2}} + \ensuremath{\gamma}\_xz\ensuremath{^{2}}))} con \texttt{M = diag(1, 1, 1, 1/2, 1/2, 1/2)}. Equivale a la norma de Frobenius del tensor deformación 3D y es la extensión natural de la convención 2D (\texttt{M\_2D = diag(1, 1, 1/2)}).
  \item \textbf{Evolución}: \ensuremath{\kappa} máxima histórica (Kuhn-Tucker \texttt{\ensuremath{\kappa}\_\{n+1\} = max(\ensuremath{\kappa}\_n, \ensuremath{\varepsilon}\_eq)}); ley exponencial centralizada en \texttt{solidum.materials.\_softening.evaluate\_exponential\_damage} (compartida con \texttt{IsotropicDamage1D} e \texttt{IsotropicDamage2D}); saturación a \texttt{DAMAGE\_MAX}.
  \item \textbf{Tangente}: \textbf{algorítmica consistente} en carga activa con daño no saturado; \textbf{secante} \texttt{(1-d)\ensuremath{\cdot}C\_e} en descarga, sin daño activo (\texttt{\ensuremath{\kappa} \ensuremath{\le} \ensuremath{\kappa}\_0}) o al saturar.
  \item Fórmula consistente: \texttt{C\_alg = (1-d)\ensuremath{\cdot}C\_e - [(1-d)\ensuremath{\cdot}(1/\ensuremath{\kappa} + \ensuremath{\alpha})/\ensuremath{\varepsilon}\_eq] \ensuremath{\cdot} (C\_e\ensuremath{\cdot}\ensuremath{\varepsilon}) ⊗ (M\ensuremath{\cdot}\ensuremath{\varepsilon})}.
  \item Derivación: idéntica al 2D extendida a 6D. \texttt{\ensuremath{\partial}d/\ensuremath{\partial}\ensuremath{\kappa} = (1-d)\ensuremath{\cdot}(1/\ensuremath{\kappa} + \ensuremath{\alpha})}; \texttt{\ensuremath{\partial}\ensuremath{\kappa}/\ensuremath{\partial}\ensuremath{\varepsilon} = M\ensuremath{\cdot}\ensuremath{\varepsilon}/\ensuremath{\varepsilon}\_eq} en carga activa, 0 en descarga.
  \item \textbf{No simétrica en general}: \texttt{(C\_e\ensuremath{\cdot}\ensuremath{\varepsilon})} y \texttt{(M\ensuremath{\cdot}\ensuremath{\varepsilon})} no son proporcionales salvo en estados de deformación muy particulares (uniaxial puro alineado con ejes). Recupera convergencia cuadrática del Newton global frente a la secante (Simó-Ju 1987).
  \item \textbf{Relación con \texttt{IsotropicDamage2D} plane\_strain}: bajo restricción \texttt{\ensuremath{\varepsilon}\_zz = \ensuremath{\gamma}\_yz = \ensuremath{\gamma}\_xz = 0}, Damage3D se reduce \textbf{exactamente} a Damage2D plane\_strain (misma \texttt{C\_e}, misma \texttt{\ensuremath{\varepsilon}\_eq}, misma ley centralizada). Test cruzado \texttt{TestIsotropicDamage3DvsPlaneStrain} blindado a \textbf{14 decimales en \texttt{d} y \texttt{\ensuremath{\kappa}}}, 10 decimales en \ensuremath{\sigma}. Plane stress 2D usa un \texttt{C\_e} proyectado distinto y no es comparable bajo restricción cinemática.
  \item \textbf{STRAIN\_DIM}: 6 \ensuremath{\cdot} \textbf{PRIMARY\_STATE\_VAR}: \texttt{'damage'} \ensuremath{\cdot} \textbf{IS\_SYMMETRIC}: \texttt{False} (la tangente consistente no es simétrica; el despachador algebraico ADR 0003 elige LU para el sistema global).
  \item \textbf{Parámetros}: \texttt{E} (>0), \texttt{nu} \ensuremath{\in} (-1, 0.5), \texttt{kappa\_0} (>0), \texttt{alpha} (>0), \texttt{density} (opcional, ADR 0008).
  \item \textbf{Variables internas}: \texttt{kappa}, \texttt{damage}. \textbf{Sin variables tensoriales} (a diferencia de los modelos plásticos): el daño es escalar, el historial es escalar.
  \item \textbf{Admisibilidad (ADR 0006)}: \texttt{admissibility\_scale = E \ensuremath{\cdot} \ensuremath{\kappa}\_0} --- idéntica al caso 1D/2D, constante respecto al estado.
  \item \textbf{Limitaciones} (idénticas al 2D): la \texttt{\ensuremath{\varepsilon}\_eq} simétrica no distingue daño en tensión vs compresión; sin regularización por longitud característica \ensuremath{\to} mesh-dependency en régimen de ablandamiento (localización en una banda de elementos). Para hormigón con resistencia diferente en tracción/compresión usar split tipo Mazars o de Vree (out-of-scope). Locking volumétrico en \texttt{Hex8}/\texttt{Tet4} con \texttt{\ensuremath{\nu}} cercano a 0.5 \ensuremath{\to} mismo régimen ya documentado para los demás materiales 3D, blindado por test del elemento.
  \item \textbf{Compatible con}: \texttt{Hex8}, \texttt{Tet4}.
  \item \textbf{Referencia}: ver \texttt{docs/specs/IsotropicDamage3D.md}. Simó \& Ju (1987, IJSS 23) tangente consistente para daño isótropo. Lemaitre \& Chaboche (1990) marco de continuum damage mechanics. de Souza Neto, Perić \& Owen (2008) §12.
  \item \textbf{Archivo}: solidum/materials/damage\_3d.py
\end{itemize}


\section{IsotropicDamage2D --- daño isótropo 2D con softening exponencial}

\begin{itemize}
  \item \textbf{Modelo}: extensión 2D de \texttt{IsotropicDamage1D}. Esfuerzo nominal \texttt{\ensuremath{\sigma} = (1 − d) \ensuremath{\cdot} C\_e \ensuremath{\cdot} \ensuremath{\varepsilon}}; daño escalar \texttt{d \ensuremath{\in} [0, DAMAGE\_MAX]}.
  \item \textbf{Deformación equivalente}: \texttt{\ensuremath{\varepsilon}\_eq = \ensuremath{\surd}(\ensuremath{\varepsilon}\_xx\ensuremath{^{2}} + \ensuremath{\varepsilon}\_yy\ensuremath{^{2}} + ½\ensuremath{\cdot}\ensuremath{\gamma}\_xy\ensuremath{^{2}})} (norma simple sin distinguir tensión/compresión).
  \item \textbf{Evolución}: \ensuremath{\kappa} máxima histórica (Kuhn-Tucker \texttt{\ensuremath{\kappa}\_n+1 = max(\ensuremath{\kappa}\_n, \ensuremath{\varepsilon}\_eq)}), ley exponencial \texttt{d(\ensuremath{\kappa}) = 1 - (\ensuremath{\kappa}\_0/\ensuremath{\kappa})\ensuremath{\cdot}exp(-\ensuremath{\alpha}(\ensuremath{\kappa}-\ensuremath{\kappa}\_0))} para \texttt{\ensuremath{\kappa} > \ensuremath{\kappa}\_0}, saturado a \texttt{DAMAGE\_MAX}.
  \item \textbf{Tangente}: \textbf{algorítmica consistente} en carga activa con daño no saturado; \textbf{secante} \texttt{(1-d)\ensuremath{\cdot}C\_e} en descarga, sin daño activo (\texttt{\ensuremath{\kappa} \ensuremath{\le} \ensuremath{\kappa}\_0}) o al saturar.
  \item Fórmula consistente: \texttt{C\_alg = (1-d)\ensuremath{\cdot}C\_e - [(1-d)\ensuremath{\cdot}(1/\ensuremath{\kappa} + \ensuremath{\alpha})/\ensuremath{\varepsilon}\_eq] \ensuremath{\cdot} (C\_e\ensuremath{\cdot}\ensuremath{\varepsilon}) ⊗ (M\ensuremath{\cdot}\ensuremath{\varepsilon})} con \texttt{M = diag(1, 1, 1/2)}.
  \item Derivación: \texttt{\ensuremath{\partial}d/\ensuremath{\partial}\ensuremath{\kappa} = (1-d)\ensuremath{\cdot}(1/\ensuremath{\kappa} + \ensuremath{\alpha})} por la ley exponencial; \texttt{\ensuremath{\partial}\ensuremath{\kappa}/\ensuremath{\partial}\ensuremath{\varepsilon} = M\ensuremath{\cdot}\ensuremath{\varepsilon}/\ensuremath{\varepsilon}\_eq} en carga activa, 0 en descarga.
  \item \textbf{No simétrica en general} (\texttt{(C\_e\ensuremath{\cdot}\ensuremath{\varepsilon})} y \texttt{(M\ensuremath{\cdot}\ensuremath{\varepsilon})} no son proporcionales). Recupera convergencia cuadrática del Newton global frente a la secante (que daba convergencia lineal).
  \item \textbf{STRAIN\_DIM}: 3 \ensuremath{\cdot} \textbf{PRIMARY\_STATE\_VAR}: \texttt{'damage'} \ensuremath{\cdot} \textbf{IS\_SYMMETRIC}: \texttt{False} (la tangente consistente no es simétrica; el despachador algebraico ADR 0003 elige LU en lugar de Cholesky para el sistema global).
  \item \textbf{Parámetros}: \texttt{E}, \texttt{nu}, \texttt{kappa\_0}, \texttt{alpha}, \texttt{hypothesis \ensuremath{\in} \{'plane\_stress' (default), 'plane\_strain'\}}, \texttt{density} (opcional, ADR 0008).
  \item \textbf{Variables internas}: \texttt{kappa}, \texttt{damage}.
  \item \textbf{Admisibilidad (ADR 0006)}: \texttt{admissibility\_scale = E \ensuremath{\cdot} \ensuremath{\kappa}\_0} --- idéntica al caso 1D, el criterio de daño tiene el mismo umbral característico independientemente de la dimensión.
  \item \textbf{Limitación}: la \texttt{\ensuremath{\varepsilon}\_eq} simétrica no distingue daño en tensión vs compresión; para hormigón usar modelo de Mazars o split tensión/compresión (out-of-scope). Sin regularización por longitud característica \ensuremath{\to} mesh-dependency en régimen de ablandamiento (banda localizada en una fila de elementos).
  \item \textbf{Compatible con}: \texttt{Quad4}, \texttt{Tri3}, \texttt{Tri6}, \texttt{Quad8}, \texttt{Quad9}.
  \item \textbf{Referencia}: ver \texttt{docs/specs/IsotropicDamage2D.md}. Simó \& Ju (1987, IJSS 23) tangente consistente para daño isótropo. Lemaitre \& Chaboche (1990) marco de continuum damage mechanics.
  \item \textbf{Archivo}: solidum/materials/damage\_2d.py
\end{itemize}


\section{Elastoplastic1D --- plasticidad J2 1D con endurecimiento isótropo lineal}

\begin{itemize}
  \item \textbf{Modelo}: plasticidad asociativa con criterio \texttt{f = |\ensuremath{\sigma}\_trial| − (\ensuremath{\sigma}\_y + H\ensuremath{\cdot}\ensuremath{\alpha}) \ensuremath{\le} 0}.
  \item \textbf{Algoritmo}: return mapping clásico 1D.
  \item Predictor elástico: \texttt{\ensuremath{\sigma}\_trial = E \ensuremath{\cdot} (\ensuremath{\varepsilon} − \ensuremath{\varepsilon}\_p\_old)}.
  \item Corrector: \texttt{\ensuremath{\Delta}\ensuremath{\gamma} = f / (E + H)}, \texttt{\ensuremath{\sigma} = \ensuremath{\sigma}\_trial − \ensuremath{\Delta}\ensuremath{\gamma}\ensuremath{\cdot}E\ensuremath{\cdot}sign(\ensuremath{\sigma}\_trial)}.
  \item Tangente algorítmica consistente: \texttt{E\_t = E\ensuremath{\cdot}H / (E + H)}.
  \item \textbf{STRAIN\_DIM}: 1 \ensuremath{\cdot} \textbf{PRIMARY\_STATE\_VAR}: \texttt{'alpha'}.
  \item \textbf{Parámetros}: \texttt{E}, \texttt{sigma\_y} (fluencia inicial), \texttt{H} (módulo de endurecimiento; 0 = perfectamente plástico).
  \item \textbf{Variables internas}: \texttt{eps\_p} (deformación plástica), \texttt{alpha} (acumulada equivalente).
  \item \textbf{Admisibilidad (ADR 0006)}: \texttt{admissibility\_scale = \ensuremath{\sigma}\_y + H \ensuremath{\cdot} \ensuremath{\alpha}} --- fluencia corriente (adaptativa al endurecimiento). El check \texttt{f \ensuremath{\le} atol + rtol \ensuremath{\cdot} escala} se hace contra la superficie de fluencia \textit{en el estado de entrada del paso}.
  \item \textbf{Compatible con}: \texttt{Truss2D}, \texttt{Truss3D}, \texttt{Frame2DEuler}, \texttt{Frame2DTimoshenko} (en estos últimos solo se aplica al esfuerzo axial $\sigma$).
  \item \textbf{Referencia}: Simo \& Hughes, \textit{Computational Inelasticity}, cap. 1.
  \item \textbf{Spec}: docs/specs/Elastoplastic1D.md
  \item \textbf{Archivo}: solidum/materials/plastic\_1d.py
\end{itemize}


\section{CableMaterial1D --- cable axial 1D con elasticidad unilateral}

\begin{itemize}
  \item \textbf{Ley}: \texttt{\ensuremath{\sigma} = E\ensuremath{\cdot}\ensuremath{\varepsilon}} si \texttt{\ensuremath{\varepsilon} > 0}; \texttt{\ensuremath{\sigma} = 0} si \texttt{\ensuremath{\varepsilon} \ensuremath{\le} 0} (sin rigidez en compresión).
  \item \textbf{STRAIN\_DIM}: 1.
  \item \textbf{Parámetros}: \texttt{E} (módulo de Young en tensión, estrictamente positivo).
  \item \textbf{Variables internas}: ninguna (respuesta memoryless: la rigidez depende solo del signo instantáneo de \texttt{\ensuremath{\varepsilon}}).
  \item \textbf{Tangente}: \texttt{E\_t = E} en tensión, \texttt{E\_t = 0} en compresión o sin tensión.
  \item \textbf{Implicación numérica}: la rigidez tangente colapsa a cero cuando el cable se afloja, lo que requiere precondicionamiento o regularización en el solver para no degenerar la matriz global. El uso típico es a través del elemento \texttt{Cable2DCorot}/\texttt{Cable3DCorot}, que detecta la situación y la maneja correctamente.
  \item \textbf{Compatible con}: \texttt{Cable2DCorot}, \texttt{Cable3DCorot}.
  \item \textbf{Referencia}: ver \texttt{docs/specs/CableMaterial1D.md}.
  \item \textbf{Archivo}: solidum/materials/cable\_1d.py
\end{itemize}


\section{DruckerPrager3D --- plasticidad friccional cohesivo-friccional 3D}

\begin{itemize}
  \item \textbf{Modelo}: cono circular suave de Mohr-Coulomb 3D con cohesión y fricción interna. Criterio \texttt{f = \ensuremath{\surd}J\ensuremath{_{2}} + \ensuremath{\eta}\_f\ensuremath{\cdot}I\ensuremath{_{1}} − k(\ensuremath{\alpha}) \ensuremath{\le} 0} con \texttt{k(\ensuremath{\alpha}) = k\_0 + H\ensuremath{\cdot}\ensuremath{\alpha}} (endurecimiento isótropo lineal en cohesión). \textbf{Sin variantes de hipótesis} --- en 3D todas las componentes son activas.
  \item \textbf{Calibraciones disponibles} (3D):
  \item \texttt{outer\_cone} (default): \ensuremath{\eta}\_f = 2\ensuremath{\cdot}sin(\ensuremath{\varphi})/[\ensuremath{\surd}3\ensuremath{\cdot}(3 - sin(\ensuremath{\varphi}))], k\_0 = 6\ensuremath{\cdot}c\_0\ensuremath{\cdot}cos(\ensuremath{\varphi})/[\ensuremath{\surd}3\ensuremath{\cdot}(3 - sin(\ensuremath{\varphi}))]. Circunscribe Mohr-Coulomb en el meridiano de compresión triaxial.
  \item \texttt{inner\_cone}: \ensuremath{\eta}\_f = 2\ensuremath{\cdot}sin(\ensuremath{\varphi})/[\ensuremath{\surd}3\ensuremath{\cdot}(3 + sin(\ensuremath{\varphi}))], k\_0 = 6\ensuremath{\cdot}c\_0\ensuremath{\cdot}cos(\ensuremath{\varphi})/[\ensuremath{\surd}3\ensuremath{\cdot}(3 + sin(\ensuremath{\varphi}))]. Inscribe MC en el meridiano de extensión triaxial.
  \item \textbf{No incluye \texttt{plane\_strain\_matched}} (existente en DP2D): esa calibración es 2D-only por construcción --- MC en 3D depende del ángulo de Lode y no admite coincidencia circular global. El constructor rechaza explícitamente esa variante con mensaje claro.
  \item \textbf{Plasticidad no asociada por defecto}: ángulo de dilatancia \texttt{\ensuremath{\psi}} parámetro independiente, default \texttt{\ensuremath{\psi} = \ensuremath{\varphi}} (asociada). Para \texttt{\ensuremath{\psi} \ensuremath{\neq} \ensuremath{\varphi}} la tangente algorítmica es \textbf{asimétrica} \ensuremath{\to} el despachador algebraico ADR 0003 elige LU.
  \item \textbf{Algoritmo --- dos ramas cerradas con detección automática}:
  \item \textbf{Return regular} (cone surface): \texttt{\ensuremath{\Delta}\ensuremath{\gamma} = f\_trial / (G + 9K\ensuremath{\cdot}\ensuremath{\eta}\_f\ensuremath{\cdot}\ensuremath{\eta}\_g + H\_k)}, dirección desviadora invariante bajo carga radial. Tangente algorítmica \texttt{C\_alg = K\ensuremath{\cdot}v⊗v + 2G\ensuremath{\cdot}(1-\ensuremath{\beta})\ensuremath{\cdot}I\_dev + 4G\ensuremath{\cdot}\ensuremath{\beta}\ensuremath{\cdot}n̂⊗n̂ - (1/A)\ensuremath{\cdot}b\_g⊗b\_f} con \texttt{n̂} y \texttt{b\_f, b\_g} en Voigt 6D tensorial. Idéntica estructura algebraica al kernel 2D, solo extendida a 6 componentes activas.
  \item \textbf{Return al ápice} (vértice hidrostático): \texttt{\ensuremath{\Delta}\ensuremath{\gamma}\_apex = (I\_1\_trial\ensuremath{\cdot}\ensuremath{\eta}\_f - k(\ensuremath{\alpha}))/(9K\ensuremath{\cdot}\ensuremath{\eta}\_f\ensuremath{\cdot}\ensuremath{\eta}\_g + H\_k)}, \ensuremath{\sigma}\_n+1 = (k/(3\ensuremath{\eta}\_f))\ensuremath{\cdot}I puramente hidrostático en las tres componentes diagonales. Tangente reducida \texttt{(K\ensuremath{\cdot}H\_k/(9K\ensuremath{\cdot}\ensuremath{\eta}\_f\ensuremath{\cdot}\ensuremath{\eta}\_g + H\_k))\ensuremath{\cdot}v⊗v} (sin rigidez desviadora; salvaguarda numérica \texttt{K\ensuremath{\cdot}1e-6\ensuremath{\cdot}v⊗v} si la rigidez colapsa con H\_k=0 y \ensuremath{\eta}\_g\ensuremath{\approx}0).
  \item Detección regular \ensuremath{\leftrightarrow} apex: si tras \texttt{\ensuremath{\Delta}\ensuremath{\gamma}\_regular} resulta \texttt{\ensuremath{\surd}J\ensuremath{_{2}}\_new < 0} \ensuremath{\to} cae en apex; cambia a la fórmula apex.
  \item \textbf{Relación con \texttt{DruckerPrager2D} plane\_strain}: DP2D plane\_strain con \texttt{variant='outer\_cone'} (o \texttt{'inner\_cone'}) es la restricción de DP3D bajo \texttt{\ensuremath{\varepsilon}\_zz = \ensuremath{\gamma}\_yz = \ensuremath{\gamma}\_xz = 0}. Test cruzado \texttt{TestDruckerPrager3DvsPlaneStrain} blindado a 10 decimales sobre \texttt{\ensuremath{\sigma}\_xx, \ensuremath{\sigma}\_yy, \ensuremath{\sigma}\_xy, \ensuremath{\alpha}} y componentes plásticas planas; las componentes 3D ausentes en plane strain (\texttt{\ensuremath{\varepsilon}\textasciicircum{}p\_yz, \ensuremath{\varepsilon}\textasciicircum{}p\_xz}) permanecen nulas en VM3D.
  \item \textbf{STRAIN\_DIM}: 6 \ensuremath{\cdot} \textbf{PRIMARY\_STATE\_VAR}: \texttt{'alpha'} \ensuremath{\cdot} \textbf{IS\_SYMMETRIC}: \texttt{False} declarativo (override por instancia a \texttt{True} cuando \texttt{\ensuremath{\psi} = \ensuremath{\varphi}}, asociada).
  \item \textbf{Parámetros} (físicos):
  \item \texttt{E}, \texttt{nu}: elásticos (\texttt{\ensuremath{\nu} \ensuremath{\in} (-1, 0.5)} estricto).
  \item \texttt{cohesion}: cohesión inicial \texttt{c\_0} (esfuerzo, >0).
  \item \texttt{phi\_deg}: ángulo de fricción interna (grados, \texttt{0 \ensuremath{\le} \ensuremath{\varphi} < 90}).
  \item \texttt{psi\_deg}: ángulo de dilatancia (grados, \texttt{0 \ensuremath{\le} \ensuremath{\psi} \ensuremath{\le} \ensuremath{\varphi}}; default \texttt{\ensuremath{\varphi}} \ensuremath{\Rightarrow} asociada).
  \item \texttt{H} (\ensuremath{\ge}0, default 0): endurecimiento isótropo lineal en cohesión.
  \item \texttt{variant}: \texttt{'outer\_cone'} (default) o \texttt{'inner\_cone'}.
  \item \texttt{density} (opcional, ADR 0008).
  \item \textbf{Variables internas}: \texttt{eps\_p} (6 componentes Voigt 6D con cortantes tensoriales \texttt{[xx, yy, zz, xy, yz, xz]}), \texttt{alpha} (multiplicador plástico acumulado, adimensional). \textbf{Invariante cinemático}: \texttt{tr(\ensuremath{\varepsilon}\textasciicircum{}p) = 3\ensuremath{\cdot}\ensuremath{\eta}\_g\ensuremath{\cdot}\ensuremath{\alpha}} en cualquier estado plástico (distinto de J2 donde \texttt{tr(\ensuremath{\varepsilon}\textasciicircum{}p) = 0}).
  \item \textbf{Admisibilidad (ADR 0006)}: \texttt{admissibility\_scale = k(\ensuremath{\alpha}) = k\_0 + H\ensuremath{\cdot}\ensuremath{\alpha}}.
  \item \textbf{Limitaciones declaradas} (\texttt{out\_of\_scope} en spec): MC con aristas (pirámide hexagonal --- DP no captura efectos de Lode angle, diferencia hasta 15-20\% en frontera fuera de los meridianos puros), variantes con dependencia explícita del Lode angle (Matsuoka-Nakai, Lade-Duncan), endurecimiento por fricción/dilatancia, endurecimiento cinemático, cap-cone, ablandamiento (\texttt{H<0} requiere regularización), grandes deformaciones. Locking volumétrico en \texttt{Hex8}/\texttt{Tet4} cuando \texttt{\ensuremath{\psi} > 0} (flujo dilatante) + \texttt{\ensuremath{\nu}} moderado-alto --- mismo régimen ya documentado para \texttt{Elastic3D}/\texttt{VonMises3D}, sin mitigación implementada.
  \item \textbf{Compatible con}: \texttt{Hex8}, \texttt{Tet4}.
  \item \textbf{Referencia}: ver \texttt{docs/specs/DruckerPrager3D.md}. Drucker \& Prager (1952). Simó \& Hughes (1998) §6 (cone plasticity, apex return). de Souza Neto, Perić \& Owen (2008) §8 (calibraciones outer/inner 3D, tangentes algorítmicas detalladas).
  \item \textbf{Archivo}: solidum/materials/drucker\_prager\_3d.py
\end{itemize}


\section{DruckerPrager2D --- plasticidad friccional cohesivo-friccional 2D plane strain}

\begin{itemize}
  \item \textbf{Modelo}: cono circular suave de Mohr-Coulomb con cohesión y fricción interna. Criterio \texttt{f = \ensuremath{\surd}J\ensuremath{_{2}} + \ensuremath{\eta}\_f\ensuremath{\cdot}I\ensuremath{_{1}} − k(\ensuremath{\alpha}) \ensuremath{\le} 0} con \texttt{k(\ensuremath{\alpha}) = k\_0 + H\ensuremath{\cdot}\ensuremath{\alpha}} (endurecimiento isótropo lineal en cohesión).
  \item \textbf{Hipótesis}: solo \texttt{plane\_strain} en esta entrega. \texttt{plane\_stress} declarado \textit{out-of-scope} (proyección con \ensuremath{\sigma}\_zz=0 acoplada a flujo dilatante es notoriamente delicada).
  \item \textbf{Plasticidad no asociada por defecto}: ángulo de dilatancia \texttt{\ensuremath{\psi}} parámetro independiente, default \texttt{\ensuremath{\psi} = \ensuremath{\varphi}} (asociada). Para \texttt{\ensuremath{\psi} \ensuremath{\neq} \ensuremath{\varphi}} la tangente algorítmica es \textbf{asimétrica} \ensuremath{\to} \texttt{IS\_SYMMETRIC = False}.
  \item \textbf{Calibración con Mohr-Coulomb} (variant):
  \item \texttt{plane\_strain\_matched} (default): \ensuremath{\eta}\_f = tan(\ensuremath{\varphi})/\ensuremath{\surd}(9 + 12tan\ensuremath{^{2}}(\ensuremath{\varphi})), k\_0 = 3c\_0/\ensuremath{\surd}(9 + 12tan\ensuremath{^{2}}(\ensuremath{\varphi})). Coincide exactamente con MC en plane strain.
  \item \texttt{outer\_cone}: \ensuremath{\eta}\_f = 2\ensuremath{\cdot}sin(\ensuremath{\varphi})/[\ensuremath{\surd}3\ensuremath{\cdot}(3 - sin(\ensuremath{\varphi}))]. Circunscribe MC.
  \item \texttt{inner\_cone}: \ensuremath{\eta}\_f = 2\ensuremath{\cdot}sin(\ensuremath{\varphi})/[\ensuremath{\surd}3\ensuremath{\cdot}(3 + sin(\ensuremath{\varphi}))]. Inscribe MC.
  \item \textbf{Algoritmo --- dos ramas cerradas con detección automática}:
  \item \textbf{Return regular} (cone surface): \texttt{\ensuremath{\Delta}\ensuremath{\gamma} = f\_trial / (G + 9K\ensuremath{\cdot}\ensuremath{\eta}\_f\ensuremath{\cdot}\ensuremath{\eta}\_g + H)}, dirección desviadora preservada. Tangente algorítmica \texttt{C\_alg = K\ensuremath{\cdot}v⊗v + 2G\ensuremath{\cdot}(1-\ensuremath{\beta})\ensuremath{\cdot}I\_dev + 4G\ensuremath{\cdot}\ensuremath{\beta}\ensuremath{\cdot}n̂⊗n̂ - (1/A)\ensuremath{\cdot}b\_g⊗b\_f} con \texttt{\ensuremath{\beta} = G\ensuremath{\cdot}\ensuremath{\Delta}\ensuremath{\gamma}/\ensuremath{\surd}J\ensuremath{_{2}}\_trial}, \texttt{b\_g = 2G\ensuremath{\cdot}n̂ + 3K\ensuremath{\cdot}\ensuremath{\eta}\_g\ensuremath{\cdot}v}, \texttt{b\_f = 2G\ensuremath{\cdot}n̂ + 3K\ensuremath{\cdot}\ensuremath{\eta}\_f\ensuremath{\cdot}v}.
  \item \textbf{Return al ápice} (vértice del cono, zona puramente hidrostática): \texttt{\ensuremath{\Delta}\ensuremath{\gamma}\_apex = (I\ensuremath{_{1}}\_trial\ensuremath{\cdot}\ensuremath{\eta}\_f - k(\ensuremath{\alpha}))/(9K\ensuremath{\cdot}\ensuremath{\eta}\_f\ensuremath{\cdot}\ensuremath{\eta}\_g + H)}, \ensuremath{\sigma}\_n+1 = (k/3\ensuremath{\eta}\_f)\ensuremath{\cdot}I puramente hidrostático. Tangente reducida a \texttt{K\ensuremath{\cdot}H/(9K\ensuremath{\cdot}\ensuremath{\eta}\_f\ensuremath{\cdot}\ensuremath{\eta}\_g + H) \ensuremath{\cdot} v⊗v} (simétrica, sin rigidez desviadora).
  \item Detección regular \ensuremath{\leftrightarrow} apex: si tras \texttt{\ensuremath{\Delta}\ensuremath{\gamma}\_regular} resulta \texttt{\ensuremath{\surd}J\ensuremath{_{2}}\_new < 0} \ensuremath{\to} cae en apex; cambia a la fórmula apex.
  \item \textbf{STRAIN\_DIM}: 3 \ensuremath{\cdot} \textbf{PRIMARY\_STATE\_VAR}: \texttt{'alpha'} \ensuremath{\cdot} \textbf{IS\_SYMMETRIC}: \texttt{False} (declarativo, conservador; el despachador algebraico ADR 0003 elige LU).
  \item \textbf{Parámetros} (físicos):
  \item \texttt{E}, \texttt{nu}: elásticos.
  \item \texttt{cohesion}: cohesión inicial \texttt{c\_0} (esfuerzo).
  \item \texttt{phi\_deg}: ángulo de fricción interna (grados, \texttt{0 \ensuremath{\le} \ensuremath{\varphi} < 90}).
  \item \texttt{psi\_deg}: ángulo de dilatancia (grados, \texttt{0 \ensuremath{\le} \ensuremath{\psi} \ensuremath{\le} \ensuremath{\varphi}}; default \texttt{\ensuremath{\varphi}} \ensuremath{\Rightarrow} asociada).
  \item \texttt{H} (\ensuremath{\ge}0, default 0): endurecimiento isótropo lineal en cohesión.
  \item \texttt{variant}: calibración con MC (default \texttt{'plane\_strain\_matched'}).
  \item \texttt{density} (opcional, ADR 0008).
  \item \textbf{Variables internas}: \texttt{eps\_p} (tensorial 4 componentes \texttt{[xx, yy, zz, xy\_tensorial]}), \texttt{alpha} (multiplicador plástico acumulado, adimensional).
  \item \textbf{Admisibilidad (ADR 0006)}: \texttt{admissibility\_scale = k(\ensuremath{\alpha}) = k\_0 + H\ensuremath{\cdot}\ensuremath{\alpha}} (cohesión efectiva corriente, unidades de esfuerzo igual que \texttt{f}).
  \item \textbf{Limitaciones} (out-of-scope):
  \item Sin \texttt{plane\_stress}, endurecimiento por fricción/dilatancia, endurecimiento cinemático, cap (modelo cap-cone), ablandamiento (\texttt{H<0} requiere regularización), grandes deformaciones.
  \item Validación inicial cubierta por tests unitarios (incluyendo FD numérica de la tangente) y un benchmark de pipeline Quad4 (rama elástica + rama apex + caso asociado). El régimen "rama regular pura" en geometría confinada es difícil de calibrar sin caer en apex; se cubre por los tests unitarios de cortante puro.
  \item \textbf{Compatible con}: \texttt{Quad4}, \texttt{Tri3}, \texttt{Tri6}, \texttt{Quad8}, \texttt{Quad9} (todos con material 2D plane strain).
  \item \textbf{Referencia}: ver \texttt{docs/specs/DruckerPrager2D.md}. Drucker \& Prager (1952). de Souza Neto, Perić \& Owen (2008) cap. 8 (Drucker-Prager, tangentes algorítmicas).
  \item \textbf{Archivo}: solidum/materials/drucker\_prager\_2d.py
\end{itemize}


\section{VonMises3D --- plasticidad J2 3D con endurecimiento isótropo lineal}

\begin{itemize}
  \item \textbf{Modelo}: J2 (Von Mises) con regla de flujo asociada y endurecimiento isótropo lineal, formulado íntegramente en Voigt 6D del proyecto (ADR 0012). \textbf{Sin variantes de hipótesis} --- en 3D todas las componentes de \texttt{\ensuremath{\varepsilon}} y \texttt{\ensuremath{\sigma}} son activas.
  \item \textbf{Algoritmo} (Simó-Hughes §3.3): return mapping radial cerrado único sobre la parte desviadora 6D. Predictor elástico \texttt{s\_trial = 2G\ensuremath{\cdot}(e\_dev − e\_p)} con presión \texttt{p = K\ensuremath{\cdot}tr(\ensuremath{\varepsilon})}; criterio \texttt{f = \ensuremath{\|}s\ensuremath{\|} − \ensuremath{\surd}(2/3)\ensuremath{\cdot}(\ensuremath{\sigma}\_y + H\ensuremath{\cdot}\ensuremath{\alpha}) \ensuremath{\le} 0} con norma Frobenius tensorial \texttt{\ensuremath{\|}s\ensuremath{\|}\ensuremath{^{2}} = \ensuremath{\Sigma}\_diag s\_ii\ensuremath{^{2}} + 2\ensuremath{\cdot}\ensuremath{\Sigma}\_off s\_ij\ensuremath{^{2}}}; corrector radial \texttt{\ensuremath{\Delta}\ensuremath{\gamma} = f\_trial / (2G + ⅔\ensuremath{\cdot}H)} con \texttt{N = s\_trial/\ensuremath{\|}s\_trial\ensuremath{\|}}; actualización \texttt{\ensuremath{\alpha}\_new = \ensuremath{\alpha} + \ensuremath{\surd}(2/3)\ensuremath{\cdot}\ensuremath{\Delta}\ensuremath{\gamma}}. Tangente algorítmica consistente cerrada \texttt{C\_alg = K\ensuremath{\cdot}v⊗v + 2G(1-\ensuremath{\beta})\ensuremath{\cdot}I\_dev - 2G\ensuremath{\cdot}\ensuremath{\gamma}̄\ensuremath{\cdot}N⊗N}. Idéntica estructura al kernel \texttt{plane\_strain} de \texttt{VonMises2D} ampliada a 6 componentes activas; comparte el principio físico pero el código kernel se mantiene separado por simplicidad (zero-padding sistemático en Numba sería contraproducente).
  \item \textbf{Relación con \texttt{VonMises2D} plane strain}: VM2D plane strain es la restricción de este modelo bajo \texttt{\ensuremath{\varepsilon}\_zz = \ensuremath{\gamma}\_yz = \ensuremath{\gamma}\_xz = 0} (deformación total impuesta) con \texttt{\ensuremath{\varepsilon}\textasciicircum{}p\_zz, \ensuremath{\varepsilon}\textasciicircum{}p\_yz, \ensuremath{\varepsilon}\textasciicircum{}p\_xz} libres por incompresibilidad/isotropía. Test cruzado \texttt{TestVonMises3DvsPlaneStrain} blindado a 10 decimales sobre \texttt{\ensuremath{\sigma}\_xx, \ensuremath{\sigma}\_yy, \ensuremath{\sigma}\_xy, \ensuremath{\alpha}} y componentes plásticas planas.
  \item \textbf{STRAIN\_DIM}: 6 \ensuremath{\cdot} \textbf{PRIMARY\_STATE\_VAR}: \texttt{'alpha'} \ensuremath{\cdot} \textbf{IS\_SYMMETRIC}: \texttt{True} (J2 asociado).
  \item \textbf{Parámetros}: \texttt{E} (>0), \texttt{nu \ensuremath{\in} (-1, 0.5)}, \texttt{sigma\_y} (>0), \texttt{H} (\ensuremath{\ge}0, default 0), \texttt{density} (opcional, ADR 0008).
  \item \textbf{Variables internas}: \texttt{eps\_p} (6 componentes Voigt 6D con cortantes tensoriales \texttt{[xx, yy, zz, xy, yz, xz]}), \texttt{alpha} (acumulada equivalente, adimensional). La parte plástica es \textbf{incompresible exacta}: \texttt{tr(\ensuremath{\varepsilon}\textasciicircum{}p) = 0} por construcción del flujo desviador.
  \item \textbf{Admisibilidad (ADR 0006)}: \texttt{admissibility\_scale = \ensuremath{\surd}(2/3)\ensuremath{\cdot}(\ensuremath{\sigma}\_y + H\ensuremath{\cdot}\ensuremath{\alpha})} --- idéntica a plane strain (misma norma desviadora en la frontera J2).
  \item \textbf{Limitaciones declaradas} (\texttt{out\_of\_scope} en spec): no endurecimiento cinemático/Voce, no ablandamiento (\texttt{H<0} requiere regularización), no regla de flujo no asociada, no daño acoplado, no viscoplasticidad, no grandes deformaciones, no anisotropía. Locking volumétrico en \texttt{Hex8}/\texttt{Tet4} cuando la plasticidad domina + \texttt{\ensuremath{\nu}} moderado-alto --- mismo régimen ya documentado para \texttt{Elastic3D}, blindado por tests, sin mitigación implementada (B-bar/F-bar diferidas).
  \item \textbf{Implementación}: kernel \texttt{\_compute\_j2\_3d} con \texttt{@njit}. Sin despacho por hipótesis (no aplica en 3D).
  \item \textbf{Compatible con}: \texttt{Hex8}, \texttt{Tet4} (todos los sólidos 3D con \texttt{STRAIN\_DIM = 6}).
  \item \textbf{Referencia}: ver \texttt{docs/specs/VonMises3D.md}. Simó \& Hughes, \textit{Computational Inelasticity} (1998), §3.3 (3D J2 return mapping radial cerrado). de Souza Neto, Perić \& Owen, \textit{Computational Methods for Plasticity} (2008), §7.
  \item \textbf{Archivo}: solidum/materials/von\_mises\_3d.py
\end{itemize}


\section{VonMises2D --- plasticidad J2 2D con endurecimiento isótropo lineal}

\begin{itemize}
  \item \textbf{Modelo}: J2 (Von Mises) con regla de flujo asociada y endurecimiento isótropo lineal. Soporta \textbf{dos hipótesis cinemáticas} mutuamente excluyentes (selección por kwarg \texttt{hypothesis} al construir): \texttt{plane\_strain} (\texttt{\ensuremath{\varepsilon}\_zz = 0}) y \texttt{plane\_stress} (\texttt{\ensuremath{\sigma}\_zz = 0}). Cada hipótesis usa un kernel Numba especializado.
  \item \textbf{Algoritmo plane strain} (Simó-Hughes §3.3): return mapping radial cerrado sobre la parte desviadora 3D extendida. Predictor elástico desviador \texttt{s\_trial = 2G\ensuremath{\cdot}(e\_dev − e\_p)} con presión \texttt{p = K\ensuremath{\cdot}tr(\ensuremath{\varepsilon})}; criterio \texttt{f = \ensuremath{\|}s\_trial\ensuremath{\|} − \ensuremath{\surd}(2/3)\ensuremath{\cdot}(\ensuremath{\sigma}\_y + H\ensuremath{\cdot}\ensuremath{\alpha}) \ensuremath{\le} 0}; corrector radial \texttt{\ensuremath{\Delta}\ensuremath{\gamma} = f\_trial / (2G + ⅔\ensuremath{\cdot}H)} con \texttt{N = s\_trial/\ensuremath{\|}s\_trial\ensuremath{\|}}; actualización \texttt{\ensuremath{\alpha}\_new = \ensuremath{\alpha} + \ensuremath{\surd}(2/3)\ensuremath{\cdot}\ensuremath{\Delta}\ensuremath{\gamma}}. Tangente algorítmica consistente cerrada.
  \item \textbf{Algoritmo plane stress} (Simó-Hughes §3.4.1, "plane stress projected algorithm"): predictor \texttt{\ensuremath{\sigma}\_trial = C\_e\_ps \ensuremath{\cdot} (\ensuremath{\varepsilon} − e\_p)} en subespacio plane stress; función de fluencia proyectada \texttt{f̄ = ½\ensuremath{\cdot}\ensuremath{\sigma}\ensuremath{\cdot}P\ensuremath{\cdot}\ensuremath{\sigma} − R\ensuremath{^{2}}/3} con operador \texttt{P = (1/3)\ensuremath{\cdot}[[2,-1,0],[-1,2,0],[0,0,6]]}. Newton local escalar sobre \texttt{\ensuremath{\Delta}\ensuremath{\gamma}} en base de autovectores ortonormales \texttt{v\_1=[1,1,0]/\ensuremath{\surd}2}, \texttt{v\_2=[1,-1,0]/\ensuremath{\surd}2}, \texttt{v\_3=[0,0,1]} de \texttt{C\_e\_ps\ensuremath{\cdot}P} con autovalores \texttt{\ensuremath{\mu}\_1=E/(3(1-\ensuremath{\nu}))}, \texttt{\ensuremath{\mu}\_2=\ensuremath{\mu}\_3=2G}. Converge en 3-6 iteraciones (tangente cerrada). Actualización de \texttt{\ensuremath{\alpha}} con la regla físicamente correcta \texttt{\ensuremath{\alpha}\_new = \ensuremath{\alpha} + \ensuremath{\Delta}\ensuremath{\gamma}\ensuremath{\cdot}\ensuremath{\surd}(2\ensuremath{\sigma}P\ensuremath{\sigma}/3)} (la heurística \texttt{\ensuremath{\surd}(2/3)\ensuremath{\cdot}\ensuremath{\Delta}\ensuremath{\gamma}} válida en plane strain rompe la invariancia bajo cambio de unidades en plane stress, donde \texttt{\ensuremath{\Delta}\ensuremath{\gamma}} tiene unidades \texttt{[1/esfuerzo]}). Incompresibilidad plástica cierra \texttt{e\_p\_zz = -(e\_p\_xx + e\_p\_yy)}. Tangente algorítmica consistente cerrada con corrección por \texttt{d\ensuremath{\alpha}/d\ensuremath{\Delta}\ensuremath{\gamma} \ensuremath{\neq} \ensuremath{\surd}(2/3)}.
  \item \textbf{STRAIN\_DIM}: 3 \ensuremath{\cdot} \textbf{PRIMARY\_STATE\_VAR}: \texttt{'alpha'}.
  \item \textbf{Parámetros}: \texttt{E}, \texttt{nu}, \texttt{sigma\_y}, \texttt{H} (\ensuremath{\ge}0, default 0), \texttt{hypothesis \ensuremath{\in} \{'plane\_strain' (default), 'plane\_stress'\}}, \texttt{density} (opcional, ADR 0008).
  \item \textbf{Variables internas}: \texttt{eps\_p} (tensorial 4 componentes \texttt{[xx, yy, zz, xy\_tensorial]}), \texttt{alpha} (acumulada equivalente, adimensional en ambas hipótesis).
  \item \textbf{Admisibilidad (ADR 0006)}:
  \item plane\_strain: \texttt{admissibility\_scale = \ensuremath{\surd}(2/3)\ensuremath{\cdot}(\ensuremath{\sigma}\_y + H\ensuremath{\cdot}\ensuremath{\alpha})} --- radio desviador de la superficie J2.
  \item plane\_stress: \texttt{admissibility\_scale = (\ensuremath{\sigma}\_y + H\ensuremath{\cdot}\ensuremath{\alpha})\ensuremath{^{2}}/3} --- escala de \texttt{f̄} proyectada.
  \item En ambos, la tolerancia se precomputa fuera del kernel \texttt{@njit} y se pasa como \texttt{float}.
  \item \textbf{Limitaciones declaradas} (ver \texttt{out\_of\_scope} en spec): no endurecimiento cinemático/Voce, no ablandamiento (\texttt{H<0} requiere regularización), no regla de flujo no asociada, no daño acoplado, no viscoplasticidad, no grandes deformaciones, no anisotropía. Locking volumétrico en plane strain con elementos de bajo orden cuando \texttt{\ensuremath{\nu} \ensuremath{\to} 0.5} o dominio plástico (mitigaciones B-bar/F-bar diferidas).
  \item \textbf{Implementación}: kernels \texttt{\_compute\_j2\_plane\_strain} y \texttt{\_compute\_j2\_plane\_stress} con \texttt{@njit}. Despacho por \texttt{hypothesis} en construcción (sin coste runtime).
  \item \textbf{Compatible con}: \texttt{Quad4}, \texttt{Tri3}, \texttt{Tri6}, \texttt{Quad8}, \texttt{Quad9} (configurados en cualquiera de las dos hipótesis).
  \item \textbf{Referencia}: ver \texttt{docs/specs/VonMises2D.md}. Simó \& Hughes, \textit{Computational Inelasticity} (1998), §3.3 (plane strain), §3.4 (plane stress projected, Box 3.1). de Souza Neto, Perić \& Owen, \textit{Computational Methods for Plasticity} (2008), §9.4.
  \item \textbf{Archivo}: solidum/materials/von\_mises\_2d.py
\end{itemize}



Los materiales cohesivos son una \textbf{jerarquía paralela e independiente} de los materiales continuos: operan sobre el salto de desplazamientos \texttt{[[u]]} sobre una superficie de discontinuidad \texttt{Γ\_d} y devuelven tracciones \texttt{t}, no esfuerzos sobre \texttt{\ensuremath{\varepsilon}}. Viven en \texttt{solidum/cohesive\_materials/}, heredan de \texttt{solidum.core.cohesive\_material.CohesiveMaterial}, se registran vía \texttt{CohesiveMaterialRegistry} y se declaran en YAML bajo la sección \texttt{cohesive\_materials} (separada de \texttt{materials}). El bulk del elemento sigue siendo un \texttt{Material} continuo; el cohesivo entra al sistema sólo cuando el elemento activa una discontinuidad embebida (ver ADR 0010).

\textit{\textbf{Convenciones de la familia}: \texttt{JUMP\_DIM} = dimensión del vector de salto (2 en 2D, 3 en 3D); \texttt{PRIMARY\_STATE\_VAR} = variable interna que se exporta al post-proceso; \texttt{IS\_SYMMETRIC} = simetría de la contribución a la tangente del elemento. Parámetros físicos típicos: \texttt{sigma\_t0} (Pa), \texttt{G\_f} (N/m), \texttt{K\_e} (Pa/m). No tienen densidad --- la inercia es del bulk.}



\section{CohesiveDamageIsotropic --- daño cohesivo isótropo Modo-I con softening lineal/exponencial}

\begin{itemize}
  \item \textbf{Modelo}: daño escalar \texttt{\ensuremath{\omega} \ensuremath{\in} [0, 1]} sobre el salto. Tracción \texttt{t\_n = (1 − \ensuremath{\omega})\ensuremath{\cdot}K\_e\ensuremath{\cdot}[[u\_n]]} en dirección normal, \texttt{t\_s = 0} en tangencial (Modo-I puro). Activación tipo Rankine (\texttt{f = \ensuremath{\langle}[[u\_n]]\ensuremath{\rangle} − \ensuremath{\kappa}}); historial monótono \texttt{\ensuremath{\kappa}\_\{n+1\} = max(\ensuremath{\kappa}\_n, \ensuremath{\langle}[[u\_n]]\ensuremath{\rangle})} con \texttt{\ensuremath{\kappa}\_0 = \ensuremath{\sigma}\_t0/K\_e}. Energía de fractura \texttt{G\_f} cierra la curva analíticamente:
  \item \textbf{Lineal}: \texttt{T\_soft(\ensuremath{\kappa}) = \ensuremath{\sigma}\_t0\ensuremath{\cdot}(w\_c − \ensuremath{\kappa})/(w\_c − \ensuremath{\kappa}\_0)} con apertura crítica \texttt{w\_c = 2\ensuremath{\cdot}G\_f/\ensuremath{\sigma}\_t0}. \texttt{\ensuremath{\omega}(\ensuremath{\kappa}) = 1 − \ensuremath{\sigma}\_t0\ensuremath{\cdot}(w\_c − \ensuremath{\kappa})/[K\_e\ensuremath{\cdot}\ensuremath{\kappa}\ensuremath{\cdot}(w\_c − \ensuremath{\kappa}\_0)]}. \texttt{\ensuremath{\omega} = 1} en \texttt{\ensuremath{\kappa} \ensuremath{\ge} w\_c}.
  \item \textbf{Exponencial}: \texttt{T\_soft(\ensuremath{\kappa}) = \ensuremath{\sigma}\_t0\ensuremath{\cdot}exp[−\ensuremath{\sigma}\_t0\ensuremath{\cdot}(\ensuremath{\kappa} − \ensuremath{\kappa}\_0)/H]} con \texttt{H = G\_f − \ensuremath{\sigma}\_t0\ensuremath{\cdot}\ensuremath{\kappa}\_0/2}. \texttt{\ensuremath{\omega}(\ensuremath{\kappa}) = 1 − T\_soft(\ensuremath{\kappa})/(K\_e\ensuremath{\cdot}\ensuremath{\kappa})}. Asintótica a 1.
  \item \textbf{Tangente}: simétrica por construcción (rank-1 sobre \texttt{n⊗n}). Algorítmica consistente en carga activa: \texttt{T\_tan = K\_e\ensuremath{\cdot}[(1−\ensuremath{\omega}) − [[u\_n]]\ensuremath{\cdot}d\ensuremath{\omega}/d\ensuremath{\kappa}]\ensuremath{\cdot}(n⊗n)}. Secante reducida \texttt{(1−\ensuremath{\omega})\ensuremath{\cdot}K\_e\ensuremath{\cdot}(n⊗n)} en descarga y sin daño. Cap por \texttt{DAMAGE\_MAX} aplicado \textbf{sólo a la rigidez tangente} cuando \texttt{\ensuremath{\omega} \ensuremath{\ge} DAMAGE\_MAX}; el valor de \texttt{\ensuremath{\omega}} reportado y la tracción usan el valor físico (puede llegar a 1.0 exacto) --- esto distingue al material de \texttt{IsotropicDamage2D}, donde \texttt{K\_e ≡ E} permite truncar \ensuremath{\omega} sin sesgo.
  \item \textbf{JUMP\_DIM}: 2 \ensuremath{\cdot} \textbf{PRIMARY\_STATE\_VAR}: \texttt{'damage'} \ensuremath{\cdot} \textbf{IS\_SYMMETRIC}: \texttt{True}.
  \item \textbf{Parámetros}: \texttt{sigma\_t0} (resistencia a tracción, Pa), \texttt{G\_f} (energía de fractura, N/m), \texttt{K\_e} (rigidez del salto / penalty, Pa/m; sin default automático, ver §12 de la spec para guía \texttt{K\_e \ensuremath{\approx} 10\ensuremath{\cdot}E\_bulk/ℓ\_c}), \texttt{softening \ensuremath{\in} \{'linear', 'exponential'\}}.
  \item \textbf{Variables internas}: \texttt{kappa} (historial del salto equivalente, [m]), \texttt{damage} (\ensuremath{\omega}).
  \item \textbf{Validación energética}: por construcción \texttt{\ensuremath{\int}\_0\textasciicircum{}\{w\_c\} t\ensuremath{\cdot}d[[u\_n]] = G\_f} (lineal) o \texttt{\ensuremath{\int}\_0\textasciicircum{}\ensuremath{\infty} \ensuremath{\approx} G\_f} (exponencial); cubierto por tests con cuadratura trapezoidal.
  \item \textbf{Limitaciones declaradas} (\texttt{out\_of\_scope} en la spec): sólo Modo-I (mixto I--II diferido a fase G del ADR 0010), sin contacto unilateral en compresión (la grieta dañada transmite compresión con rigidez \texttt{(1−\ensuremath{\omega})\ensuremath{\cdot}K\_e}, no rígida), sin anisotropía del daño, sin acoplamiento viscoso/cíclico, sin regularización para mesh-objectivity (se aborda a nivel del elemento \texttt{CST\_Embedded2D}, no del material).
  \item \textbf{Compatible con}: pendiente de fase 2 del ADR 0010 (elemento \texttt{CST\_Embedded2D}). En aislamiento, testeable con historial prescrito de \texttt{[[u]]}.
  \item \textbf{Referencia}: ver \texttt{docs/specs/CohesiveDamageIsotropic.md}. Retama (2010) Cap. 3; Hillerborg, Modéer \& Petersson (1976); Simó \& Ju (1987).
  \item \textbf{Archivo}: solidum/cohesive\_materials/damage\_isotropic.py
\end{itemize}


Familia paralela a los constitutivos mecánicos: relacionan el \textbf{flujo de calor} \texttt{q} con el \textbf{gradiente de temperatura} \texttt{\ensuremath{\nabla}T}, no \texttt{\ensuremath{\sigma}} con \texttt{\ensuremath{\varepsilon}}. Ni el gradiente ni el flujo son tensores simétricos, así que \textbf{no usan notación Voigt}. Viven en \texttt{ThermalMaterialRegistry} y sólo los consumen elementos térmicos; se declaran en el bloque \texttt{materials} del YAML con su \texttt{type} propio.

\section{ThermalConduction --- conducción de calor de Fourier con conductividad tensorial}

\begin{itemize}
  \item \textbf{Ley}: \texttt{q = -k\ensuremath{\cdot}\ensuremath{\nabla}T}. Lineal y \textbf{sin variables internas} --- el análogo térmico de \texttt{Elastic2D}/\texttt{Elastic3D}. El signo negativo es física (segunda ley), no convención elegible.
  \item \textbf{FLUX\_DIM}: 2 ó 3 \ensuremath{\cdot} \textbf{PRIMARY\_STATE\_VAR}: \texttt{None} \ensuremath{\cdot} \textbf{IS\_SYMMETRIC}: \texttt{True}.
  \item A diferencia de \texttt{STRAIN\_DIM}/\texttt{JUMP\_DIM}, \textbf{no es \texttt{ClassVar}}: la dimensión la fija el tensor con que se construye, así que la misma clase sirve en 2D y en 3D.
  \item \textbf{Conductividad tensorial}: \texttt{k} se guarda siempre como matriz. El constructor acepta \textbf{escalar} (isótropo, se expande a \texttt{k\ensuremath{\cdot}I} con la dimensión de \texttt{dim}) o \textbf{matriz} 2\ensuremath{\times}2 / 3\ensuremath{\times}3 (ortótropo/anisótropo), en cuyo caso el tamaño del tensor manda sobre \texttt{dim}.
  \item Validada \textbf{simétrica} (relaciones recíprocas de Onsager) con tolerancia escalada \texttt{1e-12\ensuremath{\cdot}max|k|}, adimensional respecto a las unidades del usuario (ADR 0006).
  \item Validada \textbf{definida positiva} por autovalores (\texttt{eigvalsh}), no por determinante: en dimensión \ensuremath{\ge} 2 un determinante positivo no implica definición positiva.
  \item \textbf{Parámetros}: \texttt{k} (conductividad, W/(m\ensuremath{\cdot}K); escalar o matriz), \texttt{c} (calor específico, J/(kg\ensuremath{\cdot}K), opcional), \texttt{density} (kg/m\ensuremath{^{3}}, opcional), \texttt{dim} (2 ó 3, sólo para expandir un \texttt{k} escalar).
  \item \textbf{\texttt{c} y \texttt{density} opcionales} siguiendo el criterio del ADR 0008: un análisis \textbf{estacionario} sólo usa \texttt{k}, y exigir propiedades que no entran en ninguna ecuación obligaría a inventar valores. Si un consumidor transitorio las encuentra ausentes, \texttt{ValueError} \textbf{accionable} --- nombra el material, los parámetros que faltan y la salida (declararlos o usar un solver estacionario).
  \item \textbf{Magnitud derivada}: \texttt{thermal\_diffusivity} = \texttt{\ensuremath{\alpha} = k/(\ensuremath{\rho}c)} [m\ensuremath{^{2}}/s], que gobierna la velocidad del frente térmico y sirve para estimar \texttt{\ensuremath{\Delta}t \textasciitilde{} h\ensuremath{^{2}}/\ensuremath{\alpha}}. \textbf{Rechaza el caso anisótropo} en vez de devolver un escalar plausible: con \texttt{k} tensorial la difusividad es un tensor.
  \item \textbf{Escala de temperatura}: la conducción pura sólo involucra \textbf{diferencias} de \texttt{T}, luego K y \ensuremath{^{\circ}}C son equivalentes. La distinción importará al entrar radiación (\texttt{T⁴} exige escala absoluta).
  \item \textbf{Caveat de interoperabilidad}: no aplica el caveat de Voigt de ABAQUS --- \texttt{k} es un tensor de orden 2 en ejes físicos, sin reordenamiento.
  \item \textbf{Limitaciones declaradas} (\texttt{out\_of\_scope} en la spec): conductividad dependiente de la temperatura \texttt{k(T)} (introduciría no linealidad), convección y radiación (son condiciones de frontera, no del material), cambio de fase / calor latente, acoplamiento termomecánico.
  \item \textbf{Compatible con}: \texttt{Quad4Thermal}, \texttt{Hex8Thermal} (elementos térmicos con \texttt{field='temperature'}).
  \item \textbf{Referencia}: ver \texttt{docs/specs/ThermalConduction.md}. Incropera \& DeWitt (2007) §2; Carslaw \& Jaeger (1959) §1; Lewis, Nithiarasu \& Seetharamu (2004) §3.
  \item \textbf{Archivo}: solidum/materials/thermal\_conduction.py \ensuremath{\cdot} clase base en solidum/core/thermal\_material.py
\end{itemize}


\section{Cómo añadir un material nuevo}

\texttt{/solidum-new material <Name>} --- genera archivo en \texttt{solidum/materials/}, decorador \texttt{@MaterialRegistry.register}, esqueleto de test.
Declarar \textbf{\texttt{STRAIN\_DIM}} y, si tiene historia, \textbf{\texttt{PRIMARY\_STATE\_VAR}} (la variable que aparecerá en el VTK).
Tras implementar el modelo constitutivo, \textbf{añadir una entrada a este catálogo} siguiendo el formato de arriba.

Para un material \textbf{cohesivo} nuevo: el patrón es análogo pero el archivo vive en \texttt{solidum/cohesive\_materials/}, la clase hereda de \texttt{CohesiveMaterial} (no \texttt{Material}) y se registra con \texttt{@CohesiveMaterialRegistry.register}. Declarar \texttt{JUMP\_DIM}, \texttt{PRIMARY\_STATE\_VAR} y \texttt{IS\_SYMMETRIC}. Añadir la entrada en la sección "Materiales cohesivos" de este mismo catálogo.

Para un material \textbf{térmico} nuevo: la clase hereda de \texttt{ThermalMaterial} (en \texttt{solidum/core/thermal\_material.py}) y se registra con \texttt{@ThermalMaterialRegistry.register}. El método a implementar es \texttt{compute\_flux(\ensuremath{\nabla}T) -> (q, k)}, no \texttt{compute\_state}. Declarar \texttt{FLUX\_DIM} (2 ó 3 --- como propiedad de instancia si la clase sirve en ambas dimensiones), \texttt{PRIMARY\_STATE\_VAR} e \texttt{IS\_SYMMETRIC}. La validación de \texttt{c}/\texttt{density} con mensaje accionable ya la aporta \texttt{volumetric\_capacity()} de la clase base; no reimplementarla. En la spec, el contrato lleva \texttt{kind: thermal\_material} e \texttt{interface.field: temperature}. Añadir la entrada en la sección "Materiales térmicos" de este mismo catálogo.


\newpage

\chapter{Solvers — catálogo}



\textit{Referencia rápida de los métodos de solución implementados. Una entrada por solver. Para detalles del algoritmo \ensuremath{\to} código fuente.}



\section{LinearSolver --- solución lineal directa en un paso}

\begin{itemize}
  \item \textbf{Propósito}: resolver \texttt{K \ensuremath{\cdot} U = F} en un único \texttt{spsolve}.
  \item \textbf{Esquema}: ensamblaje único \ensuremath{\to} eliminación directa de DOFs prescritos (ADR 0004) \ensuremath{\to} \texttt{scipy.sparse.linalg.spsolve}.
  \item \textbf{Parámetros}: \texttt{linear\_algebra} (default \texttt{auto}; ADR 0003 §4).
  \item \textbf{Cuándo usarlo}: problemas estrictamente lineales (todos los materiales con tangente constante, sin grandes desplazamientos, sin contacto).
  \item \textbf{No converge / no aplica}: cualquier no-linealidad material (daño, plasticidad) o geométrica (corotacional).
  \item \textbf{Spec}: docs/specs/LinearSolver.md
  \item \textbf{Archivo}: solidum/math/solvers/linear.py
\end{itemize}


\section{NonlinearSolver --- Newton-Raphson incremental con paso adaptativo}

\begin{itemize}
  \item \textbf{Propósito}: resolver problemas no lineales con control de carga, dividiendo la carga total en pasos y aplicando Newton-Raphson dentro de cada paso.
  \item \textbf{Esquema}:
  \item Bucle externo: incrementar factor de carga \texttt{\ensuremath{\lambda}} desde 0 hasta 1 en \texttt{num\_steps} pasos.
  \item Bucle interno: Newton-Raphson con \texttt{K\_t \ensuremath{\cdot} \ensuremath{\Delta}U = R = \ensuremath{\lambda}\ensuremath{\cdot}F\_ext − F\_int}.
  \item \textbf{Criterio de convergencia dual}: \texttt{max(err\_disp, err\_force) < tol}, con
\end{itemize}
    \texttt{err\_disp = \ensuremath{\|}\ensuremath{\Delta}U\ensuremath{\|} / (\ensuremath{\|}U\ensuremath{\|} + \ensuremath{\varepsilon})} y \texttt{err\_force = \ensuremath{\|}R\ensuremath{\|} / max(\ensuremath{\|}\ensuremath{\lambda}\ensuremath{\cdot}F\_ext\ensuremath{\|}, \ensuremath{\|}F\_int\ensuremath{\|}, \ensuremath{\varepsilon})}.
\begin{itemize}
  \item \textbf{Adaptatividad} (\texttt{adaptive=True}): si converge en < 5 iteraciones, agranda el siguiente paso (\ensuremath{\times}1.5); si no converge, biseca (÷2). Falla si \texttt{\ensuremath{\Delta}\ensuremath{\lambda} < min\_delta\_lambda}.
  \item Tras converger un paso \ensuremath{\to} \texttt{assembler.commit\_all\_states()} (trial \ensuremath{\to} committed en todos los elementos).
  \item \textbf{Parámetros}: \texttt{convergence}, \texttt{max\_iter}, \texttt{num\_steps}, \texttt{adaptive}, \texttt{min\_delta\_lambda}, \texttt{linear\_algebra}, \texttt{freeze\_tangent\_after\_iter}, \textbf{\texttt{line\_search}} (default \texttt{False}, ADR 0011).
  \item \textbf{Cuándo usarlo}: la opción por defecto para no-linealidad material o geométrica suave (sin snap-back).
  \item \textbf{Cuándo activar \texttt{line\_search=True}}: cuando se observe oscilación del residuo entre iteraciones (síntoma clásico: ratios alternados sin descenso monótono). Caso documentado: plasticidad perfecta con carga cerca/sobre la capacidad. \textbf{Default \texttt{False}} porque en problemas con tangente consistente cuasi-cuadrática (daño activo, plasticidad estándar) el line search rechaza pasos correctos del Newton donde el residuo sube transitoriamente --- ver ADR 0011 §"Enmiendas".
  \item \textbf{Diagnóstico al diverger}: lanza una subclase tipada de \texttt{RuntimeError} (\texttt{OscillatingNewtonError}, \texttt{SingularTangentError}, \texttt{LoadExceedsCapacityError}, \texttt{UnknownDivergenceError}) con métricas (último \ensuremath{\|}R\ensuremath{\|}, último \ensuremath{\|}\ensuremath{\delta}U\ensuremath{\|}, factor de carga, bisecciones consumidas) y \texttt{hint} textual. Ver \texttt{solidum/math/solvers/diagnostics.py}.
  \item \textbf{Limitación}: con bisección adaptativa y cinemática no lineal (corotacional) puede atravesar puntos límite suaves; no atraviesa snap-back con \texttt{du/d\ensuremath{\lambda} < 0}. Para eso \ensuremath{\to} \texttt{ArcLengthSolver}. Hallazgo de la auditoría fase A: el solver es más capaz de lo asumido históricamente (ver \texttt{docs/auditorias/solvers\_robustez\_fase\_A.md}).
  \item \textbf{Referencia}: Crisfield, \textit{Non-linear Finite Element Analysis of Solids and Structures}, vol. 1, cap. 9. Line search: Grippo-Lampariello-Lucidi 1986 (variante de descenso no monótono).
  \item \textbf{Spec}: docs/specs/NonlinearSolver.md
  \item \textbf{Archivo}: solidum/math/solvers/nonlinear.py
\end{itemize}


\section{ArcLengthSolver --- método de longitud de arco cilíndrico (Crisfield)}

\begin{itemize}
  \item \textbf{Propósito}: trazar curvas de equilibrio con snap-through, snap-back o pérdida de unicidad de carga, controlando simultáneamente desplazamientos y factor de carga.
  \item \textbf{Esquema}:
  \item \textbf{Predictor tangente}: \texttt{du\_t = K\_t\ensuremath{^{-}}¹ \ensuremath{\cdot} F\_ext\_ref}; \texttt{d\ensuremath{\lambda} = sign \ensuremath{\cdot} dl / \ensuremath{\|}du\_t\ensuremath{\|}}.
\end{itemize}
    El \texttt{sign} se elige por proyección con el incremento del paso anterior (evitar regresar).
\begin{itemize}
  \item \textbf{Corrector iterativo}: en cada iteración resuelve dos sistemas (\texttt{du\_R = K\ensuremath{^{-}}¹\ensuremath{\cdot}R} y \texttt{du\_t = K\ensuremath{^{-}}¹\ensuremath{\cdot}F\_ext}) y aplica la restricción cuadrática cilíndrica de Crisfield: \texttt{\ensuremath{\|}dU\ensuremath{\|}\ensuremath{^{2}} = dl\ensuremath{^{2}}}.
  \item De las dos raíces se elige la que mantiene el ángulo positivo con el incremento previo.
  \item \textbf{Caso especial --- final\_step}: si el predictor sobrepasaría \texttt{max\_lambda}, fija \ensuremath{\lambda} exactamente a \texttt{max\_lambda} y resuelve solo desplazamientos (Newton-Raphson puro).
  \item \textbf{Auto-ajuste de \texttt{dl}}: < 4 iter \ensuremath{\to} ampliar (\ensuremath{\times}1.5, tope \texttt{5\ensuremath{\cdot}dl\_initial}); > 8 iter \ensuremath{\to} reducir (\ensuremath{\times}0.6); no converge \ensuremath{\to} biseca (÷2).
  \item Convergencia: mismo criterio dual que \texttt{NonlinearSolver}.
  \item \textbf{Parámetros}: \texttt{tol}, \texttt{max\_iter}, \texttt{max\_lambda}, \texttt{initial\_dl}, \texttt{max\_steps}, \texttt{dl\_grow\_factor}, \texttt{dl\_max\_factor}, \texttt{dl\_shrink\_factor}, \texttt{dl\_grow\_iter\_threshold}, \texttt{dl\_shrink\_iter\_threshold}.
  \item \textbf{Cuándo usarlo}: problemas con softening pronunciado (daño, post-pandeo, snap-through de cúpulas), o cuando \texttt{NonlinearSolver} diverge cerca de un punto límite.
  \item \textbf{Limitación}: más caro por paso (dos \texttt{spsolve} por iteración); requiere ajuste de \texttt{initial\_dl} para problemas nuevos.
  \item \textbf{Referencia}: Crisfield, "A fast incremental/iterative solution procedure that handles snap-through" (Computers \& Structures, 1981); Crisfield vol. 1, cap. 9.
  \item \textbf{Spec}: docs/specs/ArcLengthSolver.md
  \item \textbf{Archivo}: solidum/math/solvers/arclength.py
\end{itemize}


\section{DissipationArcLengthSolver --- variante de arc-length por disipación (Gutiérrez 2004)}

\begin{itemize}
  \item \textbf{Propósito}: trazar la rama post-pico de problemas con \textbf{softening severo} controlando la disipación incremental de energía por paso en lugar de la longitud de arco euclídea. Más robusto que el cilíndrico cuando la curva U-\ensuremath{\lambda} tiene tramos casi-verticales en el plano (U, \ensuremath{\lambda}).
  \item \textbf{Esquema}:
  \item Subclase de \texttt{ArcLengthSolver}; reusa predictor tangente, corrector Newton, ajuste de paso, Dirichlet/MPC y backend algebraico.
  \item \textbf{Restricción de paso}: \texttt{g(\ensuremath{\Delta}U, \ensuremath{\Delta}\ensuremath{\lambda}) = ½\ensuremath{\cdot}(\ensuremath{\lambda}\_n\ensuremath{\cdot}F\ensuremath{\cdot}\ensuremath{\Delta}U − \ensuremath{\Delta}\ensuremath{\lambda}\ensuremath{\cdot}F\ensuremath{\cdot}U\_n) = \ensuremath{\tau}}. \textbf{Lineal} en \texttt{(\ensuremath{\Delta}U, \ensuremath{\Delta}\ensuremath{\lambda})} \ensuremath{\Rightarrow} corrector con una sola raíz en \texttt{dd\ensuremath{\lambda}}, sin selección por menor ángulo ni patología de raíces imaginarias.
  \item \textbf{Switching automático cilíndrico\ensuremath{\leftrightarrow}disipación} según se detecte disipación neta sobre el umbral relativo \texttt{dissipation\_threshold\ensuremath{\cdot}\ensuremath{\|}F\_ref\ensuremath{\|}\ensuremath{\cdot}\ensuremath{\|}U\ensuremath{\|}}. Arranque obligado en cilíndrico (\texttt{\ensuremath{\alpha} = ½\ensuremath{\cdot}(\ensuremath{\lambda}\_n\ensuremath{\cdot}F\ensuremath{\cdot}du\_t − F\ensuremath{\cdot}U\_n)} es 0 cuando \texttt{\ensuremath{\lambda}\_n = U\_n = 0}).
  \item \textbf{Sign-of-pivot tracking aproximado} vía signo del \texttt{slogdet(K\_t)}. Distingue 0 vs número impar de pivots negativos --- diagnóstico de paso por punto límite simple. Tracking exacto requiere LDL\ensuremath{^{\top}} Bunch-Kaufman (deuda técnica \#7 STATUS.md).
  \item \textbf{Salvaguarda contra \texttt{final\_step} prematuro}: si \texttt{|d\ensuremath{\lambda}\_pred| > 3\ensuremath{\cdot}(max\_lambda − \ensuremath{\lambda}\_curr)}, bisecta dl o \ensuremath{\tau} antes de aceptar el paso.
  \item \textbf{Detección de \ensuremath{\alpha}\ensuremath{\approx}0}: threshold relativo \texttt{|\ensuremath{\alpha}| < 1e-6\ensuremath{\cdot}escala}; revierte temporalmente a cilíndrico (en problemas lineales monotónicos \texttt{\ensuremath{\alpha} ≡ 0} exactamente por construcción).
  \item \textbf{Parámetros heredados}: todos los de \texttt{ArcLengthSolver} (\texttt{max\_iter}, \texttt{max\_lambda}, \texttt{initial\_dl}, \texttt{max\_steps}, factores \texttt{dl\_*}, \texttt{linear\_algebra}). \textbf{Nuevos}: \texttt{initial\_tau} (obligatorio), \texttt{tau\_grow\_factor}, \texttt{tau\_max\_factor}, \texttt{tau\_shrink\_factor}, \texttt{tau\_grow\_iter\_threshold}, \texttt{tau\_shrink\_iter\_threshold}, \texttt{dissipation\_threshold}.
  \item \textbf{Cuándo usarlo}: problemas con softening continuo donde \texttt{ArcLengthSolver} cilíndrico bisecta excesivamente cerca del pico --- daño bulk 1D/2D, plasticidad con localización progresiva, post-pandeo con descarga elástica.
  \item \textbf{Cuándo NO usarlo (limitación documentada)}: cohesivo+embedded discontinuity con \texttt{K\_e} stiff (\texttt{\textasciitilde{}1e13}). La activación discreta del Rankine introduce un salto en \texttt{F\_int}/\texttt{K\_t} que el \texttt{sign} por producto escalar no maneja graciosamente; el predictor del paso siguiente reorienta hacia descarga. Resolverlo requiere LDL\ensuremath{^{\top}} Bunch-Kaufman real (item \#7 deuda STATUS.md) o control por CMOD/CTOD del salto (spec separada). En problemas lineales puramente elásticos coincide con \texttt{ArcLengthSolver} cilíndrico a paridad de bits (regresión validada).
  \item \textbf{Referencias}: Gutiérrez (2004), \textit{Communications in Numerical Methods in Engineering} 20, 19-29; Verhoosel, Remmers, Gutiérrez (2009), \textit{IJNME} 77, 1290-1321; Crisfield (1991) vol. 1 §9.4.
  \item \textbf{Spec}: docs/specs/DissipationArcLengthSolver.md (status \texttt{implemented}, validación parcial).
  \item \textbf{Archivo}: solidum/math/solvers/dissipation\_arclength.py
\end{itemize}


\section{ModalSolver --- análisis modal por autovalores generalizados (ADR 0009)}

\begin{itemize}
  \item \textbf{Propósito}: calcular frecuencias naturales \texttt{\ensuremath{\omega}\_n} y modos \texttt{\ensuremath{\varphi}\_n} de vibración no amortiguados resolviendo \texttt{K \ensuremath{\cdot} \ensuremath{\varphi} = \ensuremath{\omega}\ensuremath{^{2}} \ensuremath{\cdot} M \ensuremath{\cdot} \ensuremath{\varphi}}.
  \item \textbf{Esquema}:
  \item Ensamblaje de \texttt{K} (rigidez lineal en \texttt{u = 0}) y \texttt{M} (masa consistente, ADR 0009 §1) reusando la topología COO cacheada del \texttt{Assembler}.
  \item Reducción simultánea por Dirichlet \texttt{Assembler.reduce\_pair(K, M)} \ensuremath{\to} \texttt{K\_red, M\_red, T}.
  \item \texttt{EigenSolver} envuelve \texttt{scipy.sparse.linalg.eigsh} con \textbf{shift-invert} en \texttt{sigma} (\texttt{sigma=0} \ensuremath{\to} frecuencias más bajas) y \texttt{which="LM"}. ARPACK factoriza \texttt{(K\_red − \ensuremath{\sigma}\ensuremath{\cdot}M\_red)} una sola vez y la reusa en todas las iteraciones Lanczos.
  \item Expansión de modos al espacio completo: \texttt{\ensuremath{\varphi} = T \ensuremath{\cdot} \ensuremath{\varphi}\_red}. En DOFs prescritos por Dirichlet la componente es 0.
  \item Salida en \texttt{ModalResult}: \texttt{frequencies\_rad}, \texttt{frequencies\_hz}, \texttt{periods}, \texttt{modes} (M-ortonormales), \texttt{n\_modes}, \texttt{converged}.
  \item \textbf{Parámetros}: \texttt{n\_modes} (obligatorio), \texttt{sigma} (default 0.0), \texttt{which} (default \texttt{"LM"}), \texttt{tolerance} (default 1.0e-9), \texttt{lumping} (default \texttt{"consistent"}; fase 1 solo admite ese valor), \texttt{linear\_algebra} (reservado; ADR 0003).
  \item \textbf{Firma del solver}: \texttt{solve()} sin argumentos (no consume \texttt{F\_applied}). El entrypoint público es \texttt{solidum.run\_modal(domain, n\_modes=N)} o YAML con \texttt{solver: type: ModalSolver}. \texttt{solidum.run\_yaml} despacha automáticamente al detectar el tipo.
  \item \textbf{Cuándo usarlo}: caracterización dinámica de la estructura (frecuencias naturales, formas modales) antes de pasar a análisis transitorio o sísmico. Validación de la matriz de masa contra fórmulas analíticas (barra empotrada-libre, viga Bernoulli-Euler).
  \item \textbf{Limitaciones}: lineal (\texttt{K} evaluada en \texttt{u = 0}); sin amortiguamiento; modos complejos no soportados. Cómputos derivados (factores de participación, masas efectivas, combinación espectral SRSS/CQC) viven en \texttt{solidum/math/modal\_response.py} y los consume \texttt{ResponseSpectrumSolver} (ADR 0009 fase 7 cerrada).
  \item \textbf{Referencia}: Bathe, \textit{Finite Element Procedures}, cap. 9 (masa) y cap. 11 (algoritmos de autovalor); Hughes, \textit{The Finite Element Method}, cap. 7; ARPACK Users' Guide (Lehoucq--Sorensen--Yang).
  \item \textbf{Spec}: docs/specs/ModalSolver.md
  \item \textbf{Archivos}: solidum/math/solvers/modal.py (clase \texttt{ModalSolver}), solidum/math/linalg/eigen.py (clase \texttt{EigenSolver}), solidum/results.py (\texttt{ModalResult}).
\end{itemize}


\section{NewmarkSolver --- análisis dinámico transitorio lineal (ADR 0009 fase 3)}

\begin{itemize}
  \item \textbf{Propósito}: integrar en el tiempo la ecuación de movimiento semidiscreta \texttt{M\ensuremath{\cdot}ü + C\ensuremath{\cdot}u̇ + K\ensuremath{\cdot}u = F(t)} con amortiguamiento Rayleigh \texttt{C = \ensuremath{\alpha}\ensuremath{\cdot}M + \ensuremath{\beta}\ensuremath{\cdot}K} y apoyos Dirichlet constantes.
  \item \textbf{Esquema}:
  \item Familia Newmark-\ensuremath{\beta} con \texttt{(\ensuremath{\beta}, \ensuremath{\gamma})} parametrizables; default \texttt{(1/4, 1/2)} --- average acceleration, incondicionalmente estable, sin amortiguamiento numérico, error \texttt{O(\ensuremath{\Delta}t\ensuremath{^{2}})}.
  \item \textbf{Predictores}: \texttt{ũ = u\_n + \ensuremath{\Delta}t\ensuremath{\cdot}u̇\_n + (\ensuremath{\Delta}t\ensuremath{^{2}}/2)(1−2\ensuremath{\beta})\ensuremath{\cdot}ü\_n}, \texttt{ũ̇ = u̇\_n + \ensuremath{\Delta}t\ensuremath{\cdot}(1−\ensuremath{\gamma})\ensuremath{\cdot}ü\_n}.
  \item \textbf{Sistema efectivo}: \texttt{A\_eff\ensuremath{\cdot}ü\_\{n+1\} = F\_\{n+1\} − C\ensuremath{\cdot}ũ̇ − K\ensuremath{\cdot}ũ} con \texttt{A\_eff = M + \ensuremath{\gamma}\ensuremath{\Delta}t\ensuremath{\cdot}C + \ensuremath{\beta}\ensuremath{\Delta}t\ensuremath{^{2}}\ensuremath{\cdot}K}.
  \item \textbf{Correctores}: \texttt{u\_\{n+1\} = ũ + \ensuremath{\beta}\ensuremath{\Delta}t\ensuremath{^{2}}\ensuremath{\cdot}ü\_\{n+1\}}, \texttt{u̇\_\{n+1\} = ũ̇ + \ensuremath{\gamma}\ensuremath{\Delta}t\ensuremath{\cdot}ü\_\{n+1\}}.
  \item Reducción por Dirichlet (operador \texttt{T}) aplicada a \texttt{(K, M, C)} simultáneamente; factorización única de \texttt{A\_eff\_red} reusada en todos los pasos.
  \item Aceleración inicial consistente con la ecuación de movimiento en \texttt{t=0}.
  \item \textbf{Parámetros}: \texttt{t\_end}, \texttt{dt} (obligatorios), \texttt{beta} (default 0.25), \texttt{gamma} (default 0.5), \texttt{rayleigh} (None | \texttt{\{alpha, beta\}} directos | \texttt{\{xi1, omega1, xi2, omega2\}} para calibración modal), \texttt{u0}, \texttt{u0\_dot} (default ceros), \texttt{F\_func} (callback Python \texttt{t \ensuremath{\to} F(t)}; default vibración libre), \texttt{linear\_algebra}, \texttt{lumping}.
  \item \textbf{Firma del solver}: \texttt{solve()} sin argumentos, retorna \texttt{TransientResult} con historiales \texttt{t\_history}, \texttt{u\_history}, \texttt{udot\_history}, \texttt{uddot\_history} (shape \texttt{(n\_dof, n\_steps + 1)}).
  \item \textbf{Cuándo usarlo}: respuesta a cargas dependientes del tiempo (sísmica time-history, impacto, vibración forzada), evolución transitoria desde condiciones iniciales arbitrarias, validación de respuesta libre/forzada de estructuras lineales.
  \item \textbf{Limitaciones}:
  \item Solo lineal (\texttt{K}, \texttt{M} constantes). La extensión a no-linealidad (Newton dentro de cada paso) es fase 4 del ADR 0009.
  \item Apoyos Dirichlet constantes en el tiempo. Multi-support excitation sísmica diferida.
  \item Paso \texttt{\ensuremath{\Delta}t} fijo. Adaptativo diferido.
  \item HHT-\ensuremath{\alpha} y generalized-\ensuremath{\alpha} (con amortiguamiento numérico controlado) no incluidos.
  \item \textbf{Referencia}: Newmark (1959), J. Eng. Mech. Div. ASCE; Bathe, \textit{FEP} cap. 9.4; Hughes, \textit{FEM} cap. 9; Chopra, \textit{Dynamics of Structures} cap. 5 y 15.
  \item \textbf{Spec}: docs/specs/NewmarkSolver.md
  \item \textbf{Archivos}: solidum/math/solvers/newmark.py (clase \texttt{NewmarkSolver}), solidum/math/damping.py (Rayleigh), solidum/results.py (\texttt{TransientResult}).
\end{itemize}


\section{NewtonNewmarkSolver --- análisis dinámico transitorio no lineal (ADR 0009 fase 4)}

\begin{itemize}
  \item \textbf{Propósito}: integrar en el tiempo \texttt{M\ensuremath{\cdot}ü + C\ensuremath{\cdot}u̇ + F\_int(u) = F\_ext(t)} con materiales no lineales (plasticidad, daño) o no linealidad geométrica. \textbf{Variante} de \texttt{NewmarkSolver} (Reglas §4): hereda predictores/correctores y reducción de Dirichlet; añade Newton-Raphson dentro de cada paso temporal.
  \item \textbf{Esquema}:
  \item Predictores Newmark idénticos a la fase 3.
  \item \textbf{Bucle interno de Newton-Raphson} sobre el residuo dinámico
\end{itemize}
    \texttt{R(ü\_\{n+1\}) = F\_ext(t\_\{n+1\}) − F\_int(u\_\{n+1\}) − C\ensuremath{\cdot}u̇\_\{n+1\} − M\ensuremath{\cdot}ü\_\{n+1\}},
    con jacobiano \texttt{J = M + \ensuremath{\gamma}\ensuremath{\Delta}t\ensuremath{\cdot}C + \ensuremath{\beta}\ensuremath{\Delta}t\ensuremath{^{2}}\ensuremath{\cdot}K\_t} y \texttt{K\_t} la rigidez tangente corriente.
\begin{itemize}
  \item \textbf{Convergencia dual} (ADR 0007): \texttt{R/(atol + rtol\ensuremath{\cdot}F\_ref) \ensuremath{\le} 1} y \texttt{\ensuremath{\delta}u/(atol + rtol\ensuremath{\cdot}\ensuremath{\|}u\ensuremath{\|}) \ensuremath{\le} 1}.
  \item \textbf{Commit de estado}: tras converger cada paso, \texttt{assembler.commit\_all\_states()} (trial \ensuremath{\to} committed). Si no converge en \texttt{max\_iter}, lanza \texttt{RuntimeError} y se descarta el state trial.
  \item \textbf{Amortiguamiento Rayleigh constante en el tiempo}, calibrado con la \textbf{rigidez elástica de referencia} \texttt{K\_0} (al inicio del análisis, \texttt{u = 0}). Convención estándar (Abaqus, ANSYS, OpenSees); evita acoplamiento ad-hoc entre disipación viscosa y plástica.
  \item \textbf{Newton modificado opcional} (\texttt{freeze\_tangent\_after\_iter}, ADR 0003 fase 2): factoriza fresco las primeras N iter de cada paso y reusa la factorización en las siguientes.
  \item \textbf{Recuperación del caso lineal}: con materiales lineales (\texttt{K\_t ≡ K\_0}), el residuo se anula en 1 iter y el resultado coincide exactamente con \texttt{NewmarkSolver}. Validado en tests.
  \item \textbf{Parámetros}: heredados de \texttt{NewmarkSolver} (\texttt{t\_end}, \texttt{dt}, \texttt{beta}, \texttt{gamma}, \texttt{rayleigh}, \texttt{u0}, \texttt{u0\_dot}, \texttt{F\_func}, \texttt{linear\_algebra}, \texttt{lumping}) más \texttt{convergence} (\texttt{ConvergenceCriterion}, ADR 0007), \texttt{max\_iter} (default 20), \texttt{freeze\_tangent\_after\_iter} (default \texttt{None}), \textbf{\texttt{line\_search}} (default \texttt{False}, ADR 0011 --- mismo criterio que \texttt{NonlinearSolver}).
  \item \textbf{Diagnóstico al diverger}: lanza una subclase tipada de \texttt{RuntimeError} con métricas y \texttt{hint} (mismo patrón que \texttt{NonlinearSolver}, ADR 0011).
  \item \textbf{Firma del solver}: \texttt{solve()} sin argumentos, retorna \texttt{TransientResult} (mismo formato que \texttt{NewmarkSolver}).
  \item \textbf{Cuándo usarlo}: respuesta dinámica de estructuras con plasticidad transitoria, fatiga de bajo ciclo, impacto en hormigón friccional, vibraciones de marcos con disipación plástica.
  \item \textbf{Limitaciones}: igual que \texttt{NewmarkSolver} en lo lineal (apoyos constantes, paso fijo). Además: no resuelve snap-back en problemas con softening severo --- combinar con \texttt{ArcLengthSolver} si emerge.
  \item \textbf{Despacho YAML}: \texttt{solver.type: NewtonNewmarkSolver}. Hereda \texttt{PIPELINE\_KIND = "transient"} de \texttt{NewmarkSolver}; \texttt{entry.run\_yaml} despacha por ese atributo declarativo (regla C del ADR 0009, aplicada 2026-05-18) --- sin \texttt{isinstance} ni cadenas hardcoded.
  \item \textbf{Referencia}: ADR 0009 §Fase 4; Hughes, \textit{FEM} cap. 7-9 (Newton dentro de paso temporal); Crisfield vol. 2 cap. 24 (dinámica no lineal).
  \item \textbf{Spec}: docs/specs/NewtonNewmarkSolver.md
  \item \textbf{Archivo}: solidum/math/solvers/newmark.py (clase \texttt{NewtonNewmarkSolver}, junto al padre).
\end{itemize}


\section{HHTSolver --- análisis dinámico transitorio lineal con disipación numérica (Hilber-Hughes-Taylor \ensuremath{\alpha})}

\begin{itemize}
  \item \textbf{Propósito}: integrar \texttt{M\ensuremath{\cdot}ü + C\ensuremath{\cdot}u̇ + K\ensuremath{\cdot}u = F(t)} con \textbf{amortiguamiento numérico controlado} que disipa selectivamente los modos de alta frecuencia espurios sin afectar los modos resueltos. \textbf{Variante} de \texttt{NewmarkSolver} (Reglas §4): hereda predictores/correctores, reducción Dirichlet y factorización única de \texttt{A\_eff}.
  \item \textbf{Esquema}:
  \item Ecuación de equilibrio temporal: \texttt{M\ensuremath{\cdot}ü\_\{n+1\} + (1+\ensuremath{\alpha})\ensuremath{\cdot}C\ensuremath{\cdot}u̇\_\{n+1\} − \ensuremath{\alpha}\ensuremath{\cdot}C\ensuremath{\cdot}u̇\_n + (1+\ensuremath{\alpha})\ensuremath{\cdot}K\ensuremath{\cdot}u\_\{n+1\} − \ensuremath{\alpha}\ensuremath{\cdot}K\ensuremath{\cdot}u\_n = (1+\ensuremath{\alpha})\ensuremath{\cdot}F\_\{n+1\} − \ensuremath{\alpha}\ensuremath{\cdot}F\_n}.
  \item \textbf{Auto-derivación de \ensuremath{\beta}, \ensuremath{\gamma} desde \ensuremath{\alpha}} (Hilber 1977): \texttt{\ensuremath{\beta} = (1−\ensuremath{\alpha})\ensuremath{^{2}}/4}, \texttt{\ensuremath{\gamma} = (1−2\ensuremath{\alpha})/2}. Preserva orden 2 y estabilidad incondicional.
  \item \textbf{Sistema efectivo}: \texttt{A\_eff = M + (1+\ensuremath{\alpha})\ensuremath{\cdot}\ensuremath{\gamma}\ensuremath{\Delta}t\ensuremath{\cdot}C + (1+\ensuremath{\alpha})\ensuremath{\cdot}\ensuremath{\beta}\ensuremath{\Delta}t\ensuremath{^{2}}\ensuremath{\cdot}K}. Constante en el tiempo \ensuremath{\to} factorización única reutilizable (como Newmark).
  \item \textbf{Radio espectral en alta frecuencia}: \texttt{\ensuremath{\rho}\_\ensuremath{\infty} = (1+\ensuremath{\alpha})/(1−\ensuremath{\alpha})}. Con \texttt{\ensuremath{\alpha}=0} \ensuremath{\to} \texttt{\ensuremath{\rho}\_\ensuremath{\infty}=1} (sin disipación, recupera Newmark trapezoidal). Con \texttt{\ensuremath{\alpha}=−0.05} \ensuremath{\to} \texttt{\ensuremath{\rho}\_\ensuremath{\infty}\ensuremath{\approx}0.905}. Con \texttt{\ensuremath{\alpha}=−1/3} \ensuremath{\to} \texttt{\ensuremath{\rho}\_\ensuremath{\infty}=0.5} (disipación máxima manteniendo orden 2).
  \item \textbf{Disipación selectiva}: modos con \texttt{\ensuremath{\omega}\ensuremath{\Delta}t > \ensuremath{\pi}} se atenúan significativamente; modos con \texttt{\ensuremath{\omega}\ensuremath{\Delta}t « 1} casi intactos. Es lo que distingue HHT-\ensuremath{\alpha} de subir \ensuremath{\gamma} Newmark (que disipa en todas las frecuencias y baja a orden 1).
  \item \textbf{Parámetros}: \texttt{t\_end}, \texttt{dt} obligatorios. \textbf{\texttt{alpha}} (default \texttt{-0.05}, rango \texttt{[−1/3, 0]}). \texttt{beta}, \texttt{gamma} autoderivados desde \texttt{alpha} salvo override explícito (no recomendado fuera de experimentación). Resto heredado de \texttt{NewmarkSolver}: \texttt{rayleigh}, \texttt{u0}, \texttt{u0\_dot}, \texttt{F\_func}, \texttt{linear\_algebra}, \texttt{lumping}.
  \item \textbf{Cuándo usarlo}: cuando la respuesta transitoria contenga modos de alta frecuencia espurios (mallas con modos altos no de interés, impacto, contacto, propagación con \texttt{\ensuremath{\omega}\ensuremath{\Delta}t} cerca o por encima de \ensuremath{\pi}) que enmascaren la señal. Si quieres exactamente Newmark trapezoidal sin disipación, usa \texttt{NewmarkSolver}.
  \item \textbf{Default de \texttt{alpha}}: \texttt{-0.05} es el valor canónico (Hilber 1977; Abaqus, OpenSees, ANSYS): disipación >9\% por ciclo en altas frecuencias, <0.5\% en modos de interés. Si pides \texttt{HHTSolver} con \texttt{alpha=0.0}, es funcionalmente idéntico a \texttt{NewmarkSolver(beta=0.25, gamma=0.5)}.
  \item \textbf{Despacho YAML}: \texttt{solver.type: HHTSolver}. Hereda \texttt{PIPELINE\_KIND = "transient"} de \texttt{NewmarkSolver}; \texttt{entry.run\_yaml} despacha por el atributo declarativo (regla C ADR 0009).
  \item \textbf{Referencia}: Hilber, Hughes \& Taylor (1977), \textit{Earthquake Eng. Struct. Dyn.} 5(3) 283-292; Hughes \textit{FEM} §9.3; Chopra \textit{Dynamics of Structures} §15.4.
  \item \textbf{Spec}: docs/specs/HHTSolver.md
  \item \textbf{Archivo}: solidum/math/solvers/newmark.py (clase \texttt{HHTSolver}).
\end{itemize}


\section{NewtonHHTSolver --- análisis dinámico transitorio no lineal HHT-\ensuremath{\alpha}}

\begin{itemize}
  \item \textbf{Propósito}: combinar HHT-\ensuremath{\alpha} (disipación numérica) con Newton-Raphson interno para problemas con materiales no lineales (plasticidad, daño) o no linealidad geométrica. \textbf{Variante} de \texttt{NewtonNewmarkSolver} (Reglas §4) que aplica el esquema temporal HHT-\ensuremath{\alpha} al residuo dinámico no lineal.
  \item \textbf{Esquema}:
  \item Residuo dinámico: \texttt{R(ü\_\{n+1\}) = (1+\ensuremath{\alpha})\ensuremath{\cdot}F\_\{n+1\} − \ensuremath{\alpha}\ensuremath{\cdot}F\_n − [(1+\ensuremath{\alpha})\ensuremath{\cdot}F\_int(u\_\{n+1\}) − \ensuremath{\alpha}\ensuremath{\cdot}F\_int(u\_n)] − [(1+\ensuremath{\alpha})\ensuremath{\cdot}C\ensuremath{\cdot}u̇\_\{n+1\} − \ensuremath{\alpha}\ensuremath{\cdot}C\ensuremath{\cdot}u̇\_n] − M\ensuremath{\cdot}ü\_\{n+1\}}.
  \item Jacobiano: \texttt{J = M + (1+\ensuremath{\alpha})\ensuremath{\cdot}\ensuremath{\gamma}\ensuremath{\Delta}t\ensuremath{\cdot}C + (1+\ensuremath{\alpha})\ensuremath{\cdot}\ensuremath{\beta}\ensuremath{\Delta}t\ensuremath{^{2}}\ensuremath{\cdot}K\_t}.
  \item Todo lo demás idéntico a \texttt{NewtonNewmarkSolver}: convergencia dual (ADR 0007), commit/rollback, Rayleigh con K\_0, Newton modificado opcional, \textbf{\texttt{line\_search}} opt-in (ADR 0011), telemetría tipada al diverger.
  \item \textbf{Recuperación}: con materiales lineales, el residuo de Newton se anula en una iteración y el resultado coincide con \texttt{HHTSolver}. Validado en tests.
  \item \textbf{Parámetros}: \texttt{alpha} (default \texttt{-0.05}) + heredados de \texttt{NewtonNewmarkSolver} (\texttt{convergence}, \texttt{max\_iter}, \texttt{freeze\_tangent\_after\_iter}, \texttt{line\_search}, \ldots{}).
  \item \textbf{Cuándo usarlo}: respuesta dinámica de estructuras con plasticidad transitoria + presencia de modos altos espurios que el Newmark trapezoidal no atenuaría.
  \item \textbf{Despacho YAML}: \texttt{solver.type: NewtonHHTSolver}. Vía atributo \texttt{PIPELINE\_KIND="transient"} heredado de \texttt{NewmarkSolver} \ensuremath{\to} \texttt{run\_transient}.
  \item \textbf{Spec}: docs/specs/NewtonHHTSolver.md (variante corta sobre la spec padre \texttt{HHTSolver.md}).
  \item \textbf{Archivo}: solidum/math/solvers/newmark.py (clase \texttt{NewtonHHTSolver}).
\end{itemize}


\section{CentralDifferenceSolver --- integración explícita por diferencias centradas (ADR 0009 fase 5)}

\begin{itemize}
  \item \textbf{Propósito}: integrar \texttt{M\ensuremath{\cdot}ü + C\ensuremath{\cdot}u̇ + F\_int(u) = F(t)} con esquema \textbf{explícito} de segundo orden --- la aceleración se actualiza en cada paso por \texttt{M\ensuremath{^{-}}¹\ensuremath{\cdot}rhs} con \texttt{M} diagonal lumped, sin sistema lineal ni Newton interno. Cubre análisis lineal y no lineal con un solo solver (parámetro \texttt{nonlinear: bool}).
  \item \textbf{Esquema (leapfrog Belytschko-Liu-Moran)}:
  \item Inicialización: \texttt{a\_0 = M\ensuremath{^{-}}¹\ensuremath{\cdot}(F(0) − F\_int(u\_0) − C\ensuremath{\cdot}v\_0 − F\_dir)}, \texttt{v\_\{1/2\} = v\_0 + (\ensuremath{\Delta}t/2)\ensuremath{\cdot}a\_0}.
  \item Paso \texttt{n \ensuremath{\to} n+1}: \texttt{u\_\{n+1\} = u\_n + \ensuremath{\Delta}t\ensuremath{\cdot}v\_\{n+1/2\}}; reevaluar \texttt{F\_int(u\_\{n+1\})}; \texttt{a\_\{n+1\} = M\ensuremath{^{-}}¹\ensuremath{\cdot}(F\_\{n+1\} − F\_int − C\ensuremath{\cdot}v\_\{n+1/2\} − F\_dir)}; \texttt{v\_\{n+1\} = v\_\{n+1/2\} + (\ensuremath{\Delta}t/2)\ensuremath{\cdot}a\_\{n+1\}}.
  \item Modo lineal: \texttt{F\_int = K\ensuremath{\cdot}u} (K constante, ensamblada una vez). Modo no lineal: \texttt{Assembler.assemble\_non\_linear\_system(u)} cada paso (descarta la tangente).
  \item \textbf{Estabilidad condicional CFL}: \texttt{\ensuremath{\Delta}t < 2/\ensuremath{\omega}\_max}. Para barras: \texttt{\ensuremath{\Delta}t < L\_el/c\_p = L\_el\ensuremath{\cdot}\ensuremath{\surd}(\ensuremath{\rho}/E)}. Si se excede, el solver detecta la explosión exponencial a posteriori y aborta con \texttt{RuntimeError} informativo. \textbf{El solver no calcula \texttt{\ensuremath{\Delta}t\_crit} automáticamente} --- responsabilidad del usuario en esta versión inicial (power iteration sobre \texttt{M\ensuremath{^{-}}¹K} queda diferida).
  \item \textbf{Parámetros}: \texttt{t\_end}, \texttt{dt}, \texttt{nonlinear} (default \texttt{False}), \texttt{rayleigh}, \texttt{u0}, \texttt{u0\_dot}, \texttt{F\_func}, \texttt{lumping} (default \texttt{"lumped"}; \texttt{"consistent"} rechazado), \texttt{divergence\_threshold} (default \texttt{1e8}).
  \item \textbf{Cuándo usarlo}: wave propagation, impacto, dinámica de respuesta rápida --- situaciones donde el \texttt{\ensuremath{\Delta}t} natural del problema cae por debajo del paso de estabilidad incondicional de Newmark. Para problemas con plasticidad transitoria de respuesta moderadamente rápida, sigue siendo competitivo frente a Newton-Newmark por evitar la factorización de la tangente cada iteración.
  \item \textbf{Rechazos defensivos}: \texttt{lumping="consistent"} \ensuremath{\to} \texttt{ValueError} (la inversión de M consistente cada paso anula la ventaja); Frame3D con eje oblicuo a globales \ensuremath{\to} \texttt{ValueError} (M\_red no es estrictamente diagonal; ADR 0009 fase 2 documenta la bloque-diagonalidad inherente).
  \item \textbf{Despacho YAML}: \texttt{solver.type: CentralDifferenceSolver}. Vía atributo \texttt{PIPELINE\_KIND="transient"} \ensuremath{\to} \texttt{run\_transient}. \textbf{Primer solver fuera de la jerarquía de Newmark}, lo que disparó la regla C de refactor (atributo declarativo en lugar de cadena de isinstance).
  \item \textbf{Spec}: docs/specs/CentralDifferenceSolver.md.
  \item \textbf{Archivo}: solidum/math/solvers/central\_difference.py.
\end{itemize}


\section{HarmonicSolver --- respuesta forzada armónica en el dominio de la frecuencia (ADR 0009 fase 6)}

\begin{itemize}
  \item \textbf{Propósito}: resolver \texttt{(−\ensuremath{\omega}\ensuremath{^{2}}M + i\ensuremath{\omega}C + K)\ensuremath{\cdot}û = F̂} para un barrido de frecuencias y devolver la amplitud compleja \texttt{û(\ensuremath{\omega})} en cada frecuencia. Lineal, en estado estacionario --- sin transitorio. Apropiado para FRFs, diagnóstico de resonancias y respuesta forzada periódica.
  \item \textbf{Esquema}: ensamble único de \texttt{K}, \texttt{M}, \texttt{C = \ensuremath{\alpha}\ensuremath{\cdot}M + \ensuremath{\beta}\ensuremath{\cdot}K}; reducción de Dirichlet; barrido \texttt{for \ensuremath{\omega} in omega: spsolve(Z(\ensuremath{\omega}), F\_red)} con \texttt{Z(\ensuremath{\omega}) = -\ensuremath{\omega}\ensuremath{^{2}}M + i\ensuremath{\omega}C + K} compleja. Factorización LU compleja por frecuencia (no se cachea --- \texttt{Z} cambia con \texttt{\ensuremath{\omega}}).
  \item \textbf{Barrido configurable}: lineal (\texttt{scale="linear"}), logarítmico (\texttt{scale="log"}), o vector explícito (\texttt{omega=array}). Default \texttt{n\_omega=100}.
  \item \textbf{Parámetros}: \texttt{omega} o (\texttt{omega\_min}, \texttt{omega\_max}, \texttt{n\_omega}, \texttt{scale}); \texttt{F\_amplitude} (default ceros; YAML inyecta automáticamente las \texttt{point\_loads}); \texttt{rayleigh}; \texttt{lumping}.
  \item \textbf{Cuándo usarlo}: diagnóstico de resonancias, FRFs, respuesta forzada armónica de máquinas rotantes, viento turbulento descompuesto en armónicos. Si el problema requiere no linealidad (geométrica o material), usar \texttt{NewtonNewmarkSolver} con excitación armónica.
  \item \textbf{Caveats}: resonancia exacta sin amortiguamiento \ensuremath{\Rightarrow} \texttt{Z(\ensuremath{\omega}\_n)} singular (\texttt{spsolve} puede fallar). Mitigación: añadir Rayleigh. Cargas con fase compleja requieren construcción desde código --- YAML no soporta tipos complejos nativos.
  \item \textbf{Despacho YAML}: \texttt{solver.type: HarmonicSolver}. Vía atributo \texttt{PIPELINE\_KIND="harmonic"} \ensuremath{\to} \texttt{run\_harmonic}. Tercer valor del literal (después de \texttt{"static"}, \texttt{"modal"}, \texttt{"transient"}).
  \item \textbf{Spec}: docs/specs/HarmonicSolver.md.
  \item \textbf{Archivo}: solidum/math/solvers/harmonic.py.
\end{itemize}


\section{ResponseSpectrumSolver --- análisis sísmico por combinación modal (ADR 0009 fase 7)}

\begin{itemize}
  \item \textbf{Propósito}: combinar las respuestas máximas de los primeros \texttt{n\_modes} modos del sistema bajo un espectro de respuesta dado (\texttt{S\_d(\ensuremath{\omega})} o \texttt{S\_a(\ensuremath{\omega})}) en una envolvente máxima. Análisis estadístico, no temporal --- apropiado para diseño sísmico normativo (ASCE 7, Eurocódigo 8, NCh 433, NTC-DS).
  \item \textbf{Esquema}:
  \item Análisis modal interno vía \texttt{ModalSolver} para obtener \texttt{(modes, frequencies)}.
  \item Factores de participación \texttt{Γ\_n = \ensuremath{\varphi}\_n\ensuremath{^{\top}}\ensuremath{\cdot}M\ensuremath{\cdot}r} con \texttt{r} el vector de excitación rígida en la dirección sísmica.
  \item Respuesta máxima por modo: \texttt{u\_n\textasciicircum{}max = Γ\_n \ensuremath{\cdot} \ensuremath{\varphi}\_n \ensuremath{\cdot} S\_d(\ensuremath{\omega}\_n)}.
  \item Combinación SRSS (default) o CQC con coeficientes de correlación de Der Kiureghian.
  \item \textbf{Parámetros}: \texttt{n\_modes}, \texttt{direction} (ndarray o \texttt{\{"dof\_name": "ux"\}}), \texttt{spectrum} (callable o dict tabulado/constante), \texttt{combination} (\texttt{"SRSS"}/\texttt{"CQC"}), \texttt{damping} (sólo CQC, default 0.05), \texttt{sigma}, \texttt{lumping}.
  \item \textbf{Cuándo usarlo}: diseño sísmico contra espectro normativo; diagnóstico de modos dominantes en una dirección de excitación. Para no-linealidad usar transitorio (Newton-Newmark con acelerograma).
  \item \textbf{Caveats}: precisión depende del número de modos incluidos --- verifica \texttt{cumulative\_effective\_mass\_ratio} \ensuremath{\ge} 0.9. CQC requerido con modos cercanos en frecuencia (cociente < 1.3). Excitación multi-direccional simultánea (CQC3, 100/30/30) requiere combinación adicional fuera del solver.
  \item \textbf{Despacho YAML}: \texttt{solver.type: ResponseSpectrumSolver}. Vía atributo \texttt{PIPELINE\_KIND="spectrum"} \ensuremath{\to} \texttt{run\_response\_spectrum}. Cuarto valor del literal (después de \texttt{"static"}, \texttt{"modal"}, \texttt{"transient"}, \texttt{"harmonic"}).
  \item \textbf{Algoritmos centralizados}: \texttt{solidum/math/modal\_response.py} --- \texttt{response\_spectrum\_srss}, \texttt{response\_spectrum\_cqc}, \texttt{participation\_factors}, \texttt{spectrum\_from\_sa}, \texttt{spectrum\_tabulated}. Compartido con \texttt{ModalResult.free\_vibration} (wrapper delgado tras regla D 2026-05-18).
  \item \textbf{Spec}: docs/specs/ResponseSpectrumSolver.md.
  \item \textbf{Archivo}: solidum/math/solvers/response\_spectrum.py.
\end{itemize}


\section{Cómo añadir un solver nuevo}

\texttt{/solidum-new solver <Name>} --- genera archivo en \texttt{solidum/math/solvers/<snake>.py}, decorador \texttt{@SolverRegistry.register}, esqueleto de test.

Convenciones de interfaz:

\begin{itemize}
  \item \textbf{Solvers estáticos} (lineales, no lineales, arc-length): \texttt{PIPELINE\_KIND = "static"}. Constructor recibe \texttt{assembler} + parámetros; método \texttt{solve(F\_ext\_global, step\_callback=None) \ensuremath{\to} U\_final}. Comprometen los estados internos vía \texttt{assembler.commit\_all\_states()} al converger cada paso. Retornan el campo de desplazamientos completo.
  \item \textbf{Solvers modales / autovalor} (modal --- ADR 0009 --- y futuros pandeo lineal): \texttt{PIPELINE\_KIND = "modal"}. Constructor recibe \texttt{assembler} + parámetros; método \texttt{solve() \ensuremath{\to} ModalResult} (u otro tipo específico). No consumen vector de cargas. \texttt{run\_yaml} despacha a \texttt{run\_modal}.
  \item \textbf{Solvers transitorios} (Newmark, HHT, central differences): \texttt{PIPELINE\_KIND = "transient"}. Constructor recibe \texttt{assembler}, \texttt{t\_end}, \texttt{dt} + parámetros; método \texttt{solve() \ensuremath{\to} TransientResult}. \texttt{run\_yaml} despacha a \texttt{run\_transient}.
  \item \textbf{Solvers en frecuencia / espectrales} (harmonic con \texttt{PIPELINE\_KIND = "harmonic"}, response spectrum con \texttt{"spectrum"}). Constructor recibe \texttt{assembler} + parámetros del análisis; método \texttt{solve()} devuelve el resultado específico (\texttt{HarmonicResult}, \texttt{ResponseSpectrumResult}). \texttt{run\_yaml} despacha a \texttt{run\_harmonic} o \texttt{run\_response\_spectrum} según el atributo.
\end{itemize}

El \textbf{dispatch en \texttt{run\_yaml}} se hace por el atributo de clase \texttt{PIPELINE\_KIND} (regla C, 2026-05-18). Solvers nuevos no clásicos (que no hereden de los existentes) sólo declaran su \texttt{PIPELINE\_KIND} y quedan automáticamente cableados --- no requieren tocar \texttt{entry.py}.

Tras implementar, \textbf{añadir una entrada a este catálogo} siguiendo el formato de arriba y promover la spec a \texttt{status: validated}.


\newpage

\chapter{Anexos técnicos}



\textit{Esta sección es referencia técnica del subsistema interno que resuelve el sistema lineal \texttt{K\ensuremath{\cdot}x = b} que aparece en el corazón de cada solver no lineal. La elección del backend es \textbf{automática}; el usuario no decide. Esta sección documenta cómo se decide y cuáles son los puntos de override para diagnóstico.}

\textit{> Diseño completo: ver ADR 0003 en \texttt{docs/adr/0003-despachador-de-solver-algebraico.md}.}


\section{Dos capas que se llaman ambas «solver»}

Conviene distinguirlas:

\begin{itemize}
  \item \textbf{Solver no lineal} --- \texttt{LinearSolver}, \texttt{NonlinearSolver}, \texttt{ArcLengthSolver}. Orquesta la estrategia de paso, las iteraciones de Newton, el control de longitud de arco, los criterios de convergencia. Es lo que el usuario elige en el YAML con \texttt{solver.type}.
  \item \textbf{Capa algebraica} --- \texttt{solidum.math.linalg}. Resuelve el sistema lineal \texttt{K\ensuremath{\cdot}\ensuremath{\delta}U = R} que aparece dentro de cada iteración del solver no lineal. Tiene varios \textit{backends} numéricos y un despachador interno que elige el adecuado según las propiedades de \texttt{K}. Es plumbing automático: el usuario no la ve salvo que pida diagnóstico.
\end{itemize}

\section{Backends disponibles y criterio de selección}

El despachador (\texttt{solidum.math.linalg.dispatcher.select\_solver}) elige el backend en función de tres flags declarativos sobre \texttt{K}: simetría, positividad y talla.

\begin{center}
\begin{tabularx}{\textwidth}{YYY}
\hline
\textbf{Régimen de \texttt{K}} & \textbf{Backend elegido} & \textbf{Notas} \\
\hline
Simétrica + positiva definida & \textbf{Cholesky} (CHOLMOD) & \textasciitilde{}2\ensuremath{\times} más rápido y \textasciitilde{}2\ensuremath{\times} menos memoria que LU. Requiere la dependencia opcional \texttt{scikit-sparse}. Si falta, degrada a LU con un \textit{warning} una sola vez. \\
Simétrica indefinida & LU general (placeholder LDL\ensuremath{^{\top}}) & Ideal para LDL\ensuremath{^{\top}} con conteo de pivotes negativos (Sturm sequence) en pandeo y snap-through. La interfaz está reservada; mientras no haya backend nativo, se usa LU sin pérdida de corrección, solo sin diagnóstico de bifurcación. \\
No simétrica & \textbf{LU general} (SuperLU) & Backend universal. Cubre plasticidad no asociada, follower loads, contacto con fricción. \\
\hline
\end{tabularx}
\end{center}
\textbf{Origen de los flags}, todos derivables del modelo:
\begin{itemize}
  \item \texttt{is\_symmetric}: AND lógico de \texttt{material.IS\_SYMMETRIC} y \texttt{element.PRESERVES\_SYMMETRY} sobre todos los componentes del dominio. Defaults \texttt{True}.
  \item \texttt{is\_positive\_definite}: arranca en \texttt{True} para \texttt{LinearSolver} y \texttt{NonlinearSolver}; en \texttt{False} para \texttt{ArcLengthSolver}. Se degrada automáticamente si Cholesky reporta no-positividad.
  \item \texttt{size}: número total de DOFs.
\end{itemize}

Ningún flag se pide al usuario.

\textbf{\texttt{EigenSolver} (ADR 0009, problema generalizado)}. La capa algebraica incluye además \texttt{solidum.math.linalg.eigen.EigenSolver}, que resuelve el problema generalizado simétrico \texttt{K \ensuremath{\cdot} \ensuremath{\varphi} = \ensuremath{\omega}\ensuremath{^{2}} \ensuremath{\cdot} M \ensuremath{\cdot} \ensuremath{\varphi}} envolviendo \texttt{scipy.sparse.linalg.eigsh} (ARPACK Lanczos con shift-invert centrado en \texttt{\ensuremath{\sigma}}). No comparte el \texttt{Protocol} \texttt{LinearAlgebraSolver} de los backends de \texttt{K\ensuremath{\cdot}x = b} porque su firma natural es \texttt{solve(K, M, n\_modes) \ensuremath{\to} (\ensuremath{\lambda}, \ensuremath{\varphi})}. El \texttt{ModalSolver} lo invoca para el análisis modal; los cálculos internos de shift-invert generan una factorización de \texttt{(K − \ensuremath{\sigma} \ensuremath{\cdot} M)} reutilizada en cada iteración Lanczos, así que la palanca de optimización es la misma --- Cholesky enchufado en lugar de SuperLU bajaría 2\ensuremath{\times} el coste del modal en problemas SPD (pendiente, ver memoria de cierre).

\section{Fallback automático SPD \ensuremath{\to} LU}

Si en algún paso \texttt{K} deja de ser positiva definida (paso a régimen postcrítico, daño con reblandecimiento), Cholesky lanza \texttt{CholeskyNotPositiveDefiniteError}. Los tres solvers no lineales lo capturan y reinstancian el backend con LU general para el resto del análisis. La transición es silenciosa --- solo se imprime un mensaje informativo en stdout.

Esto significa que \textbf{forzar Cholesky desde YAML como diagnóstico nunca rompe el análisis}: si la matriz se vuelve incompatible, el sistema se autocorrige.

\section{Factorización reusable y Newton modificado}

La interfaz expone \texttt{solver.factorize(K)} que retorna un objeto \texttt{FactorizedSolver} con \texttt{solve(b)} reutilizable. Habilita:

\begin{itemize}
  \item \textbf{Newton modificado}: el \texttt{NonlinearSolver} admite el parámetro \texttt{freeze\_tangent\_after\_iter: int} que congela la factorización tras N iteraciones del paso. Útil cuando la tangente cambia poco y el coste dominante es la factorización.
  \item \textbf{Dinámica implícita lineal} (ADR 0009 fase 3): el \texttt{NewmarkSolver} factoriza una sola vez la matriz efectiva \texttt{A\_eff = M + \ensuremath{\gamma}\ensuremath{\Delta}t \ensuremath{\cdot} C + \ensuremath{\beta}\ensuremath{\Delta}t\ensuremath{^{2}} \ensuremath{\cdot} K} al inicio del análisis y reutiliza la factorización en cada paso temporal. Con \texttt{\ensuremath{\Delta}t} constante y problema lineal, los cientos o miles de pasos del análisis se reducen a sustituciones triangulares baratas. Es el mismo \texttt{FactorizedSolver} del despachador.
  \item \textbf{Análisis futuros} (pandeo lineal, dinámica implícita no lineal con factorización congelada entre pasos cuando la tangente cambia poco) reutilizan la misma interfaz sin extensiones adicionales.
\end{itemize}

Ejemplo YAML:

\begin{lstlisting}[language=yaml]
solver:
  type: NonlinearSolver
  num_steps: 10
  tol: 1.0e-8
  freeze_tangent_after_iter: 2

\end{lstlisting}

\section{Override desde YAML --- herramienta de diagnóstico}

Solo en caso de necesidad de \textit{diagnóstico} o \textit{benchmarking}, el bloque \texttt{solver} admite el campo opcional \texttt{linear\_algebra}:

\begin{lstlisting}[language=yaml]
solver:
  type: LinearSolver
  linear_algebra: auto         # default; el despachador decide
  # linear_algebra: cholesky   # forzar Cholesky (requiere scikit-sparse)
  # linear_algebra: lu         # forzar LU (siempre disponible)
  # linear_algebra: ldlt       # placeholder; degrada a LU con warning

\end{lstlisting}

\textbf{Cuándo usar el override:}

\begin{itemize}
  \item Comparar tiempos entre Cholesky y LU en un mismo problema.
  \item Forzar LU si se sospecha un bug en la auto-detección de simetría tras añadir un material nuevo.
  \item Reproducir un caso de regresión bit-a-bit con un backend específico.
\end{itemize}

\textbf{Cuándo NO usarlo}: como parte del setup habitual de un caso de análisis. No es decisión de modelado; el default \texttt{auto} cubre todos los regímenes correctamente.

\section{Activación de Cholesky (opcional)}

\begin{lstlisting}[language=bash]
conda install -c conda-forge scikit-sparse

\end{lstlisting}

Una vez instalado, los problemas estáticos lineales y los Newton estables empiezan a usar Cholesky automáticamente --- sin tocar YAML, sin recompilar, sin reescribir specs. Si la dependencia no está disponible, Solidum funciona idéntico con LU.


\newpage


\end{document}
